% Auto-generated from data/segments.json + layout.json (+ transforms) — do not edit by hand.
% volume: 08Di03
% mode: reading
% segments: 1615
% transform_rules: 0
% toc_marks: 300

% Volume layout defaults (layout.json ``layout``).
\csromanlayoutapply{1}{1.5}{21.6pt}{5pt}{6.3pt}{65pt}{2.5em}%

% Reading physical-page layout (layout.json ``page_layout_reading_mode``).
\csromanreadingpagelayoutvolume{1}{1.5}{21.6pt}{5pt}{6.3pt}{65pt}{2.5em}%
\csromanreadingpagelayoutenable

\csromanheader{2}{\textbf{ทีฆนิกาย}}
\csromanmatikaanchor{mtk.0}
\csromantocmarknopage{chapter}{ปาถิกวคฺคปาฬิ}{mtk.0}
\bookmark[dest={mtk.0},level=0]{ปาถิกวคฺคปาฬิ}
\gambhira{\textbf{ปาถิกวคฺคปาฬิ}}
\namakkaram{นโม ตสฺส ภควโต อรหโต สมฺมาสมฺพุทฺธสฺส.}
\csromanmatikaanchor{mtk.1}
\csromanmatikaanchor{mtk.2}
\csromantocmarknopage{chapter}{1. ปาถิกสุตฺต}{mtk.1}
\bookmark[dest={mtk.1},level=0]{1. ปาถิกสุตฺต}
\csromantocmarkref{section}{สุนกฺขตฺตวตฺถุ}{mtk.2}
\bookmark[dest={mtk.2},level=1]{สุนกฺขตฺตวตฺถุ}
\csromanheader{1}{\textbf{สุนกฺขตฺตวตฺถุ}}
\proseitem{1}{\csromanfolio{1}เอวํ เม สุตํ\csromandash{}เอกํ สมยํ ภควา มลฺเลสุ วิหรติ อนุปิยํ นาม\footnote{อนุปฺปิยํ นาม (สฺยา)} มลฺลานํ นิคโม.\csromanspacer{} อถ โข ภควา ปุพฺพณฺหสมยํ นิวาเสตฺวา ปตฺตจีวรมาทาย อนุปิยํ ปิณฺฑาย ปาวิสิ.\csromanspacer{} อถ โข ภควโต เอตทโหสิ “อติปฺปโค โข ตาว อนุปิยายํ\footnote{อนุปิยํ (ก)} ปิณฺฑาย จริตุํ, ยํนูนาหํ เยน ภคฺควโคตฺตสฺส ปริพฺพาชกสฺส อาราโม, เยน ภคฺควโคตฺโต ปริพฺพาชโก เตนุปสงฺกเมยฺยนฺ”ติ.}

\proseitem{2}{อถ โข ภควา เยน ภคฺควโคตฺตสฺส ปริพฺพาชกสฺส อาราโม, เยน ภคฺควโคตฺโต ปริพฺพาชโก เตนุปสงฺกมิ.\csromanspacer{} อถ โข ภคฺควโคตฺโต ปริพฺพาชโก ภควนฺตํ เอตทโวจ “เอตุ โข ภนฺเต ภควา, สฺวาคตํ ภนฺเต ภควโต, จิรสฺสํ โข ภนฺเต ภควา อิมํ ปริยายมกาสิ ยทิทํ อิธาคมนาย, นิสีทตุ ภนฺเต ภควา, อิทมาสนํ ปญฺญตฺตนฺ”ติ.\csromanspacer{} นิสีทิ ภควา ปญฺญตฺเต อาสเน.\csromanspacer{} ภคฺควโคตฺโตปิ โข ปริพฺพาชโก อญฺญตรํ นีจํ อาสนํ คเหตฺวา เอกมนฺตํ นิสีทิ.\csromanspacer{} เอกมนฺตํ นิสินฺโน โข ภคฺควโคตฺโต ปริพฺพาชโก ภควนฺตํ เอตทโวจ\csromandash{}ปุริมานิ ภนฺเต \csromanfolio{2}ทิวสานิ ปุริมตรานิ สุนกฺขตฺโต ลิจฺฉวิปุตฺโต เยนาหํ เตนุปสงฺกมิ, อุปสงฺกมิตฺวา มํ เอตทโวจ “ปจฺจกฺขาโต ทานิ มยา ภคฺคว ภควา, น ทานาหํ ภควนฺตํ อุทฺทิสฺส วิหรามี”ติ.\csromanspacer{} กจฺเจตํ ภนฺเต ตเถว, ยถา สุนกฺขตฺโต ลิจฺฉวิปุตฺโต อวจาติ.\csromanspacer{} ตเถว โข เอตํ ภคฺคว, ยถา สุนกฺขตฺโต ลิจฺฉวิปุตฺโต อวจ.}

\proseitem{3}{ปุริมานิ ภคฺคว ทิวสานิ ปุริมตรานิ สุนกฺขตฺโต ลิจฺฉวิปุตฺโต เยนาหํ เตนุปสงฺกมิ, อุปสงฺกมิตฺวา มํ อภิวาเทตฺวา เอกมนฺตํ นิสีทิ.\csromanspacer{} เอกมนฺตํ นิสินฺโน โข ภคฺคว สุนกฺขตฺโต ลิจฺฉวิปุตฺโต มํ เอตทโวจ “ปจฺจกฺขามิ ทานาหํ ภนฺเต ภควนฺตํ, น ทานาหํ ภนฺเต ภควนฺตํ อุทฺทิสฺส วิหริสฺสามี”ติ.\csromanspacer{} เอวํ วุตฺเต อหํ ภคฺคว สุนกฺขตฺตํ ลิจฺฉวิปุตฺตํ เอตทโวจํ “อปิ นุ ตาหํ สุนกฺขตฺต เอวํ อวจํ, เอหิ ตฺวํ สุนกฺขตฺต มมํ อุทฺทิสฺส วิหราหี”ติ.\csromanspacer{} โน เหตํ ภนฺเต.\csromanspacer{} ตฺวํ วา ปน มํ เอวํ อวจ “อหํ ภนฺเต ภควนฺตํ อุทฺทิสฺส วิหริสฺสามี”ติ.\csromanspacer{} โน เหตํ ภนฺเต.\csromanspacer{} อิติ กิร สุนกฺขตฺต เนวาหํ ตํ วทามิ “เอหิ ตฺวํ สุนกฺขตฺต มมํ อุทฺทิสฺส วิหราหี”ติ.\csromanspacer{} นปิ กิร มํ ตฺวํ วเทสิ “อหํ ภนฺเต ภควนฺตํ อุทฺทิสฺส วิหริสฺสามี”ติ.\csromanspacer{} เอวํ สนฺเต โมฆปุริส โก สนฺโต กํ ปจฺจาจิกฺขสิ.\csromanspacer{} ปสฺส โมฆปุริส ยาวจ\footnote{ยาวญฺจ (สี, สฺยา, อิ)} เต อิทํ อปรทฺธนฺติ.}

\proseitem{4}{น หิ ปน เม ภนฺเต ภควา อุตฺตริ มนุสฺสธมฺมา อิทฺธิปาฏิหาริยํ กโรตีติ.\csromanspacer{} อปิ นุ ตาหํ สุนกฺขตฺต เอวํ อวจํ “เอหิ ตฺวํ สุนกฺขตฺต มมํ อุทฺทิสฺส วิหราหิ, อหํ เต อุตฺตริ มนุสฺสธมฺมา อิทฺธิปาฏิหาริยํ กริสฺสามี”ติ.\csromanspacer{} โน เหตํ ภนฺเต.\csromanspacer{} ตฺวญฺจ ปน มํ เอวํ อวจ “อหํ ภนฺเต ภควนฺตํ อุทฺทิสฺส วิหริสฺสามิ, ภควา เม อุตฺตริ มนุสฺสธมฺมา อิทฺธิปาฏิหาริยํ กริสฺสตี”ติ.\csromanspacer{} โน เหตํ ภนฺเต.\csromanspacer{} อิติ กิร สุนกฺขตฺต เนวาหํ ตํ วทามิ “เอหิ ตฺวํ สุนกฺขตฺต มมํ อุทฺทิสฺส วิหราหิ, อหํ เต อุตฺตริ มนุสฺสธมฺมา อิทฺธิปาฏิหาริยํ กริสฺสามี”ติ.\csromanspacer{} นปิ กิร มํ ตฺวํ วเทสิ “อหํ ภนฺเต ภควนฺตํ อุทฺทิสฺส วิหริสฺสามิ, ภควา เม อุตฺตริ มนุสฺสธมฺมา อิทฺธิปาฏิหาริยํ กริสฺสตี”ติ.\csromanspacer{} เอวํ สนฺเต โมฆปุริส โก สนฺโต กํ ปจฺจาจิกฺขสิ.\csromanspacer{} ตํ กิํมญฺญสิ สุนกฺขตฺต, กเต วา อุตฺตริ มนุสฺสธมฺมา อิทฺธิปาฏิหาริเย อกเต วา \csromanfolio{3}อุตฺตริ มนุสฺสธมฺมา อิทฺธิปาฏิหาริเย ยสฺสตฺถาย มยา ธมฺโม เทสิโต, โส นิยฺยาติ ตกฺกรสฺส สมฺมา ทุกฺขกฺขยายาติ.\csromanspacer{} กเต วา ภนฺเต อุตฺตริ มนุสฺสธมฺมา อิทฺธิปาฏิหาริเย อกเต วา อุตฺตริ มนุสฺสธมฺมา อิทฺธิปาฏิหาริเย ยสฺสตฺถาย ภควตา ธมฺโม เทสิโต, โส นิยฺยาติ ตกฺกรสฺส สมฺมา ทุกฺขกฺขยายาติ.\csromanspacer{} อิติ กิร สุนกฺขตฺต กเต วา อุตฺตริ มนุสฺสธมฺมา อิทฺธิปาฏิหาริเย อกเต วา อุตฺตริ มนุสฺสธมฺมา อิทฺธิปาฏิหาริเย ยสฺสตฺถาย มยา ธมฺโม เทสิโต, โส นิยฺยาติ ตกฺกรสฺส สมฺมา ทุกฺขกฺขยาย.\csromanspacer{} ตตฺร สุนกฺขตฺต กิํ อุตฺตริ มนุสฺสธมฺมา อิทฺธิปาฏิหาริยํ กตํ กริสฺสติ.\csromanspacer{} ปสฺส โมฆปุริส ยาวจ เต อิทํ อปรทฺธนฺติ.}

\proseitem{5}{น หิ ปน เม ภนฺเต ภควา อคฺคญฺญํ ปญฺญเปตีติ\footnote{ปญฺญาเปตีติ (อิ)}.\csromanspacer{} อปิ นุ ตาหํ สุนกฺขตฺต เอวํ อวจํ “เอหิ ตฺวํ สุนกฺขตฺต มมํ อุทฺทิสฺส วิหราหิ, อหํ เต อคฺคญฺญํ ปญฺญเปสฺสามี”ติ.\csromanspacer{} โน เหตํ ภนฺเต.\csromanspacer{} ตฺวํ วา ปน มํ เอวํ อวจ “อหํ ภนฺเต ภควนฺตํ อุทฺทิสฺส วิหริสฺสามิ, ภควา เม อคฺคญฺญํ ปญฺญเปสฺสตี”ติ.\csromanspacer{} โน เหตํ ภนฺเต.\csromanspacer{} อิติ กิร สุนกฺขตฺต เนวาหํ ตํ วทามิ “เอหิ ตฺวํ สุนกฺขตฺต มมํ อุทฺทิสฺส วิหราหิ, อหํ เต อคฺคญฺญํ ปญฺญเปสฺสามี”ติ.\csromanspacer{} นปิ กิร มํ ตฺวํ วเทสิ “อหํ ภนฺเต ภควนฺตํ อุทฺทิสฺส วิหริสฺสามิ, ภควา เม อคฺคญฺญํ ปญฺญเปสฺสตี”ติ.\csromanspacer{} เอวํ สนฺเต โมฆปุริส โก สนฺโต กํ ปจฺจาจิกฺขสิ.\csromanspacer{} ตํ กิํมญฺญสิ สุนกฺขตฺต, ปญฺญตฺเต วา อคฺคญฺเญ อปญฺญตฺเต วา อคฺคญฺเญ ยสฺสตฺถาย มยา ธมฺโม เทสิโต, โส นิยฺยาติ ตกฺกรสฺส สมฺมา ทุกฺขกฺขยายาติ.\csromanspacer{} ปญฺญตฺเต วา ภนฺเต อคฺคญฺเญ อปญฺญตฺเต วา อคฺคญฺเญ ยสฺสตฺถาย ภควตา ธมฺโม เทสิโต, โส นิยฺยาติ ตกฺกรสฺส สมฺมา ทุกฺขกฺขยายาติ.\csromanspacer{} อิติ กิร สุนกฺขตฺต ปญฺญตฺเต วา อคฺคญฺเญ อปญฺญตฺเต วา อคฺคญฺเญ ยสฺสตฺถาย มยา ธมฺโม เทสิโต, โส นิยฺยาติ ตกฺกรสฺส สมฺมา ทุกฺขกฺขยาย.\csromanspacer{} ตตฺร สุนกฺขตฺต กิํ อคฺคญฺญํ ปญฺญตฺตํ กริสฺสติ.\csromanspacer{} ปสฺส โมฆปุริส ยาว จ เต อิทํ อปรทฺธํ.}

\proseitem{6}{อเนกปริยาเยน โข เต สุนกฺขตฺต มม วณฺโณ ภาสิโต วชฺชิคาเม “อิติปิ โส ภควา อรหํ สมฺมาสมฺพุทฺโธ วิชฺชาจรณสมฺปนฺโน \csromanfolio{4}สุคโต โลกวิทู อนุตฺตโร ปุริสทมฺมสารถิ สตฺถา เทวมนุสฺสานํ พุทฺโธ ภควา”ติ.\csromanspacer{} อิติ โข เต สุนกฺขตฺต อเนกปริยาเยน มม วณฺโณ ภาสิโต วชฺชิคาเม.}

\prose{อเนกปริยาเยน โข เต สุนกฺขตฺต ธมฺมสฺส วณฺโณ ภาสิโต วชฺชิคาเม “สฺวากฺขาโต ภควตา ธมฺโม สนฺทิฏฺฐิโก อกาลิโก เอหิปสฺสิโก โอปเนยฺยิโก ปจฺจตฺตํ เวทิตพฺโพ วิญฺญูหี”ติ.\csromanspacer{} อิติ โข เต สุนกฺขตฺต อเนกปริยาเยน ธมฺมสฺส วณฺโณ ภาสิโต วชฺชิคาเม.}

\prose{อเนกปริยาเยน โข เต สุนกฺขตฺต สํฆสฺส วณฺโณ ภาสิโต วชฺชิคาเม “สุปฺปฏิปนฺโน ภควโต สาวกสํโฆ, อุชุปฺปฏิปนฺโน ภควโต สาวกสํโฆ, ญายปฺปฏิปนฺโน ภควโต สาวกสํโฆ, สามีจิปฺปฏิปนฺโน ภควโต สาวกสํโฆ, ยทิทํ จตฺตาริ ปุริสยุคานิ อฏฺฐ ปุริสปุคฺคลา, เอส ภควโต สาวกสํโฆ อาหุเนยฺโย ปาหุเนยฺโย ทกฺขิเณยฺโย อญฺชลิกรณีโย อนุตฺตรํ ปุญฺญกฺเขตฺตํ โลกสฺสา”ติ.\csromanspacer{} อิติ โข เต สุนกฺขตฺต อเนกปริยาเยน สํฆสฺส วณฺโณ ภาสิโต วชฺชิคาเม.}

\prose{อาโรจยามิ โข เต สุนกฺขตฺต, ปฏิเวทยามิ โข เต สุนกฺขตฺต, ภวิสฺสนฺติ โข เต สุนกฺขตฺต วตฺตาโร, “โน วิสหิ สุนกฺขตฺโต ลิจฺฉวิปุตฺโต สมเณ โคตเม พฺรหฺมจริยํ จริตุํ, โส อวิสหนฺโต สิกฺขํ ปจฺจกฺขาย หีนายาวตฺโต”ติ.\csromanspacer{} อิติ โข เต สุนกฺขตฺต ภวิสฺสนฺติ วตฺตาโรติ.\csromanspacer{} เอวํ โข ภคฺคว สุนกฺขตฺโต ลิจฺฉวิปุตฺโต มยา วุจฺจมาโน อปกฺกเมว อิมสฺมา ธมฺมวินยา, ยถา ตํ อาปายิโก เนรยิโก.}

\csromanmatikaanchor{mtk.3}
\csromantocmarkref{section}{โกรกฺขตฺติยวตฺถุ}{mtk.3}
\bookmark[dest={mtk.3},level=1]{โกรกฺขตฺติยวตฺถุ}
\csromanheader{1}{\textbf{โกรกฺขตฺติยวตฺถุ}}
\proseitem{7}{เอกมิทาหํ ภคฺคว สมยํ ถูลูสุ\footnote{พุมูสุ (สี, อิ)} วิหรามิ อุตฺตรกา นาม ถูลูนํ นิคโม.\csromanspacer{} อถ ขฺวาหํ ภคฺคว ปุพฺพณฺหสมยํ นิวาเสตฺวา ปตฺตจีวรมาทาย สุนกฺขตฺเตน ลิจฺฉวิปุตฺเตน ปจฺฉาสมเณน อุตฺตรกํ ปิณฺฑาย ปาวิสิํ.\csromanspacer{} เตน โข ปน สมเยน อเจโล โกรกฺขตฺติโย กุกฺกุรวติโก จตุโกณฺฑิโก\footnote{จตุกุณฺฑิโก(สี, อิ)} ฉมานิกิณฺณํ ภกฺขสํ มุเขเนว ขาทติ, มุเขเนว ภุญฺชติ, อทฺทสา โข ภคฺคว สุนกฺขตฺโต ลิจฺฉวิปุตฺโต อเจลํ โกรกฺขตฺติยํ \csromanfolio{5}กุกฺกุรวติกํ จตุโกณฺฑิกํ ฉมานิกิณฺณํ ภกฺขสํ มุเขเนว ขาทนฺตํ มุเขเนว ภุญฺชนฺตํ.\csromanspacer{} ทิสฺวานสฺส เอตทโหสิ “สาธุรูโป วต โภ อยํ\footnote{อรหํ (สี, สฺยา, อิ)} สมโณ จตุโกณฺฑิโก ฉมานิกิณฺณํ ภกฺขสํ มุเขเนว ขาทติ, มุเขเนว ภุญฺชตี”ติ.}

\prose{อถ ขฺวาหํ ภคฺคว สุนกฺขตฺตสฺส ลิจฺฉวิปุตฺตสฺส เจตสา เจโตปริวิตกฺกมญฺญาย สุนกฺขตฺตํ ลิจฺฉวิปุตฺตํ เอตทโวจํ “ตฺวํปิ นาม โมฆปุริส สมโณ สกฺยปุตฺติโย\footnote{โมฆปุริส สกฺยปุตฺติโย (สี, สฺยา, อิ)} ปฏิชานิสฺสสี”ติ.\csromanspacer{} กิํ ปน มํ ภนฺเต ภควา เอวมาห “ตฺวํปิ นาม โมฆปุริส สมโณ สกฺยปุตฺติโย2 ปฏิชานิสฺสสี”ติ.\csromanspacer{} นนุ เต สุนกฺขตฺต อิมํ อเจลํ โกรกฺขตฺติยํ กุกฺกุรวติกํ จตุโกณฺฑิกํ ฉมานิกิณฺณํ ภกฺขสํ มุเขเนว ขาทนฺตํ มุเขเนว ภุญฺชนฺตํ ทิสฺวาน เอตทโหสิ “สาธุรูโป วต โภ อยํ สมโณ จตุโกณฺฑิโก ฉมานิกิณฺณํ ภกฺขสํ มุเขเนว ขาทติ, มุเขเนว ภุญฺชตี”ติ.\csromanspacer{} เอวํ ภนฺเต.\csromanspacer{} กิํ ปน ภนฺเต ภควา อรหตฺตสฺส มจฺฉรายตีติ.\csromanspacer{} น โข อหํ โมฆปุริส อรหตฺตสฺส มจฺฉรายามิ.\csromanspacer{} อปิ จ ตุเยฺหเวตํ ปาปกํ ทิฏฺฐิคตํ อุปฺปนฺนํ, ตํ ปชห.\csromanspacer{} มา เต อโหสิ ทีฆรตฺตํ อหิตาย ทุกฺขาย.\csromanspacer{} ยํ โข ปเนตํ สุนกฺขตฺต มญฺญสิ อเจลํ โกรกฺขตฺติยํ “สาธุรูโป อยํ สมโณ”ติ\footnote{มญฺญสิ “อเจโล โกรขตฺติโย สาธุรูโป อรหํ สมโณติ” (สฺยา)}.\csromanspacer{} โส สตฺตมํ ทิวสํ อลสเกน กาลงฺกริสฺสติ, กาลงฺกโต\footnote{กาลกโต (สี, สฺยา, อิ)} จ กาลกญฺจิกา\footnote{กาลกญฺชา (สี, อิ), กาลกญฺชิกา (สฺยา)} นาม อสุรา สพฺพนิหีโน อสุรกาโย, ตตฺร อุปปชฺชิสฺสติ.\csromanspacer{} กาลงฺกตญฺจ นํ พีรณตฺถมฺพเก สุสาเน ฉฑฺเฑสฺสนฺติ.\csromanspacer{} อากงฺขมาโน จ ตฺวํ สุนกฺขตฺต อเจลํ โกรกฺขตฺติยํ อุปสงฺกมิตฺวา ปุจฺเฉยฺยาสิ “ชานาสิ อาวุโส โกรกฺขตฺติย\footnote{อเจล โกรขตฺติย (ก)} อตฺตโน คตินฺ”ติ.\csromanspacer{} ฐานํ โข ปเนตํ สุนกฺขตฺต วิชฺชติ, ยํ เต อเจโล โกรกฺขตฺติโย พฺยากริสฺสติ “ชานามิ อาวุโส สุนกฺขตฺต อตฺตโน คติํ, กาลกญฺจิกา นาม อสุรา สพฺพนิหีโน อสุรกาโย, ตตฺรามฺหิ อุปปนฺโน”ติ.}

\prose{อถ โข ภคฺคว สุนกฺขตฺโต ลิจฺฉวิปุตฺโต เยน อเจโล โกรกฺขตฺติโย เตนุปสงฺกมิ, อุปสงฺกมิตฺวา อเจลํ โกรกฺขตฺติยํ \csromanfolio{6}เอตทโวจ “พฺยากโต โขสิ อาวุโส โกรกฺขตฺติย สมเณน โคตเมน ‘อเจโล โกรกฺขตฺติโย สตฺตมํ ทิวสํ อลสเกน กาลงฺกริสฺสติ, กาลงฺกโต จ กาลกญฺจิกา นาม อสุรา สพฺพนิหีโน อสุรกาโย, ตตฺร อุปปชฺชิสฺสติ.\csromanspacer{} กาลงฺกตญฺจ นํ พีรณตฺถมฺพเก สุสาเน ฉฑฺเฑสฺสนฺตี’ติ.\csromanspacer{} เยน ตฺวํ อาวุโส โกรกฺขตฺติย มตฺตํ มตฺตญฺจ ภตฺตํ ภุญฺเชยฺยาสิ, มตฺตํ มตฺตญฺจ ปานียํ ปิเวยฺยาสิ.\csromanspacer{} ยถา สมณสฺส โคตมสฺส มิจฺฉา อสฺส วจนนฺ”ติ.}

\proseitem{8}{อถ โข ภคฺคว สุนกฺขตฺโต ลิจฺฉวิปุตฺโต เอกทฺวีหิกาย สตฺตรตฺตินฺทิวานิ คเณสิ, ยถา ตํ ตถาคตสฺส อสทฺทหมาโน.\csromanspacer{} อถ โข ภคฺคว อเจโล โกรกฺขตฺติโย สตฺตมํ ทิวสํ อลสเกน กาลมกาสิ, กาลงฺกโต จ กาลกญฺจิกา นาม อสุรา สพฺพนิหีโน อสุรกาโย, ตตฺร อุปปชฺชิ.\csromanspacer{} กาลงฺกตญฺจ นํ พีรณตฺถมฺพเก สุสาเน ฉฑฺเฑสุํ.}

\proseitem{9}{อสฺโสสิ โข ภคฺคว สุนกฺขตฺโต ลิจฺฉวิปุตฺโต “อเจโล กิร โกรกฺขตฺติโย อลสเกน กาลงฺกโต พีรณตฺถมฺพเก สุสาเน ฉฑฺฑิโต”ติ.\csromanspacer{} อถ โข ภคฺคว สุนกฺขตฺโต ลิจฺฉวิปุตฺโต เยน พีรณตฺถมฺพกํ สุสานํ, เยน อเจโล โกรกฺขตฺติโย เตนุปสงฺกมิ, อุปสงฺกมิตฺวา อเจลํ โกรกฺขตฺติยํ ติกฺขตฺตุํ ปาณินา อาโกเฏสิ “ชานาสิ อาวุโส โกรกฺขตฺติย อตฺตโน คตินฺ”ติ.\csromanspacer{} อถ โข ภคฺคว อเจโล โกรกฺขตฺติโย ปาณินา ปิฏฺฐิํ ปริปุญฺฉนฺโต วุฏฺฐาสิ, “ชานามิ อาวุโส สุนกฺขตฺต อตฺตโน คติํ, กาลกญฺจิกา นาม อสุรา สพฺพนิหีโน อสุรกาโย, ตตฺรามฺหิ อุปปนฺโน”ติ วตฺวา ตตฺเถว อุตฺตาโน ปปติ\footnote{ปริปติ (สฺยา, ก)}.}

\proseitem{10}{อถ โข ภคฺคว สุนกฺขตฺโต ลิจฺฉวิปุตฺโต เยนาหํ เตนุปสงฺกมิ, อุปสงฺกมิตฺวา มํ อภิวาเทตฺวา เอกมนฺตํ นิสีทิ.\csromanspacer{} เอกมนฺตํ นิสินฺนํ โข อหํ ภคฺคว สุนกฺขตฺตํ ลิจฺฉวิปุตฺตํ เอตทโวจํ “ตํ กิํมญฺญสิ สุนกฺขตฺต, ยเถว เต อหํ อเจลํ โกรกฺขตฺติยํ อารพฺภ พฺยากาสิํ, ตเถว ตํ วิปากํ, อญฺญถา วา”ติ.\csromanspacer{} ยเถว เม ภนฺเต ภควา อเจลํ โกรกฺขตฺติยํ อารพฺภ พฺยากาสิ, ตเถว ตํ วิปากํ, โน อญฺญถาติ. \csromanfolio{7}ตํ กิํมญฺญสิ สุนกฺขตฺต, ยทิ เอวํ สนฺเต กตํ วา โหติ อุตฺตริ มนุสฺสธมฺมา อิทฺธิปาฏิหาริยํ, อกตํ วาติ.\csromanspacer{} อทฺธา โข ภนฺเต เอวํ สนฺเต กตํ โหติ อุตฺตริ มนุสฺสธมฺมา อิทฺธิปาฏิหาริยํ, โน อกตนฺติ.\csromanspacer{} เอวมฺปิ โข มํ ตฺวํ โมฆปุริส อุตฺตริ มนุสฺสธมฺมา อิทฺธิปาฏิหาริยํ กโรนฺตํ เอวํ วเทสิ “น หิ ปน เม ภนฺเต ภควา อุตฺตริ มนุสฺสธมฺมา อิทฺธิปาฏิหาริยํ กโรตี”ติ.\csromanspacer{} ปสฺส โมฆปุริส ยาวจ เต อิทํ อปรทฺธนฺติ.\csromanspacer{} เอวมฺปิ โข ภคฺคว สุนกฺขตฺโต ลิจฺฉวิปุตฺโต มยา วุจฺจมาโน อปกฺกเมว อิมสฺมา ธมฺมวินยา, ยถา ตํ อาปายิโก เนรยิโก.}

\csromanmatikaanchor{mtk.4}
\csromantocmarkref{section}{อเจลกฬารมฏฺฏกวตฺถุ}{mtk.4}
\bookmark[dest={mtk.4},level=1]{อเจลกฬารมฏฺฏกวตฺถุ}
\csromanheader{1}{\textbf{อเจลกฬารมฏฺฏกวตฺถุ}}
\proseitem{11}{เอกมิทาหํ ภคฺคว สมยํ เวสาลิยํ วิหรามิ มหาวเน กูฏาคารสาลายํ.\csromanspacer{} เตน โข ปน สมเยน อเจโล กฬารมฏฺฏโก เวสาลิยํ ปฏิวสติ ลาภคฺคปฺปตฺโต เจว ยสคฺคปฺปตฺโต จ วชฺชิคาเม.\csromanspacer{} ตสฺส สตฺตวตปทานิ\footnote{สตฺตวตฺตปทานิ (สฺยา, อิ) 2. พหุปุตฺตกํ นาม (สฺยา) 3. อสาทิยิมฺหเส (สฺยา)} สมตฺตานิ สมาทินฺนานิ โหนฺติ.\csromanspacer{} ยาวชีวํ อเจลโก อสฺสํ, น วตฺถํ ปริทเหยฺยํ.\csromanspacer{} ยาวชีวํ พฺรหฺมจารี อสฺสํ, น เมถุนํ ธมฺมํ ปฏิเสเวยฺยํ.\csromanspacer{} ยาวชีวํ สุรามํเสเนว ยาเปยฺยํ, น โอทนกุมฺมาสํ ภุญฺเชยฺยํ.\csromanspacer{} ปุรตฺถิเมน เวสาลิํ อุเทนํ นาม เจติยํ, ตํ นาติกฺกเมยฺยํ.\csromanspacer{} ทกฺขิเณน เวสาลิํ โคตมกํ นาม เจติยํ, ตํ นาติกฺกเมยฺยํ.\csromanspacer{} ปจฺฉิเมน เวสาลิํ สตฺตมฺพํ นาม เจติยํ, ตํ นาติกฺกเมยฺยํ.\csromanspacer{} อุตฺตเรน เวสาลิํ พหุปุตฺตํ นาม2 เจติยํ, ตํ นาติกฺกเมยฺยนฺติ.\csromanspacer{} โส อิเมสํ สตฺตนฺนํ วตปทานํ สมาทานเหตุ ลาภคฺคปฺปตฺโต เจว ยสคฺคปฺปตฺโต จ วชฺชิคาเม.}

\proseitem{12}{อถ โข ภคฺคว สุนกฺขตฺโต ลิจฺฉวิปุตฺโต เยน อเจโล กฬารมฏฺฏโก เตนุปสงฺกมิ, อุปสงฺกมิตฺวา อเจลํ กฬารมฏฺฏกํ ปญฺหํ อปุจฺฉิ.\csromanspacer{} ตสฺส อเจโล กฬารมฏฺฏโก ปญฺหํ ปุฏฺโฐ น สมฺปายาสิ.\csromanspacer{} อสมฺปายนฺโต โกปญฺจ โทสญฺจ อปฺปจฺจยญฺจ ปาตฺวากาสิ.\csromanspacer{} อถ โข ภคฺคว สุนกฺขตฺตสฺส ลิจฺฉวิปุตฺตสฺส เอตทโหสิ “สาธุรูปํ วต โภ อรหนฺตํ สมณํ อาสาทิมฺหเส3.\csromanspacer{} มา วต โน อโหสิ ทีฆรตฺตํ อหิตาย ทุกฺขายา”ติ.}

\proseitem{13}{\csromanfolio{8}อถ โข ภคฺคว สุนกฺขตฺโต ลิจฺฉวิปุตฺโต เยนาหํ เตนุปสงฺกมิ, อุปสงฺกมิตฺวา มํ อภิวาเทตฺวา เอกมนฺตํ นิสีทิ.\csromanspacer{} เอกมนฺตํ นิสินฺนํ โข อหํ ภคฺคว สุนกฺขตฺตํ ลิจฺฉวิปุตฺตํ เอตทโวจํ “ตฺวํปิ นาม โมฆปุริส สมโณ สกฺยปุตฺติโย ปฏิชานิสฺสสี”ติ.\csromanspacer{} กิํ ปน มํ ภนฺเต ภควา เอวมาห “ตฺวํปิ นาม โมฆปุริส สมโณ สกฺยปุตฺติโย ปฏิชานิสฺสสี”ติ.\csromanspacer{} นนุ ตฺวํ สุนกฺขตฺต อเจลํ กฬารมฏฺฏกํ อุปสงฺกมิตฺวา ปญฺหํ อปุจฺฉิ, ตสฺส เต อเจโล กฬารมฏฺฏโก ปญฺหํ ปุฏฺโฐ น สมฺปายาสิ, อสมฺปายนฺโต โกปญฺจ โทสญฺจ อปฺปจฺจยญฺจ ปาตฺวากาสิ.\csromanspacer{} ตสฺส เต เอตทโหสิ “สาธุรูปํ วต โภ อรหนฺตํ สมณํ อาสาทิมฺหเส.\csromanspacer{} มา วต โน อโหสิ ทีฆรตฺตํ อหิตาย ทุกฺขายา”ติ.\csromanspacer{} เอวํ ภนฺเต.\csromanspacer{} กิํ ปน ภนฺเต ภควา อรหตฺตสฺส มจฺฉรายตีติ.}

\prose{น โข อหํ โมฆปุริส อรหตฺตสฺส มจฺฉรายามิ, อปิ จ ตุเยฺหเวตํ ปาปกํ ทิฏฺฐิคตํ อุปฺปนฺนํ, ตํ ปชห.\csromanspacer{} มา เต อโหสิ ทีฆรตฺตํ อหิตาย ทุกฺขายาติ.\csromanspacer{} ยํ โข ปเนตํ สุนกฺขตฺต มญฺญสิ อเจลํ กฬารมฏฺฏกํ “สาธุรูโป อยํ\footnote{อรหํ (สฺยา)} สมโณ”ติ.\csromanspacer{} โส นจิรสฺเสว ปริหิโต สานุจาริโก วิจรนฺโต โอทนกุมฺมาสํ ภุญฺชมาโน สพฺพาเนว เวสาลิยานิ เจติยานิ สมติกฺกมิตฺวา ยสา นิหีโน\footnote{ยสานิกิณฺโณ (ก)} กาลํ กริสฺสตีติ.\csromanspacer{} อถ โข ภคฺคว อเจโล กฬารมฏฺฏโก นจิรสฺเสว ปริหิโต สานุจาริโก วิจรนฺโต โอทนกุมฺมาสํ ภุญฺชมาโน สพฺพาเนว เวสาลิยานิ เจติยานิ สมติกฺกมิตฺวา ยสา นิหีโน กาลมกาสิ.}

\proseitem{14}{อสฺโสสิ โข ภคฺคว สุนกฺขตฺโต ลิจฺฉวิปุตฺโต “อเจโล กิร กฬารมฏฺฏโก ปริหิโต สานุจาริโก วิจรนฺโต โอทนกุมฺมาสํ ภุญฺชมาโน สพฺพาเนว เวสาลิยานิ เจติยานิ สมติกฺกมิตฺวา ยสา นิหีโน กาลงฺกโต”ติ.\csromanspacer{} อถ โข ภคฺคว สุนกฺขตฺโต ลิจฺฉวิปุตฺโต เยนาหํ เตนุปสงฺกมิ, อุปสงฺกมิตฺวา มํ อภิวาเทตฺวา เอกมนฺตํ นิสีทิ.\csromanspacer{} เอกมนฺตํ นิสินฺนํ โข อหํ ภคฺคว สุนกฺขตฺตํ ลิจฺฉวิปุตฺตํ เอตทโวจํ “ตํ กิํมญฺญสิ สุนกฺขตฺต, ยเถว เต อหํ อเจลํ กฬารมฏฺฏกํ อารพฺภ พฺยากาสิํ, ตเถว ตํ วิปากํ, อญฺญถา วา”ติ.\csromanspacer{} ยเถว เม ภนฺเต ภควา อเจลํ กฬารมฏฺฏกํ อารพฺภ พฺยากาสิ, ตเถว ตํ วิปากํ, \csromanfolio{9}โน อญฺญถาติ.\csromanspacer{} ตํ กิํมญฺญสิ สุนกฺขตฺต, ยทิ เอวํ สนฺเต กตํ วา โหติ อุตฺตริ มนุสฺสธมฺมา อิทฺธิปาฏิหาริยํ อกตํ วาติ.\csromanspacer{} อทฺธา โข ภนฺเต เอวํ สนฺเต กตํ โหติ อุตฺตริ มนุสฺสธมฺมา อิทฺธิปาฏิหาริยํ, โน อกตนฺติ.\csromanspacer{} เอวํปิ โข มํ ตฺวํ โมฆปุริส อุตฺตริ มนุสฺสธมฺมา อิทฺธิปาฏิหาริยํ กโรนฺตํ เอวํ วเทสิ “น หิ ปน เม ภนฺเต ภควา อุตฺตริ มนุสฺสธมฺมา อิทฺธิปาฏิหาริยํ กโรตี”ติ.\csromanspacer{} ปสฺส โมฆปุริส ยาวจ เต อิทํ อปรทฺธนฺติ.\csromanspacer{} เอวมฺปิ โข ภคฺคว สุนกฺขตฺโต ลิจฺฉวิปุตฺโต มยา วุจฺจมาโน อปกฺกเมว อิมสฺมา ธมฺมวินยา, ยถา ตํ อาปายิโก เนรยิโก.}

\csromanmatikaanchor{mtk.5}
\csromantocmarkref{section}{อเจลปาถิกปุตฺตวตฺถุ}{mtk.5}
\bookmark[dest={mtk.5},level=1]{อเจลปาถิกปุตฺตวตฺถุ}
\csromanheader{1}{\textbf{อเจลปาถิกปุตฺตวตฺถุ}}
\proseitem{15}{เอกมิทาหํ ภคฺคว สมยํ ตตฺเถว เวสาลิยํ วิหรามิ มหาวเน กูฏาคารสาลายํ.\csromanspacer{} เตน โข ปน สมเยน อเจโล ปาถิกปุตฺโต\footnote{ปาฏิกปุตฺโต (สี, สฺยา, อิ)} เวสาลิยํ ปฏิวสติ ลาภคฺคปฺปตฺโต เจว ยสคฺคปฺปตฺโต จ วชฺชิคาเม.\csromanspacer{} โส จ เวสาลิยํ ปริสติ เอวํ วาจํ ภาสติ “สมโณ โคตโม ญาณวาโท, อหมฺปิ ญาณวาโท.\csromanspacer{} ญาณวาโท โข ปน ญาณวาเทน อรหติ อุตฺตริ มนุสฺสธมฺมา อิทฺธิปาฏิหาริยํ ทสฺเสตุํ.\csromanspacer{} สมโณ โคตโม อุปฑฺฒปถํ อาคจฺเฉยฺย, อหมฺปิ อุปฑฺฒปถํ คจฺเฉยฺยํ.\csromanspacer{} เต ตตฺถ อุโภปิ อุตฺตริ มนุสฺสธมฺมา อิทฺธิปาฏิหาริยํ กเรยฺยาม.\csromanspacer{} เอกํ เจ สมโณ โคตโม อุตฺตริ มนุสฺสธมฺมา อิทฺธิปาฏิหาริยํ กริสฺสติ, ทฺวาหํ กริสฺสามิ.\csromanspacer{} เทฺว เจ สมโณ โคตโม อุตฺตริ มนุสฺสธมฺมา อิทฺธิปาฏิหาริยานิ กริสฺสติ, จตฺตาราหํ กริสฺสามิ.\csromanspacer{} จตฺตาริ เจ สมโณ โคตโม อุตฺตริ มนุสฺสธมฺมา อิทฺธิปาฏิหาริยานิ กริสฺสติ, อฏฺฐาหํ กริสฺสามิ.\csromanspacer{} อิติ ยาวตกํ ยาวตกํ สมโณ โคตโม อุตฺตริ มนุสฺสธมฺมา อิทฺธิปาฏิหาริยํ กริสฺสติ, ตทฺทิคุณํ ตทฺทิคุณาหํ กริสฺสามี”ติ.}

\proseitem{16}{อถ โข ภคฺคว สุนกฺขตฺโต ลิจฺฉวิปุตฺโต เยนาหํ เตนุปสงฺกมิ, อุปสงฺกมิตฺวา มํ อภิวาเทตฺวา เอกมนฺตํ นิสีทิ.\csromanspacer{} เอกมนฺตํ นิสินฺโน โข ภคฺคว สุนกฺขตฺโต ลิจฺฉวิปุตฺโต มํ เอตทโวจ\csromandash{}อเจโล ภนฺเต ปาถิกปุตฺโต เวสาลิยํ ปฏิวสติ ลาภคฺคปฺปตฺโต เจว ยสคฺคปฺปตฺโต จ วชฺชิคาเม, โส เวสาลิยํ ปริสติ เอวํ วาจํ ภาสติ “สมโณ \csromanfolio{10}โคตโม ญาณวาโท, อหมฺปิ ญาณวาโท.\csromanspacer{} ญาณวาโท โข ปน ญาณวาเทน อรหติ อุตฺตริ มนุสฺสธมฺมา อิทฺธิปาฏิหาริยํ ทสฺเสตุํ.\csromanspacer{} สมโณ โคตโม อุปฑฺฒปถํ อาคจฺเฉยฺย, อหมฺปิ อุปฑฺฒปถํ คจฺเฉยฺยํ.\csromanspacer{} เต ตตฺถ อุโภปิ อุตฺตริ มนุสฺสธมฺมา อิทฺธิปาฏิหาริยํ กเรยฺยาม.\csromanspacer{} เอกํ เจ สมโณ โคตโม อุตฺตริ มนุสฺสธมฺมา อิทฺธิปาฏิหาริยํ กริสฺสติ, ทฺวาหํ กริสฺสามิ.\csromanspacer{} เทฺว เจ สมโณ โคตโม อุตฺตริ มนุสฺสธมฺมา อิทฺธิปาฏิหาริยานิ กริสฺสติ, จตฺตาราหํ กริสฺสามิ.\csromanspacer{} จตฺตาริ เจ สมโณ โคตโม อุตฺตริ มนุสฺสธมฺมา อิทฺธิปาฏิหาริยานิ กริสฺสติ, อฏฺฐาหํ กริสฺสามิ.\csromanspacer{} อิติ ยาวตกํ ยาวตกํ สมโณ โคตโม อุตฺตริ มนุสฺสธมฺมา อิทฺธิปาฏิหาริยํ กริสฺสติ, ตทฺทิคุณํ ตทฺทิคุณาหํ กริสฺสามี”ติ.}

\prose{เอวํ วุตฺเต อหํ ภคฺคว สุนกฺขตฺตํ ลิจฺฉวิปุตฺตํ เอตทโวจํ “อภพฺโพ โข สุนกฺขตฺต อเจโล ปาถิกปุตฺโต ตํ วาจํ อปฺปหาย ตํ จิตฺตํ อปฺปหาย ตํ ทิฏฺฐิํ อปฺปฏินิสฺสชฺชิตฺวา มม สมฺมุขีภาวํ อาคนฺตุํ.\csromanspacer{} สเจปิสฺส เอวมสฺส ‘อหํ ตํ วาจํ อปฺปหาย ตํ จิตฺตํ อปฺปหาย ตํ ทิฏฺฐิํ อปฺปฏินิสฺสชฺชิตฺวา สมณสฺส โคตมสฺส สมฺมุขีภาวํ คจฺเฉยฺยนฺ’ติ, มุทฺธาปิ ตสฺส วิปเตยฺยา”ติ.}

\proseitem{17}{รกฺขเตตํ ภนฺเต ภควา วาจํ, รกฺขเตตํ สุคโต วาจนฺติ.\csromanspacer{} กิํ ปน มํ ตฺวํ สุนกฺขตฺต เอวํ วเทสิ “รกฺขเตตํ ภนฺเต ภควา วาจํ, รกฺขเตตํ สุคโต วาจนฺ”ติ.\csromanspacer{} ภควตา จสฺส ภนฺเต เอสา วาจา เอกํเสน โอธาริตา\footnote{โอวาทิตา (ก)} “อภพฺโพ อเจโล ปาถิกปุตฺโต ตํ วาจํ อปฺปหาย ตํ จิตฺตํ อปฺปหาย ตํ ทิฏฺฐิํ อปฺปฏินิสฺสชฺชิตฺวา มม สมฺมุขีภาวํ อาคนฺตุํ.\csromanspacer{} สเจปิสฺส เอวมสฺส ‘อหํ ตํ วาจํ อปฺปหาย ตํ จิตฺตํ อปฺปหาย ตํ ทิฏฺฐิํ อปฺปฏินิสฺสชฺชิตฺวา สมณสฺส โคตมสฺส สมฺมุขีภาวํ คจฺเฉยฺยนฺ’ติ, มุทฺธาปิ ตสฺส วิปเตยฺยา”ติ.\csromanspacer{} อเจโล จ ภนฺเต ปาถิกปุตฺโต วิรูปรูเปน ภควโต สมุมฺขีภาวํ อาคจฺเฉยฺย, ตทสฺส ภควโต มุสาติ.}

\proseitem{18}{อปิ นุ สุนกฺขตฺต ตถาคโต ตํ วาจํ ภาเสยฺย, ยา สา วาจา ทฺวยคามินีติ.\csromanspacer{} กิํ ปน ภนฺเต ภควตา อเจโล ปาถิกปุตฺโต \csromanfolio{11}เจตสา เจโต ปริจฺจ วิทิโต “อภพฺโพ อเจโล ปาถิกปุตฺโต ตํ วาจํ อปฺปหาย ตํ จิตฺตํ อปฺปหาย ตํ ทิฏฺฐิํ อปฺปฏินิสฺสชฺชิตฺวา มม สมฺมุขีภาวํ อาคนฺตุํ.\csromanspacer{} สเจปิสฺส เอวมสฺส ‘อหํ ตํ วาจํ อปฺปหาย ตํ จิตฺตํ อปฺปหาย ตํ ทิฏฺฐิํ อปฺปฏินิสฺสชฺชิตฺวา สมณสฺส โคตมสฺส สมฺมุขีภาวํ คจฺเฉยฺยนฺ’ติ, มุทฺธาปิ ตสฺส วิปเตยฺยา”ติ.}

\prose{อุทาหุ เทวตา ภควโต เอตมตฺถํ อาโรเจสุํ “อภพฺโพ ภนฺเต อเจโล ปาถิกปุตฺโต ตํ วาจํ อปฺปหาย ตํ จิตฺตํ อปฺปหาย ตํ ทิฏฺฐิํ อปฺปฏินิสฺสชฺชิตฺวา ภควโต สมฺมุขีภาวํ อาคนฺตุํ.\csromanspacer{} สเจปิสฺส เอวมสฺส ‘อหํ ตํ วาจํ อปฺปหาย ตํ จิตฺตํ อปฺปหาย ตํ ทิฏฺฐิํ อปฺปฏินิสฺสชฺชิตฺวา สมณสฺส โคตมสฺส สมฺมุขีภาวํ คจฺเฉยฺยนฺ’ติ, มุทฺธาปิ ตสฺส วิปเตยฺยา”ติ.}

\proseitem{19}{เจตสา เจโต ปริจฺจ วิทิโต เจว เม สุนกฺขตฺต “อภพฺโพ อเจโล ปาถิกปุตฺโต ตํ วาจํ อปฺปหาย ตํ จิตฺตํ อปฺปหาย ตํ ทิฏฺฐิํ อปฺปฏินิสฺสชฺชิตฺวา มม สมฺมุขีภาวํ อาคนฺตุํ.\csromanspacer{} สเจปิสฺส เอวมสฺส ‘อหํ ตํ วาจํ อปฺปหาย ตํ จิตฺตํ อปฺปหาย ตํ ทิฏฺฐิํ อปฺปฏินิสฺสชฺชิตฺวา สมณสฺส โคตมสฺส สมฺมุขีภาวํ คจฺเฉยฺยนฺ’ติ, มุทฺธาปิ ตสฺส วิปเตยฺยา”ติ.}

\prose{เทวตาปิ เม เอตมตฺถํ อาโรเจสุํ “อภพฺโพ ภนฺเต อเจโล ปาถิกปุตฺโต ตํ วาจํ อปฺปหาย ตํ จิตฺตํ อปฺปหาย ตํ ทิฏฺฐิํ อปฺปฏินิสฺสชฺชิตฺวา ภควโต สมฺมุขีภาวํ อาคนฺตุํ.\csromanspacer{} สเจปิสฺส เอวมสฺส ‘อหํ ตํ วาจํ อปฺปหาย ตํ จิตฺตํ อปฺปหาย ตํ ทิฏฺฐิํ อปฺปฏินิสฺสชฺชิตฺวา สมณสฺส โคตมสฺส สมฺมุขีภาวํ คจฺเฉยฺยนฺ’ติ, มุทฺธาปิ ตสฺส วิปเตยฺยา”ติ.}

\prose{อชิโตปิ นาม ลิจฺฉวีนํ เสนาปติ อธุนา กาลงฺกโต ตาวติํสกายํ อุปปนฺโน.\csromanspacer{} โสปิ มํ อุปสงฺกมิตฺวา เอวมาโรเจสิ “อลชฺชี ภนฺเต อเจโล ปาถิกปุตฺโต, มุสาวาที ภนฺเต อเจโล ปาถิกปุตฺโต, มํปิ ภนฺเต อเจโล ปาถิกปุตฺโต พฺยากาสิ วชฺชิคาเม ‘อชิโต ลิจฺฉวีนํ เสนาปติ มหานิรยํ อุปปนฺโน’ติ.\csromanspacer{} น โข ปนาหํ ภนฺเต มหานิรยํ อุปปนฺโน, ตาวติํสกายมฺหิ อุปปนฺโน.\csromanspacer{} อลชฺชี ภนฺเต อเจโล ปาถิกปุตฺโต, มุสาวาที ภนฺเต อเจโล ปาถิกปุตฺโต, อภพฺโพ จ ภนฺเต อเจโล ปาถิกปุตฺโต ตํ วาจํ อปฺปหาย ตํ จิตฺตํ อปฺปหาย ตํ ทิฏฺฐิํ อปฺปฏินิสฺสชฺชิตฺวา ภควโต สมฺมุขีภาวํ อาคนฺตุํ. \csromanfolio{12}สเจปิสฺส เอวมสฺส ‘อหํ ตํ วาจํ อปฺปหาย ตํ จิตฺตํ อปฺปหาย ตํ ทิฏฺฐิํ อปฺปฏินิสฺสชฺชิตฺวา สมณสฺส โคตมสฺส สมฺมุขีภาวํ คจฺเฉยฺยนฺ’ติ, มุทฺธาปิ ตสฺส วิปเตยฺยา”ติ.}

\prose{อิติ โข สุนกฺขตฺต เจตสา เจโต ปริจฺจ วิทิโต เจว เม อเจโล ปาถิกปุตฺโต “อภพฺโพ อเจโล ปาถิกปุตฺโต ตํ วาจํ อปฺปหาย ตํ จิตฺตํ อปฺปหาย ตํ ทิฏฺฐิํ อปฺปฏินิสฺสชฺชิตฺวา มม สมฺมุขีภาวํ อาคนฺตุํ.\csromanspacer{} สเจปิสฺส เอวมสฺส ‘อหํ ตํ วาจํ อปฺปหาย ตํ จิตฺตํ อปฺปหาย ตํ ทิฏฺฐิํ อปฺปฏินิสฺสชฺชิตฺวา สมณสฺส โคตมสฺส สมฺมุขีภาวํ คจฺเฉยฺยนฺ’ติ, มุทฺธาปิ ตสฺส วิปเตยฺยา”ติ.\csromanspacer{} เทวตาปิ เม เอตมตฺถํ อาโรเจสุํ “อภพฺโพ ภนฺเต อเจโล ปาถิกปุตฺโต ตํ วาจํ อปฺปหาย ตํ จิตฺตํ อปฺปหาย ตํ ทิฏฺฐิํ อปฺปฏินิสฺสชฺชิตฺวา ภควโต สมฺมุขีภาวํ อาคนฺตุํ.\csromanspacer{} สเจปิสฺส เอวมสฺส “อหํ ตํ วาจํ อปฺปหาย ตํ จิตฺตํ อปฺปหาย ตํ ทิฏฺฐิํ อปฺปฏินิสฺสชฺชิตฺวา สมณสฺส โคตมสฺส สมฺมุขีภาวํ คจฺเฉยฺยนฺ’ติ, มุทฺธาปิ ตสฺส วิปเตยฺยา”ติ.}

\prose{โส โข ปนาหํ สุนกฺขตฺต เวสาลิยํ ปิณฺฑาย จริตฺวา ปจฺฉาภตฺตํ ปิณฺฑปาตปฏิกฺกนฺโต เยน อเจลสฺส ปาถิกปุตฺตสฺส อาราโม เตนุปสงฺกมิสฺสามิ ทิวาวิหาราย, ยสฺสทานิ ตฺวํ สุนกฺขตฺต อิจฺฉสิ, ตสฺส อาโรเจหีติ.}

\csromanmatikaanchor{mtk.6}
\csromantocmarkref{section}{อิทฺธิปาฏิหาริยกถา}{mtk.6}
\bookmark[dest={mtk.6},level=1]{อิทฺธิปาฏิหาริยกถา}
\csromanheader{1}{\textbf{อิทฺธิปาฏิหาริยกถา}}
\proseitem{20}{อถ ขฺวาหํ\footnote{อถ โข สฺวาหํ (สฺยา)} ภคฺคว ปุพฺพณฺหสมยํ นิวาเสตฺวา ปตฺตจีวรมาทาย เวสาลิํ ปิณฺฑาย ปาวิสิํ, เวสาลิยํ ปิณฺฑาย จริตฺวา ปจฺฉาภตฺตํ ปิณฺฑปาตปฏิกฺกนฺโต เยน อเจลสฺส ปาถิกปุตฺตสฺส อาราโม เตนุปสงฺกมิํ ทิวาวิหาราย.\csromanspacer{} อถ โข ภคฺคว สุนกฺขตฺโต ลิจฺฉวิปุตฺโต ตรมานรูโป เวสาลิํ ปวิสิตฺวา เยน อภิญฺญาตา อภิญฺญาตา ลิจฺฉวี เตนุปสงฺกมิ, อุปสงฺกมิตฺวา อภิญฺญาเต อภิญฺญาเต ลิจฺฉวี เอตทโวจ “เอสาวุโส ภควา เวสาลิยํ ปิณฺฑาย จริตฺวา ปจฺฉาภตฺตํ ปิณฺฑปาตปฏิกฺกนฺโต เยน อเจลสฺส ปาถิกปุตฺตสฺส อาราโม เตนุปสงฺกมิ ทิวาวิหาราย.\csromanspacer{} อภิกฺกมถายสฺมนฺโต อภิกฺกมถายสฺมนฺโต, สาธุรูปานํ สมณานํ อุตฺตริ มนุสฺสธมฺมา อิทฺธิปาฏิหาริยํ \csromanfolio{13}ภวิสฺสตี”ติ.\csromanspacer{} อถ โข ภคฺคว อภิญฺญาตานํ อภิญฺญาตานํ ลิจฺฉวีนํ เอตทโหสิ “สาธุรูปานํ กิร โภ สมณานํ อุตฺตริ มนุสฺสธมฺมา อิทฺธิปาฏิหาริยํ ภวิสฺสติ, หนฺท วต โภ คจฺฉามา”ติ.}

\prose{เยน จ อภิญฺญาตา อภิญฺญาตา พฺราหฺมณมหาสาลา.\csromanspacer{} คหปติเนจยิกา.\csromanspacer{} นานาติตฺถิยา\footnote{นานาติตฺถิย (สฺยา)} สมณพฺราหฺมณา เตนุปสงฺกมิ, อุปสงฺกมิตฺวา อภิญฺญาเต อภิญฺญาเต นานาติตฺถิเย1 สมณพฺราหฺมเณ เอตทโวจ “เอสาวุโส ภควา เวสาลิยํ ปิณฺฑาย จริตฺวา ปจฺฉาภตฺตํ ปิณฺฑปาตปฏิกฺกนฺโต เยน อเจลสฺส ปาถิกปุตฺตสฺส อาราโม เตนุปสงฺกมิ ทิวาวิหาราย.\csromanspacer{} อภิกฺกมถายสฺมนฺโต อภิกฺกมถายสฺมนฺโต, สาธุรูปานํ สมณานํ อุตฺตริ มนุสฺสธมฺมา อิทฺธิปาฏิหาริยํ ภวิสฺสตี”ติ.\csromanspacer{} อถ โข ภคฺคว อภิญฺญาตานํ อภิญฺญาตานํ นานาติตฺถิยานํ สมณพฺราหฺมณานํ เอตทโหสิ “สาธุรูปานํ กิร โภ สมณานํ อุตฺตริ มนุสฺสธมฺมา อิทฺธิปาฏิหาริยํ ภวิสฺสติ, หนฺท วต โภ คจฺฉามา”ติ.}

\prose{อถ โข ภคฺคว อภิญฺญาตา อภิญฺญาตา ลิจฺฉวี, อภิญฺญาตา อภิญฺญาตา จ พฺราหฺมณมหาสาลา คหปติเนจยิกา นานาติตฺถิยา สมณพฺราหฺมณา เยน อเจลสฺส ปาถิกปุตฺตสฺส อาราโม เตนุปสงฺกมิํสุ, สา เอสา ภคฺคว ปริสา มหา โหติ\footnote{ปริสา โหติ (สี, สฺยา, อิ)} อเนกสตา อเนกสหสฺสา.}

\proseitem{21}{อสฺโสสิ โข ภคฺคว อเจโล ปาถิกปุตฺโต “อภิกฺกนฺตา กิร อภิญฺญาตา อภิญฺญาตา ลิจฺฉวี, อภิกฺกนฺตา อภิญฺญาตา อภิญฺญาตา จ พฺราหฺมณมหาสาลา คหปติเนจยิกา นานาติตฺถิยา สมณพฺราหฺมณา.\csromanspacer{} สมโณปิ โคตโม มยฺหํ อาราเม ทิวาวิหารํ นิสินฺโน”ติ.\csromanspacer{} สุตฺวานสฺส ภยํ ฉมฺภิตตฺตํ โลมหํโส อุทปาทิ.\csromanspacer{} อถ โข ภคฺคว อเจโล ปาถิกปุตฺโต ภีโต สํวิคฺโค โลมหฏฺฐชาโต เยน ตินฺทุกขาณุปริพฺพาชการาโม เตนุปสงฺกมิ.}

\prose{อสฺโสสิ โข ภคฺคว สา ปริสา “อเจโล กิร ปาถิกปุตฺโต ภีโต สํวิคฺโค โลมหฏฺฐชาโต เยน ตินฺทุกขาณุปริพฺพาชการาโม \csromanfolio{14}เตนุปสงฺกนฺโต”ติ\footnote{เตนุปสงฺกมนฺโต (สี, อิ, ก)}.\csromanspacer{} อถ โข ภคฺคว สา ปริสา อญฺญตรํ ปุริสํ อามนฺเตสิ\csromandash{}}

\prose{เอหิ ตฺวํ โภ ปุริส เยน ตินฺทุกขาณุปริพฺพาชการาโม, เยน อเจโล ปาถิกปุตฺโต เตนุปสงฺกม, อุปสงฺกมิตฺวา อเจลํ ปาถิกปุตฺตํ เอวํ วเทหิ “อภิกฺกมาวุโส ปาถิกปุตฺต, อภิกฺกนฺตา อภิญฺญาตา อภิญฺญาตา ลิจฺฉวี, อภิกฺกนฺตา อภิญฺญาตา อภิญฺญาตา จ พฺราหฺมณมหาสาลา คหปติเนจยิกา นานาติตฺถิยา สมณพฺราหฺมณา.\csromanspacer{} สมโณปิ โคตโม อายสฺมโต อาราเม ทิวาวิหารํ นิสินฺโน.\csromanspacer{} ภาสิตา โข ปน เต เอสา อาวุโส ปาถิกปุตฺต เวสาลิยํ ปริสติ วาจา ‘สมโณปิ โคตโม ญาณวาโท, อหมฺปิ ญาณวาโท.\csromanspacer{} ญาณวาโท โข ปน ญาณวาเทน อรหติ อุตฺตริ มนุสฺสธมฺมา อิทฺธิปาฏิหาริยํ ทสฺเสตุํ.\csromanspacer{} สมโณ (เจ) โคตโม อุปฑฺฒปถํ อาคจฺเฉยฺย, อหมฺปิ อุปฑฺฒปถํ คจฺเฉยฺยํ.\csromanspacer{} เต ตตฺถ อุโภปิ อุตฺตริ มนุสฺสธมฺมา อิทฺธิปาฏิหาริยํ กเรยฺยาม.\csromanspacer{} เอกํ เจ สมโณ โคตโม อุตฺตริ มนุสฺสธมฺมา อิทฺธิปาฏิหาริยํ กริสฺสติ, ทฺวาหํ กริสฺสามิ.\csromanspacer{} เทฺว เจ สมโณ โคตโม อุตฺตริ มนุสฺสธมฺมา อิทฺธิปาฏิหาริยานิ กริสฺสติ, จตฺตาราหํ กริสฺสามิ.\csromanspacer{} จตฺตาริ เจ สมโณ โคตโม อุตฺตริ มนุสฺสธมฺมา อิทฺธิปาฏิหาริยานิ กริสฺสติ, อฏฺฐาหํ กริสฺสามิ.\csromanspacer{} อิติ ยาวตกํ ยาวตกํ สมโณ โคตโม อุตฺตริ มนุสฺสธมฺมา อิทฺธิปาฏิหาริยํ กริสฺสติ, ตทฺทิคุณํ ตทฺทิคุณาหํ กริสฺสามี’ติ.\csromanspacer{} อภิกฺกมสฺเสว\footnote{อภิกฺกมเยว (สี, สฺยา, อิ)} โข อาวุโส ปาถิกปุตฺต อุปฑฺฒปถํ.\csromanspacer{} สพฺพปฐมํเยว อาคนฺตฺวา สมโณ โคตโม อายสฺมโต อาราเม ทิวาวิหารํ นิสินฺโน”ติ.}

\proseitem{22}{“เอวํ โภ”ติ โข ภคฺคว โส ปุริโส ตสฺสา ปริสาย ปฏิสฺสุตฺวา เยน ตินฺทุกขาณุปริพฺพาชการาโม, เยน อเจโล ปาถิกปุตฺโต เตนุปสงฺกมิ, อุปสงฺกมิตฺวา อเจลํ ปาถิกปุตฺตํ เอตทโวจ “อภิกฺกมาวุโส ปาถิกปุตฺต อภิกฺกนฺตา อภิญฺญาตา อภิญฺญาตา ลิจฺฉวี, อภิกฺกนฺตา อภิญฺญาตา อภิญฺญาตา จ พฺราหฺมณมหาสาลา คหปติเนจยิกา นานาติตฺถิยา สมณพฺราหฺมณา.\csromanspacer{} สมโณปิ โคตโม อายสฺมโต อาราเม ทิวาวิหารํ นิสินฺโน. \csromanfolio{15}ภาสิตา โข ปน เต เอสา อาวุโส ปาถิกปุตฺต เวสาลิยํ ปริสติ วาจา ‘สมโณ โคตโม ญาณวาโท, อหมฺปิ ญาณวาโท.\csromanspacer{} ญาณวาโท โข ปน ญาณวาเทน อรหติ อุตฺตริ มนุสฺสธมฺมา อิทฺธิปาฏิหาริยํ ทสฺเสตุํ ฯเปฯ ตทฺทิคุณํ ตทฺทิคุณาหํ กริสฺสามี’ติ.\csromanspacer{} อภิกฺกมสฺเสว โข อาวุโส ปาถิกปุตฺต อุปฑฺฒปถํ.\csromanspacer{} สพฺพปฐมํ เยว อาคนฺตฺวา สมโณ โคตโม อายสฺมโต อาราเม ทิวาวิหารํ นิสินฺโน”ติ.}

\prose{เอวํ วุตฺเต ภคฺคว อเจโล ปาถิกปุตฺโต “อายามิ อาวุโส อายามิ อาวุโส”ติ วตฺวา ตตฺเถว สํสปฺปติ\footnote{สํสพฺพติ (ก)}, น สกฺโกติ อาสนาปิ วุฏฺฐาตุํ.\csromanspacer{} อถ โข โส ภคฺคว ปุริโส อเจลํ ปาถิกปุตฺตํ เอตทโวจ “กิํ สุ นาม เต อาวุโส ปาถิกปุตฺต ปาวฬา สุ นาม เต ปีฐกสฺมิํ อลฺลีนา, ปีฐกํ สุ นาม เต ปาวฬาสุ อลฺลีนํ. ‘อายามิ อาวุโส อายามิ อาวุโส’ติ วตฺวา ตตฺเถว สํสปฺปสิ, น สกฺโกสิ อาสนาปิ วุฏฺฐาตุนฺ”ติ.\csromanspacer{} เอวํปิ โข ภคฺคว วุจฺจมาโน อเจโล ปาถิกปุตฺโต “อายามิ อาวุโส อายามิ อาวุโส”ติ วตฺวา ตตฺเถว สํสปฺปติ, น สกฺโกติ อาสนาปิ วุฏฺฐาตุํ.}

\proseitem{23}{ยทา โข โส ภคฺคว ปุริโส อญฺญาสิ “ปราภูตรูโป อยํ อเจโล ปาถิกปุตฺโต ‘อายามิ อาวุโส อายามิ อาวุโส’ติ วตฺวา ตตฺเถว สํสปฺปติ, น สกฺโกติ อาสนาปิ วุฏฺฐาตุนฺ”ติ.\csromanspacer{} อถ ตํ ปริสํ อาคนฺตฺวา เอวมาโรเจสิ “ปราภูตรูโป โภ\footnote{ปราภูตรูโป โภ อยํ (สฺยา, ก), ปราภูตรูโป (สี, อิ)} อเจโล ปาถิกปุตฺโต ‘อายามิ อาวุโส อายามิ อาวุโส’ติ วตฺวา ตตฺเถว สํสปฺปติ, น สกฺโกติ อาสนาปิ วุฏฺฐาตุนฺ”ติ.\csromanspacer{} เอวํ วุตฺเต อหํ ภคฺคว ตํ ปริสํ เอตทโวจํ “อภพฺโพ โข อาวุโส อเจโล ปาถิกปุตฺโต ตํ วาจํ อปฺปหาย ตํ จิตฺตํ อปฺปหาย ตํ ทิฏฺฐิํ อปฺปฏินิสฺสชฺชิตฺวา มม สมฺมุขีภาวํ อาคนฺตุํ.\csromanspacer{} สเจปิสฺส เอวมสฺส ‘อหํ ตํ วาจํ อปฺปหาย ตํ จิตฺตํ อปฺปหาย ตํ ทิฏฺฐิํ อปฺปฏินิสฺสชฺชิตฺวา สมณสฺส โคตมสฺส สมฺมุขีภาวํ คจฺเฉยฺยนฺ’ติ, มุทฺธาปิ ตสฺส วิปเตยฺยา”ติ.}

\nitthitam{ปฐมภาณวาโร นิฏฺฐิโต.}
\proseitem{24}{\csromanfolio{16}อถ โข ภคฺคว อญฺญตโร ลิจฺฉวิมหามตฺโต อุฏฺฐายาสนา ตํ ปริสํ เอตทโวจ “เตน หิ โภ มุหุตฺตํ ตาว อาคเมถ, ยาวาหํ คจฺฉามิ\footnote{ปจฺจาคจฺฉามิ (?)}, อปฺเปว นาม อหมฺปิ สกฺกุเณยฺยํ อเจลํ ปาถิกปุตฺตํ อิมํ ปริสํ อาเนตุนฺ”ติ.}

\prose{อถ โข โส ภคฺคว ลิจฺฉวิมหามตฺโต เยน ตินฺทุกขาณุปริพฺพาชการาโม, เยน อเจโล ปาถิกปุตฺโต เตนุปสงฺกมิ, อุปสงฺกมิตฺวา อเจลํ ปาถิกปุตฺตํ เอตทโวจ\csromandash{}อภิกฺกมาวุโส ปาถิกปุตฺต, อภิกฺกนฺตํ เต เสยฺโย, อภิกฺกนฺตา อภิญฺญาตา อภิญฺญาตา ลิจฺฉวี, อภิกฺกนฺตา อภิญฺญาตา อภิญฺญาตา จ พฺราหฺมณมหาสาลา คหปติเนจยิกา นานาติตฺถิยา สมณพฺราหฺมณา.\csromanspacer{} สมโณปิ โคตโม อายสฺมโต อาราเม ทิวาวิหารํ นิสินฺโน.\csromanspacer{} ภาสิตา โข ปน เต เอสา อาวุโส ปาถิกปุตฺต เวสาลิยํ ปริสติ วาจา “สมโณปิ โคตโม ญาณวาโท ฯเปฯ ตทฺทิคุณํ ตทฺทิคุณาหํ กริสฺสามี”ติ.\csromanspacer{} อภิกฺกมสฺเสว โข อาวุโส ปาถิกปุตฺต อุปฑฺฒปถํ.\csromanspacer{} สพฺพปฐมํเยว อาคนฺตฺวา สมโณ โคตโม อายสฺมโต อาราเม ทิวาวิหารํ นิสินฺโน.\csromanspacer{} ภาสิตา โข ปเนสา อาวุโส ปาถิกปุตฺต สมเณน โคตเมน ปริสติ วาจา “อภพฺโพ โข อเจโล ปาถิกปุตฺโต ตํ วาจํ อปฺปหาย ตํ จิตฺตํ อปฺปหาย ตํ ทิฏฺฐิํ อปฺปฏินิสฺสชฺชิตฺวา มม สมฺมุขีภาวํ อาคนฺตุํ.\csromanspacer{} สเจปิสฺส เอวมสฺส ‘อหํ ตํ วาจํ อปฺปหาย ตํ จิตฺตํ อปฺปหาย ตํ ทิฏฺฐิํ อปฺปฏินิสฺสชฺชิตฺวา สมณสฺส โคตมสฺส สมฺมุขีภาวํ คจฺเฉยฺยนฺ’ติ, มุทฺธาปิ ตสฺส วิปเตยฺยา”ติ.\csromanspacer{} อภิกฺกมาวุโส ปาถิกปุตฺต, อภิกฺกมเนเนว เต ชยํ กริสฺสาม, สมณสฺส โคตมสฺส ปราชยนฺติ.}

\prose{เอวํ วุตฺเต ภคฺคว อเจโล ปาถิกปุตฺโต “อายามิ อาวุโส อายามิ อาวุโส”ติ วตฺวา ตตฺเถว สํสปฺปติ, น สกฺโกติ อาสนาปิ วุฏฺฐาตุํ.\csromanspacer{} อถ โข โส ภคฺคว ลิจฺฉวิมหามตฺโต อเจลํ ปาถิกปุตฺตํ เอตทโวจ “กิํ สุ นาม เต อาวุโส ปาถิกปุตฺต ปาวฬา สุ นาม เต ปีฐกสฺมิํ อลฺลีนา, ปีฐกํ สุ นาม เต ปาวฬาสุ \csromanfolio{17}อลฺลีนํ. ‘อายามิ อาวุโส อายามิ อาวุโส’ติ วตฺวา ตตฺเถว สํสปฺปสิ, น สกฺโกสิ อาสนาปิ วุฏฺฐาตุนฺ”ติ.\csromanspacer{} เอวํปิ โข ภคฺคว วุจฺจมาโน อเจโล ปาถิกปุตฺโต “อายามิ อาวุโส อายามิ อาวุโส”ติ วตฺวา ตตฺเถว สํสปฺปติ, น สกฺโกติ อาสนาปิ วุฏฺฐาตุํ.}

\proseitem{25}{ยทา โข โส ภคฺคว ลิจฺฉวิมหามตฺโต อญฺญาสิ “ปราภูตรูโป อยํ อเจโล ปาถิกปุตฺโต ‘อายามิ อาวุโส อายามิ อาวุโส’ติ วตฺวา ตตฺเถว สํสปฺปติ, น สกฺโกติ อาสนาปิ วุฏฺฐาตุนฺ”ติ.\csromanspacer{} อถ ตํ ปริสํ อาคนฺตฺวา เอวมาโรเจสิ “ปราภูตรูโป โภ\footnote{ปราภูตรูโป (สี, อิ), ปราภูตรูโป อยํ (สฺยา)} อเจโล ปาถิกปุตฺโต ‘อายามิ อาวุโส อายามิ อาวุโส’ติ วตฺวา ตตฺเถว สํสปฺปติ, น สกฺโกติ อาสนาปิ วุฏฺฐาตุนฺ”ติ.\csromanspacer{} เอวํ วุตฺเต อหํ ภคฺคว ตํ ปริสํ เอตทโวจํ “อภพฺโพ โข อาวุโส อเจโล ปาถิกปุตฺโต ตํ วาจํ อปฺปหาย ตํ จิตฺตํ อปฺปหาย ตํ ทิฏฺฐิํ อปฺปฏินิสฺสชฺชิตฺวา มม สมฺมุขีภาวํ อาคนฺตุํ.\csromanspacer{} สเจปิสฺส เอวมสฺส ‘อหํ ตํ วาจํ อปฺปหาย ตํ จิตฺตํ อปฺปหาย ตํ ทิฏฺฐิํ อปฺปฏินิสฺสชฺชิตฺวา สมณสฺส โคตมสฺส สมฺมุขีภาวํ คจฺเฉยฺยนฺ’ติ, มุทฺธาปิ ตสฺส วิปเตยฺย.\csromanspacer{} สเจปายสฺมนฺตานํ ลิจฺฉวีนํ เอวมสฺส ‘มยํ อเจลํ ปาถิกปุตฺตํ วรตฺตาหิ\footnote{ยาหิ วรตฺตาหิ (สฺยา, ก)} พนฺธิตฺวา โคยุเคหิ อาวิญฺเฉยฺยามา’ติ\footnote{อาวิญฺเชยฺยามาติ (สฺยา), อาวิชฺเฌยฺยามาติ (สี, อิ)}, ตา วรตฺตา ฉิชฺเชยฺยุํ ปาถิกปุตฺโต วา.\csromanspacer{} อภพฺโพ ปน อเจโล ปาถิกปุตฺโต ตํ วาจํ อปฺปหาย ตํ จิตฺตํ อปฺปหาย ตํ ทิฏฺฐิํ อปฺปฏินิสฺสชฺชิตฺวา มม สมฺมุขีภาวํ อาคนฺตุํ.\csromanspacer{} สเจปิสฺส เอวมสฺส ‘อหํ ตํ วาจํ อปฺปหาย ตํ จิตฺตํ อปฺปหาย ตํ ทิฏฺฐิํ อปฺปฏินิสฺสชฺชิตฺวา สมณสฺส โคตมสฺส สมฺมุขีภาวํ คจฺเฉยฺยนฺ’ติ, มุทฺธาปิ ตสฺส วิปเตยฺยา”ติ.}

\proseitem{26}{อถ โข ภคฺคว ชาลิโย ทารุปตฺติกนฺเตวาสี อุฏฺฐายาสนา ตํ ปริสํ เอตทโวจ “เตน หิ โภ มุหุตฺตํ ตาว อาคเมถ, ยาวาหํ คจฺฉามิ, อปฺเปว นาม อหมฺปิ สกฺกุเณยฺยํ อเจลํ ปาถิกปุตฺตํ อิมํ ปริสํ อาเนตุนฺ”ติ.}

\prose{\csromanfolio{18}อถ โข ภคฺคว ชาลิโย ทารุปตฺติกนฺเตวาสี เยน ตินฺทุกขาณุปริพฺพาชการาโม, เยน อเจโล ปาถิกปุตฺโต เตนุปสงฺกมิ, อุปสงฺกมิตฺวา อเจลํ ปาถิกปุตฺตํ เอตทโวจ\csromandash{}อภิกฺกมาวุโส ปาถิกปุตฺต, อภิกฺกนฺตํ เต เสยฺโย.\csromanspacer{} อภิกฺกนฺตา อภิญฺญาตา อภิญฺญาตา ลิจฺฉวี, อภิกฺกนฺตา อภิญฺญาตา อภิญฺญาตา จ พฺราหฺมณมหาสาลา คหปติเนจยิกา นานาติตฺถิยา สมณพฺราหฺมณา.\csromanspacer{} สมโณปิ โคตโม อายสฺมโต อาราเม ทิวาวิหารํ นิสินฺโน.\csromanspacer{} ภาสิตา โข ปน เต เอสา อาวุโส ปาถิกปุตฺต เวสาลิยํ ปริสติ วาจา “สมโณปิ โคตโม ญาณวาโท ฯเปฯ ตทฺทิคุณํ ตทฺทิคุณาหํ กริสฺสามี”ติ.\csromanspacer{} อภิกฺกมสฺเสว โข อาวุโส ปาถิกปุตฺต อุปฑฺฒปถํ, สพฺพปฐมํเยว อาคนฺตฺวา สมโณ โคตโม อายสฺมโต อาราเม ทิวาวิหารํ นิสินฺโน.\csromanspacer{} ภาสิตา โข ปเนสา อาวุโส ปาถิกปุตฺต สมเณน โคตเมน ปริสติ วาจา “อภพฺโพ อเจโล ปาถิกปุตฺโต ตํ วาจํ อปฺปหาย ตํ จิตฺตํ อปฺปหาย ตํ ทิฏฺฐิํ อปฺปฏินิสฺสชฺชิตฺวา มม สมฺมุขีภาวํ อาคนฺตุํ.\csromanspacer{} สเจปิสฺส เอวมสฺส ‘อหํ ตํ วาจํ อปฺปหาย ตํ จิตฺตํ อปฺปหาย ตํ ทิฏฺฐิํ อปฺปฏินิสฺสชฺชิตฺวา สมณสฺส โคตมสฺส สมฺมุขีภาวํ คจฺเฉยฺยนฺ’ติ, มุทฺธาปิ ตสฺส วิปเตยฺย.\csromanspacer{} สเจปายสฺมนฺตานํ ลิจฺฉวีนํ เอวมสฺส ‘มยํ อเจลํ ปาถิกปุตฺตํ วรตฺตาหิ พนฺธิตฺวา โคยุเคหิ อาวิญฺเฉยฺยามา’ติ.\csromanspacer{} ตา วรตฺตา ฉิชฺเชยฺยุํ ปาถิกปุตฺโต วา.\csromanspacer{} อภพฺโพ ปน อเจโล ปาถิกปุตฺโต ตํ วาจํ อปฺปหาย ตํ จิตฺตํ อปฺปหาย ตํ ทิฏฺฐิํ อปฺปฏินิสฺสชฺชิตฺวา มม สมฺมุขีภาวํ อาคนฺตุํ.\csromanspacer{} สเจปิสฺส เอวมสฺส ‘อหํ ตํ วาจํ อปฺปหาย ตํ จิตฺตํ อปฺปหาย ตํ ทิฏฺฐิํ อปฺปฏินิสฺสชฺชิตฺวา สมณสฺส โคตมสฺส สมฺมุขีภาวํ อาคจฺเฉยฺยนฺ’ติ, มุทฺธาปิ ตสฺส วิปเตยฺยา’ติ.\csromanspacer{} อภิกฺกมาวุโส ปาถิกปุตฺต, อภิกฺกมเนเนว เต ชยํ กริสฺสาม, สมณสฺส โคตมสฺส ปราชยนฺติ.}

\prose{เอวํ วุตฺเต ภคฺคว อเจโล ปาถิกปุตฺโต “อายามิ อาวุโส อายามิ อาวุโส”ติ วตฺวา ตตฺเถว สํสปฺปติ, น สกฺโกติ อาสนาปิ วุฏฺฐาตุํ.\csromanspacer{} อถ โข ภคฺคว ชาลิโย ทารุปตฺติกนฺเตวาสี อเจลํ ปาถิกปุตฺตํ เอตทโวจ “กิํ สุ นาม เต อาวุโส ปาถิกปุตฺต ปาวฬา สุ นาม เต ปีฐกสฺมิํ อลฺลีนา, ปีฐกํ สุ นาม เต ปาวฬาสุ อลฺลีนํ. ‘อายามิ อาวุโส อายามิ อาวุโส’ติ วตฺวา \csromanfolio{19}ตตฺเถว สํสปฺปสิ, น สกฺโกสิ อาสนาปิ วุฏฺฐาตุนฺ”ติ.\csromanspacer{} เอวมฺปิ โข ภคฺคว วุจฺจมาโน อเจโล ปาถิกปุตฺโต “อายามิ อาวุโส อายามิ อาวุโส”ติ วตฺวา ตตฺเถว สํสปฺปติ, น สกฺโกติ อาสนาปิ วุฏฺฐาตุนฺติ.}

\proseitem{27}{ยทา โข ภคฺคว ชาลิโย ทารุปตฺติกนฺเตวาสี อญฺญาสิ “ปราภูตรูโป อยํ อเจโล ปาถิกปุตฺโต ‘อายามิ อาวุโส อายามิ อาวุโส’ติ วตฺวา ตตฺเถว สํสปฺปติ, น สกฺโกติ อาสนาปิ วุฏฺฐาตุนฺ”ติ.\csromanspacer{} อถ นํ เอตทโวจ\csromandash{}}

\prose{ภูตปุพฺพํ อาวุโส ปาถิกปุตฺต สีหสฺส มิครญฺโญ เอตทโหสิ “ยํนูนาหํ อญฺญตรํ วนสณฺฑํ นิสฺสาย อาสยํ กปฺเปยฺยํ, ตตฺราสยํ กปฺเปตฺวา สายนฺหสมยํ อาสยา นิกฺขเมยฺยํ, อาสยา นิกฺขมิตฺวา วิชมฺเภยฺยํ, วิชมฺภิตฺวา สมนฺตา จตุทฺทิสา อนุวิโลเกยฺยํ, สมนฺตา จตุทฺทิสา อนุวิโลเกตฺวา ติกฺขตฺตุํ สีหนาทํ นเทยฺยํ, ติกฺขตฺตุํ สีหนาทํ นทิตฺวา โคจราย ปกฺกเมยฺยํ.\csromanspacer{} โส วรํ วรํ มิคสํเฆ\footnote{มิค สํฆํ (สฺยา, ก)} วธิตฺวา มุทุมํสานิ มุทุมํสานิ ภกฺขยิตฺวา ตเมว อาสยํ อชฺฌุเปยฺยนฺ”ติ.}

\prose{อถ โข อาวุโส โส สีโห มิคราชา อญฺญตรํ วนสณฺฑํ นิสฺสาย อาสยํ กปฺเปสิ, ตตฺราสยํ กปฺเปตฺวา สายนฺหสมยํ อาสยา นิกฺขมิ, อาสยา นิกฺขมิตฺวา วิชมฺภิ, วิชมฺภิตฺวา สมนฺตา จตุทฺทิสา อนุวิโลเกสิ, สมนฺตา จตุทฺทิสา อนุวิโลเกตฺวา ติกฺขตฺตุํ สีหนาทํ นทิ, ติกฺขตฺตุํ สีหนาทํ นทิตฺวา โคจราย ปกฺกมิ.\csromanspacer{} โส วรํ วรํ มิคสํเฆ วธิตฺวา มุทุมํสานิ มุทุมํสานิ ภกฺขยิตฺวา ตเมว อาสยํ อชฺฌุเปสิ.}

\proseitem{28}{ตสฺเสว โข อาวุโส ปาถิกปุตฺต สีหสฺส มิครญฺโญ วิฆาสสํวฑฺโฒ ชรสิงฺคาโล\footnote{ชรสิคาโล (สี, สฺยา, อิ)} ทิตฺโต เจว พลวา จ.\csromanspacer{} อถ โข อาวุโส ตสฺส ชรสิงฺคาลสฺส เอตทโหสิ “โก จาหํ โก สีโห มิคราชา, ยํนูนาหํปิ อญฺญตรํ วนสณฺฑํ นิสฺสาย อาสยํ กปฺเปยฺยํ, ตตฺราสยํ กปฺเปตฺวา สายนฺหสมยํ อาสยา นิกฺขเมยฺยํ, \csromanfolio{20}อาสยา นิกฺขมิตฺวา วิชมฺเภยฺยํ, วิชมฺภิตฺวา สมนฺตา จตุทฺทิสา อนุวิโลเกยฺยํ, สมนฺตา จตุทฺทิสา อนุวิโลเกตฺวา ติกฺขตฺตุํ สีหนาทํ นเทยฺยํ, ติกฺขตฺตุํ สีหนาทํ นทิตฺวา โคจราย ปกฺกเมยฺยํ.\csromanspacer{} โส วรํ วรํ มิคสํเฆ วธิตฺวา มุทุมํสานิ มุทุมํสานิ ภกฺขยิตฺวา ตเมว อาสยํ อชฺฌุเปยฺยนฺ”ติ.}

\prose{อถ โข โส อาวุโส ชรสิงฺคาโล อญฺญตรํ วนสณฺฑํ นิสฺสาย อาสยํ กปฺเปสิ, ตตฺราสยํ กปฺเปตฺวา สายนฺหสมยํ อาสยา นิกฺขมิ, อาสยา นิกฺขมิตฺวา วิชมฺภิ, วิชมฺภิตฺวา สมนฺตา จตุทฺทิสา อนุวิโลเกสิ, สมนฺตา จตุทฺทิสา อนุวิโลเกตฺวา ติกฺขตฺตุํ “สีหนาทํ นทิสฺสามี”ติ สิงฺคาลกํเยว อนทิ, เภรณฺฑกํเยว\footnote{เภทณฺฑกํเยว (ก)} อนทิ, เก จ ฉเว สิงฺคาเล, เก ปน สีหนาเทติ\footnote{สีหนาเท (?)}.\csromanspacer{} เอวเมว โข ตฺวํ อาวุโส ปาถิกปุตฺต สุคตาปทาเนสุ ชีวมาโน สุคตาติริตฺตานิ ภุญฺชมาโน ตถาคเต อรหนฺเต สมฺมาสมฺพุทฺเธ อาสาเทตพฺพํ มญฺญสิ.\csromanspacer{} เก จ ฉเว ปาถิกปุตฺเต, กา จ ตถาคตานํ อรหนฺตานํ สมฺมาสมฺพุทฺธานํ อาสาทนาติ.}

\proseitem{29}{ยโต โข ภคฺคว ชาลิโย ทารุปตฺติกนฺเตวาสี อิมินา โอปมฺเมน เนว อสกฺขิ อเจลํ ปาถิกปุตฺตํ ตมฺหา อาสนา จาเวตุํ.\csromanspacer{} อถ นํ เอตทโวจ\csromandash{}}

\csromangathameasure{“สีโหติ อตฺตานํ สเมกฺขิยาน, \\ อมญฺญิ โกตฺถุ มิคราชาหมสฺมิ. \\ ตเถว โส สิงฺคาลกํ อนทิ, \\ เก จ ฉเว สิงฺคาเล เก ปน สีหนาเท”ติ.}
\csromangathagroup{\csromangathastanza{“สีโหติ อตฺตานํ สเมกฺขิยาน, \\ อมญฺญิ โกตฺถุ มิคราชาหมสฺมิ. \\ ตเถว\footnote{ตเมว (สฺยา)} โส สิงฺคาลกํ อนทิ, \\ เก จ ฉเว สิงฺคาเล เก ปน สีหนาเท”ติ.}}

\prose{เอวเมว โข ตฺวํ อาวุโส ปาถิกปุตฺต สุคตาปทาเนสุ ชีวมาโน สุคตาติริตฺตานิ ภุญฺชมาโน ตถาคเต อรหนฺเต สมฺมาสมฺพุทฺเธ อาสาเทตพฺพํ มญฺญสิ, เก จ ฉเว ปาถิกปุตฺเต, กา จ ตถาคตานํ อรหนฺตานํ สมฺมาสมฺพุทฺธานํ อาสาทนาติ.}

\proseitem{30}{ยโต โข ภคฺคว ชาลิโย ทารุปตฺติกนฺเตวาสี อิมินาปิ โอปมฺเมน เนว อสกฺขิ อเจลํ ปาถิกปุตฺตํ ตมฺหา อาสนา จาเวตุํ.\csromanspacer{} อถ นํ เอตทโวจ\csromandash{}}

\csromangathameasure{“อญฺญํ อนุจงฺกมนํ, อตฺตานํ วิฆาเส สเมกฺขิย. \\ ยาว อตฺตานํ น ปสฺสติ, โกตฺถุ ตาว พฺยคฺโฆติ มญฺญติ. \\ ตเถว โส สิงฺคาลกํ อนทิ, เก จ ฉเว สิงฺคาเล เก ปน สีหนาเท”ติ.}
\csromangathagroup{\csromangathastanza{\csromanfolio{21}“อญฺญํ อนุจงฺกมนํ, อตฺตานํ วิฆาเส สเมกฺขิย. \\ ยาว อตฺตานํ น ปสฺสติ, โกตฺถุ ตาว พฺยคฺโฆติ มญฺญติ. \\ ตเถว โส สิงฺคาลกํ อนทิ, เก จ ฉเว สิงฺคาเล เก ปน สีหนาเท”ติ.}}

\prose{เอวเมว โข ตฺวํ อาวุโส ปาถิกปุตฺต สุคตาปทาเนสุ ชีวมาโน สุคตาติริตฺตานิ ภุญฺชมาโน ตถาคเต อรหนฺเต สมฺมาสมฺพุทฺเธ อาสาเทตพฺพํ มญฺญสิ “เก จ ฉเว ปาถิกปุตฺเต, กา จ ตถาคตานํ อรหนฺตานํ สมฺมาสมฺพุทฺธานํ อาสาทนา”ติ.}

\proseitem{31}{ยโต โข ภคฺคว ชาลิโย ทารุปตฺติกนฺเตวาสี อิมินาปิ โอปมฺเมน เนว อสกฺขิ อเจลํ ปาถิกปุตฺตํ ตมฺหา อาสนา จาเวตุํ.\csromanspacer{} อถ นํ เอตทโวจ\csromandash{}}

\csromangathameasure{“ภุตฺวาน เภเก ขลมูสิกาโย, \\ กฏสีสุ ขิตฺตานิ จ โกณปานิ. \\ มหาวเน สุญฺญวเน วิวฑฺโฒ, \\ อมญฺญิ โกตฺถุ มิคราชาหมสฺมิ. \\ ตเถว โส สิงฺคาลกํ อนทิ, \\ เก จ ฉเว สิงฺคาเล เก ปน สีหนาเท”ติ.}
\csromangathagroup{\csromangathastanza{“ภุตฺวาน เภเก\footnote{ภิงฺเค (ก)} ขลมูสิกาโย, \\ กฏสีสุ ขิตฺตานิ จ โกณปานิ\footnote{ภิงฺเค (ก)}. \\ มหาวเน สุญฺญวเน วิวฑฺโฒ, \\ อมญฺญิ โกตฺถุ มิคราชาหมสฺมิ. \\ ตเถว โส สิงฺคาลกํ อนทิ, \\ เก จ ฉเว สิงฺคาเล เก ปน สีหนาเท”ติ.}}

\prose{เอวเมว โข ตฺวํ อาวุโส ปาถิกปุตฺต สุคตาปทาเนสุ ชีวมาโน สุคตาติริตฺตานิ ภุญฺชมาโน ตถาคเต อรหนฺเต สมฺมาสมฺพุทฺเธ อาสาเทตพฺพํ มญฺญสิ, เก จ ฉเว ปาถิกปุตฺเต, กา จ ตถาคตานํ อรหนฺตานํ สมฺมาสมฺพุทฺธานํ อาสาทนาติ.}

\proseitem{32}{ยโต โข ภคฺคว ชาลิโย ทารุปตฺติกนฺเตวาสี อิมินาปิ โอปมฺเมน เนว อสกฺขิ อเจลํ ปาถิกปุตฺตํ ตมฺหา อาสนา จาเวตุํ.\csromanspacer{} อถ ตํ ปริสํ อาคนฺตฺวา เอวมาโรเจสิ “ปราภูตรูโป โภ อเจโล ปาถิกปุตฺโต ‘อายามิ อาวุโส อายามิ อาวุโส’ติ วตฺวา ตตฺเถว สํสปฺปติ, น สกฺโกติ อาสนาปิ วุฏฺฐาตุนฺ”ติ.}

\proseitem{33}{\csromanfolio{22}เอวํ วุตฺเต อหํ ภคฺคว ตํ ปริสํ เอตทโวจํ “อภพฺโพ โข อาวุโส อเจโล ปาถิกปุตฺโต ตํ วาจํ อปฺปหาย ตํ จิตฺตํ อปฺปหาย ตํ ทิฏฺฐิํ อปฺปฏินิสฺสชฺชิตฺวา มม สมฺมุขีภาวํ อาคนฺตุํ.\csromanspacer{} สเจปิสฺส เอวมสฺส ‘อหํ ตํ วาจํ อปฺปหาย ตํ จิตฺตํ อปฺปหาย ตํ ทิฏฺฐิํ อปฺปฏินิสฺสชฺชิตฺวา สมณสฺส โคตมสฺส สมฺมุขีภาวํ คจฺเฉยฺยนฺ’ติ, มุทฺธาปิ ตสฺส วิเปเตยฺย.\csromanspacer{} สเจปายสฺมนฺตานํ ลิจฺฉวีนํ เอวมสฺส ‘มยํ อเจลํ ปาถิกปุตฺตํ วรตฺตาหิ พนฺธิตฺวา โคยุเคหิ อาวิญฺเฉยฺยามา’ติ.\csromanspacer{} ตา วรตฺตา ฉิชฺเชยฺยุํ ปาถิกปุตฺโต วา.\csromanspacer{} อภพฺโพ ปน อเจโล ปาถิกปุตฺโต ตํ วาจํ อปฺปหาย ตํ จิตฺตํ อปฺปหาย ตํ ทิฏฺฐิํ อปฺปฏินิสฺสชฺชิตฺวา มม สมฺมุขีภาวํ อาคนฺตุํ.\csromanspacer{} สเจปิสฺส เอวมสฺส ‘อหํ ตํ วาจํ อปฺปหาย ตํ จิตฺตํ อปฺปหาย ตํ ทิฏฺฐิํ อปฺปฏินิสฺสชฺชิตฺวา สมณสฺส โคตมสฺส สมฺมุขีภาวํ คจฺเฉยฺยนฺ’ติ, มุทฺธาปิ ตสฺส วิปเตยฺยา”ติ.}

\proseitem{34}{อถ ขฺวาหํ ภคฺคว ตํ ปริสํ ธมฺมิยา กถาย สนฺทสฺเสสิํ สมาทเปสิํ สมุตฺเตเชสิํ สมฺปหํเสสิํ, ตํ ปริสํ ธมฺมิยา กถาย สนฺทสฺเสตฺวา สมาทเปตฺวา สมุตฺเตเชตฺวา สมฺปหํเสตฺวา มหาพนฺธนา โมกฺขํ กริตฺวา จตุราสีติปาณสหสฺสานิ มหาวิทุคฺคา อุทฺธริตฺวา เตโชธาตุํ สมาปชฺชิตฺวา สตฺตตาลํ เวหาสํ อพฺภุคฺคนฺตฺวา อญฺญํ สตฺตตาลมฺปิ อจฺจิํ\footnote{อคฺคิํ (สฺยา)} อภินิมฺมินิตฺวา ปชฺชลิตฺวา ธูมายิตฺวา\footnote{ธูปายิตฺวา (สี, อิ)} มหาวเน กูฏคารสาลายํ ปจฺจุฏฺฐาสิํ.}

\proseitem{35}{อถ โข ภคฺคว สุนกฺขตฺโต ลิจฺฉวิปุตฺโต เยนาหํ เตนุปสงฺกมิ, อุปสงฺกมิตฺวา มํ อภิวาเทตฺวา เอกมนฺตํ นิสีทิ.\csromanspacer{} เอกมนฺตํ นิสินฺนํ โข อหํ ภคฺคว สุนกฺขตฺตํ ลิจฺฉวิปุตฺตํ เอตทโวจํ “ตํ กิํมญฺญสิ สุนกฺขตฺต, ยเถว เต อหํ อเจลํ ปาถิกปุตฺตํ อารพฺภ พฺยากาสิํ, ตเถว ตํ วิปากํ อญฺญถา วา”ติ.\csromanspacer{} ยเถว เม ภนฺเต ภควา อเจลํ ปาถิกปุตฺตํ อารพฺภ พฺยากาสิ, ตเถว ตํ วิปากํ โน อญฺญถาติ.}

\prose{ตํ กิํ มญฺญสิ สุนกฺขตฺต, ยทิ เอวํ สนฺเต กตํ วา โหติ อุตฺตริ มนุสฺสธมฺมา อิทฺธิปาฏิหาริยํ อกตํ วาติ.\csromanspacer{} อทฺธา โข ภนฺเต เอวํ สนฺเต กตํ โหติ อุตฺตริ มนุสฺสธมฺมา อิทฺธิปาฏิหาริยํ, โน อกตนฺติ.\csromanspacer{} เอวมฺปิ \csromanfolio{23}โข มํ ตฺวํ โมฆปุริส อุตฺตริ มนุสฺสธมฺมา อิทฺธิปาฏิหาริยํ กโรนฺตํ เอวํ วเทสิ “น หิ ปน เม ภนฺเต ภควา อุตฺตริ มนุสฺสธมฺมา อิทฺธิปาฏิหาริยํ กโรตี”ติ.\csromanspacer{} ปสฺส โมฆปุริส ยาวญฺจ เต อิทํ อปรทฺธนฺติ.\csromanspacer{} เอวมฺปิ โข ภคฺคว สุนกฺขตฺโต ลิจฺฉวิปุตฺโต มยา วุจฺจมาโน อปกฺกเมว อิมสฺมา ธมฺมวินยา, ยถา ตํ อาปายิโก เนรยิโก.}

\csromanmatikaanchor{mtk.7}
\csromantocmarkref{section}{อคฺคญฺญปญฺญตฺติกถา}{mtk.7}
\bookmark[dest={mtk.7},level=1]{อคฺคญฺญปญฺญตฺติกถา}
\csromanheader{1}{\textbf{อคฺคญฺญปญฺญตฺติกถา}}
\proseitem{36}{อคฺคญฺญญฺจาหํ ภคฺคว ปชานามิ.\csromanspacer{} ตญฺจ ปชานามิ\footnote{“ตญฺจ ปชานามี”ติ อิทํ สฺยา-โปตฺถเก นตฺถิ.}, ตโต จ อุตฺตริตรํ ปชานามิ, ตญฺจ ปชานนํ น ปรามสามิ, อปรามสโต จ เม ปจฺจตฺตญฺเญว นิพฺพุติ วิทิตา, ยทภิชานํ ตถาคโต โน อนยํ อาปชฺชติ.}

\proseitem{37}{สนฺติ ภคฺคว เอเก สมณพฺราหฺมณา อิสฺสรกุตฺตํ พฺรหฺมกุตฺตํ อาจริยกํ อคฺคญฺญํ ปญฺญเปนฺติ.\csromanspacer{} ตฺยาหํ อุปสงฺกมิตฺวา เอวํ วทามิ “สจฺจํ กิร ตุเมฺห อายสฺมนฺโต อิสฺสรกุตฺตํ พฺรหฺมกุตฺตํ อาจริยกํ อคฺคญฺญํ ปญฺญเปถา”ติ.\csromanspacer{} เต จ เม เอวํ ปุฏฺฐา “อาโม”ติ\footnote{อามาติ (สฺยา)} ปฏิชานนฺติ.\csromanspacer{} ตฺยาหํ เอวํ วทามิ “กถํวิหิตกํ ปน\footnote{กถํ วิหิตกํ โน ปน (ก)} ตุเมฺห อายสฺมนฺโต อิสฺสรกุตฺตํ พฺรหฺมกุตฺตํ อาจริยกํ อคฺคญฺญํ ปญฺญเปถา”ติ.\csromanspacer{} เต มยา ปุฏฺฐา น สมฺปายนฺติ, อสมฺปายนฺตา มมํเยว ปฏิปุจฺฉนฺติ.\csromanspacer{} เตสาหํ ปุฏฺโฐ พฺยากโรมิ\csromandash{}}

\proseitem{38}{โหติ โข โส อาวุโส สมโย ยํ กทาจิ กรหจิ ทีฆสฺส อทฺธุโน อจฺจเยน อยํ โลโก สํวฏฺฏติ, สํวฏฺฏมาเน โลเก เยภุยฺเยน สตฺตา อาภสฺสรสํวตฺตนิกา โหนฺติ, เต ตตฺถ โหนฺติ มโนมยา ปีติภกฺขา สยํปภา อนฺตลิกฺขจรา สุภฏฺฐายิโน จิรํ ทีฆมทฺธานํ ติฏฺฐนฺติ.}

\prose{โหติ โข โส อาวุโส สมโย ยํ กทาจิ กรหจิ ทีฆสฺส อทฺธุโน อจฺจเยน อยํ โลโก วิวฏฺฏติ, วิวฏฺฏมาเน โลเก สุญฺญํ พฺรหฺมวิมานํ ปาตุภวติ.\csromanspacer{} อถ โข\footnote{อถ (สี, สฺยา, อิ)} อญฺญตโร สตฺโต อายุกฺขยา วา ปุญฺญกฺขยา วา อาภสฺสรกายา จวิตฺวา สุญฺญํ พฺรหฺมวิมานํ อุปปชฺชติ, \csromanfolio{24}โส ตตฺถ โหติ มโนมโย ปีติภกฺโข สยํปโภ อนฺตลิกฺขจโร สุภฏฺฐายี, จิรํ ทีฆมทฺธานํ ติฏฺฐติ.}

\prose{ตสฺส ตตฺถ เอกกสฺส ทีฆรตฺตํ นิวุสิตตฺตา อนภิรติ ปริตสฺสนา อุปฺปชฺชติ “อโห วต อญฺเญปิ สตฺตา อิตฺถตฺตํ อาคจฺเฉยฺยุนฺ”ติ.\csromanspacer{} อถ อญฺเญปิ สตฺตา อายุกฺขยา วา ปุญฺญกฺขยา วา อาภสฺสรกายา จวิตฺวา สุญฺญํ พฺรหฺมวิมานํ อุปปชฺชนฺติ ตสฺส สตฺตสฺส สหพฺยตํ, เตปิ ตตฺถ โหนฺติ มโนมยา ปีติภกฺขา สยํปภา อนฺตลิกฺขจรา สุภฏฺฐายิโน, จิรํ ทีฆมทฺธานํ ติฏฺฐนฺติ.}

\proseitem{39}{ตตฺราวุโส โย โส สตฺโต ปฐมํ อุปปนฺโน, ตสฺส เอวํ โหติ “อหมสฺมิ พฺรหฺมา มหาพฺรหฺมา อภิภู อนภิภูโต อญฺญทตฺถุทโส วสวตฺตี อิสฺสโร กตฺตา นิมฺมาตา เสฏฺโฐ สชิตา\footnote{สญฺชิตา (สี, อิ), สชฺชิตา (สฺยา, กํ)} วสี ปิตา ภูตภพฺยานํ, มยา อิเม สตฺตา นิมฺมิตา.\csromanspacer{} ตํ กิสฺส เหตุ, มมํ หิ ปุพฺเพ เอตทโหสิ ‘อโห วต อญฺเญปิ สตฺตา อิตฺถตฺตํ อาคจฺเฉยฺยุนฺ’ติ, อิติ มม จ มโนปณิธิ, อิเม จ สตฺตา อิตฺถตฺตํ อาคตา”ติ.}

\prose{เยปิ เต สตฺตา ปจฺฉา อุปปนฺนา, เตสมฺปิ เอวํ โหติ “อยํ โข ภวํ พฺรหฺมา มหาพฺรหฺมา อภิภู อนภิภูโต อญฺญทตฺถุทโส วสวตฺตี อิสฺสโร กตฺตา นิมฺมาตา เสฏฺโฐ สชิตา วสี ปิตา ภูตภพฺยานํ, อิมินา มยํ โภตา พฺรหฺมุนา นิมฺมิตา.\csromanspacer{} ตํ กิสฺส เหตุ, อิมํ หิ มยํ อทฺทสาม อิธ ปฐมํ อุปปนฺนํ, มยํ ปนามฺห ปจฺฉา อุปปนฺนา”ติ.}

\proseitem{40}{ตตฺราวุโส โย โส สตฺโต ปฐมํ อุปปนฺโน, โส ทีฆายุกตโร จ โหติ วณฺณวนฺตตโร จ มเหสกฺขตโร จ.\csromanspacer{} เย ปน เต สตฺตา ปจฺฉา อุปปนฺนา, เต อปฺปายุกตรา จ โหนฺติ ทุพฺพณฺณตรา จ อปฺเปสกฺขตรา จ.}

\prose{ฐานํ โข ปเนตํ อาวุโส วิชฺชติ, ยํ อญฺญตโร สตฺโต ตมฺหา กายา จวิตฺวา อิตฺถตฺตํ อาคจฺฉติ.\csromanspacer{} อิตฺถตฺตํ อาคโต สมาโน อคารสฺมา อนคาริยํ ปพฺพชติ, อคารสฺมา อนคาริยํ ปพฺพชิโต \csromanfolio{25}สมาโน อาตปฺปมนฺวาย ปธานมนฺวาย อนุโยคมนฺวาย อปฺปมาทมนฺวาย สมฺมามนสิการมนฺวาย ตถารูปํ เจโตสมาธิํ ผุสติ, ยถาสมาหิเต จิตฺเต ตํ ปุพฺเพนิวาสํ อนุสฺสรติ, ตโต ปรํ นานุสฺสรติ.}

\prose{โส เอวมาห “โย โข โส ภวํ พฺรหฺมา มหาพฺรหฺมา อภิภู อนภิภูโต อญฺญทตฺถุทโส วสวตฺตี อิสฺสโร กตฺตา นิมฺมาตา เสฏฺโฐ สชิตา วสี ปิตา ภูตภพฺยานํ, เยน มยํ โภตา พฺรหฺมุนา นิมฺมิตา, โส นิจฺโจ ธุโว สสฺสโต\footnote{สสฺสโต ทีฆายุโก (สฺยา, ก)} อวิปริณามธมฺโม สสฺสติสมํ ตเถว ฐสฺสติ.\csromanspacer{} เย ปน มยํ อหุมฺหา เตน โภตา พฺรหฺมุนา นิมฺมิตา, เต มยํ อนิจฺจา อทฺธุวา\footnote{อทฺธุวา อสสฺสตา (สฺยา, ก)} อปฺปายุกา จวนธมฺมา อิตฺถตฺตํ อาคตา”ติ.\csromanspacer{} เอวํวิหิตกํ โน ตุเมฺห อายสฺมนฺโต อิสฺสรกุตฺตํ พฺรหฺมกุตฺตํ อาจริยกํ อคฺคญฺญํ ปญฺญเปถาติ.\csromanspacer{} เต เอวมาหํสุ “เอวํ โข โน อาวุโส โคตม สุตํ, ยเถวายสฺมา โคตโม อาหา”ติ.\csromanspacer{} อคฺคญฺญญฺจาหํ ภคฺคว ปชานามิ.\csromanspacer{} ตญฺจ ปชานามิ, ตโต จ อุตฺตริตรํ ปชานามิ, ตญฺจ ปชานนํ น ปรามสามิ, อปรามสโต จ เม ปจฺจตฺตญฺเญว นิพฺพุติ วิทิตา, ยทภิชานํ ตถาคโต โน อนยํ อาปชฺชติ.}

\proseitem{41}{สนฺติ ภคฺคว เอเก สมณพฺราหฺมณา ขิฑฺฑาปโทสิกํ อาจริยกํ อคฺคญฺญํ ปญฺญเปนฺติ, ตฺยาหํ อุปสงฺกมิตฺวา เอวํ วทามิ “สจฺจํ กิร ตุเมฺห อายสฺมนฺโต ขิฑฺฑาปโทสิกํ อาจริยกํ อคฺคญฺญํ ปญฺญเปถา”ติ.\csromanspacer{} เต จ เม เอวํ ปุฏฺฐา “อาโม”ติ ปฏิชานนฺติ.\csromanspacer{} ตฺยาหํ เอวํ วทามิ “กถํวิหิตกํ ปน ตุเมฺห อายสฺมนฺโต ขิฑฺฑาปโทสิกํ อาจริยกํ อคฺคญฺญํ ปญฺญเปถา”ติ.\csromanspacer{} เต มยา ปุฏฺฐา น สมฺปายนฺติ, อสมฺปายนฺตา, มมํเยว ปฏิปุจฺฉนฺติ, เตสาหํ ปุฏฺโฐ พฺยากโรมิ\csromandash{}}

\proseitem{42}{สนฺตาวุโส ขิฑฺฑาปโทสิกา นาม เทวา, เต อติเวลํ หสฺสาขิฑฺฑารติธมฺมสมาปนฺนา\footnote{หสขิฑฺฑารติธมฺมสมาปนฺนา (ก)} วิหรนฺติ, เตสํ อติเวลํ หสฺสขิฑฺฑารติธมฺมสมาปนฺนานํ วิหรตํ สติ สมฺมุสฺสติ, สติยา สมฺโมสา\footnote{สติยา สมฺโมสาย (สฺยา)} เต เทวา ตมฺหา กายา จวนฺติ.}

\prose{\csromanfolio{26}ฐานํ โข ปเนตํ อาวุโส วิชฺชติ, ยํ อญฺญตโร สตฺโต ตมฺหา กายา จวิตฺวา อิตฺถตฺตํ อาคจฺฉติ, อิตฺถตฺตํ อาคโต สมาโน อคารสฺมา อนคาริยํ ปพฺพชติ, อคารสฺมา อนคาริยํ ปพฺพชิโต สมาโน อาตปฺปมนฺวาย ปธานมนฺวาย อนุโยคมนฺวาย อปฺปมาทมนฺวาย สมฺมามนสิการมนฺวาย ตถารูปํ เจโตสมาธิํ ผุสติ, ยถาสมาหิเต จิตฺเต ตํ ปุพฺเพนิวาสํ อนุสฺสรติ, ตโต ปรํ นานุสฺสรติ.}

\prose{โส เอวมาห “เย โข เต โภนฺโต เทวา น ขิฑฺฑาปโทสิกา, เต น อติเวลํ หสฺสขิฑฺฑารติธมฺมสมาปนฺนา วิหรนฺติ, เตสํ นาติเวลํ หสฺสขิฑฺฑารติธมฺมสมาปนฺนานํ วิหรตํ สติ น สมฺมุสฺสติ, สติยา อสมฺโมสา เต เทวา ตมฺหา กายา น จวนฺติ, นิจฺจา ธุวา สสฺสตา อวิปริณามธมฺมา สสฺสติสมํ ตเถว ฐสฺสนฺติ.\csromanspacer{} เย ปน มยํ อหุมฺหา ขิฑฺฑาปโทสิกา, เต มยํ อติเวลํ หสฺสขิฑฺฑารติธมฺมสมฺปนฺนา วิหริมฺหา, เตสํ โน อติเวลํ หสฺสขิฑฺฑารติธมฺมสมาปนฺนานํ วิหรตํ สติ สมฺมุสฺสติ, สติยา สมฺโมสา เต\footnote{สมฺโมสา เอว (สี, อิ)} มยํ ตมฺหา กายา จุตา, อนิจฺจา อทฺธุวา อปฺปายุกา จวนธมฺมา อิตฺถตฺตํ อาคตา”ติ.\csromanspacer{} เอวํวิหิตกํ โน ตุเมฺห อายสฺมนฺโต ขิฑฺฑาปโทสิกํ อาจริยกํ อคฺคญฺญํ ปญฺญเปถาติ.\csromanspacer{} เต เอวมาหํสุ “เอวํ โข โน อาวุโส โคตม สุตํ, ยเถวายสฺมา โคตโม อาหา”ติ.\csromanspacer{} อคฺคญฺญญฺจาหํ ภคฺคว ปชานามิ ฯเปฯ ยทภิชานํ ตถาคโต โน อนยํ อาปชฺชติ.}

\proseitem{43}{สนฺติ ภคฺคว เอเก สมณพฺราหฺมณา มโนปโทสิกํ อาจริยกํ อคฺคญฺญํ ปญฺญเปนฺติ, ตฺยาหํ อุปสงฺกมิตฺวา เอวํ วทามิ “สจฺจํ กิร ตุเมฺห อายสฺมนฺโต มโนปโทสิกํ อาจริยกํ อคฺคญฺญํ ปญฺญเปถา”ติ.\csromanspacer{} เต จ เม เอวํ ปุฏฺฐา “อาโม”ติ ปฏิชานนฺติ.\csromanspacer{} ตฺยาหํ เอวํ วทามิ “กถํวิหิตกํ ปน ตุเมฺห อายสฺมนฺโต มโนปโทสิกํ อาจริยกํ อคฺคญฺญํ ปญฺญเปถา”ติ.\csromanspacer{} เต มยา ปุฏฺฐา น สมฺปายนฺติ, อสมฺปายนฺตา มมํเยว ปฏิปุจฺฉนฺติ.\csromanspacer{} เตสาหํ ปุฏฺโฐ พฺยากโรมิ\csromandash{}}

\proseitem{44}{สนฺตาวุโส มโนปโทสิกา นาม เทวา, เต อติเวลํ อญฺญมญฺญํ อุปนิชฺฌายนฺติ, เต อติเวลํ อญฺญมญฺญํ อุปนิชฺฌายนฺตา อญฺญมญฺญมฺหิ \csromanfolio{27}จิตฺตานิ ปทูเสนฺติ, เต อญฺญมญฺญํ ปทุฏฺฐจิตฺตา กิลนฺตกายา กิลนฺตจิตฺตา เต เทวา ตมฺหา กายา จวนฺติ.}

\prose{ฐานํ โข ปเนตํ อาวุโส วิชฺชติ.\csromanspacer{} ยํ อญฺญตโร สตฺโต ตมฺหา กายา จวิตฺวา อิตฺถตฺตํ อาคจฺฉติ, อิตฺถตฺตํ อาคโต สมาโน อคารสฺมา อนคาริยํ ปพฺพชติ, อคารสฺมา อนคาริยํ ปพฺพชิโต สมาโน อาตปฺปมนฺวาย ปธานมนฺวาย อนุโยคมนฺวาย อปฺปมาทมนฺวาย สมฺมามนสิการมนฺวาย ตถารุปํ เจโตสมาธิํ ผุสติ, ยถาสมาหิเต จิตฺเต ตํ ปุพฺเพนิวาสํ อนุสฺสรติ, ตโต ปรํ นานุสฺสรติ.}

\prose{โส เอวมาห “เย โข เต โภนฺโต เทวา น มโนปโทสิกา เต นาติเวลํ อญฺญมญฺญํ อุปนิชฺฌายนฺติ, เต นาติเวลํ อญฺญมญฺญํ อุปนิชฺฌายนฺตา อญฺญมญฺญมฺหิ จิตฺตานิ นปฺปทูเสนฺติ, เต อญฺญมญฺญํ อปฺปทุฏฺฐจิตฺตา อกิลนฺตกายา อกิลนฺตจิตฺตา\footnote{อกิลนฺตจิตฺตา เต เทวา (สี, สฺยา, อิ)} ตมฺหา กายา น จวนฺติ, นิจฺจา ธุวา สสฺสตา อวิปริณามธมฺมา สสฺสติสมํ ตเถว ฐสฺสนฺติ.\csromanspacer{} เย ปน มยํ อหุมฺหา มโนปโทสิกา, เต มยํ อติเวลํ อญฺญมญฺญํ อุปนิชฺฌายิมฺหา, เต มยํ อติเวลํ อญฺญมญฺญํ อุปนิชฺฌายนฺตา อญฺญมญฺญมฺหิ จิตฺตานิ ปทูสิมฺหา\footnote{ปโทสิยิมฺหา (สฺยา), ปทูสยิมฺหา (?)}, เต มยํ อญฺญมญฺญํ ปทุฏฺฐจิตฺตา กิลนฺตกายา กิลนฺตจิตฺตา\footnote{กิลนฺตจิตฺตา-เอว มยํ (สี, อิ), กิลนฺตจิตฺตา เอวํ มยํ (สฺยา)} ตมฺหา กายา จุตา, อนิจฺจา อทฺธุวา อสสฺสตา อปฺปายุกา จวนธมฺมา อิตฺถตฺตํ อาคตา”ติ.\csromanspacer{} เอวํวิหิตกํ โน ตุเมฺห อายสฺมนฺโต มโนปโทสิกํ อาจริยกํ อคฺคญฺญํ ปญฺญเปถาติ.\csromanspacer{} เต เอวมาหํสุ “เอวํ โข โน อาวุโส โคตม สุตํ, ยเถวายสฺมา โคตโม อาหา”ติ.\csromanspacer{} อคฺคญฺญญฺจาหํ ภคฺคว ปชานามิ ฯเปฯ ยทภิชานํ ตถาคโต โน อนยํ อาปชฺชติ.}

\proseitem{45}{สนฺติ ภคฺคว เอเก สมณพฺราหฺมณา อธิจฺจสมุปฺปนฺนํ อาจริยกํ อคฺคญฺญํ ปญฺญเปนฺติ, ตฺยาหํ อุปสงฺกมิตฺวา เอวํ วทามิ “สจฺจํ กิร ตุเมฺห อายสฺมนฺโต อธิจฺจสมุปฺปนฺนํ อาจริยกํ อคฺคญฺญํ ปญฺญเปถา”ติ.\csromanspacer{} เต จ เม เอวํ ปุฏฺฐา “อาโม”ติ ปฏิชานนฺติ.\csromanspacer{} ตฺยาหํ เอวํ วทามิ “กถํวิหิตกํ ปน ตุเมฺห อายสฺมนฺโต อธิจฺจสมุปฺปนฺนํ อาจริยกํ อคฺคญฺญํ ปญฺญเปถา”ติ. \csromanfolio{28}เต จ มยา ปุฏฺฐา น สมฺปายนฺติ, อสมฺปายนฺตา มมํเยว ปฏิปุจฺฉนฺติ, เตสาหํ ปุฏฺโฐ พฺยากโรมิ\csromandash{}}

\proseitem{46}{สนฺตาวุโส อสญฺญสตฺตา นาม เทวา, สญฺญุปฺปาทา จ ปน เต เทวา ตมฺหา กยา จวนฺติ.}

\prose{ฐานํ โข ปเนตํ อาวุโส วิชฺชติ.\csromanspacer{} ยํ อญฺญตโร สตฺโต ตมฺหา กายา จวิตฺวา อิตฺถตฺตํ อาคจฺฉติ, อิตฺถตฺตํ อาคโต สมาโน อคารสฺมา อนคาริยํ ปพฺพชติ, อคารสฺมา อนคาริยํ ปพฺพชิโต สมาโน อาตปฺปมนฺวาย ปธานมนฺวาย อนุโยคมนฺวาย อปฺปมาทมนฺวาย สมฺมามนสิการมนฺวาย ตถารูปํ เจโตสมาธิํ ผุสติ, ยถาสมาหิเต จิตฺเต ตํ\footnote{อิทํ ปทํ พฺรหฺมชาลสุตฺเต น ทิสฺสติ. เอวํ (อิ, ก) 2. สตฺตตาย (สี, อิ), สตฺตาย (ก-สี)} สญฺญุปฺปาทํ อนุสฺสรติ, ตโต ปรํ นานุสฺสรติ.}

\prose{โส เอวมาห “อธิจฺจสมุปฺปนฺโน อตฺตา จ โลโก จ.\csromanspacer{} ตํ กิสฺส เหตุ, อหํ หิ ปุพฺเพ นาโหสิํ, โสมฺหิ เอตรหิ อหุตฺวา สนฺตตาย2 ปริณโต”ติ.\csromanspacer{} เอวํวิหิตกํ โน ตุเมฺห อายสฺมนฺโต อธิจฺจสมุปฺปนฺนํ อาจริยกํ อคฺคญฺญํ ปญฺญเปถาติ.\csromanspacer{} เต เอวมาหํสุ “เอวํ โข โน อาวุโส โคตม สุตํ ยเถวายสฺมา โคตโม อาหา”ติ.\csromanspacer{} อคฺคญฺญญฺจาหํ ภคฺคว ปชานามิ.\csromanspacer{} ตญฺจ ปชานามิ, ตโต จ อุตฺตริตรํ ปชานามิ, ตญฺจ ปชานํ น ปรามสามิ, อปรามสโต จ เม ปจฺจตฺตญฺเญว นิพฺพุติ วิทิตา, ยทภิชานํ ตถาคโต โน อนยํ อาปชฺชติ.}

\proseitem{47}{เอวํวาทิํ โข มํ ภคฺคว เอวมกฺขายิํ เอเก สมณพฺราหฺมณา อสตา ตุจฺฉา มุสา อภูเตน อพฺภาจิกฺขนฺติ “วิปรีโต สมโณ โคตโม ภิกฺขโว จ, สมโณ โคตโม เอวมาห ‘ยสฺมิํ สมเย สุภํ วิโมกฺขํ อุปสมฺปชฺช วิหรติ, สพฺพํ ตสฺมิํ สมเย อสุภนฺเตฺวว\footnote{อสุภนฺเตว (สี, สฺยา, อิ)} ปชานาตี’ติ\footnote{สญฺชานาตีติ (สี, อิ)}”.\csromanspacer{} น โข ปนาหํ ภคฺคว เอวํ วทามิ “ยสฺมิํ สมเย สุภํ วิโมกฺขํ อุปสมฺปชฺช วิหรติ, สพฺพํ ตสฺมิํ สมเย อสุภนฺเตฺวว ปชานาตี”ติ.\csromanspacer{} เอวญฺจ ขฺวาหํ ภคฺคว วทามิ “ยสฺมิํ สมเย สุภํ วิโมกฺขํ อุปสมฺปชฺช วิหรติ, สุภนฺเตฺวว ตสฺมิํ สมเย ปชานาตี”ติ.}

\prose{\csromanfolio{29}เต จ ภนฺเต วิปรีตา, เย ภควนฺตํ วิปรีตโต ทหนฺติ ภิกฺขโว จ.\csromanspacer{} เอวํปสนฺโน อหํ ภนฺเต ภควติ “ปโหติ เม ภควา ตถา ธมฺมํ เทเสตุํ, ยถา อหํ สุภํ วิโมกฺขํ อุปสมฺปชฺช วิหเรยฺยนฺ”ติ.}

\proseitem{48}{ทุกฺกรํ โข เอตํ ภคฺคว ตยา อญฺญทิฏฺฐิเกน อญฺญขนฺติเกน อญฺญรุจิเกน อญฺญตฺราโยเคน อญฺญตฺราจริยเกน สุภํ วิโมกฺขํ อุปสมฺปชฺช วิหริตุํ.\csromanspacer{} อิงฺฆ ตฺวํ ภคฺคว โย จ เต อยํ มยิ ปสาโท, ตเมว ตฺวํ สาธุกมนุรกฺขาติ.\csromanspacer{} สเจตํ ภนฺเต มยา ทุกฺกรํ อญฺญทิฏฺฐิเกน อญฺญขนฺติเกน อญฺญรุจิเกน อญฺญตฺราโยเคน อญฺญตฺราจริยเกน สุภํ วิโมกฺขํ อุปสมฺปชฺช วิหริตุํ.\csromanspacer{} โย จ เม อยํ ภนฺเต ภควติ ปสาโท, ตเมวาหํ สาธุกมนุรกฺขิสฺสามีติ.\csromanspacer{} อิทมโวจ ภควา.\csromanspacer{} อตฺตมโน ภคฺควโคตฺโต ปริพฺพาชโก ภควโต ภาสิตํ อภินนฺทีติ.}

\nitthitam{\textbf{ปาถิกสุตฺตํ}\footnote{ปาฏิกสุตฺตนฺตํ (สี, สฺยา, กํ, อิ)} \textbf{นิฏฺฐิตํ ปฐมํ.}}
\csromanmatikaanchor{mtk.8}
\csromantocmarknopage{chapter}{2. อุทุมฺพริกสุตฺต}{mtk.8}
\bookmark[dest={mtk.8},level=0]{2. อุทุมฺพริกสุตฺต}
\chapterheadpageread{2. อุทุมฺพริกสุตฺต}
\csromanmatikaanchor{mtk.9}
\csromantocmarkref{section}{นิโคฺรธปริพฺพาชกวตฺถุ}{mtk.9}
\bookmark[dest={mtk.9},level=1]{นิโคฺรธปริพฺพาชกวตฺถุ}
\csromanheader{1}{\textbf{นิโคฺรธปริพฺพาชกวตฺถุ}}
\proseitem{49}{\csromanfolio{30}เอวํ เม สุตํ\csromandash{}เอกํ สมยํ ภควา ราชคเห วิหรติ คิชฺฌกูเฏ ปพฺพเต.\csromanspacer{} เตน โข ปน สมเยน นิโคฺรโธ ปริพฺพาชโก อุทุมฺพริกาย ปริพฺพาชการาเม ปฏิวสติ มหติยา ปริพฺพาชกปริสาย สทฺธิํ ติํสมตฺเตหิ ปริพฺพาชกสเตหิ.\csromanspacer{} อถ โข สนฺธาโน คหปติ ทิวา ทิวสฺส\footnote{ทิวาทิวสฺเสว (สี, สฺยา, อิ)} ราชคหา นิกฺขมิ ภควนฺตํ ทสฺสนาย.\csromanspacer{} อถ โข สนฺธานสฺส คหปติสฺส เอตทโหสิ “อกาโล โข ภควนฺตํ ทสฺสนาย, ปฏิสลฺลีโน ภควา, มโนภาวนียานํปิ ภิกฺขูนํ อสมโย ทสฺสนาย, ปฏิสลฺลีนา มโนภาวนียา ภิกฺขู, ยํนูนาหํ เยน อุทุมฺพริกาย ปริพฺพาชการาโม, เยน นิโคฺรโธ ปริพฺพาชโก เตนุปสงฺกเมยฺยนฺติ.\csromanspacer{} อถ โข สนฺธาโน คหปติ เยน อุทุมฺพริกาย ปริพฺพาชการาโม, เยน นิโคฺรโธ ปริพฺพาชโก เตนุปสงฺกมิ.}

\proseitem{50}{เตน โข ปน สมเยน นิโคฺรโธ ปริพฺพาชโก มหติยา ปริพฺพาชกปริสาย สทฺธิํ นิสินฺโน โหติ อุนฺนาทินิยา อุจฺจาสทฺทมหาสทฺทาย อเนกวิหิตํ ติรจฺฉานกถํ กเถนฺติยา.\csromanspacer{} เสยฺยถิทํ, ราชกถํ โจรกถํ มหามตฺตกถํ เสนากถํ ภยกถํ ยุทฺธกถํ อนฺนกถํ ปานกถํ วตฺถกถํ สยนกถํ มาลากถํ คนฺธกถํ ญาติกถํ ยานกถํ คามกถํ นิคมกถํ นครกถํ ชนปทกถํ อิตฺถิกถํ สูรกถํ วิสิขากถํ กุมฺภฏฺฐานกถํ ปุพฺพเปตกถํ นานตฺตกถํ โลกกฺขายิกํ สมุทฺทกฺขายิกํ อิติภวาภวกถํ อิติ วา.}

\proseitem{51}{อทฺทสา โข นิโคฺรโธ ปริพฺพาชโก สนฺธานํ คหปติํ ทูรโตว อาคจฺฉนฺตํ, ทิสฺวา สกํ ปริสํ สณฺฐาเปสิ “อปฺปสทฺทา โภนฺโต โหนฺตุ, มา โภนฺโต สทฺทมกตฺถ.\csromanspacer{} อยํ สมณสฺส โคตมสฺส สาวโก อาคจฺฉติ สนฺธาโน คหปติ.\csromanspacer{} ยาวตา โข ปน สมณสฺส โคตมสฺส สาวกา คิหี โอทาตวสนา ราชคเห ปฏิวสนฺติ, อยํ เตสํ อญฺญตโร สนฺธาโน คหปติ, อปฺปสทฺทกามา โข ปเนเต อายสฺมานฺโต}

\vspace{\tipitakaparskip}
\csromancenter{อุทุมฺพริกสุตฺต อปฺปสทฺทวินีตา, อปฺปสทฺทสฺส วณฺณวาทิโน, อปฺเปว นาม อปฺปสทฺทํ ปริสํ วิทิตฺวา อุปสงฺกมิตพฺพํ มญฺเญยฺยา”ติ.\csromanspacer{} เอวํ วุตฺเต เต ปริพฺพาชกา ตุณฺหี อเหสุํ.}
\proseitem{52}{\csromanfolio{31}อถ โข สนฺธาโน คหปติ เยน นิโคฺรโธ ปริพฺพาชโก เตนุปสงฺกมิ, อุปสงฺกมิตฺวา นิโคฺรเธน ปริพฺพาชเกน สทฺธิํ สมฺโมทิ, สมฺโมทนียํ กถํ สารณียํ วีติสาเรตฺวา เอกมนฺตํ นิสีทิ.\csromanspacer{} เอกมนฺตํ นิสินฺโน โข สนฺธาโน คหปติ นิโคฺรธํ ปริพฺพาชกํ เอตทโวจ “อญฺญถา โข อิเม โภนฺโต อญฺญติตฺถิยา ปริพฺพาชกา สงฺคมฺม สมาคมฺม อุนฺนาทิโน อุจฺจาสทฺทมหาสทฺทา อเนกวิหิตํ ติรจฺฉานกถํ อนุยุตฺตา วิหรนฺติ.\csromanspacer{} เสยฺยถิทํ, ราชกถํ ฯเปฯ อิติภวาภวกถํ อิติ วา.\csromanspacer{} อญฺญถา โข\footnote{จ (สี, อิ)} ปน โส ภควา อรญฺญวนปตฺถานิ ปนฺตานิ เสนาสนานิ ปฏิเสวติ อปฺปสทฺทานิ อปฺปนิคฺโฆสานิ วิชนวาตานิ มนุสฺสราหสฺเสยฺยกานิ ปฏิสลฺลานสารุปฺปานี”ติ.}

\proseitem{53}{เอวํ วุตฺเต นิโคฺรโธ ปริพฺพาชโก สนฺธานํ คหปติํ เอตทโวจ “ยคฺเฆ คหปติ ชาเนยฺยาสิ, เกน สมโณ โคตโม สทฺธิํ สลฺลปติ, เกน สากจฺฉํ สมาปชฺชติ, เกน ปญฺญาเวยฺยตฺติยํ อาปชฺชติ.\csromanspacer{} สุญฺญาคารหตา สมณสฺส โคตมสฺส ปญฺญา, อปริสาวจโร สมโณ โคตโม, นาลํ สลฺลาปาย.\csromanspacer{} โส อนฺตมนฺตาเนว\footnote{อนฺตปนฺตาเนว (สฺยา)} เสวติ.\csromanspacer{} เสยฺยถาปิ นาม โคกาณา ปริยนฺตจารินี อนฺตมนฺตาเนว เสวติ.\csromanspacer{} เอวเมว สุญฺญาคารหตา สมณสฺส โคตมสฺส ปญฺญา, อปริสาวจโร สมโณ โคตโม, นาลํ สลฺลาปาย.\csromanspacer{} โส อนฺตมนฺตาเนว เสวติ.\csromanspacer{} อิงฺฆ คหปติ สมโณ โคตโม อิมํ ปริสํ อาคจฺเฉยฺย, เอกปเญฺหเนว นํ สํสาเทยฺยาม\footnote{สํหเรยฺยาม (ก)}, ตุจฺฉกุมฺภีว นํ มญฺเญ โอโรเธยฺยามา”ติ.}

\proseitem{54}{อสฺโสสิ โข ภควา ทิพฺพาย โสตธาตุยา วิสุทฺธาย อติกฺกนฺตมานุสิกาย สนฺธานสฺส คหปติสฺส นิโคฺรเธน ปริพฺพาชเกน สทฺธิํ อิมํ กถาสลฺลาปํ.\csromanspacer{} อถ โข ภควา คิชฺฌกูฏา ปพฺพตา โอโรหิตฺวา เยน สุมาคธาย ตีเร โมรนิวาโป เตนุปสงฺกมิ, \csromanfolio{32}อุปสงฺกมิตฺวา สุมาคธาย ตีเร โมรนิวาเป อพฺโภกาเส จงฺกมิ.\csromanspacer{} อทฺทสา โข นิโคฺรโธ ปริพฺพาชโก ภควนฺตํ สุมาคธาย ตีเร โมรนิวาเป อพฺโภกาเส จงฺกมนฺตํ, ทิสฺวาน สกํ ปริสํ สณฺฐาเปสิ, “อปฺปสทฺทา โภนฺโต โหนฺตุ, มา โภนฺโต สทฺทมกตฺถ, อยํ สมโณ โคตโม สุมาคธาย ตีเร โมรนิวาเป อพฺโภกาเส จงฺกมติ.\csromanspacer{} อปฺปสทฺทกาโม โข ปน โส อายสฺมา, อปฺปสทฺทสฺส วณฺณวาที, อปฺเปว นาม อปฺปสทฺทํ ปริสํ วิทิตฺวา อุปสงฺกมิตพฺพํ มญฺเญยฺย.\csromanspacer{} สเจ สมโณ โคตโม อิมํ ปริสํ อาคจฺเฉยฺย, อิมํ ตํ ปญฺหํ ปุจฺเฉยฺยาม “โก นาม โส ภนฺเต ภควโต ธมฺโม, เยน ภควา สาวเก วิเนติ, เยน ภควตา สาวกา วินีตา อสฺสาสปฺปตฺตา ปฏิชานนฺติ อชฺฌาสยํ อาทิพฺรหฺมจริยนฺ”ติ.\csromanspacer{} เอวํ วุตฺเต เต ปริพฺพาชกา ตุณฺหี อเหสุํ.}

\csromanmatikaanchor{mtk.10}
\csromantocmarkref{section}{ตโปชิคุจฺฉาวาท}{mtk.10}
\bookmark[dest={mtk.10},level=1]{ตโปชิคุจฺฉาวาท}
\csromanheader{1}{\textbf{ตโปชิคุจฺฉาวาท}}
\proseitem{55}{อถ โข ภควา เยน นิโคฺรโธ ปริพฺพาชโก เตนุปสงฺกมิ.\csromanspacer{} อถ โข นิโคฺรโธ ปริพฺพาชโก ภควนฺตํ เอตทโวจ “เอตุ โข ภนฺเต ภควา, สฺวาคตํ ภนฺเต ภควโต, จิรสฺสํ โข ภนฺเต ภควา อิมํ ปริยายมกาสิ ยทิทํ อิธาคมนาย.\csromanspacer{} นิสีทตุ ภนฺเต ภควา, อิทมาสนํ ปญฺญตฺตนฺ”ติ.\csromanspacer{} นิสีทิ ภควา ปญฺญตฺเต อาสเน.\csromanspacer{} นิโคฺรโธปิ โข ปริพฺพาชโก อญฺญตรํ นีจาสนํ คเหตฺวา เอกมนฺตํ นิสีทิ.\csromanspacer{} เอกมนฺตํ นิสินฺนํ โข นิโคฺรธํ ปริพฺพาชกํ ภควา เอตทโวจ “กายนุตฺถ นิโคฺรธ เอตรหิ กถาย สนฺนิสินฺนา, กา จ ปน โว อนฺตรากถา วิปฺปกตา”ติ.\csromanspacer{} เอวํ วุตฺเต นิโคฺรโธ ปริพฺพาชโก ภควนฺตํ เอตทโวจ\csromandash{} อิธ มยํ ภนฺเต อทฺทสาม ภควนฺตํ สุมาคธาย ตีเร โมรนิวาเป อพฺโภกาเส จงฺกมนฺตํ, ทิสฺวาน เอวํ อโวจุมฺหา “สเจ สมโณ โคตโม อิมํ ปริสํ อาคจฺเฉยฺย, อิมํ ตํ ปญฺหํ ปุจฺเฉยฺยาม ‘โก นาม โส ภนฺเต ภควโต ธมฺโม, เยน ภควา สาวเก วิเนติ, เยน ภควตา สาวกา วินีตา อสฺสาสปฺปตฺตา ปฏิชานนฺติ อชฺฌาสยํ อาทิพฺรหฺมจริยนฺ’ติ.\csromanspacer{} อยํ โข โน ภนฺเต อนฺตรากถา วิปฺปกตา, อถ ภควา อนุปฺปตฺโต”ติ.}

\proseitem{56}{ทุชฺชานํ โข เอตํ นิโคฺรธ ตยา อญฺญทิฏฺฐิเกน อญฺญขนฺติเกน อญฺญรุจิเกน อญฺญตฺราโยเคน อญฺญตฺราจริยเกน, เยนาหํ สาวเก}

\vspace{\tipitakaparskip}
\csromancenter{อุทุมฺพริกสุตฺต วิเนมิ, เยน มยา สาวกา วินีตา อสฺสาสปฺปตฺตา ปฏิชานนฺติ อชฺฌาสยํ อาทิพฺรหฺมจริยํ.\csromanspacer{} อิงฺฆ ตฺวํ มํ นิโคฺรธ สเก อาจริยเก อธิเชคุจฺเฉ ปญฺหํ ปุจฺฉ “กถํ สนฺตา นุ โข ภนฺเต ตโปชิคุจฺฉา ปริปุณฺณา โหติ, กถํ อปริปุณฺณา”ติ.\csromanspacer{} เอวํ วุตฺเต เต ปริพฺพาชกา อุนฺนาทิโน อุจฺจาสทฺทมหาสทฺทา อเหสุํ “อจฺฉริยํ วต โภ, อพฺภุตํ วต โภ, สมณสฺส โคตมสฺส มหิทฺธิกตา มหานุภาวตา, ยตฺร หิ นาม สกวาทํ ฐเปสฺสติ, ปรวาเทน ปวาเรสฺสตี”ติ.}
\proseitem{57}{\csromanfolio{33}อถ โข นิโคฺรโธ ปริพฺพาชโก เต ปริพฺพาชเก อปฺปสทฺเท กตฺวา ภควนฺตํ เอตทโวจ “มยํ โข ภนฺเต ตโปชิคุจฺฉาวาทา ตโปชิคุจฺฉาสารา\footnote{ตโปชิคุจฺฉํสาโรทา (ก)} ตโปชิคุจฺฉา-อลฺลีนา วิหราม.\csromanspacer{} กฺถํ สนฺตา นุ โข ภนฺเต ตโปชิคุจฺฉา ปริปุณฺณา โหติ, กถํ อปริปุณฺณา”ติ.}

\prose{อิธ นิโคฺรธ ตปสฺสี อเจลโก โหติ มุตฺตาจาโร, หตฺถาปเลขโน\footnote{หตฺถาวเลขโน (สฺยา)}, น-เอหิภทฺทนฺติโก, นติฏฺฐภทฺทนฺติโก, นาภิหฏํ, น-อุทฺทิสฺสกตํ, นนิมนฺตนํ สาทิยติ, โส น กุมฺภิมุขา ปฏิคฺคณฺหาติ, น กโฬปิมุขา ปฏิคฺคณฺหาติ, น เอฬกมนฺตรํ, น ทณฺฑมนฺตรํ, น มุสลมนฺตรํ, น ทฺวินฺนํ ภุญฺชมานานํ, น คพฺภินิยา, น ปายมานาย, น ปุริสนฺตรคตาย, น สงฺกตฺตีสุ, น ยตฺถ สา อุปฏฺฐิโต โหติ, น ยตฺถ มกฺขิกา สณฺฑสณฺฑจารินี, น มจฺฉํ, น มํสํ, น สุรํ, น เมรยํ, น ถุโสทกํ ปิวติ, โส เอกาคาริโก วา โหติ เอกาโลปิโก, ทฺวาคาริโก วา โหติ ทฺวาโลปิโก, สตฺตาคาริโก วา โหติ สตฺตาโลปิโก, เอกิสฺสาปิ ทตฺติยา ยาเปติ, ทฺวีหิปิ ทตฺตีหิ ยาเปติ, สตฺตหิปิ ทตฺตีหิ ยาเปติ, เอกาหิกมฺปิ อาหารํ อาหาเรติ, ทฺวีหิกมฺปิ\footnote{ทฺวาหิกํปิ ((สี, สฺยา)} อาหารํ อาหาเรติ, สตฺตาหิกมฺปิ อาหารํ อาหาเรติ, อิติ เอวรูปํ อทฺธมาสิกํปริยายภตฺตโภชนานุโยคมนุยุตฺโต วิหรติ.\csromanspacer{} โส สากภกฺโข วา โหติ, สามากภกฺโข วา โหติ, นีวารภกฺโข วา โหติ, ททฺทุลภกฺโข วา โหติ, หฏภกฺโข วา โหติ, กณภกฺโข วา โหติ, อาจามภกฺโข วา โหติ, ปิญฺญากภกฺโข วา โหติ, ติณภกฺโข วา โหติ, โคมยภกฺโข วา โหติ, วนมูลผลาหาโร ยาเปติ ปวตฺตผลโภชี.\csromanspacer{} โส สาณานิปิ \csromanfolio{34}ธาเรติ, มสาณานิปิ ธาเรติ, ฉวทุสฺสานิปิ ธาเรติ, ปํสุกูลานิปิ ธาเรติ, ติรีฏานิปิ ธาเรติ, อชินานิปิ ธาเรติ, อชินกฺขิปมฺปิ ธาเรติ, กุสจีรมฺปิ ธาเรติ, วากจีรมฺปิ ธาเรติ, ผลกจีรมฺปิ ธาเรติ, เกสกมฺพลมฺปิ ธาเรติ, วาฬกมฺพลมฺปิ ธาเรติ, อุลูกปกฺขมฺปิ ธาเรติ, เกสมสฺสุโลจโกปิ โหติ เกสมสฺสุโลจนานุโยคมนุยุตฺโต, อุพฺภฏฺฐโกปิ\footnote{อุภฏฺฐโกปิ (สฺยา), อุพฺภฏฺฐิโกปิ (ก)} โหติ อาสนปฏิกฺขิตฺโต, อุกฺกุฏิโกปิ โหติ อุกฺกุฏิกปฺปธานมนุยุตฺโต, กณฺฏกาปสฺสยิโกปิ โหติ กณฺฏกาปสฺสเย เสยฺยํ กปฺเปติ, ผลกเสยฺยมฺปิ กปฺเปติ, ถณฺฑิลเสยฺยมฺปิ กปฺเปติ, เอกปสฺสยิโกปิ โหติ รโชชลฺลธโร, อพฺโภกาสิโกปิ โหติ ยถาสนฺถติโก, เวกฏิโกปิ โหติ วิกฏโภชนานุโยคมนุยุตฺโต, อปานโกปิ โหติ อปานกตฺตมนุยุตฺโต, สายตติยกํปิ อุทโกโรหนานุโยคมนุยุตฺโต วิหรติ.\csromanspacer{} ตํ กิํมญฺญสิ นิโคฺรธ, ยทิ เอวํ สนฺเต ตโปชิคุจฺฉา ปริปุณฺณา วา โหติ อปริปุณฺณา วาติ.\csromanspacer{} อทฺธา โข ภนฺเต เอวํ สนฺเต ตโปชิคุจฺฉา ปริปุณฺณา โหติ, โน อปริปุณฺณาติ.\csromanspacer{} เอวํ ปริปุณฺณายปิ โข อหํ นิโคฺรธ ตโปชิคุจฺฉาย อเนกวิหิเต อุปกฺกิเลเส วทามีติ.}

\csromanmatikaanchor{mtk.11}
\csromantocmarkref{section}{อุปกฺกิเลส}{mtk.11}
\bookmark[dest={mtk.11},level=1]{อุปกฺกิเลส}
\csromanheader{1}{\textbf{อุปกฺกิเลส}}
\proseitem{58}{ยถา กถํ ปน ภนฺเต ภควา เอวํ ปริปุณฺณาย ตโปชิคุจฺฉาย อเนกวิหิเต อุปกฺกิเลเส วทตีติ.\csromanspacer{} อิธ นิโคฺรธ ตปสฺสี ตปํ สมาทิยติ, โส เตน ตปสา อตฺตมโน โหติ ปริปุณฺณสงฺกปฺโป.\csromanspacer{} ยมฺปิ นิโคฺรธ ตปสฺสี ตปํ สมาทิยติ, โส เตน ตปสา อตฺตมโน โหติ ปริปุณฺณสงฺกปฺโป, อยมฺปิ โข นิโคฺรธ ตปสฺสิโน อุปกฺกิเลโส โหติ.}

\prose{ปุน จปรํ นิโคฺรธ ตปสฺสี ตปํ สมาทิยติ, โส เตน ตปสา อตฺตานุกฺกํเสติ ปรํ วมฺเภติ.\csromanspacer{} ยมฺปิ นิโคฺรธ ตปสฺสี ตปํ สมาทิยติ, โส เตน ตปสา อตฺตานุกฺกํเสติ ปรํ วมฺเภติ, อยมฺปิ โข นิโคฺรธ ตปสฺสิโน อุปกฺกิเลโส โหติ.}

\prose{ปุน จปรํ นิโคฺรธ ตปสฺสี ตปํ สมาทิยติ, โส เตน ตปสา มชฺชติ มุจฺฉติ ปมาทมาปชฺชติ\footnote{มทมาปชฺชติ (สฺยา)}.\csromanspacer{} ยมฺปิ นิโคฺรธ ตปสฺสี ตปํ สมาทิยติ,}

\vspace{\tipitakaparskip}
\csromancenter{อุทุมฺพริกสุตฺต โส เตน ตปสา มชฺชติ มุจฺฉติ ปมาทมาปชฺชติ, อยมฺปิ โข นิโคฺรธ ตปสฺสิโน อุปกฺกิเลโส โหติ.}
\proseitem{59}{\csromanfolio{35}ปุน จปรํ นิโคฺรธ ตปสฺสี ตปํ สมาทิยติ, โส เตน ตปสา ลาภสกฺการสิโลกํ อภินิพฺพตฺเตติ, โส เตน ลาภสกฺการสิโลเกน อตฺตมโน โหติ ปริปุณฺณสงฺกปฺโป.\csromanspacer{} ยมฺปิ นิโคฺรธ ตปสฺสี ตปํ สมาทิยติ, โส เตน ตปสา ลาภสกฺการสิโลกํ อภินิพฺพตฺเตติ, โส เตน ลาภสกฺการสิโลเกน อตฺตมโน โหติ ปริปุณฺณสงฺกปฺโป, อยมฺปิ โข นิโคฺรธ ตปสฺสิโน อุปกฺกิเลโส โหติ.}

\prose{ปุน จปรํ นิโคฺรธ ตปสฺสี ตปํ สมาทิยติ, โส เตน ตปสา ลาภสกฺการสิโลกํ อภินิพฺพตฺเตติ, โส เตน ลาภสกฺการสิโลเกน อตฺตานุกฺกํเสติ ปรํ วมฺเภติ.\csromanspacer{} ยมฺปิ นิโคฺรธ ตปสฺสี ตปํ สมาทิยติ, โส เตน ตปสา ลาภสกฺการสิโลกํ อภินิพฺพตฺเตติ, โส เตน ลาภสกฺการสิโลเกน อตฺตานุกฺกํเสติ ปรํ วมฺเภติ, อยมฺปิ โข นิโคฺรธ ตปสฺสิโน อุปกฺกิเลโส โหติ.}

\prose{ปุน จปรํ นิโคฺรธ ตปสฺสี ตปํ สมาทิยติ, โส เตน ตปสา ลาภสกฺการสิโลกํ อภินิพฺพตฺเตติ, โส เตน ลาภสกฺการสิโลเกน มชฺชติ มุจฺฉติ ปมาทมาปชฺชติ.\csromanspacer{} ยมฺปิ โข นิโคฺรธ ตปสฺสี ตปํ สมาทิยติ, โส เตน ตปสา ลาภสกฺการสิโลกํ อภินิพฺพตฺเตติ, โส เตน ลาภสกฺการสิโลเกน มชฺชติ มุจฺฉติ ปมาทมาปชฺชติ, อยมฺปิ โข นิโคฺรธ ตปสฺสิโน อุปกฺกิเลโส โหติ.}

\proseitem{60}{ปุน จปรํ นิโคฺรธ ตปสฺสี ตปํ สมาทิยติ, โภชเนสุ โวทาสํ อาปชฺชติ “อิทํ เม ขมติ, อิทํ เม นกฺขมตี”ติ.\csromanspacer{} โส ยญฺจ\footnote{ยํ หิ (สี, อิ)} ขฺวสฺส นกฺขมติ, ตํ สาเปกฺโข ปชหติ.\csromanspacer{} ยํ ปนสฺส ขมติ, ตํ คธิโต\footnote{คถิโต (สี, อิ)} มุจฺฉิโต อชฺฌาปนฺโน อนาทีนวทสฺสาวี อนิสฺสรณปญฺโญ ปริภุญฺชติ ฯเปฯ อยมฺปิ โข นิโคฺรธ ตปสฺสิโน อุปกฺกิเลโส โหติ.}

\prose{ปุน จปรํ นิโคฺรธ ตปสฺสี ตปํ สมาทิยติ ลาภสกฺการสิโลกนิกนฺติเหตุ “สกฺกริสฺสนฺติ มํ ราชาโน ราชมหามตฺตา ขตฺติยา \csromanfolio{36}พฺราหฺมณา คหปติกา ติตฺถิยา”ติ ฯเปฯ อยมฺปิ โข นิโคฺรธ ตปสฺสิโน อุปกฺกิเลโส โหติ.}

\proseitem{61}{ปุน จปรํ นิโคฺรธ ตปสฺสี อญฺญตรํ สมณํ วา พฺราหฺมณํ วา อปสาเทตา\footnote{อปสาเรตา (ก)} โหติ “กิํ ปนายํ สมฺพหุลาชีโว\footnote{พหุลาชีโว (สี, อิ)} สพฺพํ สํภกฺเขติ.\csromanspacer{} เสยฺยถิทํ, มูลพีชํ ขนฺธพีชํ ผลุพีชํ อคฺคพีชํ พีชพีชเมว ปญฺจมํ, อสนิวิจกฺกํ ทนฺตกูฏํ, สมณปฺปวาเทนา”ติ ฯเปฯ อยมฺปิ โข นิโคฺรธ ตปสฺสิโน อุปกฺกิเลโส โหติ.}

\prose{ปุน จปรํ นิโคฺรธ ตปสฺสี ปสฺสติ อญฺญตรํ สมณํ วา พฺราหฺมณํ วา กุเลสุ สกฺกริยมานํ ครุกริยมานํ มานิยมานํ ปูชิยมานํ, ทิสฺวา ตสฺส เอวํ โหติ “อิมญฺหิ นาม สมฺพหุลาชีวํ กุเลสุ สกฺกโรนฺติ ครุํ กโรนฺติ มาเนนฺติ ปูเชนฺติ, มํ ปน ตปสฺสิํ ลูขาชีวิํ กุเลสุ น สกฺกโรนฺติ น ครุํ กโรนฺติ น มาเนนฺติ น ปูเชนฺตี”ติ, อิติ โส อิสฺสามจฺฉริยํ กุเลสุ อุปฺปาเทตา โหติ ฯเปฯ อยมฺปิ โข นิโคฺรธ ตปสฺสิโน อุปกฺกิเลโส โหติ.}

\proseitem{62}{ปุน จปรํ นิโคฺรธ ตปสฺสี อาปาถกนิสาที โหติ ฯเปฯ อยมฺปิ โข นิโคฺรธ ตปสฺสิโน อุปกฺกิเลโส โหติ.}

\prose{ปุน จปรํ นิโคฺรธ ตปสฺสี อตฺตานํ อทสฺสยมาโน กุเลสุ จรติ “อิทมฺปิ เม ตปสฺมิํ อิทมฺปิ เม ตปสฺมินฺ”ติ ฯเปฯ อยมฺปิ โข นิโคฺรธ ตปสฺสิโน อุปกฺกิเลโส โหติ.}

\prose{ปุน จปรํ นิโคฺรธ ตปสฺสี กิญฺจิเทว ปฏิจฺฉนฺนํ เสวติ, โส “ขมติ เต อิทนฺ”ติ ปุฏฺโฐ สมาโน อกฺขมมานํ อาห “ขมตี”ติ.\csromanspacer{} ขมมานํ อาห “นกฺขมตี”ติ.\csromanspacer{} อิติ โส สมฺปชานมุสา ภาสิตา โหติ ฯเปฯ อยมฺปิ โข นิโคฺรธ ตปสฺสิโน อุปกฺกิเลโส โหติ.}

\prose{ปุน จปรํ นิโคฺรธ ตปสฺสี ตถาคตสฺส วา ตถาคตสาวกสฺส วา ธมฺมํ เทเสนฺตสฺส สนฺตํเยว ปริยายํ อนุญฺเญยฺยํ นานุชานาติ ฯเปฯ อยมฺปิ โข นิโคฺรธ ตปสฺสิโน อุปกฺกิเลโส โหติ.}

\csromanheader{2}{2. อุทุมฺพริกสุตฺต}
\proseitem{63}{\csromanfolio{37}ปุน จปรํ นิโคฺรธ ตปสฺสี โกธโน โหติ อุปนาหี.\csromanspacer{} ยมฺปิ นิโคฺรธ ตปสฺสี โกธโน โหติ อุปนาหี.\csromanspacer{} อยมฺปิ โข นิโคฺรธ ตปสฺสิโน อุปกฺกิเลโส โหติ.}

\prose{ปุน จปรํ นิโคฺรธ ตปสฺสี มกฺขี โหติ ปฬาสี\footnote{ปลาสี (สี, สฺยา, อิ)} ฯเปฯ อิสฺสุกี โหติ มจฺฉรี.\csromanspacer{} สโฐ โหติ มายาวี.\csromanspacer{} ถทฺโธ โหติ อติมานี.\csromanspacer{} ปาปิจฺโฉ โหติ ปาปิกานํ อิจฺฉานํ วสํ คโต.\csromanspacer{} มิจฺฉาทิฏฺฐิโก โหติ อนฺตคฺคาหิกาย ทิฏฺฐิยา สมนฺนาคโต.\csromanspacer{} สนฺทิฏฺฐิปรามาสี โหติ อาธานคฺคาหี ทุปฺปฏินิสฺสคฺคี.\csromanspacer{} ยมฺปิ นิโคฺรธ ตปสฺสี สนฺทิฏฺฐิปรามาสี โหติ อาธานคฺคาหี ทุปฺปฏินิสฺสคฺคี.\csromanspacer{} อยมฺปิ โข นิโคฺรธ ตปสฺสิโน อุปกฺกิเลโส โหติ.}

\prose{ตํ กิํ มญฺญสิ นิโคฺรธ “ยทิเม ตโปชิคุจฺฉา\footnote{ตโปชิคุจฺฉาย (?) 3. อุปกฺกิเลสา โหติ (ก)} อุปกฺกิเลสา วา อนุปกฺกิเลสา วา”ติ.\csromanspacer{} อทฺธา โข อิเม ภนฺเต ตโปชิคุจฺฉา2 อุปกฺกิเลสา3, โน อนุปกฺกิเลสา.\csromanspacer{} ฐานํ โข ปเนตํ ภนฺเต วิชฺชติ “ยํ อิเธกจฺโจ ตปสฺสี สพฺเพเหว อิเมหิ อุปกฺกิเลเสหิ สมนฺนาคโต อสฺส, โก ปน วาโท อญฺญตรญฺญตเรนา”ติ.}

\csromanmatikaanchor{mtk.12}
\csromantocmarkref{section}{ปริสุทฺธปปฏิกปฺปตฺตกถา}{mtk.12}
\bookmark[dest={mtk.12},level=1]{ปริสุทฺธปปฏิกปฺปตฺตกถา}
\csromanheader{1}{\textbf{ปริสุทฺธปปฏิกปฺปตฺตกถา}}
\proseitem{64}{อิธ นิโคฺรธ ตปสฺสี ตปํ สมาทิยติ, โส เตน ตปสา น อตฺตมโน โหติ น ปริปุณฺณาสงฺกปฺโป.\csromanspacer{} ยมฺปิ นิโคฺรธ ตปสฺสี ตปํ สมาทิยติ, โส เตน ตปสา น อตฺตมโน โหติ น ปริปุณฺณสงฺกปฺโป.\csromanspacer{} เอวํ โส ตสฺมิํ ฐาเน ปริสุทฺโธ โหติ.}

\prose{ปุน จปรํ นิโคฺรธ ตปสฺสี ตปํ สมาทิยติ, โส เตน ตปสา น อตฺตานุกฺกํเสติ น ปรํ วมฺเภติ ฯเปฯ.\csromanspacer{} เอวํ โส ตสฺมิํ ฐาเน ปริสุทฺโธ โหติ.}

\prose{ปุน จปรํ นิโคฺรธ ตปสฺสี ตปํ สมาทิยติ, โส เตน ตปสา น มชฺชติ น มุจฺฉติ น ปมาทมาปชฺชติ ฯเปฯ.\csromanspacer{} เอวํ โส ตสฺมิํ ฐาเน ปริสุทฺโธ โหติ.}

\proseitem{65}{\csromanfolio{38}ปุน จปรํ นิโคฺรธ ตปสฺสี ตปํ สมาทิยติ, โส เตน ตปสา ลาภสกฺการสิโลกํ อภินิพฺพตฺเตติ, โส เตน ลาภสกฺการสิโลเกน น อตฺตมโน โหติ น ปริปุณฺณสงฺกปฺโป ฯเปฯ.\csromanspacer{} เอวํ โส ตสฺมิํ ฐาเน ปริสุทฺโธ โหติ.}

\prose{ปุน จปรํ นิโคฺรธ ตปสฺสี ตปํ สมาทิยติ, โส เตน ตปสา ลาภสกฺการสิโลกํ อภินิพฺพตฺเตติ, โส เตน ลาภสกฺการสิโลเกน น อตฺตานุกฺกํเสติ น ปรํ วมฺเภติ ฯเปฯ.\csromanspacer{} เอวํ โส ตสฺมิํ ฐาเน ปริสุทฺโธ โหติ.}

\prose{ปุน จปรํ นิโคฺรธ ตปสฺสี ตปํ สมาทิยติ, โส เตน ตปสา ลาภสกฺการสิโลกํ อภินิพฺพตฺเตติ, โส เตน ลาภสกฺการสิโลเกน น มชฺชติ น มุจฺฉติ น ปมาทมาปชฺชติ ฯเปฯ.\csromanspacer{} เอวํ โส ตสฺมิํ ฐาเน ปริสุทฺโธ โหติ.}

\proseitem{66}{ปุน จปรํ นิโคฺรธ ตปสฺสี โภชเนสุ น โวทาสํ อาปชฺชติ “อิทํ เม ขมติ, อิทํ เม นกฺขมตี”ติ.\csromanspacer{} โส ยญฺจ ขฺวสฺส นกฺขมติ, ตํ อนเปกฺโข ปชหติ.\csromanspacer{} ยํ ปนสฺส ขมติ, ตํ อคธิโต อมุจฺฉิโต อนชฺฌาปนฺโน อาทีนวทสฺสาวี นิสฺสรณปญฺโญ ปริภุญฺชติ ฯเปฯ.\csromanspacer{} เอวํ โส ตสฺมิํ ฐาเน ปริสุทฺโธ โหติ.}

\prose{ปุน จปรํ นิโคฺรธ ตปสฺสี น ตปํ สมาทิยติ ลาภสกฺการสิโลกนิกนฺติเหตุ “สกฺกริสฺสนฺติ มํ ราชาโน ราชมหามตฺตา ขตฺติยา พฺราหฺมณา คหปติกา ติตฺถิยา”ติ.\csromanspacer{} เอวํ โส ตสฺมิํ ฐาเน ปริสุทฺโธ โหติ.}

\proseitem{67}{ปุน จปรํ นิโคฺรธ ตปสฺสี อญฺญตรํ สมณํ วา พฺราหฺมณํ วา นาปสาเทตา โหติ “กิํ ปนายํ สมฺพหุลาชีโว สพฺพํ สํภกฺเขติ.\csromanspacer{} เสยฺยถิทํ, มูลพีชํ ขนฺธพีชํ ผลุพีชํ อคฺคพีชํ พีชพีชเมว ปญฺจมํ, อสนิวิจกฺกํ ทนฺตกูฏํ, สมณปฺปวาเทนา”ติ ฯเปฯ.\csromanspacer{} เอวํ โส ตสฺมิํ ฐาเน ปริสุทฺโธ โหติ.}

\prose{ปุน จปรํ นิโคฺรธ ตปสฺสี ปสฺสติ อญฺญตรํ สมณํ วา พฺราหฺมณํ วา กุเลสุ สกฺกริยมานํ ครุํ กริยมานํ มานิยมานํ ปูชิยมานํ, ตสฺส น}

\vspace{\tipitakaparskip}
\csromancenter{อุทุมฺพริกสุตฺต เอวํ โหติ “อิมญฺหิ นาม สมฺพหุลาชีวํ กุเลสุ สกฺกโรนฺติ ครุํ กโรนฺติ มาเนนฺติ ปูเชนฺติ, มํ ปน ตปสฺสิํ ลูขาชีวิํ กุเลสุ น สกฺกโรนฺติ น ครุํ กโรนฺติ น มาเนนฺติ น ปูเชนฺตี”ติ, อิติ โส อิสฺสามจฺฉริยํ กุเลสุ นุปฺปาเทตา โหติ ฯเปฯ.\csromanspacer{} เอวํ โส ตสฺมิํ ฐาเน ปริสุทฺโธ โหติ.}
\proseitem{68}{\csromanfolio{39}ปุน จปรํ นิโคฺรธ ตปสฺสี น อาปาถกนิสาที โหติ ฯเปฯ.\csromanspacer{} เอวํ โส ตสฺมิํ ฐาเน ปริสุทฺโธ โหติ.}

\prose{ปุน จปรํ นิโคฺรธ ตปสฺสี น อตฺตานํ อทสฺสยมาโน กุเลสุ จรติ “อิทมฺปิ เม ตปสฺมิํ อิทมฺปิ เม ตปสฺมินฺ”ติ ฯเปฯ.\csromanspacer{} เอวํ โส ตสฺมิํ ฐาเน ปริสุทฺโธ โหติ.}

\prose{ปุน จปรํ นิโคฺรธ ตปสฺสี น กิญฺจิเทว ปฏิจฺฉนฺนํ เสวติ, โส “ขมติ เต อิทนฺ”ติ ปุฏฺโฐ สมาโน อกฺขมมานํ อาห “นกฺขมตี”ติ.\csromanspacer{} ขมมานํ อาห “ขมตี”ติ.\csromanspacer{} อิติ โส สมฺปชานมุสา น ภาสิตา โหติ ฯเปฯ.\csromanspacer{} เอวํ โส ตสฺมิํ ฐาเน ปริสุทฺโธ โหติ.}

\prose{ปุน จปรํ นิโคฺรธ ตปสฺสี ตถาคตสฺส วา ตถาคตสาวกสฺส วา ธมฺมํ เทเสนฺตสฺส สนฺตํเยว ปริยายํ อนุญฺเญยฺยํ อนุชานาติ ฯเปฯ.\csromanspacer{} เอวํ โส ตสฺมิํ ฐาเน ปริสุทฺโธ โหติ.}

\proseitem{69}{ปุน จปรํ นิโคฺรธ ตปสฺสี อกฺโกธโน โหติ อนุปนาหี.\csromanspacer{} ยมฺปิ นิโคฺรธ ตปสฺสี อกฺโกธโน โหติ อนุปนาหี ฯเปฯ.\csromanspacer{} เอวํ โส ตสฺมิํ ฐาเน ปริสุทฺโธ โหติ.}

\prose{ปุน จปรํ นิโคฺรธ ตปสฺสี อมกฺขี โหติ อปฬาสี ฯเปฯ อนิสฺสุกี โหติ อมจฺฉรี.\csromanspacer{} อสโฐ โหติ อมายาวี.\csromanspacer{} อตฺถทฺโธ โหติ อนติมานี.\csromanspacer{} น ปาปิจฺโฉ โหติ น ปาปิกานํ อิจฺฉานํ วสํ คโต.\csromanspacer{} น มิจฺฉาทิฏฺฐิโก โหติ น อนฺตคฺคาหิกาย ทิฏฺฐิยา สมนฺนาคโต.\csromanspacer{} น สนฺทิฏฺฐิปรามาสี โหติ น อาธานคฺคาหี สุปฺปฏินิสฺสคฺคี.\csromanspacer{} ยมฺปิ นิโคฺรธ ตปสฺสี น สนฺทิฏฺฐิปรามาสี โหติ น อาธานคฺคาหี สุปฺปฏินิสฺสคฺคี.\csromanspacer{} เอวํ โส ตสฺมิํ ฐาเน ปริสุทฺโธ โหติ.}

\prose{ตํ กิํ มญฺญสิ นิโคฺรธ “ยทิ เอวํ สนฺเต ตโปชิคุจฺฉา ปริสุทฺธา วา โหติ อปริสุทฺธา วา”ติ.\csromanspacer{} อทฺธา โข ภนฺเต เอวํ สนฺเต ตโปชิคุจฺฉา \csromanfolio{40}ปริสุทฺธา โหติ โน อปริสุทฺธา, อคฺคปฺปตฺตา จ สารปฺปตฺตา จาติ.\csromanspacer{} น โข นิโคฺรธ เอตฺตาวตา ตโปชิคุจฺฉา อคฺคปฺปตฺตา จ สารปฺปตฺตา จ โหติ, อปิ จ โข ปปฏิกปฺปตฺตา\footnote{ปปฺปฏิกปตฺตา (ก)} โหติ.}

\csromanmatikaanchor{mtk.13}
\csromantocmarkref{section}{ปริสุทฺธตจปฺปตฺตกถา}{mtk.13}
\bookmark[dest={mtk.13},level=1]{ปริสุทฺธตจปฺปตฺตกถา}
\csromanheader{1}{\textbf{ปริสุทฺธตจปฺปตฺตกถา}}
\proseitem{70}{กิตฺตาวตา ปน ภนฺเต ตโปชิคุจฺฉา อคฺคปฺปตฺตา จ สารปฺปตฺตา จ โหติ, สาธุ ภนฺเต ภควา ตโปชิคุจฺฉาย อคฺคญฺเญว ปาเปตุ, สารญฺเญว ปาเปตูติ.\csromanspacer{} อิธ นิโคฺรธ ตปสฺสี จาตุยามสํวรสํวุโต โหติ.\csromanspacer{} กถญฺจ นิโคฺรธ ตปสฺสี จาตุยามสํวรสํวุโต โหติ.\csromanspacer{} อิธ นิโคฺรธ ตปสฺสี น ปาณํ อติปาเตติ\footnote{อติปาเปติ (ก-สี, อิ ก)}, น ปาณํ อติปาตยติ, น ปาณมติปาตยโต สมนุญฺโญ โหติ.\csromanspacer{} น อทินฺนํ อาทิยติ, น อทินฺนํ อาทิยาเปติ, น อทินฺนํ อาทิยโต สมนุญฺโญ โหติ.\csromanspacer{} น มุสา ภณติ, น มุสา ภณาเปติ, น มุสา ภณโต สมนุญฺโญ โหติ.\csromanspacer{} น ภาวิตมาสีสติ\footnote{น ภาวิตมาสิํสติ (สี, สฺยา, อิ)}, น ภาวิตมาสีสาเปติ, น ภาวิตมาสีสโต สมนุญฺโญ โหติ.\csromanspacer{} เอวํ โข นิโคฺรธ ตปสฺสี จาตุยามสํวรสํวุโต โหติ.}

\prose{ยโต โข นิโคฺรธ ตปสฺสี จาตุยามสํวรสํวุโต โหติ, อทุํ จสฺส โหติ ตปสฺสิตาย.\csromanspacer{} โส อภิหรติ โน หีนายาวตฺตติ.\csromanspacer{} โส วิวิตฺตํ เสนาสนํ ภชติ อรญฺญํ รุกฺขมูลํ ปพฺพตํ กนฺทรํ คิริคุหํ สุสานํ วนปตฺถํ อพฺโภกาสํ ปลาลปุญฺชํ.\csromanspacer{} โส ปจฺฉาภตฺตํ ปิณฺฑปาตปฏิกฺกนฺโต นิสีทติ ปลฺลงฺกํ อาภุชิตฺวา อุชุํ กายํ ปณิธาย ปริมุขํ สติํ อุปฏฺฐเปตฺวา.\csromanspacer{} โส อภิชฺฌํ โลเก ปหาย วิคตาภิชฺเฌน เจตสา วิหรติ, อภิชฺฌาย จิตฺตํ ปริโสเธติ.\csromanspacer{} พฺยาปาทปฺปโทสํ ปหาย อพฺยาปนฺนจิตฺโต วิหรติ สพฺพปาณภูตหิตานุกมฺปี, พฺยาปาทปฺปโทสา จิตฺตํ ปริโสเธติ.\csromanspacer{} ถินมิทฺธํ\footnote{ถีนมิทฺธํ (สี, สฺยา, อิ)} ปหาย วิคตถินมิทฺโธ วิหรติ อาโลกสญฺญี สโต สมฺปชาโน, ถินมิทฺธา จิตฺตํ ปริโสเธติ.\csromanspacer{} อุทฺธจฺจกุกฺกุจฺจํ ปหาย อนุทฺธโต วิหรติ อชฺฌตฺตํ วูปสนฺตจิตฺโต, อุทฺธจฺจกุกฺกุจฺจา จิตฺตํ ปริโสเธติ.\csromanspacer{} วิจิกิจฺฉํ ปหาย ติณฺณวิจิกิจฺโฉ วิหรติ อกถํกถี กุสเลสุ ธมฺเมสุ, วิจิกิจฺฉาย จิตฺตํ ปริโสเธติ.}

\csromanheader{2}{2. อุทุมฺพริกสุตฺต}
\proseitem{71}{\csromanfolio{41}โส อิเม ปญฺจ นีวรเณ ปหาย เจตโส อุปกฺกิเลเส ปญฺญาย ทุพฺพลีกรเณ เมตฺตาสหคเตน เจตสา เอกํ ทิสํ ผริตฺวา วิหรติ.\csromanspacer{} ตถา ทุติยํ.\csromanspacer{} ตถา ตติยํ.\csromanspacer{} ตถา จตุตฺถํ.\csromanspacer{} อิติ อุทฺธมโธ ติริยํ สพฺพธิ สพฺพตฺตตาย สพฺพาวนฺตํ โลกํ เมตฺตาสหคเตน เจตสา วิปุเลน มหคฺคเตน อปฺปมาเณน อเวเรน อพฺยาปชฺเชน ผริตฺวา วิหรติ.\csromanspacer{} กรุณาสหคเตน เจตสา ฯเปฯ.\csromanspacer{} มุทิตาสหคเตน เจตสา ฯเปฯ.\csromanspacer{} อุเปกฺขาสหคเตน เจตสา เอกํ ทิสํ ผริตฺวา วิหรติ.\csromanspacer{} ตถา ทุติยํ.\csromanspacer{} ตถา ตติยํ.\csromanspacer{} ตถา จตุตฺถํ.\csromanspacer{} อิติ อุทฺธมโธ ติริยํ สพฺพธิ สพฺพตฺตตาย สพฺพาวนฺตํ โลกํ อุเปกฺขาสหคเตน เจตสา วิปุเลน มหคฺคเตน อปฺปมาเณน อเวเรน อพฺยาปชฺเชน ผริตฺวา วิหรติ.}

\prose{ตํ กิํ มญฺญสิ นิโคฺรธ “ยทิ เอวํ สนฺเต ตโปชิคุจฺฉา ปริสุทฺธา วา โหติ อปริสุทฺธา วา”ติ.\csromanspacer{} อทฺธา โข ภนฺเต เอวํ สนฺเต ตโปชิคุจฺฉา ปริสุทฺธา โหติ โน อปริสุทฺธา, อคฺคปฺปตฺตา จ สารปฺปตฺตา จาติ.\csromanspacer{} น โข นิโคฺรธ เอตฺตาวตา ตโปชิคุจฺฉา อคฺคปฺปตฺตา จ โหติ สารปฺปตฺตา จ, อปิ จ โข ตจปฺปตฺตา โหตีติ.}

\csromanmatikaanchor{mtk.14}
\csromantocmarkref{section}{ปริสุทฺธเผคฺคุปฺปตฺตกถา}{mtk.14}
\bookmark[dest={mtk.14},level=1]{ปริสุทฺธเผคฺคุปฺปตฺตกถา}
\csromanheader{1}{\textbf{ปริสุทฺธเผคฺคุปฺปตฺตกถา}}
\proseitem{72}{กิตฺตาวตา โข ปน ภนฺเต ตโปชิคุจฺฉา อคฺคปฺปตฺตา จ โหติ สารปฺปตฺตา จ, สาธุ เม ภนฺเต ภควา ตโปชิคุจฺฉาย อคฺคญฺเญว ปาเปตุ, สารญฺเญว ปาเปตูติ.\csromanspacer{} อิธ นิโคฺรธ ตปสฺสี จาตุยามสํวรสํวุโต โหติ.\csromanspacer{} กถญฺจ ปน นิโคฺรธ ตปสฺสี จาตุยามสํวรสํวุโต โหติ ฯเปฯ.\csromanspacer{} ยโต โข นิโคฺรธ ตปสฺสี เอวํ จาตุยามสํวรสํวุโต โหติ, อทุํ จสฺส โหติ ตปสฺสิตาย.\csromanspacer{} โส อภิหรติ โน หีนายาวตฺตติ.\csromanspacer{} โส วิวิตฺตํ เสนาสนํ ภชติ ฯเปฯ.\csromanspacer{} โส อิเม ปญฺจ นีวรเณ ปหาย เจตโส อุปกฺกิเลเส ปญฺญาย ทุพฺพลีกรเณ เมตฺตาสหคเตน เจตสา ฯเปฯ.\csromanspacer{} กรุณาสหคเตน เจตสา ฯเปฯ.\csromanspacer{} มุทิตาสหคเตน เจตสา ฯเปฯ.\csromanspacer{} อุเปกฺขาสหคเตน เจตสา วิปุเลน มหคฺคเตน อปฺปมาเณน อเวเรน อพฺยาปชฺเชน ผริตฺวา วิหรติ.\csromanspacer{} โส อเนกวิหิตํ ปุพฺเพนิวาสํ อนุสฺสรติ.\csromanspacer{} เสยฺยถิทํ, เอกมฺปิ ชาติํ เทฺวปิ ชาติโย ติสฺโสปิ ชาติโย จตสฺโสปิ ชาติโย ปญฺจปิ ชาติโย ทสปิ ชาติโย วีสมฺปิ ชาติโย ติํสมฺปิ ชาติโย \csromanfolio{42}จตฺตาลีสมฺปิ ชาติโย ปญฺญาสมฺปิ ชาติโย ชาติสตมฺปิ ชาติสหสฺสมฺปิ ชาติสตสหสฺสมฺปิ อเนเกปิ สํวฏฺฏกปฺเป อเนเกปิ วิวฏฺฏกปฺเป อเนเกปิ สํวฏฺฏวิวฏฺฏกปฺเป “อมุตฺราสิํ เอวํนาโม เอวํโคตฺโต เอวํวณฺโณ เอวมาหาโร เอวํสุขทุกฺขปฺปฏิสํเวที เอวมายุปริยนฺโต, โส ตโต จุโต อมุตฺร อุทปาทิํ, ตตฺราปาสิํ เอวํนาโม เอวํโคตฺโต เอวํวณฺโณ เอวมาหาโร เอวํสุขทุกฺขปฺปฏิสํเวที เอวมายุปริยนฺโต, โส ตโต จุโต อิธูปปนฺโน”ติ, อิติ สาการํ ส-อุทฺเทสํ อเนกวิหิตํ ปุพฺเพนิวาสํ อนุสฺสรติ.}

\prose{ตํ กิํ มญฺญสิ นิโคฺรธ “ยทิ เอวํ สนฺเต ตโปชิคุจฺฉา ปริสุทฺธา วา โหติ อปริสุทฺธา วา”ติ.\csromanspacer{} อทฺธา โข ภนฺเต เอวํ สนฺเต ตโปชิคุจฺฉา ปริสุทฺธา โหติ โน อปริสุทฺธา, อคฺคปฺปตฺตา จ สารปฺปตฺตา จาติ.\csromanspacer{} น โข นิโคฺรธ เอตฺตาวตา ตโปชิคุจฺฉา อคฺคปฺปตฺตา จ โหติ สารปฺปตฺตา จ, อปิ จ โข เผคฺคุปฺปตฺตา โหตีติ.}

\csromanmatikaanchor{mtk.15}
\csromantocmarkref{section}{ปริสุทฺธ-อคฺคปฺปตฺตสารปฺปตฺตกถา}{mtk.15}
\bookmark[dest={mtk.15},level=1]{ปริสุทฺธ-อคฺคปฺปตฺตสารปฺปตฺตกถา}
\csromanheader{1}{\textbf{ปริสุทฺธ-อคฺคปฺปตฺตสารปฺปตฺตกถา}}
\proseitem{73}{กิตฺตาวตา ปน ภนฺเต ตโปชิคุจฺฉา อคฺคปฺปตฺตา จ โหติ สารปฺปตฺตา จ, สาธุ เม ภนฺเต ภควา ตโปชิคุจฺฉาย อคฺคญฺเญว ปาเปตุ, สารญฺเญว ปาเปตูติ.\csromanspacer{} อิธ นิโคฺรธ ตปสฺสี จาตุยามสํวรสํวุโต โหติ.\csromanspacer{} กถญฺจ นิโคฺรธ ตปสฺสี จาตุยามสํวรสํวุโต โหติ ฯเปฯ.\csromanspacer{} ยโต โข นิโคฺรธ ตปสฺสี จาตุยามสํวรสํวุโต โหติ, อทุํ จสฺส โหติ ตปสฺสิตาย.\csromanspacer{} โส อภิหรติ โน หีนายาวตฺตติ.\csromanspacer{} โส วิวิตฺตํ เสนาสนํ ภชติ ฯเปฯ.\csromanspacer{} โส อิเม ปญฺจ นีวรเณ ปหาย เจตโส อุปกฺกิเลเส ปญฺญาย ทุพฺพลีกรเณ เมตฺตาสหคเตน เจตสา ฯเปฯ อุเปกฺขาสหคเตน เจตสา วิปุเลน มหคฺคเตน อปฺปมาเณน อเวเรน อพฺยาปชฺเชน ผริตฺวา วิหรติ.\csromanspacer{} โส อเนกวิหิตํ ปุพฺเพนิวาสํ อนุสฺสรติ.\csromanspacer{} เสยฺยถิทํ ฯเปฯ เทฺวปิ ชาติโย ติสฺโสปิ ชาติโย จตสฺโสปิ ชาติโย ปญฺจปิ ชาติโย ฯเปฯ อิติ สาการํ ส-อุทฺเทสํ อเนกวิหิตํ ปุพฺเพนิวาสํ อนุสฺสรติ.\csromanspacer{} โส ทิพฺเพน จกฺขุนา วิสุทฺเธน อติกฺกนฺตมานุสเกน สตฺเต ปสฺสติ จวมาเน อุปปชฺชมาเน หีเน ปณีเต สุวณฺเณ ทุพฺพณฺเณ สุคเต ทุคฺคเต, ยถากมฺมูปเค สตฺเต ปชานาติ “อิเม วต โภนฺโต สตฺตา}

\vspace{\tipitakaparskip}
\csromancenter{อุทุมฺพริกสุตฺต กายทุจฺจริเตน สมนฺนาคตา วจีทุจฺจริเตน สมนฺนาคตา มโนทุจฺจริเตน สมนฺนาคตา อริยานํ อุปวาทกา มิจฺฉาทิฏฺฐิกา มิจฺฉาทิฏฺฐิกมฺมสมาทานา, เต กายสฺส เภทา ปรํ มรณา อปายํ ทุคฺคติํ วินิปาตํ นิรยํ อุปปนฺนา.\csromanspacer{} อิเม วา ปน โภนฺโต สตฺตา กายสุจริเตน สมนฺนาคตา วจีสุจริเตน สมนฺนาคตา มโนสุจริเตน สมนฺนาคตา อริยานํ อนุปวาทกา สมฺมาทิฏฺฐิกา สมฺมาทิฏฺฐิกมฺมสมาทานา, เต กายสฺส เภทา ปรํ มรณา สุคติํ สคฺคํ โลกํ อุปปนฺนา”ติ.\csromanspacer{} อิติ ทิพฺเพน จกฺขุนา วิสุทฺเธน อติกฺกนฺตมานุสเกน สตฺเต ปสฺสติ จวมาเน อุปปชฺชมาเน หีเน ปณีเต สุวณฺเณ ทุพฺพณฺเณ สุคเต ทุคฺคเต, ยถากมฺมูปเค สตฺเต ปชานาติ.}
\prose{\csromanfolio{43}ตํ กิํ มญฺญสิ นิโคฺรธ “ยทิ เอวํ สนฺเต ตโปชิคุจฺฉา ปริสุทฺธา วา โหติ อปริสุทฺธา วา”ติ.\csromanspacer{} อทฺธา โข ภนฺเต เอวํ สนฺเต ตโปชิคุจฺฉา ปริสุทฺธา โหติ โน อปริสุทฺธา, อคฺคปฺปตฺตา จ สารปฺปตฺตา จาติ.}

\proseitem{74}{เอตฺตาวตา โข นิโคฺรธ ตโปชิคุจฺฉา อคฺคปฺปตฺตา จ โหติ สารปฺปตฺตา จ.\csromanspacer{} อิติ โข นิโคฺรธ\footnote{อิติ นิโคฺรธ (สฺยา)} ยํ มํ ตฺวํ อวจาสิ “โก นาม โส ภนฺเต ภควโต ธมฺโม, เยน ภควา สาวเก วิเนติ, เยน ภควตา สาวกา วินีตา อสฺสาสปฺปตฺตา ปฏิชานนฺติ อชฺฌาสยํ อาทิพฺรหฺมจริยนฺ”ติ.\csromanspacer{} อิติ โข (ตํ) นิโคฺรธ ฐานํ อุตฺตริตรญฺจ ปณีตตรญฺจ, เยนาหํ สาวเก วิเนมิ, เยน มยา สาวกา วินีตา อสฺสาสปฺปตฺตา ปฏิชานนฺติ อชฺฌาสยํ อาทิพฺรหฺมจริยนฺติ.}

\prose{เอวํ วุตฺเต เต ปริพฺพาชกา อุนฺนาทิโน อุจฺจาสทฺทมหาสทฺทา อเหสุํ “เอตฺถ มยํ อนสฺสาม สาจริยกา, น มยํ อิโต ภิยฺโย อุตฺตริตรํ ปชานามา”ติ.}

\csromanmatikaanchor{mtk.16}
\csromantocmarkref{section}{นิโคฺรธสฺส ปชฺฌายน}{mtk.16}
\bookmark[dest={mtk.16},level=1]{นิโคฺรธสฺส ปชฺฌายน}
\csromanheader{1}{\textbf{นิโคฺรธสฺสปชฺฌายน}}
\proseitem{75}{ยทา อญฺญาสิ สนฺธาโน คหปติ “อญฺญทตฺถุ โขทานิเม อญฺญติตฺถิยา ปริพฺพาชกา ภควโต ภาสิตํ สุสฺสูสนฺติ, โสตํ โอทหนฺติ, อญฺญาจิตฺตํ อุปฏฺฐาเปนฺตี”ติ.\csromanspacer{} อถ\footnote{อถ นํ (ก)} นิโคฺรธํ ปริพฺพาชกํ \csromanfolio{44}เอตทโวจ “อิติ โข ภนฺเต นิโคฺรธ ยํ มํ ตฺวํ อวจาสิ ‘ยคฺเฆ คหปติ ชาเนยฺยาสิ, เกน สมโณ โคตโม สทฺธิํ สลฺลปติ, เกน สากจฺฉํ สมาปชฺชติ, เกน ปญฺญาเวยฺยตฺติยํ สมาปชฺชติ, สุญฺญาคารหตา สมณสฺส โคตมสฺส ปญฺญา, อปริสาวจโร สมโณ โคตโม นาลํ สลฺลาปาย, โส อนฺตมนฺตาเนว เสวติ.\csromanspacer{} เสยฺยถาปิ นาม โคกาณา ปริยนฺตจารินี อนฺตมนฺตาเนว เสวติ.\csromanspacer{} เอวเมว สุญฺญาคารหตา สมณสฺส โคตมสฺส ปญฺญา, อปริสาวจโร สมโณ โคตโม นาลํ สลฺลาปาย, โส อนฺตมนฺตาเนว เสวติ.\csromanspacer{} อิงฺฆ จ คหปติ สมโณ โคตโม อิมํ ปริสํ อาคจฺเฉยฺย, เอกปเญฺหเนว นํ สํสาเทยฺยาม, ตุจฺฉกุมฺภีว นํ มญฺเญ โอโรเธยฺยามา’ติ.\csromanspacer{} อยํ โข โส ภนฺเต ภควา อรหํ สมฺมาสมฺพุทฺโธ อิธานุปฺปตฺโต, อปริสาวจรํ ปน นํ กโรถ, โคกาณํ ปริยนฺตจารินิํ กโรถ, เอกปเญฺหเนว นํ สํสาเทถ, ตุจฺฉกุมฺภีว นํ โอโรเธถา”ติ.\csromanspacer{} เอวํ วุตฺเต นิโคฺรโธ ปริพฺพาชโก ตุณฺหีภูโต มงฺกุภูโต ปตฺตกฺขนฺโธ อโธมุโข ปชฺฌายนฺโต อปฺปฏิภาโน นิสีทิ.}

\proseitem{76}{อถ โข ภควา นิโคฺรธํ ปริพฺพาชกํ ตุณฺหีภูตํ มงฺกุภูตํ ปตฺตกฺขนฺธํ อโธมุขํ ปชฺฌายนฺตํ อปฺปฏิภานํ วิทิตฺวา นิโคฺรธํ ปริพฺพาชกํ เอตทโวจ “สจฺจํ กิร นิโคฺรธ ภาสิตา เต เอสา วาจา”ติ.\csromanspacer{} สจฺจํ ภนฺเต ภาสิตา เม เอสา วาจา ยถาพาเลน ยถามูเฬฺหน ยถา-อกุสเลนาติ.\csromanspacer{} ตํ กิํ มญฺญสิ นิโคฺรธ, กินฺติ เต สุตํ ปริพฺพาชกานํ วุฑฺฒานํ มหลฺลกานํ อาจริยปาจริยานํ ภาสมานานํ “เย เต อเหสุํ อตีตมทฺธานํ อรหนฺโต สมฺมาสมฺพุทฺธา, เอวํ สุ เต ภควนฺโต สํคมฺม สมาคมฺม อุนฺนาทิโน อุจฺจาสทฺทมหาสทฺทา อเนกวิหิตํ ติรจฺฉานกถํ อนุยุตฺตา วิหริํสุ.\csromanspacer{} เสยฺยถิทํ, ราชกถํ โจรกถํ ฯเปฯ อิติภวาภวกถํ อิติ วา.\csromanspacer{} เสยฺยถาปิ ตฺวํ เอตรหิ สาจริยโก.\csromanspacer{} อุทาหุ เอวํ สุ เต ภควนฺโต อรญฺญวนปตฺถานิ ปนฺตานิ เสนาสนานิ ปฏิเสวนฺติ อปฺปสทฺทานิ อปฺปนิคฺโฆสานิ วิชนวาตานิ มนุสฺสราหสฺเสยฺยกานิ ปฏิสลฺลานสารุปฺปานิ, เสยฺยถาปาหํ เอตรหี”ติ.}

\prose{สุตํ เมตํ ภนฺเต ปริพฺพาชกานํ วุฑฺฒานํ มหลฺลกานํ อาจริยปาจริยานํ ภาสมานานํ “เย เต อเหสุํ อตีตมทฺธานํ อรหนฺโต}

\vspace{\tipitakaparskip}
\csromancenter{อุทุมฺพริกสุตฺต สมฺมาสมฺพุทฺธา, น เอวํ สุ\footnote{นาสฺสุ (สี, อิ)} เต ภควนฺโต สํคมฺม สมาคมฺม อุนฺนาทิโน อุจฺจาสทฺทมหาสทฺทา อเนกวิหิตํ ติรจฺฉานกถํ อนุยุตฺตา วิหรนฺติ.\csromanspacer{} เสยฺยถิทํ ราชกถํ โจรกถํ ฯเปฯ อิติภวาภวกถํ อิติ วา, เสยฺยถาปาหํ เอตรหิ สาจริยโก.\csromanspacer{} เอวํ สุ เต ภควนฺโต อรญฺญวนปตฺถานิ ปนฺตานิ เสนาสนานิ ปฏิเสวนฺติ อปฺปสทฺทานิ อปฺปนิคฺโฆสานิ วิชนวาตานิ มนุสฺสราหสฺเสยฺยกานิ ปฏิสลฺลานสารุปฺปานิ, เสยฺยถาปิ ภควา เอตรหี”ติ.}
\prose{\csromanfolio{45}ตสฺส เต นิโคฺรธ วิญฺญุสฺส สโต มหลฺลกสฺส น เอตทโหสิ “พุทฺโธ โส ภควา โพธาย ธมฺมํ เทเสติ, ทนฺโต โส ภควา ทมถาย ธมฺมํ เทเสติ, สนฺโต โส ภควา สมถาย ธมฺมํ เทเสติ, ติณฺโณ โส ภควา ตรณาย ธมฺมํ เทเสติ, ปรินิพฺพุโต โส ภควา ปรินิพฺพานาย ธมฺมํ เทเสตี”ติ.}

\csromanmatikaanchor{mtk.17}
\csromantocmarkref{section}{พฺรหฺมจริยปริโยสานสจฺฉิกิริยา}{mtk.17}
\bookmark[dest={mtk.17},level=1]{พฺรหฺมจริยปริโยสานสจฺฉิกิริยา}
\csromanheader{1}{\textbf{พฺรหฺมจริยปริโยสานสจฺฉิกิริยา}}
\proseitem{77}{เอวํ วุตฺเต นิโคฺรโธ ปริพฺพาชโก ภควนฺตํ เอตทโวจ “อจฺจโย มํ ภนฺเต อจฺจคมา ยถาพาลํ ยถามูฬฺหํ ยถา-อกุสลํ.\csromanspacer{} ยฺวาหํ เอวํ ภควนฺตํ อวจาสิํ, ตสฺส เม ภนฺเต ภควา อจฺจยํ อจฺจยโต ปฏิคฺคณฺหาตุ อายติํ สํวรายา”ติ.\csromanspacer{} ตคฺฆ ตฺวํ\footnote{ตํ (สี, สฺยา, อิ)} นิโคฺรธ อจฺจโย อจฺจคมา ยถาพาลํ ยถามูฬฺหํ ยถา-อกุสลํ, โย มํ ตฺวํ เอวํ อวจาสิ.\csromanspacer{} ยโต จ โข ตฺวํ นิโคฺรธ อจฺจยํ อจฺจยโต ทิสฺวา ยถาธมฺมํ ปฏิกโรสิ, ตํ เต มยํ ปฏิคฺคณฺหาม.\csromanspacer{} วุทฺธิ เหสา นิโคฺรธ อริยสฺส วินเย, โย อจฺจยํ อจฺจยโต ทิสฺวา ยถาธมฺมํ ปฏิกโรติ อายติํ สํวรํ อาปชฺชติ.\csromanspacer{} อหํ โข ปน นิโคฺรธ เอวํ วทามิ\csromandash{}}

\csromanmatikaanchor{mtk.18}
\csromantocmarkref{section}{ปริพฺพาชกานํปชฺฌายน}{mtk.18}
\bookmark[dest={mtk.18},level=1]{ปริพฺพาชกานํปชฺฌายน}
\prose{เอตุ วิญฺญู ปุริโส อสโฐ อมายาวี อุชุชาติโก, อหมนุสาสามิ อหํ ธมฺมํ เทเสมิ.\csromanspacer{} ยถานุสิฏฺฐํ ตถา\footnote{ยถานุสิฏฺฐํ (?)} ปฏิปชฺชมาโน, ยสฺสตฺถาย กุลปุตฺตา สมฺมเทว อคารสฺมา อนคาริยํ ปพฺพชนฺติ, ตทนุตฺตรํ พฺรหฺมจริยปริโยสานํ ทิฏฺเฐว ธมฺเม สยํ อภิญฺญา สจฺฉิกตฺวา อุปสมฺปชฺช วิหริสฺสติ สตฺตวสฺสานิ.\csromanspacer{} ติฏฺฐนฺตุ นิโคฺรธ สตฺตวสฺสานิ.\csromanspacer{} เอตุ วิญฺญู ปุริโส อสโฐ อมายาวี อุชุชาติโก, อหมนุสาสามิ อหํ ธมฺมํ เทเสมิ.\csromanspacer{} ยถานุสิฏฺฐํ ตถา ปฏิปชฺชมาโน, ยสฺสตฺถาย \csromanfolio{46}กุลปุตฺตา สมฺมเทว อคารสฺมา อนคาริยํ ปพฺพชนฺติ, ตทนุตฺตรํ พฺรหฺมจริยปริโยสานํ ทิฏฺเฐว ธมฺเม สยํ อภิญฺญา สจฺฉิกตฺวา อุปสมฺปชฺช วิหริสฺสติ ฉ วสฺสานิ.\csromanspacer{} ปญฺจ วสฺสานิ.\csromanspacer{} จตฺตาริ วสฺสานิ.\csromanspacer{} ตีณิ วสฺสานิ.\csromanspacer{} เทฺว วสฺสานิ.\csromanspacer{} เอกํ วสฺสํ.\csromanspacer{} ติฏฺฐตุ นิโคฺรธ เอกํ วสฺสํ.\csromanspacer{} เอตุ วิญฺญู ปุริโส อสโฐ อมายาวี อุชุชาติโก, อหมนุสาสามิ อหํ ธมฺมํ เทเสมิ.\csromanspacer{} ยถานุสิฏฺฐํ ตถา ปฏิปชฺชมาโน, ยสฺสตฺถาย กุลปุตฺตา สมฺมเทว อคารสฺมา อนคาริยํ ปพฺพชนฺติ, ตทนุตฺตรํ พฺรหฺมจริยปริโยสานํ ทิฏฺเฐว ธมฺเม สยํ อภิญฺญา สจฺฉิกตฺวา อุปสมฺปชฺช วิหริสฺสติ สตฺตมาสานิ.\csromanspacer{} ติฏฺฐนฺตุ นิโคฺรธ สตฺต มาสานิ.\csromanspacer{} ฉ มาสานิ.\csromanspacer{} ปญฺจ มาสานิ.\csromanspacer{} จตฺตาริ มาสานิ.\csromanspacer{} ตีณิ มาสานิ.\csromanspacer{} เทฺว มาสานิ.\csromanspacer{} เอกํ มาสํ.\csromanspacer{} อฑฺฒมาสํ.\csromanspacer{} ติฏฺฐตุ นิโคฺรธ อฑฺฒมาโส, เอตุ วิญฺญู ปุริโส อสโฐ อมายาวี อุชุชาติโก, อหมนุสาสามิ อหํ ธมฺมํ เทเสมิ.\csromanspacer{} ยถานุสิฏฺฐํ ตถา ปฏิปชฺชมาโน, ยสฺสตฺถาย กุลปุตฺตา สมฺมเทว อคารสฺมา อนคาริยํ ปพฺพชนฺติ, ตทนุตฺตรํ พฺรหฺมจริยปริโยสานํ ทิฏฺเฐว ธมฺเม สยํ อภิญฺญา สจฺฉิกตฺวา อุปสมฺปชฺช วิหริสฺสติ สตฺตาหํ.}

\csromanheader{2}{\textbf{ปริพฺพาชกานํ ปชฺฌายน}}
\proseitem{78}{สิยา โข ปเนตํ นิโคฺรธ เอวมสฺส “อนฺเตวาสิกมฺยตา โน สมโณ โคตโม เอวมาหา”ติ.\csromanspacer{} น โข ปเนตํ นิโคฺรธ เอวํ ทฏฺฐพฺพํ.\csromanspacer{} โย เอว โว อาจริโย, โส เอว โว อาจริโย โหตุ.\csromanspacer{} สิยา โข ปน เต นิโคฺรธ เอวมสฺส “อุทฺเทสา โน จาเวตุกาโม สมโณ โคตโม เอวมาหา”ติ.\csromanspacer{} น โข ปเนตํ นิโคฺรธ เอวํ ทฏฺฐพฺพํ, โย เอว โว อุทฺเทโส, โส เอว โว อุทฺเทโส โหตุ.\csromanspacer{} สิยา โข ปน เต นิโคฺรธ เอวมสฺส “อาชีวา โน จาเวตุกาโม สมโณ โคตโม เอวมาหา”ติ.\csromanspacer{} น โข ปเนตํ นิโคฺรธ เอวํ ทฏฺฐพฺพํ.\csromanspacer{} โย เอว โว อาชีโว, โส เอว โว อาชีโว โหตุ.\csromanspacer{} สิยา โข ปน เต นิโคฺรธ เอวมสฺส “เย โน ธมฺมา อกุสลา อกุสลสงฺขาตา สาจริยกานํ, เตสุ ปติฏฺฐาเปตุกาโม สมโณ โคตโม เอวมาหา”ติ.\csromanspacer{} น โข ปเนตํ นิโคฺรธ เอวํ ทฏฺฐพฺพํ.\csromanspacer{} อกุสลา เจว โว ธมฺมา\footnote{โว เต ธมฺมา((สี, อิ)} โหนฺตุ อกุสลสงฺขาตา จ สาจริยกานํ.\csromanspacer{} สิยา โข ปน}

\vspace{\tipitakaparskip}
\csromancenter{อุทุมฺพริกสุตฺต เต นิโคฺรธ เอวมสฺส “เย โน ธมฺมา กุสลา กุสลสงฺขาตา สาจริยกานํ, เตหิ วิเวเจตุกาโม สมโณ โคตโม เอวมาหา”ติ.\csromanspacer{} น โข ปเนตํ นิโคฺรธ เอวํ ทฏฺฐพฺพํ.\csromanspacer{} กุสลา เจว โว ธมฺมา\footnote{โว เต ธมฺมา (สี, อิ)} โหนฺตุ กุสลสงฺขาตา จ สาจริยกานํ.\csromanspacer{} อิติ ขฺวาหํ นิโคฺรธ เนว อนฺเตวาสิกมฺยตา เอวํ วทามิ, นปิ อุทฺเทสา จาเวตุกาโม เอวํ วทามิ, นปิ อาชีวา จาเวตุกาโม เอวํ วทามิ, นปิ เย จ โว ธมฺมา\footnote{นปิ เย โข ธมฺมา (สี), นปิ เย เต ธมฺมา (สฺยา), นปิ เย โว ธมฺมา (อิ)} อกุสลา อกุสลสงฺขาตา สาจริยกานํ, เตสุ ปติฏฺฐาเปตุกาโม เอวํ วทามิ, นปิ เย จ โว ธมฺมา2 กุสลา กุสลสงฺขาตา สาจริยกานํ, เตหิ วิเวเจตุกาโม เอวํ วทามิ.\csromanspacer{} สนฺติ จ โข นิโคฺรธ อกุสลา ธมฺมา อปฺปหีนา สํกิเลสิกา โปโนพฺภวิกา\footnote{โปโนภวิกา (ก)} สทรา\footnote{สทฺทรา (อิ, ก), สทรถา (สฺยา, ก)} ทุกฺขวิปากา อายติํ ชาติชรามรณิยา, เยสาหํ ปหานาย ธมฺมํ เทเสมิ.\csromanspacer{} ยถาปฏิปนฺนานํ โว สํกิเลสิกา ธมฺมา ปหียิสฺสนฺติ.\csromanspacer{} โวทานียา ธมฺมา อภิวฑฺฒิสฺสนฺติ ปญฺญาปาริปูริํ เวปุลฺลตฺตญฺจ ทิฏฺเฐว ธมฺเม สยํ อภิญฺญา สจฺฉิกตฺวา อุปสมฺปชฺช วิหริสฺสถา”ติ.}
\proseitem{79}{\csromanfolio{47}เอวํ วุตฺเต เต ปริพฺพาชกา ตุณฺหีภูตา มงฺกุภูตา ปตฺตกฺขนฺธา อโธมุขา ปชฺฌายนฺตา อปฺปฏิภานา นิสีทิํสุ ยถา ตํ มาเรน ปริยุฏฺฐิตจิตฺตา.\csromanspacer{} อถ โข ภควโต เอตทโหสิ “สพฺเพปิเม โมฆปุริสา ผุฏฺฐา ปาปิมตา.\csromanspacer{} ยตฺร หิ นาม เอกสฺสปิ น เอวํ ภวิสฺสติ ‘หนฺท มยํ อญฺญาณตฺถมฺปิ สมเณ โคตเม พฺรหฺมจริยํ จราม, กิํ กริสฺสติ สตฺตาโห’ติ”.\csromanspacer{} อถ โข ภควา อุทุมฺพริกาย ปริพฺพาชการาเม สีหนาทํ นทิตฺวา เวหาสํ อพฺภุคฺคนฺตฺวา คิชฺฌกูเฏ ปพฺพเต ปจฺจุปฏฺฐาสิ\footnote{ปจฺจุฏฺฐาสิ (สี, สฺยา, อิ)}.\csromanspacer{} สนฺธาโน ปน คหปติ ตาวเทว ราชคหํ ปาวิสีติ.}

\nitthitam{\textbf{อุทุมฺพริกสุตฺตํ นิฏฺฐิตํ ทุติยํ.}}
\csromanmatikaanchor{mtk.19}
\csromantocmarknopage{chapter}{3. จกฺกวตฺติสุตฺต}{mtk.19}
\bookmark[dest={mtk.19},level=0]{3. จกฺกวตฺติสุตฺต}
\chapterheadpageread{3. จกฺกวตฺติสุตฺต}
\csromanmatikaanchor{mtk.20}
\csromantocmarkref{section}{อตฺตทีปสรณตา}{mtk.20}
\bookmark[dest={mtk.20},level=1]{อตฺตทีปสรณตา}
\csromanheader{1}{\textbf{อตฺตทีปสรณตา}}
\proseitem{80}{\csromanfolio{48}เอวํ เม สุตํ\csromandash{}เอกํ สมยํ ภควา มคเธสุ วิหรติ มาตุลายํ.\csromanspacer{} ตตฺร โข ภควา ภิกฺขู อามนฺเตสิ “ภิกฺขโว”ติ.\csromanspacer{} ภทฺทนฺเตติ เต ภิกฺขู ภควโต ปจฺจสฺโสสุํ.\csromanspacer{} ภควา เอตทโวจ “อตฺตทีปา ภิกฺขเว วิหรถ อตฺตสรณา อนญฺญสรณา ธมฺมทีปา ธมฺมสรณา อนญฺญสรณา.\csromanspacer{} กถญฺจ ปน ภิกฺขเว ภิกฺขุ อตฺตทีโป วิหรติ อตฺตสรโณ อนญฺญสรโณ ธมฺมทีโป ธมฺมสรโณ อนญฺญสรโณ.\csromanspacer{} อิธ ภิกฺขเว ภิกฺขุ กาเย กายานุปสฺสี วิหรติ อาตาปี สมฺปชาโน สติมา วิเนยฺย โลเก อภิชฺฌาโทมนสฺสํ.\csromanspacer{} เวทนาสุ เวทนานุปสฺสี ฯเปฯ.\csromanspacer{} จิตฺเต จิตฺตานุปสฺสี ฯเปฯ.\csromanspacer{} ธมฺเมสุ ธมฺมานุปสฺสี วิหรติ อาตาปี สมฺปชาโน สติมา เวเนยฺย โลเก อภิชฺฌาโทมนสฺสํ.\csromanspacer{} เอวํ โข ภิกฺขเว ภิกฺขุ อตฺตทีโป วิหรติ อตฺตสรโณ อนญฺญสรโณ ธมฺมทีโป ธมฺมสรโณ อนญฺญสรโณ.}

\prose{โคจเร ภิกฺขเว จรถ สเก เปตฺติเก วิสเย.\csromanspacer{} โคจเร ภิกฺขเว จรตํ สเก เปตฺติเก วิสเย น ลจฺฉติ มาโร โอตารํ, น ลจฺฉติ มาโร อารมฺมณํ\footnote{อารมณํ (?)}.\csromanspacer{} กุสลานํ ภิกฺขเว ธมฺมานํ สมาทานเหตุ เอวมิทํ ปุญฺญํ ปวฑฺฒติ.}

\csromanmatikaanchor{mtk.21}
\csromantocmarkref{section}{ทฬฺหเนมิจกฺกวตฺติราชา}{mtk.21}
\bookmark[dest={mtk.21},level=1]{ทฬฺหเนมิจกฺกวตฺติราชา}
\csromanheader{1}{\textbf{ทฬฺหเนมิจกฺกวตฺติราชา}}
\proseitem{81}{ภูตปุพฺพํ ภิกฺขเว ราชา ทฬฺหเนมิ นาม อโหสิ จกฺกวตฺตี\footnote{จกฺกวตฺติ (สฺยา, อิ)} ธมฺมิโก ธมฺมราชา จาตุรนฺโต วิชิตาวี ชนปทตฺถาวริยปฺปตฺโต สตฺตรตนสมนฺนาคโต.\csromanspacer{} ตสฺสิมานิ สตฺต รตนานิ อเหสุํ.\csromanspacer{} เสยฺยถิทํ, จกฺกรตนํ หตฺถิรตนํ อสฺสรตนํ มณิรตนํ อิตฺถิรตนํ คหปติรตนํ ปริณายกรตนเมว สตฺตมํ.\csromanspacer{} ปโรสหสฺสํ โข ปนสฺส ปุตฺตา อเหสุํ สูรา วีรงฺครูปา ปรเสนปฺปมทฺทนา, โส อิมํ ปถวิํ สาครปริยนฺตํ อทณฺเฑน อสตฺเถน ธมฺเมน\footnote{ธมฺเมน สเมน (สฺยา, ก)} อภิวิชิย อชฺฌาวสิ.}

\csromanheader{2}{3. จกฺกวตฺติสุตฺต}
\proseitem{82}{\csromanfolio{49}อถ โข ภิกฺขเว ราชา ทฬฺหเนมิ พหุนฺนํ วสฺสานํ พหุนฺนํ วสฺสสตานํ พหุนฺนํ วสฺสสหสฺสานํ อจฺจเยน อญฺญตรํ ปุริสํ อามนฺเตสิ “ยทา ตฺวํ อมฺโภ ปุริส ปสฺเสยฺยาสิ ทิพฺพํ จกฺกรตนํ โอสกฺกิตํ ฐานา จุตํ, อถ เม อาโรเจยฺยาสี”ติ. “เอวํ เทวา”ติ โข ภิกฺขเว โส ปุริโส รญฺโญ ทฬฺหเนมิสฺส ปจฺจสฺโสสิ.\csromanspacer{} อทฺทสา โข ภิกฺขเว โส ปุริโส พหุนฺนํ วสฺสานํ พหุนฺนํ วสฺสสตานํ พหุนฺนํ วสฺสสหสฺสานํ อจฺจเยน ทิพฺพํ จกฺกรตนํ โอสกฺกิตํ ฐานา จุตํ, ทิสฺวาน เยน ราชา ทฬฺหเนมิ เตนุปสงฺกมิ, อุปสงฺกมิตฺวา ราชานํ ทฬฺหเนมิํ เอตทโวจ “ยคฺเฆ เทว ชาเนยฺยาสิ, ทิพฺพํ เต จกฺกรตนํ โอสกฺกิตํ ฐานา จุตนฺ”ติ.\csromanspacer{} อถ โข ภิกฺขเว ราชา ทฬฺหเนมิ เชฏฺฐปุตฺตํ กุมารํ อามนฺตาเปตฺวา\footnote{อามนฺเตตฺวา (สฺยา, ก)} เอตทโวจ “ทิพฺพํ กิร เม ตาต กุมาร จกฺกรตนํ โอสกฺกิตํ ฐานา จุตํ.\csromanspacer{} สุตํ โข ปน เมตํ ‘ยสฺส รญฺโญ จกฺกวตฺติสฺส ทิพฺพํ จกฺกรตนํ โอสกฺกติ ฐานา จวติ, น ทานิ เตน รญฺญา จิรํ ชีวิตพฺพํ โหตี’ติ.\csromanspacer{} ภุตฺตา โข ปน เม มานุสกา กามา, สมโยทานิ เม ทิพฺเพ กาเม ปริเยสิตุํ.\csromanspacer{} เอหิ ตฺวํ ตาต กุมาร อิมํ สมุทฺทปริยนฺตํ ปถวิํ ปฏิปชฺช, อหํ ปน เกสมสฺสุํ โอหาเรตฺวา กาสายานิ วตฺถานิ อจฺฉาเทตฺวา อคารสฺมา อนคาริยํ ปพฺพชิสฺสามี”ติ.}

\proseitem{83}{อถ โข ภิกฺขเว ราชา ทฬฺหเนมิ เชฏฺฐปุตฺตํ กุมารํ สาธุกํ รชฺเช สมนุสาสิตฺวา เกสมสฺสุํ โอหาเรตฺวา กาสายานิ วตฺถานิ อจฺฉาเทตฺวา อคารสฺมา อนคาริยํ ปพฺพชิ, สตฺตาหปพฺพชิเต โข ปน ภิกฺขเว ราชิสิมฺหิ ทิพฺพํ จกฺกรตนํ อนฺตรธายิ.}

\prose{อถ โข ภิกฺขเว อญฺญตโร ปุริโส เยน ราชา ขตฺติโย มุทฺธาภิสิตฺโต\footnote{มุทฺธาวสิตฺโต (สี, สฺยา, อิ) เอวมุปริปิ} เตนุปสงฺกมิ, อุปสงฺกมิตฺวา ราชานํ ขตฺติยํ มุทฺธาภิสิตฺตํ เอตทโวจ “ยคฺเฆ เทว ชาเนยฺยาสิ ทิพฺพํ จกฺกรตนํ อนฺตรหิตนฺ”ติ.\csromanspacer{} อถ โข ภิกฺขเว ราชา ขตฺติโย มุทฺธาภิสิตฺโต ทิพฺเพ จกฺกรตเน อนฺตรหิเต อนตฺตมโน อโหสิ, อนตฺตมนตญฺจ ปฏิสํเวเทสิ.\csromanspacer{} โส เยน ราชิสิ เตนุปสงฺกมิ, อุปสงฺกมิตฺวา ราชิสิํ เอตทโวจ “ยคฺเฆ เทว ชาเนยฺยาสิ ทิพฺพํ จกฺกรตนํ อนฺตรหิตนฺ”ติ.\csromanspacer{} เอวํ วุตฺเต ภิกฺขเว ราชิสิ ราชานํ ขตฺติยํ มุทฺธาภิสิตฺตํ เอตทโวจ “มา โข ตฺวํ ตาต \csromanfolio{50}ทิพฺเพ จกฺกรตเน อนฺตรหิเต อนตฺตมโน อโหสิ, มา อนตฺตมนตญฺจ ปฏิสํเวเทสิ, น หิ เต ตาต ทิพฺพํ จกฺกรตนํ เปตฺติกํ ทายชฺชํ, อิงฺฆ ตฺวํ ตาต อริเย จกฺกวตฺติวตฺเต วตฺตาหิ.\csromanspacer{} ฐานํ โข ปเนตํ วิชฺชติ, ยํ เต อริเย จกฺกวตฺติวตฺเต วตฺตมานสฺส ตทหุโปสเถ ปนฺนรเส สีสํนฺหาตสฺส\footnote{สีสํ นหาตสฺส (สี, อิ), สีสนฺหาตสฺส (สฺยา)} อุโปสถิกสฺส อุปริปาสาทวรคตสฺส ทิพฺพํ จกฺกรตนํ ปาตุภวิสฺสติ สหสฺสารํ สเนมิกํ สนาภิกํ สพฺพาการปริปูรนฺ”ติ.}

\csromanmatikaanchor{mtk.22}
\csromantocmarkref{section}{จกฺกวตฺติ-อริยวตฺต}{mtk.22}
\bookmark[dest={mtk.22},level=1]{จกฺกวตฺติ-อริยวตฺต}
\csromanheader{1}{\textbf{จกฺกวตฺติ-อริยวตฺต}}
\proseitem{84}{กตมํ ปน ตํ เทว อริยํ จกฺกวตฺติวตฺตนฺติ.\csromanspacer{} เตน หิ ตฺวํ ตาต ธมฺมํเยว นิสฺสาย ธมฺมํ สกฺกโรนฺโต ธมฺมํ ครุํ กโรนฺโต\footnote{ครุกโรนฺโต (สี, สฺยา, อิ)} ธมฺมํ มาเนนฺโต ธมฺมํ ปูเชนฺโต ธมฺมํ อปจายมาโน ธมฺมทฺธโช ธมฺมเกตุ ธมฺมาธิปเตยฺโย ธมฺมิกํ รกฺขาวรณคุตฺติํ สํวิทหสฺสุ อนฺโตชนสฺมิํ พลกายสฺมิํ ขตฺติเยสุ อนุยนฺเตสุ\footnote{อนุยุตฺเตสุ (สี, อิ)} พฺราหฺมณคหปติเกสุ เนคมชานปเทสุ สมณพฺราหฺมเณสุ มิคปกฺขีสุ.\csromanspacer{} มา จ เต ตาต วิชิเต อธมฺมกาโร ปวตฺติตฺถ.\csromanspacer{} เย จ เต ตาต วิชิเต อธนา อสฺสุ, เตสญฺจ ธนมนุปฺปเทยฺยาสิ\footnote{ธนมนุปฺปทชฺเชยฺยาสิ (สี, สฺยา, อิ)}.\csromanspacer{} เย จ เต ตาต วิชิเต สมณพฺราหฺมณา มทปฺปมาทา ปฏิวิรตา ขนฺติโสรจฺเจ นิวิฏฺฐา เอกมตฺตานํ ทเมนฺติ, เอกมตฺตานํ สเมนฺติ, เอกมตฺตานํ ปรินิพฺพาเปนฺติ.\csromanspacer{} เต กาเลน กาลํ อุปสงฺกมิตฺวา ปริปุจฺเฉยฺยาสิ ปริคฺคเณฺหยฺยาสิ “กิํ ภนฺเต กุสลํ กิํ อกุสลํ, กิํ สาวชฺชํ กิํ อนวชฺชํ, กิํ เสวิตพฺพํ กิํ น เสวิตพฺพํ, กิํ เม กรียมานํ ทีฆรตฺตํ อหิตาย ทุกฺขาย อสฺส, กิํ วา ปน เม กรียมานํ ทีฆรตฺตํ หิตาย สุขาย อสฺสา”ติ.\csromanspacer{} เตสํ สุตฺวา ยํ อกุสลํ ตํ อภินิวชฺเชยฺยาสิ, ยํ กุสลํ ตํ สมาทาย วตฺเตยฺยาสิ.\csromanspacer{} อิทํ โข ตาต ตํ อริยํ จกฺกวตฺติวตฺตนฺติ.}

\csromanmatikaanchor{mtk.23}
\csromantocmarkref{section}{จกฺกรตนปาตุภาว}{mtk.23}
\bookmark[dest={mtk.23},level=1]{จกฺกรตนปาตุภาว}
\csromanheader{1}{\textbf{จกฺกรตนปาตุภาว}}
\proseitem{85}{“เอวํ เทวา”ติ โข ภิกฺขเว ราชา ขตฺติโย มุทฺธาภิสิตฺโต ราชิสิสฺส ปฏิสฺสุตฺวา อริเย จกฺกวตฺติวตฺเต\footnote{อริยํ จกฺกวตฺติวตฺตํ (ก)} วตฺติ.\csromanspacer{} ตสฺส อริเย จกฺกวตฺติวตฺเต วตฺตมานสฺส ตทหุโปสเถ ปนฺนรเส สีสํนฺหาตสฺส}

\vspace{\tipitakaparskip}
\csromancenter{จกฺกวตฺติสุตฺต อุโปสถิกสฺส อุปริปาสาทวรคตสฺส ทิพฺพํ จกฺกรตนํ ปาตุรโหสิ สหสฺสารํ สเนมิกํ สนาภิกํ สพฺพาการปริปูรํ.\csromanspacer{} ทิสฺวาน รญฺโญ ขตฺติยสฺส มุทฺธาภิสิตฺตสฺส เอตทโหสิ “สุตํ โข ปน เมตํ ‘ยสฺส รญฺโญ ขตฺติยสฺส มุทฺธาภิสิตฺตสฺส ตทหุโปสเถ ปนฺนรเส สีสํนฺหาตสฺส อุโปสถิกสฺส อุปริปาสาทวรคตสฺส ทิพฺพํ จกฺกรตนํ ปาตุภวติ สหสฺสารํ สเนมิกํ สนาภิกํ สพฺพาการปริปูรํ, โส โหติ ราชา จกฺกวตฺตี’ติ, อสฺสํ นุ โข อหํ ราชา จกฺกวตฺตี”ติ.}
\prose{\csromanfolio{51}อถ โข ภิกฺขเว ราชา ขตฺติโย มุทฺธาภิสิตฺโต อุฏฺฐายาสนา เอกํสํ อุตฺตราสงฺคํ กริตฺวา วาเมน หตฺเถน ภิงฺการํ คเหตฺวา ทกฺขิเณน หตฺเถน จกฺกรตนํ อพฺภุกฺกิริ “ปวตฺตตุ ภวํ จกฺกรตนํ, อภิวิชินาตุ ภวํ จกฺกรตนนฺ”ติ.}

\prose{อถ โข ตํ ภิกฺขเว จกฺกรตนํ ปุรตฺถิมํ ทิสํ ปวตฺติ, อนฺวเทว ราชา จกฺกวตฺตี สทฺธิํ จตุรงฺคินิยา เสนาย.\csromanspacer{} ยสฺมิํ โข ปน ภิกฺขเว ปเทเส จกฺกรตนํ ปติฏฺฐาสิ, ตตฺถ ราชา จกฺกวตฺตี วาสํ อุปคจฺฉิ สทฺธิํ จตุรงฺคินิยา เสนาย.\csromanspacer{} เย โข ปน ภิกฺขเว ปุรตฺถิมาย ทิสาย ปฏิราชาโน, เต ราชานํ จกฺกวตฺติํ อุปสงฺกมิตฺวา เอวมาหํสุ “เอหิ โข มหาราช, สฺวาคตํ เต\footnote{สาคตํ (สี, อิ)} มหาราช, สกํ เต มหาราช, อนุสาส มหาราชา”ติ.\csromanspacer{} ราชา จกฺกวตฺตี เอวมาห “ปาโณ น หนฺตพฺโพ, อทินฺนํ นาทาตพฺพํ, กาเมสุมิจฺฉา น จริตพฺพา, มุสา น ภาสิตพฺพา, มชฺชํ น ปาตพฺพํ, ยถาภุตฺตญฺจ ภุญฺชถา”ติ.\csromanspacer{} เย โข ปน ภิกฺขเว ปุรตฺถิมาย ทิสาย ปฏิราชาโน, เต รญฺโญ จกฺกวตฺติสฺส อนุยนฺตา\footnote{อนุยุตฺตา (สี, อิ)} อเหสุํ.}

\proseitem{86}{อถ โข ตํ ภิกฺขเว จกฺกรตนํ ปุรตฺถิมํ สมุทฺทํ อชฺโฌคาเหตฺวา\footnote{อชฺโฌคเหตฺวา (สี, สฺยา, อิ)} ปจฺจุตฺตริตฺวา ทกฺขิณํ ทิสํ ปวตฺติ ฯเปฯ ทกฺขิณํ สมุทฺทํ อชฺโฌคาเหตฺวา ปจฺจุตฺตริตฺวา ปจฺฉิมํ ทิสํ ปวตฺติ, อนฺวเทว ราชา จกฺกวตฺตี สทฺธิํ จตุรงฺคินิยา เสนาย.\csromanspacer{} ยสฺมิํ โข ปน ภิกฺขเว ปเทเส ทิพฺพํ จกฺกรตนํ ปติฏฺฐาสิ, ตตฺถ ราชา จกฺกวตฺตี วาสํ อุปคจฺฉิ สทฺธิํ จตุรงฺคินิยา เสนาย.\csromanspacer{} เย โข ปน ภิกฺขเว ปจฺฉิมาย ทิสาย ปฏิราชาโน, เต ราชานํ จกฺกวตฺติํ อุปสงฺกมิตฺวา เอวมาหํสุ “เอหิ โข มหาราช, สฺวาคตํ เต มหาราช, สกํ เต มหาราช, อนุสาส มหาราชา”ติ.\csromanspacer{} ราชา \csromanfolio{52}จกฺกวตฺตี เอวมาห “ปาโณ น หนฺตพฺโพ, อทินฺนํ นาทาตพฺพํ, กาเมสุมิจฺฉา น จริตพฺพา, มุสา น ภาสิตพฺพา, มชฺชํ น ปาตพฺพํ, ยถาภุตฺตญฺจ ภุญฺชถา”ติ.\csromanspacer{} เย โข ปน ภิกฺขเว ปจฺฉิมาย ทิสาย ปฏิราชาโน, เต รญฺโญ จกฺกวตฺติสฺส อนุยนฺตา อเหสุํ.}

\proseitem{87}{อถ โข ตํ ภิกฺขเว จกฺกรตนํ ปจฺฉิมํ สมุทฺทํ อชฺโฌคาเหตฺวา ปจฺจุตฺตริตฺวา อุตฺตรํ ทิสํ ปวตฺติ, อนฺวเทว ราชา จกฺกวตฺตี สทฺธิํ จตุรงฺคินิยา เสนาย.\csromanspacer{} ยสฺมิํ โข ปน ภิกฺขเว ปเทเส ทิพฺพํ จกฺกรตนํ ปติฏฺฐาสิ, ตตฺถ ราชา จกฺกวตฺตี วาสํ อุปคจฺฉิ สทฺธิํ จตุรงฺคินิยา เสนาย.\csromanspacer{} เย โข ปน ภิกฺขเว อุตฺตราย ทิสาย ปฏิราชาโน, เต ราชานํ จกฺกวตฺติํ อุปสงฺกมิตฺวา เอวมาหํสุ “เอหิ โข มหาราช, สฺวาคตํ เต มหาราช, สกํ เต มหาราช, อนุสาส มหาราชา”ติ.\csromanspacer{} ราชา จกฺกวตฺตี เอวมาห “ปาโณ น หนฺตพฺโพ, อทินฺนํ นาทาตพฺพํ, กาเมสุมิจฺฉา น จริตพฺพา, มุสา น ภาสิตพฺพา, มชฺชํ น ปาตพฺพํ, ยถาภุตฺตญฺจ ภุญฺชถา”ติ.\csromanspacer{} เย โข ปน ภิกฺขเว อุตฺตราย ทิสาย ปฏิราชาโน, เต รญฺโญ จกฺกวตฺติสฺส อนุยนฺตา อเหสุํ.}

\prose{อถ โข ตํ ภิกฺขเว จกฺกรตนํ สมุทฺทปริยนฺตํ ปถวิํ อภิวิชินิตฺวา ตเมว ราชธานิํ ปจฺจาคนฺตฺวา รญฺโญ จกฺกวตฺติสฺส อนฺเตปุรทฺวาเร อตฺถกรณปมุเข\footnote{อฑฺฑกรณปมุเข (ก)} อกฺขาหตํ มญฺเญ อฏฺฐาสิ รญฺโญ จกฺกวตฺติสฺส อนฺเตปุรํ อุปโสภยมานํ.}

\csromanmatikaanchor{mtk.24}
\csromantocmarkref{section}{ทุติยาทิจกฺกวตฺติกถา}{mtk.24}
\bookmark[dest={mtk.24},level=1]{ทุติยาทิจกฺกวตฺติกถา}
\csromanheader{1}{\textbf{ทุติยาทิจกฺกวตฺติกถา}}
\proseitem{88}{ทุติโยปิ โข ภิกฺขเว ราชา จกฺกวตฺตี ฯเปฯ.\csromanspacer{} ตติโยปิ โข ภิกฺขเว ราชา จกฺกวตฺตี.\csromanspacer{} จตุตฺโถปิ โข ภิกฺขเว ราชา จกฺกวตฺตี.\csromanspacer{} ปญฺจโมปิ โข ภิกฺขเว ราชา จกฺกวตฺตี.\csromanspacer{} ฉฏฺโฐปิ โข ภิกฺขเว ราชา จกฺกวตฺตี.\csromanspacer{} สตฺตโมปิ โข ภิกฺขเว ราชา จกฺกวตฺตี พหุนฺนํ วสฺสานํ พหุนฺนํ วสฺสสตานํ พหุนฺนํ วสฺสสหสฺสานํ อจฺจเยน อญฺญตรํ ปุริสํ อามนฺเตสิ “ยทา ตฺวํ อมฺโภ ปุริส ปสฺเสยฺยาสิ ทิพฺพํ จกฺกรตนํ โอสกฺกิตํ ฐานา จุตํ, อถ เม อาโรเจยฺยาสี”ติ. “เอวํ เทวา”ติ โข ภิกฺขเว โส ปุริโส รญฺโญ จกฺกวตฺติสฺส ปจฺจสฺโสสิ.\csromanspacer{} อทฺทสา}

\vspace{\tipitakaparskip}
\csromancenter{จกฺกวตฺติสุตฺต โข ภิกฺขเว โส ปุริโส พหุนฺนํ วสฺสานํ พหุนฺนํ วสฺสสตานํ พหุนฺนํ วสฺสสหสฺสานํ อจฺจเยน ทิพฺพํ จกฺกรตนํ โอสกฺกิตํ ฐานา จุตํ, ทิสฺวาน เยน ราชา จกฺกวตฺตี เตนุปสงฺกมิ, อุปสงฺกมิตฺวา ราชานํ จกฺกวตฺติํ เอตทโวจ “ยคฺเฆ เทว ชาเนยฺยาสิ, ทิพฺพํ เต จกฺกรตนํ โอสกฺกิตํ ฐานา จุตนฺ”ติ.}
\proseitem{89}{\csromanfolio{53}อถ โข ภิกฺขเว ราชา จกฺกวตฺตี เชฏฺฐปุตฺตํ กุมารํ อามนฺตาเปตฺวา เอตทโวจ “ทิพฺพํ กิร เม ตาต กุมาร จกฺกรตนํ โอสกฺกิตํ ฐานา จุตํ, สุตํ โข ปน เมตํ ‘ยสฺส รญฺโญ จกฺกวตฺติสฺส ทิพฺพํ จกฺกรตนํ โอสกฺกติ, ฐานา จวติ, น ทานิ เตน รญฺญา จิรํ ชีวิตพฺพํ โหตี’ติ, ภุตฺตา โข ปน เม มานุสกา กามา, สมโยทานิ เม ทิพฺเพ กาเม ปริเยสิตุํ, เอหิ ตฺวํ ตาต กุมาร อิมํ สมุทฺทปริยนฺตํ ปถวิํ ปฏิปชฺช, อหํ ปน เกสมสฺสุํ โอหาเรตฺวา กาสายานิ วตฺถานิ อจฺฉาเทตฺวา อคารสฺมา อนคาริยํ ปพฺพชิสฺสามี”ติ.}

\prose{อถ โข ภิกฺขเว ราชา จกฺกวตฺตี เชฏฺฐปุตฺตํ กุมารํ สาธุกํ รชฺเช สมนุสาสิตฺวา เกสมสฺสุํ โอหาเรตฺวา กาสายานิ วตฺถานิ อจฺฉาเทตฺวา อคารสฺมา อนคาริยํ ปพฺพชิ.\csromanspacer{} สตฺตาหปพฺพชิเต โข ปน ภิกฺขเว ราชิสิมฺหิ ทิพฺพํ จกฺกรตนํ อนฺตรธายิ.}

\proseitem{90}{อถ โข ภิกฺขเว อญฺญตโร ปุริโส เยน ราชา ขตฺติโย มุทฺธาภิสิตฺโต เตนุปสงฺกมิ, อุปสงฺกมิตฺวา ราชานํ ขตฺติยํ มุทฺธาภิสิตฺตํ เอตทโวจ “ยคฺเฆ เทว ชาเนยฺยาสิ ทิพฺพํ จกฺกรตนํ อนฺตรหิตนฺ”ติ.\csromanspacer{} อถ โข ภิกฺขเว ราชา ขตฺติโย มุทฺธาภิสิตฺโต ทิพฺเพ จกฺกรตเน อนฺตรหิเต อนตฺตมโน อโหสิ, อนตฺตมนตญฺจ ปฏิสํเวเทสิ, โน จ โข ราชิสิํ อุปสงฺกมิตฺวา อริยํ จกฺกวตฺติวตฺตํ ปุจฺฉิ, โส สมเตเนว สุทํ ชนปทํ ปสาสติ, ตสฺส สมเตน ชนปทํ ปสาสโต ปุพฺเพนาปรํ ชนปทา น ปพฺพนฺติ, ยถา ตํ ปุพฺพกานํ ราชูนํ อริเย จกฺกวตฺติวตฺเต วตฺตมานานํ.}

\prose{อถ โข ภิกฺขเว อมจฺจา ปาริสชฺชา คณกมหามตฺตา อนีกฏฺฐา โทวาริกา มนฺตสฺสาชีวิโน สนฺนิปติตฺวา ราชานํ ขตฺติยํ มุทฺธาภิสิตฺตํ เอตทโวจุํ “น โข เต เทว สมเตน สุทํ ชนปทํ ปสาสโต \csromanfolio{54}ปุพฺเพนาปรํ ชนปทา ปพฺพนฺติ, ยถา ตํ ปุพฺพกานํ ราชูนํ อริเย จกฺกวตฺติวตฺเต วตฺตมานานํ, สํวิชฺชนฺติ โข เต เทว วิชิเต อมจฺจา ปาริสชฺชา คณกมหามตฺตา อนีกฏฺฐา โทวาริกา มนฺตสฺสาชีวิโน มยญฺเจว อญฺเญ จ\footnote{อญฺเญ จ ปณฺฑิเต สมณพฺราหฺมเณ ปุจฺเฉยฺยาสิ (ก)} เย มยํ อริยํ จกฺกวตฺติวตฺตํ ธาเรม, อิงฺฆ ตฺวํ เทว อเมฺห อริยํ จกฺกวตฺติวตฺตํ ปุจฺฉ, ตสฺส เต มยํ อริยํ จกฺกวตฺติวตฺตํ ปุฏฺฐา พฺยากริสฺสามา”ติ.}

\csromanmatikaanchor{mtk.25}
\csromantocmarkref{section}{อายุวณฺณาทิปริหานิกถา}{mtk.25}
\bookmark[dest={mtk.25},level=1]{อายุวณฺณาทิปริหานิกถา}
\csromanheader{1}{\textbf{อายุวณฺณาทิปริหานิกถา}}
\proseitem{91}{อถ โข ภิกฺขเว ราชา ขตฺติโย มุทฺธาภิสิตฺโต อมจฺเจ ปาริสชฺเช คณกมหามตฺเต อนีกฏฺเฐ โทวาริเก มนฺตสฺสาชีวิโน สนฺนิปาเตตฺวา อริยํ จกฺกวตฺติวตฺตํ ปุจฺฉิ.\csromanspacer{} ตสฺส เต อริยํ จกฺกวตฺติวตฺตํ ปุฏฺฐา พฺยากริํสุ.\csromanspacer{} เตสํ สุตฺวา ธมฺมิกญฺหิ โข รกฺขาวรณคุตฺติํ สํวิทหิ, โน จ โข อธนานํ ธนมนุปฺปทาสิ.\csromanspacer{} อธนานํ ธเน อนนุปฺปทิยมาเน ทาลิทฺทิยํ เวปุลฺลมคมาสิ.\csromanspacer{} ทาลิทฺทิเย เวปุลฺลํ คเต อญฺญตโร ปุริโส ปเรสํ อทินฺนํ เถยฺยสงฺขาตํ อาทิยิ.\csromanspacer{} ตเมนํ อคฺคเหสุํ.\csromanspacer{} คเหตฺวา รญฺโญ ขตฺติยสฺส มุทฺธาภิสิตฺตสฺส ทสฺเสสุํ “อยํ เทว ปุริโส ปเรสํ อทินฺนํ เถยฺยสงฺขาตํ อาทิยี”ติ.\csromanspacer{} เอวํ วุตฺเต ภิกฺขเว ราชา ขตฺติโย มุทฺธาภิสิตฺโต ตํ ปุริสํ เอตทโวจ “สจฺจํ กิร ตฺวํ อมฺโภ ปุริส ปเรสํ อทินฺนํ เถยฺยสงฺขาตํ อาทิยี”ติ\footnote{อาทิยสีติ (สฺยา)}.\csromanspacer{} สจฺจํ เทวาติ.\csromanspacer{} กิํ การณาติ.\csromanspacer{} น หิ เทว ชีวามีติ.\csromanspacer{} อถ โข ภิกฺขเว ราชา ขตฺติโย มุทฺธาภิสิตฺโต ตสฺส ปุริสสฺส ธนมนุปฺปทาสิ “อิมินา ตฺวํ อมฺโภ ปุริส ธเนน อตฺตนา จ ชีวาหิ, มาตาปิตโร จ โปเสหิ, ปุตฺตทารญฺจ โปเสหิ, กมฺมนฺเต จ ปโยเชหิ, สมณพฺราหฺมเณสุ\footnote{สมเณสุ พฺราหฺมเณสุ (พหูสุ)} อุทฺธคฺคิกํ ทกฺขิณํ ปติฏฺฐาเปหิ โสวคฺคิกํ สุขวิปากํ สคฺคสํวตฺตนิกนฺ”ติ. “เอวํ เทวา”ติ โข ภิกฺขเว โส ปุริโส รญฺโญ ขตฺติยสฺส มุทฺธาภิสิตฺตสฺส ปจฺจสฺโสสิ.}

\prose{อญฺญตโรปิ โข ภิกฺขเว ปุริโส ปเรสํ อทินฺนํ เถยฺยสงฺขาตํ อาทิยิ.\csromanspacer{} ตเมนํ อคฺคเหสุํ.\csromanspacer{} คเหตฺวา รญฺโญ ขตฺติยสฺส มุทฺธาภิสิตฺตสฺส ทสฺเสสุํ “อยํ เทว ปุริโส ปเรสํ อทินฺนํ เถยฺยสงฺขาตํ อาทิยี”ติ.\csromanspacer{} เอวํ วุตฺเต ภิกฺขเว ราชา ขตฺติโย มุทฺธาภิสิตฺโต ตํ ปุริสํ เอตทโวจ}

\vspace{\tipitakaparskip}
\csromancenter{จกฺกวตฺติสุตฺต “สจฺจํ กิร ตฺวํ อมฺโภ ปุริส ปเรสํ อทินฺนํ เถยฺยสงฺขาตํ อาทิยี”ติ.\csromanspacer{} สจฺจํ เทวาติ.\csromanspacer{} กิํ การณาติ.\csromanspacer{} น หิ เทว ชีวามีติ.\csromanspacer{} อถ โข ภิกฺขเว ราชา ขตฺติโย มุทฺธาภิสิตฺโต ตสฺส ปุริสสฺส ธนมนุปฺปทาสิ “อิมินา ตฺวํ อมฺโภ ปุรส ธเนน อตฺตนา จ ชีวาหิ, มาตาปิตโร จ โปเสหิ, ปุตฺตทารญฺจ โปเสหิ, กมฺมนฺเต จ ปโยเชหิ, สมณพฺราหฺมเณสุ อุทฺธคฺคิกํ ทกฺขิณํ ปติฏฺฐาเปหิ โสวคฺคิกํ สุขวิปากํ สคฺคสํวตฺตนิกนฺ”ติ. “เอวํ เทวา”ติ โข ภิกฺขเว โส ปุริโส รญฺโญ ขตฺติยสฺส มุทฺธาภิสิตฺตสฺส ปจฺจสฺโสสิ.}
\proseitem{92}{\csromanfolio{55}อสฺโสสุํ โข ภิกฺขเว มนุสฺสา “เย กิร โภ ปเรสํ อทินฺนํ เถยฺยสงฺขาตํ อาทิยนฺติ, เตสํ ราชา ธนมนุปฺปเทตี”ติ, สุตฺวาน เตสํ เอตทโหสิ “ยํนูน มยํปิ ปเรสํ อทินฺนํ เถยฺยสงฺขาตํ อาทิเยยฺยามา”ติ.\csromanspacer{} อถ โข ภิกฺขเว อญฺญตโร ปุริโส ปเรสํ อทินฺนํ เถยฺยสงฺขาตํ อาทิยิ.\csromanspacer{} ตเมนํ อคฺคเหสุํ.\csromanspacer{} คเหตฺวา รญฺโญ ขตฺติยสฺส มุทฺธาภิสิตฺตสฺส ทสฺเสสุํ “อยํ เทว ปุริโส ปเรสํ อทินฺนํ เถยฺยสงฺขาตํ อาทิยี”ติ.\csromanspacer{} เอวํ วุตฺเต ภิกฺขเว ราชา ขตฺติโย มุทฺทาภิสิตฺโต ตํ ปุริสํ เอตทโวจ “สจฺจํ กิร ตฺวํ อมฺโภ ปุริส ปเรสํ อทินฺนํ เถยฺยสงฺขาตํ อาทิยี”ติ.\csromanspacer{} สจฺจํ เทวาติ.\csromanspacer{} กิํ การณาติ.\csromanspacer{} น หิ เทว ชีวามีติ.\csromanspacer{} อถ โข ภิกฺขเว รญฺโญ ขตฺติยสฺส มุทฺธาภิสิตฺตสฺส เอตทโหสิ “สเจ โข อหํ โย โย ปเรสํ อทินฺนํ เถยฺยสงฺขาตํ อาทิยิสฺสติ, ตสฺส ตสฺส ธนมนุปฺปทสฺสามิ, เอวมิทํ อทินฺนาทานํ ปวฑฺฒิสฺสติ, ยํนูนาหํ อิมํ ปุริสํ สุนิเสธํ นิเสเธยฺยํ, มูลฆจฺจํ\footnote{มูลฆจฺฉํ (สฺยา), มูลเฉชฺชํ (ก)} กเรยฺยํ, สีสมสฺส ฉินฺเทยฺยนฺ”ติ.\csromanspacer{} อถ โข ภิกฺขเว ราชา ขตฺติโย มุทฺธาภิสิตฺโต ปุริเส อาณาเปสิ “เตน หิ ภเณ อิมํ ปุริสํ ทฬฺหาย รชฺชุยา ปจฺฉาพาหํ\footnote{ปจฺฉาพาหุํ (สฺยา)} คาฬฺหพนฺธนํ พนฺธิตฺวา ขุรมุณฺฑํ กริตฺวา ขรสฺสเรน ปณเวน รถิกาย รถิกํ สิงฺฆาฏเกน สิงฺฆาฏกํ ปริเนตฺวา ทกฺขิเณน ทฺวาเรน นิกฺขมิตฺวา ทกฺขิณโต นครสฺส สุนิเสธํ นิเสเธถ, มูลฆจฺจํ กโรถ, สีสมสฺส ฉินฺทถา”ติ. “เอวํ เทวา”ติ โข ภิกฺขเว เต ปุริสา รญฺโญ ขตฺติยสฺส มุทฺธาภิสิตฺตสฺส ปฏิสฺสุตฺวา ตํ ปุริสํ ทฬฺหาย รชฺชุยา ปจฺฉาพาหํ คาฬฺหพนฺธนํ พนฺธิตฺวา ขุรมุณฺฑํ กริตฺวา ขรสฺสเรน ปณเวน รถิกาย รถิกํ สิงฺฆฏเกน สิงฺฆาฏกํ \csromanfolio{56}ปริเนตฺวา ทกฺขิเณน ทฺวาเรน นิกฺขมิตฺวา ทกฺขิณโต นครสฺส สุนิเสธํ นิเสเธสุํ, มูลฆจฺจํ อกํสุ, สีสมสฺส ฉินฺทิํสุ.}

\proseitem{93}{อสฺโสสุํ โข ภิกฺขเว มนุสฺสา “เย กิร โภ ปเรสํ อทินฺนํ เถยฺยสงฺขาตํ อาทิยนฺติ, เต ราชา สุนิเสธํ นิเสเธติ, มูลฆจฺจํ กโรติ, สีสานิ เตสํ ฉินฺทตี”ติ, สุตฺวาน เตสํ เอตทโหสิ “ยํนูน มยํปิ ติณฺหานิ สตฺถานิ การาเปสฺสาม, ติณฺหานิ สตฺถานิ การาเปตฺวา เยสํ อทินฺนํ เถยฺยสงฺขาตํ อาทิยิสฺสาม, เต สุนิเสธํ นิเสเธสฺสาม, มูลฆจฺจํ กริสฺสาม, สีสานิ เตสํ ฉินฺทิสฺสามา”ติ.\csromanspacer{} เต ติณฺหานิ สตฺถานิ การาเปสุํ, ติณฺหานิ สตฺถานิ การาเปตฺวา คามฆาตํปิ อุปกฺกมิํสุ กาตุํ, นิคมฆาตํปิ อุปกฺกมิํสุ กาตุํ, นครฆาตํปิ อุปกฺกมิํสุ กาตุํ, ปนฺถทุหนํปิ\footnote{ปนฺถทูหนํปิ (สี, สฺยา, อิ)} อุปกฺกมิํสุ กาตุํ, เยสํ เต อทินฺนํ เถยฺยสงฺขาตํ อาทิยนฺติ, เต สุนิเสธํ นิเสเธนฺติ, มูลฆจฺจํ กโรนฺติ, สีสานิ เตสํ ฉินฺทนฺติ.}

\proseitem{94}{อิติ โข ภิกฺขเว อธนานํ ธเน อนนุปฺปทิยมาเน ทาลิทฺทิยํ เวปุลฺลมคมาสิ, ทาลิทฺทิเย เวปุลฺลํ คเต อทินฺนาทานํ เวปุลฺลมคมาสิ, อทินฺนาทาเน เวปุลฺลํ คเต สตฺถํ เวปุลฺลมคมาสิ, สตฺเถ เวปุลฺลํ คเต ปาณาติปาโต เวปุลฺลมคมาสิ, ปาณาติปาเต เวปุลฺลํ คเต เตสํ สตฺตานํ อายุปิ ปริหายิ, วณฺโณปิ ปริหายิ.\csromanspacer{} เตสํ อายุนาปิ ปริหายมานานํ วณฺเณนปิ ปริหายมานานํ อสีติวสฺสสหสฺสายุกานํ มนุสฺสานํ จตฺตารีสวสฺสสหสฺสายุกา ปุตฺตา อเหสุํ.}

\prose{จตฺตารีสวสฺสสหสฺสายุเกสุ ภิกฺขเว มนุสฺเสสุ อญฺญตโร ปุริโส ปเรสํ อทินฺนํ เถยฺยสงฺขาตํ อาทิยิ.\csromanspacer{} ตเมนํ อคฺคเหสุํ.\csromanspacer{} คเหตฺวา รญฺโญ ขตฺติยสฺส มุทฺธาภิสิตฺตสฺส ทสฺเสสุํ “อยํ เทว ปุริโส ปเรสํ อทินฺนํ เถยฺยสงฺขาตํ อาทิยี”ติ.\csromanspacer{} เอวํ วุตฺเต ภิกฺขเว ราชา ขตฺติโย มุทฺธาภิสิตฺโต ตํ ปุริสํ เอตทโวจ “สจฺจํ กิร ตฺวํ อมฺโภ ปุริส ปเรสํ อทินฺนํ เถยฺยสงฺขาตํ อาทิยี”ติ. “น หิ เทวา”ติ สมฺปชานมุสา อภาสิ.}

\csromanheader{2}{3. จกฺกวตฺติสุตฺต}
\proseitem{95}{\csromanfolio{57}อิติ โข ภิกฺขเว อธนานํ ธเน อนนุปฺปทิยมาเน ทาลิทฺทิยํ เวปุลฺลมคมาสิ, ทาลิทฺทิเย เวปุลฺลํ คเต อทินฺนาทานํ เวปุลฺลมคมาสิ, อทินฺนาทาเน เวปุลฺลํ คเต สตฺถํ เวปุลฺลมคมาสิ, สตฺเถ เวปุลฺลํ คเต ปาณาติปาโต เวปุลฺลมคมาสิ, ปาณาติปาเต เวปุลฺลํ คเต มุสาวาโท เวปุลฺลมคมาสิ, มุสาวาเท เวปุลฺลํ คเต เตสํ สตฺตานํ อายุปิ ปริหายิ, วณฺโณปิ ปริหายิ.\csromanspacer{} เตสํ อายุนาปิ ปริหายมานานํ วณฺเณนปิ ปริหายมานานํ จตฺตารีสวสฺสสหสฺสายุกานํ มนุสฺสานํ วีสติวสฺสสหสฺสายุกา ปุตฺตา อเหสุํ.}

\prose{วีสติวสฺสสหสฺสายุเกสุ ภิกฺขเว มนุสฺเสสุ อญฺญตโร ปุริโส ปเรสํ อทินฺนํ เถยฺยสงฺขาตํ อาทิยิ.\csromanspacer{} ตเมนํ อญฺญตโร ปุริโส รญฺโญ ขตฺติยสฺส มุทฺธาภิสิตฺตสฺส อาโรเจสิ “อิตฺถนฺนาโม เทว ปุริโส ปเรสํ อทินฺนํ เถยฺยสงฺขาตํ อาทิยี”ติ เปสุญฺญมกาสิ.}

\proseitem{96}{อิติ โข ภิกฺขเว อธนานํ ธเน อนนุปฺปทิยมาเน ทาลิทฺทิยํ เวปุลฺลมคมาสิ, ทาลิทฺทิเย เวปุลฺลํ คเต อทินฺนาทานํ เวปุลฺลมคมาสิ อทินฺนาทาเน เวปุลฺลํ คเต สตฺถํ เวปุลฺลมคมาสิ, สตฺเถ เวปุลฺลํ คเต ปาณาติปาโต เวปุลฺลมคมาสิ, ปาณาติปาเต เวปุลฺลํ คเต มุสาวาโท เวปุลฺลมคมาสิ, มุสาวาเท เวปุลฺลํ คเต ปิสุณา วาจา เวปุลฺลมคมาสิ, ปิสุณาย วาจาย เวปุลฺลํ คตาย เตสํ สตฺตานํ อายุปิ ปริหายิ, วณฺโณปิ ปริหายิ.\csromanspacer{} เตสํ อายุนาปิ ปริหายมานานํ วณฺเณนปิ ปริหายมานานํ วีสติวสฺสสหสฺสายุกานํ มนุสฺสานํ ทสวสฺสสหสฺสายุกา ปุตฺตา อเหสุํ.}

\prose{ทสวสฺสสหสฺสายุเกสุ ภิกฺขเว มนุสฺเสสุ เอกิทํ สตฺตา วณฺณวนฺโต โหนฺติ, เอกิทํ สตฺตา ทุพฺพณฺณา.\csromanspacer{} ตตฺถ เย เต สตฺตา ทุพฺพณฺณา, เต วณฺณวนฺเต สตฺเต อภิชฺฌายนฺตา ปเรสํ ทาเรสุ จาริตฺตํ อาปชฺชิํสุ.}

\proseitem{97}{อิติ โข ภิกฺขเว อธนานํ ธเน อนนุปฺปทิยมาเน ทาลิทฺทิยํ เวปุลฺลมคมาสิ, ทาลิทฺทิเย เวปุลฺลํ คเต ฯเปฯ กาเมสุมิจฺฉาจาโร เวปุลฺลมคมาสิ, กาเมสุมิจฺฉาจาเร เวปุลฺลํ คเต เตสํ สตฺตานํ อายุปิ ปริหายิ, วณฺโณปิ ปริหายิ.\csromanspacer{} เตสํ อายุนาปิ ปริหายมานานํ \csromanfolio{58}วณฺเณนปิ ปริหายมานานํ ทสวสฺสสหสฺสายุกานํ มนุสฺสานํ ปญฺจวสฺสสหสฺสายุกา ปุตฺตา อเหสุํ.}

\proseitem{98}{ปญฺจวสฺสสหสฺสายุเกสุ ภิกฺขเว มนุสฺเสสุ เทฺว ธมฺมา เวปุลฺลมคมํสุ, ผรุสาวาจา สมฺผปฺปลาโป จ.\csromanspacer{} ทฺวีสุ ธมฺเมสุ เวปุลฺลํ คเตสุ เตสํ สตฺตานํ อายุปิ ปริหายิ, วณฺโณปิ ปริหายิ.\csromanspacer{} เตสํ อายุนาปิ ปริหายมานานํ วณฺเณนปิ ปริหายมานานํ ปญฺจวสฺสสหสฺสายุกานํ มนุสฺสานํ อปฺเปกจฺเจ อฑฺฒเตยฺยวสฺสสหสฺสายุกา, อปฺเปกจฺเจ เทฺววสฺสสหสฺสายุกา ปุตฺตา อเหสุํ.}

\proseitem{99}{อฑฺฒเตยฺยวสฺสสหสฺสายุเกสุ ภิกฺขเว มนุสฺเสสุ อภิชฺฌาพฺยาปาทา เวปุลฺลมคมํสุ, อภิชฺฌาพฺยาปาเทสุ เวปุลฺลํ คเตสุ เตสํ สตฺตานํ อายุปิ ปริหายิ, วณฺโณปิ ปริหายิ.\csromanspacer{} เตสํ อายุนาปิ ปริหายมานานํ วณฺเณนปิ ปริหายมานานํ อฑฺฒเตยฺยวสฺสสหสฺสายุกานํ มนุสฺสานํ วสฺสสหสฺสายุกา ปุตฺตา อเหสุํ.}

\proseitem{100}{วสฺสสหสฺสายุเกสุ ภิกฺขเว มนุสฺเสสุ มิจฺฉาทิฏฺฐิ เวปุลฺลมคมาสิ, มิจฺฉาทิฏฺฐิยา เวปุลฺลํ คตาย เตสํ สตฺตานํ อายุปิ ปริหายิ, วณฺโณปิ ปริหายิ.\csromanspacer{} เตสํ อายุนาปิ ปริหายมานานํ วณฺเณนปิ ปริหายมานานํ วสฺสสหสฺสายุกานํ มนุสฺสานํ ปญฺจวสฺสสตายุกา ปุตฺตา อเหสุํ.}

\proseitem{101}{ปญฺจวสฺสสตายุเกสุ ภิกฺขเว มนุสฺเสสุ ตโย ธมฺมา เวปุลฺลมคมํสุ, อธมฺมราโค วิสมโลโภ มิจฺฉาธมฺโม.\csromanspacer{} ตีสุ ธมฺเมสุ เวปุลฺลํ คเตสุ เตสํ สตฺตานํ อายุปิ ปริหายิ, วณฺโณปิ ปริหายิ.\csromanspacer{} เตสํ อายุนาปิ ปริหายมานานํ วณฺเณนปิ ปริหายมานานํ ปญฺจวสฺสสตายุกานํ มนุสฺสานํ อปฺเปกจฺเจ อฑฺฒเตยฺยวสฺสสตายุกา, อปฺเปกจฺเจ เทฺววสฺสสตายุกา ปุตฺตา อเหสุํ.}

\prose{อฑฺฒเตยฺยวสฺสสตายุเกสุ ภิกฺขเว มนุสฺเสสุ อิเม ธมฺมา เวปุลฺลมคมํสุ, อมตฺเตยฺยตา อเปตฺเตยฺยตา อสามญฺญตา อพฺรหฺมญฺญตา น กุเล เชฏฺฐาปจายิตา.}

\csromanheader{2}{3. จกฺกวตฺติสุตฺต}
\proseitem{102}{\csromanfolio{59}อิติ โข ภิกฺขเว อธนานํ ธเน อนนุปฺปทิยมาเน ทาลิทฺทิยํ เวปุลฺลมคมาสิ.\csromanspacer{} ทาลิทฺทิเย เวปุลฺลํ คเต อทินฺนาทานํ เวปุลฺลมคมาสิ.\csromanspacer{} อทินฺนาทาเน เวปุลฺลํ คเต สตฺถํ เวปุลฺลมคมาสิ.\csromanspacer{} สตฺเถ เวปุลฺลํ คเต ปาณาติปาโต เวปุลฺลมคมาสิ.\csromanspacer{} ปาณาติปาเต เวปุลฺลํ คเต มุสาวาโท เวปุลฺลมคมาสิ.\csromanspacer{} มุสาวาเท เวปุลฺลํ คเต ปิสุณา วาจา เวปุลฺลมคมาสิ.\csromanspacer{} ปิสุณาย วาจาย เวปุลฺลํ คตาย กาเมสุมิจฺฉาจาโร เวปุลฺลมคมาสิ.\csromanspacer{} กาเมสุมิจฺฉาจาเร เวปุลฺลํ คเต เทฺว ธมฺมา เวปุลฺลมคมํสุ, ผรุสา วาจา สมฺผปฺปลาโป จ.\csromanspacer{} ทฺวีสุ ธมฺเมสุ เวปุลฺลํ คเตสุ อภิชฺฌาพฺยาปาทา เวปุลฺลมคมํสุ.\csromanspacer{} อภิชฺฌาพฺยาปาเทสุ เวปุลฺลํ คเตสุ มิจฺฉาทิฏฺฐิ เวปุลฺลมคมาสิ.\csromanspacer{} มิจฺฉาทิฏฺฐิยา เวปุลฺลํ คตาย ตโย ธมฺมา เวปุลฺลมคมํสุ, อธมฺมราโค วิสมโลโภ มิจฺฉาธมฺโม.\csromanspacer{} ตีสุ ธมฺเมสุ เวปุลฺลํ คเตสุ อิเม ธมฺมา เวปุลฺลมคมํสุ, อมตฺเตยฺยตา อเปตฺเตยฺยตา อสามญฺญตา อพฺรหฺมญฺญตา น กุเล เชฏฺฐาปจายิตา.\csromanspacer{} อิเมสุ ธมฺเมสุ เวปุลฺลํ คเตสุ เตสํ สตฺตานํ อายุปิ ปริหายิ, วณฺโณปิ ปริหายิ.\csromanspacer{} เตสํ อายุนาปิ ปริหายมานานํ วณฺเณนปิ ปริหายมานานํ อฑฺฒเตยฺยวสฺสสตายุกานํ มนุสฺสานํ วสฺสสตายุกา ปุตฺตา อเหสุํ.}

\csromanmatikaanchor{mtk.26}
\csromantocmarkref{section}{ทสวสฺสายุกสมย}{mtk.26}
\bookmark[dest={mtk.26},level=1]{ทสวสฺสายุกสมย}
\csromanheader{1}{\textbf{ทสวสฺสายุกสมย}}
\proseitem{103}{ภวิสฺสติ ภิกฺขเว โส สมโย, ยํ อิเมสํ มนุสฺสานํ ทสวสฺสายุกา ปุตฺตา ภวิสฺสนฺติ.\csromanspacer{} ทสวสฺสายุเกสุ ภิกฺขเว มนุสฺเสสุ ปญฺจวสฺสิกา\footnote{ปญฺจมาสิกา (ก-สี)} กุมาริกา อลํปเตยฺยา ภวิสฺสนฺติ.\csromanspacer{} ทสวสฺสายุเกสุ ภิกฺขเว มนุสฺเสสุ อิมานิ รสานิ อนฺตรธายิสฺสนฺติ, เสยฺยถิทํ, สปฺปิ นวนีตํ เตลํ มธุ ผาณิตํ โลณํ.\csromanspacer{} ทสวสฺสายุเกสุ ภิกฺขเว มนุสฺเสสุ กุทฺรูสโก อคฺคํ โภชนานํ\footnote{อคฺคโภชนํ (สฺยา)} ภวิสฺสติ, เสยฺยถาปิ ภิกฺขเว เอตรหิ สาลิมํโสทโน อคฺคํ โภชนานํ, เอวเมว โข ภิกฺขเว ทสวสฺสายุเกสุ มนุสฺเสสุ กุทฺรูสโก อคฺคํ โภชนานํ ภวิสฺสติ.}

\prose{ทสวสฺสายุเกสุ ภิกฺขเว มนุสฺเสสุ ทส กุสลกมฺมปถา สพฺเพน สพฺพํ อนฺตรธายิสฺสนฺติ, ทส อกุสลกมฺมปถา อติพฺยาทิปฺปิสฺสนฺติ\footnote{อติวิย ทิปฺปิสฺสนฺติ (สฺยา, อิ), อติฌาทิปฺปิสฺสนฺติ (สี)}. \csromanfolio{60}ทสวสฺสายุเกสุ ภิกฺขเว มนุสฺเสสุ กุสลนฺติปิ น ภวิสฺสติ, กุโต ปน กุสลสฺส การโก.\csromanspacer{} ทสวสฺสายุเกสุ ภิกฺขเว มนุสฺเสสุ เย เต ภวิสฺสนฺติ อมตฺเตยฺยา อเปตฺเตยฺยา อสามญฺญา อพฺรหฺมญฺญา น กุเล เชฏฺฐาปจายิโน, เต ปุชฺชา จ ภวิสฺสนฺติ ปาสํสา จ, เสยฺยถาปิ ภิกฺขเว เอตรหิ มตฺเตยฺยา เปตฺเตยฺยา สามญฺญา พฺรหฺมญฺญา กุเล เชฏฺฐาปจายิโน ปุชฺชา จ ปาสํสา จ, เอวเมว โข ภิกฺขเว ทสวสฺสายุเกสุ มนุสฺเสสุ เย เต ภวิสฺสนฺติ อมตฺเตยฺยา อเปตฺเตยฺยา อสามญฺญา อพฺรหฺมญฺญา น กุเล เชฏฺฐาปจายิโน, เต ปุชฺชา จ ภวิสฺสนฺติ ปาสํสา จ. ทสวสฺสายุเกสุ ภิกฺขเว มนุสฺเสสุ น ภวิสฺสติ มาตาติ วา มาตุจฺฉาติ วา มาตุลานีติ วา อาจริยภริยาติ วา ครูนํ ทาราติ วา, สมฺเภทํ โลโก คมิสฺสติ ยถา อเชฬกา กุกฺกุฏสูกรา โสณสิงฺคาลา\footnote{โสณสิคาลา (สี, อิ)}. ทสวสฺสายุเกสุ ภิกฺขเว มนุสฺเสสุ เตสํ สตฺตานํ อญฺญมญฺญมฺหิ ติพฺโพ อาฆาโต ปจฺจุปฏฺฐิโต ภวิสฺสติ ติพฺโพ พฺยาปาโท ติพฺโพ มโนปโทโส ติพฺพํ วธกจิตฺตํ.\csromanspacer{} มาตุปิ ปุตฺตมฺหิ ปุตฺตสฺสปิ มาตริ, ปิตุปิ ปุตฺตมฺหิ ปุตฺตสฺสปิ ปิตริ, ภาตุปิ ภคินิยา ภคินิยาปิ ภาตริ ติพฺโพ อาฆาโต ปจฺจุปฏฺฐิโต ภวิสฺสติ ติพฺโพ พฺยาปาโท ติพฺโพ มโนปโทโส ติพฺพํ วธกจิตฺตํ.\csromanspacer{} เสยฺยถาปิ ภิกฺขเว มาควิกสฺส มิคํ ทิสฺวา ติพฺโพ อาฆาโต ปจฺจุปฏฺฐิโต โหติ ติพฺโพ พฺยาปาโท ติพฺโพ มโนปโทโส ติพฺพํ วธกจิตฺตํ, เอวเมว โข ภิกฺขเว ทสวสฺสายุเกสุ มนุสฺเสสุ เตสํ สตฺตานํ อญฺญมญฺญมฺหิ ติพฺโพ อาฆาโต ปจฺจุปฏฺฐิโต ภวิสฺสติ ติพฺโพ พฺยาปาโท ติพฺโพ มโนปโทโส ติพฺพํ วธกจิตฺตํ.\csromanspacer{} มาตุปิ ปุตฺตมฺหิ ปุตฺตสฺสปิ มาตริ, ปิตุปิ ปุตฺตมฺหิ ปุตฺตสฺสปิ ปิตริ, ภาตุปิ ภคินิยา ภคินิยาปิ ภาตริ ติพฺโพ อาฆาโต ปจฺจุปฏฺฐิโต ภวิสฺสติ ติพฺโพ พฺยาปาโท ติพฺโพ มโนปโทโส ติพฺพํ วธกจิตฺตํ.}

\proseitem{104}{ทสวสฺสายุเกสุ ภิกฺขเว มนุสฺเสสุ สตฺตาหํ สตฺถนฺตรกปฺโป ภวิสฺสติ.\csromanspacer{} เต อญฺญมญฺญมฺหิ มิคสญฺญํ ปฏิลภิสฺสนฺติ.\csromanspacer{} เตสํ ติณฺหานิ}

\vspace{\tipitakaparskip}
\csromancenter{จกฺกวตฺติสุตฺต สตฺถานิ หตฺเถสุ ปาตุภวิสฺสนฺติ.\csromanspacer{} เต ติเณฺหน สตฺเถน “เอส มิโค เอส มิโค”ติ อญฺญมญฺญํ ชีวิตา โวโรเปสฺสนฺติ.}
\prose{\csromanfolio{61}อถ โข เตสํ ภิกฺขเว สตฺตานํ เอกจฺจานํ เอวํ ภวิสฺสติ “มา จ มยํ กญฺจิ\footnote{กิญฺจิ (ก) 2. เต ติณคหนํ วนคหนํ รุกฺขคหนํ นทีวิทุคฺคํ ปพฺพตวิสมํ (สี, อิ)}, มา จ อเมฺห โกจิ, ยํนูน มยํ ติณคหนํ วา วนคหนํ วา รุกฺขคหนํ วา นทีวิทุคฺคํ วา ปพฺพตวิสมํ วา ปวิสิตฺวา วนมูลผลาหารา ยาเปยฺยามา”ติ.\csromanspacer{} เต ติณคหนํ วา วนคหนํ วา รุกฺขคหนํ วา นทีวิทุคฺคํ วา ปพฺพตวิสมํ วา2 ปวิสิตฺวา สตฺตาหํ วนมูลผลาหารา ยาเปสฺสนฺติ.\csromanspacer{} เต ตสฺส สตฺตาหสฺส อจฺจเยน ติณคหนา วนคหนา รุกฺขคหนา นทีวิทุคฺคา ปพฺพตวิสมา นิกฺขมิตฺวา อญฺญมญฺญํ อาลิงฺคิตฺวา สภาคายิสฺสนฺติ สมสฺสาสิสฺสนฺติ “ทิฏฺฐา โภ สตฺตา ชีวสิ, ทิฏฺฐา โภ สตฺตา ชีวสี”ติ.}

\csromanmatikaanchor{mtk.27}
\csromantocmarkref{section}{อายุวณฺณาทิวฑฺฒนกถา}{mtk.27}
\bookmark[dest={mtk.27},level=1]{อายุวณฺณาทิวฑฺฒนกถา}
\csromanheader{1}{\textbf{อายุวณฺณาทิวฑฺฒนกถา}}
\proseitem{105}{อถ โข เตสํ ภิกฺขเว สตฺตานํ เอวํ ภวิสฺสติ “มยํ โข อกุสลานํ ธมฺมานํ สมาทานเหตุ เอวรูปํ อายตํ ญาติกฺขยํ ปตฺตา, ยํนูน มยํ กุสลํ กเรยฺยาม, กิํ กุสลํ กเรยฺยาม, ยํนูน มยํ ปาณาติปาตา วิรเมยฺยาม, อิทํ กุสลํ ธมฺมํ สมาทาย วตฺเตยฺยามา”ติ.\csromanspacer{} เต ปาณาติปาตา วิรมิสฺสนฺติ, อิทํ กุสลํ ธมฺมํ สมาทาย วตฺติสฺสนฺติ.\csromanspacer{} เต กุสลานํ ธมฺมานํ สมาทานเหตุ อายุนาปิ วฑฺฒิสฺสนฺติ, วณฺเณนปิ วฑฺฒิสฺสนฺติ.\csromanspacer{} เตสํ อายุนาปิ วฑฺฒมานานํ วณฺเณนปิ วฑฺฒมานานํ ทสวสฺสายุกานํ มนุสฺสานํ วีสติวสฺสายุกา ปุตฺตา ภวิสฺสนฺติ.}

\prose{อถ โข เตสํ ภิกฺขเว สตฺตานํ เอวํ ภวิสฺสติ “มยํ โข กุสลานํ ธมฺมานํ สมาทานเหตุ อายุนาปิ วฑฺฒาม, วณฺเณนปิ วฑฺฒาม, ยํนูน มยํ ภิยฺโยโส มตฺตาย กุสลํ กเรยฺยาม, กิํ กุสลํ กเรยฺยาม, ยํนูน มยํ อทินฺนาทานา วิรเมยฺยาม.\csromanspacer{} กาเมสุมิจฺฉาจารา วิรเมยฺยาม.\csromanspacer{} มุสาวาทา วิรเมยฺยาม.\csromanspacer{} ปิสุณาย วาจาย วิรเมยฺยาม.\csromanspacer{} ผรุสาย วาจาย วิรเมยฺยาม.\csromanspacer{} สมฺผปฺปลาปา วิรเมยฺยาม.\csromanspacer{} อภิชฺฌํ ปชเหยฺยาม.\csromanspacer{} พฺยาปาทํ ปชเหยฺยาม.\csromanspacer{} มิจฺฉาทิฏฺฐิํ ปชเหยฺยาม.\csromanspacer{} ตโย ธมฺเม ปชเหยฺยาม, อธมฺมราคํ วิสมโลภํ มิจฺฉาธมฺมํ.\csromanspacer{} ยํนูน มยํ มตฺเตยฺยา อสฺสาม เปตฺเตยฺยา สามญฺญา พฺรหฺมญฺญา กุเล เชฏฺฐาปจายิโน, อิทํ กุสลํ ธมฺมํ สมาทาย วตฺเตยฺยามา”ติ.\csromanspacer{} เต มตฺเตยฺยา ภวิสฺสนฺติ เปตฺเตยฺยา \csromanfolio{62}สามญฺญา พฺรหฺมญฺญา กุเล เชฏฺฐาปจายิโน, อิทํ กุสลํ ธมฺมํ สมาทาย วตฺติสฺสนฺติ.}

\prose{เต กุสลานํ ธมฺมานํ สมาทานเหตุ อายุนาปิ วฑฺฒิสฺสนฺติ, วณฺเณนปิ วฑฺฒิสฺสนฺติ.\csromanspacer{} เตสํ อายุนาปิ วฑฺฒมานานํ วณฺเณนปิ วฑฺฒมานานํ วีสติวสฺสายุกานํ มนุสฺสานํ จตฺตารีสวสฺสายุกา ปุตฺตา ภวิสฺสนฺติ.\csromanspacer{} จตฺตารีสวสฺสายุกานํ มนุสฺสานํ อสีติวสฺสายุกา ปุตฺตา ภวิสฺสนฺติ.\csromanspacer{} อสีติวสฺสายุกานํ มนุสฺสานํ สฏฺฐิวสฺสสตายุกา ปุตฺตา ภวิสฺสนฺติ.\csromanspacer{} สฏฺฐิวสฺสาสตายุกานํ มนุสฺสานํ วีสติ-ติวสฺสสตายุกา ปุตฺตา ภวิสฺสนฺติ.\csromanspacer{} วีสติ-ติวสฺสสตายุกานํ มนุสฺสานํ จตฺตารีส- ฉพฺพสฺสสตายุกา ปุตฺตา ภวิสฺสนฺติ.\csromanspacer{} จตฺตารีส-ฉพฺพสฺสสตายุกานํ มนุสฺสานํ เทฺววสฺสสหสฺสายุกา ปุตฺตา ภวิสฺสนฺติ.\csromanspacer{} เทฺววสฺสสหสฺสายุกานํ มนุสฺสานํ จตฺตาริวสฺสสหสฺสายุกา ปุตฺตา ภวิสฺสนฺติ.\csromanspacer{} จตฺตาริวสฺสสหสฺสายุกานํ มนุสฺสานํ อฏฺฐวสฺสสหสฺสายุกา ปุตฺตา ภวิสฺสนฺติ.\csromanspacer{} อฏฺฐวสฺสสหสฺสายุกานํ มนุสฺสานํ วีสติวสฺสสหสฺสายุกา ปุตฺตา ภวิสฺสนฺติ.\csromanspacer{} วีสติวสฺสสหสฺสายุกานํ มนุสฺสานํ จตฺตารีสวสฺสสหสฺสายุกา ปุตฺตา ภวิสฺสนฺติ.\csromanspacer{} จตฺตารีสวสฺสสหสฺสายุกานํ มนุสฺสานํ อสีติวสฺสสหสฺสายุกา ปุตฺตา ภวิสฺสนฺติ.\csromanspacer{} อสีติวสฺส-สหสฺสายุเกสุ ภิกฺขเว มนุสฺเสสุ ปญฺจวสฺสสติกา กุมาริกา อลํปเตยฺยา ภวิสฺสนฺติ.}

\csromanmatikaanchor{mtk.28}
\csromantocmarkref{section}{สงฺขราช-อุปฺปตฺติ}{mtk.28}
\bookmark[dest={mtk.28},level=1]{สงฺขราช-อุปฺปตฺติ}
\csromanheader{1}{\textbf{สงฺขราช-อุปฺปตฺติ}}
\proseitem{106}{อสีติวสฺสสหสฺสายุเกสุ ภิกฺขเว มนุสฺเสสุ ตโย อาพาธา ภวิสฺสนฺติ, อิจฺฉา อนสนํ ชรา.\csromanspacer{} อสีติวสฺสสหสฺสายุเกสุ ภิกฺขเว มนุสฺเสสุ อยํ ชมฺพุทีโป อิทฺโธ เจว ภวิสฺสติ ผีโต จ, กุกฺกุฏสมฺปาติกา คามนิคมราชธานิโย\footnote{คามนิคมชนปทา ราชธานิโย (ก)}.\csromanspacer{} อสีติวสฺสสหสฺสายุเกสุ ภิกฺขเว มนุสฺเสสุ อยํ ชมฺพุทีโป อวีจิ มญฺเญ ผุโฏ ภวิสฺสติ มนุสฺเสหิ, เสยฺยถาปิ นฬวนํ วา สรวนํ\footnote{สารวนํ (สฺยา)} วา.\csromanspacer{} อสีติวสฺสสหสฺสายุเกสุ ภิกฺขเว มนุสฺเสสุ อยํ พาราณสี เกตุมตี นาม ราชธานี ภวิสฺสติ อิทฺธา เจว ผีตา จ พหุชนา จ อากิณฺณมนุสฺสา จ สุภิกฺขา จ.\csromanspacer{} อสีติวสฺสสหสฺสายุเกสุ ภิกฺขเว}

\vspace{\tipitakaparskip}
\csromancenter{จกฺกวตฺติสุตฺต มนุสฺเสสุ อิมสฺมิํ ชมฺพุทีเป จตุราสีติ นครสหสฺสานิ ภวิสฺสนฺติ เกตุมตีราชธานีปมุขานิ.\csromanspacer{} อสีติวสฺสสหสฺสายุเกสุ ภิกฺขเว มนุสฺเสสุ เกตุมติยา ราชธานิยา สงฺโข นาม ราชา อุปฺปชฺชิสฺสติ จกฺกวตฺตี ธมฺมิโก ธมฺมราชา จาตุรนฺโต วิชิตาวี ชนปทตฺถาวริยปฺปตฺโต สตฺตรตนสมนฺนาคโต.\csromanspacer{} ตสฺสิมานิ สตฺต รตนานิ ภวิสฺสนฺติ, เสยฺยถิทํ, จกฺกรตนํ หตฺถิรตนํ อสฺสรตนํ มณิรตนํ อิตฺถิรตนํ คหปติรตนํ ปริณายกรตนเมว สตฺตมํ.\csromanspacer{} ปโรสหสฺสํ โข ปนสฺส ปุตฺตา ภวิสฺสนฺติ สูรา วีรงฺครูปา ปรเสนปฺปมทฺทนา, โส อิมํ ปถวิํ สาครปริยนฺตํ อทณฺเฑน อสตฺเถน ธมฺเมน อภิวิชิย อชฺฌาวสิสฺสติ.}
\csromanmatikaanchor{mtk.29}
\csromantocmarkref{section}{เมตฺเตยฺยพุทฺธุปฺปาท}{mtk.29}
\bookmark[dest={mtk.29},level=1]{เมตฺเตยฺยพุทฺธุปฺปาท}
\csromanheader{1}{\textbf{เมตฺเตยฺยพุทฺธุปฺปาท}}
\proseitem{107}{\csromanfolio{63}อสีติวสฺสสหสฺสายุเกสุ ภิกฺขเว มนุสฺเสสุ เมตฺเตยฺโย นาม ภควา โลเก อุปฺปชฺชิสฺสติ อรหํ สมฺมาสมฺพุทฺโธ วิชฺชาจรณสมฺปนฺโน สุคโต โลกวิทู อนุตฺตโร ปุริสทมฺมสารถิ สตฺถา เทวมนุสฺสานํ พุทฺโธ ภควา เสยฺยถาปาหเมตรหิ โลเก อุปฺปนฺโน อรหํ สมฺมาสมฺพุทฺโธ วิชฺชาจรณสมฺปนฺโน สุคโต โลกวิทู อนุตฺตโร ปุริสทมฺมสารถิ สตฺถา เทวมนุสฺสานํ พุทฺโธ ภควา.\csromanspacer{} โส อิมํ โลกํ สเทวกํ สมารกํ สพฺรหฺมกํ สสฺสมณพฺราหฺมณิํ ปชํ สเทวมนุสฺสํ สยํ อภิญฺญา สจฺฉิกตฺวา ปเวเทสฺสติ, เสยฺยถาปาหเมตรหิ อิมํ โลกํ สเทวกํ สมารกํ สพฺรหฺมกํ สสฺสมณพฺราหฺมณิํ ปชํ สเทวมนุสฺสํ สยํ อภิญฺญา สจฺฉิกตฺวา ปเวเทมิ.\csromanspacer{} โส ธมฺมํ เทเสสฺสติ อาทิกลฺยาณํ มชฺเฌกลฺยาณํ ปริโยสานกลฺยาณํ สาตฺถํ สพฺยญฺชนํ เกวลปริปุณฺณํ ปริสุทฺธํ พฺรหฺมจริยํ ปกาเสสฺสติ, เสยฺยถาปาหเมตรหิ ธมฺมํ เทเสมิ อาทิกลฺยาณํ มชฺเฌกลฺยาณํ ปริโยสานกลฺยาณํ สาตฺถํ สพฺยญฺชนํ เกวลปริปุณฺณํ ปริสุทฺธํ พฺรหฺมจริยํ ปกาเสมิ.\csromanspacer{} โส อเนกสหสฺสํ\footnote{อเนกสตสหสฺสํ (ก)} ภิกฺขุสํฆํ ปริหริสฺสติ, เสยฺยถาปาหเมตรหิ อเนกสตํ ภิกฺขุสํฆํ ปริหรามิ.}

\proseitem{108}{อถ โข ภิกฺขเว สงฺโข นาม ราชา โย โส ยูโป รญฺญา มหาปนาเทน การาปิโต, ตํ ยูปํ อุสฺสาเปตฺวา อชฺฌาวสิตฺวา ตํ \csromanfolio{64}ทตฺวา วิสฺสชฺชิตฺวา สมณพฺราหฺมณกปณทฺธิกวณิพฺพกยาจกานํ ทานํ ทตฺวา เมตฺเตยฺยสฺส ภควโต อรหโต สมฺมาสมฺพุทฺธสฺส สนฺติเก เกสมสฺสุํ โอหาเรตฺวา กาสายานิ วตฺถานิ อจฺฉาเทตฺวา อคารสฺมา อนคาริยํ ปพฺพชิสฺสติ.\csromanspacer{} โส เอวํ ปพฺพชิโต สมาโน เอโก วูปกฏฺโฐ อปฺปมตฺโต อาตาปี ปหิตตฺโต วิหรนฺโต นจิรสฺเสว ยสฺสตฺถาย กุลปุตฺตา สมฺมเทว อคารสฺมา อนคาริยํ ปพฺพชนฺติ, ตทนุตฺตรํ พฺรหฺมจริยปริโยสานํ ทิฏฺเฐว ธมฺเม สยํ อภิญฺญา สจฺฉิกตฺวา อุปสมฺปชฺช วิหริสฺสติ.}

\proseitem{109}{อตฺตทีปา ภิกฺขเว วิหรถ อตฺตสรณา อนญฺญสรณา ธมฺมทีปา ธมฺมสรณา อนญฺญสรณา.\csromanspacer{} กถญฺจ ภิกฺขเว ภิกฺขุ อตฺตทีโป วิหรติ อตฺตสรโณ อนญฺญสรโณ ธมฺมทีโป ธมฺมสรโณ อนญฺญสรโณ.\csromanspacer{} อิธ ภิกฺขเว ภิกฺขุ กาเย กายานุปสฺสี วิหรติ อาตาปี สมฺปชาโน สติมา วิเนยฺย โลเก อภิชฺฌาโทมนสฺสํ.\csromanspacer{} เวทนาสุ เวทนานุปสฺสี ฯเปฯ.\csromanspacer{} จิตฺเต จิตฺตานุปสฺสี ฯเปฯ.\csromanspacer{} ธมฺเมสุ ธมฺมานุปสฺสี วิหรติ อาตาปี สมฺปชาโน สติมา วิเนยฺย โลเก อภิชฺฌาโทมนสฺสํ.\csromanspacer{} เอวํ โข ภิกฺขเว ภิกฺขุ อตฺตทีโป วิหรติ อตฺตสรโณ อนญฺญสรโณ ธมฺมทีโป ธมฺมสรโณ อนญฺญสรโณ.}

\csromanmatikaanchor{mtk.30}
\csromantocmarkref{section}{ภิกฺขุโน-อายุวณฺณาทิวฑฺฒนกถา}{mtk.30}
\bookmark[dest={mtk.30},level=1]{ภิกฺขุโน-อายุวณฺณาทิวฑฺฒนกถา}
\csromanheader{1}{\textbf{ภิกฺขุโน-อายุวณฺณาทิวฑฺฒนกถา}}
\proseitem{110}{โคจเร ภิกฺขเว จรถ สเก เปตฺติเก วิสเย.\csromanspacer{} โคจเร ภิกฺขเว จรนฺตา สเก เปตฺติเก วิสเย อายุนาวิ วฑฺฒิสฺสถ, วณฺเณนปิ วฑฺฒิสฺสถ, สุเขนปิ วฑฺฒิสฺสถ, โภเคนปิ วฑฺฒิสฺสถ, พเลนปิ วฑฺฒิสฺสถ.}

\prose{กิญฺจ ภิกฺขเว ภิกฺขุโน อายุสฺมิํ.\csromanspacer{} อิธ ภิกฺขเว ภิกฺขุ ฉนฺทสมาธิปธานสงฺขารสมนฺนาคตํ อิทฺธิปาทํ ภาเวติ, วีริยสมาธิปธานสงฺขารสมนฺนาคตํ อิทฺธิปาทํ ภาเวติ, จิตฺตสมาธิปธานสงฺขารสมนฺนาคตํ อิทฺธิปาทํ ภาเวติ, วีมํสาสมาธิปธานสงฺขารสมนฺนาคตํ อิทฺธิปาทํ ภาเวติ.\csromanspacer{} โส อิเมสํ จตุนฺนํ อิทฺธิปาทานํ ภาวิตตฺตา พหุลีกตตฺตา อากงฺขมาโน กปฺปํ วา ติฏฺเฐยฺย กปฺปาวเสสํ วา.\csromanspacer{} อิทํ โข ภิกฺขเว ภิกฺขุโน อายุสฺมิํ.}

\prose{กิญฺจ ภิกฺขเว ภิกฺขุโน วณฺณสฺมิํ.\csromanspacer{} อิธ ภิกฺขเว ภิกฺขุ สีลวา โหติ, ปาติโมกฺขสํวรสํวุโต วิหรติ อาจารโคจรสมฺปนฺโน อณุมตฺเตสุ}

\vspace{\tipitakaparskip}
\csromancenter{จกฺกวตฺติสุตฺต วชฺเชสุ ภยทสฺสาวี, สมาทาย สิกฺขติ สิกฺขาปเทสุ.\csromanspacer{} อิทํ โข ภิกฺขเว ภิกฺขุโน วณฺณสฺมิํ.}
\prose{\csromanfolio{65}กิญฺจ ภิกฺขเว ภิกฺขุโน สุขสฺมิํ.\csromanspacer{} อิธ ภิกฺขเว ภิกฺขุ วิวิจฺเจว กาเมหิ วิวิจฺจ อกุสเลหิ ธมฺเมหิ สวิตกฺกํ สวิจารํ วิเวกชํ ปีติสุขํ ปฐมํ ฌานํ อุปสมฺปชฺช วิหรติ.\csromanspacer{} วิตกฺกวิจารานํ วูปสมา ฯเปฯ ทุติยํ ฌานํ.\csromanspacer{} ตติยํ ฌานํ.\csromanspacer{} จตุตฺถํ ฌานํ อุปสมฺปชฺช วิหรติ.\csromanspacer{} อิทํ โข ภิกฺขเว ภิกฺขุโน สุขสฺมิํ.}

\prose{กิญฺจ ภิกฺขเว ภิกฺขุโน โภคสฺมิํ.\csromanspacer{} อิธ ภิกฺขเว ภิกฺขุ เมตฺตาสหคเตน เจตสา เอกํ ทิสํ ผริตฺวา วิหรติ.\csromanspacer{} ตถา ทุติยํ.\csromanspacer{} ตถา ตติยํ.\csromanspacer{} ตถา จตุตฺถํ.\csromanspacer{} อิติ อุทฺธมโธ ติริยํ สพฺพธิ สพฺพตฺตตาย สพฺพาวนฺตํ โลกํ เมตฺตาสหคเตน เจตสา วิปุเลน มหคฺคเตน อปฺปมาเณน อเวเรน อพฺยาปชฺเชน ผริตฺวา วิหรติ.\csromanspacer{} กรุณาสหคเตน เจตสา ฯเปฯ.\csromanspacer{} มุทิตาสหคเตน เจตสา ฯเปฯ.\csromanspacer{} อุเปกฺขาสหคเตน เจตสา เอกํ ทิสํ ผริตฺวา วิหรติ.\csromanspacer{} ตถา ทุติยํ.\csromanspacer{} ตถา ตติยํ.\csromanspacer{} ตถา จตุตฺถํ.\csromanspacer{} อิติ อุทฺธมโธ ติริยํ สพฺพธิ สพฺพตฺตตาย สพฺพาวนฺตํ โลกํ อุเปกฺขาสหคเตน เจตสา วิปุเลน มหคฺคเตน อปฺปมาเณน อเวเรน อพฺยาปชฺเชน ผริตฺวา วิหรติ.\csromanspacer{} อิทํ โข ภิกฺขเว ภิกฺขุโน โภคสฺมิํ.}

\prose{กิญฺจ ภิกฺขเว ภิกฺขุโน พลสฺมิํ.\csromanspacer{} อิธ ภิกฺขเว ภิกฺขุ อาสวานํ ขยา อนาสวํ เจโตวิมุตฺติํ ปญฺญวิมุตฺติํ ทิฏฺเฐว ธมฺเม สยํ อภิญฺญา สจฺฉิกตฺวา อุปสมฺปชฺช วิหรติ.\csromanspacer{} อิทํ โข ภิกฺขเว ภิกฺขุโน พลสฺมิํ.}

\prose{นาหํ ภิกฺขเว อญฺญํ เอกพลํปิ สมนุปสฺสามิ, ยํ เอวํ ทุปฺปสหํ, ยถยิทํ ภิกฺขเว มารพลํ.\csromanspacer{} กุสลานํ ภิกฺขเว ธมฺมานํ สมาทานเหตุ เอวมิทํ ปุญฺญํ ปวฑฺฒตีติ.\csromanspacer{} อิทมโวจ ภควา.\csromanspacer{} อตฺตมนา เต ภิกฺขู ภควโต ภาสิตํ อภินนฺทุนฺติ.}

\nitthitam{\textbf{จกฺกวตฺติสุตฺตํ นิฏฺฐิตํ ตติยํ.}}
\csromanmatikaanchor{mtk.31}
\csromantocmarknopage{chapter}{4. อคฺคญฺญสุตฺต}{mtk.31}
\bookmark[dest={mtk.31},level=0]{4. อคฺคญฺญสุตฺต}
\chapterheadpageread{4. อคฺคญฺญสุตฺต}
\csromanmatikaanchor{mtk.32}
\csromantocmarkref{section}{วาเสฏฺฐภารทฺวาช}{mtk.32}
\bookmark[dest={mtk.32},level=1]{วาเสฏฺฐภารทฺวาช}
\csromanheader{1}{\textbf{วาเสฏฺฐภารทฺวาช}}
\proseitem{111}{\csromanfolio{66}เอวํ เม สุตํ\csromandash{}เอกํ สมยํ ภควา สาวตฺถิยํ วิหรติ ปุพฺพาราเม มิคารมาตุปาสาเท.\csromanspacer{} เตน โข ปน สมเยน วาเสฏฺฐภารทฺวาชา ภิกฺขูสุ ปริวสนฺติ ภิกฺขุภาวํ อากงฺขมานา.\csromanspacer{} อถ โข ภควา, สายนฺหสมยํ ปฏิสลฺลานา วุฏฺฐิโต ปาสาทา โอโรหิตฺวา ปาสาทปจฺฉายายํ\footnote{ปาสาทจฺฉายายํ (ก)} อพฺโภกาเส จงฺกมติ.}

\proseitem{112}{อทฺทสา โข วาเสฏฺโฐ ภควนฺตํ สยนฺหสมยํ ปฏิสลฺลานา วุฏฺฐิตํ ปาสาทา โอโรหิตฺวา ปาสาทปจฺฉายายํ อพฺโภกาเส จงฺกมนฺตํ, ทิสฺวาน ภารทฺวาชํ อามนฺเตสิ “อยํ อาวุโส ภารทฺวาช ภควา สายนฺหสมยํ ปฏิสลฺลานา วุฏฺฐิโต ปาสาทา โอโรหิตฺวา ปาสาทปจฺฉายายํ อพฺโภกาเส จงฺกมติ.\csromanspacer{} อายามาวุโส ภารทฺวาช เยน ภควา เตนุปสงฺกมิสฺสาม, อปฺเปว นาม ลเภยฺยาม ภควโต สนฺติกา\footnote{สมฺมุขา (สฺยา, ก)} ธมฺมิํ กถํ สวนายา”ติ. “เอวมาวุโส”ติ โข ภารทฺวาโช วาเสฏฺฐสฺส ปจฺจสฺโสสิ.}

\proseitem{113}{อถ โข วาเสฏฺฐภารทฺวาชา เยน ภควา เตนุปสงฺกมิํสุ, อุปสงฺกมิตฺวา ภควนฺตํ อภิวาเทตฺวา ภควนฺตํ จงฺกมนฺตํ อนุจงฺกมิํสุ.\csromanspacer{} อถ โข ภควา วาเสฏฺฐํ อามนฺเตสิ “ตุเมฺห ขฺวตฺถ วาเสฏฺฐ พฺราหฺมณชจฺจา พฺราหฺมณกุลีนา พฺราหฺมณกุลา อคารสฺมา อนคาริยํ ปพฺพชิตา, กจฺจิ โว วาเสฏฺฐ พฺราหฺมณา น อกฺโกสนฺติ น ปริภาสนฺตี”ติ.\csromanspacer{} ตคฺฆ โน ภนฺเต พฺราหฺมณา อกฺโกสนฺติ ปริภาสนฺติ อตฺตรูปาย ปริภาสาย ปริปุณฺณาย โน อปริปุณฺณายาติ.\csromanspacer{} ยถา กถํ ปน โว วาเสฏฺฐ พฺราหฺมณา อกฺโกสนฺติ ปริภาสนฺติ อตฺตรูปาย ปริภาสาย ปริปุณฺณาย โน อปริปุณฺณายาติ.\csromanspacer{} พฺราหฺมณา ภนฺเต เอวมาหํสุ “พฺราหฺมโณว เสฏฺโฐ วณฺโณ, หีนา อญฺเญ วณฺณา\footnote{หีโน อญฺโญ วณฺโณ (สี, อิ, ม 2 มธุรสุตฺต)}.\csromanspacer{} พฺราหฺมโณว สุกฺโก}

\vspace{\tipitakaparskip}
\csromancenter{อคฺคญฺญสุตฺต วณฺโณ, กณฺหา อญฺเญ วณฺณา\footnote{กโณฺห อญฺโญ วณฺโณ (สี, อิ, ม. 2 มธุรสุตฺต)}.\csromanspacer{} พฺราหฺมณาว สุชฺฌนฺติ, โน อพฺราหฺมณา.\csromanspacer{} พฺราหฺมณาว\footnote{พฺราหฺมณา (สฺยา)} พฺรหฺมุโน ปุตฺตา โอรสา มุขโต ชาตา พฺรหฺมชา พฺรหฺมนิมฺมิตา พฺรหฺมทายาทา.\csromanspacer{} เต ตุเมฺห เสฏฺฐํ วณฺณํ หิตฺวา หีนมตฺถ วณฺณํ อชฺฌุปคตา ยทิทํ มุณฺฑเก สมณเก อิพฺเภ กเณฺห พนฺธุปาทาปจฺเจ.\csromanspacer{} ตยิทํ น สาธุ, ตยิทํ นปฺปติรูปํ, ยํ ตุเมฺห เสฏฺฐํ วณฺณํ หิตฺวา หีนมตฺถ วณฺณํ อชฺฌุปคตา ยทิทํ มุณฺฑเก สมณเก อิพฺเภ กเณฺห พนฺธุปาทาปจฺเจติ.\csromanspacer{} เอวํ โข โน ภนฺเต พฺราหฺมณา อกฺโกสนฺติ ปริภาสนฺติ อตฺตรูปาย ปริภาสาย ปริปุณฺณาย, โน อปริปุณฺณายาติ.}
\proseitem{114}{\csromanfolio{67}ตคฺฆ โว วาเสฏฺฐ พฺราหฺมณา โปราณํ อสฺสรนฺตา เอวมาหํสุ “พฺราหฺมโณว เสฏฺโฐ วณฺโณ, หีนา อญฺเญ วณฺณา.\csromanspacer{} พฺราหฺมโณว สุกฺโก วณฺโณ, กณฺหา อญฺเญ วณฺณา.\csromanspacer{} พฺราหฺมณาว สุชฺฌนฺติ, โน อพฺราหฺมณา.\csromanspacer{} พฺราหฺมณาว พฺรหฺมุโน ปุตฺตา โอรสา มุขโต ชาตา พฺรหฺมชา พฺรหฺมนิมฺมิตา พฺรหฺมทายาทา”ติ.\csromanspacer{} ทิสฺสนฺติ โข ปน วาเสฏฺฐ พฺราหฺมณานํ พฺราหฺมณิโย อุตุนิโยปิ คพฺภินิโยปิ วิชายมานาปิ ปายมานาปิ.\csromanspacer{} เต จ พฺราหฺมณา โยนิชาว สมานา เอวมาหํสุ “พฺราหฺมโณว เสฏฺโฐ วณฺโณ, หีนา อญฺเญ วณฺณา.\csromanspacer{} พฺราหฺมโณว สุกฺโก วณฺโณ, กณฺหา อญฺเญ วณฺณา.\csromanspacer{} พฺราหฺมณาว สุชฺฌนฺติ, โน อพฺราหฺมณา.\csromanspacer{} พฺราหฺมณาว พฺรหฺมุโน ปุตฺตา โอรสา มุขโต ชาตา พฺรหฺมชา พฺรหฺมนิมฺมิตา พฺรหฺมทายาทา”ติ.\csromanspacer{} เต\footnote{เต จ (สฺยา, ก)} พฺรหฺมานญฺเจว อพฺภาจิกฺขนฺติ, มุสา จ ภาสนฺติ, พหุญฺจ อปุญฺญํ ปสวนฺติ.}

\csromanmatikaanchor{mtk.33}
\csromantocmarkref{section}{จตุวณฺณสุทฺธิ}{mtk.33}
\bookmark[dest={mtk.33},level=1]{จตุวณฺณสุทฺธิ}
\csromanheader{1}{\textbf{จตุวณฺณสุทฺธิ}}
\proseitem{115}{จตฺตาโรเม วาเสฏฺฐ วณฺณา ขตฺติยา พฺราหฺมณา เวสฺสา สุทฺทา.\csromanspacer{} ขตฺติโยปิ โข วาเสฏฺฐ อิเธกจฺโจ ปาณาติปาตี โหติ อทินฺนาทายี กาเมสุมิจฺฉาจารี มุสาวาที ปิสุณวาโจ ผรุสวาโจ สมฺผปฺปลาปี อภิชฺฌาลุ พฺยาปนฺนจิตฺโต มิจฺฉาทิฏฺฐี.\csromanspacer{} อิติ โข วาเสฏฺฐ เยเม ธมฺมา อกุสลา อกุสลสงฺขาตา สาวชฺชา สาวชฺชสงฺขาตา อเสวิตพฺพา อเสวิตพฺพสงฺขาตา น-อลมริยา น-อลมริยสงฺขาตา กณฺหา กณฺหวิปากา วิญฺญุครหิตา, ขตฺติเยปิ เต\footnote{โข วาเสฏฺฐ (ก)} อิเธกจฺเจ สนฺทิสฺสนฺติ.\csromanspacer{} พฺราหฺมโณปิ โข วาเสฏฺฐ ฯเปฯ.\csromanspacer{} เวสฺโสปิ โข วาเสฏฺฐ ฯเปฯ.\csromanspacer{} สุทฺโทปิ โข วาเสฏฺฐ \csromanfolio{68}อิเธกจฺโจ ปาณาติปาตี โหติ อทินฺนาทายี กาเมสุมิจฺฉาจารี มุสาวาที ปิสุณวาโจ ผรุสวาโจ สมฺผปฺปลาปี อภิชฺฌาลุ พฺยาปนฺนจิตฺโต มิจฺฉาทิฏฺฐี.\csromanspacer{} อิติ โข วาเสฏฺฐ เยเม ธมฺมา อกุสลา อกุสลสงฺขาตา ฯเปฯ กณฺหา กณฺหวิปากา วิญฺญุครหิตา, สุทฺเทปิ เต อิเธกจฺเจ สนฺทิสฺสนฺติ.}

\prose{ขตฺติโยปิ โข วาเสฏฺฐ อิเธกจฺโจ ปาณาติปาตา ปฏิวิรโต โหติ, อทินฺนาทานา ปฏิวิรโต, กาเมสุมิจฺฉาจารา ปฏิวิรโต, มุสาวาทา ปฏิวิรโต, ปิสุณาย วาจาย ปฏิวิรโต, ผรุสาย วาจาย ปฏิวิรโต, สมฺผปฺปลาปา ปฏิวิรโต, อนภิชฺฌาลุ, อพฺยาปนฺนจิตฺโต, สมฺมาทิฏฺฐี.\csromanspacer{} อิติ โข วาเสฏฺฐ เยเม ธมฺมา กุสลา กุสลสงฺขาตา อนวชฺชา อนวชฺชสงฺขาตา เสวิตพฺพา เสวิตพฺพสงฺขาตา อลมริยา อลมริยสงฺขาตา สุกฺกา สุกฺกวิปากา วิญฺญุปฺปสตฺถา, ขตฺติเยปิ เต อิเธกจฺเจ สนฺทิสฺสนฺติ.\csromanspacer{} พฺราหฺมโณปิ โข วาเสฏฺฐ ฯเปฯ.\csromanspacer{} เวสฺโสปิ โข วาเสฏฺฐ ฯเปฯ.\csromanspacer{} สุทฺโทปิ โข วาเสฏฺฐ อิเธกจฺโจ ปาณาติปาตา ปฏิวิรโต โหติ ฯเปฯ อนภิชฺฌาลุ, อพฺยาปนฺนจิตฺโต, สมฺมาทิฏฺฐี.\csromanspacer{} อิติ โข วาเสฏฺฐ เยเม ธมฺมา กุสลา กุสลสงฺขาตา อนวชฺชา อนวชฺชสงฺขาตา เสวิตพฺพา เสวิตพฺพสงฺขาตา อลมริยา อลมริยสงฺขาตา สุกฺกา สุกฺกวิปากา วิญฺญุปฺปสตฺถา, สุทฺเทปิ เต อิเธกจฺเจ สนฺทิสฺสนฺติ.}

\proseitem{116}{อิเมสุ โข วาเสฏฺฐ จตูสุ วณฺเณสุ เอวํ อุภยโวกิณฺเณสุ วตฺตมาเนสุ กณฺห สุกฺเกสุ ธมฺเมสุ วิญฺญุครหิเตสุ เจว วิญฺญุปฺปสตฺเถสุ จ ยเทตฺถ พฺราหฺมณา เอวมาหํสุ “พฺราหฺมโณว เสฏฺโฐ วณฺโณ, หีนา อญฺเญ วณฺณา.\csromanspacer{} พฺราหฺมโณว สุกฺโก วณฺโณ, กณฺหา อญฺเญ วณฺณา.\csromanspacer{} พฺราหฺมณาว สุชฺฌนฺติ, โน อพฺราหฺมณา.\csromanspacer{} พฺราหฺมณาว พฺรหฺมุโน ปุตฺตา โอรสา มุขโต ชาตา พฺรหฺมชา พฺรหฺมนิมฺมิตา พฺรหฺมทายาทา”ติ.\csromanspacer{} ตํ เตสํ วิญฺญู นานุชานนฺติ.\csromanspacer{} ตํ กิสฺส เหตุ, อิเมสํ หิ วาเสฏฺฐ จตุนฺนํ วณฺณานํ โย โหติ ภิกฺขุ อรหํ ขีณาสโว วุสิตวา กตกรณีโย โอหิตภาโร อนุปฺปตฺตสทตฺโถ ปริกฺขีณภวสํโยชโน สมฺมทญฺญาวิมุตฺโต, โส เนสํ อคฺคมกฺขายติ ธมฺเมเนว, โน อธมฺเมน.\csromanspacer{} ธมฺโม หิ วาเสฏฺฐ เสฏฺโฐ ชเนตสฺมิํ ทิฏฺเฐ เจว ธมฺเม อภิสมฺปรายญฺจ.}

\proseitem{117}{ตทมินาเปตํ วาเสฏฺฐ ปริยาเยน เวทิตพฺพํ, ยถา ธมฺโมว เสฏฺโฐ ชเนตสฺมิํ ทิฏฺเฐ เจว ธมฺเม อภิสมฺปรายญฺจ.}

\csromanheader{2}{4. อคฺคญฺญสุตฺต}
\prose{\csromanfolio{69}ชานาติ โข\footnote{โข ปน (ก)} วาเสฏฺฐ ราชา ปเสนทิ โกสโล “สมโณ โคตโม อนนฺตรา\footnote{อนุตฺตโร (พหูสุ)} สกฺยกุลา ปพฺพชิโต”ติ.\csromanspacer{} สกฺยา โข ปน วาเสฏฺฐ รญฺโญ ปเสนทิสฺส โกสลสฺส อนุยุตฺตา\footnote{อนนฺตรา อนุยนฺตา (สฺยา), อนนฺตรา อนุยุตฺตา (ก)} ภวนฺติ.\csromanspacer{} กโรนฺติ โข วาเสฏฺฐ สกฺยา รญฺเญ ปเสนทิมฺหิ โกสเล นิปจฺจการํ อภิวาทนํ ปจฺจุฏฺฐานํ อญฺชลิกมฺมํ สามีจิกมฺมํ.\csromanspacer{} อิติ โข วาเสฏฺฐ ยํ กโรนฺติ สกฺยา รญฺเญ ปเสนทิมฺหิ โกสเล นิปจฺจการํ อภิวาทนํ ปจฺจุฏฺฐานํ อญฺชลิกมฺมํ สามีจิกมฺมํ.\csromanspacer{} กโรติ ตํ ราชา ปเสนทิ โกสโล ตถาคเต นิปจฺจการํ อภิวาทนํ ปจฺจุฏฺฐานํ อญฺชลิกมฺมํ สามีจิกมฺมํ, น นํ\footnote{นนุ (พหูสุ)} “สุชาโต สมโณ โคตโม, ทุชฺชาโตหมสฺมิ.\csromanspacer{} พลวา สมโณ โคตโม, ทุพฺพโลหมสฺมิ.\csromanspacer{} ปาสาทิโก สมโณ โคตโม, ทุพฺพณฺโณหมสฺมิ.\csromanspacer{} มเหสกฺโข สมโณ โคตโม, อปฺเปสกฺโขหมสฺมี”ติ.\csromanspacer{} อถ โข นํ ธมฺมํเยว สกฺกโรนฺโต ธมฺมํ ครุํ กโรนฺโต ธมฺมํ มาเนนฺโต ธมฺมํ ปูเชนฺโต ธมฺมํ อปจายมาโน เอวํ ราชา ปเสนทิ โกสโล ตถาคเต นิปจฺจการํ กาโรติ อภิวาทนํ ปจฺจุฏฺฐานํ อญฺชลิกมฺมํ สามีจิกมฺมํ.\csromanspacer{} อิมินาปิ โข เอตํ วาเสฏฺฐ ปริยาเยน เวทิตพฺพํ, ยถา ธมฺโม เสฏฺโฐ ชเนตสฺมิํ ทิฏฺเฐ เจว ธมฺเม อภิสมฺปรายญฺจ.}

\proseitem{118}{ตุเมฺห ขฺวตฺถ วาเสฏฺฐ นานาชจฺจา นานานามา นานาโคตฺตา นานากุลา อคารสฺมา อนคาริยํ ปพฺพชิตา, “เก ตุเมฺห”ติ ปุฏฺฐา สมานา “สมณา สกฺยปุตฺติยามฺหา”ติ ปฏิชานาถ.\csromanspacer{} ยสฺส โข ปนสฺส วาเสฏฺฐ ตถาคเต สทฺธา นิวิฏฺฐา มูลชาตา ปติฏฺฐิตา ทฬฺหา อสํหาริยา สมเณน วา พฺราหฺมเณน วา เทเวน วา มาเรน วา พฺรหฺมุนา วา เกนจิ วา โลกสฺมิํ, ตสฺเสตํ กลฺลํ วจนาย “ภควโตมฺหิ ปุตฺโต โอรโส มุขโต ชาโต ธมฺมโช ธมฺมนิมฺมิโต ธมฺมทายาโท”ติ.\csromanspacer{} ตํ กิสฺส เหตุ, ตถาคตสฺส เหตํ วาเสฏฺฐ อธิวจนํ “ธมฺมกาโย” อิติปิ “พฺรหฺมกาโย” อิติปิ “ธมฺมภูโต” อิติปิ “พฺรหฺมภูโต” อิติปิ.}

\proseitem{119}{โหติ โข โส วาเสฏฺฐ สมโย ยํ กทาจิ กรหจิ ทีฆสฺส อทฺธุโน อจฺจเยน อยํ โลโก สํวฏฺฏติ, สํวฏฺฏมาเน โลเก เยภุยฺเยน สตฺตา อาภสฺสรสํวตฺตนิกา โหนฺติ.\csromanspacer{} เต ตตฺถ \csromanfolio{70}โหนฺติ มโนมยา ปีติภกฺขา สยํปภา อนฺตลิกฺขจรา สุภฏฺฐายิโน จิรํ ทีฆมทฺธานํ ติฏฺฐนฺติ.}

\prose{โหติ โข โส วาเสฏฺฐ สมโย ยํ กทาจิ กรหจิ ทีฆสฺส อทฺธุโน อจฺจเยน อยํ โลโก วิวฏฺฏติ, วิวฏฺฏมาเน โลเก เยภุยฺเยน สตฺตา อาภสฺสรกายา จวิตฺวา อิตฺถตฺตํ อาคจฺฉนฺติ.\csromanspacer{} เตธ โหนฺติ มโนมยา ปีติภกฺขา สยํปภา อนฺตลิกฺขจรา สุภฏฺฐายิโน จิรํ ทีฆมทฺธานํ ติฏฺฐนฺติ.}

\csromanmatikaanchor{mtk.34}
\csromantocmarkref{section}{รสปถวิปาตุภาว}{mtk.34}
\bookmark[dest={mtk.34},level=1]{รสปถวิปาตุภาว}
\csromanheader{1}{\textbf{รสปถวิปาตุภาว}}
\proseitem{120}{เอโกทกีภูตํ โข ปน วาเสฏฺฐ เตน สมเยน โหติ อนฺธกาโร อนฺธการติมิสา.\csromanspacer{} น จนฺทิมสูริยา ปญฺญายนฺติ, น นกฺขตฺตานิ ตารกรูปานิ ปญฺญายนฺติ, น รตฺตินฺทิวา ปญฺญายนฺติ, น มาสฑฺฒมาสา ปญฺญายนฺติ, น อุตุสํวจฺฉารา ปญฺญายนฺติ, น อิตฺถิปุมา ปญฺญายนฺติ, สตฺตา สตฺตาเตฺวว สงฺขฺยํ คจฺฉนฺติ.\csromanspacer{} อถ โข เตสํ วาเสฏฺฐ สตฺตานํ กทาจิ กรหจิ ทีฆสฺส อทฺธุโน อจฺจเยน รสปถวี อุทกสฺมิํ สมตนิ\footnote{สมตานิ (พหูสุ)}, เสยฺยถาปิ นาม ปยโส ตตฺตสฺส\footnote{ปยตตฺตสฺส (สฺยา)} นิพฺพายมานสฺส อุปริ สนฺตานกํ โหติ, เอวเมว ปาตุรโหสิ.\csromanspacer{} สา อโหสิ วณฺณสมฺปนฺนา คนฺธสมฺปนฺนา รสสมฺปนฺนา, เสยฺยถาปิ นาม สมฺปนฺนํ วา สปฺปิ สมฺปนฺนํ วา นวนีตํ เอวํวณฺณา อโหสิ.\csromanspacer{} เสยฺยถาปิ นาม ขุทฺทมธุํ\footnote{ขุทฺทํ มธุํ (ก-สี)} อเนฬกํ\footnote{อเนลกํ (สี, อิ)}, เอวมสฺสาทา อโหสิ.\csromanspacer{} อถ โข วาเสฏฺฐ อญฺญตโร สตฺโต โลลชาติโก “อมฺโภ กิเมวิทํ ภวิสฺสตี”ติ รสปถวิํ องฺคุลิยา สายิ.\csromanspacer{} ตสฺส รสปถวิํ องฺคุลิยา สายโต อจฺฉาเทสิ, ตณฺหา จสฺส โอกฺกมิ.\csromanspacer{} อญฺเญปิ โข วาเสฏฺฐ สตฺตา ตสฺส สตฺตสฺส ทิฏฺฐานุคติํ อาปชฺชมานา รสปถวิํ องฺคุลิยา สายิํสุ.\csromanspacer{} เตสํ รสปถวิํ องฺคุลิยา สายตํ อจฺฉาเทสิ, ตณฺหา จ เตสํ โอกฺกมิ.}

\csromanmatikaanchor{mtk.35}
\csromantocmarkref{section}{จนฺทิมสูริยาทิปาตุภาว}{mtk.35}
\bookmark[dest={mtk.35},level=1]{จนฺทิมสูริยาทิปาตุภาว}
\csromanheader{1}{\textbf{จนฺทิมสูริยาทิปาตุภาว}}
\proseitem{121}{อถ โข เต วาเสฏฺฐ สตฺตา รสปถวิํ หตฺเถหิ อาลุปฺปการกํ อุปกฺกมิํสุ ปริภุญฺชิตุํ.\csromanspacer{} ยโต โข เต\footnote{ยโต โข (สี, สฺยา, อิ)} วาเสฏฺฐ สตฺตา}

\vspace{\tipitakaparskip}
\csromancenter{อคฺคญฺญสุตฺต รสปถวิํ หตฺเถหิ อาลุปฺปการกํ อุปกฺกมิํสุ ปริภุญฺชิตุํ.\csromanspacer{} อถ เตสํ สตฺตานํ สยํปภา อนฺตรธายิ.\csromanspacer{} สยํปภาย อนฺตรหิตาย จนฺทิมสูริยา ปาตุรเหสุํ.\csromanspacer{} จนฺทิมสูริเยสุ ปาตุภูเตสุ นกฺขตฺตานิ ตารกรูปานิ ปาตุรเหสุํ.\csromanspacer{} นกฺขตฺเตสุ ตารกรูเปสุ ปาตุภูเตสุ รตฺตินฺทิวา ปญฺญายิํสุ.\csromanspacer{} รตฺตินฺทิเวสุ ปญฺญายมาเนสุ มาสฑฺฒมาสา ปญฺญายิํสุ.\csromanspacer{} มาสฑฺฒมาเสสุ ปญฺญายมาเนสุ อุตุสํวจฺฉรา ปญฺญายิํสุ.\csromanspacer{} เอตฺตาวตา โข วาเสฏฺฐ อยํ โลโก ปุน วิวฏฺโฏ โหติ.}
\proseitem{122}{\csromanfolio{71}อถ โข เต วาเสฏฺฐ สตฺตา รสปถวิํ ปริภุญฺชนฺตา ตมฺภกฺขา\footnote{ตพฺภกฺขา (สฺยา)} ตทาหารา จิรํ ทีฆมทฺธานํ อฏฺฐํสุ.\csromanspacer{} ยถา ยถา โข เต วาเสฏฺฐ สตฺตา รสปถวิํ ปริภุญฺชนฺตา ตมฺภกฺขา ตทาหารา จิรํ ทีฆมทฺธานํ อฏฺฐํสุ, ตถา ตถา เตสํ สตฺตานํ (รสปถวิํ ปริภุญฺชนฺตานํ)\footnote{( ) สี-สฺยา-อิ-โปตฺถเกสุ นตฺถิ,} ขรตฺตญฺเจว กายสฺมิํ โอกฺกมิ, วณฺณเววณฺณตา\footnote{วณฺณเววชฺชตา (ฏีกา)} จ ปญฺญายิตฺถ.\csromanspacer{} เอกิทํ สตฺตา วณฺณวนฺโต โหนฺติ, เอกิทํ สตฺตา ทุพฺพณฺณา.\csromanspacer{} ตตฺถ เย เต สตฺตา วณฺณวนฺโต, เต ทุพฺพณฺเณ สตฺเต อติมญฺญนฺติ “มยเมเตหิ วณฺณวนฺตตรา, อเมฺหเหเต ทุพฺพณฺณตรา”ติ.\csromanspacer{} เตสํ วณฺณาติมานปจฺจยา มานาติมานชาติกานํ รสปถวี อนฺตรธาติ, รสาย ปถวิยา อนฺตรหิตาย สนฺนิปติํสุ, สนฺนิปติตฺวา อนุตฺถุนิํสุ “อโหรสํ อโหรสนฺ”ติ.\csromanspacer{} ตเทตรหิปิ มนุสฺสา กญฺจิเทว สุรสํ\footnote{สาธุรสํ (สี, สฺยา, อิ)} ลภิตฺวา เอวมาหํสุ “อโหรสํ อโหรสนฺ”ติ.\csromanspacer{} ตเทว โปราณํ อคฺคญฺญํ อกฺขรํ อนุสรนฺติ, น เตฺววสฺส อตฺถํ อาชานนฺติ.}

\csromanmatikaanchor{mtk.36}
\csromantocmarkref{section}{ภูมิปปฺปฏกปาตุภาว}{mtk.36}
\bookmark[dest={mtk.36},level=1]{ภูมิปปฺปฏกปาตุภาว}
\csromanheader{1}{\textbf{ภูมิปปฺปฏกปาตุภาว}}
\proseitem{123}{อถ โข เตสํ วาเสฏฺฐ สตฺตานํ รสาย ปถวิยา อนฺตรหิตาย ภูมิปปฺปฏโก ปาตุรโหสิ.\csromanspacer{} เสยฺยถาปิ นาม อหิจฺฉตฺตโก, เอวเมว ปาตุรโหสิ.\csromanspacer{} โส อโหสิ วณฺณสมฺปนฺโน คนฺธสมฺปนฺโน รสสมฺปนฺโน, เสยฺยถาปิ นาม สมฺปนฺนํ วา สปฺปิ สมฺปนฺนํ วา นวนีตํ เอวํวณฺโณ อโหสิ.\csromanspacer{} เสยฺยถาปิ นาม ขุทฺทมธุํ อเนฬกํ, เอวมสฺสาโท อโหสิ.}

\prose{\csromanfolio{72}อถ โข เต วาเสฏฺฐ สตฺตา ภูมิปปฺปฏกํ อุปกฺกมิํสุ ปริภุญฺชิตุํ.\csromanspacer{} เต ตํ ปริภุญฺชนฺตา ตมฺภกฺขา ตทาหารา จิรํ ทีฆมทฺธานํ อฏฺฐํสุ.\csromanspacer{} ยถา ยถา โข เต วาเสฏฺฐ สตฺตา ภูมิปปฺปฏกํ ปริภุญฺชนฺตา ตมฺภกฺขา ตทาหารา จิรํ ทีฆมทฺธานํ อฏฺฐํสุ, ตถา ตถา เตสํ สตฺตานํ ภิยฺโยโส มตฺตาย ขรตฺตญฺเจว กายสฺมิํ โอกฺกมิ, วณฺณเววณฺณตา จ ปญฺญายิตฺถ.\csromanspacer{} เอกิทํ สตฺตา วณฺณวนฺโต โหนฺติ, เอกิทํ สตฺตา ทุพฺพณฺณา.\csromanspacer{} ตตฺถ เย เต สตฺตา วณฺณวนฺโต, เต ทุพฺพณฺเณ สตฺเต อติมญฺญนฺติ “มยเมเตหิ วณฺณวนฺตตรา, อเมฺหเหเต ทุพฺพณฺณตรา”ติ.\csromanspacer{} เตสํ วณฺณาติมานปจฺจยา มานาติมานชาติกานํ ภูมิปปฺปฏโก อนฺตรธายิ.}

\csromanmatikaanchor{mtk.37}
\csromantocmarkref{section}{ปทาลตาปาตุภาว}{mtk.37}
\bookmark[dest={mtk.37},level=1]{ปทาลตาปาตุภาว}
\csromanheader{1}{\textbf{ปทาลตาปาตุภาว}}
\proseitem{124}{ภูมิปปฺปฏเก อนฺตรหิเต ปทาลตา\footnote{ภทฺทาลตา (สี)} ปาตุรโหสิ, เสยฺยถาปิ นาม กลมฺพุกา\footnote{กลมฺพกา (สฺยา)}, เอวเมว ปาตุรโหสิ.\csromanspacer{} สา อโหสิ วณฺณสมฺปนฺนา คนฺธสมฺปนฺนา รสสมฺปนฺนา, เสยฺยถาปิ นาม สมฺปนฺนํ วา สปฺปิ สมฺปนฺนํ วา นวนีตํ เอวํวณฺณา อโหสิ.\csromanspacer{} เสยฺยถาปิ นาม ขุทฺทมธุํ อเนฬกํ, เอวมสฺสาทา อโหสิ.}

\prose{อถ โข เต วาเสฏฺฐ สตฺตา ปทาลตํ อุปกฺกมิํสุ ปริภุญฺชิตุํ.\csromanspacer{} เต ตํ ปริภุญฺชนฺตา ตมฺภกฺขา ตทาหารา จิรํ ทีฆมทฺธานํ อฏฺฐํสุ.\csromanspacer{} ยถา ยถา โข เต วาเสฏฺฐ สตฺตา ปทาลตํ ปริภุญฺชนฺตา ตมฺภกฺขา ตทาหารา จิรํ ทีฆมทฺธานํ อฏฺฐํสุ, ตถา ตถา เตสํ สตฺตานํ ภิยฺโยโส มตฺตาย ขรตฺตญฺเจว กายสฺมิํ โอกฺกมิ, วณฺณเววณฺณตา จ ปญฺญายิตฺถ.\csromanspacer{} เอกิทํ สตฺตา วณฺณวนฺโต โหนฺติ, เอกิทํ สตฺตา ทุพฺพณฺณา.\csromanspacer{} ตตฺถ เย เต สตฺตา วณฺณวนฺโต, เต ทุพฺพณฺเณ สตฺเต อติมญฺญนฺติ “มยเมเตหิ วณฺณวนฺตตรา, อเมฺหเหเต ทุพฺพณฺณตรา”ติ.\csromanspacer{} เตสํ วณฺณาติมานปจฺจยา มานาติมานชาติกานํ ปทาลตา อนฺตรธายิ.}

\prose{ปทาลตาย อนฺตรหิตาย สนฺนิปติํสุ, สนฺนิปติตฺวา อนุตฺถุนิํสุ “อหุ วต โน, อหายิ วต โน ปทาลตา”ติ.\csromanspacer{} ตเทตรหิปิ มนุสฺสา เกนจิ\footnote{เกนจิเทว (สี, สฺยา, อิ)} ทุกฺขธมฺเมน ผุฏฺฐา เอวมาหํสุ “อหุ วต โน,}

\vspace{\tipitakaparskip}
\csromancenter{อคฺคญฺญสุตฺต อหายิ วต โน”ติ.\csromanspacer{} ตเทว โปราณํ อคฺคญฺญํ อกฺขรํ อนุสรนฺติ, น เตฺววสฺส อตฺถํ อาชานนฺติ.}
\csromanmatikaanchor{mtk.38}
\csromantocmarkref{section}{อกฏฺฐปากสาลิปาตุภาว}{mtk.38}
\bookmark[dest={mtk.38},level=1]{อกฏฺฐปากสาลิปาตุภาว}
\csromanheader{1}{\textbf{อกฏฺฐปากสาลิปาตุภาว}}
\proseitem{125}{\csromanfolio{73}อถ โข เตสํ วาเสฏฺฐ สตฺตานํ ปทาลตาย อนฺตรหิตาย อกฏฺฐปาโก สาลิ ปาตุรโหสิ อกโณ อถุโส สุทฺโธ สุคนฺโธ ตณฺฑุลปฺผโล.\csromanspacer{} ยํ ตํ สายํ สายมาสาย อาหรนฺติ, ปาโต ตํ โหติ ปกฺกํ ปฏิวิรุฬฺหํ.\csromanspacer{} ยํ ตํ ปาโต ปาตราสาย อาหรนฺติ, สายํ ตํ โหติ ปกฺกํ ปฏิวิรูฬฺหํ, นาปทานํ ปญฺญายติ.\csromanspacer{} อถ โข เต วาเสฏฺฐ สตฺตา อกฏฺฐปากํ สาลิํ ปริภุญฺชนฺตา ตมฺภกฺขา ตทาหารา จิรํ ทีฆมทฺธานํ อฏฺฐํสุ.}

\csromanmatikaanchor{mtk.39}
\csromantocmarkref{section}{อิตฺถิปุริสลิงฺคปาตุภาว}{mtk.39}
\bookmark[dest={mtk.39},level=1]{อิตฺถิปุริสลิงฺคปาตุภาว}
\csromanheader{1}{\textbf{อิตฺถิปุริสลิงฺคปาตุภาว}}
\proseitem{126}{ยถา ยถา โข เต วาเสฏฺฐ สตฺตา อกฏฺฐปากํ สาลิํ ปริภุญฺชนฺตา ตมฺภกฺขา ตทาหารา จิรํ ทีฆมทฺธานํ อฏฺฐํสุ, ตถา ตถา เตสํ สตฺตานํ ภิยฺโยโส มตฺตาย ขรตฺตญฺเจว กายสฺมิํ โอกฺกมิ, วณฺณเววณฺณตา จ ปญฺญายิตฺถ, อิตฺถิยา จ อิตฺถิลิงฺคํ ปาตุรโหสิ ปุริสสฺส จ ปุริสลิงฺคํ, อิตฺถี จ ปุริสํ อติเวลํ อุปนิชฺฌายติ ปุริโส จ อิตฺถิํ.\csromanspacer{} เตสํ อติเวลํ อญฺญมญฺญํ อุปนิชฺฌายตํ สาราโค อุทปาทิ, ปริฬาโห กายสฺมิํ โอกฺกมิ.\csromanspacer{} เต ปริฬาหปจฺจยา เมถุนํ ธมฺมํ ปฏิเสวิํสุ.}

\prose{เย โข ปน เต วาเสฏฺฐ เตน สมเยน สตฺตา ปสฺสนฺติ เมถุนํ ธมฺมํ ปฏิเสวนฺเต, อญฺเญ ปํสุํ ขิปนฺติ, อญฺเญ เสฏฺฐิํ ขิปนฺติ, อญฺเญ โคมยํ ขิปนฺติ “นสฺส อสุจิ\footnote{วสลิ (สฺยา), วสลี (ก)} นสฺส อสุจิ1”ติ, “กถํ หิ นาม สตฺโต สตฺตสฺส เอวรูปํ กริสฺสตี”ติ.\csromanspacer{} ตเทตรหิปิ มนุสฺสา เอกจฺเจสุ ชนปเทสุ วธุยา นิพฺพุยฺหมานาย\footnote{นิวยฺหมานาย, นิคฺคยฺหมานาย (ก)} อญฺเญ ปํสุํ ขิปนฺติ, อญฺเญ เสฏฺฐิํ ขิปนฺติ, อญฺเญ โคมยํ ขิปนฺติ.\csromanspacer{} ตเทว โปราณํ อคฺคญฺญํ อกฺขรํ อนุสรนฺติ, น เตฺววสฺส อตฺถํ อาชานนฺติ.}

\csromanmatikaanchor{mtk.40}
\csromantocmarkref{section}{เมถุนธมฺมสมาจาร}{mtk.40}
\bookmark[dest={mtk.40},level=1]{เมถุนธมฺมสมาจาร}
\csromanheader{1}{\textbf{เมถุนธมฺมสมาจาร}}
\proseitem{127}{\csromanfolio{74}อธมฺมสมฺมตํ โข ปน\footnote{อธมฺมสมฺมตํ ตํ โข ปน (สฺยา), อธมฺมสมฺมตํ โข ปน ตํ (?)} วาเสฏฺฐ เตน สมเยน โหติ, ตเทตรหิ ธมฺมสมฺมตํ.\csromanspacer{} เย โข ปน วาเสฏฺฐ เตน สมเยน สตฺตา เมถุนํ ธมฺมํ ปฏิเสวนฺติ, เต มาสมฺปิ เทฺวมาสมฺปิ น ลภนฺติ คามํ วา นิคมํ วา ปวิสิตุํ.\csromanspacer{} ยโต โข เต วาเสฏฺฐ สตฺตา ตสฺมิํ อสทฺธมฺเม อติเวลํ ปาตพฺยตํ อาปชฺชิํสุ.\csromanspacer{} อถ อคารานิ อุปกฺกมิํสุ กาตุํ ตสฺเสว อสทฺธมฺมสฺส ปฏิจฺฉาทนตฺถํ.\csromanspacer{} อถ โข วาเสฏฺฐ อญฺญตรสฺส สตฺตสฺส อลสชาติกสฺส เอตทโหสิ “อมฺโภ กิเมวาหํ วิหญฺญามิ สาลิํ อาหรนฺโต สายํ สายมาสาย ปาโต ปาตราสาย, ยํนูนาหํ สาลิํ อาหเรยฺยํ สกิเทว\footnote{สกิํเทว (ก)} สายปาตราสายา”ติ.}

\prose{อถ โข โส วาเสฏฺฐ สตฺโต สาลิํ อาหาสิ สกิเทว สายปาตราสาย.\csromanspacer{} อถ โข วาเสฏฺฐ อญฺญตโร สตฺโต เยน โส สตฺโต เตนุปสงฺกมิ, อุปสงฺกมิตฺวา ตํ สตฺตํ เอตทโวจ “เอหิ โภ สตฺต สาลาหารํ คมิสฺสามา”ติ.\csromanspacer{} อลํ โภ สตฺต อาหโต\footnote{อาหโฏ (สฺยา)} เม สาลิ สกิเทว สายปาตราสายาติ.\csromanspacer{} อถ โข โส วาเสฏฺฐ สตฺโต ตสฺส สตฺตสฺส ทิฏฺฐานุคติํ อาปชฺชมาโน สาลิํ อาหาสิ สกิเทว ทฺวีหาย.\csromanspacer{} เอวํปิ กิร โภ สาธูติ.}

\prose{อถ โข วาเสฏฺฐ อญฺญตโร สตฺโต เยน โส สตฺโต เตนุปสงฺกมิ, อุปสงฺกมิตฺวา ตํ สตฺตํ เอตทโวจ “เอหิ โภ สตฺต สาลาหารํ คมิสฺสามา”ติ.\csromanspacer{} อลํ โภ สตฺต อาหโต เม สาลิ สกิเทว ทฺวีหายาติ.\csromanspacer{} อถ โข โส วาเสฏฺฐ สตฺโต ตสฺส สตฺตสฺส ทิฏฺฐานุคติํ อาปชฺชมาโน สาลิํ อาหาสิ สกิเทว จตูหาย, เอวํปิ กิร โภ สาธูติ.}

\prose{อถ โข วาเสฏฺฐ อญฺญตโร สตฺโต เยน โส สตฺโต เตนุปสงฺกมิ, อุปสงฺกมิตฺวา ตํ สตฺตํ เอตทโวจ “เอหิ โภ สตฺต สาลาหารํ คมิสฺสามา”ติ.\csromanspacer{} อลํ โภ สตฺต อาหโต เม สาลิ สกิเทว จตูหายาติ.\csromanspacer{} อถ โข โส วาเสฏฺฐ สตฺโต ตสฺส}

\vspace{\tipitakaparskip}
\csromancenter{อคฺคญฺญสุตฺต สตฺตสฺส ทิฏฺฐานุคติํ อาปชฺชมาโน สาลิํ อาหาสิ สกิเทว อฏฺฐาหาย, เอวํปิ กิร โภ สาธูติ.}
\prose{\csromanfolio{75}ยโต โข เต วาเสฏฺฐ สตฺตา สนฺนิธิการกํ สาลิํ อุปกฺกมิํสุ ปริภุญฺชิตุํ.\csromanspacer{} อถ กโณปิ ตณฺฑุลํ ปริโยนนฺธิ, ถุโสปิ ตณฺฑุลํ ปริโยนนฺธิ, ลูนมฺปิ นปฺปฏิวิรูฬฺหํ, อปทานํ ปญฺญายิตฺถ, สณฺฑสณฺฑา สาลโย อฏฺฐํสุ.}

\csromanmatikaanchor{mtk.41}
\csromantocmarkref{section}{สาลิวิภาค}{mtk.41}
\bookmark[dest={mtk.41},level=1]{สาลิวิภาค}
\csromanheader{1}{\textbf{สาลิวิภาค}}
\proseitem{128}{อถ โข เต วาเสฏฺฐ สตฺตา สนฺนิปติํสุ, สนฺนิปติตฺวา อนุตฺถุนิํสุ “ปาปกา วต โภ ธมฺมา สตฺเตสุ ปาตุภูตา, มยํ หิ ปุพฺเพ มโนมยา อหุมฺหา ปีติภกฺขา สยํปภา อนฺตลิกฺขจรา สุภฏฺฐายิโน, จิรํ ทีฆมทฺธานํ อฏฺฐมฺหา, เตสํ โน อมฺหากํ กทาจิ กรหจิ ทีฆสฺส อทฺธุโน อจฺจเยน รสปถวี อุทกสฺมิํ สมตนิ, สา อโหสิ วณฺณสมฺปนฺนา คนฺธสมฺปนฺนา รสสมฺปนฺนา.\csromanspacer{} เต มยํ รสปถวิํ หตฺเถหิ อาลุปฺปการกํ อุปกฺกมิมฺห ปริภุญฺชิตุํ, เตสํ โน รสปถวิํ หตฺเถหิ อาลุปฺปการกํ อุปกฺกมตํ ปริภุญฺชิตุํ สยํปภา อนฺตรธายิ, ตาย อนฺตรหิตาย จนฺทิมสูริยา ปาตุรเหสุํ, จนฺทิมสูริเยสุ ปาตุภูเตสุ นกฺขตฺตานิ ตารกรูปานิ ปาตุรเหสุํ, นกฺขตฺเตสุ ตารกรูเปสุ ปาตุภูเตสุ รตฺตินฺทิวา ปญฺญายิํสุ, รตฺตินฺทิเวสุ ปญฺญายมาเนสุ มาสฑฺฒมาสา ปญฺญายิํสุ, มาสฑฺฒมาเสสุ ปญฺญายมาเนสุ อุตุสํวจฺฉรา ปญฺญายิํสุ.\csromanspacer{} เต มยํ รสปถวิํ ปริภุญฺชนฺตา ตมฺภกฺขา ตทาหารา จิรํ ทีฆมทฺธานํ อฏฺฐมฺหา.\csromanspacer{} เตสํ โน ปาปกานํเยว อกุสลานํ ธมฺมานํ ปาตุภาวา รสปถวี อนฺตรธายิ, รสปถวิยา อนฺตรหิตาย ภูมิปปฺปฏโก ปาตุรโหสิ, โส อโหสิ วณฺณสมฺปนฺโน คนฺธสมฺปนฺโน รสสมฺปนฺโน.\csromanspacer{} เต มยํ ภูมิปปฺปฏกํ อุปกฺกมิมฺห ปริภุญฺชิตุํ, เต มยํ ตํ ปริภุญฺชนฺตา ตมฺภกฺขา ตทาหารา จิรํ ทีฆมทฺธานํ อฏฺฐมฺหา.\csromanspacer{} เตสํ โน ปาปกานํเยว อกุสลานํ ธมฺมานํ ปาตุภาวา ภูมิปปฺปฏโก อนฺตรธายิ, ภูมิปปฺปฏเก อนฺตรหิเต ปทาลตา ปาตุรโหสิ, สา อโหสิ วณฺณสมฺปนฺนา คนฺธสมฺปนฺนา รสสมฺปนฺนา.\csromanspacer{} เต มยํ ปทาลตํ อุปกฺกมิมฺห ปริภุญฺชิตุํ, เต มยํ ตํ ปริภุญฺชนฺตา ตมฺภกฺขา ตทาหารา จิรํ ทีฆมทฺธานํ อฏฺฐมฺหา, เตสํ โน ปาปกานํเยว อกุสลานํ ธมฺมานํ \csromanfolio{76}ปาตุภาวา ปทาลตา อนฺตรธายิ, ปทาลตาย อนฺตรหิตาย อกฏฺฐปาโก สาลิ ปาตุรโหสิ อกโณ อถุโส สุทฺโธ สุคนฺโธ ตณฺฑุลปฺผโล.\csromanspacer{} ยํ ตํ สายํ สายมาสาย อาหราม, ปาโต ตํ โหติ ปกฺกํ ปฏิวิรูฬฺหํ.\csromanspacer{} ยํ ตํ ปาโต ปาตราสาย อาหราม, สายํ ตํ โหติ ปกฺกํ ปฏิวิรูฬฺหํ.\csromanspacer{} นาปทานํ ปญฺญายิตฺถ.\csromanspacer{} เต มยํ อกฏฺฐปากํ สาลิํ ปริภุญฺชนฺตา ตมฺภกฺขา ตทาหารา จิรํ ทีฆมทฺธานํ อฏฺฐมฺหา.\csromanspacer{} เตสํ โน ปาปกานํเยว อกุสลานํ ธมฺมานํ ปาตุภาวา กโณปิ ตณฺฑุลํ ปริโยนนฺธิ, ถุโสปิ ตณฺฑุลํ ปริโยนนฺธิ, ลูนมฺปิ นปฺปฏิวิรูฬฺหํ, อปทานํ ปญฺญายิตฺถ, สณฺฑสณฺฑา สาลโย ฐิตา.\csromanspacer{} ยํนูน มยํ สาลิํ วิภเชยฺยาม, มริยาทํ ฐเปยฺยามา”ติ.\csromanspacer{} อถ โข เต วาเสฏฺฐ สตฺตา สาลิํ วิภชิํสุ, มริยาทํ ฐเปสุํ.}

\proseitem{129}{อถ โข วาเสฏฺฐ อญฺญตโร สตฺโต โลลชาติโก สกํ ภาคํ ปริรกฺขนฺโต อญฺญตรํ\footnote{อญฺญสฺส (?)} ภาคํ อทินฺนํ อาทิยิตฺวา ปริภุญฺชิ, ตเมนํ อคฺเคเหสุํ, คเหตฺวา เอตทโวจุํ “ปาปกํ วต โภ สตฺต กโรสิ, ยตฺร หิ นาม สกํ ภาคํ ปริรกฺขนฺโต อญฺญตรํ ภาคํ อทินฺนํ อาทิยิตฺวา ปริภุญฺชสิ, มาสฺสุ โภ สตฺต ปุนปิ เอวรูปมกาสี”ติ. “เอวํ โภ”ติ โข วาเสฏฺฐ โส สตฺโต เตสํ สตฺตานํ ปจฺจสฺโสสิ.\csromanspacer{} ทุติยมฺปิ โข วาเสฏฺฐ โส สตฺโต ฯเปฯ.\csromanspacer{} ตติยมฺปิ โข วาเสฏฺฐ โส สตฺโต สกํ ภาคํ ปริรกฺขนฺโต อญฺญตรํ ภาคํ อทินฺนํ อาทิยิตฺวา ปริภุญฺชิ, ตเมนํ อคฺคเหสุํ, คเหตฺวา เอตทโวจุํ “ปาปกํ วต โภ สตฺต กโรสิ, ยตฺร หิ นาม สกํ ภาคํ ปริรกฺขนฺโต อญฺญตรํ ภาคํ อทินฺนํ อาทิยิตฺวา ปริภุญฺชสิ, มาสฺสุ โภ สตฺต ปุนปิ เอวรูปมกาสี”ติ.\csromanspacer{} อญฺเญ ปาณินา ปหริํสุ, อญฺเญ เลฑฺฑุนา ปหริํสุ, อญฺเญ ทณฺเฑน ปหริํสุ.\csromanspacer{} ตทคฺเค โข วาเสฏฺฐ อทินฺนาทานํ ปญฺญายติ, ครหา ปญฺญายติ, มุสาวาโท ปญฺญายติ, ทณฺฑาทานํ ปญฺญายติ.}

\csromanmatikaanchor{mtk.42}
\csromantocmarkref{section}{มหาสมฺมตราชา}{mtk.42}
\bookmark[dest={mtk.42},level=1]{มหาสมฺมตราชา}
\csromanheader{1}{\textbf{มหาสมฺมตราชา}}
\proseitem{130}{อถ โข เต วาเสฏฺฐ สตฺตา สนฺนิปติํสุ, สนฺนิปติตฺวา อนุตฺถุนิํสุ “ปาปกา วต โภ ธมฺมา สตฺเตสุ ปาตุภูตา, ยตฺร หิ}

\vspace{\tipitakaparskip}
\csromancenter{อคฺคญฺญสุตฺต นาม อทินฺนาทานํ ปญฺญายิสฺสติ, ครหา ปญฺญายิสฺสติ, มุสาวาโท ปญฺญายิสฺสติ, ทณฺฑาทานํ ปญฺญายิสฺสติ.\csromanspacer{} ยํนูน มยํ เอกํ สตฺตํ สมฺมนฺเนยฺยาม, โย โน สมฺมา ขียิตพฺพํ ขีเยยฺย, สมฺมา ครหิตพฺพํ ครเหยฺย, สมฺมา ปพฺพาเชตพฺพํ ปพฺพาเชยฺย.\csromanspacer{} มยํ ปนสฺส สาลีนํ ภาคํ อนุปฺปทสฺสามา”ติ.}
\prose{\csromanfolio{77}อถ โข เต วาเสฏฺฐ สตฺตา โย เนสํ สตฺโต อภิรูปตโร จ ทสฺสนียตโร จ ปาสาทิกตโร จ มเหสกฺขตโร จ.\csromanspacer{} ตํ สตฺตํ อุปสงฺกมิตฺวา เอตทโวจุํ “เอหิ โภ สตฺต สมฺมา ขียิตพฺพํ ขีย, สมฺมา ครหิตพฺพํ ครห, สมฺมา ปพฺพาเชตพฺพํ ปพฺพาเชหิ.\csromanspacer{} มยํ ปน โว สาลีนํ ภาคํ อนุปฺปทสฺสามา”ติ. “เอวํ โภ”ติ โข วาเสฏฺฐ โส สตฺโต เตสํ สตฺตานํ ปฏิสฺสุณิตฺวา สมฺมา ขียิตพฺพํ ขียิ, สมฺมา ครหิตพฺพํ ครหิ, สมฺมา ปพฺพาเชตพฺพํ ปพฺพาเชสิ.\csromanspacer{} เต ปนสฺส สาลีนํ ภาคํ อนุปฺปทํสุ.}

\proseitem{131}{มหาชนสมฺมโตติ โข วาเสฏฺฐ “มหาสมฺมโต มหาสมฺมโต”เตฺวว ปฐมํ อกฺขรํ อุปนิพฺพตฺตํ.\csromanspacer{} เขตฺตานํ อธิปตีติ โข วาเสฏฺฐ “ขตฺติโย ขตฺติโย”เตฺวว ทุติยํ อกฺขรํ อุปนิพฺพตฺตํ.\csromanspacer{} ธมฺเมน ปเร รญฺเชตีติ โข วาเสฏฺฐ “ราชา ราชา”เตฺวว ตติยํ อกฺขรํ อุปนิพฺพตฺตํ.\csromanspacer{} อิติ โข วาเสฏฺฐ เอวเมตสฺส ขตฺติยมณฺฑลสฺส โปราเณน อคฺคญฺเญน อกฺขเรน อภินิพฺพตฺติ อโหสิ เตสํเยว สตฺตานํ, อนญฺเญสํ.\csromanspacer{} สทิสานํเยว, โน อสทิสานํ.\csromanspacer{} ธมฺเมเนว, โน อธมฺเมน.\csromanspacer{} ธมฺโม หิ วาเสฏฺฐ เสฏฺโฐ ชเนตสฺมิํ ทิฏฺเฐ เจว ธมฺเม อภิสมฺปรายญฺจ.}

\csromanmatikaanchor{mtk.43}
\csromantocmarkref{section}{พฺราหฺมณมณฺฑล}{mtk.43}
\bookmark[dest={mtk.43},level=1]{พฺราหฺมณมณฺฑล}
\csromanheader{1}{\textbf{พฺราหฺมณมณฺฑล}}
\proseitem{132}{อถ โข เตสํ วาเสฏฺฐ สตฺตานํเยว\footnote{เตสํ เยว โข วาเสฏฺฐ สตฺตานํ (สี, อิ)} เอกจฺจานํ เอตทโหสิ “ปาปกา วต โภ ธมฺมา สตฺเตสุ ปาตุภูตา, ยตฺร หิ นาม อทินฺนาทานํ ปญฺญายิสฺสติ, ครหา ปญฺญายิสฺสติ มุสาวาโท ปญฺญายิสฺสติ, ทณฺฑาทานํ ปญฺญายิสฺสติ, ปพฺพาชนํ ปญฺญายิสฺสติ.\csromanspacer{} ยํนูน มยํ ปาปเก อกุสเล ธมฺเม วาเหยฺยามา”ติ.\csromanspacer{} เต ปาปเก อกุสเล ธมฺเม วาเหสุํ.\csromanspacer{} ปาปเก อกุสเล ธมฺเม วาเหนฺตีติ โข วาเสฏฺฐ “พฺราหฺมณา พฺราหฺมณา”เตฺวว ปฐมํ อกฺขรํ อุปนิพฺพตฺตํ.\csromanspacer{} เต อรญฺญายตเน \csromanfolio{78}ปณฺณกุฏิโย กริตฺวา ปณฺณกุฏีสุ ฌายนฺติ วีตงฺคารา วีตธูมา ปนฺนมุสลา สายํ สายมาสาย ปาโต ปาตราสาย คามนิคมราชธานิโย โอสรนฺติ ฆาสเมสมานา\footnote{ฆาสเมสนา (สี, สฺยา, อิ)}.\csromanspacer{} เต ฆาสํ ปฏิลภิตฺวา ปุนเทว อรญฺญายตเน ปณฺณกุฏีสุ ฌายนฺติ.\csromanspacer{} ตเมนํ มนุสฺสา ทิสฺวา เอวมาหํสุ “อิเม โข โภ สตฺตา อรญฺญายตเน ปณฺณกุฏิโย กริตฺวา ปณฺณกุฏีสุ ฌายนฺติ, วีตงฺคารา วีตธูมา ปนฺนมุสลา สายํ สายมาสาย ปาโต ปาตราสาย คามนิคมราชธานิโย โอสรนฺติ ฆาสเมสมานา.\csromanspacer{} เต ฆาสํ ปฏิลภิตฺวา ปุนเทว อรญฺญายตเน ปณฺณกุฏีสุ ฌายนฺตี”ติ, ฌายนฺตีติ โข\footnote{ปณฺณกุฏีสุ ฌายนฺติ ฌายนฺตีติ โข (สี, อิ), ปณฺณกุฏีสุ ฌายนฺตีติ โข (ก)} วาเสฏฺฐ “ฌายกา ฌายกา”เตฺวว ทุติยํ อกฺขรํ อุปนิพฺพตฺตํ.\csromanspacer{} เตสํเยว โข วาเสฏฺฐ สตฺตานํ เอกจฺเจ สตฺตา อรญฺญายตเน ปณฺณกุฏีสุ ตํ ฌานํ อนภิสมฺภุณมานา\footnote{อนภิสํภูนมานา (กตฺถจิ)} คามสามนฺตํ นิคมสามนฺตํ โอสริตฺวา คนฺเถ กโรนฺตา อจฺฉนฺติ.\csromanspacer{} ตเมนํ มนุสฺสา ทิสฺวา เอวมาหํสุ “อิเม โข โภ สตฺตา อรญฺญายตเน ปณฺณกุฏีสุ ตํ ฌานํ อนภิสมฺภุณมานา คามสามนฺตํ นิคมสามนฺตํ โอสริตฺวา คนฺเถ กโรนฺตา อจฺฉนฺติ, น ทานิเม ฌายนฺตี”ติ.\csromanspacer{} น ทานิเม\footnote{น ทานิเม ฌายนฺติ น ทานิเม (สี, อิ, ก)} ฌายนฺตีติ โข วาเสฏฺฐ “อชฺฌายกา อชฺฌายกา”เตฺวว ตติยํ อกฺขรํ อุปนิพฺพตฺตํ.\csromanspacer{} หีนสมฺมตํ โข ปน วาเสฏฺฐ เตน สมเยน โหติ, ตเทตรหิ เสฏฺฐสมฺมตํ.\csromanspacer{} อิติ โข วาเสฏฺฐ เอวเมตสฺส พฺราหฺมณมณฺฑลสฺส โปราเณน อคฺคญฺเญน อกฺขเรน อภินิพฺพตฺติ อโหสิ เตสํเยว สตฺตานํ, อนญฺเญสํ.\csromanspacer{} สทิสานํเยว, โน อสทิสานํ.\csromanspacer{} ธมฺเมเนว, โน อธมฺเมน.\csromanspacer{} ธมฺโม หิ วาเสฏฺฐ เสฏฺโฐ ชเนตสฺมิํ ทิฏฺเฐ เจว ธมฺเม อภิสมฺปรายญฺจ.}

\csromanmatikaanchor{mtk.44}
\csromantocmarkref{section}{เวสฺสมณฺฑล}{mtk.44}
\bookmark[dest={mtk.44},level=1]{เวสฺสมณฺฑล}
\csromanheader{1}{\textbf{เวสฺสมณฺฑล}}
\proseitem{133}{เตสํเยว โข วาเสฏฺฐ สตฺตานํ เอกจฺเจ สกฺกา เมถุนํ ธมฺมํ สมาทาย วิสุกมฺมนฺเต\footnote{วิสฺสุตกมฺมนฺเต (สี, อิ), วิสฺสุกมฺมนฺเต (ก-สี), วิสุํ กมฺมนฺเต (สฺยา, ก)} ปโยเชสุํ.\csromanspacer{} เมถุนํ ธมฺมํ สมาทาย วิสุกมฺมนฺเต ปโยเชนฺตีติ โข วาเสฏฺฐ “เวสฺสา เวสฺสา”เตฺวว อกฺขรํ อุปนิพฺพตฺตํ.\csromanspacer{} อิติ โข วาเสฏฺฐ เอวเมตสฺส เวสฺสมณฺฑลสฺส โปราเณน อคฺคญฺเญน อกฺขเรน อภินิพฺพตฺติ อโหสิ เตสํเยว สตฺตานํ, อนญฺเญสํ.}

\vspace{\tipitakaparskip}
\csromancenter{อคฺคญฺญสุตฺต สทิสานํเยว, โน อสทิสานํ.\csromanspacer{} ธมฺเมเนว, โน อธมฺเมน.\csromanspacer{} ธมฺโม หิ วาเสฏฺฐ เสฏฺโฐ ชเนตสฺมิํ ทิฏฺเฐ เจว ธมฺเม อภิสมฺปรายญฺจ.}
\csromanmatikaanchor{mtk.45}
\csromantocmarkref{section}{สุทฺทมณฺฑล}{mtk.45}
\bookmark[dest={mtk.45},level=1]{สุทฺทมณฺฑล}
\csromanheader{1}{\textbf{สุทฺทมณฺฑล}}
\proseitem{134}{\csromanfolio{79}เตสํเยว โข วาเสฏฺฐ สตฺตานํ เย เต สตฺตา อวเสสา, เต ลุทฺทาจารา ขุทฺทาจารา อเหสุํ.\csromanspacer{} ลุทฺทาจารา ขุทฺทาจาราติ โข วาเสฏฺฐ “สุทฺทา สุทฺทา”เตฺวว อกฺขรํ อุปนิพฺพตฺตํ.\csromanspacer{} อิติ โข วาเสฏฺฐ เอวเมตสฺส สุทฺทมณฺฑลสฺส โปราเณน อคฺคญฺเญน อกฺขเรน อภินิพฺพตฺติ อโหสิ เตสํเยว สตฺตานํ, อนญฺเญสํ.\csromanspacer{} สทิสานํเยว, โน อสทิสานํ.\csromanspacer{} ธมฺเมเนว, โน อธมฺเมน.\csromanspacer{} ธมฺโม หิ วาเสฏฺฐ เสฏฺโฐ ชเนตสฺมิํ ทิฏฺเฐ เจว ธมฺเม อภิสมฺปรายญฺจ.}

\proseitem{135}{อหุ โข โส วาเสฏฺฐ สมโย, ยํ ขตฺติโยปิ สกํ ธมฺมํ ครหมาโน อคารสฺมา อนคาริยํ ปพฺพชติ “สมโณ ภวิสฺสามี”ติ.\csromanspacer{} พฺราหฺมโณปิ โข วาเสฏฺฐ ฯเปฯ.\csromanspacer{} เวสฺโสปิ โข วาเสฏฺฐ ฯเปฯ.\csromanspacer{} สุทฺโทปิ โข วาเสฏฺฐ สกํ ธมฺมํ ครหมาโน อคารสฺมา อนคาริยํ ปพฺพชติ “สมโณ ภวิสฺสามี”ติ.\csromanspacer{} อิเมหิ โข วาเสฏฺฐ จตูหิ มณฺฑเลหิ สมณมณฺฑลสฺส อภินิพฺพตฺติ อโหสิ เตสํเยว สตฺตานํ, อนญฺเญสํ.\csromanspacer{} สทิสานํเยว, โน อสทิสานํ.\csromanspacer{} ธมฺเมเนว, โน อธมฺเมน.\csromanspacer{} ธมฺโม หิ วาเสฏฺฐ เสฏฺโฐ ชเนตสฺมิํ ทิฏฺเฐ เจว ธมฺเม อภิสมฺปรายญฺจ.}

\csromanmatikaanchor{mtk.46}
\csromantocmarkref{section}{ทุจฺจริตาทิกถา}{mtk.46}
\bookmark[dest={mtk.46},level=1]{ทุจฺจริตาทิกถา}
\csromanheader{1}{\textbf{ทุจฺจริตาทิกถา}}
\proseitem{136}{ขตฺติโยปิ โข วาเสฏฺฐ กาเยน ทุจฺจริตํ จริตฺวา วาจาย ทุจฺจริตํ จริตฺวา มนสา ทิจฺจริตํ จริตฺวา มิจฺฉาทิฏฺฐิโก มิจฺฉาทิฏฺฐิกมฺมสมาทาโน\footnote{อิทํ ปทํ สี-อิ-โปตฺถเกสุ นตฺถิ,} มิจฺฉาทิฏฺฐิกมฺมสมาทานเหตุ กายสฺส เภทา ปรํ มรณา อปายํ ทุคฺคติํ วินิปาตํ นิรยํ อุปปชฺชติ.\csromanspacer{} พฺราหฺมโณปิ โข วาเสฏฺฐ ฯเปฯ.\csromanspacer{} เวสฺโสปิ โข วาเสฏฺฐ.\csromanspacer{} สุทฺโทปิ โข วาเสฏฺฐ.\csromanspacer{} สมโณปิ โข วาเสฏฺฐ กาเยน ทุจฺจริตํ จริตฺวา วาจาย ทุจฺจริตํ จริตฺวา มนสา ทุจฺจริตํ จริตฺวา มิจฺฉาทิฏฺฐิโก มิจฺฉาทิฏฺฐิกมฺมสมาทาโน มิจฺฉาทิฏฺฐิกมฺมสมาทานเหตุ กายสฺส เภทา ปรํ มรณา อปายํ ทุคฺคติํ วินิปาตํ นิรยํ อุปปชฺชติ.}

\prose{\csromanfolio{80}ขตฺติโยปิ โข วาเสฏฺฐ กาเยน สุจริตํ จริตฺวา วาจาย สุจริตํ จริตฺวา มนสา สุจริตํ จริตฺวา สมฺมาทิฏฺฐิโก สมฺมาทิฏฺฐิกมฺมสมาทาโน\footnote{อิทํ ปทํ สี-อิ-โปตฺถเกสุ นตฺถิ,} สมฺมาทิฏฺฐิกมฺมสมาทานเหตุ กายสฺส เภทา ปรํ มรณา สุคติํ สคฺคํ โลกํ อุปปชฺชติ.\csromanspacer{} พฺราหฺมโณปิ โข วาเสฏฺฐ ฯเปฯ.\csromanspacer{} เวสฺโสปิ โข วาเสฏฺฐ.\csromanspacer{} สุทฺโทปิ โข วาเสฏฺฐ.\csromanspacer{} สมโณปิ โข วาเสฏฺฐ กาเยน สุจริตํ จริตฺวา วาจาย สุจริตํ จริตฺวา มนสา สุจริตํ จริตฺวา สมฺมาทิฏฺฐิโก สมฺมาทิฏฺฐิกมฺมสมาทาโน สมฺมาทิฏฺฐิกมฺมสมาทานเหตุ กายสฺส เภทา ปรํ มรณา สุคติํ สคฺคํ โลกํ อุปปชฺชติ.}

\proseitem{137}{ขตฺติโยปิ โข วาเสฏฺฐ กาเยน ทฺวยการี, วาจาย ทฺวยการี, มนสา ทฺวยการี, วิมิสฺสทิฏฺฐิโก วิมิสฺสทิฏฺฐิกมฺมสมาทาโน วิมิสฺสทิฏฺฐิกมฺมสมาทานเหตุ\footnote{วิมิสฺสทิฏฺฐิโก วิมิสฺสกมฺมสมาทาโน วิมิสฺสกมฺมสมาทานเหตุ (สฺยา), วีติมิสฺสทิฏฺฐิโก วีติมิสฺสทิฏฺฐิกมฺมสมาทานเหตุ (สี, อิ)} กายสฺส เภทา ปรํ มรณา สุขทุกฺขปฺปฏิสํเวที โหติ.\csromanspacer{} พฺราหฺมโณปิ โข วาเสฏฺฐ ฯเปฯ.\csromanspacer{} เวสฺโสปิ โข วาเสฏฺฐ.\csromanspacer{} สุทฺโทปิ โข วาเสฏฺฐ.\csromanspacer{} สมโณปิ โข วาเสฏฺฐ กาเยน ทฺวยการี, วาจาย ทฺวยการี, มนสา ทฺวยการี, วิมิสฺสทิฏฺฐิโก วิมิสฺสทิฏฺฐิกมฺมสมาทาโน วิมิสฺสทิฏฺฐิกมฺมสมาทานเหตุ กายสฺส เภทา ปรํ มรณา สุขทุกฺขปฺปฏิสํเวที โหติ.}

\csromanmatikaanchor{mtk.47}
\csromantocmarkref{section}{โพธิปกฺขิยภาวนา}{mtk.47}
\bookmark[dest={mtk.47},level=1]{โพธิปกฺขิยภาวนา}
\csromanheader{1}{\textbf{โพธิปกฺขิยภาวนา}}
\proseitem{138}{ขตฺติโยปิ โข วาเสฏฺฐ กาเยน สํวุโต วาจาย สํวุโต มนสา สํวุโต สตฺตนฺนํ โพธิปกฺขิยานํ ธมฺมานํ ภาวนมนฺวาย ทิฏฺเฐว ธมฺเม ปรินิพฺพายติ\footnote{ปรินิพฺพาติ (ก)}.\csromanspacer{} พฺราหฺมโณปิ โข วาเสฏฺฐ ฯเปฯ.\csromanspacer{} เวสฺโสปิ โข วาเสฏฺฐ.\csromanspacer{} สุทฺโทปิ โข วาเสฏฺฐ.\csromanspacer{} สมโณปิ โข วาเสฏฺฐ กาเยน สํวุโต วาจาย สํวุโต มนสา สํวุโต สตฺตนฺนํ โพธิปกฺขิยานํ ธมฺมานํ ภาวนมนฺวาย ทิฏฺเฐว ธมฺเม ปรินิพฺพายติ.}

\proseitem{139}{อิเมสํ หิ วาเสฏฺฐ จตุนฺนํ วณฺณานํ โย โหติ ภิกฺขุ อรหํ ขีณาสโว วุสิตวา กตกรณีโย โอหิตภาโร อนุปฺปตฺตสทตฺโถ}

\vspace{\tipitakaparskip}
\csromancenter{อคฺคญฺญสุตฺต ปริกฺขีณภวสํโยชโน สมฺมทญฺญา วิมุตฺโต.\csromanspacer{} โส เนสํ อคฺคมกฺขายติ ธมฺเมเนว, โน อธมฺเมน.\csromanspacer{} ธมฺโม หิ วาเสฏฺฐ เสฏฺโฐ ชเนตสฺมิํ ทิฏฺเฐ เจว ธมฺเม อภิสมฺปรายญฺจ.}
\csromanheader{2}{140. พฺรหฺมุนาเปสา วาเสฏฺฐ สนงฺกุมาเรน คาถา ภาสิตา\csromandash{}}
\csromangathameasure{“ขตฺติโย เสฏฺโฐ ชเนตสฺมิํ, เย โคตฺตปฏิสาริโน. \\ วิชฺชาจรณสมฺปนฺโน, โส เสฏฺโฐ เทวมานุเส”ติ.}
\csromangathagroup{\csromangathastanza{\csromanfolio{81}“ขตฺติโย เสฏฺโฐ ชเนตสฺมิํ, เย โคตฺตปฏิสาริโน. \\ วิชฺชาจรณสมฺปนฺโน, โส เสฏฺโฐ เทวมานุเส”ติ.}}

\prose{สา โข ปเนสา วาเสฏฺฐ คาถา พฺรหฺมุนา สนงฺกุมาเรน สุคีตา, โน ทุคฺคีตา.\csromanspacer{} สุภาสิตา, โน ทุพฺภาสิตา.\csromanspacer{} อตฺถสํหิตา, โน อนตฺถสํหิตา.\csromanspacer{} อนุมตา มยา.\csromanspacer{} อหมฺปิ วาเสฏฺฐ เอวํ วทามิ\csromandash{}}

\csromangathameasure{“ขตฺติโย เสฏฺโฐ ชเนตสฺมิํ, เย โคตฺตปฏิสาริโน. \\ วิชฺชาจรณสมฺปนฺโน, โส เสฏฺโฐ เทวมานุเส”ติ.}
\csromangathagroup{\csromangathastanza{“ขตฺติโย เสฏฺโฐ ชเนตสฺมิํ, เย โคตฺตปฏิสาริโน. \\ วิชฺชาจรณสมฺปนฺโน, โส เสฏฺโฐ เทวมานุเส”ติ.}}

\prose{อิทมโวจ ภควา.\csromanspacer{} อตฺตมนา วาเสฏฺฐภารทฺวาชา ภควโต ภาสิตํ อภินนฺทุนฺติ.}

\nitthitam{\textbf{อคฺคญฺญสุตฺตํ นิฏฺฐิตํ จตุตฺถํ.}}
\csromanmatikaanchor{mtk.48}
\csromantocmarknopage{chapter}{5. สมฺปสาทนียสุตฺต}{mtk.48}
\bookmark[dest={mtk.48},level=0]{5. สมฺปสาทนียสุตฺต}
\chapterheadpageread{5. สมฺปสาทนียสุตฺต}
\csromanmatikaanchor{mtk.49}
\csromantocmarkref{section}{สาริปุตฺตสีหนาท}{mtk.49}
\bookmark[dest={mtk.49},level=1]{สาริปุตฺตสีหนาท}
\csromanheader{1}{\textbf{สาริปุตฺตสีหนาท}}
\proseitem{141}{\csromanfolio{82}เอวํ เม สุตํ\csromandash{}เอกํ สมยํ ภควา นาฬนฺทายํ วิหรติ ปาวาริกมฺพวเน.\csromanspacer{} อถ โข อายสฺมา สาริปุตฺโต เยน ภควา เตนุปสงฺกมิ, อุปสงฺกมิตฺวา ภควนฺตํ อภิวาเทตฺวา เอกมนฺตํ นิสีทิ, เอกมนฺตํ นิสินฺโน โข อายสฺมา สาริปุตฺโต ภควนฺตํ เอตทโวจ “เอวํปสนฺโน อหํ ภนฺเต ภควติ, น จาหุ น จ ภวิสฺสติ น เจตรหิ วิชฺชติ อญฺโญ สมโณ วา พฺราหฺมโณ วา ภควตา ภิยฺโยภิญฺญตโร ยทิทํ สมฺโพธิยนฺ”ติ.}

\proseitem{142}{อุฬารา โข เต อยํ สาริปุตฺต อาสภี วาจา ภาสิตา, เอกํโส คหิโต, สีหนาโท นทิโต, เอวํปสนฺโน อหํ ภนฺเต ภควติ, น จาหุ น จ ภวิสฺสติ น เจตรหิ วิชฺชติ อญฺโญ สมโณ วา พฺราหฺมโณ วา ภควตา ภิยฺโยภิญฺญตโร ยทิทํ สมฺโพธิยนฺติ.\csromanspacer{} กิํ เต\footnote{กิํ นุ (สี, อิ), กิํ นุ โข เต (สฺยา)} สาริปุตฺต เย เต อเหสุํ อตีตมทฺธานํ อรหนฺโต สมฺมาสมฺพุทฺธา, สพฺเพ เต ภควนฺโต เจตสา เจโต ปริจฺจ วิทิตา “เอวํสีลา เต ภควนฺโต อเหสุํ” อิติปิ, “เอวํธมฺมา เต ภควนฺโต อเหสุํ” อิติปิ, “เอวํปญฺญา เต ภควนฺโต อเหสุํ” อิติปิ, “เอวํวิหารี เต ภควนฺโต อเหสุํ” อิติปิ, “เอวํวิมุตฺตา เต ภควนฺโต อเหสุํ” อิติปีติ.\csromanspacer{} โน เหตํ ภนฺเต.}

\prose{กิํ ปน เต\footnote{กิํ ปน (สี, อิ)} สาริปุตฺต เย เต ภวิสฺสนฺติ อนาคตมทฺธานํ อรหนฺโต สมฺมาสมฺพุทฺธา, สพฺเพ เต ภควนฺโต เจตสา เจโต ปริจฺจ วิทิตา, “เอวํสีลา เต ภควนฺโต ภวิสฺสนฺติ” อิติปิ, “เอวํธมฺมา.\csromanspacer{} เอวํปญฺญา.\csromanspacer{} เอวํวิหารี.\csromanspacer{} เอวํวิมุตฺตา เต ภควนฺโต ภวิสฺสนฺติ” อิติปีติ.\csromanspacer{} โน เหตํ ภนฺเต.}

\prose{กิํ ปน เต2 สาริปุตฺต อหํ เอตรหิ อรหํ สมฺมาสมฺพุทฺโธ เจตสา เจโต ปริจฺจ วิทิโต “เอวํสีโล ภควา” อิติปิ, “เอวํธมฺโม.}

\vspace{\tipitakaparskip}
\csromancenter{สมฺปสาทนียสุตฺต เอวํปญฺโญ.\csromanspacer{} เอวํวิหารี.\csromanspacer{} เอวํวิมุตฺโต ภควา” อิติปีติ.\csromanspacer{} โน เหตํ ภนฺเต.}
\prose{\csromanfolio{83}เอตฺถ จ หิ เต สาริปุตฺต อตีตานาคตปจฺจุปฺปนฺเนสุ อรหนฺเตสุ สมฺมาสมฺพุทฺเธสุ เจโตปริยญาณํ นตฺถิ.\csromanspacer{} อถ กิญฺจรหิ เต อยํ สาริปุตฺต อุฬารา อาสภี วาจา ภาสิตา, เอกํโส คหิโต, สีหนาโท นทิโต “เอวํปสนฺโน อหํ ภนฺเต ภควติ, น จาหุ น จ ภวิสฺสติ น เจตรหิ วิชฺชติ อญฺโญ สมโณ วา พฺราหฺมโณ วา ภควตา ภิยฺโยภิญฺญตโร ยทิทํ สมฺโพธิยนฺ”ติ.}

\proseitem{143}{น โข เม\footnote{น โข ปเนตํ (สฺยา, ก)} ภนฺเต อตีตานาคตปจฺจุปฺปนฺเนสุ อรหนฺเตสุ สมฺมาสมฺพุทฺเธสุ เจโตปริยญาณํ อตฺถิ.\csromanspacer{} อปิ จ เม\footnote{เม ภนฺเต (สี, อิ, ก)} ธมฺมนฺวโย วิทิโต.\csromanspacer{} เสยฺยถาปิ ภนฺเต รญฺโญ ปจฺจนฺติมํ นครํ ทฬฺหุทฺธาปํ\footnote{ทฬฺหุทฺทาปํ (สี, อิ, ก)} ทฬฺหปาการโตรณํ เอกทฺวารํ.\csromanspacer{} ตตฺรสฺส โทวาริโก ปณฺฑิโต พฺยาตฺโต เมธาวี อญฺญาตานํ นิวาเรตา, ญาตานํ ปเวเสตา.\csromanspacer{} โส ตสฺส นครสฺส สมนฺตา อนุปริยายปถํ อนุกฺกมมาโน น ปสฺเสยฺย ปาการสนฺธิํ วา ปาการวิวรํ วา อนฺตมโส พิฬารนิกฺขมนมตฺตมฺปิ.\csromanspacer{} ตสฺส เอวมสฺส “เย โข เกจิ โอฬาริกา ปาณา อิมํ นครํ ปวิสนฺติ วา นิกฺขมนฺติ วา, สพฺเพ เต อิมินาว ทฺวาเรน ปวิสนฺติ วา นิกฺขมนฺติ วา”ติ.\csromanspacer{} เอวเมว โข เม ภนฺเต ธมฺมนฺวโย วิทิโต.\csromanspacer{} เย เต ภนฺเต อเหสุํ อตีตมทฺธานํ อรหนฺโต สมฺมาสมฺพุทฺธา, สพฺเพ เต ภควนฺโต ปญฺจ นีวรเณ ปหาย เจตโส อุปกฺกิเลเส ปญฺญาย ทุพฺพลีกรเณ จตูสุ สติปฏฺฐาเนสุ สุปฺปติฏฺฐิตจิตฺตา สตฺต สมฺโพชฺฌงฺเค ยถาภูตํ ภาเวตฺวา อนุตฺตรํ สมฺมาสมฺโพธิํ อภิสมฺพุชฺฌิํสุ.\csromanspacer{} เยปิ เต ภนฺเต ภวิสฺสนฺติ อนาคตมทฺธานํ อรหนฺโต สมฺมาสมฺพุทฺธา, สพฺเพ เต ภควนฺโต ปญฺจ นีวรเณ ปหาย เจตโส อุปกฺกิเลเส ปญฺญาย ทุพฺพลีกรเณ จตูสุ สติปฏฺฐาเนสุ สุปฺปติฏฺฐิตจิตฺตา สตฺต สมฺโพชฺฌงฺเค ยถาภูตํ ภาเวตฺวา อนุตฺตรํ สมฺมาสมฺโพธิํ อภิสมฺพุชฺฌิสฺสนฺติ.\csromanspacer{} ภควาปิ ภนฺเต เอตรหิ อรหํ สมฺมาสมฺพุทฺโธ ปญฺจ นีวรเณ ปหาย เจตโส อุปกฺกิเลเส ปญฺญาย ทุพฺพลีกรเณ จตูสุ สติปฏฺฐาเนสุ \csromanfolio{84}สุปฺปติฏฺฐิตจิตฺโต สตฺต สมฺโพชฺฌงฺเค ยถาภูตํ ภาเวตฺวา อนุตฺตรํ สมฺมาสมฺโพธิํ อภิสมฺพุทฺโธ.}

\proseitem{144}{อิธาหํ ภนฺเต เยน ภควา เตนุปสงฺกมิํ ธมฺมสฺสวนาย.\csromanspacer{} ตสฺส เม ภนฺเต ภควา ธมฺมํ เทเสติ อุตฺตรุตฺตรํ ปณีตปณีตํ กณฺหสุกฺกสปฺปฏิภาคํ.\csromanspacer{} ยถา ยถา เม ภนฺเต ภควา ธมฺมํ เทเสสิ อุตฺตรุตฺตรํ ปณีตปณีตํ กณฺหสุกฺกสปฺปฏิภาคํ, ตถา ตถาหํ ตสฺมิํ ธมฺเม อภิญฺญา อิเธกจฺจํ ธมฺมํ ธมฺเมสุ นิฏฺฐมคมํ สตฺถริ ปสีทิํ “สมฺมาสมฺพุทฺโธ ภควา, สฺวากฺขาโต ภควตา ธมฺโม, สุปฺปฏิปนฺโน สาวกสํโฆ”ติ.}

\csromanmatikaanchor{mtk.50}
\csromantocmarkref{section}{กุสลธมฺมเทสนา}{mtk.50}
\bookmark[dest={mtk.50},level=1]{กุสลธมฺมเทสนา}
\csromanheader{1}{\textbf{กุสลธมฺมเทสนา}}
\proseitem{145}{อปรํ ปน ภนฺเต เอตทานุตฺตริยํ, ยถา ภควา ธมฺมํ เทเสติ กุสเลสุ ธมฺเมสุ.\csromanspacer{} ตตฺริเม กุสลา ธมฺมา.\csromanspacer{} เสยฺยถิทํ, จตฺตาโร สติปฏฺฐานา จตฺตาโร สมฺมปฺปธานา จตฺตาโร อิทฺธิปาทา ปญฺจินฺทฺริยานิ ปญฺจ พลานิ สตฺต โพชฺฌงฺคา อริโย อฏฺฐงฺคิโก มคฺโค.\csromanspacer{} อิธ ภนฺเต ภิกฺขุ อาสวานํ ขยา อนาสวํ เจโตวิมุตฺติํ ปญฺญาวิมุตฺติํ ทิฏฺเฐวธมฺเม สยํ อภิญฺญา สจฺฉิกตฺวา อุปสมฺปชฺช วิหรติ.\csromanspacer{} เอตทานุตฺตริยํ ภนฺเต กุสเลสุ ธมฺเมสุ.\csromanspacer{} ตํ ภควา อเสสมภิชานาติ, ตํ ภควโต อเสสมภิชานโต อุตฺตริ อภิญฺเญยฺยํ นตฺถิ, ยทภิชานํ อญฺโญ สมโณ วา พฺราหฺมโณ วา ภควตา ภิยฺโยภิญฺญตโร อสฺส ยทิทํ กุสเลสุ ธมฺเมสุ.}

\csromanmatikaanchor{mtk.51}
\csromantocmarkref{section}{อายตนปณฺณตฺติเทสนา}{mtk.51}
\bookmark[dest={mtk.51},level=1]{อายตนปณฺณตฺติเทสนา}
\csromanheader{1}{\textbf{อายตนปณฺณตฺติเทสนา}}
\proseitem{146}{อปรํ ปน ภนฺเต เอตทานุตฺตริยํ, ยถา ภควา ธมฺมํ เทเสติ อายตนปณฺณตฺตีสุ.\csromanspacer{} ฉยิมานิ ภนฺเต อชฺฌตฺติกพาหิรานิ อายตนานิ.\csromanspacer{} จกฺขุญฺเจว รูปา\footnote{รูปานิ (ก)} จ, โสตญฺเจว สทฺทา จ, ฆานญฺเจว คนฺธา จ, ชิวฺหา เจว รสา จ, กาโย เจว โผฏฺฐพฺพา จ, มโน เจว ธมฺมา จ.\csromanspacer{} เอตทานุตฺตริยํ ภนฺเต อายตนปณฺณตฺตีสุ.\csromanspacer{} ตํ ภควา อเสสมภิชานาติ, ตํ ภควโต อเสสมภิชานโต อุตฺตริ อภิญฺเญยฺยํ นตฺถิ, ยทภิชานํ อญฺโญ สมโณ วา พฺราหฺมโณ วา ภควตา ภิยฺโยภิญฺญตโร อสฺส ยทิทํ อายตนปณฺณตฺตีสุ.}

\csromanmatikaanchor{mtk.52}
\csromantocmarkref{section}{คพฺภาวกฺกนฺติเทสนา}{mtk.52}
\bookmark[dest={mtk.52},level=1]{คพฺภาวกฺกนฺติเทสนา}
\csromanheader{1}{5. สมฺปสาทนียสุตฺต \textbf{คพฺภาวกฺกนฺติเทสนา}}
\proseitem{147}{\csromanfolio{85}อปรํ ปน ภนฺเต เอตทานุตฺตริยํ, ยถา ภควา ธมฺมํ เทเสติ คพฺภาวกฺกนฺตีสุ.\csromanspacer{} จตสฺโส อิมา ภนฺเต คพฺภาวกฺกนฺติโย.\csromanspacer{} อิธ ภนฺเต เอกจฺโจ อสมฺปชาโน มาตุกุจฺฉิํ โอกฺกมติ, อสมฺปชาโน มาตุกุจฺฉิสฺมิํ ฐาติ, อสมฺปชาโน มาตุกุจฺฉิมฺหา นิกฺขมติ.\csromanspacer{} อยํ ปฐมา คพฺภาวกฺกนฺติ.}

\prose{ปุน จปรํ ภนฺเต อิเธกจฺโจ สมฺปชาโน มาตุกุจฺฉิํ โอกฺกมติ, อสมฺปชาโน มาตุกุจฺฉิสฺมิํ ฐาติ, อสมฺปชาโน มาตุกุจฺฉิมฺหา นิกฺขมติ.\csromanspacer{} อยํ ทุติยา คพฺภาวกฺกนฺติ.}

\prose{ปุน จปรํ ภนฺเต อิเธกจฺโจ สมฺปชาโน มาตุกุจฺฉิํ โอกฺกมติ, สมฺปชาโน มาตุกุจฺฉิสฺมิํ ฐาติ, อสมฺปชาโน มาตุกุจฺฉิมฺหา นิกฺขมติ อยํ ตติยา คพฺภาวกฺกนฺติ.}

\prose{ปุน จปรํ ภนฺเต อิเธกจฺโจ สมฺปชาโน มาตุกุจฺฉิํ โอกฺกมติ, สมฺปชาโน มาตุกุจฺฉิสฺมิํ ฐาติ, สมฺปชาโน มาตุกุจฺฉิมฺหา นิกฺขมติ.\csromanspacer{} อยํ จตุตฺถา คพฺภาวกฺกนฺติ.\csromanspacer{} เอตทานุตฺตริยํ ภนฺเต คพฺภาวกฺกนฺตีสุ.}

\csromanmatikaanchor{mtk.53}
\csromantocmarkref{section}{อาเทสนวิธาเทสนา}{mtk.53}
\bookmark[dest={mtk.53},level=1]{อาเทสนวิธาเทสนา}
\csromanheader{1}{\textbf{อาเทสนวิธาเทสนา}}
\proseitem{148}{อปรํ ปน ภนฺเต เอตทานุตฺตริยํ, ยถา ภควา ธมฺมํ เทเสติ อาเทสนวิธาสุ.\csromanspacer{} จตสฺโส อิมา ภนฺเต อาเทสนวิธา.\csromanspacer{} อิธ ภนฺเต เอกจฺโจ นิมิตฺเตน อาทิสติ “เอวมฺปิ เต มโน, อิตฺถํปิ เต มโน, อิติปิ เต จิตฺตนฺ”ติ.\csromanspacer{} โส พหุํ เจปิ อาทิสติ, ตเถว ตํ โหติ, โน อญฺญถา.\csromanspacer{} อยํ ปฐมา อาเทสนวิธา.}

\prose{ปุน จปรํ ภนฺเต อิเธกจฺโจ น เหว โข นิมิตฺเตน อาทิสติ.\csromanspacer{} อปิ จ โข มนุสฺสานํ วา อมนุสฺสานํ วา เทวตานํ วา สทฺทํ สุตฺวา อาทิสติ “เอวมฺปิ เต มโน, อิตฺถํปิ เต มโน, อิติปิ เต จิตฺตนฺ”ติ.\csromanspacer{} โส พหุํ เจปิ อาทิสติ, ตเถว ตํ โหติ, โน อญฺญถา.\csromanspacer{} อยํ ทุติยา อาเทสนวิธา.}

\prose{ปุน จปรํ ภนฺเต อิเธกจฺโจ น เหว โข นิมิตฺเตน อาทิสติ, นาปิ มนุสฺสานํ วา อมนุสฺสานํ วา เทวตานํ วา สทฺทํ สุตฺวา อาทิสติ.\csromanspacer{} อปิ จ โข วิตกฺกยโต วิจารยโต วิตกฺกวิปฺผารสทฺทํ สุตฺวา \csromanfolio{86}อาทิสติ “เอวมฺปิ เต มโน, อิตฺถํปิ เต มโน, อิติปิ เต จิตฺตนฺ”ติ.\csromanspacer{} โส พหุํ เจปิ อาทิสติ, ตเถว ตํ โหติ, โน อญฺญถา.\csromanspacer{} อยํ ตติยา อาเทสนวิธา.}

\prose{ปุน จปรํ ภนฺเต อิเธกจฺโจ น เหว โข นิมิตฺเตน อาทิสติ, นาปิ มนุสฺสานํ วา อมนุสฺสานํ วา เทวตานํ วา สทฺทํ สุตฺวา อาทิสติ, นาปิ วิตกฺกยโต วิจารยโต วิตกฺกวิปฺผารสทฺทํ สุตฺวา อาทิสติ.\csromanspacer{} อปิ จ โข วิตกฺกวิจารสมาธิสมาปนฺนสฺส\footnote{อวิตกฺกํ อวิจารํ สมาธิํ สมาปนฺนสฺส (สี, อิ).} เจตสา เจโต ปริจฺจ ปชานาติ “ยถา อิมสฺส โภโต มโนสงฺขารา ปณิหิตา.\csromanspacer{} ตถา อิมสฺส จิตฺตสฺส อนนฺตรา อิมํ นาม วิตกฺกํ วิตกฺเกสฺสตี”ติ.\csromanspacer{} โส พหุํ เจปิ อาทิสติ, ตเถว ตํ โหติ, โน อญฺญถาติ.\csromanspacer{} อยํ จตุตฺถา อาเทสนวิธา.\csromanspacer{} เอตทานุตฺตริยํ ภนฺเต อาเทสนวิธาสุ.}

\csromanmatikaanchor{mtk.54}
\csromantocmarkref{section}{ทสฺสนสมาปตฺติเทสนา}{mtk.54}
\bookmark[dest={mtk.54},level=1]{ทสฺสนสมาปตฺติเทสนา}
\csromanheader{1}{\textbf{ทสฺสนสมาปตฺติเทสนา}}
\proseitem{149}{อปรํ ปน ภนฺเต เอตทานุตฺตริยํ, ยถา ภควา ธมฺมํ เทเสติ ทสฺสนสมาปตฺตีสุ.\csromanspacer{} จตสฺโส อิมา ภนฺเต ทสฺสนสมาปตฺติโย.\csromanspacer{} อิธ ภนฺเต เอกจฺโจ สมโณ วา พฺราหฺมโณ วา อาตปฺปมนฺวาย ปธานมนฺวาย อนุโยคมนฺวาย อปฺปมาทมนฺวาย สมฺมามนสิการมนฺวาย ตถารูปํ เจโตสมาธิํ ผุสติ, ยถาสมาหิเต จิตฺเต อิมเมว กายํ อุทฺธํ ปาทตลา อโธ เกสมตฺถกา ตจปริยนฺตํ ปูรํ นานปฺปการสฺส อสุจิโน ปจฺจเวกฺขติ “อตฺถิ อิมสฺมิํ กาเย เกสา โลมา นขา ทนฺตา ตโจ มํสํ นฺหารุ อฏฺฐิ อฏฺฐิมิญฺชํ วกฺกํ หทยํ ยกนํ กิโลมกํ ปิหกํ ปปฺผาสํ อนฺตํ อนฺตคุณํ อุทริยํ กรีสํ ปิตฺตํ เสมฺหํ ปุพฺโพ โลหิตํ เสโท เมโท อสฺสุ วสา เขโฬ สิงฺฆานิกา ลสิกา มุตฺตนฺ”ติ.\csromanspacer{} อยํ ปฐมา ทสฺสนสมาปตฺติ.}

\prose{ปุน จปรํ ภนฺเต อิเธกจฺโจ สมโณ วา พฺราหฺมโณ วา อาตปฺปมนฺวาย ฯเปฯ ตถารูปํ เจโตสมาธิํ ผุสติ, ยถาสมาหิเต จิตฺเต อิมเมว กายํ อุทฺธํ ปาทตลา อโธ เกสมตฺถกา ตจปริยนฺตํ ปูรํ นานปฺปการสฺส อสุจิโน ปจฺจเวกฺขติ “อตฺถิ อิมสฺมิํ กาเย}

\vspace{\tipitakaparskip}
\csromancenter{สมฺปสาทนียสุตฺต เกสา โลมา ฯเปฯ ลสิกา มุตฺตนฺ”ติ.\csromanspacer{} อติกฺกมฺม จ ปุริสสฺส ฉวิมํสโลหิตํ อฏฺฐิํ ปจฺจเวกฺขติ.\csromanspacer{} อยํ ทุติยา ทสฺสนสมาปตฺติ.}
\prose{\csromanfolio{87}ปุน จปรํ ภนฺเต อิเธกจฺโจ สมโณ วา พฺราหฺมโณ วา อาตปฺปมนฺวาย ฯเปฯ ตถารูปํ เจโตสมาธิํ ผุสติ, ยถาสมาหิเต จิตฺเต อิมเมว กายํ อุทฺธํ ปาทตลา อโธ เกสมตฺถกา ตจปริยนฺตํ ปูรํ นานปฺปการสฺส อสุจิโน ปจฺจเวกฺขติ “อตฺถิ อิมสฺมิํ กาเย เกสา โลมา ฯเปฯ ลสิกา มุตฺตนฺ”ติ.\csromanspacer{} อติกฺกมฺม จ ปุริสสฺส ฉวิมํสโลหิตํ อฏฺฐิํ ปจฺจเวกฺขติ.\csromanspacer{} ปุริสสฺส จ วิญฺญาณโสตํ ปชานาติ, อุภยโต อพฺโพจฺฉินฺนํ อิธ โลเก ปติฏฺฐิตญฺจ ปรโลเก ปติฏฺฐิตญฺจ.\csromanspacer{} อยํ ตติยา ทสฺสนสมาปตฺติ.}

\prose{ปุน จปรํ ภนฺเต อิเธกจฺโจ สมโณ วา พฺราหฺมโณ วา อาตปฺปมนฺวาย ฯเปฯ ตถารูปํ เจโตสมาธิํ ผุสติ, ยถาสมาหิเต จิตฺเต อิมเมว กายํ อุทฺธํ ปาทตลา อโธ เกสมตฺถกา ตจปริยนฺต ปูรํ นานปฺปการสฺส อสุจิโน ปจฺจเวกฺขติ “อตฺถิ อิมสฺมิํ กาเย เกสา โลมา ฯเปฯ ลสิกา มุตฺตนฺ”ติ.\csromanspacer{} อติกฺกมฺม จ ปุริสสฺส ฉวิมํสโลหิตํ อฏฺฐิํ ปจฺจเวกฺขติ.\csromanspacer{} ปุริสสฺส จ วิญฺญาณโสตํ ปชานาติ, อุภยโต อพฺโพจฺฉินฺนํ อิธ โลเก อปฺปติฏฺฐิตญฺจ ปรโลเก อปฺปติฏฺฐิตญฺจ.\csromanspacer{} อยํ จตุตฺถา ทสฺสนสมาปตฺติ.\csromanspacer{} เอตทานุตฺตริยํ ภนฺเต ทสฺสนสมาปตฺตีสุ.}

\csromanmatikaanchor{mtk.55}
\csromantocmarkref{section}{ปุคฺคลปณฺณตฺติเทสนา}{mtk.55}
\bookmark[dest={mtk.55},level=1]{ปุคฺคลปณฺณตฺติเทสนา}
\csromanheader{1}{\textbf{ปุคฺคลปณฺณตฺติเทสนา}}
\proseitem{150}{อปรํ ปน ภนฺเต เอตทานุตฺตริยํ, ยถา ภควา ธมฺมํ เทเสติ ปุคฺคลปณฺณตฺตีสุ.\csromanspacer{} สตฺติเม ภนฺเต ปุคฺคลา.\csromanspacer{} อุภโตภาควิมุตฺโต ปญฺญาวิมุตฺโต กายสกฺขี ทิฏฺฐิปฺปตฺโต สทฺธาวิมุตฺโต ธมฺมานุสารี สทฺธานุสารี.\csromanspacer{} เอตทานุตฺตริยํ ภนฺเต ปุคฺคลปณฺณตฺตีสุ.}

\csromanmatikaanchor{mtk.56}
\csromantocmarkref{section}{ปธานเทสนา}{mtk.56}
\bookmark[dest={mtk.56},level=1]{ปธานเทสนา}
\csromanheader{1}{\textbf{ปธานเทสนา}}
\proseitem{151}{อปรํ ปน ภนฺเต เอตทานุตฺตริยํ, ยถา ภควา ธมฺมํ เทเสติ ปธาเนสุ.\csromanspacer{} สตฺติเม ภนฺเต สมฺโมชฺฌงฺคา.\csromanspacer{} สติสมฺโพชฺฌงฺโค ธมฺมวิจยสมฺโพชฺฌงฺโค วีริยสมฺโพชฺฌงฺโค ปีติสมฺโพชฺฌงฺโค ปสฺสทฺธิสมฺโพชฺฌงฺโค \csromanfolio{88}สมาธิสมฺโพชฺฌงฺโค อุเปกฺขาสมฺโพชฺฌงฺโค.\csromanspacer{} เอตทานุตฺตริยํ ภนฺเต ปธาเนสุ.}

\csromanmatikaanchor{mtk.57}
\csromantocmarkref{section}{ปฏิปทาเทสนา}{mtk.57}
\bookmark[dest={mtk.57},level=1]{ปฏิปทาเทสนา}
\csromanheader{1}{\textbf{ปฏิปทาเทสนา}}
\proseitem{152}{อปรํ ปน ภนฺเต เอตทานุตฺตริยํ, ยถา ภควา ธมฺมํ เทเสติ ปฏิปทาสุ.\csromanspacer{} จตสฺโส อิมา ภนฺเต ปฏิปทา.\csromanspacer{} ทุกฺขาปฏิปทา ทนฺธาภิญฺญา, ทุกฺขาปฏิปทา ขิปฺปาภิญฺญา, สุขาปฏิปทา ทนฺธาภิญฺญา, สุขาปฏิปทา ขิปฺปาภิญฺญาติ.\csromanspacer{} ตตฺร ภนฺเต ยายํ ปฏิปทา ทุกฺขา ทนฺธาภิญฺญา, อยํ ภนฺเต ปฏิปทา อุภเยเนว หีนา อกฺขายติ ทุกฺขตฺตา จ ทนฺธตฺตา จ.\csromanspacer{} ตตฺร ภนฺเต ยายํ ปฏิปทา ทุกฺขา ขิปฺปาภิญฺญา, อยํ ปน ภนฺเต ปฏิปทา ทุกฺขตฺตา หีนา อกฺขายติ.\csromanspacer{} ตตฺร ภนฺเต ยายํ ปฏิปทา สุขา ทนฺธาภิญฺญา, อยํ ปน ภนฺเต ปฏิปทา ทนฺธตฺตา หีนา อกฺขายติ.\csromanspacer{} ตตฺร ภนฺเต ยายํ ปฏิปทา สุขา ขิปฺปาภิญฺญา, อยํ ปน ภนฺเต ปฏิปทา อุภเยเนว ปณีตา อกฺขายติ สุขตฺตา จ ขิปฺปตฺตา จ.\csromanspacer{} เอตทานุตฺตริยํ ภนฺเต ปฏิปทาสุ.}

\csromanmatikaanchor{mtk.58}
\csromantocmarkref{section}{ภสฺสสมาจาราทิเทสนา}{mtk.58}
\bookmark[dest={mtk.58},level=1]{ภสฺสสมาจาราทิเทสนา}
\csromanheader{1}{\textbf{ภสฺสสมาจาราทิเทสนา}}
\proseitem{153}{อปรํ ปน ภนฺเต เอตทานุตฺตริยํ, ยถา ภควา ธมฺมํ เทเสติ ภสฺสสมาจาเร.\csromanspacer{} อิธ ภนฺเต เอกจฺโจ น เจว มุสาวาทุปสญฺหิตํ วาจํ ภาสติ น จ เวภูติยํ น จ เปสุณิยํ น จ สารมฺภชํ ชยาเปกฺโข, มนฺตา มนฺตา จ วาจํ ภาสติ นิธานวติํ กาเลน.\csromanspacer{} เอตทานุตฺตริยํ ภนฺเต ภสฺสสมาจาเร.}

\prose{อปรํ ปน ภนฺเต เอตทานุตฺตริยํ, ยถา ภควา ธมฺมํ เทเสติ ปุริสสีลสมาจาเร.\csromanspacer{} อิธ ภนฺเต เอกจฺโจ สจฺโจ จสฺส สทฺโธ จ, น จ กุหโก, น จ ลปโก, น จ เนมิตฺติโก, น จ นิปฺเปสิโก, น จ ลาเภน ลาภํ นิชิคีสนโก\footnote{นิชิคิํสนโก (สฺยา), นิชิคิํสิตา (สี, อิ)}, อินฺทฺริเยสุ คุตฺตทฺวาโร, โภชเน มตฺตญฺญู, สมการี, ชาคริยานุโยคมนุยุตฺโต, อตนฺทิโต, อารทฺธวีริโย, ฌายี, สติมา, กลฺยาณปฏิภาโน, คติมา, ธิติมา, มติมา, น จ กาเมสุ คิทฺโธ, สโต จ นิปโก จ.\csromanspacer{} เอตทานุตฺตริยํ ภนฺเต ปุริสสีลสมาจาเร.}

\csromanmatikaanchor{mtk.59}
\csromantocmarkref{section}{อนุสาสนวิธาเทสนา}{mtk.59}
\bookmark[dest={mtk.59},level=1]{อนุสาสนวิธาเทสนา}
\csromanheader{1}{5. สมฺปสาทนียสุตฺต \textbf{อนุสาสนวิธาเทสนา}}
\proseitem{154}{\csromanfolio{89}อปรํ ปน ภนฺเต เอตทานุตฺตริยํ, ยถา ภควา ธมฺมํ เทเสติ อนุสาสนวิธาสุ.\csromanspacer{} จตสฺโส อิมา ภนฺเต อนุสาสนวิธา.\csromanspacer{} ชานาติ ภนฺเต ภควา อปรํ ปุคฺคลํ ปจฺจตฺตํ โยนิโสมนสิการา “อยํ ปุคฺคโล ยถานุสิฏฺฐํ ตถา ปฏิปชฺชมาโน ติณฺณํ สํโยชนานํ ปริกฺขยา โสตาปนฺโน ภวิสฺสติ อวินิปาตธมฺโม นิยโต สมฺโพธิปรายโณ”ติ.\csromanspacer{} ชานาติ ภนฺเต ภควา ปรํ ปุคฺคลํ ปจฺจตฺตํ โยนิโสมนสิการา “อยํ ปุคฺคโล ยถานุสิฏฺฐํ ตถา ปฏิปชฺชมาโน ติณฺณํ สํโยชนานํ ปริกฺขยา ราคโทสโมหานํ ตนุตฺตา สกทาคามี ภวิสฺสติ, สกิเทว อิมํ โลกํ อาคนฺตฺวา ทุกฺขสฺสนฺตํ กริสฺสตี”ติ.\csromanspacer{} ชานาติ ภนฺเต ภควา ปรํ ปุคฺคลํ ปจฺจตฺตํ โยนิโสมนสิการา “อยํ ปุคฺคโล ยถานุสิฏฺฐํ ตถา ปฏิปชฺชมาโน ปญฺจนฺนํ โอรมฺภาคิยานํ สํโยชนานํ ปริกฺขยา โอปปาติโก ภวิสฺสติ ตตฺถ ปรินิพฺพายี อนาวตฺติธมฺโม ตสฺมา โลกา”ติ.\csromanspacer{} ชานาติ ภนฺเต ภควา ปรํ ปุคฺคลํ ปจฺจตฺตํ โยนิโสมนสิการา “อยํ ปุคฺคโล ยถานุสิฏฺฐํ ตถา ปฏิปชฺชมาโน อาสวานํ ขยา อนาสวํ เจโตวิมุตฺติํ ปญฺญาวิมุตฺติํ ทิฏฺเฐวธมฺเม สยํ อภิญฺญา สจฺฉิกตฺวา อุปสมฺปชฺช วิหริสฺสตี”ติ.\csromanspacer{} เอตทานุตฺตริยํ ภนฺเต อนุสาสนวิธาสุ.}

\csromanmatikaanchor{mtk.60}
\csromantocmarkref{section}{ปรปุคฺคลวิมุตฺติญาณเทสนา}{mtk.60}
\bookmark[dest={mtk.60},level=1]{ปรปุคฺคลวิมุตฺติญาณเทสนา}
\csromanheader{1}{\textbf{ปรปุคฺคลวิมุตฺติญาณเทสนา}}
\proseitem{155}{อปรํ ปน ภนฺเต เอตทานุตฺตริยํ, ยถา ภควา ธมฺมํ เทเสติ ปรปุคฺคลวิมุตฺติญาเณ.\csromanspacer{} ชานาติ ภนฺเต ภควา ปรํ ปุคฺคลํ ปจฺจตฺตํ โยนิโสมนสิการา “อยํ ปุคฺคโล ติณฺณํ สํโยชนานํ ปริกฺขยา โสตาปนฺโน ภวิสฺสติ อวินิปาตธมฺโม นิยโต สมฺโพธิปรายโณ”ติ, ชานาติ ภนฺเต ภควา ปรํ ปุคฺคลํ ปจฺจตฺตํ โยนิโสมนสิการา “อยํ ปุคฺคโล ติณฺณํ สํโยชนานํ ปริกฺขยา ราคโทสโมหานํ ตนุตฺตา สกทาคามี ภวิสฺสติ, สกิเทว อิมํ โลกํ อาคนฺตฺวา ทุกฺขสฺสนฺตํ กริสฺสตี”ติ.\csromanspacer{} ชานาติ ภนฺเต ภควา ปรํ ปุคฺคลํ ปจฺจตฺตํ โยนิโสมนสิการา “อยํ ปุคฺคโล ปญฺจนฺนํ โอรมฺภาคิยานํ สํโยชนานํ ปริกฺขยา โอปปาติโก ภวิสฺสติ ตตฺถ ปรินิพฺพายี อนาวตฺติธมฺโม ตสฺมา โลกา”ติ.\csromanspacer{} ชานาติ ภนฺเต ภควา ปรํ ปุคฺคลํ ปจฺจตฺตํ โยนิโสมนสิการา “อยํ ปุคฺคโล อาสวานํ ขยา อนาสวํ เจโตวิมุตฺติํ \csromanfolio{90}ปญฺญาวิมุตฺติํ ทิฏฺเฐวธมฺเม สยํ อภิญฺญา สจฺฉิกตฺวา อุปสมฺปชฺช วิหริสฺสตี”ติ.\csromanspacer{} เอตทานุตฺตริยํ ภนฺเต ปรปุคฺคลวิมุตฺติญาเณ.}

\csromanmatikaanchor{mtk.61}
\csromantocmarkref{section}{สสฺสตวาทเทสนา}{mtk.61}
\bookmark[dest={mtk.61},level=1]{สสฺสตวาทเทสนา}
\csromanheader{1}{\textbf{สสฺสตวาทเทสนา}}
\proseitem{156}{อปรํ ปน ภนฺเต เอตทานุตฺตริยํ, ยถา ภควา ธมฺมํ เทเสติ สสฺสตวาเทสุ.\csromanspacer{} ตโยเม ภนฺเต สสฺสตวาทา.\csromanspacer{} อิธ ภนฺเต เอกจฺโจ สมโณ วา พฺราหฺมโณ วา อาตปฺปมนฺวาย ฯเปฯ ตถารูปํ เจโตสมาธิํ ผุสติ, ยถาสมาหิเต จิตฺเต อเนกวิหิตํ ปุพฺเพนิวาสํ อนุสฺสรติ.\csromanspacer{} เสยฺยถิทํ, เอกํปิ ชาติํ เทฺวปิ ชาติโย ติสฺโสปิ ชาติโย จตสฺโสปิ ชาติโย ปญฺจปิ ชาติโย ทสปิ ชาติโย วีสํปิ ชาติโย ติํสํปิ ชาติโย จตฺตาลีสํปิ ชาติโย ปญฺญาสํปิ ชาติโย ชาติสตํปิ ชาติสหสฺสํปิ ชาติสตสหสฺสํปิ อเนกานิปิ ชาติสตานิ อเนกานิปิ ชาติสหสฺสานิ อเนกานิปิ ชาติสตสหสฺสานิ, “อมุตฺราสิํ เอวํนาโม เอวํโคตฺโต เอวํวณฺโณ เอวมาหาโร เอวํสุขทุกฺขปฺปฏิสํเวที เอวมายุปริยนฺโต, โส ตโต จุโต อมุตฺร อุทปาทิํ, ตตฺราปาสิํ เอวํนาโม เอวํโคตฺโต เอวํวณฺโณ เอวมาหาโร เอวํสุขทุกฺขปฺปฏิสํเวที เอวมายุปริยนฺโต, โส ตโต จุโต อิธูปปนฺโน”ติ.\csromanspacer{} อิติ สาการํ ส-อุทฺเทสํ อเนกวิหิตํ ปุพฺเพนิวาสํ อนุสฺสรติ.\csromanspacer{} โส เอวมาห “อตีตํปาหํ อทฺธานํ ชานามิ, ‘สํวฏฺฏิ วา โลโก วิวฏฺฏิ วา’ติ.\csromanspacer{} อนาคตํปาหํ อทฺธานํ ชานามิ ‘สํวฏฺฏิสฺสติ วา โลโก วิวฏฺฏิสฺสติ วา’ติ.\csromanspacer{} สสฺสโต อตฺตา จ โลโก จ วญฺโฌ กูฏฏฺโฐ เอสิกฏฺฐายิฏฺฐิโต.\csromanspacer{} เต จ สตฺตา สนฺธาวนฺติ สํสรนฺติ จวนฺติ อุปปชฺชนฺติ, อตฺถิเตฺวว สสฺสติสมนฺ”ติ.\csromanspacer{} อยํ ปฐโม สสฺสตวาโท.}

\prose{ปุน จปรํ ภนฺเต อิเธกจฺโจ สมโณ วา พฺราหฺมโณ วา อาตปฺปมนฺวาย ฯเปฯ ตถารูปํ เจโตสมาธิํ ผุสติ, ยถาสมาหิเต จิตฺเต อเนกวิหิตํ ปุพฺเพนิวาสํ อนุสฺสรติ.\csromanspacer{} เสยฺยถิทํ, เอกํปิ สํวฏฺฏวิวฏฺฏํ เทฺวปิ สํวฏฺฏวิวฏฺฏานิ ตีณิปิ สํวฏฺฏวิวฏฺฏานิ จตฺตาริปิ สํวฏฺฏวิวฏฺฏานิ ปญฺจปิ สํวฏฺฏวิวฏฺฏานิ ทสปิ สํวฏฺฏวิวฏฺฏานิ, “อมุตฺราสิํ เอวํนาโม เอวํโคตฺโต เอวํวณฺโณ เอวมาหาโร เอวํสุขทุกฺขปฺปฏิสํเวที เอวมายุปริยนฺโต, โส ตโต จุโต อมุตฺร อุทปาทิํ, ตตฺราปาสิํ เอวํนาโม}

\vspace{\tipitakaparskip}
\csromancenter{สมฺปสาทนียสุตฺต เอวํโคตฺโต เอวํวณฺโณ เอวมาหาโร เอวํสุขทุกฺขปฺปฏิสํเวที เอวมายุปริยนฺโต, โส ตโต จุโต อิธูปปนฺโน”ติ.\csromanspacer{} อิติ สาการํ ส-อุทฺเทสํ อเนกวิหิตํ ปุพฺเพนิวาสํ อนุสฺสรติ.\csromanspacer{} โส เอวมาห “อตีตํปาหํ อทฺธานํ ชานามิ ‘สํวฏฺฏิ วา โลโก วิวฏฺฏิ วา’ติ.\csromanspacer{} อนาคตญฺจ ขฺวาหํ อทฺธานํ ชานามิ ‘สํวฏฺฏิสฺสติ วา โลโก วิวฏฺฏิสฺสติ วา’ติ.\csromanspacer{} สสฺสโต อตฺตา จ โลโก จ วญฺโฌ กูฏฏฺโฐ เอสิกฏฺฐายิฏฺฐิโต, เต จ สตฺตา สนฺธาวนฺติ สํสรนฺติ จวนฺติ อุปปชฺชนฺติ, อตฺถิเตฺวว สสฺสติสมนฺ”ติ.\csromanspacer{} อยํ ภนฺเต ทุติโย สสฺสตวาโท.}
\prose{\csromanfolio{91}ปุน จปรํ ภนฺเต อิเธกจฺโจ สมโณ วา พฺราหฺมโณ วา อาตปฺปมนฺวาย ฯเปฯ ตถารูปํ เจโตสมาธิํ ผุสติ, ยถาสมาหิเต จิตฺเต อเนกวิหิตํ ปุพฺเพนิวาสํ อนุสฺสรติ.\csromanspacer{} เสยฺยถิทํ, ทสปิ สํวฏฺฏวิวฏฺฏานิ วีสํปิ สํวฏฺฏวิวฏฺฏานิ ติํสํปิ สํวฏฺฏวิวฏฺฏานิ จตฺตาลีสํปิ สํวฏฺฏวิวฏฺฏานิ, “อมุตฺราสิํ เอวํนาโม เอวํโคตฺโต เอวํวณฺโณ เอวมาหาโร เอวํสุขทุกฺขปฺปฏิสํเวที เอวมายุปริยนฺโต, โส ตโต จุโต อมุตฺร อุทปาทิํ, ตตฺราปาสิํ เอวํนาโม เอวํโคตฺโต เอวํวณฺโณ เอวมาหาโร เอวํสุขทุกฺขปฺปฏิสํเวที เอวมายุปริยนฺโต, โส ตโต จุโต อิธูปปนฺโนติ.\csromanspacer{} อิติ สาการํ ส-อุทฺเทสํ อเนกวิหิตํ ปุพฺเพนิวาสํ อนุสฺสรติ.\csromanspacer{} โส เอวมาห “อตีตํปาหํ อทฺธานํ ชานามิ ‘สํวฏฺฏิปิ โลโก วิวฏฺฏิปี’ติ, อนาคตํปาหํ อทฺธานํ ชานามิ ‘สํวฏฺฏิสฺสติปิ โลโก วิวฏฺฏิสฺสติปี’ติ, สสฺสโต อตฺตา จ โลโก จ วญฺโฌ กูฏฏฺโฐ เอสิกฏฺฐายิฏฺฐิโต, เต จ สตฺตา สนฺธาวนฺติ สํสรนฺติ จวนฺติ อุปปชฺชนฺติ, อตฺถิเตฺวว สสฺสติสมนฺ”ติ.\csromanspacer{} อยํ ภนฺเต ตติโย สสฺสตวาโท, เอตทานุตฺตริยํ ภนฺเต สสฺสตวาเทสุ.}

\csromanmatikaanchor{mtk.62}
\csromantocmarkref{section}{ปุพฺเพนิวาสานุสฺสติญาณเทสนา}{mtk.62}
\bookmark[dest={mtk.62},level=1]{ปุพฺเพนิวาสานุสฺสติญาณเทสนา}
\csromanheader{1}{\textbf{ปุพฺเพนิวาสานุสฺสติญาณเทสนา}}
\csromanmatikaanchor{mtk.63}
\csromantocmarkref{section}{จุตูปปาตญณเทสนา}{mtk.63}
\bookmark[dest={mtk.63},level=1]{จุตูปปาตญณเทสนา}
\proseitem{157}{อปรํ ปน ภนฺเต เอตทานุตฺตริยํ, ยถา ภควา ธมฺมํ เทเสติ ปุพฺเพนิวาสานุสฺสติญาเณ.\csromanspacer{} อิธ ภนฺเต เอกจฺโจ สมโณ วา พฺราหฺมโณ วา อาตปฺปมนฺวาย ฯเปฯ ตถารูปํ เจโตสมาธิํ ผุสติ, ยถาสมาหิเต จิตฺเต อเนกวิหิตํ ปุพฺเพนิวาสํ อนุสฺสรติ.\csromanspacer{} เสยฺยถิทํ, เอกํปิ ชาติํ เทฺวปิ ชาติโย ติสฺโสปิ ชาติโย จตสฺโสปิ ชาติโย ปญฺจปิ ชาติโย ทสปิ ชาติโย วีสํปิ ชาติโย ติํสํปิ ชาติโย \csromanfolio{92}จตฺตาลีสํปิ ชาติโย ปญฺญาสํปิ ชาติโย ชาติสตํปิ ชาติสหสฺสํปิ ชาติสตสหสฺสํปิ อเนเกปิ สํวฏฺฏกปฺเป อเนเกปิ วิวฏฺฏกปฺเป อเนเกปิ สํวฏฺฏวิวฏฺฏกปฺเป, “อมุตฺราสิํ เอวํนาโม เอวํโคตฺโต เอวํวณฺโณ เอวมาหาโร เอวํสุขทุกฺขปฺปฏิสํเวที เอวมายุปริยนฺโต, โส ตโต จุโต อมุตฺร อุทปาทิํ, ตตฺราปาสิํ เอวํนาโม เอวํโคตฺโต เอวํวณฺโณ เอวมาหาโร เอวํสุขทุกฺขปฺปฏิสํเวที เอวมายุปริยนฺโต, โส ตโต จุโต อิธูปปนฺโน”ติ.\csromanspacer{} อิติ สาการํ ส-อุทฺเทสํ อเนกวิหิตํ ปุพฺเพนิวาสํ อนุสฺสรติ.\csromanspacer{} สนฺติ ภนฺเต เทวา\footnote{สตฺตา (สฺยา)}, เยสํ น สกฺกา คณนาย วา สงฺขาเนน วา อายุ สงฺขาตุํ.\csromanspacer{} อปิ จ ยสฺมิํ ยสฺมิํ อตฺตภาเว อภินิวุฏฺฐปุพฺโพ\footnote{อภินิวุตฺถปุพฺโพ (สี, สฺยา, อิ)} โหติ ยทิ วา รูปีสุ ยทิ วา อรูปีสุ ยทิ วา สญฺญีสุ ยทิ วา อสญฺญีสุ ยทิ วา เนวสญฺญีนาสญฺญีสุ.\csromanspacer{} อิติ สาการํ ส-อุทฺเทสํ อเนกวิหิตํ ปุพฺเพนิวาสํ อนุสฺสรติ.\csromanspacer{} เอตทานุตฺตริยํ ภนฺเต ปุพฺเพนิวาสานุสฺสติญาเณ.}

\csromanheader{2}{\textbf{จุตูปปาตญาณเทสนา}}
\proseitem{158}{อปรํ ปน ภนฺเต เอตทานุตฺตริยํ, ยถา ภควา ธมฺมํ เทเสติ สตฺตานํ จุตูปปาตญาเณ.\csromanspacer{} อิธ ภนฺเต เอกจฺโจ สมโณ วา พฺราหฺมโณ วา อาตปฺปมนฺวาย ฯเปฯ ตถารูปํ เจโตสมาธิํ ผุสติ, ยถาสมาหิเต จิตฺเต ทิพฺเพน จกฺขุนา วิสุทฺเธน อติกฺกนฺตมานุสเกน สตฺเต ปสฺสติ จวมาเน อุปปชฺชมาเน หีเน ปณีเต สุวณฺเณ ทุพฺพณฺเณ สุคเต ทุคฺคเต ยถากมฺมูปเค สตฺเต ปชานาติ “อิเม วต โภนฺโต สตฺตา กายทุจฺจริเตน สมนฺนาคตา วจีทุจฺจริเตน สมนฺนาคตา มโนทุจฺจริเตน สมนฺนาคตา อริยานํ อุปวาทกา มิจฺฉาทิฏฺฐิกา มิจฺฉาทิฏฺฐิกมฺมสมาทานา, เต กายสฺส เภทา ปรํ มรณา อปายํ ทุคฺคติํ วินิปาตํ นิรยํ อุปปนฺนา.\csromanspacer{} อิเม วา ปน โภนฺโต สตฺตา กายสุจริเตน สมนฺนาคตา วจีสุจริเตน สมนฺนาคตา มโนสุจริเตน สมนฺนาคตา อริยานํ อนุปวาทกา สมฺมาทิฏฺฐิกา สมฺมาทิฏฺฐิกมฺมสมาทานา, เต กายสฺส เภทา ปรํ มรณา สุคติํ สคฺคํ โลกํ อุปปนฺนา”ติ.\csromanspacer{} อิติ ทิพฺเพน จกฺขุนา วิสุทฺเธน อติกฺกนฺตมานุสเกน สตฺเต ปสฺสติ จวมาเน อุปปชฺชมาเน หีเน}

\vspace{\tipitakaparskip}
\csromancenter{สมฺปสาทนียสุตฺต ปณีเต สุวณฺเณ ทุพฺพณฺเณ สุคเต ทุคฺคเต ยถากมฺมูปเค สตฺเต ปชานาติ.\csromanspacer{} เอตทานุตฺตริยํ ภนฺเต สตฺตานํ จุตูปปาตญาเณ.}
\csromanmatikaanchor{mtk.64}
\csromantocmarkref{section}{อิทฺธิวิธเทสนา}{mtk.64}
\bookmark[dest={mtk.64},level=1]{อิทฺธิวิธเทสนา}
\csromanheader{1}{\textbf{อิทฺธิวิธเทสนา}}
\proseitem{159}{\csromanfolio{93}อปรํ ปน ภนฺเต เอตทานุตฺตริยํ, ยถา ภควา ธมฺมํ เทเสติ อิทฺธิวิธาสุ.\csromanspacer{} เทฺวมา ภนฺเต อิทฺธิวิธาโย.\csromanspacer{} อตฺถิ ภนฺเต อิทฺธิ “สาสวา ส-อุปธิกา โน อริยา”ติ วุจฺจติ, อตฺถิ ภนฺเต อิทฺธิ “อนาสวา อนุปธิกา อริยา”ติ วุจฺจติ.\csromanspacer{} กตมา จ ภนฺเต อิทฺธิ “สาสวา ส-อุปธิกา โน อริยา”ติ วุจฺจติ.\csromanspacer{} อิธ ภนฺเต เอกจฺโจ สมโณ วา พฺราหฺมโณ วา อาตปฺปมนฺวาย ฯเปฯ ตถารูปํ เจโตสมาธิํ ผุสติ, ยถาสมาหิเต จิตฺเต อเนกวิหิตํ อิทฺธิวิธํ ปจฺจนุโภติ.\csromanspacer{} เอโกปิ หุตฺวา พหุธา โหติ, พหุธาปิ หุตฺวา เอโก โหติ, อาวิภาวํ ติโรภาวํ ติโรกุฏฺฏํ ติโรปาการํ ติโรปพฺพตํ อสชฺชมาโน คจฺฉติ เสยฺยถาปิ อากาเส.\csromanspacer{} ปถวิยาปิ อุมฺมุชฺชนิมุชฺชํ กโรติ เสยฺยถาปิ อุทเก.\csromanspacer{} อุทเกปิ อภิชฺชมาเน คจฺฉติ เสยฺยถาปิ ปถวิยํ.\csromanspacer{} อากาเสปิ ปลฺลงฺเกน กมติ เสยฺยถาปิ ปกฺขี สกุโณ.\csromanspacer{} อิเมปิ จนฺทิมสูริเย เอวํมหิทฺธิเก เอวํมหานุภาเว ปาณินา ปรามสติ ปริมชฺชติ.\csromanspacer{} ยาว พฺรหฺมโลกาปิ กาเยน วสํ วตฺเตติ.\csromanspacer{} อยํ ภนฺเต อิทฺธิ “สาสวา ส-อุปธิกา โน อริยา”ติ วุจฺจติ.}

\prose{กตมา ปน ภนฺเต อิทฺธิ “อนาสวา อนุปธิกา อริยา”ติ วุจฺจติ.\csromanspacer{} อิธ ภนฺเต ภิกฺขุ สเจ อากงฺขติ “ปฏิกูเล อปฺปฏิกูลสญฺญี วิหเรยฺยนฺ”ติ, อปฺปฏิกูลสญฺญี ตตฺถ วิหรติ.\csromanspacer{} สเจ อากงฺขติ “อปฺปฏิกูเล ปฏิกูลสญฺญี วิหเรยฺยนฺ”ติ, ปฏิกูลสญฺญี ตตฺถ วิหรติ.\csromanspacer{} สเจ อากงฺขติ “ปฏิกูเล จ อปฺปฏิกูเล จ อปฺปฏิกูลสญฺญี วิหเรยฺยนฺ”ติ, อปฺปฏิกูลสญฺญี ตตฺถ วิหรติ.\csromanspacer{} สเจ อากงฺขติ “ปฏิกูเล จ อปฺปฏิกูเล จ ปฏิกูลสญฺญี วิหเรยฺยนฺ”ติ, ปฏิกูลสญฺญี ตตฺถ วิหรติ.\csromanspacer{} สเจ อากงฺขติ “ปฏิกูลญฺจ อปฺปฏิกูลญฺจ ตทุภยํ อภินิวชฺเชตฺวา อุเปกฺขโก วิหเรยฺยํ สโต สมฺปชาโน”ติ, อุเปกฺขโก ตตฺถ วิหรติ สโต สมฺปชาโน.\csromanspacer{} อยํ ภนฺเต อิทฺธิ “อนาสวา อนุปธิกา อริยา”ติ วุจฺจติ.\csromanspacer{} เอตทานุตฺตริยํ ภนฺเต อิทฺธิวิธาสุ.\csromanspacer{} ตํ ภควา อเสสมภิชานาติ, ตํ ภควโต อเสสมภิชานโต อุตฺตริ อภิญฺเญยฺยํ นตฺถิ, ยทภิชานํ อญฺโญ \csromanfolio{94}สมโณ วา พฺราหฺมโณ วา ภควตา ภิยฺโยภิญฺญตโร อสฺส ยทิทํ อิทฺธิวิธาสุ.}

\csromanmatikaanchor{mtk.65}
\csromantocmarkref{section}{อญฺญถาสตฺถุคุณทสฺสน}{mtk.65}
\bookmark[dest={mtk.65},level=1]{อญฺญถาสตฺถุคุณทสฺสน}
\csromanheader{1}{\textbf{อญฺญถาสตฺถุคุณทสฺสน}}
\proseitem{160}{ยํ ตํ ภนฺเต สทฺเธน กุลปุตฺเตน ปตฺตพฺพํ อารทฺธวีริเยน ถามวตา ปุริสถาเมน ปุริสวีริเยน ปุริสปรกฺกเมน ปุริสโธรเยฺหน, อนุปฺปตฺตํ ตํ ภควตา.\csromanspacer{} น จ ภนฺเต ภควา กาเมสุ กามสุขลฺลิกานุโยคมนุยุตฺโต หีนํ คมฺมํ โปถุชฺชนิกํ อนริยํ อนตฺถสํหิตํ, น จ อตฺตกิลมถานุโยคมนุยุตฺโต ทุกฺขํ อนริยํ อนตฺถสํหิตํ.\csromanspacer{} จตุนฺนญฺจ ภควา ฌานานํ อาภิเจตสิกานํ ทิฏฺฐธมฺมสุขวิหารานํ นิกามลาภี อกิจฺฉลาภี อกสิรลาภี.}

\csromanmatikaanchor{mtk.66}
\csromantocmarkref{section}{อนุโยคทานปฺปการ}{mtk.66}
\bookmark[dest={mtk.66},level=1]{อนุโยคทานปฺปการ}
\csromanheader{1}{\textbf{อนุโยคทานปฺปการ}}
\proseitem{161}{สเจ มํ ภนฺเต เอวํ ปุจฺเฉยฺย “กิํ นุ โข อาวุโส สาริปุตฺต อเหสุํ อตีตมทฺธานํ อญฺเญ สมณา วา พฺราหฺมณา วา ภควตา ภิยฺโยภิญฺญตรา สมฺโพธิยนฺ”ติ, เอวํ ปุฏฺโฐ อหํ ภนฺเต “โน”ติ วเทยฺยํ.\csromanspacer{} กิํ ปนาวุโส สาริปุตฺต ภวิสฺสนฺติ อนาคตมทฺธานํ อญฺเญ สมณา วา พฺราหฺมณา วา ภควตา ภิยฺโยหิญฺญตรา สมฺโพธิยนฺติ, เอวํ ปุฏฺโฐ อหํ ภนฺเต “โน”ติ วเทยฺยํ.\csromanspacer{} กิํ ปนาวุโส สาริปุตฺต อตฺเถตรหิ อญฺโญ สมโณ วา พฺราหฺมโณ วา ภควตา ภิยฺโยภิญฺญตโร สมฺโพธิยนฺติ, เอวํ ปุฏฺโฐ อหํ ภนฺเต “โน”ติ วเทยฺยํ.}

\prose{สเจ ปน มํ ภนฺเต เอวํ ปุจฺเฉยฺย “กิํ นุ โข อาวุโส สาริปุตฺต อเหสุํ อตีตมทฺธานํ อญฺเญ สมณา วา พฺราหฺมณา วา ภควตา สมสมา สมฺโพธิยนฺ”ติ, เอวํ ปุฏฺโฐ อหํ ภนฺเต “เอวนฺ”ติ วเทยฺยํ.\csromanspacer{} กิํ ปนาวุโส สาริปุตฺต ภวิสฺสนฺติ อนาคตมทฺธานํ อญฺเญ สมณา วา พฺราหฺมณา วา ภควตา สมสมา สมฺโพธิยนฺติ, เอวํ ปุฏฺโฐ อหํ ภนฺเต “เอวนฺ”ติ วเทยฺยํ.\csromanspacer{} กิํ ปนาวุโส สาริปุตฺต อตฺเถตรหิ อญฺเญ สมณา วา พฺราหฺมณา วา ภควตา สมสมา สมฺโพธิยนฺติ, เอวํ ปุฏฺโฐ อหํ ภนฺเต “โน”ติ วเทยฺยํ.}

\prose{สเจ ปน ภนฺเต เอวํ ปุจฺเฉยฺย “กิํ ปนายสฺมา สาริปุตฺโต เอกจฺจํ อพฺภนุชานาติ, เอกจฺจํ น อพฺภนุชานาตี”ติ.\csromanspacer{} เอวํ ปุฏฺโฐ อหํ ภนฺเต เอวํ}

\vspace{\tipitakaparskip}
\csromancenter{สมฺปสาทนียสุตฺต พฺยากเรยฺยํ “สมฺมุขา เมตํ อาวุโส ภควโต สุตํ, สมฺมุขา ปฏิคฺคหิตํ ‘อเหสุํ อตีตมทฺธานํ อรหนฺโต สมฺมาสมฺพุทฺธา มยา สมสมา สมฺโพธิยนฺ’ติ.\csromanspacer{} สมฺมุขา เมตํ อาวุโส ภควโต สุตํ, สมฺมุขา ปฏิคฺคหิตํ ‘ภวิสฺสนฺติ อนาคตมทฺธานํ อรหนฺโต สมฺมาสมฺพุทฺธา มยา สมสมา สมฺโพธิยนฺ’ติ.\csromanspacer{} สมฺมุขา เมตํ อาวุโส ภควโต สุตํ, สมฺมุขา ปฏิคฺคหิตํ ‘อฏฺฐานเมตํ อนวกาโส ยํ เอกิสฺสา โลกธาตุยา เทฺว อรหนฺโต สมฺมาสมฺพุทฺธา อปุพฺพํ อจริมํ อุปฺปชฺเชยฺยุํ, เนตํ ฐานํ วิชฺชตี’ติ”.}
\prose{\csromanfolio{95}กจฺจาหํ ภนฺเต เอวํ ปุฏฺโฐ เอวํ พฺยากรมาโน วุตฺตวาที เจว ภควโต โหมิ, น จ ภควนฺตํ อภูเตน อพฺภาจิกฺขามิ, ธมฺมสฺส จานุธมฺมํ พฺยากโรมิ, น จ โกจิ สหธมฺมิโก วาทานุวาโท\footnote{วาทานุปาโต (สี)} คารยฺหํ ฐานํ อาคจฺฉตีติ.\csromanspacer{} ตคฺฆ ตฺวํ สาริปุตฺต เอวํ ปุฏฺโฐ เอวํ พฺยากรมาโน วุตฺตวาที เจว เม โหสิ, น จ มํ อภูเตน อพฺภาจิกฺขสิ, ธมฺมสฺส จานุธมฺมํ พฺยากโรสิ, น จ โกจิ สหธมฺมิโก วาทานุวาโท คารยฺหํ ฐานํ อาคจฺฉตีติ.}

\csromanmatikaanchor{mtk.67}
\csromantocmarkref{section}{อจฺฉริย-อพฺภุต}{mtk.67}
\bookmark[dest={mtk.67},level=1]{อจฺฉริย-อพฺภุต}
\csromanheader{1}{\textbf{อจฺฉริย-อพฺภุต}}
\proseitem{162}{เอวํ วุตฺเต อายสฺมา อุทายี ภควนฺตํ เอตทโวจ “อจฺฉริยํ ภนฺเต, อพฺภุตํ ภนฺเต, ตถาคตสฺส อปฺปิจฺฉตา สนฺตุฏฺฐิตา สลฺเลขตา.\csromanspacer{} ยตฺร หิ นาม ตถาคโต เอวํมหิทฺธิโก เอวํมหานุภาโว, อถ จ ปน เนวตฺตานํ ปาตุกริสฺสติ.\csromanspacer{} เอกเมกญฺเจปิ อิโต ภนฺเต ธมฺมํ อญฺญติตฺถิยา ปริพฺพาชกา อตฺตนิ สมนุปสฺเสยฺยุํ, เต ตาวตเกเนว ปฏากํ ปริหเรยฺยุํ.\csromanspacer{} อจฺฉริยํ ภนฺเต, อพฺภุตํ ภนฺเต, ตถาคตสฺส อปฺปิจฺฉตา สนฺตุฏฺฐิตา สลฺเลขตา.\csromanspacer{} ยตฺร หิ นาม ตถาคโต เอวํมหิทฺธิโก เอวํมหานุภาโว, อถ จ ปน เนวตฺตานํ ปาตุกริสฺสตี”ติ.}

\prose{ปสฺส โข ตฺวํ อุทายิ ตถาคตสฺส อปฺปิจฺฉตา สนฺตุฏฺฐิตา สลฺเลขตา.\csromanspacer{} ยตฺร หิ นาม ตถาคโต เอวํมหิทฺธิโก เอวํมหานุภาโว, อถ จ ปน เนวตฺตานํ ปาตุกริสฺสติ.\csromanspacer{} เอกเมกญฺเจปิ อิโต อุทายิ ธมฺมํ อญฺญติตฺถิยา ปริพฺพาชกา อตฺตนิ สมนุปสฺเสยฺยุํ, เต ตาวตเกเนว \csromanfolio{96}ปฏากํ ปริหเรยฺยุํ.\csromanspacer{} ปสฺส โข ตฺวํ อุทายิ ตถาคตสฺส อปฺปิจฺฉตา สนฺตุฏฺฐิตา สลฺเลขตา.\csromanspacer{} ยตฺร หิ นาม ตถาคโต เอวํมหิทฺธิโก เอวํมหานุภาโว, อถ จ ปน เนวตฺตานํ ปาตุกริสฺสตีติ.}

\proseitem{163}{อถ โข ภควา อายสฺมนฺตํ สาริปุตฺตํ อามนฺเตสิ “ตสฺมา ติห ตฺวํ สาริปุตฺต อิมํ ธมฺมปริยายํ อภิกฺขณํ ภาเสยฺยาสิ ภิกฺขูนํ ภิกฺขุนีนํ อุปาสกานํ อุปาสิกานํ.\csromanspacer{} เยสํปิ หิ สาริปุตฺต โมฆปุริสานํ ภวิสฺสติ ตถาคเต กงฺขา วา วิมติ วา, เตสมิมํ ธมฺมปริยายํ สุตฺวา ตถาคเต กงฺขา วา วิมติ วา, สา ปหียิสฺสตี”ติ.\csromanspacer{} อิติ หิทํ อายสฺมา สาริปุตฺโต ภควโต สมฺมุขา สมฺปสาทํ ปเวเทสิ, ตสฺมา อิมสฺส เวยฺยากรณสฺส สมฺปสาทนียนฺเตฺวว อธิวจนนฺติ.}

\nitthitam{\textbf{สมฺปสาทนียสุตฺตํ นิฏฺฐิตํ ปญฺจมํ.}}
\csromanmatikaanchor{mtk.68}
\csromantocmarknopage{chapter}{6. ปาสาทิกสุตฺต}{mtk.68}
\bookmark[dest={mtk.68},level=0]{6. ปาสาทิกสุตฺต}
\chapterheadpageread{6. ปาสาทิกสุตฺต}
\proseitem{164}{\csromanfolio{97}เอวํ เม สุตํ\csromandash{}เอกํ สมยํ ภควา สกฺเกสุ วิหรติ เวธญฺญา นาม สกฺยา, เตสํ อมฺพวเน ปาสาเท.}

\csromanmatikaanchor{mtk.69}
\csromantocmarkref{section}{นิคณฺฐนาฏปุตฺตกาลงฺกิริยา}{mtk.69}
\bookmark[dest={mtk.69},level=1]{นิคณฺฐนาฏปุตฺตกาลงฺกิริยา}
\csromanheader{1}{\textbf{นิคณฺฐนาฏปุตฺตกาลงฺกิริยา}}
\prose{เตน โข ปน สมเยน นิคณฺโฐ นาฏปุตฺโต\footnote{นาถปุตฺโต (สี, อิ)} ปาวายํ อธุนากาลงฺกโต โหติ.\csromanspacer{} ตสฺส กาลงฺกิริยาย ภินฺนา นิคณฺฐา เทฺวธิกชาตา ภณฺฑนชาตา กลหชาตา วิวาทาปนฺนา อญฺญมญฺญํ มุขสตฺตีหิ วิตุทนฺตา วิหรนฺติ “น ตฺวํ อิมํ ธมฺมวินยํ อาชานาสิ, อหํ อิมํ ธมฺมวินยํ อาชานามิ, กิํ ตฺวํ อิมํ ธมฺมวินยํ อาชานิสฺสสิ.\csromanspacer{} มิจฺฉาปฏิปนฺโน ตฺวมสิ, อหมสฺมิ สมฺมาปฏิปนฺโน.\csromanspacer{} สหิตํ เม, อสหิตํ เต.\csromanspacer{} ปุเรวจนียํ ปจฺฉา อวจ, ปจฺฉาวจนียํ ปุเร อวจ.\csromanspacer{} อธิจิณฺณํ เต วิปราวตฺตํ, อาโรปิโต เต วาโท, นิคฺคหิโต ตฺวมสิ, จร วาทปฺปโมกฺขาย, นิพฺเพเฐหิ วา สเจ ปโหสี”ติ.\csromanspacer{} วโธเยว โข\footnote{วโธเยเวโก (ก)} มญฺเญ นิคณฺเฐสุ นาฏปุตฺติเยสุ วตฺตติ\footnote{อนุวตฺตติ (สฺยา, ก)}.\csromanspacer{} เยปิ นิคณฺฐสฺส นาฏปุตฺตสฺส สาวกา คิหี โอทาตวสนา, เตปิ\footnote{เต เตสุ (ก)} นิคณฺเฐสุ นาฏปุตฺติเยสุ นิพฺพินฺนรูปา\footnote{นิพฺพินฺทรูปา (ก)} วิรตฺตรูปา ปฏิวานรูปา, ยถา ตํ ทุรกฺขาเต ธมฺมวินเย ทุปฺปเวทิเต อนิยฺยานิเก อนุปสมสํวตฺตนิเก อสมฺมาสมฺพุทฺธปฺปเวทิเต ภินฺนถูเป อปฺปฏิสรเณ.}

\proseitem{165}{อถ โข จุนฺโท สมณุทฺเทโส ปาวายํ วสฺสํวุฏฺโฐ\footnote{วสฺสํวุตฺโถ (สี, สฺยา, อิ)} เยน สามคาโม, เยนายสฺมา อานนฺโท เตนุปสงฺกมิ, อุปสงฺกมิตฺวา อายสฺมนฺตํ อานนฺทํ อภิวาเทตฺวา เอกมนฺตํ นิสีทิ, เอกมนฺตํ นิสินฺโน โข จุนฺโท สมณุทฺเทโส อายสฺมนฺตํ อานนฺทํ เอตทโวจ “นิคณฺโฐ ภนฺเต นาฏปุตฺโต ปาวายํ อธุนากาลงฺกโต.\csromanspacer{} ตสฺส กาลงฺกิริยาย ภินฺนา นิคณฺฐา เทฺวธิกชาตา ฯเปฯ ภินฺนถูเป อปฺปฏิสรเณ”ติ.}

\prose{เอวํ วุตฺเต อายสฺมา อานนฺโท จุนฺทํ สมณุทฺเทสํ เอตทโวจ “อตฺถิ โข อิทํ อาวุโส จุนฺท กถาปาภตํ ภควนฺตํ ทสฺสนาย. \csromanfolio{98}อายามาวุโส จุนฺท เยน ภควา เตนุปสงฺกมิสฺสาม, อุปสงฺกมิตฺวา เอตมตฺถํ ภควโต อาโรเจสฺสามา”ติ\footnote{อาโรเจยฺยามาติ (สฺยา)}. “เอวํ ภนฺเต”ติ โข จุนฺโท สมณุทฺเทโส อายสฺมโต อานนฺทสฺส ปจฺจสฺโสสิ.}

\prose{อถ โข อายสฺมา จ อานนฺโท จุนฺโท จ สมณุทฺเทโส เยน ภควา เตนุปสงฺกมิํสุ, อุปสงฺกมิตฺวา ภควนฺตํ อภิวาเทตฺวา เอกมนฺตํ นิสีทิํสุ, เอกมนฺตํ นิสินฺโน โข อายสฺมา อานนฺโท ภควนฺตํ เอตทโวจ “อยํ ภนฺเต จุนฺโท สมณุทฺเทโส เอวมาห, นิคณฺโฐ ภนฺเต นาฏปุตฺโต ปาวายํ อธุนากาลงฺกโต, ตสฺส กาลงฺกิริยาย ภินฺนา นิคณฺฐา ฯเปฯ ภินฺนถูเป อปฺปฏิสรเณ”ติ.}

\csromanmatikaanchor{mtk.70}
\csromantocmarkref{section}{อสมฺมาสมฺพุทฺธปฺปเวทิตธมฺมวินย}{mtk.70}
\bookmark[dest={mtk.70},level=1]{อสมฺมาสมฺพุทฺธปฺปเวทิตธมฺมวินย}
\csromanheader{1}{\textbf{อสมฺมาสมฺพุทฺธปฺปเวทิตธมฺมวินย}}
\proseitem{166}{เอวํ เหตํ จุนฺท โหติ ทุรกฺขาเต ธมฺมวินเย ทุปฺปเวทิเต อนิยฺยานิเก อนุปสมสํวตฺตนิเก อสมฺมาสมฺพุทฺธปฺปเวทิเต.\csromanspacer{} อิธ จุนฺท สตฺถา จ โหติ อสมฺมาสมฺพุทฺโธ, ธมฺโม จ ทุรกฺขาโต ทุปฺปเวทิโต อนิยฺยานิโก อนุปสมสํวตฺตนิโก อสมฺมาสมฺพุทฺธปฺปเวทิโต, สาวโก จ ตสฺมิํ ธมฺเม น ธมฺมานุธมฺมปฺปฏิปนฺโน วิหรติ, น สามีจิปฺปฏิปนฺโน, น อนุธมฺมจารี, โวกฺกมฺม จ ตมฺหา ธมฺมา วตฺตติ.\csromanspacer{} โส เอวมสฺส วจนีโย “ตสฺส เต อาวุโส ลาภา, ตสฺส เต สุลทฺธํ, สตฺถา จ เต อสมฺมาสมฺพุทฺโธ, ธมฺโม จ ทุรกฺขาโต ทุปฺปเวทิโต อนิยฺยานิโก อนุปสมสํวตฺตนิโก อสมฺมาสมฺพุทฺธปฺปเวทิโต.\csromanspacer{} ตฺวญฺจ ตสฺมิํ ธมฺเม น ธมฺมานุธมฺมปฺปฏิปนฺโน วิหรสิ, น สามีจิปฺปฏิปนฺโน, น อนุธมฺมจารี, โวกฺกมฺม จ ตมฺหา ธมฺมา วตฺตสี”ติ.\csromanspacer{} อิติ โข จุนฺท สตฺถาปิ ตตฺถ คารโยฺห, ธมฺโมปิ ตตฺถ คารโยฺห, สาวโก จ ตตฺถ เอวํ ปาสํโส.\csromanspacer{} โย โข จุนฺท เอวรูปํ สาวกํ เอวํ วเทยฺย “เอตายสฺมา ตถา ปฏิปชฺชตุ, ยถา เต สตฺถารา ธมฺโม เทสิโต ปญฺญตฺโต”ติ.\csromanspacer{} โย จ สมาทเปติ\footnote{สมาทาเปติ (สี-ฏฺฐ)}, ยญฺจ สมาทเปติ, โย จ สมาทปิโต\footnote{สมาทาปิโต (สี-ฏฺฐ)} ตถตฺตาย ปฏิปชฺชติ.\csromanspacer{} สพฺเพ เต พหุํ อปุญฺญํ ปสวนฺติ.\csromanspacer{} ตํ กิสฺส เหตุ, เอวเญฺหตํ จุนฺท โหติ ทุรกฺขาเต ธมฺมวินเย ทุปฺปเวทิเต อนิยฺยานิเก อนุปสมสํวตฺตนิเก อสมฺมาสมฺพุทฺธปฺปเวทิเต.}

\csromanheader{2}{6. ปาสาทิกสุตฺต}
\proseitem{167}{\csromanfolio{99}อิธ ปน จุนฺท สตฺถา จ โหติ อสมฺมาสมฺพุทฺโธ, ธมฺโม จ ทุรกฺขาโต ทุปฺปเวทิโต อนิยฺยานิโก อนุปสมสํวตฺตนิโก อสมฺมาสมฺพุทฺธปฺปเวทิโต, สาวโก จ ตสฺมิํ ธมฺเม ธมฺมานุธมฺมปฺปฏิปนฺโน วิหรติ สามีจิปฺปฏิปนฺโน อนุธมฺมจารี, สมาทาย ตํ ธมฺมํ วตฺตติ.\csromanspacer{} โส เอวมสฺส วจนีโย “ตสฺส เต อาวุโส อลาภา, ตสฺส เต ทุลฺลทฺธํ, สตฺถา จ เต อสมฺมาสมฺพุทฺโธ, ธมฺโม จ ทุรกฺขาโต ทุปฺปเวทิโต อนิยฺยานิโก อนุปสมสํวตฺตนิโก อสมฺมาสมฺพุทฺธปฺปเวทิโต.\csromanspacer{} ตฺวญฺจ ตสฺมิํ ธมฺเม ธมฺมานุธมฺมปฏิปนฺโน วิหรสิ สามีจิปฺปฏิปนฺโน อนุธมฺมจารี, สมาทาย ตํ ธมฺมํ วตฺตสี”ติ.\csromanspacer{} อิติ โข จุนฺท สตฺถาปิ ตตฺถ คารโยฺห, ธมฺโมปิ ตตฺถ คารโยฺห, สาวโกปิ ตตฺถ เอวํ คารโยฺห.\csromanspacer{} โย โข จุนฺท เอวรูปํ สาวกํ เอวํ วเทยฺย “อทฺธายสฺมา ญายปฺปฏิปนฺโน ญายมาราเธสฺสตี”ติ.\csromanspacer{} โย จ ปสํสติ, ยญฺจ ปสํสติ, โย จ ปสํสิโต ภิยฺโยโส มตฺตาย วีริยํ อารภติ.\csromanspacer{} สพฺเพ เต พหุํ อปุญฺญํ ปสวนฺติ.\csromanspacer{} ตํ กิสฺส เหตุ, เอวเญฺหตํ จุนฺท โหติ ทุรกฺขาเต ธมฺมวินเย ทุปฺปเวทิเต อนิยฺยานิเก อนุปสมสํวตฺตนิเก อสมฺมาสมฺพุทฺธปฺปเวทิเต.}

\csromanmatikaanchor{mtk.71}
\csromantocmarkref{section}{สมฺมาสมฺพุทฺธปฺปเวทิตธมฺมวินย}{mtk.71}
\bookmark[dest={mtk.71},level=1]{สมฺมาสมฺพุทฺธปฺปเวทิตธมฺมวินย}
\csromanheader{1}{\textbf{สมฺมาสมฺพุทฺธปฺปเวทิตธมฺมวินย}}
\proseitem{168}{อิธ ปน จุนฺท สตฺถา จ โหติ สมฺมาสมฺพุทฺโธ, ธมฺโม จ สฺวากฺขาโต สุปฺปเวทิโต นิยฺยานิโก อุปสมสํวตฺตนิโก สมฺมาสมฺพุทฺธปฺปเวทิโต, สาวโก จ ตสฺมิํ ธมฺเม น ธมฺมานุธมฺมปฺปฏิปนฺโน วิหรติ, น สามีจิปฺปฏิปนฺโน, น อนุธมฺมจารี, โวกฺกมฺม จ ตมฺหา ธมฺมา วตฺตติ.\csromanspacer{} โส เอวมสฺส วจนีโย “ตสฺส เต อาวุโส อลาภา, ตสฺส เต ทุลฺลทฺธํ, สตฺถา จ เต สมฺมาสมฺพุทฺโธ, ธมฺโม จ สฺวากฺขาโต สุปฺปเวทิโต นิยฺยานิโก อุปสมสํวตฺตนิโก สมฺมาสมฺพุทฺธปฺปเวทิโต.\csromanspacer{} ตฺวญฺจ ตสฺมิํ ธมฺเม น ธมฺมานุธมฺมปฺปฏิปนฺโน วิหรสิ, น สามีจิปฺปฏิปนฺโน, น อนุธมฺมจารี, โวกฺกมฺม จ ตมฺหา ธมฺมา วตฺตสี”ติ.\csromanspacer{} อิติ โข จุนฺท สตฺถาปิ ตตฺถ ปาสํโส, ธมฺโมปิ ตตฺถ ปาสํโส, สาวโก จ ตตฺถ เอวํ คารโยฺห.\csromanspacer{} โย โข จุนฺท เอวรูปํ สาวกํ เอวํ วเทยฺย “เอตายสฺมา ตถา ปฏิปชฺชตุ, ยถา เต สตฺถารา ธมฺโม เทสิโต ปญฺญตฺโต”ติ.\csromanspacer{} โย จ สมาทเปติ, ยญฺจ สมาทเปติ, โย จ สมาทปิโต ตถตฺตาย ปฏิปชฺชติ.\csromanspacer{} สพฺเพ เต พหุํ ปุญฺญํ ปสวนฺติ.\csromanspacer{} ตํ กิสฺส เหตุ, \csromanfolio{100}เอวเญฺหตํ จุนฺท โหติ สฺวากฺขาเต ธมฺมวินเย สุปฺปเวทิเต นิยฺยานิเก อุปสมสํวตฺตนิเก สมฺมาสมฺพุทฺธปฺปเวทิเต.}

\proseitem{169}{อิธ ปน จุนฺท สตฺถา จ โหติ สมฺมาสมฺพุทฺโธ, ธมฺโม จ สฺวากฺขาโต สุปฺปเวทิโต นิยฺยานิโก อุปสมสํวตฺตนิโก สมฺมาสมฺพุทฺธปฺปเวทิโต, สาวโก จ ตสฺมิํ ธมฺเม ธมฺมานุธมฺมปฺปฏิปนฺโน วิหรติ สามีจิปฺปฏิปนฺโน อนุธมฺมจารี, สมาทาย ตํ ธมฺมํ วตฺตติ.\csromanspacer{} โส เอวมสฺส วจนีโย “ตสฺส เต อาวุโส ลาภา, ตสฺส เต สุลทฺธํ, สตฺถา จ เต\footnote{สตฺถา จ เต อรหํ (สฺยา)} สมฺมาสมฺพุทฺโธ, ธมฺโม จ สฺวากฺขาโต สุปฺปเวทิโต นิยฺยานิโก อุปสมสํวตฺตนิโก สมฺมาสมฺพุทฺธปฺปเวทิโต.\csromanspacer{} ตฺวญฺจ ตสฺมิํ ธมฺเม ธมฺมานุธมฺมปฺปฏิปนฺโน วิหรสิ สามีจิปฺปฏิปนฺโน อนุธมฺมจารี, สมาทาย ตํ ธมฺมํ วตฺตสี”ติ.\csromanspacer{} อิติ โข จุนฺท สตฺถาปิ ตตฺถ ปาสํโส, ธมฺโมปิ ตตฺถ ปาสํโส, สาวโกปิ ตตฺถ เอวํ ปาสํโส.\csromanspacer{} โย โข จุนฺท เอวรูปํ สาวกํ เอวํ วเทยฺย “อทฺธายสฺมา ญายปฺปฏิปนฺโน ญายมาราเธสฺสตี”ติ.\csromanspacer{} โย จ ปสํสติ, ยญฺจ ปสํสติ, โย จ ปสํสิโต\footnote{ปสตฺโถ (สฺยา)} ภิยฺโยโส มตฺตาย วีริยํ อารภติ.\csromanspacer{} สพฺเพ เต พหุํ ปุญฺญํ ปสวนฺติ.\csromanspacer{} ตํ กิสฺส เหตุ, เอวเญฺหตํ จุนฺท โหติ สฺวากฺขาเต ธมฺมวินเย สุปฺปเวทิเต นิยฺยานิเก อุปสมสํวตฺตนิเก สมฺมาสมฺพุทฺธปฺปเวทิเต.}

\csromanmatikaanchor{mtk.72}
\csromantocmarkref{section}{สาวกานุตปฺปสตฺถุ}{mtk.72}
\bookmark[dest={mtk.72},level=1]{สาวกานุตปฺปสตฺถุ}
\csromanheader{1}{\textbf{สาวกานุตปฺปสตฺถุ}}
\proseitem{170}{อิธ ปน จุนฺท สตฺถา จ โลเก อุทปาทิ อรหํ สมฺมาสมฺพุทฺโธ, ธมฺโม จ สฺวากฺขาโต สุปฺปเวทิโต นิยฺยานิโก อุปสมสํวตฺตนิโก สมฺมาสมฺพุทฺธปฺปเวทิโต, อวิญฺญาปิตตฺถา จสฺส โหนฺติ สาวกา สทฺธมฺเม, น จ เตสํ เกวลํ ปริปูรํ พฺรหฺมจริยํ อาวิกตํ โหติ อุตฺตานีกตํ สพฺพสงฺคาหปทกตํ สปฺปาฏิหีรกตํ ยาว เทวมนุสฺเสหิ สุปฺปกาสิตํ.\csromanspacer{} อถ เนสํ สตฺถุโน อนฺตรธานํ โหติ.\csromanspacer{} เอวรูโป โข จุนฺท สตฺถา สาวกานํ กาลงฺกโต อนุตปฺโป โหติ.\csromanspacer{} ตํ กิสฺส เหตุ, สตฺถา จ โลเก อุทปาทิ อรหํ สมฺมาสมฺพุทฺโธ, ธมฺโม จ สฺวากฺขาโต สุปฺปเวทิโต นิยฺยานิโก อุปสมสํวตฺตนิโก สมฺมาสมฺพุทฺธปฺปเวทิโต, อวิญฺญาปิตตฺถา จมฺห สทฺธมฺเม, น จ โน เกวลํ ปริปูรํ พฺรหฺมจริยํ อาวิกตํ}

\vspace{\tipitakaparskip}
\csromancenter{ปาสาทิกสุตฺต โหติ อุตฺตานีกตํ สพฺพสงฺคาหปทกตํ สปฺปาฏิหีรกตํ ยาว เทวมนุสฺเสหิ สุปฺปกาสิตํ.\csromanspacer{} อถ โน สตฺถุโน อนฺตรธานํ โหตีติ.\csromanspacer{} เอวรูโป โข จุนฺท สตฺถา สาวกานํ กาลงฺกโต อนุตปฺโป โหติ.}
\csromanmatikaanchor{mtk.73}
\csromantocmarkref{section}{สาวกานนุตปฺปสตฺถุ}{mtk.73}
\bookmark[dest={mtk.73},level=1]{สาวกานนุตปฺปสตฺถุ}
\csromanheader{1}{\textbf{สาวกานนุตปฺปสตฺถุ}}
\proseitem{171}{\csromanfolio{101}อิธ ปน จุนฺท สตฺถา จ โลเก อุทปาทิ อรหํ สมฺมาสมฺพุทฺโธ, ธมฺโม จ สฺวากฺขาโต สุปฺปเวทิโต นิยฺยานิโก อุปสมสํวตฺตนิโก สมฺมาสมฺพุทฺธปฺปเวทิโต, วิญฺญาปิตตฺถา จสฺส โหนฺติ สาวกา สทฺธมฺเม, เกวลญฺจ เตสํ ปริปูรํ พฺรหฺมจริยํ อาวิกตํ โหติ อุตฺตานีกตํ สพฺพสงฺคาหปทกตํ สปฺปาฏิหีรกตํ ยาว เทวมนุสฺเสหิ สุปฺปกาสิตํ.\csromanspacer{} อถ เนสํ สตฺถุโน อนฺตรธานํ โหติ.\csromanspacer{} เอวรูโป โข จุนฺท สตฺถา สาวกานํ กาลงฺกโต อนนุตปฺโป โหติ.\csromanspacer{} ตํ กิสฺส เหตุ, สตฺถา จ โน โลเก อุทปาทิ อรหํ สมฺมาสมฺพุทฺโธ, ธมฺโม จ สฺวากฺขาโต สุปฺปเวทิโต นิยฺยานิโก อุปสมสํวตฺตนิโก สมฺมาสมฺพุทฺธปฺปเวทิโต, วิญฺญาปิตตฺถา จมฺห สทฺธมฺเม, เกวลญฺจ โน ปริปูรํ พฺรหฺมจริยํ อาวิกตํ โหติ อุตฺตานีกตํ สพฺพสงฺคาหปทกตํ สปฺปาฏิหีรกตํ ยาว เทวมนุสฺเสหิ สุปฺปกาสิตํ.\csromanspacer{} อถ โน สตฺถุโน อนฺตรธานํ โหตีติ.\csromanspacer{} เอวรูโป โข จุนฺท สตฺถา สาวกานํ กาลงฺกโต อนนุตปฺโป โหติ.}

\csromanmatikaanchor{mtk.74}
\csromantocmarkref{section}{พฺรหฺมจริย-อปริปูราทิกถา}{mtk.74}
\bookmark[dest={mtk.74},level=1]{พฺรหฺมจริย-อปริปูราทิกถา}
\csromanheader{1}{\textbf{พฺรหฺมจริย-อปริปูราทิกถา}}
\proseitem{172}{เอเตหิ เจปิ จุนฺท องฺเคหิ สมนฺนาคตํ พฺรหฺมจริยํ โหติ, โน จ โข สตฺถา โหติ เถโร รตฺตญฺญู จิรปพฺพชิโต อทฺธคโต วโย-อนุปฺปตฺโต.\csromanspacer{} เอวํ ตํ พฺรหฺมจริยํ อปริปูรํ โหติ เตนงฺเคน.}

\prose{ยโต จ โข จุนฺท เอเตหิ เจว องฺเคหิ สมนฺนาคตํ พฺรหฺมจริยํ โหติ, สตฺถา จ โหติ เถโร รตฺตญฺญู จิรปพฺพชิโต อทฺธคโต วโย-อนุปฺปตฺโต.\csromanspacer{} เอวํ ตํ พฺรหฺมจริยํ ปริปูรํ โหติ เตนงฺเคน.}

\proseitem{173}{เอเตหิ เจปิ จุนฺท องฺเคหิ สมนฺนาคตํ พฺรหฺมจริยํ โหติ, สตฺถา จ โหติ เถโร รตฺตญฺญู จิรปพฺพชิโต อทฺธคโต วโย-อนุปฺปตฺโต, โน จ ขฺวสฺส เถรา ภิกฺขู สาวกา โหนฺติ วิยตฺตา วินีตา วิสารทา ปตฺตโยคกฺเขมา, อลํ สมกฺขาตุํ สทฺธมฺมสฺส, อลํ \csromanfolio{102}อุปฺปนฺนํ ปรปฺปวาทํ สหธมฺเมหิ สุนิคฺคหิตํ นิคฺคเหตฺวา สปฺปาฏิหาริยํ ธมฺมํ เทเสตุํ.\csromanspacer{} เอวํ ตํ พฺรหฺมจริยํ อปริปูรํ โหติ เตนงฺเคน.}

\prose{ยโต จ โข จุนฺท เอเตหิ เจว องฺเคหิ สมนฺนาคตํ พฺรหฺมจริยํ โหติ, สตฺถา จ โหติ เถโร รตฺตญฺญู จิรปพฺพชิโต อทฺธคโต วโย-อนุปฺปตฺโต, เถรา จสฺส ภิกฺขู สาวกา โหนฺติ วิยตฺตา วินีตา วิสารทา ปตฺตโยคกฺเขมา, อลํ สมกฺขาตุํ สทฺธมฺมสฺส, อลํ อุปฺปนฺนํ ปรปฺปวาทํ สหธมฺเมหิ สุนิคฺคหิตํ นิคฺคเหตฺวา สปฺปาฏิหาริยํ ธมฺมํ เทเสตุํ.\csromanspacer{} เอวํ ตํ พฺรหฺมจริยํ ปริปูรํ โหติ เตนงฺเคน.}

\proseitem{174}{เอเตหิ เจปิ จุนฺท องฺเคหิ สมนฺนาคตํ พฺรหฺมจริยํ โหติ, สตฺถา จ โหติ เถโร รตฺตญฺญู จิรปพฺพชิโต อทฺธคโต วโย-อนุปฺปตฺโต, เถรา จสฺส ภิกฺขู สาวกา โหนฺติ วิยตฺตา วินีตา วิสารทา ปตฺตโยคกฺเขมา, อลํ สมกฺขาตุํ สทฺธมฺมสฺส, อลํ อุปฺปนฺนํ ปรปฺปวาทํ สหธมฺเมหิ สุนิคฺคหิตํ นิคฺคเหตฺวา สปฺปาฏิหาริยํ ธมฺมํ เทเสตุํ.\csromanspacer{} โน จ ขฺวสฺส มชฺฌิมา ภิกฺขู สาวกา โหนฺติ ฯเปฯ มชฺฌิมา จสฺส ภิกฺขู สาวกา โหนฺติ, โน จ ขฺวสฺส นวา ภิกฺขู สาวกา โหนฺติ ฯเปฯ นวา จสฺส ภิกฺขู สาวกา โหนฺติ, โน จ ขฺวสฺส เถรา ภิกฺขุนิโย สาวิกา โหนฺติ ฯเปฯ เถรา จสฺส ภิกฺขุนิโต สาวิกา โหนฺติ, โน จ ขฺวสฺส มชฺฌิมา ภิกฺขุนิโย สาวิกา โหนฺติ ฯเปฯ มชฺฌิมา จสฺส ภิกฺขุนิโย สาวิกา โหนฺติ, โน จ ขฺวสฺส นวา ภิกฺขุนิโย สาวิกา โหนฺติ ฯเปฯ นวา จสฺส ภิกฺขุนิโย สาวิกา โหนฺติ, โน จ ขฺวสฺส อุปาสกา สาวกา โหนฺติ คิหี โอทาตวสนา พฺรหฺมจาริโน ฯเปฯ อุปาสกา จสฺส สาวกา โหนฺติ คิหี โอทาตวสนา พฺรหฺมจาริโน, โน จ ขฺวสฺส อุปาสกา สาวกา โหนฺติ คิหี โอทาตวสนา กามโภคิโน ฯเปฯ อุปาสกา จสฺส สาวกา โหนฺติ คิหี โอทาตวสนา กามโภคิโน, โน จ ขฺวสฺส อุปาสิกา สาวิกา โหนฺติ คิหินิโย โอทาตวสนา พฺรหฺมจารินิโย ฯเปฯ อุปาสิกา จสฺส สาวิกา โหนฺติ คิหินิโย โอทาตวสนา พฺรหฺมจารินิโย, โน จ ขฺวสฺส อุปาสิกา สาวิกา โหนฺติ คิหินิโย โอทาตวสนา กามโภคินิโย ฯเปฯ อุปาสิกา จสฺส สาวิกา โหนฺติ คิหินิโย โอทาตวสนา กามโภคินิโย, โน จ ขฺวสฺส พฺรหฺมจริยํ โหติ อิทฺธญฺเจว ผีตญฺจ วิตฺถาริกํ พาหุชญฺญํ ปุถุภูตํ ยาว เทวมนุสฺเสหิ สุปฺปกาสิตํ ฯเปฯ พฺรหฺมจริยญฺจสฺส โหติ อิทฺธญฺเจว ผีตญฺจ วิตฺถาริกํ พาหุชญฺญํ ปุถุภูตํ}

\vspace{\tipitakaparskip}
\csromancenter{ปาสาทิกสุตฺต ยาว เทวมนุสฺเสหิ สุปฺปกาสิตํ, โน จ โข ลาภคฺคยสคฺคปฺปตฺตํ.\csromanspacer{} เอวํ ตํ พฺรหฺมจริยํ อปริปูรํ โหติ เตนงฺเคน.}
\prose{\csromanfolio{103}ยโต จ โข จุนฺท เอเตหิ เจว องฺเคหิ สมนฺนาคตํ พฺรหฺมจริยํ โหติ, สตฺถา จ โหติ เถโร รตฺตญฺญู จิรปพฺพชิโต อทฺธคโต วโย-อนุปฺปตฺโต, เถรา จสฺส ภิกฺขู สาวกา โหนฺติ วิยตฺตา วินีตา วิสารทา ปตฺตโยคกฺเขมา, อลํ สมกฺขาตุํ สทฺธมฺมสฺส, อลํ อุปฺปนฺนํ ปรปฺปวาทํ สหธมฺเมหิ สุนิคฺคหิตํ นิคฺคเหตฺวา สปฺปาฏิหาริยํ ธมฺมํ เทเสตุํ.\csromanspacer{} มชฺฌิมา จสฺส ภิกฺขู สาวกา โหนฺติ.\csromanspacer{} นวา จสฺส ภิกฺขู สาวกา โหนฺติ.\csromanspacer{} เถรา จสฺส ภิกฺขุนิโย สาวิกา โหนฺติ.\csromanspacer{} มชฺฌิมา จสฺส ภิกฺขุนิโย สาวิกา โหนฺติ.\csromanspacer{} นวา จสฺส ภิกฺขุนิโย สาวิกา โหนฺติ.\csromanspacer{} อุปาสกา จสฺส สาวกา โหนฺติ คิหี โอทาตวสนา พฺรหฺมจาริโน.\csromanspacer{} อุปาสกา จสฺส สาวกา โหนฺติ คิหี โอทาตวสนา กามโภคิโน.\csromanspacer{} อุปาสิกา จสฺส สาวิกา โหนฺติ คิหินิโย โอทาตวสนา พฺรหฺมจารินิโย.\csromanspacer{} อุปาสิกา จสฺส สาวิกา โหนฺติ คิหินิโย โอทาตวสนา กามโภคินิโย.\csromanspacer{} พฺรหฺมจริยญฺจสฺส โหติ อิทฺธญฺจว ผีตญฺจ วิตฺถาริกํ พาหุชญฺญํ ปุถุภูตํ ยาว เทวมนุสฺเสหิ สุปฺปกาสิตํ, ลาภคฺคปฺปตฺตญฺจ ยสคฺคปฺปตฺตญฺจ.\csromanspacer{} เอวํ ตํ พฺรหฺมจริยํ ปริปูรํ โหติ เตนงฺเคน.}

\proseitem{175}{อหํ โข ปน จุนฺท เอตรหิ สตฺถา โลเก อุปฺปนฺโน อรหํ สมฺมาสมฺพุทฺโธ, ธมฺโม จ สฺวากฺขาโต สุปฺปเวทิโต นิยฺยานิโก อุปสมสํวตฺตนิโก สมฺมาสมฺพุทฺธปฺปเวทิโต, วิญฺญาปิตตฺถา จ เม สาวกา สทฺธมฺเม, เกวลญฺจ เตสํ ปริปูรํ พฺรหฺมจริยํ อาวิกตํ อุตฺตานีกตํ สพฺพสงฺคาหปทกตํ สปฺปาฏิหีรกตํ ยาว เทวมนุสฺเสหิ สุปฺปกาสิตํ.\csromanspacer{} อหํ โข ปน จุนฺท เอตรหิ สตฺถา เถโร รตฺตญฺญู จิรปพฺพชิโต อทฺธคโต วโย-อนุปฺปตฺโต.}

\prose{สนฺติ โข ปน เม จุนฺท เอตรหิ เถรา ภิกฺขู สาวกา โหนฺติ วิยตฺตา วินีตา วิสารทา ปตฺตโยคกฺเขมา, อลํ สมกฺขาตุํ สทฺธมฺมสฺส, อลํ อุปฺปนฺนํ ปรปฺปวาทํ สหธมฺเมหิ สุนิคฺคหิตํ นิคฺคเหตฺวา สปฺปาฏิหาริยํ ธมฺมํ เทเสตุํ.\csromanspacer{} สนฺติ โข ปน เม จุนฺท เอตรหิ มชฺฌิมา ภิกฺขู สาวกา.\csromanspacer{} สนฺติ โข ปน เม จุนฺท เอตรหิ นวา ภิกฺขู สาวกา.\csromanspacer{} สนฺติ \csromanfolio{104}โข ปน เม จุนฺท เอตรหิ เถรา ภิกฺขุนิโย สาวิกา.\csromanspacer{} สนฺติ โข ปน เม จุนฺท เอตรหิ มชฺฌิมา ภิกฺขุนิโย สาวิกา.\csromanspacer{} สนฺติ โข ปน เม จุนฺท เอตรหิ นวา ภิกฺขุนิโย สาวิกา.\csromanspacer{} สนฺติ โข ปน เม จุนฺท เอตรหิ อุปาสกา สาวกา คิหี โอทาตวสนา พฺรหฺมจาริโน.\csromanspacer{} สนฺติ โข ปน เม จุนฺท เอตรหิ อุปาสกา สาวกา คิหี โอทาตวสนา กามโภคิโน.\csromanspacer{} สนฺติ โข ปน เม จุนฺท เอตรหิ อุปาสิกา สาวิกา คิหินิโย โอทาตวสนา พฺรหฺมจารินิโย.\csromanspacer{} สนฺติ โข ปน เม จุนฺท เอตรหิ อุปาสิกา สาวิกา คิหินิโย โอทาตวสนา กามโภคินิโย.\csromanspacer{} เอตรหิ โข ปน เม จุนฺท พฺรหฺมจริยํ อิทฺธญฺเจว ผีตญฺจ วิตฺถาริกํ พาหุชญฺญํ ปุถุภูตํ ยาว เทวมนุสฺเสหิ สุปฺปกาสิตํ.}

\proseitem{176}{ยาวตา โข จุนฺท เอตรหิ สตฺถาโร โลเก อุปฺปนฺนา, นาหํ จุนฺท อญฺญํ เอกสตฺถารมฺปิ สมนุปสฺสามิ เอวํลาภคฺคยสคฺคปฺปตฺตํ ยถริวาหํ.\csromanspacer{} ยาวตา โข ปน จุนฺท เอตรหิ สํโฆ วา คโณ วา โลเก อุปฺปนฺโน, นาหํ จุนฺท อญฺญํ เอกํ สํฆมฺปิ สมนุปสฺสามิ เอวํลาภคฺคยสคฺคปฺปตฺตํ ยถริวายํ จุนฺท ภิกฺขุสํโฆ.\csromanspacer{} ยํ โข ตํ จุนฺท สมฺมา วทมาโน วเทยฺย “สพฺพาการสมฺปนฺนํ สพฺพาการปริปูรํ อนูนมนธิกํ สฺวากฺขาตํ เกวลํ ปริปูรํ พฺรหฺมจริยํ สุปฺปกาสิตนฺ”ติ.\csromanspacer{} อิทเมว ตํ สมฺมา วทมาโน วเทยฺย “สพฺพาการสมฺปนฺนํ ฯเปฯ สุปฺปกาสิตนฺ”ติ.}

\prose{อุทโก\footnote{อุทฺทโก (สี, สฺยา, อิ)} สุทํ จุนฺท รามปุตฺโต เอวํ วาจํ ภาสติ “ปสฺสํ น ปสฺสตี”ติ.\csromanspacer{} กิญฺจ ปสฺสํ น ปสฺสตีติ.\csromanspacer{} ขุรสฺส สาธุนิสิตสฺส ตลมสฺส ปสฺสติ, ธารญฺจ ขฺวสฺส น ปสฺสติ.\csromanspacer{} อิทํ วุจฺจติ จุนฺท “ปสฺสํ น ปสฺสตี”ติ.\csromanspacer{} ยํ โข ปเนตํ จุนฺท อุทเกน รามปุตฺเตน ภาสิตํ หีนํ คมฺมํ โปถุชฺชนิกํ อนริยํ อนตฺถสํหิตํ ขุรเมว สนฺธาย.\csromanspacer{} ยญฺจ ตํ\footnote{ยญฺเจตํ (สฺยา, ก)} จุนฺท สมฺมา วทมาโน วเทยฺย “ปสฺสํ น ปสฺสตี”ติ.\csromanspacer{} อิทเมว ตํ\footnote{อิทเมเวตํ (ก)} สมฺมา วทมาโน วเทยฺย “ปสฺสํ น ปสฺสตี”ติ.\csromanspacer{} กิญฺจ ปสฺสํ น ปสฺสตีติ.\csromanspacer{} เอวํ สพฺพาการสมฺปนฺนํ สพฺพาการปริปูรํ อนูนมนธิกํ สฺวากฺขาตํ เกวลํ ปริปูรํ พฺรหฺมจริยํ สุปฺปกาสิตนฺติ, อิติ เหตํ ปสฺสติ\footnote{สุปฺปกาสิตํ, อิติ เหตํ น ปสฺสตีติ (สฺยา, ก)}.\csromanspacer{} อิทเมตฺถ อปกฑฺเฒยฺย, เอวํ ตํ ปริสุทฺธตรํ อสฺสาติ, อิติ เหตํ น ปสฺสติ\footnote{น ปสฺสตีติ (สฺยา, ก)}.\csromanspacer{} อิทเมตฺถ อุปกฑฺเฒยฺย, เอวํ ตํ ปริปูรํ\footnote{ปริสุทฺธตรํ (สฺยา, ก), ปริปูรตรํ (?)} อสฺสาติ,}

\vspace{\tipitakaparskip}
\csromancenter{ปาสาทิกสุตฺต อิติ เหตํ น ปสฺสติ.\csromanspacer{} อิทํ วุจฺจติ “ปสฺสํ น ปสฺสตี”ติ.\csromanspacer{} ยํ โข ตํ จุนฺท สมฺมา วทมาโน วเทยฺย “สพฺพาการสมฺปนฺนํ ฯเปฯ พฺรหฺมจริยํ สุปฺปกาสิตนฺ”ติ.\csromanspacer{} อิทเมว ตํ สมฺมา วทมาโน วเทยฺย “สพฺพาการสมฺปนฺนํ สพฺพาการปริปูรํ อนูนมนธิกํ สฺวากฺขาตํ เกวลํ ปริปูรํ พฺรหฺมจริยํ สุปฺปกาสิตนฺ”ติ.}
\csromanmatikaanchor{mtk.75}
\csromantocmarkref{section}{สงฺคายิตพฺพธมฺม}{mtk.75}
\bookmark[dest={mtk.75},level=1]{สงฺคายิตพฺพธมฺม}
\csromanheader{1}{\textbf{สงฺคายิตพฺพธมฺม}}
\proseitem{177}{\csromanfolio{105}ตสฺมาติห จุนฺท เย โว มยา ธมฺมา อภิญฺญา เทสิตา, ตตฺถ สพฺเพเหว สงฺคมฺม สมาคมฺม อตฺเถน อตฺถํ พฺยญฺชเนน พฺยญฺชนํ สงฺคายิตพฺพํ น วิวทิตพฺพํ, ยถยิทํ พฺรหฺมจริยํ อทฺธนิยํ อสฺส จิรฏฺฐิติกํ, ตทสฺส พหุชนหิตาย พหุชนสุขาย โลกานุกมฺปาย อตฺถาย หิตาย สุขาย เทวมนุสฺสานํ.\csromanspacer{} กตเม จ เต จุนฺท ธมฺมา มยา อภิญฺญา เทสิตา, ยตฺถ สพฺเพเหว สงฺคมฺม สมาคมฺม อตฺเถน อตฺถํ พฺยญฺชเนน พฺยญฺชนํ สงฺคายิตพฺพํ น วิวทิตพฺพํ, ยถยิทํ พฺรหฺมจริยํ อทฺธนิยํ อสฺส จิรฏฺฐิติกํ, ตทสฺส พหุชนหิตาย พหุชนสุขาย โลกานุกมฺปาย อตฺถาย หิตาย สุขาย เทวมนุสฺสานํ.\csromanspacer{} เสยฺยถิทํ, จตฺตาโร สติปฏฺฐานา จตฺตาโร สมฺมปฺปธานา จตฺตาโร อิทฺธิปาทา ปญฺจินฺทฺริยานิ ปญฺจ พลานิ สตฺต โพชฺฌงฺคา อริโย อฏฺฐงฺคิโก มคฺโค.\csromanspacer{} อิเม โข เต จุนฺท ธมฺมา มยา อภิญฺญา เทสิตา, ยตฺถ สพฺเพเหว สงฺคมฺม สมาคมฺม อตฺเถน อตฺถํ พฺยญฺชเนน พฺยญฺชนํ สงฺคายิตพฺพํ น วิวทิตพฺพํ, ยถยิทํ พฺรหฺมจริยํ อทฺธนิยํ อสฺส จิรฏฺฐิติกํ, ตทสฺส พหุชนหิตาย พหุชนสุขาย โลกานุกมฺปาย อตฺถาย หิตาย สุขาย เทวมนุสฺสานํ.}

\csromanmatikaanchor{mtk.76}
\csromantocmarkref{section}{สญฺญาเปตพฺพวิธิ}{mtk.76}
\bookmark[dest={mtk.76},level=1]{สญฺญาเปตพฺพวิธิ}
\csromanheader{1}{\textbf{สญฺญาเปตพฺพวิธิ}}
\proseitem{178}{เตสญฺจ โว จุนฺท สมคฺคานํ สมฺโมทมานานํ อวิวทมานานํ สิกฺขตํ\footnote{สิกฺขิตพฺพํ (พหูสุ)} อญฺญตโร สพฺรหฺมจารี สํเฆ ธมฺมํ ภาเสยฺย.\csromanspacer{} ตตฺร เจ ตุมฺหากํ เอวมสฺส “อยํ โข อายสฺมา อตฺถญฺเจว มิจฺฉา คณฺหาติ, พฺยญฺชนานิ จ มิจฺฉา โรเปตี”ติ.\csromanspacer{} ตสฺส เนว อภินนฺทิตพฺพํ น ปฏิกฺโกสิตพฺพํ, อนภินนฺทิตฺวา อปฺปฏิกฺโกสิตฺวา โส เอวมสฺส วจนีโย “อิมสฺส นุ โข อาวุโส อตฺถสฺส อิมานิ วา พฺยญฺชนานิ เอตานิ วา พฺยญฺชนานิ กตมานิ โอปายิกตรานิ, อิเมสญฺจ\footnote{อิเมสํ วา (สฺยา, อิ, ก), อิเมสํ (สี)} พฺยญฺชนานํ อยํ วา อตฺโถ เอโส วา \csromanfolio{106}อตฺโถ กตโม โอปายิกตโร”ติ.\csromanspacer{} โส เจ เอวํ วเทยฺย “อิมสฺส โข อาวุโส อตฺถสฺส อิมาเนว พฺยญฺชนานิ โอปายิกตรานิ, ยา เจว\footnote{ยญฺเจว (สี, ก), ฏีกา โอโลเกตพฺพา.} เอตานิ, อิเมสญฺจ\footnote{อิเมสํ (สพฺพตฺถ)} พฺยญฺชนานํ อยเมว อตฺโถ โอปายิกตโร, ยา เจว1 เอโสติ.\csromanspacer{} โส เนว อุสฺสาเทตพฺโพ น อปสาเทตพฺโพ, อนุสฺสาเทตฺวา อนปสาเทตฺวา เสฺวว สาธุกํ สญฺญาเปตพฺโพ ตสฺส จ อตฺถสฺส เตสญฺจ พฺยญฺชนานํ นิสนฺติยา.}

\proseitem{179}{อปโรปิ เจ จุนฺท สพฺรหฺมจารี สํเฆ ธมฺมํ ภาเสยฺย.\csromanspacer{} ตตฺร เจ ตุมฺหากํ เอวมสฺส “อยํ โข อายสฺมา อตฺถํ หิ โข มิจฺฉา คณฺหาติ พฺยญฺชนานิ สมฺมา โรเปตี”ติ.\csromanspacer{} ตสฺส เนว อภินนฺทิตพฺพํ น ปฏิกฺโกสิตพฺพํ, อนภินนฺทิตฺวา อปฺปฏิกฺโกสิตฺวา โส เอวมสฺส วจนีโย “อิเมสํ นุ โข อาวุโส พฺยญฺชนานํ อยํ วา อตฺโถ เอโส วา อตฺโถ กตโม โอปายิกตโร”ติ.\csromanspacer{} โส เจ เอวํ วเทยฺย “เตสํ โข อาวุโส พฺยญฺชนานํ อยเมว อตฺโถ โอปายิกตโร, ยา เจว เอโส”ติ.\csromanspacer{} โส เนว อุสฺสาเทตพฺโพ น อปสาเทตพฺโพ, อนุสฺสาเทตฺวา อนปสาเทตฺวา เสฺวว สาธุกํ สญฺญาเปตพฺโพ ตสฺเสว อตฺถสฺส นิสนฺติยา.}

\proseitem{180}{อปโรปิ เจ จุนฺท สพฺรหฺมจารี สํเฆ ธมฺมํ ภาเสยฺย.\csromanspacer{} ตตฺร เจ ตุมฺหากํ เอวมสฺส “อยํ โข อายสฺมา อตฺถํ หิ โข สมฺมา คณฺหาติ พฺยญฺชนานิ มิจฺฉา โรเปตี”ติ.\csromanspacer{} ตสฺส เนว อภินนฺทิตพฺพํ น ปฏิกฺโกสิตพฺพํ, อนภินนฺทิตฺวา อปฺปฏิกฺโกสิตฺวา โส เอวมสฺส วจนีโย “อิมสฺส นุ โข อาวุโส อตฺถสฺส อิมานิ วา พฺยญฺชนานิ เอตานิ วา พฺยญฺชนานิ กตมานิ โอปายิกตรานี”ติ.\csromanspacer{} โส เจ เอวํ วเทยฺย “อิมสฺส โข อาวุโส อตฺถสฺส อิมาเนว พฺยญฺชนานิ โอปายิกตรานิ, ยา เจว เอตานี”ติ.\csromanspacer{} โส เนว อุสฺสาเทตพฺโพ น อปสาเทตพฺโพ, อนุสฺสาเทตฺวา อนปสาเทตฺวา เสฺวว สาธุกํ สญฺญาเปตพฺโพ เตสญฺเญว พฺยญฺชนานํ นิสนฺติยา.}

\proseitem{181}{อปโรปิ เจ จุนฺท สพฺรหฺมจารี สํเฆ ธมฺมํ ภาเสยฺย.\csromanspacer{} ตตฺร เจ ตุมฺหากํ เอวมสฺส “อยํ โข อายสฺมา อตฺถญฺเจว สมฺมา คณฺหาติ}

\vspace{\tipitakaparskip}
\csromancenter{ปาสาทิกสุตฺต พฺยญฺชนานิ จ สมฺมา โรเปตี”ติ.\csromanspacer{} ตสฺส สาธูติ ภาสิตํ อภินนฺทิตพฺพํ อนุโมทิตพฺพํ, ตสฺส สาธูติ ภาสิตํ อภินนฺทิตฺวา อนุโมทิตฺวา โส เอวมาสฺส วจนีโย “ลาภา โน อาวุโส สุลทฺธํ โน อาวุโส, เย มยํ อายสฺมนฺตํ ตาทิสํ สพฺรหฺมจาริํ ปสฺสาม เอวํ อตฺถุเปตํ พฺยญฺชนุเปตนฺ”ติ.}
\csromanmatikaanchor{mtk.77}
\csromantocmarkref{section}{ปจฺจยานุญฺญาตการณ}{mtk.77}
\bookmark[dest={mtk.77},level=1]{ปจฺจยานุญฺญาตการณ}
\csromanheader{1}{\textbf{ปจฺจยานุญฺญาตการณ}}
\proseitem{182}{\csromanfolio{107}น โว อหํ จุนฺท ทิฏฺฐธมฺมิกานํเยว อาสวานํ สํวราย ธมฺมํ เทเสมิ, น ปนาหํ จุนฺท สมฺปรายิกานํเยว อาสวานํ ปฏิฆาตาย ธมฺมํ เทเสมิ.\csromanspacer{} ทิฏฺฐธมฺมิกานญฺเจวาหํ จุนฺท อาสวานํ สํวราย ธมฺมํ เทเสมิ, สมฺปรายิกานญฺจ อาสวานํ ปฏิฆาตาย.\csromanspacer{} ตสฺมาติห จุนฺท ยํ โว มยา จีวรํ อนุญฺญาตํ, อลํ โว ตํ ยาวเทว สีตสฺส ปฏิฆาตาย, อุณฺหสฺส ปฏิฆาตาย, ฑํส มกส วาตาตป สรีสป\footnote{สิริํสป (สฺยา)} สมฺผสฺสานํ ปฏิฆาตาย, ยาวเทว หิริโกปีน ปฏิจฺฉาทนตฺถํ.\csromanspacer{} โย โว มยา ปิณฺฑปาโต อนุญฺญาโต, อลํ โว โส ยาวเทว อิมสฺส กายสฺส ฐิติยา ยาปนาย วิหิํสูปรติยา พฺรหฺมจริยานุคฺคหาย, อิติ ปุราณญฺจ เวทนํ ปฏิหงฺขามิ, นวญฺจ เวทนํ น อุปฺปาเทสฺสามิ, ยาตฺรา จ เม ภวิสฺสติ อนวชฺชตา จ ผาสุวิหาโร จ\footnote{จาติ (พหูสุ) 3. อพฺยาปชฺฌปรมตายาติ (สี, สฺยา, อิ), อพฺยาพชฺฌปรมตาย (?)}.\csromanspacer{} ยํ โว มยา เสนาสนํ อนุญฺญาตํ, อลํ โว ตํ ยาวเทว สีตสฺส ปฏิฆาตาย, อุณฺหสฺส ปฏิฆาตาย, ฑํสมกสวาตาตปสรีสปสมฺผสฺสานํ ปฏิฆาตาย ยาวเทว อุตุปริสฺสยวิโนทน ปฏิสลฺลานารามตฺถํ.\csromanspacer{} โย โว มยา คิลานปจฺจยเภสชฺช ปริกฺขาโร อนุญฺญาโต, อลํ โว โส ยาวเทว อุปฺปนฺนานํ เวยฺยาพาธิกานํ เวทนานํ ปฏิทนานํ ปฏิฆาตาย อพฺยาปชฺชปรมตาย3.}

\csromanmatikaanchor{mtk.78}
\csromantocmarkref{section}{สุขลฺลิกานุโยค}{mtk.78}
\bookmark[dest={mtk.78},level=1]{สุขลฺลิกานุโยค}
\csromanheader{1}{\textbf{สุขลฺลิกานุโยค}}
\proseitem{183}{ฐานํ โข ปเนตํ จุนฺท วิชฺชติ ยํ อญฺญติตฺถิยา ปริพฺพาชกา เอวํ วเทยฺยุํ “สุขลฺลิกานุโยคมนุยุตฺตา สมณา สกฺยปุตฺติยา วิหรนฺตี”ติ.\csromanspacer{} เอวํวาทิโน\footnote{วทมานา (สฺยา)} จุนฺท อญฺญติตฺถิยา ปริพฺพาชกา เอวมสฺสุ วจนียา “กตโม \csromanfolio{108}โส อาวุโส สุขลฺลิกานุโยโค.\csromanspacer{} สุขลฺลิกานุโยคา หิ พหู อเนกวิหิตา นานปฺปการกา”ติ.}

\prose{จตฺตาโรเม จุนฺท สุขลฺลิกานุโยคา หีนา คมฺมา โปถุชฺชนิกา อนริยา อนตฺถสํหิตา น นิพฺพิทาย น วิราคาย น นิโรธาย น อุปสมาย น อภิญฺญาย น สมฺโพธาย น นิพฺพานาย สํวตฺตนฺติ.\csromanspacer{} กตเม จตฺตาโร.}

\prose{อิธ จุนฺท เอกจฺโจ พาโล ปาเณ วธิตฺวา วธิตฺวา อตฺตานํ สุเขติ ปีเณติ, อยํ ปฐโม สุขลฺลิกานุโยโค.}

\prose{ปุน จปรํ จุนฺท อิเธกจฺโจ อทินฺนํ อาทิยิตฺวา อาทิยิตฺวา อตฺตานํ สุเขติ ปีเณติ, อยํ ทุติโย สุขลฺลิกานุโยโค.}

\prose{ปุน จปรํ จุนฺท อิเธกจฺโจ มุสา ภณิตฺวา ภณิตฺวา อตฺตานํ สุเขติ ปีเณติ, อยํ ตติโย สุขลฺลิกานุโยโค.}

\prose{ปุน จปรํ จุนฺท อิเธกจฺโจ ปญฺจหิ กามคุเณหิ สมปฺปิโต สมงฺคีภูโต ปริจาเรติ, อยํ จตุตฺโถ สุขลฺลิกานุโยโค.}

\prose{อิเม โข จุนฺท จตฺตาโร สุขลฺลิกานุโยคา หีนา คมฺมา โปถุชฺชนิกา อนริยา อนตฺถสํหิตา น นิพฺพิทาย น วิราคาย น นิโรธาย น อุปสมาย น อภิญฺญาย น สมฺโพธาย น นิพฺพานาย สํวตฺตนฺติ.}

\prose{ฐานํ โข ปเนตํ จุนฺท วิชฺชติ ยํ อญฺญติตฺถิยา ปริพฺพาชกา เอวํ วเทยฺยุํ “อิเม จตฺตาโร สุขลฺลิกานุโยเค อนุยุตฺตา สมณา สกฺยปุตฺติยา วิหรนฺตี”ติ.\csromanspacer{} เต โว\footnote{เต (สี, อิ)} “มาเหวนฺ”ติสฺสุ วจนียา.\csromanspacer{} น เต โว สมฺมา วทมานา วเทยฺยุํ, อพฺภาจิกฺเขยฺยุํ อสตา อภูเตน.}

\proseitem{184}{จตฺตาโรเม จุนฺท สุขลฺลิกานุโยคา เอกนฺตนิพฺพิทาย วิราคาย นิโรธาย อุปสมาย อภิญฺญาย สมฺโพธาย นิพฺพานาย สํวตฺตนฺติ.\csromanspacer{} กตเม จตฺตาโร.}

\prose{อิธ จุนฺท ภิกฺขุ วิวิจฺเจว กาเมหิ วิวิจฺจ อกุสเลหิ ธมฺเมหิ สวิตกฺกํ สวิจารํ วิเวกชํ ปีติสุขํ ปฐมํ ฌานํ อุปสมฺปชฺช วิหรติ, อยํ ปฐโม สุขลฺลิกานุโยโค.}

\csromanheader{2}{6. ปาสาทิกสุตฺต}
\prose{\csromanfolio{109}ปุน จปรํ จุนฺท ภิกฺขุ วิตกฺกวิจารานํ วูปสมา ทุติยํ ฌานํ อุปสมฺปชฺช วิหรติ, อยํ ทุติโย สุขลฺลิกานุโยโค.}

\prose{ปุน จปรํ จุนฺท ภิกฺขุ ปีติยา จ วิราคา ตติยํ ฌานํ อุปสมฺปชฺช วิหรติ, อยํ ตติโย สุขลฺลิกานุโยโค.}

\prose{ปุน จปรํ จุนฺท ภิกฺขุ สุขสฺส จ ปหานา ทุกฺขสฺส จ ปหานา จตุตฺถํ ฌานํ อุปสมฺปชฺช วิหรติ, อยํ จตุตฺโถ สุขลฺลิกานุโยโค.}

\prose{อิเม โข จุนฺท จตฺตาโร สุขลฺลิกานุโยคา เอกนฺตนิพฺพิทาย วิราคาย นิโรธาย อุปสมาย อภิญฺญาย สมฺโพธาย นิพฺพานาย สํวตฺตนฺติ.}

\prose{ฐานํ โข ปเนตํ จุนฺท วิชฺชติ ยํ อญฺญติตฺถิยา ปริพฺพาชกา เอวํ วเทยฺยุํ “อิเม จตฺตาโร สุขลฺลิกานุโยเค อนุยุตฺตา สมณา สกฺยปุตฺติยา วิหรนฺตี”ติ.\csromanspacer{} เต โว “เอวนฺ”ติสฺสุ วจนียา.\csromanspacer{} สมฺมา เต โว วทมานา วเทยฺยุํ, น เต โว อพฺภาจิกฺเขยฺยุํ อสตา อภูเตน.}

\csromanmatikaanchor{mtk.79}
\csromantocmarkref{section}{สุขลฺลิกานุโยคานิสํส}{mtk.79}
\bookmark[dest={mtk.79},level=1]{สุขลฺลิกานุโยคานิสํส}
\csromanheader{1}{\textbf{สุขลฺลิกานุโยคานิสํส}}
\proseitem{185}{ฐานํ โข ปเนตํ จุนฺท วิชฺชติ, ยํ อญฺญติตฺถิยา ปริพฺพาชกา เอวํ วเทยฺยุํ “อิเม ปนาวุโส จตฺตาโร สุขลฺลิกานุโยเค อนุยุตฺตานํ วิหรตํ กติ ผลานิ กตานิสํสา ปาฏิกงฺขา”ติ.\csromanspacer{} เอวํวาทิโน จุนฺท อญฺญติตฺถิยา ปริพฺพาชกา เอวมสฺสุ วจนียา “อิเม โข อาวุโส จตฺตาโร สุขลฺลิกานุโยเค อนุยุตฺตานํ วิหรตํ จตฺตาริ ผลานิ จตฺตาโร อานิสํสา ปาฏิกงฺขา.\csromanspacer{} กตเม จตฺตาโร.\csromanspacer{} อิธาวุโส ภิกฺขุ ติณฺณํ สํโยชนานํ ปริกฺขยา โสตาปนฺโน โหติ อวินิปาตธมฺโม นิยโต สมฺโพธิปรายโณ, อิทํ ปฐมํ ผลํ ปฐโม อานิสํโส.\csromanspacer{} ปุน จปรํ อาวุโส ภิกฺขุ ติณฺณํ สํโยชนานํ ปริกฺขยา ราคโทสโมหานํ ตนุตฺตา สกทาคามี โหติ สกิเทว อิมํ โลกํ อาคนฺตฺวา ทุกฺขสฺสนฺตํ กโรติ, อิทํ ทุติยํ ผลํ ทุติโย อานิสํโส.\csromanspacer{} ปุน จปรํ อาวุโส ภิกฺขุ ปญฺจนฺนํ โอรมฺภาคิยานํ สํโยชนานํ ปริกฺขยา โอปปาติโก โหติ ตตฺถ ปรินิพฺพายี อนาวตฺติธมฺโม ตสฺมา โลกา, อิทํ ตติยํ ผลํ ตติโย อานิสํโส.\csromanspacer{} ปุน จปรํ อาวุโส ภิกฺขุ อาสวานํ ขยา อนาสวํ เจโตวิมุตฺติํ ปญฺญาวิมุตฺติํ ทิฏฺเฐว ธมฺเม สยํ อภิญฺญา สจฺฉิกตฺวา อุปสมฺปชฺช วิหรติ, อิทํ จตุตฺถํ ผลํ จตุตฺโถ อานิสํโส. \csromanfolio{110}อิเม โข อาวุโส จตฺตาโร สุขลฺลิกานุโยเค อนุยุตฺตานํ วิหรตํ อิมานิ จตฺตาริ ผลานิ จตฺตาโร อานิสํสา ปาฏิกงฺขา”ติ.}

\csromanmatikaanchor{mtk.80}
\csromantocmarkref{section}{ขีณาสว-อภพฺพฐาน}{mtk.80}
\bookmark[dest={mtk.80},level=1]{ขีณาสว-อภพฺพฐาน}
\csromanheader{1}{\textbf{ขีณาสว-อภพฺพฐาน}}
\proseitem{186}{ฐานํ โข ปเนตํ จุนฺท วิชฺชติ, ยํ อญฺญติตฺถิยา ปริพฺพาชกา เอวํ วเทยฺยุํ “อฏฺฐิตธมฺมา สมณา สกฺยปุตฺติยา วิหรนฺตี”ติ.\csromanspacer{} เอวํวาทิโน จุนฺท อญฺญติตฺถิยา ปริพฺพาชกา เอวมสฺสุ วจนียา “อตฺถิ โข อาวุโส เตน ภควตา ชานตา ปสฺสาตา อรหตา สมฺมาสมฺพุทฺเธน สาวกานํ ธมฺมา เทสิตา ปญฺญตฺตา ยาวชีวํ อนติกฺกมนียา.\csromanspacer{} เสยฺยถาปิ อาวุโส อินฺทขีโล วา อโยขีโล วา คมฺภีรเนโม สุนิขาโต อจโล อสมฺปเวธี.\csromanspacer{} เอวเมว โข อาวุโส เตน ภควตา ชานตา ปสฺสตา อรหตา สมฺมาสมฺพุทฺเธน สาวกานํ ธมฺมา เทสิตา ปญฺญตฺตา ยาวชีวํ อนติกฺกมนียา.\csromanspacer{} โย โส อาวุโส ภิกฺขุ อรหํ ขีณาสโว วุสิตวา กตกรณีโย โอหิตภาโร อนุปฺปตฺตสทตฺโถ ปริกฺขีณภวสํโยชโน สมฺมทญฺญา วิมุตฺโต, อภพฺโพ โส นว ฐานานิ อชฺฌาจริตุํ.\csromanspacer{} อภพฺโพ อาวุโส ขีณาสโว ภิกฺขุ สญฺจิจฺจ ปาณํ ชีวิตา โวโรเปตุํ, อภพฺโพ ขีณาสโว ภิกฺขุ อทินฺนํ เถยฺยสงฺขาตํ อาทิยิตุํ, อภพฺโพ ขีณาสโว ภิกฺขุ เมถุนํ ธมฺมํ ปฏิเสวิตุํ, อภพฺโพ ขีณาสโว ภิกฺขุ สมฺปชานมุสา ภาสิตุํ, อภพฺโพ ขีณาสโว ภิกฺขุ สนฺนิธิการกํ กาเม ปริภุญฺชิตุํ เสยฺยถาปิ ปุพฺเพ อาคาริกภูโต, อภพฺโพ ขีณาสโว ภิกฺขุ ฉนฺทาคติํ คนฺตุํ, อภพฺโพ ขีณาสโว ภิกฺขุ โทสาคติํ คนฺตุํ, อภพฺโพ ขีณาสโว ภิกฺขุ โมหาคติํ คนฺตุํ, อภพฺโพ ขีณาสโว ภิกฺขุ ภยาคติํ คนฺตุํ.\csromanspacer{} โย โส อาวุโส ภิกฺขุ อรหํ ขีณาสโว วุสิตวา กตกรณีโย โอหิตภาโร อนุปฺปตฺตสทตฺโถ ปริกฺขีณภวสํโยชโน สมฺมทญฺญา วิมุตฺโต, อภพฺโพ โส อิมานิ นว ฐานานิ อชฺฌาจริตุนฺ”ติ.}

\csromanmatikaanchor{mtk.81}
\csromantocmarkref{section}{ปญฺหาพฺยากรณ}{mtk.81}
\bookmark[dest={mtk.81},level=1]{ปญฺหาพฺยากรณ}
\csromanheader{1}{\textbf{ปญฺหาพฺยากรณ}}
\proseitem{187}{ฐานํ โข ปเนตํ จุนฺท วิชฺชติ, ยํ อญฺญติตฺถิยา ปริพฺพาชกา เอวํ วเทยฺยุํ “อตีตํ โข อทฺธานํ อารพฺภ สมโณ โคตโม อติเรกํ ญาณทสฺสนํ ปญฺญเปติ, โน จ โข อนาคตํ อทฺธานํ อารพฺภ อติเรกํ ญาณทสฺสนํ ปญฺญเปติ, ตยิทํ กิํสุ ตยิทํ กถํสู”ติ.\csromanspacer{} เต จ}

\vspace{\tipitakaparskip}
\csromancenter{ปาสาทิกสุตฺต อญฺญติตฺถิยา ปริพฺพาชกา อญฺญวิหิตเกน ญาณทสฺสเนน อญฺญวิหิตกํ ญาณทสฺสนํ ปญฺญเปตพฺพํ มญฺญนฺติ ยถริว พาลา อพฺยตฺตา.\csromanspacer{} อตีตํ โข จุนฺท อทฺธานํ อารพฺภ ตถาคตสฺส สตานุสาริ ญาณํ โหติ, โส ยาวตกํ อากงฺขติ ตาวตกํ อนุสฺสรติ.\csromanspacer{} อนาคตญฺจ โข อทฺธานํ อารพฺภ ตถาคตสฺส โพธิชํ ญาณํ อุปฺปชฺชติ “อยมนฺติมา ชาติ, นตฺถิทานิ ปุนพฺภโว”ติ.\csromanspacer{} อตีตํ เจปิ โข จุนฺท โหติ อภูตํ อตจฺฉํ อนตฺถสํหิตํ, น ตํ ตถาคโต พฺยากโรติ.\csromanspacer{} อตีตํ เจปิ จุนฺท โหติ ภูตํ ตจฺฉํ อนตฺถสํหิตํ, ตมฺปิ ตถาคโต น พฺยากโรติ.\csromanspacer{} อตีตํ เจปิ โข จุนฺท โหติ ภูตํ ตจฺฉํ อตฺถสํหิตํ, ตตฺร กาลญฺญู ตถาคโต โหติ ตสฺส ปญฺหสฺส เวยฺยากรณาย.\csromanspacer{} อนาคตํ เจปิ จุนฺท โหติ อภูตํ อตจฺฉํ อนตฺถสํหิตํ น ตํ ตถาคโต พฺยากโรติ ฯเปฯ ตสฺส ปญฺหสฺส เวยฺยากรณาย.\csromanspacer{} ปจฺจุปฺปนฺนํ เจปิ จุนฺท โหติ อภูตํ อตจฺฉํ อนตฺถสํหิตํ, น ตํ ตถาคโต พฺยากโรติ.\csromanspacer{} ปจฺจุปฺปนฺนํ เจปิ จุนฺท โหติ ภูตํ ตจฺฉํ อนตฺถสํหิตํ, ตมฺปิ ตถาคโต น พฺยากโรติ.\csromanspacer{} ปจฺจุปฺปนฺนํ เจปิ จุนฺท โหติ ภูตํ ตจฺฉํ อตฺถสํหิตํ, ตตฺร กาลญฺญู ตถาคโต โหติ ตสฺส ปญฺหสฺส เวยฺยากรณาย.}
\proseitem{188}{\csromanfolio{111}อิติ โข จุนฺท อตีตานาคตปจฺจุปฺปนฺเนสุ ธมฺเมสุ ตถาคโต กาลวาที\footnote{กาลวาที สจฺจวาที (สฺยา)} ภูตวาที อตฺถวาที ธมฺมวาที วินยวาที, ตสฺมา “ตถาคโต”ติ วุจฺจติ.\csromanspacer{} ยญฺจ โข จุนฺท สเทวกสฺส โลกสฺส สมารกสฺส สพฺรหฺมกสฺส สสฺสมณพฺราหฺมณิยา ปชาย สเทวมนุสฺสาย ทิฏฺฐํ สุตํ มุตํ วิญฺญาตํ ปตฺตํ ปริเยสิตํ อนุวิจริตํ มนสา, สพฺพํ ตถาคเตน อภิสมฺพุทฺธํ, ตสฺมา “ตถาคโต”ติ วุจฺจติ.\csromanspacer{} ยญฺจ จุนฺท รตฺติํ ตถาคโต อนุตฺตรํ สมฺมาสมฺโพธิํ อภิสมฺพุชฺฌติ, ยญฺจ รตฺติํ อนุปาทิเสสาย นิพฺพานธาตุยา ปรินิพฺพายติ, ยํ เอตสฺมิํ อนฺตเร ภาสติ ลปติ นิทฺทิสติ.\csromanspacer{} สพฺพํ ตํ ตเถว โหติ โน อญฺญถา, ตสฺมา “ตถาคโต”ติ วุจฺจติ.\csromanspacer{} ยถาวาที จุนฺท ตถาคโต ตถาการี, ยถาการี ตถาวาที.\csromanspacer{} อิติ ยถาวาที ตถาการี, ยถาการี ตถาวาที, ตสฺมา “ตถาคโต”ติ วุจฺจติ.\csromanspacer{} สเทวเก โลเก จุนฺท สมารเก สพฺรหฺมเก สสฺสมณพฺราหฺมณิยา ปชาย สเทวมนุสฺสาย ตถาคโต อภิภู อนภิภูโต อญฺญทตฺถุทโส วสวตฺตี, ตสฺมา “ตถาคโต”ติ วุจฺจติ.}

\csromanmatikaanchor{mtk.82}
\csromantocmarkref{section}{อพฺยากตฏฺฐาน}{mtk.82}
\bookmark[dest={mtk.82},level=1]{อพฺยากตฏฺฐาน}
\csromanheader{1}{\textbf{อพฺยากตฏฺฐาน}}
\proseitem{189}{\csromanfolio{112}ฐานํ โข ปเนตํ จุนฺท วิชฺชติ, ยํ อญฺญติตฺถิยา ปริพฺพาชกา เอวํ วเทยฺยุํ “กิํ นุ โข อาวุโส โหติ ตถาคโต ปรํ มรณา, อิทเมว สจฺจํ โมฆมญฺญนฺ”ติ.\csromanspacer{} เอวํวาทิโน จุนฺท อญฺญติตฺถิยา ปริพฺพาชกา เอวมสฺสุ วจนียา “อพฺยากตํ โข อาวุโส ภควตา ‘โหติ ตถาคโต ปรํ มรณา, อิทเมว สจฺจํ โมฆมญฺญนฺ’ติ”.}

\prose{ฐานํ โข ปเนตํ จุนฺท วิชฺชติ, ยํ อญฺญติตฺถิยา ปริพฺพาชกา เอวํ วเทยฺยุํ “กิํ ปนาวุโส น โหติ ตถาคโต ปรํ มรณา, อิทเมว สจฺจํ โมฆมญฺญนฺ”ติ.\csromanspacer{} เอวํวาทิโน จุนฺท อญฺญติตฺถิยา ปริพฺพาชกา เอวมสฺสุ วจนียา “เอวมฺปิ โข อาวุโส ภควตา อพฺยากตํ ‘น โหติ ตถาคโต ปรํ มรณา, อิทเมว สจฺจํ โมฆมญฺญนฺ’ติ”.}

\prose{ฐานํ โข ปเนตํ จุนฺท วิชฺชติ, ยํ อญฺญติตฺถิยา ปริพฺพาชกา เอวํ วเทยฺยุํ “กิํ ปนาวุโส โหติ จ น จ โหติ ตถาคโต ปรํ มรณา, อิทเมว สจฺจํ โมฆมญฺญนฺ”ติ.\csromanspacer{} เอวํวาทิโน จุนฺท อญฺญติตฺถิยา ปริพฺพาชกา เอวมสฺสุ วจนียา “อพฺยากตํ โข ปเนตํ อาวุโส ภควตา ‘โหติ จ น จ โหติ ตถาคโต ปรํ มรณา, อิทเมว สจฺจํ โมฆมญฺญนฺ’ติ”.}

\prose{ฐานํ โข ปเนตํ จุนฺท วิชฺชติ, ยํ อญฺญติตฺถิยา ปริพฺพาชกา เอวํ วเทยฺยุํ “กิํ ปนาวุโส เนว โหติ น น โหติ ตถาคโต ปรํ มรณา, อิทเมว สจฺจํ โมฆมญฺญนฺ”ติ.\csromanspacer{} เอวํวาทิโน จุนฺท อญฺญติตฺถิยา ปริพฺพาชกา เอวมสฺสุ วจนียา “เอวมฺปิ โข อาวุโส ภควตา อพฺยากตํ ‘เนว โหติ น น โหติ ตถาคโต ปรํ มรณา, อิทเมว สจฺจํ โมฆมญฺญนฺ’ติ”.}

\prose{ฐานํ โข ปเนตํ จุนฺท วิชฺชติ, ยํ อญฺญติตฺถิยา ปริพฺพาชกา เอวํ วเทยฺยุํ “กสฺมา ปเนตํ อาวุโส สมเณน โคตเมน อพฺยากตนฺ”ติ.\csromanspacer{} เอวํวาทิโน จุนฺท อญฺญติตฺถิยา ปริพฺพาชกา เอวมสฺสุ วจนียา “น เหตํ อาวุโส อตฺถสํหิตํ น จ ธมฺมสํหิตํ น อาทิพฺรหฺมจริยกํ น นิพฺพิทาย น วิราคาย น นิโรธาย น อุปสมาย น อภิญฺญาย น สมฺโพธาย น นิพฺพานาย สํวตฺตติ, ตสฺมา ตํ ภควตา อพฺยากตนฺ”ติ.}

\csromanmatikaanchor{mtk.83}
\csromantocmarkref{section}{พฺยากตฏฺฐาน}{mtk.83}
\bookmark[dest={mtk.83},level=1]{พฺยากตฏฺฐาน}
\csromanheader{1}{6. ปาสาทิกสุตฺต \textbf{พฺยากตฏฺฐาน}}
\proseitem{190}{\csromanfolio{113}ฐานํ โข ปเนตํ จุนฺท วิชฺชติ, ยํ อญฺญติตฺถิยา ปริพฺพาชกา เอวํ วเทยฺยุํ “กิํ ปนาวุโส สมเณน โคตเมน พฺยากตนฺ”ติ.\csromanspacer{} เอวํวาทิโน จุนฺท อญฺญติตฺถิยา ปริพฺพาชกา เอวมสฺสุ วจนียา “อิทํ ทุกฺขนฺติ โข อาวุโส ภควตา พฺยากตํ, อยํ ทุกฺขสมุทโยติ โข อาวุโส ภควตา พฺยากตํ,อยํ ทุกฺขนิโรโธติ โข อาวุโส ภควตา พฺยากตํ, อยํ ทุกฺขนิโรธคามินี ปฏิปทาติ โข อาวุโส ภควตา พฺยากตนฺ”ติ.}

\prose{ฐานํ โข ปเนตํ จุนฺท วิชฺชติ, ยํ อญฺญติตฺถิยา ปริพฺพาชกา เอวํ วเทยฺยุํ “กสฺมา ปเนตํ อาวุโส สมเณน โคตเมน พฺยากตนฺ”ติ.\csromanspacer{} เอวํวาทิโน จุนฺท อญฺญติตฺถิยา ปริพฺพาชกา เอวมสฺสุ วจนียา “เอตญฺหิ อาวุโส อตฺถสํหิตํ เอตํ ธมฺมสํหิตํ เอตํ อาทิพฺรหฺมจริยกํ เอกนฺตนิพฺพิทาย วิราคาย นิโรธาย อุปสมาย อภิญฺญาย สมฺโพธาย นิพฺพานาย สํวตฺตติ, ตสฺมา ตํ ภควตา พฺยากตนฺ”ติ.}

\csromanmatikaanchor{mtk.84}
\csromantocmarkref{section}{ปุพฺพนฺตสหคตทิฏฺฐินิสฺสย}{mtk.84}
\bookmark[dest={mtk.84},level=1]{ปุพฺพนฺตสหคตทิฏฺฐินิสฺสย}
\csromanheader{1}{\textbf{ปุพฺพนฺตสหคตทิฏฺฐินิสฺสย}}
\proseitem{191}{เยปิ เต จุนฺท ปุพฺพนฺตสหคตา ทิฏฺฐินิสฺสายา, เตปิ โว มยา พฺยากตา, ยถา เต พฺยากาตพฺพา, ยถา จ เต น พฺยากาตพฺพา.\csromanspacer{} กิํ โว อหํ เต ตถา\footnote{ตตฺถ (สฺยา, ก)} พฺยากริสฺสามิ.\csromanspacer{} เยปิ เต จุนฺท อปรนฺตสหคตา ทิฏฺฐินิสฺสยา, เตปิ โว มยา พฺยากตา, ยถา เต พฺยากาตพฺพา, ยถา จ เต น พฺยากาตพฺพา.\csromanspacer{} กิํ โว อหํ เต น ตถา พฺยากริสฺสามิ.\csromanspacer{} กตเม จ เต จุนฺท ปุพฺพนฺตสหคตา ทิฏฺฐินิสฺสยา, เย โว มยา พฺยากตา ยถา เต พฺยากาตพฺพา ยถา จ เต น พฺยากาตพฺพา.\csromanspacer{} สนฺติ โข จุนฺท เอเก สมณพฺราหฺมณา เอวํวาทิโน เอวํทิฏฺฐิโน “สสฺสโต อตฺตา จ โลโก จ, อิทเมว สจฺจํ โมฆมญฺญนฺ”ติ.\csromanspacer{} สนฺติ ปน จุนฺท เอเก สมณพฺราหฺมณา เอวํวาทิโน เอวํทิฏฺฐิโน “อสสฺสโต อตฺตา จ โลโก จ ฯเปฯ.\csromanspacer{} สสฺสโต จ อสสฺสโต จ อตฺตา จ โลโก จ.\csromanspacer{} เนว สสฺสโต นาสสฺสโต อตฺตา จ โลโก จ.\csromanspacer{} สยํกโต อตฺตา จ โลโก จ.\csromanspacer{} ปรํกโต อตฺตา จ โลโก จ.\csromanspacer{} สยํกโต \csromanfolio{114}จ ปรํ กโต จ อตฺตา จ โลโก จ.\csromanspacer{} อสยํกาโร อปรํกาโร อธิจฺจสมุปฺปนฺโน อตฺตา จ โลโก จ, อิทเมว สจฺจํ โมฆมญฺญนฺ”ติ.\csromanspacer{} สสฺสตํ สุขทุกฺขํ.\csromanspacer{} อสสฺสตํ สุขทุกฺขํ.\csromanspacer{} สสฺสตญฺจ อสสฺสตญฺจ สุขทุกฺขํ.\csromanspacer{} เนวสสฺสตํ นาสสฺสตํ สุขทุกฺขํ.\csromanspacer{} สยํกตํ สุขทุกฺขํ.\csromanspacer{} ปรํกตํ สุขทุกฺขํ.\csromanspacer{} สยํกตญฺจ ปรํกตญฺจ สุขทุกฺขํ.\csromanspacer{} อสยํการํ อปรํการํ อธิจฺจสมุปฺปนฺนํ สุขทุกฺขํ, อิทเมว สจฺจํ โมฆมญฺญนฺ”ติ.}

\proseitem{192}{ตตฺร จุนฺท เย เต สมณพฺราหฺมณา เอวํวาทิโน เอวํทิฏฺฐิโน “สสฺสโต อตฺตา จ โลโก จ, อิทเมว สจฺจํ โมฆมญฺญนฺ”ติ.\csromanspacer{} ตฺยาหํ อุปสงฺกมิตฺวา เอวํ วทามิ “อตฺถิ นุ โข อิทํ อาวุโส, วุจฺจติ ‘สสฺสโต อตฺตา จ โลโก จา’ติ”.\csromanspacer{} ยญฺจ โข เต “เอวมาหํสุ อิทเมว สจฺจํ โมฆมญฺญนฺ”ติ.\csromanspacer{} ตํ เตสํ นานุชานามิ.\csromanspacer{} ตํ กิสฺส เหตุ, อญฺญถาสญฺญิโนปิ เหตฺถ จุนฺท สนฺเตเก สตฺตา.\csromanspacer{} อิมายปิ โข อหํ จุนฺท ปญฺญตฺติยา เนว อตฺตนา สมสมํ สมนุปสฺสามิ กุโต ภิยฺโย, อถ โข อหเมว ตตฺถ ภิยฺโย ยทิทํ อธิปญฺญตฺติ.}

\proseitem{193}{ตตฺร จุนฺท เย เต สมณพฺราหฺมณา เอวํวาทิโน เอวํทิฏฺฐิโน “สสฺสโต อตฺตา จ โลโก จ.\csromanspacer{} อสสฺสโต อตฺตา จ โลโก จ.\csromanspacer{} สสฺสโต จ อสสฺสโต จ อตฺตา จ โลโก จ.\csromanspacer{} เนวสสฺสโต นาสสฺสโต อตฺตา จ โลโก จ.\csromanspacer{} สยํกโต อตฺตา จ โลโก จ.\csromanspacer{} ปรํกโต อตฺตา จ โลโก จ.\csromanspacer{} สยํกโต จ ปรํกโต จ อตฺตา จ โลโก จ.\csromanspacer{} อสยํกาโร อปรํกาโร อธิจฺจสมุปฺปนฺโน อตฺตา จ โลโก จ.\csromanspacer{} สสฺสตํ สุขทุกฺขํ.\csromanspacer{} อสสฺสตํ สุขทุกฺขํ.\csromanspacer{} สสฺสตญฺจ อสสฺสตญฺจ สุขทุกฺขํ.\csromanspacer{} เนวสสฺสตํ นาสสฺสตํ สุขทุกฺขํ.\csromanspacer{} สยํกตํ สุขทุกฺขํ.\csromanspacer{} ปรํกตํ สุขทุกฺขํ.\csromanspacer{} สยํกตญฺจ ปรํกตญฺจ สุข-ทุกฺขํ.\csromanspacer{} อสยํการํ อปรํการํ อธิจฺจสมุปฺปนฺนํ สุขทุกฺขํ, อิทเมว สจฺจํ โมฆมญฺญนฺ”ติ.\csromanspacer{} ตฺยาหํ อุปสงฺกมิตฺวา เอวํ วทามิ “อตฺถิ น โข อิทํ อาวุโส, วุจฺจติ ‘อสยํการํ อปรํการํ อธิจฺจสมุปฺปนฺนํ สุขทุกฺขนฺ’ติ”.\csromanspacer{} ยญฺจ โข เต เอวมาหํสุ “อิทเมว สจฺจํ โมฆมญฺญนฺ”ติ.\csromanspacer{} ตํ เตสํ นานุชานามิ.\csromanspacer{} ตํ กิสฺส เหตุ, อญฺญถาสญฺญิโนปิ เหตฺถ จุนฺท สนฺเตเก สตฺตา.\csromanspacer{} อิมายปิ โข อหํ จุนฺท ปญฺญตฺติยา เนว อตฺตนา สมสมํ สมนุปสฺสามิ กุโต ภิยฺโย.\csromanspacer{} อถ โข อหเมว}

\vspace{\tipitakaparskip}
\csromancenter{ปาสาทิกสุตฺต ตตฺถ ภิยฺโย ยทิทํ อธิปญฺญตฺติ.\csromanspacer{} อิเม โข เต จุนฺท ปุพฺพนฺตสหคตา ทิฏฺฐินิสฺสยา, เย โว มยา พฺยากตา, ยถา เต พฺยากาตพฺพา.\csromanspacer{} ยถา จ เต น พฺยากาตพฺพา.\csromanspacer{} กิํ โว อหํ เต น ตถา พฺยากริสฺสามีติ.}
\csromanmatikaanchor{mtk.85}
\csromantocmarkref{section}{อปรนฺตสหคตทิฏฺฐินิสฺสย}{mtk.85}
\bookmark[dest={mtk.85},level=1]{อปรนฺตสหคตทิฏฺฐินิสฺสย}
\csromanheader{1}{\textbf{อปรนฺตสหคตทิฏฺฐินิสฺสย}}
\proseitem{194}{\csromanfolio{115}กตเม จ เต จุนฺท อปรนฺตสหคตา ทิฏฺฐินิสฺสยา, เย โว มยา พฺยากตา, ยถา เต พฺยากาตพฺพา, ยถา จ เต น พฺยากาตพฺพา.\csromanspacer{} กิํ โว อหํ เต น ตถา พฺยากริสฺสามีติ.\csromanspacer{} สนฺติ จุนฺท เอเก สมณพฺราหฺมณา เอวํวาทิโน เอวํทิฏฺฐิโน “รูปี อตฺตา โหติ อโรโค ปรํ มรณา, อิทเมว สจฺจํ โมฆมญฺญนฺ”ติ.\csromanspacer{} สนฺติ ปน จุนฺท เอเก สมณพฺราหฺมณา เอวํวาทิโน เอวํทิฏฺฐิโน “อรูปี อตฺตา โหติ.\csromanspacer{} รูปี จ อรูปี จ อตฺตา โหติ.\csromanspacer{} เนวรูปี นารูปี อตฺตา โหติ.\csromanspacer{} สญฺญี อตฺตา โหติ.\csromanspacer{} อสญฺญี อตฺตา โหติ.\csromanspacer{} เนวสญฺญีนาสญฺญี อตฺตา โหติ.\csromanspacer{} อตฺตา อุจฺฉิชฺชติ วินสฺสติ น โหติ ปรํ มรณา, อิทเมว สจฺจํ โมฆมญฺญนฺ”ติ.\csromanspacer{} ตตฺร จุนฺท เย เต สมณพฺราหฺมณา เอวํวาทิโน เอวํทิฏฺฐิโน “รูปี อตฺตา โหติ อโรโค ปรํ มรณา, อิทเมว สจฺจํ โมฆมญฺญนฺ”ติ.\csromanspacer{} ตฺยาหํ อุปสงฺกมิตฺวา เอวํ วทามิ “อตฺถิ นุ โข อิทํ อาวุโส, วุจฺจติ ‘รูปี อตฺตา โหติ อโรโค ปรํ มรณา’ติ”.\csromanspacer{} ยญฺจ โข เต เอวมาหํสุ “อิทเมว สจฺจํ โมฆมญฺญนฺ”ติ.\csromanspacer{} ตํ เตสํ นานุชานามิ.\csromanspacer{} ตํ กิสฺส เหตุ, อญฺญถาสญฺญิโนปิ เหตฺถ จุนฺท สนฺเตเก สตฺตา.\csromanspacer{} อิมายปิ โข อหํ จุนฺท ปญฺญตฺติยา เนว อตฺตนา สมสมํ สมนุปสฺสามิ กุโต ภิยฺโย, อถ โข อหเมว ตตฺถ ภิยฺโย ยทิทํ อธิปญฺญตฺติ.}

\proseitem{195}{ตตฺร จุนฺท เย เต สมณพฺราหฺมณา เอวํวาทิโน เอวํทิฏฺฐิโน “อรูปี อตฺตา โหติ.\csromanspacer{} รูปี จ อรูปี จ อตฺตา โหติ.\csromanspacer{} เนวรูปีนารูปี อตฺตา โหติ.\csromanspacer{} สญฺญี อตฺตา โหติ.\csromanspacer{} อสญฺญี อตฺตา โหติ.\csromanspacer{} เนวสญฺญีนาสญฺญี อตฺตา โหติ.\csromanspacer{} อตฺตา อุจฺฉิชฺชติ วินสฺสติ น โหติ ปรํ มรณา, อิทเมว สจฺจํ โมฆมญฺญนฺ”ติ.\csromanspacer{} ตฺยาหํ อุปสงฺกมิตฺวา เอวํ วทามิ “อตฺถิ นุ โข อิทํ อาวุโส, วุจฺจติ ‘อตฺตา อุจฺฉิชฺชติ วินสฺสติ น โหติ ปรํ มรณา’ติ”.\csromanspacer{} ยญฺจ โข เต จุนฺท เอวมาหํสุ “อิทเมว สจฺจํ โมฆมญฺญนฺ”ติ.\csromanspacer{} ตํ เตสํ นานุชานามิ.\csromanspacer{} ตํ กิสฺส เหตุ, อญฺญถาสญฺญิโนปิ เหตฺถ จุนฺท สนฺเตเก สตฺตา.\csromanspacer{} อิมายปิ โข อหํ จุนฺท ปญฺญตฺติยา เนว \csromanfolio{116}อตฺตนา สมสมํ สมนุปสฺสามิ กุโต ภิยฺโย, อถ โข อหเมว ตตฺถ ภิยฺโย ยทิทํ อธิปญฺญตฺติ.\csromanspacer{} อิเม โข เต จุนฺท อปรนฺตสหคตา ทิฏฺฐินิสฺสยา, เย โว มยา พฺยากตา, ยถา เต พฺยากาตพฺพา, ยถา จ เต น พฺยากาตพฺพา.\csromanspacer{} กิํ โว อหํ เต น ตถา พฺยากริสฺสามีติ.}

\proseitem{196}{อิเมสญฺจ จุนฺท ปุพฺพนฺตสหคตานํ ทิฏฺฐินิสฺสยานํ อิเมสญฺจ อปรนฺตสหคตานํ ทิฏฺฐินิสฺสยานํ ปหานาย สมติกฺกมาย เอวํ มยา จตฺตาโร สติปฏฺฐานา เทสิตา ปญฺญตฺตา.\csromanspacer{} กตเม จตฺตาโร.\csromanspacer{} อิธ จุนฺท ภิกฺขุ กาเย กายานุปสฺสี วิหรติ อาตาปี สมฺปชาโน สติมา วิเนยฺย โลเก อภิชฺฌาโทมนสฺสํ.\csromanspacer{} เวทนาสุ เวทนานุปสฺสี ฯเปฯ.\csromanspacer{} จิตฺเต จิตฺตานุปสฺสี ฯเปฯ.\csromanspacer{} ธมฺเมสุ ธมฺมานุปสฺสี วิหรติ อาตาปี สมฺปชาโน สติมา วิเนยฺย โลเก อภิชฺฌาโทมนสฺสํ.\csromanspacer{} อิเมสญฺจ จุนฺท ปุพฺพนฺตสหคตานํ ทิฏฺฐินิสฺสยานํ อิเมสญฺจ อปรนฺตสหคตานํ ทิฏฺฐินิสฺสยานํ ปหานาย สมติกฺกมาย เอวํ มยา อิเม จตฺตาโร สติปฏฺฐานา เทสิตา ปญฺญตฺตาติ.}

\proseitem{197}{เตน โข ปน สมเยน อายสฺมา อุปวาโณ ภควโต ปิฏฺฐิโต ฐิโต โหติ ภควนฺตํ พีชยมาโน.\csromanspacer{} อถ โข อายสฺมา อุปวาโณ ภควนฺตํ เอตทโวจ “อจฺฉริยํ ภนฺเต อพฺภุตํ ภนฺเต, ปาสาทิโก วตายํ ภนฺเต ธมฺมปริยาโย, สุปาสาทิโก วตายํ ภนฺเต ธมฺมปริยาโย, โกนามายํ ภนฺเต ธมฺมปริยาโย”ติ.\csromanspacer{} ตสฺมาติห ตฺวํ อุปวาณ อิมํ ธมฺมปริยายํ ปาสาทิโก เตฺวว นํ ธาเรหีติ.\csromanspacer{} อิทมโวจ ภควา.\csromanspacer{} อตฺตมโน อายสฺมา อุปวาโณ ภควโต ภาสิตํ อภินนฺทีติ.}

\nitthitam{\textbf{ปาสาทิกสุตฺตํ นิฏฺฐิตํ ฉฏฺฐํ.}}
\csromanmatikaanchor{mtk.86}
\csromantocmarknopage{chapter}{7. ลกฺขณสุตฺต}{mtk.86}
\bookmark[dest={mtk.86},level=0]{7. ลกฺขณสุตฺต}
\chapterheadpageread{7. ลกฺขณสุตฺต}
\csromanmatikaanchor{mtk.87}
\csromantocmarkref{section}{ทฺวตฺติํสมหาปุริสลกฺขณานิ}{mtk.87}
\bookmark[dest={mtk.87},level=1]{ทฺวตฺติํสมหาปุริสลกฺขณานิ}
\csromanheader{1}{\textbf{ทฺวตฺติํสมหาปุริสลกฺขณานิ}}
\proseitem{198}{\csromanfolio{117}เอวํ เม สุตํ\csromandash{}เอกํ สมยํ ภควา สาวตฺถิยํ วิหรติ เชตวเน อนาถปิณฺฑิกสฺส อาราเม.\csromanspacer{} ตตฺร โข ภควา ภิกฺขู อามนฺเตสิ “ภิกฺขโว”ติ. “ภทฺทนฺเต”ติ\footnote{ภทนฺเตติ (สี, สฺยา, อิ)} เต ภิกฺขู ภควโต ปจฺจสฺโสสุํ.\csromanspacer{} ภควา เอตทโวจ\csromandash{}}

\proseitem{199}{ทฺวตฺติํสิมานิ ภิกฺขเว มหาปุริสสฺส มหาปุริสลกฺขณานิ, เยหิ สมนฺนาคตสฺส มหาปุริสสฺส เทฺวว คติโย ภวนฺติ อนญฺญา.\csromanspacer{} สเจ อคารํ อชฺฌาวสติ, ราชา โหติ จกฺกวตฺตี ธมฺมิโก ธมฺมราชา จาตุรนฺโต วิชิตาวี ชนปทตฺถาวริยปฺปตฺโต สตฺตรตนสมนฺนาคโต.\csromanspacer{} ตสฺสิมานิ สตฺต รตนานิ ภวนฺติ, เสยฺยถิทํ, จกฺกรตนํ หตฺถิรตนํ อสฺสรตนํ มณิรตนํ อิตฺถิรตนํ คหปติรตนํ ปริณายกรตนเมว สตฺตมํ.\csromanspacer{} ปโรสหสฺสํ โข ปนสฺส ปุตฺตา ภวนฺติ สูรา วีรงฺครูปา ปรเสนปฺปมทฺทนา, โส อิมํ ปถวิํ สาครปริยนฺตํ อทณฺเฑน อสตฺเถน ธมฺเมน อภิวิชิย อชฺฌาวสติ.\csromanspacer{} สเจ โข ปน อคารสฺมา อนคาริยํ ปพฺพชติ, อรหํ โหติ สมฺมาสมฺพุทฺโธ โลเก วิวฏฺฏจฺฉโท\footnote{วิวฏจฺฉโท (สฺยา, ก), วิวตฺตจฺฉโท (สี, อิ)}.}

\proseitem{200}{กตมานิ จ ตานิ ภิกฺขเว ทฺวตฺติํส มหาปุริสสฺส มหาปุริสลกฺขณานิ, เยหิ สมนฺนาคตสฺส มหาปุริสสฺส เทฺวว คติโย ภวนฺติ อนญฺญา.\csromanspacer{} สเจ อคารํ อชฺฌาวสติ, ราชา โหติ จกฺกวตฺตี ฯเปฯ.\csromanspacer{} สเจ โข ปน อคารสฺมา อนคาริยํ ปพฺพชติ, อรหํ โหติ สมฺมาสมฺพุทฺโธ โลเก วิวฏฺฏจฺฉโท.}

\prose{อิธ ภิกฺขเว มหาปุริโส สุปฺปติฏฺฐิตปาโท โหติ, ยมฺปิ ภิกฺขเว มหาปุริโส สุปฺปติฏฺฐิตปาโท โหติ, อิทมฺปิ ภิกฺขเว มหาปุริสสฺส มหาปุริสลกฺขณํ ภวติ.}

\prose{ปุน จปรํ ภิกฺขเว มหาปุริสสฺส เหฏฺฐาปาทตเลสุ จกฺกานิ ชาตานิ โหนฺติ สหสฺสารานิ สเนมิกานิ สนาภิกานิ สพฺพาการปริปูรานิ\footnote{สพฺพาการปริปูรานิ สุวิภตฺตนฺตรานิ (สี, อิ)}, \csromanfolio{118}ยมฺปิ ภิกฺขเว มหาปุริสสฺส เหฏฺฐาปาทตเลสุ จกฺกานิ ชาตานิ โหนฺติ สหสฺสารานิ สเนมิกานิ สนาภิกานิ สพฺพาการปริปูรานิ, อิทมฺปิ ภิกฺขเว มหาปุริสสฺส มหาปุริสลกฺขณํ ภวติ.}

\prose{ปุน จปรํ ภิกฺขเว มหาปุริโส อายตปณฺหิ โหติ ฯเปฯ ทีฆงฺคุลิ โหติ.\csromanspacer{} มุทุตลุนหตฺถปาโท โหติ.\csromanspacer{} ชาลหตฺถปาโท โหติ.\csromanspacer{} อุสฺสงฺขปาโท โหติ.\csromanspacer{} เอณิชงฺโฆ โหติ.\csromanspacer{} ฐิตโกว อโนนมนฺโต อุโภหิ ปาณิตเลหิ ชณฺณุกานิ ปริมสติ ปริมชฺชติ.\csromanspacer{} โกโสหิตวตฺถคุโยฺห โหติ.\csromanspacer{} สุวณฺณวณฺโณ โหติ กญฺจนสนฺนิภตฺตโจ.\csromanspacer{} สุขุมจฺฉวิ โหติ, สุขุมตฺตา ฉวิยา รโชชลฺลํ กาเย น อุปลิมฺปติ.\csromanspacer{} เอเกกโลโม โหติ, เอเกกานิ โลมานิ โลมกูเปสุ ชาตานิ.\csromanspacer{} อุทฺธคฺคโลโม โหติ, อุทฺธคฺคานิ โลมานิ ชาตานิ นีลานิ อญฺชนวณฺณานิ กุณฺฑลาวฏฺฏานิ\footnote{กุณฺฑลาวตฺตานิ (พหูสุ)} ทกฺขิณาวฏฺฏกชาตานิ\footnote{ทกฺขิณาวตฺตกชาตานิ (สี, สฺยา, อิ)}.\csromanspacer{} พฺรหฺมุชุคตฺโต โหติ.\csromanspacer{} สตฺตุสฺสโท โหติ.\csromanspacer{} สีหปุพฺพทฺธกาโย โหติ.\csromanspacer{} จิตนฺตรํโส\footnote{ปิตนฺตรํโส (สฺยา)} โหติ.\csromanspacer{} นิโคฺรธปริมณฺฑโล โหติ, ยาวตกฺวสฺส กาโย ตาวตกฺวสฺส พฺยาโม, ยาวตกฺวสฺส พฺยาโม ตาวตกฺวสฺส กาโย.\csromanspacer{} สมวฏฺฏกฺขนฺโธ โหติ.\csromanspacer{} รสคฺคสคฺคี โหติ.\csromanspacer{} สีหหนุ โหติ.\csromanspacer{} จตฺตาลีสทนฺโต โหติ.\csromanspacer{} สมทนฺโต โหติ.\csromanspacer{} อวิรฬทนฺโต โหติ.\csromanspacer{} สุสุกฺกทาโฐ โหติ.\csromanspacer{} ปหูตชิโวฺห โหติ.\csromanspacer{} พฺรหฺมสฺสโร โหติ กรวีกภาณี.\csromanspacer{} อภินีลเนตฺโต โหติ.\csromanspacer{} โคปขุโม โหติ.\csromanspacer{} อุณฺณา ภมุกนฺตเร ชาตา โหติ โอทาตา มุทุตูลสนฺนิภา, ยมฺปิ ภิกฺขเว มหาปุริสสฺส อุณฺณา ภมุกนฺตเร ชาตา โหติ โอทาตา มุทุตูลสนฺนิภา, อิทมฺปิ ภิกฺขเว มหาปุริสสฺส มหาปุริสลกฺขณํ ภวติ.}

\prose{ปุน จปรํ ภิกฺขเว มหาปุริโส อุณฺหีสสีโส โหติ, ยมฺปิ ภิกฺขเว มหาปุริโส อุณฺหีสสีโส โหติ, อิทมฺปิ ภิกฺขเว มหาปุริสสฺส มหาปุริสลกฺขณํ ภวติ.}

\prose{อิมานิ โข ตานิ ภิกฺขเว ทฺวตฺติํส มหาปุริสสฺส มหาปุริสลกฺขณานิ, เยหิ สมนฺนาคตสฺส มหาปุริสสฺส เทฺวว คติโย ภวนฺติ อนญฺญา.}

\vspace{\tipitakaparskip}
\csromancenter{ลกฺขณสุตฺต สเจ อคารํ อชฺฌาวสติ, ราชา โหติ จกฺกวตฺตี ฯเปฯ.\csromanspacer{} สเจ โข ปน อคารสฺมา อนคาริยํ ปพฺพชติ, อรหํ โหติ สมฺมาสมฺพุทฺโธ โลเก วิวฏฺฏจฺฉโท.}
\prose{\csromanfolio{119}อิมานิ โข ภิกฺขเว ทฺวตฺติํส มหาปุริสสฺส มหาปุริสลกฺขณานิ พาหิรกาปิ อิสโย ธาเรนฺติ, โน จ โข เต ชานนฺติ “อิมสฺส กมฺมสฺส กฏตฺตา อิทํ ลกฺขณํ ปฏิลภตี”ติ.}

\csromanmatikaanchor{mtk.88}
\csromantocmarkref{section}{สุปฺปติฏฺฐิตปาทตาลกฺขณํ (1)}{mtk.88}
\bookmark[dest={mtk.88},level=1]{สุปฺปติฏฺฐิตปาทตาลกฺขณํ (1)}
\csromanheader{1}{\textbf{สุปฺปติฏฺฐิตปาทตาลกฺขณํ (1)}}
\proseitem{201}{ยมฺปิ ภิกฺขเว ตถาคโต ปุริมํ ชาติํ ปุริมํ ภวํ ปุริมํ นิเกตํ ปุพฺเพ มนุสฺสภูโต สมาโน ทฬฺหสมาทาโน อโหสิ กุสเลสุ ธมฺเมสุ อวตฺถิตสมาทาโน กายสุจริเต วจีสุจริเต มโนสุจริเต ทานสํวิภาเค สีลสมาทาเน อุโปสถุปวาเส มตฺเตยฺยตาย เปตฺเตยฺยตาย สามญฺญตาย พฺรหฺมญฺญตาย กุเล เชฏฺฐาปจายิตาย อญฺญตรญฺญตเรสุ จ อธิกุสเลสุ ธมฺเมสุ.\csromanspacer{} โส ตสฺส กมฺมสฺส กฏตฺตา อุปจิตตฺตา อุสฺสนฺนตฺตา วิปุลตฺตา กายสฺส เภทา ปรํ มรณา สุคติํ สคฺคํ โลกํ อุปปชฺชติ.\csromanspacer{} โส ตตฺถ อญฺเญ เทเว ทสหิ ฐาเนหิ อธิคฺคณฺหาติ ทิพฺเพน อายุนา ทิพฺเพน วณฺเณน ทิพฺเพน สุเขน ทิพฺเพน ยเสน ทิพฺเพน อาธิปเตยฺเยน ทิพฺเพหิ รูเปหิ ทิพฺเพหิ สทฺเทหิ ทิพฺเพหิ คนฺเธหิ ทิพฺเพหิ รเสหิ ทิพฺเพหิ โผฏฺฐพฺเพหิ.\csromanspacer{} โส ตโต จุโต อิตฺถตฺตํ อาคโต สมาโน อิมํ มหาปุริสลกฺขณํ ปฏิลภติ.\csromanspacer{} สุปฺปติฏฺฐิตปาโท โหติ, สมํ ปาทํ ภูมิยํ นิกฺขิปติ, สมํ อุทฺธรติ, สมํ สพฺพาวนฺเตหิ ปาทตเลหิ ภูมิํ ผุสติ.}

\proseitem{202}{โส เตน ลกฺขเณน สมนฺนาคโต สเจ อคารํ อชฺฌาวสติ, ราชา โหติ จกฺกวตฺตี ธมฺมิโก ธมฺมราชา จาตุรนฺโต วิชิตาวี ชนปทตฺถาวริยปฺปตฺโต สตฺตรตนสมนฺนาคโต.\csromanspacer{} ตสฺสิมานิ สตฺต รตนานิ ภวนฺติ, เสยฺยถิทํ, จกฺกรตนํ หตฺถิรตนํ อสฺสรตนํ มณิรตนํ อิตฺถิรตนํ คหปติรตนํ ปริณายกรตนเมว สตฺตมํ.\csromanspacer{} ปโรสหสฺสํ โข ปนสฺส ปุตฺตา ภวนฺติ สูรา วีรงฺครูปา ปรเสนปฺปมทฺทนา, โส อิมํ ปถวิํ สาครปริยนฺตํ อขิลมนิมิตฺตมกณฺฏกํ อิทฺธํ ผีตํ เขมํ สิวํ นิรพฺพุทํ อทณฺเฑน อสตฺเถน ธมฺเมน อภิวิชิย \csromanfolio{120}อชฺฌาวสติ.\csromanspacer{} ราชา สมาโน กิํ ลภติ.\csromanspacer{} อกฺขมฺภิโย\footnote{อวิกฺขมฺภิโย (สี, อิ)} โหติ เกนจิ มนุสฺสภูเตน ปจฺจตฺถิเกน ปจฺจามิตฺเตน.\csromanspacer{} ราชา สมาโน อิทํ ลภติ.\csromanspacer{} สเจ โข ปน อคารสฺมา อนคาริยํ ปพฺพชติ, อรหํ โหติ สมฺมาสมฺพุทฺโธ โลเก วิวฏฺฏจฺฉโท.\csromanspacer{} พุทฺโธ สมาโน กิํ ลภติ.\csromanspacer{} อกฺขมฺภิโย โหติ อพฺภนฺตเรหิ วา พาหิเรหิ วา ปจฺจตฺถิเกหิ ปจฺจามิตฺเตหิ ราเคน วา โทเสน วา โมเหน วา สมเณน วา พฺราหฺมเณน วา เทเวน วา มาเรน วา พฺรหฺมุนา วา เกนจิ วา โลกสฺมิํ.\csromanspacer{} พุทฺโธ สมาโน อิทํ ลภติ.\csromanspacer{} เอตมตฺถํ ภควา อโวจ.}

\csromanheader{2}{203. ตตฺเถตํ วุจฺจติ\csromandash{}}
\csromangathameasure{สจฺเจ จ ธมฺเม จ ทเม จ สํยเม, \\ โสเจยฺยสีลาลยุโปสเถสุ จ. \\ ทาเน อหิํสาย อสาหเส รโต, \\ ทฬฺหํ สมาทาย สมตฺตมาจริ. \\ โส เตน กมฺเมน ทิวํ สมกฺกมิ, \\ สุขญฺจ ขิฑฺฑารติโย จ อนฺวภิ. \\ ตโต จวิตฺวา ปุนราคโต อิธ, \\ สเมหิ ปาเทหิ ผุสี วสุนฺธรํ. \\ พฺยากํสุ เวยฺยญฺชนิกา สมาคตา, \\ สมปฺปติฏฺฐสฺส น โหติ ขมฺภนา. \\ คิหิสฺส วา ปพฺพชิตสฺส วา ปุน, \\ ตํ ลกฺขณํ ภวติ ตทตฺถโชตกํ. \\ อกฺขมฺภิโย โหติ อคารมาวสํ, \\ ปราภิภู สตฺตุภิ นปฺปมทฺทโน. \\ มนุสฺสภูเตนิธ โหติ เกนจิ, \\ อกฺขมฺภิโย ตสฺส ผเลน กมฺมุโน.}
\csromangathagroup{\csromangathastanza{สจฺเจ จ ธมฺเม จ ทเม จ สํยเม, \\ โสเจยฺยสีลาลยุโปสเถสุ จ. \\ ทาเน อหิํสาย อสาหเส รโต, \\ ทฬฺหํ สมาทาย สมตฺตมาจริ\footnote{สมนฺตมาจริ (สฺยา, ก)}.}\csromangathastanza{โส เตน กมฺเมน ทิวํ สมกฺกมิ\footnote{อปกฺกมิ (สฺยา, ก)}, \\ สุขญฺจ ขิฑฺฑารติโย จ อนฺวภิ\footnote{อปกฺกมิ (สฺยา, ก)}. \\ ตโต จวิตฺวา ปุนราคโต อิธ, \\ สเมหิ ปาเทหิ ผุสี วสุนฺธรํ.}\csromangathastanza{พฺยากํสุ เวยฺยญฺชนิกา สมาคตา, \\ สมปฺปติฏฺฐสฺส น โหติ ขมฺภนา. \\ คิหิสฺส วา ปพฺพชิตสฺส วา ปุน\footnote{ปน (สฺยา)}, \\ ตํ ลกฺขณํ ภวติ ตทตฺถโชตกํ.}\csromangathastanza{อกฺขมฺภิโย โหติ อคารมาวสํ, \\ ปราภิภู สตฺตุภิ นปฺปมทฺทโน. \\ มนุสฺสภูเตนิธ โหติ เกนจิ, \\ อกฺขมฺภิโย ตสฺส ผเลน กมฺมุโน.}}

\csromanheader{2}{7. ลกฺขณสุตฺต}
\csromangathameasure{สเจ จ ปพฺพชฺชมุเปติ ตาทิโส, \\ เนกฺขมฺมฉนฺทาภิรโต วิจกฺขโณ. \\ อคฺโค น โส คจฺฉติ ชาตุ ขมฺภนํ, \\ นรุตฺตโม เอส หิ ตสฺส ธมฺมตาติ.}
\csromangathagroup{\csromangathastanza{\csromanfolio{121}สเจ จ ปพฺพชฺชมุเปติ ตาทิโส, \\ เนกฺขมฺมฉนฺทาภิรโต วิจกฺขโณ. \\ อคฺโค น โส คจฺฉติ ชาตุ ขมฺภนํ, \\ นรุตฺตโม เอส หิ ตสฺส ธมฺมตาติ.}}

\csromanmatikaanchor{mtk.89}
\csromantocmarkref{section}{ปาทตลจกฺกลกฺขณํ (2)}{mtk.89}
\bookmark[dest={mtk.89},level=1]{ปาทตลจกฺกลกฺขณํ (2)}
\csromanheader{1}{\textbf{ปาทตลจกฺกลกฺขณํ (2)}}
\proseitem{204}{ยมฺปิ ภิกฺขเว ตถาคโต ปุริมํ ชาติํ ปุริมํ ภวํ ปุริมํ นิเกตํ ปุพฺเพ มนุสฺสภูโต สมาโน พหุชนสฺส สุขาวโห อโหสิ, อุพฺเพค-อุตฺตาสภยํ อปนุทิตา, ธมฺมิกญฺจ รกฺขาวรณคุตฺติํ สํวิธาตา, สปริวารญฺจ ทานํ อทาสิ.\csromanspacer{} โส ตสฺส กมฺมสฺส กฏตฺตา อุปจิตตฺตา อุสฺสนฺนตฺตา วิปุลตฺตา กายสฺส เภทา ปรํ มรณา สุคติํ สคฺคํ โลกํ อุปปชฺชติ ฯเปฯ.\csromanspacer{} โส ตโต จุโต อิตฺถตฺตํ อาคโต สมโน อิมํ มหาปุริสลกฺขณํ ปฏิลภติ.\csromanspacer{} เหฏฺฐาปาทตเลสุ จกฺกานิ ชาตานิ โหนฺติ สหสฺสารานิ สเนมิกานิ สนาภิกานิ สพฺพาการปริปูรานิ สุวิภตฺตนฺตรานิ.}

\prose{โส เตน ลกฺขเณน สมนฺนาคโต สเจ อคารํ อชฺฌาวสติ, ราชา โหติ จกฺกวตฺตี ฯเปฯ.\csromanspacer{} ราชา สมาโน กิํ ลภติ.\csromanspacer{} มหาปริวาโร โหติ มหาสฺส โหนฺติ ปริวารา พฺราหฺมณคหปติกา เนคมชานปทา คณกมหามตฺตา อนีกฏฺฐา โทวาริกา อมจฺจา ปาริสชฺชา ราชาโน โภคิยา กุมารา.\csromanspacer{} ราชา สมาโน อิทํ ลภติ.\csromanspacer{} สเจ อคารสฺมา อนคาริยํ ปพฺพชติ, อรหํ โหติ สมฺมาสมฺพุทฺโธ โลเก วิวฏฺฏจฺฉโท.\csromanspacer{} พุทฺโธ สมาโน กิํ ลภติ.\csromanspacer{} มหาปริวาโร โหติ, มหาสฺส โหนฺติ ปริวารา ภิกฺขู ภิกฺขุนิโย อุปาสกา อุปาสิกาโย เทวา มนุสฺสา อสุรา นาคา คนฺธพฺพา.\csromanspacer{} พุทฺโธ สมาโน อิทํ ลภติ.\csromanspacer{} เอตมตฺถํ ภควา อโวจ.}

\csromanheader{2}{205. ตตฺเถตํ วุจฺจติ\csromandash{}}
\csromanmatikaanchor{mtk.90}
\csromantocmarkref{section}{อายตนปณฺหิตาทิติลกฺขณํ (3-5)}{mtk.90}
\bookmark[dest={mtk.90},level=1]{อายตนปณฺหิตาทิติลกฺขณํ (3-5)}
\csromangathameasure{ปุเร ปุรตฺถา ปุริมาสุ ชาติสุ, \\ มนุสฺสภูโต พหุนํ สุขาวโห. \\ อุพฺเภค-อุตฺตาสภยาปนูทโน, \\ คุตฺตีสุ รกฺขาวรเณสุ อุสฺสุโก. \\ โส เตน กมฺเมน ทิวํ สมกฺกมิ, \\ สุขญฺจ ขิฑฺฑารติโย จ อนฺวภิ. \\ ตโต จวิตฺวา ปุนราคโต อิธ, \\ จกฺกานิ ปาเทสุ ทุเวสุ วินฺทติ. \\ สมนฺตเนมีนิ สหสฺสรานิ จ, \\ พฺยากํสุ เวยฺยญฺชนิกา สมาคตา. \\ ทิสฺวา กุมารํ สตปุญฺญลกฺขณํ, \\ ปริวารวา เหสฺสติ สตฺตุมทฺทโน. \\ ตถา หิ จกฺกานิ สมนฺเตเนมินิ, \\ สเจ น ปพฺพชฺชมุเปติ ตาทิโส. \\ วตฺเตติ จกฺกํ ปถวิํ ปสาสติ, \\ ตสฺสานุยนฺตาธ ภวนฺติ ขตฺติยา. \\ มหายสํ สํปริวารยนฺติ นํ, \\ สเจ จ ปพฺพชฺชมุเปติ ตาทิโส. \\ เนกฺขมฺมฉนฺทาภิรโต วิจกฺขโณ, \\ เทวา มนุสฺสาสุรสกฺกรกฺขาสา. \\ คนฺธพฺพนาคา วิหคา จตุปฺปทา, \\ อนุตฺตรํ เทวมนุสฺสปูชิตํ, \\ มหายสํ สํปริวารยนฺติ นนฺติ.}
\csromangathagroup{\csromangathastanza{ปุเร ปุรตฺถา ปุริมาสุ ชาติสุ, \\ มนุสฺสภูโต พหุนํ สุขาวโห. \\ อุพฺเภค-อุตฺตาสภยาปนูทโน, \\ คุตฺตีสุ รกฺขาวรเณสุ อุสฺสุโก.}\csromangathastanza{\csromanfolio{122}โส เตน กมฺเมน ทิวํ สมกฺกมิ, \\ สุขญฺจ ขิฑฺฑารติโย จ อนฺวภิ. \\ ตโต จวิตฺวา ปุนราคโต อิธ, \\ จกฺกานิ ปาเทสุ ทุเวสุ วินฺทติ.}\csromangathastanza{สมนฺตเนมีนิ สหสฺสรานิ จ, \\ พฺยากํสุ เวยฺยญฺชนิกา สมาคตา. \\ ทิสฺวา กุมารํ สตปุญฺญลกฺขณํ, \\ ปริวารวา เหสฺสติ สตฺตุมทฺทโน.}\csromangathastanza{ตถา หิ จกฺกานิ สมนฺเตเนมินิ, \\ สเจ น ปพฺพชฺชมุเปติ ตาทิโส. \\ วตฺเตติ จกฺกํ ปถวิํ ปสาสติ, \\ ตสฺสานุยนฺตาธ\footnote{ตสฺสานุยุตฺตา อิธ (สี, อิ), ตสฺสานุยนฺตา อิธ (สฺยา, ก)} ภวนฺติ ขตฺติยา.}\csromangathastanza{มหายสํ สํปริวารยนฺติ นํ, \\ สเจ จ ปพฺพชฺชมุเปติ ตาทิโส. \\ เนกฺขมฺมฉนฺทาภิรโต วิจกฺขโณ, \\ เทวา มนุสฺสาสุรสกฺกรกฺขาสา\footnote{สตฺตรกฺขสา (ก) สี-สฺยา-อฏฺฐกถา โอโลเกตพฺพา.}.}\csromangathastanza{คนฺธพฺพนาคา วิหคา จตุปฺปทา, \\ อนุตฺตรํ เทวมนุสฺสปูชิตํ, \\ มหายสํ สํปริวารยนฺติ นนฺติ.}}

\csromanheader{2}{\textbf{อายตปณฺหิตาทิติลกฺขณํ (3-5)}}
\proseitem{206}{ยมฺปิ ภิกฺขเว ตถาคโต ปุริมํ ชาติํ ปุริมํ ภวํ ปุริมํ นิเกตํ ปุพฺเพ มนุสฺสภูโต สมาโน ปาณาติปาตํ ปหาย ปาณาติปาตา ปฏิวิรโต อโหสิ นิหิตทณฺโฑ นิหิตสตฺโถ ลชฺชี ทยาปนฺโน, สพฺพปาณภูตหิตานุกมฺปี วิหาสิ.\csromanspacer{} โส ตสฺส กมฺมสฺส กฏตฺตา อุปจิตตฺตา อุสฺสนฺนตฺตา วิปุลตฺตา ฯเปฯ.\csromanspacer{} โส ตโต จุโต อิตฺถตฺตํ อาคโต สมาโน อิมานิ ตีณิ มหาปุริสลกฺขณานิ ปฏิลภติ.\csromanspacer{} อายตปณฺหิ จ โหติ ทีฆงฺคุลิ จ พฺรหฺมุชุคตฺโต จ.}

\csromanheader{2}{7. ลกฺขณสุตฺต}
\prose{\csromanfolio{123}โส เตหิ ลกฺขเณหิ สมนฺนาคโต สเจ อคารํ อชฺฌาวสติ, ราชา โหติ จกฺกวตฺตี ฯเปฯ.\csromanspacer{} ราชา สมาโน กิํ ลภติ.\csromanspacer{} ทีฆายุโก โหติ จิรฏฺฐิติโก, ทีฆมายุํ ปาเลติ, น สกฺกา โหติ อนฺตรา ชีวิตา โวโรเปตุํ เกนจิ มนุสฺสภูเตน ปจฺจตฺถิเกน ปจฺจามิตฺเตน.\csromanspacer{} ราชา สมาโน อิทํ ลภติ ฯเปฯ.\csromanspacer{} พุทฺโธ สมาโน กิํ ลภติ.\csromanspacer{} ทีฆายุโก โหติ จิรฏฺฐิติโก, ทีฆมายุํ ปาเลติ, น สกฺกา โหติ อนฺตรา ชีวิตา โวโรเปตุํ ปจฺจตฺถิเกหิ ปจฺจามิตฺเตหิ สมเณน วา พฺราหฺมเณน วา เทเวน วา มาเรน วา พฺรหฺมุนา วา เกนจิ วา โลกสฺมิํ.\csromanspacer{} พุทฺโธ สมาโน อิทํ ลภติ.\csromanspacer{} เอตมตฺถํ ภควา อโวจ.}

\csromanheader{2}{207. ตตฺเถตํ วุจฺจติ\csromandash{}}
\csromangathameasure{มารณวธภยตฺตโน วิทิตฺวา, \\ ปฏิวิรโต ปรํ มารณายโหสิ. \\ เตน สุจริเตน สคฺคมคมา, \\ สุกตผลวิปากมนุโภสิ. \\ จวิย ปุนริธาคโต สมาโน, \\ ปฏิลภติ อิธ ตีณิ ลกฺขณานิ. \\ ภวติ วิปุลทีฆปาสณฺหิโก, \\ พฺรหฺมาว สุชุ สุโภ สุชาตคตฺโต. \\ สุภุโช สุสุ สุสณฺฐิโต สุชาโต, \\ มุทุตลุนงฺคุลิยสฺส โหนฺติ. \\ ทีฆา ตีภิ ปุริสวรคฺคลกฺขเณหิ, \\ จิรยปนาย กุมารมาทิสนฺติ. \\ ภวติ ยทิ คิหี จิรํ ยเปติ, \\ จิรตรํ ปพฺพชติ ยทิ ตโต หิ. \\ ยาปยติ จ วสิทฺธิภาวนาย, \\ อิติ ทีฆายุกตาย ตํ นิมิตฺตนฺติ.}
\csromangathagroup{\csromangathastanza{มารณวธภยตฺตโน\footnote{มรณวธภยตฺตโน (สี, อิ, ก), มรณวธภยมตฺตโน (สฺยา)} วิทิตฺวา, \\ ปฏิวิรโต ปรํ มารณายโหสิ. \\ เตน สุจริเตน สคฺคมคมา\footnote{มรณวธภยตฺตโน (สี, อิ, ก), มรณวธภยมตฺตโน (สฺยา)}, \\ สุกตผลวิปากมนุโภสิ.}\csromangathastanza{จวิย ปุนริธาคโต สมาโน, \\ ปฏิลภติ อิธ ตีณิ ลกฺขณานิ. \\ ภวติ วิปุลทีฆปาสณฺหิโก, \\ พฺรหฺมาว สุชุ สุโภ สุชาตคตฺโต.}\csromangathastanza{สุภุโช สุสุ สุสณฺฐิโต สุชาโต, \\ มุทุตลุนงฺคุลิยสฺส โหนฺติ. \\ ทีฆา ตีภิ ปุริสวรคฺคลกฺขเณหิ, \\ จิรยปนาย\footnote{จิรยาปนาย (สฺยา)} กุมารมาทิสนฺติ.}\csromangathastanza{ภวติ ยทิ คิหี จิรํ ยเปติ, \\ จิรตรํ ปพฺพชติ ยทิ ตโต หิ. \\ ยาปยติ จ วสิทฺธิภาวนาย, \\ อิติ ทีฆายุกตาย ตํ นิมิตฺตนฺติ.}}

\csromanmatikaanchor{mtk.91}
\csromantocmarkref{section}{สตฺตุสฺสทตาลกฺขณํ (6)}{mtk.91}
\bookmark[dest={mtk.91},level=1]{สตฺตุสฺสทตาลกฺขณํ (6)}
\csromanheader{1}{\textbf{สตฺตุสฺสทตาลกฺขณํ (6)}}
\proseitem{208}{\csromanfolio{124}ยมฺปิ ภิกฺขเว ตถาคโต ปุริมํ ชาติํ ปุริมํ ภวํ ปุริมํ นิเกตํ ปุพฺเพ มนุสฺสภูโต สมาโน ทาตา อโหสิ ปณีตานํ รสิตานํ ขาทนียานํ โภชนียานํ สายนียานํ เลหนียานํ ปานานํ.\csromanspacer{} โส ตสฺส กมฺมสฺส กฏตฺตา ฯเปฯ.\csromanspacer{} โส ตโต จุโต อิตฺถตฺตํ อาคโต สมาโน อิมํ มหาปุริสลกฺขณํ ปฏิลภติ.\csromanspacer{} สตฺตุสฺสโท โหติ, สตฺตสฺส อุสฺสทา โหนฺติ, อุโภสุ หตฺเตสุ อุสฺสทา โหนฺติ, อุโภสุ ปาเทสุ อุสฺสทา โหนฺติ, อุโภสุ อํสกูเฏสุ อุสฺสทา โหนฺติ, ขนฺเธ อุสฺสโท โหติ.}

\prose{โส เตน ลกฺขเณน สมนฺนาคโต สเจ อคารํ อชฺฌาวสติ, ราชา โหติ จกฺกวตฺตี ฯเปฯ.\csromanspacer{} ราชา สมาโน กิํ ลภติ.\csromanspacer{} ลาภี โหติ ปณีตานํ รสิตานํ ขาทนียานํ โภชนียานํ สายนียานํ เลหนียานํ ปานานํ.\csromanspacer{} ราชา สมาโน อิทํ ลภติ ฯเปฯ.\csromanspacer{} พุทฺโธ สมาโน กิํ ลภติ.\csromanspacer{} ลาภี โหติ ปณีตานํ รสิตานํ ขาทนียานํ โภชนียานํ สายนียานํ เลหนียานํ ปานานํ.\csromanspacer{} พุทฺโธ สมาโน อิทํ ลภติ.\csromanspacer{} เอตมตฺถํ ภควา อโวจ.}

\csromanheader{2}{209. ตตฺเถตํ วุจฺจติ\csromandash{}}
\csromangathameasure{ขชฺชโภชฺชมถ เลยฺย สายิยํ, \\ อุตฺตมคฺครสทายโก อหุ. \\ เตน โส สุจริเตน กมฺมุนา, \\ นนฺทเน จิรมภิปฺปโมทติ. \\ สตฺต จุสฺสเท อิธาธิคจฺฉติ, \\ หตฺถปาทมุทุตญฺจ วินฺทติ. \\ อาหุ พฺยญฺชนนิมิตฺตโกวิทา, \\ ขชฺชโภชฺชรสลาภิตาย นํ. \\ ยํ คิหิสฺสปิ ตทตฺถโชตกํ, \\ ปพฺพชมฺปิ จ ตทาธิคจฺฉติ. \\ ขชฺชโภชฺชรสลาภิรุตฺตมํ, \\ อาหุ สพฺพคิหิพนฺธนจฺฉิทนฺติ.}
\csromangathagroup{\csromangathastanza{ขชฺชโภชฺชมถ เลยฺย สายิยํ, \\ อุตฺตมคฺครสทายโก อหุ. \\ เตน โส สุจริเตน กมฺมุนา, \\ นนฺทเน จิรมภิปฺปโมทติ.}\csromangathastanza{สตฺต จุสฺสเท อิธาธิคจฺฉติ, \\ หตฺถปาทมุทุตญฺจ วินฺทติ. \\ อาหุ พฺยญฺชนนิมิตฺตโกวิทา, \\ ขชฺชโภชฺชรสลาภิตาย นํ.}\csromangathastanza{ยํ คิหิสฺสปิ\footnote{น ตํ คิหิสฺสาปิ (สฺยา)} ตทตฺถโชตกํ, \\ ปพฺพชมฺปิ จ ตทาธิคจฺฉติ. \\ ขชฺชโภชฺชรสลาภิรุตฺตมํ, \\ อาหุ สพฺพคิหิพนฺธนจฺฉิทนฺติ.}}

\csromanmatikaanchor{mtk.92}
\csromantocmarkref{section}{กรจรณมุทุชาลตาลกฺขณานิ (7-8)}{mtk.92}
\bookmark[dest={mtk.92},level=1]{กรจรณมุทุชาลตาลกฺขณานิ (7-8)}
\csromanheader{1}{7. ลกฺขณสุตฺต \textbf{กรจรณมุทุชาลตาลกฺขณานิ (7-8)}}
\proseitem{210}{\csromanfolio{125}ยมฺปิ ภิกฺขเว ตถาคโต ปุริมํ ชาติํ ปุริมํ ภวํ ปุริมํ นิเกตํ ปุพฺเพ มนุสฺสภูโต สมาโน จตูหิ สงฺคหวตฺถูหิ ชนํ สงฺคาหโก อโหสิ ทาเนน เปยฺยวชฺเชน\footnote{ปิยวาเจน (สฺยา, ก)} อตฺถจริยาย สมานตฺตตาย.\csromanspacer{} โส ตสฺส กมฺมสฺส กฏตฺตา ฯเปฯ.\csromanspacer{} โส ตโต จุโต อิตฺถตฺตํ อาคโต สมาโน อิมานิ เทฺว มหาปุริสลกฺขณานิ ปฏิลภติ.\csromanspacer{} มุทุตลุนหตฺถปาโท จ โหติ ชาลหตฺถปาโท จ.}

\prose{โส เตหิ ลกฺขเณหิ สมนฺนาคโต สเจ อคารํ อชฺฌาวสติ, ราชา โหติ จกฺกวตฺตี ฯเปฯ.\csromanspacer{} ราชา สมาโน กิํ ลภติ.\csromanspacer{} สุสงฺคหิตปริชโน โหติ, สุสงฺคหิตาสฺส โหนฺติ พฺราหฺมณคหปติกา เนคมชานปาทา คณกมหามตฺตา อนีกฏฺฐา โทวาริกา อมจฺจา ปาริสชฺชา ราชาโน โภคิยา กุมารา.\csromanspacer{} ราชา สมาโน อิทํ ลภติ.\csromanspacer{} พุทฺโธ สมาโน กิํ ลภติ.\csromanspacer{} สุสงฺคหิตปริชโน โหติ, สุสงฺคหิตาสฺส โหนฺติ ภิกฺขู ภิกฺขุนิโย อุปาสกา อุปาสิกาโย เทวา มนุสฺสา อสุรา นาคา คนฺธพฺพา.\csromanspacer{} พุทฺโธ สมาโน อิทํ ลภติ.\csromanspacer{} เอตมตฺถํ ภควา อโวจ.}

\csromanheader{2}{211. ตตฺเถตํ วุจฺจติ\csromandash{}}
\csromangathameasure{ทานมฺปิ จตฺถจริยตญฺจ, \\ ปิยวาทิตญฺจ สมานตฺตตญฺจ. \\ กริยจริยสุสงฺคหํ พหูนํ, \\ อนวมเตน คุเณน ยาติ สคฺคํ. \\ จวิย ปุนริธาคโต สมาโน, \\ กรจรณมุทุตญฺจ ชาลิโน จ. \\ อติรุจิรสุวคฺคุทสฺสเนยฺยํ, \\ ปฏิลภติ ทหโร สุสุ กุมาโร. \\ ภวติ ปริชนสฺสโว วิเธยฺโย, \\ มหิมํ อาวสิโต สุสงฺคหิโต. \\ ปิยวทู หิตสุขตํ ชิคีสมาโน, \\ อภิรุจิตานิ คุณานิ อาจรติ. \\ ยทิ จ ชหติ สพฺพกามโภคํ, \\ กถยติ ธมฺมกถํ ชิโน ชนสฺส. \\ วจนปฏิกรสฺสาภิปฺปสนฺนา, \\ สุตฺวาน ธมฺมานุธมฺมมาจรนฺตีติ.}
\csromangathagroup{\csromangathastanza{ทานมฺปิ จตฺถจริยตญฺจ\footnote{ทานมฺปิ จ อตฺถจริยตมฺปิ จ (สี, อิ)}, \\ ปิยวาทิตญฺจ สมานตฺตตญฺจ\footnote{ทานมฺปิ จ อตฺถจริยตมฺปิ จ (สี, อิ)}. \\ กริยจริยสุสงฺคหํ พหูนํ, \\ อนวมเตน คุเณน ยาติ สคฺคํ.}\csromangathastanza{จวิย ปุนริธาคโต สมาโน, \\ กรจรณมุทุตญฺจ ชาลิโน จ. \\ อติรุจิรสุวคฺคุทสฺสเนยฺยํ, \\ ปฏิลภติ ทหโร สุสุ กุมาโร.}\csromangathastanza{\csromanfolio{126}ภวติ ปริชนสฺสโว วิเธยฺโย, \\ มหิมํ อาวสิโต สุสงฺคหิโต. \\ ปิยวทู หิตสุขตํ ชิคีสมาโน\footnote{ชิคิํสมาโน (สี, สฺยา, อิ)}, \\ อภิรุจิตานิ คุณานิ อาจรติ.}\csromangathastanza{ยทิ จ ชหติ สพฺพกามโภคํ, \\ กถยติ ธมฺมกถํ ชิโน ชนสฺส. \\ วจนปฏิกรสฺสาภิปฺปสนฺนา, \\ สุตฺวาน ธมฺมานุธมฺมมาจรนฺตีติ.}}

\csromanmatikaanchor{mtk.93}
\csromantocmarkref{section}{อุสฺสงฺขปาท-อุทฺธคฺคโลมตาลกฺขณานิ (9-10)}{mtk.93}
\bookmark[dest={mtk.93},level=1]{อุสฺสงฺขปาท-อุทฺธคฺคโลมตาลกฺขณานิ (9-10)}
\csromanheader{1}{\textbf{อุสฺสงฺขปาท-อุทฺธคฺคโลมตาลกฺขณานิ (9-10)}}
\proseitem{212}{ยมฺปิ ภิกฺขเว ตถาคโต ปุริมํ ชาติํ ปุริมํ ภวํ ปุริมํ นิเกตํ ปุพฺเพ มนุสฺสภูโต สมาโน\footnote{สมาโน พหุโน ชนสฺส (สี, อิ)} อตฺถูปสํหิตํ ธมฺมูปสํหิตํ วาจํ ภาสิตา อโหสิ, พหุชนํ นิทํเสสิ, ปาณีนํ หิตสุขาวโห ธมฺมยาคี.\csromanspacer{} โส ตสฺส กมฺมสฺส กฏตฺตา ฯเปฯ.\csromanspacer{} โส ตโต จุโต อิตฺถตฺตํ อาคโต สมาโน อิมานิ เทฺว มหาปุริสลกฺขณานิ ปฏิลภติ.\csromanspacer{} อุสฺสงฺขปาโท จ โหติ อุทฺธคฺคโลโม จ.}

\prose{โส เตหิ ลกฺขเณหิ สมนฺนาคโต สเจ อคารํ อชฺฌาวสติ, ราชา โหติ จกฺกวตฺตี ฯเปฯ.\csromanspacer{} ราชา สมาโน กิํ ลภติ.\csromanspacer{} อคฺโค จ โหติ เสฏฺโฐ จ ปาโมกฺโข จ อุตฺตโม จ ปวโร จ กามโภคีนํ.\csromanspacer{} ราชา สมาโน อิทํ ลภติ ฯเปฯ.\csromanspacer{} พุทฺโธ สมาโน กิํ ลภติ.\csromanspacer{} อคฺโค จ โหติ เสฏฺโฐ จ ปาโมกฺโข จ อุตฺตโม จ ปวโร จ สพฺพสตฺตานํ.\csromanspacer{} พุทฺโธ สมาโน อิทํ ลภติ.\csromanspacer{} เอตมตฺถํ ภควา อโวจ.}

\csromanheader{2}{213. ตตฺเถตํ วุจฺจติ\csromandash{}}
\csromangathameasure{อตฺถธมฺมสหิตํ ปุเร คิรํ, \\ เอรยํ พหุชนํ นิทํสยิ. \\ ปาณินํ หิตสุขาวโห อหุ, \\ ธมฺมยาคมยชี อมจฺฉรี.}
\csromangathagroup{\csromangathastanza{อตฺถธมฺมสหิตํ\footnote{อตฺถธมฺมสํหิตํ (ก-สี, อิ), อตฺถธมฺมุปสํหิตํ (ก)} ปุเร คิรํ, \\ เอรยํ พหุชนํ นิทํสยิ. \\ ปาณินํ หิตสุขาวโห อหุ, \\ ธมฺมยาคมยชี\footnote{อตฺถธมฺมสํหิตํ (ก-สี, อิ), อตฺถธมฺมุปสํหิตํ (ก)} อมจฺฉรี.}}

\csromanheader{2}{7. ลกฺขณสุตฺต}
\csromangathameasure{เตน โส สุจริเตน กมฺมุนา, \\ สุคฺคติํ วชติ ตตฺถ โมทติ. \\ ลกฺขณานิ จ ทุเว อิธาคโต, \\ อุตฺตมปฺปมุขตาย วินฺทติ. \\ อุพฺภมุปฺปติตโลมวา สโส, \\ ปาทคณฺฐิรหุ สาธุสณฺฐิตา. \\ มํสโลหิตาจิตา ตโจตฺถตา, \\ อุปริจรณโสภนา อหุ. \\ เคหมาวสติ เจ ตถาวิโธ, \\ อคฺคตํ วชติ กามโภคินํ. \\ เตน อุตฺตริตโร น วิชฺชติ, \\ ชมฺพุทีปมภิภุยฺย อิริยติ. \\ ปพฺพชมฺปิ จ อโนมนิกฺกโม, \\ อคฺคตํ วชติ สพฺพปาณินํ. \\ เตน อุตฺตริตโร น วิชฺชติ, \\ สพฺพโลกมภิภุยฺย วิหรตีติ.}
\csromangathagroup{\csromangathastanza{\csromanfolio{127}เตน โส สุจริเตน กมฺมุนา, \\ สุคฺคติํ วชติ ตตฺถ โมทติ. \\ ลกฺขณานิ จ ทุเว อิธาคโต, \\ อุตฺตมปฺปมุขตาย\footnote{อุตฺตมสุขตาย (สฺยา), อุตฺตมปมุกฺขตาย (ก)} วินฺทติ.}\csromangathastanza{อุพฺภมุปฺปติตโลมวา สโส, \\ ปาทคณฺฐิรหุ สาธุสณฺฐิตา. \\ มํสโลหิตาจิตา ตโจตฺถตา, \\ อุปริจรณโสภนา\footnote{อุปริชานุโสภนา (สฺยา), อุปริ จ ปน โสภนา (สี, อิ)} อหุ.}\csromangathastanza{เคหมาวสติ เจ ตถาวิโธ, \\ อคฺคตํ วชติ กามโภคินํ. \\ เตน อุตฺตริตโร น วิชฺชติ, \\ ชมฺพุทีปมภิภุยฺย อิริยติ.}\csromangathastanza{ปพฺพชมฺปิ จ อโนมนิกฺกโม, \\ อคฺคตํ วชติ สพฺพปาณินํ. \\ เตน อุตฺตริตโร น วิชฺชติ, \\ สพฺพโลกมภิภุยฺย วิหรตีติ.}}

\csromanmatikaanchor{mtk.94}
\csromantocmarkref{section}{เอณิชงฺฆลกฺขณํ (11)}{mtk.94}
\bookmark[dest={mtk.94},level=1]{เอณิชงฺฆลกฺขณํ (11)}
\csromanheader{1}{\textbf{เอณิชงฺฆลกฺขณํ (11)}}
\proseitem{214}{ยมฺปิ ภิกฺขเว ตถาคโต ปุริมํ ชาติํ ปุริมํ ภวํ ปุริมํ นิเกตํ ปุพฺเพ มนุสฺสภูโต สมาโน สกฺกจฺจํ วาเจตา อโหสิ สิปฺปํ วา วิชฺชํ วา จรณํ วา กมฺมํ วา “กินฺติ เม ขิปฺปํ วิชาเนยฺยุํ, ขิปฺปํ ปฏิปชฺเชยฺยุํ, น จิรํ กิลิสฺเสยฺยุนฺ”ติ.\csromanspacer{} โส ตสฺส กมฺมสฺส กฏตฺตา ฯเปฯ.\csromanspacer{} โส ตโต จุโต อิตฺถตฺตํ อาคโต สมาโน อิมํ มหาปุริสลกฺขณํ ปฏิลภติ.\csromanspacer{} เอณิชงฺโฆ โหติ.}

\prose{โส เตน ลกฺขเณน สมนฺนาคโต สเจ อคารํ อชฺฌาวสติ, ราชา โหติ จกฺกวตฺตี ฯเปฯ.\csromanspacer{} ราชา สมาโน กิํ ลภติ.\csromanspacer{} ยานิ ตานิ ราชารหานิ ราชงฺคานิ ราชูปโภคานิ ราชานุจฺฉวิกานิ, ตานิ ขิปฺปํ ลภติ.\csromanspacer{} ราชา สมาโน อิทํ ลภติ ฯเปฯ.\csromanspacer{} พุทฺโธ สมาโน กิํ ลภติ. \csromanfolio{128}ยานิ ตานิ สมณารหานิ สมณงฺคานิ สมณูปโภคานิ สมณานุจฺฉวิกานิ, ตานิ ขิปฺปํ ปฏิลภติ.\csromanspacer{} พุทฺโธ สมาโน อิทํ ลภติ.\csromanspacer{} เอตมตฺถํ ภควา อโวจ.}

\csromanheader{2}{215. ตตฺเถตํ วุจฺจติ\csromandash{}}
\csromangathameasure{สิปฺเปสุ วิชฺชาจรเณสุ กมฺเมสุ, \\ กถํ วิชาเนยฺยุํ ลหุนฺติ อิจฺฉติ. \\ ยทูปฆาตาย น โหติ กสฺสจิ, \\ วาเจติ ขิปฺปํ น จิรํ กิลิสฺสติ. \\ ตํ กมฺมํ กตฺวา กุสลํ สุขุทฺรยํ, \\ ชงฺฆา มนุญฺญา ลภเต สุสณฺฐิตา. \\ วฏฺฏา สุชาตา อนุปุพฺพมุคฺคตา, \\ อุทฺธคฺคโลมา สุขุมตฺตโจตฺถตา. \\ เอเณยฺยชงฺโฆติ ตมาหุ ปุคฺคลํ, \\ สมฺปตฺติยา ขิปฺปมิธาหุ ลกฺขณํ. \\ เคหานุโลมานิ ยทาภิกงฺขติ, \\ อปพฺพชํ ขิปฺปมิธาธิคจฺฉติ. \\ สเจ จ ปพฺพชฺชมุเปติ ตาทิโส, \\ เนกฺขมฺมฉนฺทาภิรโต วิจกฺขโณ. \\ อนุจฺฉวิกสฺส ยทานุโลมิกํ, \\ ตํ วินฺทติ ขิปฺปมโนมวิกฺกโม ติ.}
\csromangathagroup{\csromangathastanza{สิปฺเปสุ วิชฺชาจรเณสุ กมฺเมสุ\footnote{กมฺมสุ (สี, อิ)}, \\ กถํ วิชาเนยฺยุํ\footnote{กมฺมสุ (สี, อิ)} ลหุนฺติ อิจฺฉติ. \\ ยทูปฆาตาย น โหติ กสฺสจิ, \\ วาเจติ ขิปฺปํ น จิรํ กิลิสฺสติ.}\csromangathastanza{ตํ กมฺมํ กตฺวา กุสลํ สุขุทฺรยํ\footnote{สุขินฺทฺริยํ (ก)}, \\ ชงฺฆา มนุญฺญา ลภเต สุสณฺฐิตา. \\ วฏฺฏา สุชาตา อนุปุพฺพมุคฺคตา, \\ อุทฺธคฺคโลมา สุขุมตฺตโจตฺถตา.}\csromangathastanza{เอเณยฺยชงฺโฆติ ตมาหุ ปุคฺคลํ, \\ สมฺปตฺติยา ขิปฺปมิธาหุ ลกฺขณํ. \\ เคหานุโลมานิ ยทาภิกงฺขติ, \\ อปพฺพชํ ขิปฺปมิธาธิคจฺฉติ.}\csromangathastanza{สเจ จ ปพฺพชฺชมุเปติ ตาทิโส, \\ เนกฺขมฺมฉนฺทาภิรโต วิจกฺขโณ. \\ อนุจฺฉวิกสฺส ยทานุโลมิกํ, \\ ตํ วินฺทติ ขิปฺปมโนมวิกฺกโม\footnote{นิกฺกโม (สี, สฺยา, อิ)} ติ.}}

\csromanmatikaanchor{mtk.95}
\csromantocmarkref{section}{สุขุมจฺฉวิลกฺขณํ (12)}{mtk.95}
\bookmark[dest={mtk.95},level=1]{สุขุมจฺฉวิลกฺขณํ (12)}
\csromanheader{1}{\textbf{สุขุมจฺฉวิลกฺขณํ (12)}}
\proseitem{216}{ยมฺปิ ภิกฺขเว ตถาคโต ปุริมํ ชาติํ ปุริมํ ภวํ ปุริมํ นิเกตํ ปุพฺเพ มนุสฺสภูโต สมาโน สมณํ วา พฺราหฺมณํ วา อุปสงฺกมิตฺวา ปริปุจฺฉิตา อโหสิ “กิํ ภนฺเต กุสลํ กิํ อกุสลํ, กิํ สาวชฺชํ กิํ อนวชฺชํ, กิํ เสวิตพฺพํ กิํ น เสวิตพฺพํ, กิํ เม กรียมานํ ทีฆรตฺตํ อหิตาย ทุกฺขาย อสฺส, กิํ วา ปน เม กรียมานํ ทีฆรตฺตํ หิตาย}

\vspace{\tipitakaparskip}
\csromancenter{ลกฺขณสุตฺต สุขาย อสฺสา”ติ.\csromanspacer{} โส ตสฺส กมฺมสฺส กฏตฺตา ฯเปฯ.\csromanspacer{} โส ตโต จุโต อิตฺถตฺตํ อาคโต สมาโน อิมํ มหาปุริสลกฺขณํ ปฏิลภติ.\csromanspacer{} สุขุมจฺฉวิ โหติ, สุขุมตฺตา ฉวิยา รโชชลฺลํ กาเย น อุปลิมฺปติ.}
\prose{\csromanfolio{129}โส เตน ลกฺขเณน สมนฺนาคโต สเจ อคารํ อชฺฌาวสติ, ราชา โหติ จกฺกวตฺตี ฯเปฯ.\csromanspacer{} ราชา สมาโน กิํ ลภติ.\csromanspacer{} มหาปญฺโญ โหติ, นาสฺส โหติ โกจิ ปญฺญาย สทิโส วา เสฏฺโฐ วา กามโภคีนํ.\csromanspacer{} ราชา สมาโน อิทํ ลภติ ฯเปฯ.\csromanspacer{} พุทฺโธ สมาโน กิํ ลภติ.\csromanspacer{} มหาปญฺโญ โหติ ปุถุปญฺโญ หาสปญฺโญ\footnote{หาสุปญฺโญ (สี, อิ)} ชวนปญฺโญ ติกฺขปญฺโญ นิพฺเพธิกปญฺโญ, นาสฺส โหติ โกจิ ปญฺญาย สทิโส วา เสฏฺโฐ วา สพฺพสตฺตานํ.\csromanspacer{} พุทฺโธ สมาโน อิทํ ลภติ.\csromanspacer{} เอตมตฺถํ ภควา อโวจ.}

\csromanheader{2}{217. ตตฺเถตํ วุจฺจติ\csromandash{}}
\csromangathameasure{ปุเร ปุรตฺถา ปุริมาสุ ชาติสุ, \\ อญฺญาตุกาโม ปริปุจฺฉิตา อหุ. \\ สุสฺสูสิตา ปพฺพชิตํ อุปาสิตา, \\ อตฺถนฺตโร อตฺถกถํ นิสามยิ. \\ ปญฺญาปฏิลาภคเตน กมฺมุนา, \\ มนุสฺสภูโต สุขุมจฺฉวี อหุ. \\ พฺยากํสุ อุปฺปาทนิมิตฺตโกวิทา, \\ สุขุมานิ อตฺถานิ อเวจฺจ ทกฺขิติ. \\ สเจ น ปพฺพชฺชมุเปติ ตาทิโส, \\ วตฺเตติ จกฺกํ ปถวิํ ปสาสติ. \\ อตฺถานุสิฏฺฐีสุ ปริคฺคเหสุ จ, \\ น เตน เสยฺโย สทิโส จ วิชฺชติ. \\ สเจ จ ปพฺพชฺชมุเปติ ตาทิโส, \\ เนกฺขมฺมฉนฺทาภิรโต วิจกฺขโณ. \\ ปญฺญาวิสิฏฺฐํ ลภเต อนุตฺตรํ, \\ ปปฺโปติ โพธิํ วรภูริเมธโสติ.}
\csromangathagroup{\csromangathastanza{ปุเร ปุรตฺถา ปุริมาสุ ชาติสุ, \\ อญฺญาตุกาโม ปริปุจฺฉิตา อหุ. \\ สุสฺสูสิตา ปพฺพชิตํ อุปาสิตา, \\ อตฺถนฺตโร อตฺถกถํ นิสามยิ.}\csromangathastanza{ปญฺญาปฏิลาภคเตน\footnote{ปญฺญาปฏิลาภกเตน (สี, อิ) ฏีกา โอโลเกตพฺพา.} กมฺมุนา, \\ มนุสฺสภูโต สุขุมจฺฉวี อหุ. \\ พฺยากํสุ อุปฺปาทนิมิตฺตโกวิทา, \\ สุขุมานิ อตฺถานิ อเวจฺจ ทกฺขิติ.}\csromangathastanza{สเจ น ปพฺพชฺชมุเปติ ตาทิโส, \\ วตฺเตติ จกฺกํ ปถวิํ ปสาสติ. \\ อตฺถานุสิฏฺฐีสุ ปริคฺคเหสุ จ, \\ น เตน เสยฺโย สทิโส จ วิชฺชติ.}\csromangathastanza{สเจ จ ปพฺพชฺชมุเปติ ตาทิโส, \\ เนกฺขมฺมฉนฺทาภิรโต วิจกฺขโณ. \\ ปญฺญาวิสิฏฺฐํ ลภเต อนุตฺตรํ, \\ ปปฺโปติ โพธิํ วรภูริเมธโสติ.}}

\csromanmatikaanchor{mtk.96}
\csromantocmarkref{section}{สุวณฺณวณฺณลกฺขณํ (13)}{mtk.96}
\bookmark[dest={mtk.96},level=1]{สุวณฺณวณฺณลกฺขณํ (13)}
\csromanheader{1}{\textbf{สุวณฺณวณฺณลกฺขณํ (13)}}
\proseitem{218}{\csromanfolio{130}ยมฺปิ ภิกฺขเว ตถาคโต ปุริมํ ชาติํ ปุริมํ ภวํ ปุริมํ นิเกตํ ปุพฺเพ มนุสฺสภูโต สมาโน อกฺโกธโน อโหสิ อนุปายาสพหุโล, พหุมฺปิ วุตฺโต สมาโน นาภิสชฺชิ น กุปฺปิ น พฺยาปชฺชิ น ปติตฺถียิ, น โกปญฺจ โทสญฺจ อปฺปจฺจยญฺจ ปาตฺวากาสิ.\csromanspacer{} ทาตา จ อโหสิ สุขุมานํ มุทุกานํ อตฺถรณานํ ปาวุรณานํ\footnote{ปาปุรณานํ (สี, สฺยา, อิ)} โขมสุขุมานํ กปฺปาสิกสุขุมานํ โกเสยฺยสุขุมานํ กมฺพลสุขุมานํ.\csromanspacer{} โส ตสฺส กมฺมสฺส กฏตฺตา อุปจิตตฺตา ฯเปฯ.\csromanspacer{} โส ตโต จุโต อิตฺถตฺตํ อาคโต สมาโน อิมํ มหาปุริสลกฺขณํ ปฏิลภติ.\csromanspacer{} สุวณฺณวณฺโณ โหติ กญฺจนสนฺนิภตฺตโจ.}

\prose{โส เตน ลกฺขเณน สมนฺนาคโต สเจ อคารํ อชฺฌาวสติ, ราชา โหติ จกฺกวตฺตี ฯเปฯ.\csromanspacer{} ราชา สมาโน กิํ ลภติ.\csromanspacer{} ลาภี โหติ สุขุมานํ มุทุกานํ อตฺถรณานํ ปาวุรณานํ โขมสุขุมานํ กปฺปาสิกสุขุมานํ โกเสยฺยสุขุมานํ กมฺพลสุขุมานํ.\csromanspacer{} ราชา สมาโน อิทํ ลภติ ฯเปฯ.\csromanspacer{} พุทฺโธ สมาโน กิํ ลภติ.\csromanspacer{} ลาภี โหติ สุขุมานํ มุทุกานํ อตฺถรณานํ ปาวุรณานํ โขมสุขุมานํ กปฺปาสิกสุขุมานํ โกเสยฺยสุขุมานํ กมฺพลสุขุมานํ.\csromanspacer{} พุทฺโธ สมาโน อิทํ ลภติ.\csromanspacer{} เอตมตฺถํ ภควา อโวจ.}

\csromanheader{2}{219. ตตฺเถตํ วุจฺจติ\csromandash{}}
\csromangathameasure{อกฺโกธญฺจ อธิฏฺฐหิ อทาสิ, \\ ทานญฺจ วตฺถานิ สุขุมานิ สุจฺฉวีนิ. \\ ปุริมตรภเว ฐิโต อภิวิสฺสชิ, \\ มหิมิว สุโร อภิวสฺสํ. \\ ตํ กตฺวาน อิโต จุโต ทิพฺพํ, \\ อุปปชฺชิ สุกตผลวิปากมนุภุตฺวา. \\ กนกตนุสนฺนิโภ อิธาภิภวติ, \\ สุรวรตโรริว อินฺโท.}
\csromangathagroup{\csromangathastanza{อกฺโกธญฺจ อธิฏฺฐหิ อทาสิ\footnote{อทาสิ จ (สี, อิ)}, \\ ทานญฺจ วตฺถานิ สุขุมานิ สุจฺฉวีนิ. \\ ปุริมตรภเว ฐิโต อภิวิสฺสชิ, \\ มหิมิว สุโร อภิวสฺสํ.}\csromangathastanza{ตํ กตฺวาน อิโต จุโต ทิพฺพํ, \\ อุปปชฺชิ\footnote{อุปปชฺช (สี, อิ)} สุกตผลวิปากมนุภุตฺวา. \\ กนกตนุสนฺนิโภ อิธาภิภวติ, \\ สุรวรตโรริว อินฺโท.}}

\csromanheader{2}{7. ลกฺขณสุตฺต}
\csromangathameasure{เคหญฺจาวสติ นโร อปพฺพชฺช, \\ มิจฺฉํ มหติมหิํ อนุสาสติ. \\ ปสยฺหสหิธ สตฺตรตนํ, \\ ปฏิลภติ วิมล สุขุมจฺฉวิํ สุจิญฺจ. \\ ลาภี อจฺฉาทนวตฺถโมกฺขปาวุรณานํ, \\ ภวติ ยทิ อนาคาริยตํ อุเปติ. \\ สหิโต ปุริมกตผลํ อนุภวติ. \\ น ภวติ กตสฺส ปนาโสติ.}
\csromangathagroup{\csromangathastanza{\csromanfolio{131}เคหญฺจาวสติ นโร อปพฺพชฺช, \\ มิจฺฉํ มหติมหิํ อนุสาสติ\footnote{ปสาสติ (สฺยา)}. \\ ปสยฺหสหิธ สตฺตรตนํ\footnote{ปสาสติ (สฺยา)}, \\ ปฏิลภติ วิมล\footnote{ปสาสติ (สฺยา)} สุขุมจฺฉวิํ สุจิญฺจ.}\csromangathastanza{ลาภี อจฺฉาทนวตฺถโมกฺขปาวุรณานํ, \\ ภวติ ยทิ อนาคาริยตํ อุเปติ. \\ สหิโต\footnote{สุหิต (สฺยา), ส หิ (สี, อิ)} ปุริมกตผลํ อนุภวติ. \\ น ภวติ กตสฺส ปนาโสติ.}}

\csromanmatikaanchor{mtk.97}
\csromantocmarkref{section}{โกโสหิตวตฺถคุยฺหลกฺขณํ (14)}{mtk.97}
\bookmark[dest={mtk.97},level=1]{โกโสหิตวตฺถคุยฺหลกฺขณํ (14)}
\csromanheader{1}{\textbf{โกโสหิตวตฺถคุยฺหลกฺขณํ (14)}}
\proseitem{220}{ยมฺปิ ภิกฺขเว ตถาคโต ปุริมํ ชาติํ ปุริมํ ภวํ ปุริมํ นิเกตํ ปุพฺเพ มนุสฺสภูโต สมาโน จิรปฺปนฏฺเฐ สุจิรปฺปวาสิโน ญาติมิตฺเต สุหชฺเช สขิโน สมาเนตา อโหสิ.\csromanspacer{} มาตรมฺปิ ปุตฺเตน สมาเนตา อโหสิ, ปุตฺตมฺปิ มาตรา สมาเนตา อโหสิ, ปิตรมฺปิ ปุตฺเตน สมาเนตา อโหสิ, ปุตฺตมฺปิ ปิตรา สมาเนตา อโหสิ, ภาตรมฺปิ ภาตรา สมาเนตา อโหสิ, ภาตรมฺปิ ภคินิยา สมาเนตา อโหสิ, ภคินิมฺปิ ภาตรา สมาเนตา อโหสิ, ภคินิมฺปิ ภคินิยา สมาเนตา อโหสิ, สมงฺคีกตฺวา\footnote{สมคฺคิํ กตฺวา (สี, สฺยา, อิ)} จ อพฺภนุโมทิตา อโหสิ.\csromanspacer{} โส ตสฺส กมฺมสฺส กฏตฺตา ฯเปฯ.\csromanspacer{} โส ตโต จุโต อิตฺถตฺตํ อาคโต สมาโน อิมํ มหาปุริสลกฺขณํ ปฏิลภติ.\csromanspacer{} โกโสหิตวตฺถคุโยฺห โหติ.}

\prose{โส เตน ลกฺขเณน สมนฺนาคโต สเจ อคารํ อชฺฌาวสติ, ราชา โหติ จกฺกวตฺตี ฯเปฯ.\csromanspacer{} ราชา สมาโน กิํ ลภติ.\csromanspacer{} ปหูตปุตฺโต โหติ, ปโรสหสฺสํ โข ปนสฺส ปุตฺตา ภวนฺติ สูรา วีรงฺครูปา ปรเสนปฺปมทฺทนา.\csromanspacer{} ราชา สมาโน อิทํ ลภติ ฯเปฯ.\csromanspacer{} พุทฺโธ สมาโน กิํ ลภติ.\csromanspacer{} ปหูตปุตฺโต โหติ, อเนกสหสฺสํ โข ปนสฺส ปุตฺตา ภวนฺติ สูรา วีรงฺครูปา ปรเสนปฺปมทฺทนา.\csromanspacer{} พุทฺโธ สมาโน อิทํ ลภติ.\csromanspacer{} เอตมตฺถํ ภควา อโวจ.}

\csromanheader{2}{221. ตตฺเถตํ วุจฺจติ\csromandash{}}
\csromangathameasure{ปุเร ปุรตฺถา ปุริมาสุ ชาติสุ, \\ จิรปฺปนฏฺเฐ สุจิรปฺปวาสิโน. \\ ญาตี สุหชฺเช สขิโน สมานยิ, \\ สมงฺคิกตฺวา อนุโมทิตา อหุ. \\ โส เตน กมฺเมน ทิวํ สมกฺกมิ, \\ สุขญฺจ ขิฑฺฑารติโย จ อนฺวภิ. \\ ตโต จวิตฺวา ปุนราคโต อิธ, \\ โกโสหิตํ วินฺทติ วตฺถฉาทิยํ. \\ ปหูตปุตฺโต ภวตี ตถาวิโธ, \\ ปโรสหสฺสํ จ ภวนฺติ อตฺรชา. \\ สูรา จ วีรา จ อมิตฺตตาปนา, \\ คิหิสฺส ปีติํชนนา ปิยํวทา. \\ พหูตรา ปพฺพชิตสฺส อิริยโต, \\ ภวนฺติ ปุตฺตา วจนานุสาริโน. \\ คิหิสฺส วา ปพฺพชิตสฺส วา ปุน, \\ ตํ ลกฺขณํ ชายติ ตทตฺถโชตกนฺติ.}
\csromangathagroup{\csromangathastanza{\csromanfolio{132}ปุเร ปุรตฺถา ปุริมาสุ ชาติสุ, \\ จิรปฺปนฏฺเฐ สุจิรปฺปวาสิโน. \\ ญาตี สุหชฺเช สขิโน สมานยิ, \\ สมงฺคิกตฺวา อนุโมทิตา อหุ.}\csromangathastanza{โส เตน\footnote{ส เตน (ก)} กมฺเมน ทิวํ สมกฺกมิ, \\ สุขญฺจ ขิฑฺฑารติโย จ อนฺวภิ. \\ ตโต จวิตฺวา ปุนราคโต อิธ, \\ โกโสหิตํ วินฺทติ วตฺถฉาทิยํ.}\csromangathastanza{ปหูตปุตฺโต ภวตี ตถาวิโธ, \\ ปโรสหสฺสํ จ\footnote{ปโรสหสฺสสฺส (สี, อิ)} ภวนฺติ อตฺรชา. \\ สูรา จ วีรา จ\footnote{ปโรสหสฺสสฺส (สี, อิ)} อมิตฺตตาปนา, \\ คิหิสฺส ปีติํชนนา ปิยํวทา.}\csromangathastanza{พหูตรา ปพฺพชิตสฺส อิริยโต, \\ ภวนฺติ ปุตฺตา วจนานุสาริโน. \\ คิหิสฺส วา ปพฺพชิตสฺส วา ปุน, \\ ตํ ลกฺขณํ ชายติ ตทตฺถโชตกนฺติ.}}

\nitthitamruled{ปฐมภาณวาโร นิฏฺฐิโต.}
\csromanmatikaanchor{mtk.98}
\csromantocmarkref{section}{ปริมณฺฑล-อโนนมชณฺณุปริมสนลกฺขณานิ (15-16)}{mtk.98}
\bookmark[dest={mtk.98},level=1]{ปริมณฺฑล-อโนนมชณฺณุปริมสนลกฺขณานิ (15-16)}
\csromanheader{1}{\textbf{ปริมณฺฑส-อโนนมชณฺณุปริมสนลกฺขณานิ (15-16)}}
\proseitem{222}{ยมฺปิ ภิกฺขเว ตถาคโต ปุริมํ ชาติํ ปุริมํ ภวํ ปุริมํ นิเกตํ ปุพฺเพ มนุสฺสภูโต สมาโน มหาชนสงฺคหํ\footnote{มหาชนสงฺคาหกํ (ก)} สเมกฺขมาโน\footnote{สมเปกฺขมาโน (ก)} สมํ ชานาติ สามํ ชานาติ, ปุริสํ ชานาติ ปุริสวิเสสํ ชานาติ “อยมิทมรหติ อยมิทมรหตี”ติ.\csromanspacer{} ตตฺถ ตตฺถ ปุริสวิเสสกโร อโหสิ.\csromanspacer{} โส ตสฺสฺส กมฺมสฺส กฏตฺตาฯเปฯ โส ตโต จุโต}

\vspace{\tipitakaparskip}
\csromancenter{ลกฺขณสุตฺต อิตฺถตฺตํ อาคโต สมาโน อิมานิ เทฺว มหาปุริสลกฺขณานิ ปฏิลภติ.\csromanspacer{} นิโคฺรธปริมณฺฑโล จ โหติ, ฐิตโกเยว จ อโนนมนฺโต อุโภหิ ปาณิตเลหิ ชณฺณุกานิ ปริมสติ ปริมชฺชติ.}
\prose{\csromanfolio{133}โส เตหิ ลกฺขเณหิ สมนฺนาคโต สเจ อคารํ อชฺฌาวสติ, ราชา โหติ จกฺกวตฺตี ฯเปฯ.\csromanspacer{} ราชา สมาโน กิํ ลภติ.\csromanspacer{} อฑฺโฒ โหติ มหทฺธโน มหาโภโค ปหูตชาตรูปรชโต ปหูตวิตฺตูปกรโณ ปหูตธนธญฺโญ ปริปุณฺณโกสโกฏฺฐาคาโร.\csromanspacer{} ราชา สมาโน อิทํ ลภติ ฯเปฯ.\csromanspacer{} พุทฺโธ สมาโน กิํ ลภติ.\csromanspacer{} อฑฺโฒ โหติ มหทฺธโน มหาโภโค.\csromanspacer{} ตสฺสิมานิ ธนานิ โหนฺติ, เสยฺยถิทํ สทฺธาธนํ สีลธนํ หิริธนํ โอตฺตปฺปธนํ สุตธนํ จาคธนํ ปญฺญาธนํ.\csromanspacer{} พุทฺโธ สมาโน อิทํ ลภติ.\csromanspacer{} เอตมตฺถํ ภควา อโวจ.}

\csromanheader{2}{223. ตตฺเถตํ วุจฺจติ\csromandash{}}
\csromangathameasure{ตุลิย ปฏิวิจย จินฺตยิตฺวา, \\ มหาชนสงฺคหนํ สเมกฺขมาโน. \\ อยมิทมรหติ ตตฺถ ตตฺถ, \\ ปุริสวิเสสกโร ปุเร อโหสิ. \\ มหิํ จ ปน ฐิโต อโนนมนฺโต, \\ ผุสติ กเรหิ อุโภหิ ชณฺณุกานิ. \\ มหิรุหปริมณฺฑโล อโหสิ, \\ สุจริตกมฺมวิปากเสสเกน. \\ พหุวิวิธนิมิตฺตลกฺขณญฺญู, \\ อตินิปุณา มนุชา พฺยากริํสุ. \\ พหุวิวิธา คิหีนํ อรหานิ, \\ ปฏิลภติ ทหโร สุสุ กุมาโร. \\ อิธ จ มหีปติสฺส กามโภคี, \\ คิหิปติรูปกา พหู ภวนฺติ. \\ ยทิ จ ชหติ สพฺพกามโภคํ, \\ ลภติ อนุตฺตรํ อุตฺตมธนคฺคนฺติ.}
\csromangathagroup{\csromangathastanza{ตุลิย ปฏิวิจย จินฺตยิตฺวา, \\ มหาชนสงฺคหนํ\footnote{มหาชนํ สงฺคาหกํ (ก)} สเมกฺขมาโน. \\ อยมิทมรหติ ตตฺถ ตตฺถ, \\ ปุริสวิเสสกโร ปุเร อโหสิ.}\csromangathastanza{มหิํ จ ปน\footnote{สมา จ ปน (สฺยา), ส หิ จ ปน (สี, อิ)} ฐิโต อโนนมนฺโต, \\ ผุสติ กเรหิ อุโภหิ ชณฺณุกานิ. \\ มหิรุหปริมณฺฑโล อโหสิ, \\ สุจริตกมฺมวิปากเสสเกน.}\csromangathastanza{พหุวิวิธนิมิตฺตลกฺขณญฺญู, \\ อตินิปุณา มนุชา พฺยากริํสุ. \\ พหุวิวิธา คิหีนํ อรหานิ, \\ ปฏิลภติ ทหโร สุสุ กุมาโร.}\csromangathastanza{อิธ จ มหีปติสฺส กามโภคี, \\ คิหิปติรูปกา พหู ภวนฺติ. \\ ยทิ จ ชหติ สพฺพกามโภคํ, \\ ลภติ อนุตฺตรํ อุตฺตมธนคฺคนฺติ.}}

\csromanmatikaanchor{mtk.99}
\csromantocmarkref{section}{สีหปุพฺพทฺธกายาทิติลกฺขณํ (17-19)}{mtk.99}
\bookmark[dest={mtk.99},level=1]{สีหปุพฺพทฺธกายาทิติลกฺขณํ (17-19)}
\csromanheader{1}{\textbf{สีหปุพฺพทฺธกายาทิติลกฺขณํ (17-19)}}
\proseitem{224}{\csromanfolio{134}ยมฺปิ ภิกฺขเว ตถาคโต ปุริมํ ชาติํ ปุริมํ ภวํ ปุริมํ นิเกตํ ปุพฺเพ มนุสฺสภูโต สมาโน พหุชนสฺส อตฺถกาโม อโหสิ หิตกาโม ผาสุกาโม โยคกฺเขมกาโม “กินฺติเม สทฺธาย วฑฺเฒยฺยุํ, สีเลน วฑฺเฒยฺยุํ, สุเตน วฑฺเฒยฺยุํ\footnote{สุเตน วฑฺเฒยฺยุํ, พุทฺธิยา วฑฺเฒยฺยุํ (สฺยา)}, จาเคน วฑฺเฒยฺยุํ, ธมฺเมน วฑฺเฒยฺยุํ, ปญฺญาย วฑฺเฒยฺยุํ, ธนธญฺเญน วฑฺเฒยฺยุํ, เขตฺตวตฺถุนา วฑฺเฒยฺยุํ, ทฺวิปทจตุปฺปเทหิ วฑฺเฒยฺยุํ, ปุตฺตทาเรหิ วฑฺเฒยฺยุํ, ทาสกมฺมกรโปริเสหิ วฑฺเฒยฺยุํ, ญาตีหิ วฑฺเฒยฺยุํ, มิตฺเตหิ วฑฺเฒยฺยุํ, พนฺธเวหิ วฑฺเฒยฺยุนฺ”ติ.\csromanspacer{} โส ตสฺส กมฺมสฺส กฏตฺตา ฯเปฯ.\csromanspacer{} โส ตโต จุโต อิตฺถตฺตํ อาคโต สมาโน อิมานิ ตีณิ มหาปุริสลกฺขณานิ ปฏิลภติ, สีหปุพฺพทฺธกาโย จ โหติ จิตนฺตรํโส จ สมวฏฺฏกฺขนฺโธ จ.}

\prose{โส เตหิ ลกฺขเณหิ สมนฺนาคโต สเจ อคารํ อชฺฌาวสติ, ราชา โหติ จกฺกวตฺตี ฯเปฯ.\csromanspacer{} ราชา สมาโน กิํ ลภติ.\csromanspacer{} อปริหานธมฺโม โหติ, น ปริหายติ ธนธญฺเญน เขตฺตวตฺถุนา ทฺวิปทจตุปฺปเทหิ ปุตฺตทาเรหิ ทาสกมฺมกรโปริเสหิ ญาตีหิ มิตฺเตหิ พนฺธเวหิ, น ปริหายติ สพฺพสมฺปตฺติยา.\csromanspacer{} ราชา สมาโน อิทํ ลภติ ฯเปฯ.\csromanspacer{} พุทฺโธ สมาโน กิํ ลภติ.\csromanspacer{} อปริหานธมฺโม โหติ, น ปริหายติ สทฺธาย สีเลน สุเตน จาเคน ปญฺญาย, น ปริหายติ สพฺพสมฺปตฺติยา.\csromanspacer{} พุทฺโธ สมาโน อิทํ ลภติ.\csromanspacer{} เอตมตฺถํ ภควา อโวจ.}

\csromanheader{2}{225. ตตฺเถตํ วุจฺจติ\csromandash{}}
\csromangathameasure{สทฺธาย สีเลน สุเตน พุทฺธิยา, \\ จาเคน ธมฺเมน พหูหิ สาธุหิ. \\ ธเนน ธญฺเญน จ เขตฺตวตฺถุนา, \\ ปุตฺเตหิ ทาเรหิ จตุปฺปเทหิ จ. \\ ญาตีหิ มิตฺเตหิ จ พนฺธเวหิ จ, \\ พเลน วณฺเณน สุเขน จูภยํ. \\ กถํ น หาเยยฺยุํ ปเรติ อิจฺฉติ, \\ อตฺถสฺส มิทฺธี จ ปนาภิกงฺขติ.}
\csromangathagroup{\csromangathastanza{สทฺธาย สีเลน สุเตน พุทฺธิยา, \\ จาเคน ธมฺเมน พหูหิ สาธุหิ. \\ ธเนน ธญฺเญน จ เขตฺตวตฺถุนา, \\ ปุตฺเตหิ ทาเรหิ จตุปฺปเทหิ จ.}\csromangathastanza{ญาตีหิ มิตฺเตหิ จ พนฺธเวหิ จ, \\ พเลน วณฺเณน สุเขน จูภยํ. \\ กถํ น หาเยยฺยุํ ปเรติ อิจฺฉติ, \\ อตฺถสฺส มิทฺธี จ\footnote{อิทํ สมิทฺธญฺจ (ก), อทฺธํ สมิทฺธญฺจ (สฺยา)} ปนาภิกงฺขติ.}}

\csromanheader{2}{7. ลกฺขณสุตฺต}
\csromangathameasure{ส สีหปุพฺพทฺธสุสณฺฐิโต อหุ, \\ สมวฏฺฏขนฺโธ จ จิตนฺตรํโส. \\ ปุพฺเพ สุจิณฺเณน กเตน กมฺมุนา, \\ อหานิยํ ปุพฺพนิมิตฺตมสฺส ตํ. \\ คิหีปิ ธญฺเญน ธเนน วฑฺฒติ, \\ ปุตฺเตหิ ทาเรหิ จตุปฺปเทหิ จ. \\ อกิญฺจโน ปพฺพชิโต อนุตฺตรํ, \\ ปปฺโปติ โพธิํ อสหานธมฺมตนฺติ.}
\csromangathagroup{\csromangathastanza{\csromanfolio{135}ส สีหปุพฺพทฺธสุสณฺฐิโต อหุ, \\ สมวฏฺฏขนฺโธ จ จิตนฺตรํโส. \\ ปุพฺเพ สุจิณฺเณน กเตน กมฺมุนา, \\ อหานิยํ ปุพฺพนิมิตฺตมสฺส ตํ.}\csromangathastanza{คิหีปิ ธญฺเญน ธเนน วฑฺฒติ, \\ ปุตฺเตหิ ทาเรหิ จตุปฺปเทหิ จ. \\ อกิญฺจโน ปพฺพชิโต อนุตฺตรํ, \\ ปปฺโปติ โพธิํ อสหานธมฺมตนฺติ\footnote{สมฺโพธิมหานธมฺมตนฺติ (สฺยา, ก) ตีกา โอโลเกตพฺพา.}.}}

\csromanmatikaanchor{mtk.100}
\csromantocmarkref{section}{รสคฺคสคฺคิตาลกฺขณํ (20)}{mtk.100}
\bookmark[dest={mtk.100},level=1]{รสคฺคสคฺคิตาลกฺขณํ (20)}
\csromanheader{1}{\textbf{รสคฺคสคฺคิตาลกฺขณํ (20)}}
\proseitem{226}{ยมฺปิ ภิกฺขเว ตถาคโต ปุริมํ ชาติํ ปุริมํ ภวํ ปุริมํ นิเกตํ ปุพฺเพ มนุสฺสภูโต สมาโน สตฺตานํ อวิเหฐกชาติโก อโหสิ ปาณินา วา เลฑฺฑุนา วา ทณฺเฑน วา สตฺเถน วา.\csromanspacer{} โส ตสฺส กมฺมสฺส กฏตฺตา อุปจิตตฺตา ฯเปฯ.\csromanspacer{} โส ตโต จุโต อิตฺถตฺตํ อาคโต สมาโน อิมํ มหาปุริสลกฺขณํ ปฏิลภติ, รสคฺคสคฺคี โหติ, อุทฺธคฺคาสฺส รสหรณีโย คีวาย ชาตา โหนฺติ สมาภิวาหินิโย\footnote{สมวาหรสหรณิโย (สฺยา)}.}

\prose{โส เตน ลกฺขเณน สมนฺนาคโต สเจ อคารํ อชฺฌาวสติ, ราชา โหติ จกฺกวตฺตี ฯเปฯ.\csromanspacer{} ราชา สมาโน กิํ ลภติ.\csromanspacer{} อปฺปาพาโธ โหติ อปฺปาตงฺโก สมเวปากินิยา คหณิยา สมนฺนาคโต นาติสีตาย นาจฺจุณฺหาย.\csromanspacer{} ราชา สมาโน อิทํ ลภติ ฯเปฯ.\csromanspacer{} พุทฺโธ สมาโน กิํ ลภติ.\csromanspacer{} อปฺปาพาโธ โหติ อปฺปาตงฺโก สมเวปากินิยา คหณิยา สมนฺนาคโต นาติสีตาย นาจฺจุณฺหาย มชฺฌิมาย ปธานกฺขมาย.\csromanspacer{} พุทฺโธ สมาโน อิทํ ลภติ.\csromanspacer{} เอตมตฺถํ ภควา อโวจ.}

\csromanheader{2}{227. ตตฺเถตํ วุจฺจติ\csromandash{}}
\csromangathameasure{น ปาณิทณฺเฑหิ ปนาถ เลฑฺฑุนา, \\ สตฺเถน วา มรณวเธน วา ปน. \\ อุพฺพาธนาย ปริตชฺชนาย วา, \\ น เหฐยี ชนตมเหฐโก อหุ. \\ เตเนว โส สุคติมุเปจฺจ โมทติ, \\ สุขปฺผลํ กริย สุขานิ วินฺทติ. \\ สโมชสา รสหรณี สุสณฺฐิตา, \\ อิธาคโต ลภติ รสคฺคสคฺคิตํ. \\ เตนาหุ นํ อตินิปุณา วิจกฺขณา, \\ อยํ นโร สุขพหุโล ภวิสฺสติ. \\ คิหิสฺส วา ปพฺพชิตสฺส วา ปุน2, \\ ตํ ลกฺขณํ ภวติ ตทตฺถโชตกนฺติ.}
\csromangathagroup{\csromangathastanza{น ปาณิทณฺเฑหิ ปนาถ เลฑฺฑุนา, \\ สตฺเถน วา มรณวเธน\footnote{มารณวเธน (ก)} วา ปน. \\ อุพฺพาธนาย ปริตชฺชนาย วา, \\ น เหฐยี ชนตมเหฐโก อหุ.}\csromangathastanza{\csromanfolio{136}เตเนว โส สุคติมุเปจฺจ โมทติ, \\ สุขปฺผลํ กริย สุขานิ วินฺทติ. \\ สโมชสา\footnote{สมฺปชฺชสา (สี, อิ), ปามุญฺชสา (สฺยา), สามญฺจ สา (ก) 2. ปน (สฺยา) 3. น จ วิสาจิตํ (สี, อิ), น จ วิสาวิ (สฺยา)} รสหรณี สุสณฺฐิตา, \\ อิธาคโต ลภติ รสคฺคสคฺคิตํ.}\csromangathastanza{เตนาหุ นํ อตินิปุณา วิจกฺขณา, \\ อยํ นโร สุขพหุโล ภวิสฺสติ. \\ คิหิสฺส วา ปพฺพชิตสฺส วา ปุน2, \\ ตํ ลกฺขณํ ภวติ ตทตฺถโชตกนฺติ.}}

\csromanmatikaanchor{mtk.101}
\csromantocmarkref{section}{อภินีลเนตฺต-โคปขุมลกฺขณานิ (21-22)}{mtk.101}
\bookmark[dest={mtk.101},level=1]{อภินีลเนตฺต-โคปขุมลกฺขณานิ (21-22)}
\csromanheader{1}{\textbf{อภินีลเนตฺต-โคปขุมลกฺขณานิ (21-22)}}
\proseitem{228}{ยมฺปิ ภิกฺขเว ตถาคโต ปุริมํ ชาติํ ปุริมํ ภวํ ปุริมํ นิเกตํ ปุพฺเพ มนุสฺสภูโต สมาโน น จ วิสฏํ, น จ วิสาจิ3, น จ ปน วิเจยฺย เปกฺขิตา, อุชุํ ตถา ปสฏมุชุมโน, ปิยจกฺขุนา พหุชนํ อุทิกฺขิตา อโหสิ.\csromanspacer{} โส ตสฺส กมฺมสฺส กฏตฺตา ฯเปฯ.\csromanspacer{} โส ตโต จุโต อิตฺถตฺตํ อาคโต สมาโน อิมานิ เทฺว มหาปุริสลกฺขณานิ ปฏิลภติ, อภินีลเนตฺโต จ โหติ โคปขุโม จ.}

\prose{โส เตหิ ลกฺขเณหิ สมนฺนาคโต สเจ อคารํ อชฺฌาวสติ, ราชา โหติ จกฺกวตฺตี ฯเปฯ.\csromanspacer{} ราชา สมาโน กิํ ลภติ.\csromanspacer{} ปิยทสฺสโน โหติ พหุโน ชนสฺส, ปิโย โหติ มนาโป พฺราหฺมณคหปติกานํ เนคมชานปทานํ คณกมหามตฺตานํ อนีกฏฺฐานํ โทวาริกานํ อมจฺจานํ ปาริสชฺชานํ ราชูนํ โภคิยานํ กุมารานํ.\csromanspacer{} ราชา สมาโน อิทํ ลภติ ฯเปฯ.\csromanspacer{} พุทฺโธ สมาโน กิํ ลภติ.\csromanspacer{} ปิยทสฺสโน โหติ พหุโน ชนสฺส, ปิโย โหติ มนาโป ภิกฺขูนํ ภิกฺขุนีนํ อุปาสกานํ อุปสิกานํ เทวานํ มนุสฺสานํ อสุรานํ นาคานํ คนฺธพฺพานํ.\csromanspacer{} พุทฺโธ สมาโน อิทํ ลภติ.\csromanspacer{} เอตมตฺถํ ภควา อโวจ.}

\csromanheader{2}{229. ตตฺเถตํ วุจฺจติ\csromandash{}}
\csromangathameasure{น จ วิสฏํ น จ วิสาจิ, น จ ปน วิเจยฺยเปกฺขิตา. \\ อุชุํ ตถา ปสฏมุชุมโน, ปิยจกฺขุนา พหุชนํ อุทิกฺขิตา.}
\csromangathagroup{\csromangathastanza{น จ วิสฏํ น จ วิสาจิ\footnote{วิสาจิตํ (สี, อิ), วิสาวิ (สฺยา)}, น จ ปน วิเจยฺยเปกฺขิตา. \\ อุชุํ ตถา ปสฏมุชุมโน, ปิยจกฺขุนา พหุชนํ อุทิกฺขิตา.}}

\csromanheader{2}{7. ลกฺขณสุตฺต}
\csromangathameasure{สุคตีสุ โส ผลวิปากํ, \\ อนุภวติ ตตฺถ โมทติ. \\ อิธ จ ปน ภวติ โคปขุโม, \\ อภินีลเนตฺตนยโน สุทสฺสโน. \\ อภิโยคิโน จ นิปุณา, \\ พหู ปน นิมิตฺตโกวิทา. \\ สุขุมนยนกุสลา มนุชา, \\ ปิยทสฺสโนติ อภินิทฺทิสนฺติ นํ. \\ ปิยทสฺสโน คิหีปิ สนฺโต จ, \\ ภวติ พหุชนปิยายิโต. \\ ยทิ จ น ภวติ คิหี สมโณ โหติ, \\ ปิโย พหูนํ โสกนาสโนติ.}
\csromangathagroup{\csromangathastanza{\csromanfolio{137}สุคตีสุ โส ผลวิปากํ, \\ อนุภวติ ตตฺถ โมทติ. \\ อิธ จ ปน ภวติ โคปขุโม, \\ อภินีลเนตฺตนยโน สุทสฺสโน.}\csromangathastanza{อภิโยคิโน จ นิปุณา, \\ พหู ปน นิมิตฺตโกวิทา. \\ สุขุมนยนกุสลา มนุชา, \\ ปิยทสฺสโนติ อภินิทฺทิสนฺติ นํ.}\csromangathastanza{ปิยทสฺสโน คิหีปิ สนฺโต จ, \\ ภวติ พหุชนปิยายิโต. \\ ยทิ จ น ภวติ คิหี สมโณ โหติ, \\ ปิโย พหูนํ โสกนาสโนติ.}}

\csromanmatikaanchor{mtk.102}
\csromantocmarkref{section}{อุณฺหีสสีสลกฺขณํ (23)}{mtk.102}
\bookmark[dest={mtk.102},level=1]{อุณฺหีสสีสลกฺขณํ (23)}
\csromanheader{1}{\textbf{อุณฺหีสสีสลกฺขณํ (23)}}
\proseitem{230}{ยมฺปิ ภิกฺขเว ตถาคโต ปุริมํ ชาติํ ปุริมํ ภวํ ปุริมํ นิเกตํ ปุพฺเพ มนุสฺสภูโต สมาโน พหุชนปุพฺพงฺคโม อโหสิ กุสเลสุ ธมฺเมสุ พหุชนปาโมกฺโข กายสุจริเต วจีสุจริเต มโนสุจริเต ทานสํวิภาเค สีลสมาทาเน อุโปสถุปวาเส มตฺเตยฺยตาย เปตฺเตยฺยตาย สามญฺญตาย พฺรหฺมญฺญตาย กุเล เชฏฺฐาปจายิตาย อญฺญตรญฺญตเรสุ จ อธิกุสเลสุ ธมฺเมสุ.\csromanspacer{} โส ตสฺส กมฺมสฺส กฏตฺตา ฯเปฯ.\csromanspacer{} โส ตโต จุโต อิตฺถตฺตํ อาคโต สมาโน อิมํ มหาปุริสลกฺขณํ ปฏิลภติ, อุณฺหีสสีโส โหติ.}

\prose{โส เตน ลกฺขเณน สมนฺนาคโต สเจ อคารํ อชฺฌาวสติ, ราชา โหติ จกฺกวตฺตี ฯเปฯ.\csromanspacer{} ราชา สมาโน กิํ ลภติ.\csromanspacer{} มหาสฺส ชโน อนฺวายิโก โหติ, พฺราหฺมณคหปติกา เนคมชานปทา คณกมหามตฺตา อนีกฏฺฐ โทวาริกา อมจฺจา ปาริสชฺชา ราชาโน โภคิยา กุมารา.\csromanspacer{} ราชา สมาโน อิทํ ลภติ ฯเปฯ.\csromanspacer{} พุทฺโธ สมาโน กิํ ลภติ.\csromanspacer{} มหาสฺส ชโน อนฺวายิโก โหติ, ภิกฺขู ภิกฺขุนิโย อุปาสกา อุปาสิกาโย เทวา มนุสฺสา อสุรา นาคา คนฺธพฺพา.\csromanspacer{} พุทฺโธ สมาโน อิทํ ลภติ.\csromanspacer{} เอตมตฺถํ ภควา อโวจ.}

\csromanheader{2}{231. ตตฺเถตํ วุจฺจติ\csromandash{}}
\csromangathameasure{ปุพฺพงฺคโม สุจริเตสุ อหุ, \\ ธมฺเมสุ ธมฺมจริยาภิรโต. \\ อนฺวายิโก พหุชนสฺส อหุ, \\ สคฺเคสุ เวทยิตฺถ ปุญฺญผลํ. \\ เวทิตฺวา โส สุจริตสฺส ผลํ, \\ อุณฺหีสสีสตฺตมิธชฺฌคมา. \\ พฺยากํสุ พฺยญฺชนนิมิตฺตธรา, \\ ปุพฺพงฺคโม พหุชนํ เหสฺสติ. \\ ปฏิโภคิยา มนุเชสุ อิธ, \\ ปุพฺเพว ตสฺส อภิหรนฺติ ตทา. \\ ยทิ ขตฺติโย ภวติ ภูมิปติ, \\ ปฏิหารกํ พหุชเน ลภติ. \\ อถ เจปิ ปพฺพชติ โส มนุโช, \\ ธมฺเมสุ โหติ ปคุโณ วิสวี. \\ ตสฺสานุสาสนิคุณาภิรโต, \\ อนฺวายิโก พหุชโน ภวตีติ.}
\csromangathagroup{\csromangathastanza{\csromanfolio{138}ปุพฺพงฺคโม สุจริเตสุ อหุ, \\ ธมฺเมสุ ธมฺมจริยาภิรโต. \\ อนฺวายิโก พหุชนสฺส อหุ, \\ สคฺเคสุ เวทยิตฺถ ปุญฺญผลํ.}\csromangathastanza{เวทิตฺวา โส สุจริตสฺส ผลํ, \\ อุณฺหีสสีสตฺตมิธชฺฌคมา. \\ พฺยากํสุ พฺยญฺชนนิมิตฺตธรา, \\ ปุพฺพงฺคโม พหุชนํ เหสฺสติ.}\csromangathastanza{ปฏิโภคิยา มนุเชสุ อิธ, \\ ปุพฺเพว ตสฺส อภิหรนฺติ ตทา. \\ ยทิ ขตฺติโย ภวติ ภูมิปติ, \\ ปฏิหารกํ พหุชเน ลภติ.}\csromangathastanza{อถ เจปิ ปพฺพชติ โส มนุโช, \\ ธมฺเมสุ โหติ ปคุโณ วิสวี. \\ ตสฺสานุสาสนิคุณาภิรโต, \\ อนฺวายิโก พหุชโน ภวตีติ.}}

\csromanmatikaanchor{mtk.103}
\csromantocmarkref{section}{เอเกกโลมตา-อุณฺณาลกฺขณานิ (24-25)}{mtk.103}
\bookmark[dest={mtk.103},level=1]{เอเกกโลมตา-อุณฺณาลกฺขณานิ (24-25)}
\csromanheader{1}{\textbf{เอเกกโลมตา-อุณฺณาลกฺขณานิ (24-25)}}
\proseitem{232}{ยมฺปิ ภิกฺขเว ตถาคโต ปุริมํ ชาติํ ปุริมํ ภวํ ปุริมํ นิเกตํ ปุพฺเพ มนุสฺสภูโต สมาโน มุสาวาทํ ปหาย มุสาวาทา ปฏิวิรโต อโหสิ, สจฺจวาที สจฺจสนฺโธ เถโต ปจฺจยิโก อวิสํวาทโก โลกสฺส.\csromanspacer{} โส ตสฺส กมฺมสฺส กฏตฺตา อุปจิตตฺตา ฯเปฯ.\csromanspacer{} โส ตโต จุโต อิตฺถตฺตํ อาคโต สมาโน อิมานิ เทฺว มหาปุริสลกฺขณานิ ปฏิลภติ, เอเกกโลโม จ โหติ, อุณฺณา จ ภมุกนฺตเร ชาตา โหติ โอทาตา มุทุตูลสนฺนิภา.}

\prose{โส เตหิ ลกฺขเณหิ สมนฺนาคโต สเจ อคารํ อชฺฌาวสติ, ราชา โหติ จกฺกวตฺตี ฯเปฯ.\csromanspacer{} ราชา สมาโน กิํ ลภติ.\csromanspacer{} มหาสฺส ชโน อุปวตฺตติ, พฺราหฺมณคหปติกา เนคมชานปทา คณกมหามตฺตา อนีกฏฺฐา โทวาริกา อมจฺจา ปาริสชฺชา ราชาโน โภคิยา กุมารา.}

\vspace{\tipitakaparskip}
\csromancenter{ลกฺขณสุตฺต ราชา สมาโน อิทํ ลภติ ฯเปฯ.\csromanspacer{} พุทฺโธ สมาโน กิํ ลภติ.\csromanspacer{} มหาสฺส ชโน อุปวตฺตติ, ภิกฺขู ภิกฺขุนิโย อุปาสกา อุปาสิกาโย เทวา มนุสฺสา อสุรา นาคา คนฺธพฺพา.\csromanspacer{} พุทฺโธ สมาโน อิทํ ลภติ.\csromanspacer{} เอตมตฺถํ ภควา อโวจ.}
\csromanheader{2}{233. ตตฺเถตํ วุจฺจติ\csromandash{}}
\csromangathameasure{สจฺจปฺปฏิญฺโญ ปุริมาสุ ชาติสุ, \\ อเทฺวชฺฌวาโจ อลิกํ วิวชฺชยิ. \\ น โส วิสํวาทยิตาปิ กสฺสจิ, \\ ภูเตน ตจฺเฉน ตเถน ภาสยิ. \\ เสตา สุสุกฺกา มุทุตูลสนฺนิภา, \\ อุณฺณา สุชาตา ภมุกนฺตเร อหุ. \\ น โลมกูเปสุ ทุเว อชายิสุํ, \\ เอเกกโลมูปจิตงฺควา อหุ. \\ ตํ ลกฺขณญฺญู พหโว สมาคตา, \\ พฺยากํสุ อุปฺปาทนิมิตฺตโกวิทา. \\ อุณฺณา จ โลมา จ ยถา สุสณฺฐิตา, \\ อุปวตฺตตี อีทิสกํ พหุชฺชโน. \\ คิหิมฺปิ สนฺตํ อุปวตฺตตี ชโน, \\ พหุ ปุรตฺถาปกเตน กมฺมุนา. \\ อกิญฺจนํ ปพฺพชิตํ อนุตฺตรํ, \\ พุทฺธํปิ สนฺตํ อุปวตฺตติ ชโนติ.}
\csromangathagroup{\csromangathastanza{\csromanfolio{139}สจฺจปฺปฏิญฺโญ ปุริมาสุ ชาติสุ, \\ อเทฺวชฺฌวาโจ อลิกํ วิวชฺชยิ. \\ น โส วิสํวาทยิตาปิ กสฺสจิ, \\ ภูเตน ตจฺเฉน ตเถน ภาสยิ\footnote{โตสยิ (สี, อิ)}.}\csromangathastanza{เสตา สุสุกฺกา มุทุตูลสนฺนิภา, \\ อุณฺณา สุชาตา\footnote{อุณฺณาสฺส ชาตา (ก-สี)} ภมุกนฺตเร อหุ. \\ น โลมกูเปสุ ทุเว อชายิสุํ, \\ เอเกกโลมูปจิตงฺควา อหุ.}\csromangathastanza{ตํ ลกฺขณญฺญู พหโว สมาคตา, \\ พฺยากํสุ อุปฺปาทนิมิตฺตโกวิทา. \\ อุณฺณา จ โลมา จ ยถา สุสณฺฐิตา, \\ อุปวตฺตตี อีทิสกํ พหุชฺชโน.}\csromangathastanza{คิหิมฺปิ สนฺตํ อุปวตฺตตี ชโน, \\ พหุ ปุรตฺถาปกเตน กมฺมุนา. \\ อกิญฺจนํ ปพฺพชิตํ อนุตฺตรํ, \\ พุทฺธํปิ สนฺตํ อุปวตฺตติ ชโนติ.}}

\csromanmatikaanchor{mtk.104}
\csromantocmarkref{section}{จตฺตาลีส-อวิรฬทนฺตลกฺขณานิ (26-27)}{mtk.104}
\bookmark[dest={mtk.104},level=1]{จตฺตาลีส-อวิรฬทนฺตลกฺขณานิ (26-27)}
\csromanheader{1}{\textbf{จตฺตาลีส-อวิรฬทนฺตลกฺขณานิ (26-27)}}
\proseitem{234}{ยมฺปิ ภิกฺขเว ตถาคโต ปุริมํ ชาติํ ปุริมํ ภวํ ปุริมํ นิเกตํ ปุพฺเพ มนุสฺสภูโต สมาโน ปิสุณํ วาจํ ปหาย ปิสุณาย วาจาย ปฏิวิรโต อโหสิ, อิโต สุตฺวา น อมุตฺร อกฺขาตา อิเมสํ เภทาย, อมุตฺร วา สุตฺวา น อิเมสํ อกฺขาตา อมูสํ \csromanfolio{140}เภทาย, อิติ ภินฺนานํ วา สนฺธาตา, สหิตานํ วา อนุปฺปทาตา, สมคฺคาราโม สมคฺครโต สมคฺคนนฺที สมคฺคกรณิํ วาจํ ภาสิตา อโหสิ.\csromanspacer{} โส ตสฺส กมฺมสฺส กฏตฺตา ฯเปฯ.\csromanspacer{} โส ตโต จุโต อิตฺถตฺตํ อาคโต สมาโน อิมานิ เทฺว มหาปุริสลกฺขณานิ ปฏิลภติ, จตฺตาลีสทนฺโต จ โหติ อวิรฬทนฺโต จ.}

\prose{โส เตหิ ลกฺขเณหิ สมนฺนาคโต สเจ อคารํ อชฺฌาวสติ, ราชา โหติ จกฺกวตฺตี ฯเปฯ.\csromanspacer{} ราชา สมาโน กิํ ลภติ.\csromanspacer{} อเภชฺชปริโส โหติ, อเภชฺชาสฺส โหนฺติ ปริสา, พฺราหฺมณคหปติกา เนคมชานปทา คณกมหามตฺตา อนีกฏฺฐา โทวาริกา อมจฺจา ปาริสชฺชา ราชาโน โภคิยา กุมารา.\csromanspacer{} ราชา สมาโน อิทํ ลภติ ฯเปฯ.\csromanspacer{} พุทฺโธ สมาโน กิํ ลภติ.\csromanspacer{} อเภชฺชปริโส โหติ, อเภชฺชาสฺส โหนฺติ ปริสา, ภิกฺขู ภิกฺขุนิโย อุปาสกา อุปาสิกาโย เทวา มนุสฺสา อสุรา นาคา คนฺธพฺพา.\csromanspacer{} พุทฺโธ สมาโน อิทํ ลภติ.\csromanspacer{} เอตมตฺถํ ภควา อโวจ.}

\csromanheader{2}{235. ตตฺเถตํ วุจฺจติ\csromandash{}}
\csromangathameasure{เวภูติยํ สหิตเภทการิํ, \\ เภทปฺปวฑฺฒนวิวาทการิํ. \\ กลหปฺปวฑฺฒน-อกิจฺจการิํ, \\ สหิตานํ เภทชนนิํ น ภณิ. \\ อวิวาทวฑฺฒนกริํ สุคิตํ, \\ ภินฺนานุสนฺธิชนนิํ อภณิ. \\ กลหํ ชนสฺส ปนุที สมงฺคี, \\ สหิเตหิ นนฺทติ ปโมทติ จ. \\ สุคตีสุ โส ผลวิปากํ, \\ อนุภวติ ตตฺถ โมทติ. \\ ทนฺตา อิธ โหนฺติ อวิรฬา สหิตา, \\ จตุโร ทสสฺส มุขชา สุสณฺฐิตา.}
\csromangathagroup{\csromangathastanza{เวภูติยํ สหิตเภทการิํ, \\ เภทปฺปวฑฺฒนวิวาทการิํ. \\ กลหปฺปวฑฺฒน-อกิจฺจการิํ, \\ สหิตานํ เภทชนนิํ น ภณิ.}\csromangathastanza{อวิวาทวฑฺฒนกริํ สุคิตํ, \\ ภินฺนานุสนฺธิชนนิํ อภณิ. \\ กลหํ ชนสฺส ปนุที สมงฺคี, \\ สหิเตหิ นนฺทติ ปโมทติ จ.}\csromangathastanza{สุคตีสุ โส ผลวิปากํ, \\ อนุภวติ ตตฺถ โมทติ. \\ ทนฺตา อิธ โหนฺติ อวิรฬา สหิตา, \\ จตุโร ทสสฺส มุขชา สุสณฺฐิตา.}}

\csromanheader{2}{7. ลกฺขณสุตฺต}
\csromangathameasure{ยทิ ขตฺติโย ภวติ ภูมิปติ, \\ อวิเภทิยาสฺส ปริสา ภวติ. \\ สมโณ จ โหติ วิรโช วิมโล, \\ ปริสาสฺส โหติ อนุคตา อจลาติ.}
\csromangathagroup{\csromangathastanza{\csromanfolio{141}ยทิ ขตฺติโย ภวติ ภูมิปติ, \\ อวิเภทิยาสฺส ปริสา ภวติ. \\ สมโณ จ โหติ วิรโช วิมโล, \\ ปริสาสฺส โหติ อนุคตา อจลาติ.}}

\csromanmatikaanchor{mtk.105}
\csromantocmarkref{section}{ปหูตชิวฺหา-พฺรหฺมสฺสรลกฺขณานิ (28-29)}{mtk.105}
\bookmark[dest={mtk.105},level=1]{ปหูตชิวฺหา-พฺรหฺมสฺสรลกฺขณานิ (28-29)}
\csromanheader{1}{\textbf{ปหูตชิวฺหา-พฺรหฺมสฺสรลกฺขณานิ (28-29)}}
\proseitem{236}{ยมฺปิ ภิกฺขเว ตถาคโต ปุริมํ ชาติํ ปุริมํ ภวํ ปุริมํ นิเกตํ ปุพฺเพ มนุสฺสภูโต สมาโน ผรุสํ วาจํ ปหาย ผรุสาย วาจาย ปฏิวิรโต อโหสิ.\csromanspacer{} ยา สา วาจา เนลา กณฺณสุขา เปมนียา หทยงฺคมา โปรี พหุชนกนฺตา พหุชนมนาปา, ตถารูปิํ วาจํ ภาสิตา อโหสิ.\csromanspacer{} โส ตสฺส กมฺมสฺส กฏตฺตา อุปจิตตฺตา ฯเปฯ.\csromanspacer{} โส ตโต จุโต อิตฺถตฺตํ อาคโต สมาโน อิมานิ เทฺว มหาปุริสลกฺขณานิ ปฏิลภติ, ปหูตชิโวฺห จ โหติ พฺรหฺมสฺสโร จ กรวีกภาณี.}

\prose{โส เตหิ ลกฺขเณหิ สมนฺนาคโต สเจ อคารํ อชฺฌวสติ, ราชา โหติ จกฺกวตฺตี ฯเปฯ.\csromanspacer{} ราชา สมาโน กิํ ลภติ.\csromanspacer{} อาเทยฺยวาโจ โหติ, อาทิยนฺติสฺส วจนํ พฺราหฺมณคหปติกา เนคมชานปทา คณกมหามตฺตา อนีกฏฺฐา โทวาริกา อมจฺจา ปาริสชฺชา ราชาโน โภคิยา กุมารา.\csromanspacer{} ราชา สมาโน อิทํ ลภติ ฯเปฯ.\csromanspacer{} พุทฺโธ สมาโน กิํ ลภติ.\csromanspacer{} อาเทยฺยวาโจ โหติ, อาทิยนฺตสฺส วจนํ ภิกฺขู ภิกฺขุนิโย อุปาสกา อุปาสิกาโย เทวา มนุสฺสา อสุรา นาคา คนฺธพฺพา.\csromanspacer{} พุทฺโธ สมาโน อิทํ ลภติ.\csromanspacer{} เอตมตฺตํ ภควา อโวจ.}

\csromanheader{2}{237. ตตฺเถตํ วุจฺจติ\csromandash{}}
\csromangathameasure{อกฺโกสภณฺฑนวิเหสการิํ, \\ อุพฺพาธิกํ พหุชนปฺปมทฺทนํ. \\ อพาฬฺหํ คิรํ โส น ภณิ ผรุสํ, \\ มธุรํ ภณิ สุสํหิตํ สขิลํ. \\ มนโส ปิยา หทยคามินิโย, \\ วาจา โส เอรยติ กณฺณสุขา. \\ วาจาสุจิณฺณผลมนุภวิ, \\ สคฺเคสุ เวทยถ ปุญฺญผลํ. \\ เวทิตฺวา โส สุจริตสฺส ผลํ, \\ พฺรหฺมสฺสรตฺตมิธมชฺฌคมา. \\ ชิวฺหาสฺส โหติ วิปุลา ปุถุลา, \\ อาเทยฺยวากฺยวจโน ภวติ. \\ คิหิโนปิ อิชฺฌติ ยถา ภณโต, \\ อถ เจ ปพฺพชติ โส มนุโช. \\ อาทิยนฺติสฺส วจนํ ชนตา, \\ พหุโน พหุํ สุภณิตํ ภณโตติ.}
\csromangathagroup{\csromangathastanza{อกฺโกสภณฺฑนวิเหสการิํ, \\ อุพฺพาธิกํ\footnote{อุพฺพาธกรํ (สฺยา)} พหุชนปฺปมทฺทนํ. \\ อพาฬฺหํ คิรํ โส น ภณิ ผรุสํ, \\ มธุรํ ภณิ สุสํหิตํ\footnote{อุพฺพาธกรํ (สฺยา)} สขิลํ.}\csromangathastanza{\csromanfolio{142}มนโส ปิยา หทยคามินิโย, \\ วาจา โส เอรยติ กณฺณสุขา. \\ วาจาสุจิณฺณผลมนุภวิ, \\ สคฺเคสุ เวทยถ\footnote{เวทยติ (?) ฏีกา โอโลเกตพฺพา.} ปุญฺญผลํ.}\csromangathastanza{เวทิตฺวา โส สุจริตสฺส ผลํ, \\ พฺรหฺมสฺสรตฺตมิธมชฺฌคมา. \\ ชิวฺหาสฺส โหติ วิปุลา ปุถุลา, \\ อาเทยฺยวากฺยวจโน ภวติ.}\csromangathastanza{คิหิโนปิ อิชฺฌติ ยถา ภณโต, \\ อถ เจ ปพฺพชติ โส มนุโช. \\ อาทิยนฺติสฺส วจนํ ชนตา, \\ พหุโน พหุํ สุภณิตํ ภณโตติ.}}

\csromanmatikaanchor{mtk.106}
\csromantocmarkref{section}{สีหหนุลกฺขณํ (30)}{mtk.106}
\bookmark[dest={mtk.106},level=1]{สีหหนุลกฺขณํ (30)}
\csromanheader{1}{\textbf{สีหหนุลกฺขณํ (30)}}
\proseitem{238}{ยมฺปิ ภิกฺขเว ตถาคโต ปุริมํ ชาติํ ปุริมํ ภวํ ปุริมํ นิเกตํ ปุพฺเพ มนุสฺสภูโต สมาโน สมฺผปฺปลาปํ ปหาย สมฺผปฺปลาปา ปฏิวิรโต อโหสิ กาลวาที ภูตวาที อตฺถวาที ธมฺมวาที วินยวาที, นิธานวติํ วาจํ ภาสิตา อโหสิ กาเลน สาปเทสํ ปริยนฺตวติํ อตฺตสํหิตํ.\csromanspacer{} โส ตสฺส กมฺมสฺส กฏตฺตา ฯเปฯ.\csromanspacer{} โส ตโต จุโต อิตฺถตฺตํ อาคโต สมาโน อิมํ มหาปุริสลกฺขณํ ปฏิลภติ, สีหหนุ โหติ.}

\prose{โส เตน ลกฺขเณน สมนฺนาคโต สเจ อคารํ อชฺฌาวสติ, ราชา โหติ จกฺกวตฺตี ฯเปฯ.\csromanspacer{} ราชา สมาโน กิํ ลภติ.\csromanspacer{} อปฺปธํสิโย โหติ เกนจิ มนุสฺสภูเตน ปจฺจตฺถิเกน ปจฺจามิตฺเตน.\csromanspacer{} ราชา สมาโน อิทํ ลภติ ฯเปฯ.\csromanspacer{} พุทฺโธ สมาโน กิํ ลภติ.\csromanspacer{} อปฺปธํสิโย โหติ อพฺภนฺตเรหิ วา พาหิเรหิ วา ปจฺจตฺถิเกหิ ปจฺจามิตฺเตหิ, ราเคน วา โทเสน วา โมเหน วา สมเณน วา พฺราหฺมเณน วา เทเวน วา มาเรน วา พฺรหฺมุนา วา เกนจิ วา โลกสฺมิํ.\csromanspacer{} พุทฺโธ สมาโน อิทํ ลภติ.\csromanspacer{} เอตมตฺถํ ภควา อโวจ.}

\csromanheader{2}{7. ลกฺขณสุตฺต}
\csromanheader{2}{239. ตตฺเถตํ วุจฺจติ\csromandash{}}
\csromangathameasure{น สมฺผปฺปลาปํ น มุทฺธตํ, \\ อวิกิณฺณวจนพฺยปฺปโถ อโหสิ. \\ อหิตมปิ จ อปนุทิ, \\ หิตมปิ จ พหุชนสุขญฺจ อภณิ. \\ ตํ กตฺวา อิโต จุโต ทิวมุปปชฺชิ, \\ สุกตผลวิปากมนุโภสิ. \\ จวิย ปุนริธาคโต สมาโน, \\ ทฺวิทุคมวรตรหนุตฺตมลตฺถ. \\ ราชา โหติ สุทุปฺปธํสิโย, \\ มนุชินฺโท มนุชาธิปติ มหานุภาโว. \\ ติทิวปุรวรสโม ภวติ, \\ สุรวรตโรริว อินฺโท. \\ คนฺธพฺพาสุรยกฺขรกฺขเสภิ, \\ สุเรหิ น หิ ภวติ สุปฺปธํสิโย. \\ ตถตฺโต ยทิ ภวติ ตถาวิโธ, \\ อิธ ทิสา จ ปฏิทิสา จ วิทิสา จาติ.}
\csromangathagroup{\csromangathastanza{\csromanfolio{143}น สมฺผปฺปลาปํ น มุทฺธตํ\footnote{พุทฺธตนฺติ (ก)}, \\ อวิกิณฺณวจนพฺยปฺปโถ อโหสิ. \\ อหิตมปิ จ อปนุทิ, \\ หิตมปิ จ พหุชนสุขญฺจ อภณิ.}\csromangathastanza{ตํ กตฺวา อิโต จุโต ทิวมุปปชฺชิ, \\ สุกตผลวิปากมนุโภสิ. \\ จวิย ปุนริธาคโต สมาโน, \\ ทฺวิทุคมวรตรหนุตฺตมลตฺถ.}\csromangathastanza{ราชา โหติ สุทุปฺปธํสิโย, \\ มนุชินฺโท มนุชาธิปติ มหานุภาโว. \\ ติทิวปุรวรสโม ภวติ, \\ สุรวรตโรริว อินฺโท.}\csromangathastanza{คนฺธพฺพาสุรยกฺขรกฺขเสภิ\footnote{สุรสกฺกรกฺขเสภิ (สฺยา)}, \\ สุเรหิ น หิ ภวติ สุปฺปธํสิโย. \\ ตถตฺโต ยทิ ภวติ ตถาวิโธ, \\ อิธ ทิสา จ ปฏิทิสา จ วิทิสา จาติ.}}

\csromanmatikaanchor{mtk.107}
\csromantocmarkref{section}{สมสนฺต-สุสุกฺกทาฐาลกฺขณานิ (31-32)}{mtk.107}
\bookmark[dest={mtk.107},level=1]{สมสนฺต-สุสุกฺกทาฐาลกฺขณานิ (31-32)}
\csromanheader{1}{\textbf{สมทนฺต-สุสุกฺกทาฐาลกฺขณานิ (31-32)}}
\proseitem{240}{ยมฺปิ ภิกฺขเว ตถาคโต ปุริมํ ชาติํ ปุริมํ ภวํ ปุริมํ นิเกตํ ปุพฺเพ มนุสฺสภูโต สมาโน มิจฺฉาชีวํ ปหาย สมฺมา-อาชีเวน ชีวิกํ กปฺเปสิ, ตุลากูฏ กํสกูฏ มานกูฏ อุกฺโกฏน วญฺจน นิกติ สาจิโยค เฉทน วธ พนฺธน วิปราโมส อาโลป สหสาการา\footnote{สาหสาการา (สี, สฺยา, อิ)} ปฏิวิรโต อโหสิ.\csromanspacer{} โส ตสฺส กมฺมสฺส กฏตฺตา อุปจิตตฺตา อุสฺสนฺนตฺตา วิปุลตฺตา กายสฺส เภทา ปรํ มรณา สุคติํ สคฺคํ โลกํ อุปปชฺชติ.\csromanspacer{} โส ตตฺถ อญฺเญ เทเว ทสหิ ฐาเนหิ อธิคณฺหาติ ทิพฺเพน อายุนา ทิพฺเพน วณฺเณน ทิพฺเพน สุเขน ทิพฺเพน ยเสน ทิพฺเพน อาธิปเตยฺเยน ทิพฺเพหิ รูเปหิ ทิพฺเพหิ สทฺเทหิ ทิพฺเพหิ คนฺเธหิ ทิพฺเพหิ \csromanfolio{144}รเสหิ ทิพฺเพหิ โผฏฺฐพฺเพหิ.\csromanspacer{} โส ตโต จุโต อิตฺถตฺตํ อาคโต สมาโน อิมานิ เทฺว มหาปุริสลกฺขณานิ ปฏิลภติ, สมทนฺโต จ โหติ สุสุกฺกทาโฐ จ.}

\prose{โส เตหิ ลกฺขเณหิ สมนฺนาคโต สเจ อคารํ อชฺฌาวสติ, ราชา โหติ จกฺกวตฺตี ธมฺมิโก ธมฺมราชา จาตุรนฺโต วิชิตาวี ชนปทตฺถาวริยปฺปตฺโต สตฺตรตนสมนฺนาคโต.\csromanspacer{} ตสฺสิมานิ สตฺต รตนานิ ภวนฺติ, เสยฺยถิทํ, จกฺกรตนํ หตฺถิรตนํ อสฺสรตนํ มณิรตนํ อิตฺถิรตนํ คหปติรตนํ ปริณายกรตนเมว สตฺตมํ.\csromanspacer{} ปโรสหสฺสํ โข ปนสฺส ปุตฺตา ภวนฺติ สูรา วีรงฺครูปา ปรเสนปฺปมทฺทนา.\csromanspacer{} โส อิมํ ปถวิํ สาครปริยนฺตํ อขิลมนิมิตฺตมกณฺฏกํ อิทฺธํ ผีตํ เขมํ สิวํ นิรพฺพุทํ อทณฺเฑน อสตฺเถน ธมฺเมน อภิวิชิย อชฺฌาวสติ.\csromanspacer{} ราชา สมาโน กิํ ลภติ.\csromanspacer{} สุจิปริวาโร โหติ สุจิสฺส โหนฺติ ปริวารา พฺราหฺมณคหปติกา เนคมชานปทา คณกมหามตฺตา อนีกฏฺฐา โทวาริกา อมจฺจา ปาริสชฺชา ราชาโน โภคิยา กุมารา.\csromanspacer{} ราชา สมาโน อิทํ ลภติ.}

\prose{สเจ โข ปน อคารสฺมา อนคาริยํ ปพฺพชติ, อรหํ โหติ สมฺมาสมฺพุทฺโธ โลเก วิวฏฺฏจฺฉโท.\csromanspacer{} พุทฺโธ สมาโน กิํ ลภติ.\csromanspacer{} สุจิปริวาโร โหติ, สุจิสฺส โหนฺติ ปริวารา, ภิกฺขู ภิกฺขุนิโย อุปาสกา อุปาสิกาโย เทวา มนุสฺสา อสุรา นาคา คนฺธพฺพา.\csromanspacer{} พุทฺโธ สมาโน อิทํ ลภติ.\csromanspacer{} เอตมตฺถํ ภควา อโวจ.}

\csromanheader{2}{241. ตตฺเถตํ วุจฺจติ\csromandash{}}
\csromangathameasure{มิจฺฉาชีวญฺจ อวสฺสชิ สเมน วุตฺติํ, \\ สุจินา โส ชนยิตฺถ ธมฺมิเกน. \\ อหิตมปิ จ อปนุทิ, \\ หิตมปิ จ พหุชนสุขญฺจ อจริ. \\ สคฺเค เวทยติ นโร สุขปฺผลานิ, \\ กริตฺวา นิปุเณภิ วิทูหิ สพฺภิ. \\ วณฺณิตานิ ติทิวปุรวรสโม, \\ อภิรมติ รติขิฑฺฑาสมงฺคี.}
\csromangathagroup{\csromangathastanza{มิจฺฉาชีวญฺจ อวสฺสชิ สเมน วุตฺติํ, \\ สุจินา โส ชนยิตฺถ ธมฺมิเกน. \\ อหิตมปิ จ อปนุทิ, \\ หิตมปิ จ พหุชนสุขญฺจ อจริ.}\csromangathastanza{สคฺเค เวทยติ นโร สุขปฺผลานิ, \\ กริตฺวา นิปุเณภิ วิทูหิ สพฺภิ. \\ วณฺณิตานิ ติทิวปุรวรสโม, \\ อภิรมติ รติขิฑฺฑาสมงฺคี.}}

\csromanheader{2}{7. ลกฺขณสุตฺต}
\csromangathameasure{ลทฺธาน มานุสกํ ภวํ ตโต, \\ จวิตฺวาน สุกตผลวิปากํ. \\ เสสเกน ปฏิลภติ ลปนชํ, \\ สมมปิ สุจิสุสุกฺกํ. \\ ตํ เวยฺยญฺชนิกา สมาคตา พหโว, \\ พฺยากํสุ นิปุณสมฺมตา มนุชา. \\ สุจิชนปริวารคโณ ภวติ, \\ ทิชสมสุกฺกสุจิโสภนทนฺโต. \\ รญฺโญ โหติ พหุชโน, \\ สุจิปริวาโร มหติํ มหิํ อนุสาสโต. \\ ปสยฺห น จ ชนปทตุทนํ, \\ หิตมปิ จ พหุชน สุขญฺจ จรนฺติ. \\ อถ เจ ปพฺพชติ ภวติ วิปาโป, \\ สมโณ สมิตรโช วิวฏฺฏจฺฉโท. \\ วิคตทรถกิลมโถ, \\ อิมมปิ จ ปรมปิ จ ปสฺสติ โลกํ. \\ ตสฺโสวาทกรา พหุคิหี จ ปพฺพชิตา จ, \\ อสุจิํ ครหิตํ ธุนนฺติ ปาปํ. \\ ส หิ สุจิภิ ปริวุโต ภวติ, \\ มลขิลกลิกิเลเส ปนุเทหีติ.}
\csromangathagroup{\csromangathastanza{\csromanfolio{145}ลทฺธาน มานุสกํ ภวํ ตโต, \\ จวิตฺวาน สุกตผลวิปากํ. \\ เสสเกน ปฏิลภติ ลปนชํ, \\ สมมปิ สุจิสุสุกฺกํ\footnote{ลทฺธาน มนุสฺสกํ ภวํ ตโต จวิย, ปุน สุกตผลวิปากเสสเกน. ปฏิลภติ ลปนชํ สมมปิ, สุจิ จ สุจิสุทฺธสุสุกฺกํ. (สฺยา)}.}\csromangathastanza{ตํ เวยฺยญฺชนิกา สมาคตา พหโว, \\ พฺยากํสุ นิปุณสมฺมตา มนุชา. \\ สุจิชนปริวารคโณ ภวติ, \\ ทิชสมสุกฺกสุจิโสภนทนฺโต.}\csromangathastanza{รญฺโญ โหติ พหุชโน, \\ สุจิปริวาโร มหติํ มหิํ อนุสาสโต. \\ ปสยฺห น จ ชนปทตุทนํ, \\ หิตมปิ จ พหุชน สุขญฺจ จรนฺติ.}\csromangathastanza{อถ เจ ปพฺพชติ ภวติ วิปาโป, \\ สมโณ สมิตรโช วิวฏฺฏจฺฉโท. \\ วิคตทรถกิลมโถ, \\ อิมมปิ จ ปรมปิ จ\footnote{อิมมฺปิ จ ปรมฺปิ จ (อิ), ปรํปิ ปรมํปิ จ (สฺยา)} ปสฺสติ โลกํ.}\csromangathastanza{ตสฺโสวาทกรา พหุคิหี จ ปพฺพชิตา จ, \\ อสุจิํ ครหิตํ ธุนนฺติ ปาปํ. \\ ส หิ สุจิภิ ปริวุโต ภวติ, \\ มลขิลกลิกิเลเส ปนุเทหีติ\footnote{ตสฺโสวาทกรา พหุคิหี จ, ปพฺพชิตา จ อสุจิวิครหิต. ปนุทิปาปสฺส หิ สุจิภิปริวุโต, ภวติ มลขิลกกิเลเส ปนุเทติ. (สฺยา)}.}}

\prose{อิทมโวจ ภควา.\csromanspacer{} อตฺตมนา เต ภิกฺขู ภควโต ภาสิตํ อภินนฺทุนฺติ.}

\nitthitam{\textbf{ลกฺขณสุตฺตํ นิฏฺฐิตํ สตฺตมํ.}}
\csromanmatikaanchor{mtk.108}
\csromantocmarknopage{chapter}{8. สิงฺคาลสุตฺต}{mtk.108}
\bookmark[dest={mtk.108},level=0]{8. สิงฺคาลสุตฺต}
\chapterheadpageread{8. สิงฺคาลสุตฺต}
\csromanmatikaanchor{mtk.109}
\csromantocmarkref{section}{ฉทิสา}{mtk.109}
\bookmark[dest={mtk.109},level=1]{ฉทิสา}
\proseitem{242}{\csromanfolio{146}เอวํ เม สุตํ\csromandash{}เอกํ สมยํ ภควา ราชคเห วิหรติ เวฬุวเน กลนฺทกนิวาเป.\csromanspacer{} เตน โข ปน สมเยน สิงฺคาลโก\footnote{สิคาลโก (สี)} คหปติปุตฺโต กาลสฺเสว อุฏฺฐาย ราชคหา นิกฺขมิตฺวา อลฺลวตฺโถ อลฺลเกโส ปญฺชลิโก ปุถุทิสา\footnote{ปุถุทฺทิสา (สี, สฺยา, อิ)} นมสฺสติ ปุรตฺถิมํ ทิสํ ทกฺขิณํ ทิสํ ปจฺฉิมํ ทิสํ อุตฺตรํ ทิสํ เหฏฺฐิมํ ทิสํ อุปริมํ ทิสํ.}

\proseitem{243}{อถ โข ภควา ปุพฺพณฺหสมยํ นิวาเสตฺวา ปตฺตจีวรมาทาย ราชคหํ ปิณฺฑาย ปาวิสิ.\csromanspacer{} อทฺทสา โข ภควา สิงฺคาลกํ คหปติปุตฺตํ กาลสฺเสว วุฏฺฐาย ราชคหา นิกฺขมิตฺวา อลฺลวตฺถํ อลฺลเกสํ ปญฺชลิกํ ปุถุทิสา นมสฺสนฺตํ ปุรตฺถิมํ ทิสํ ทกฺขิณํ ทิสํ ปจฺฉิมํ ทิสํ อุตฺตรํ ทิสํ เหฏฺฐิมํ ทิสํ อุปริมํ ทิสํ, ทิสฺวา สิงฺคาลกํ คหปติปุตฺตํ เอตทโวจ “กิํ นุ โข ตฺวํ คหปติปุตฺต กาลสฺเสว วุฏฺฐาย ราชคหา นิกฺขมิตฺวา อลฺลวตฺโถ อลฺลเกโส ปญฺชลิโก ปุถุทิสา นมสฺสสิ ปุรตฺถิมํ ทิสํ ทกฺขิณํ ทิสํ ปจฺฉิมํ ทิสํ อุตฺตรํ ทิสํ เหฏฺฐิมํ ทิสํ อุปริมํ ทิสนฺ”ติ.\csromanspacer{} ปิตา มํ ภนฺเต กาลํ กโรนฺโต เอวํ อวจ “ทิสา ตาต นมสฺเสยฺยาสี”ติ.\csromanspacer{} โส โข อหํ ภนฺเต ปิตุวจนํ สกฺกโรนฺโต ครุํ กโรนฺโต มาเนนฺโต ปูเชนฺโต กาลสฺเสว วุฏฺฐาย ราชคหา นิกฺขมิตฺวา อลฺลวตฺโถ อลฺลเกโส ปญฺชลิโก ปุถุทิสา นมสฺสามิ ปุรตฺถิมํ ทิสํ ทกฺขิณํ ทิสํ ปจฺฉิมํ ทิสํ อุตฺตรํ ทิสํ เหฏฺฐิมํ ทิสํ อุปริมํ ทิสนฺติ.}

\csromanheader{2}{\textbf{ฉ ทิสา}}
\proseitem{244}{น โข คหปติปุตฺต อริยสฺส วินเย เอวํ ฉ ทิสา\footnote{ฉทฺทิสา (สี, อิ)} นมสฺสิตพฺพาติ.\csromanspacer{} ยถา กถํ ปน ภนฺเต อริยสฺส วินเย ฉ ทิสา3 นมสฺสิตพฺพา, สาธุ เม ภนฺเต ภควา ตถา ธมฺมํ เทเสตุ, ยถา อริยสฺส วินเย ฉ ทิสา3 นมสฺสิตพฺพาติ.}

\prose{เตน หิ คหปติปุตฺต สุโณหิ สาธุกํ มนสิกโรหิ ภาสิสฺสามีติ. “เอวํ ภนฺเต”ติ โข สิงฺคาลโก คหปติปุตฺโต ภควโต ปจฺจสฺโสสิ.\csromanspacer{} ภควา เอตทโวจ\csromandash{}}

\csromanheader{2}{8. สิงฺคาลสุตฺต}
\prose{\csromanfolio{147}ยโต โข คหปติปุตฺต อริยสาวกสฺส จตฺตาโร กมฺมกิเลสา ปหีนา โหนฺติ, จตูหิ จ ฐาเนหิ ปาปกมฺมํ น กโรติ, ฉ จ โภคานํ อปายมุขานิ น เสวติ.\csromanspacer{} โส เอวํ จุทฺทส ปาปกาปคโต ฉทฺทิสาปฏิจฺฉาที\footnote{ปฏิจฺฉาที โหติ (สฺยา)} อุโภโลกวิชยาย ปฏิปนฺโน โหติ.\csromanspacer{} ตสฺส อยญฺจว โลโก อารทฺโธ โหติ ปโร จ โลโก, โส กายสฺส เภทา ปรํ มรณา สุคติํ สคฺคํ โลกํ อุปปชฺชติ.}

\csromanmatikaanchor{mtk.110}
\csromantocmarkref{section}{จตฺตาโร กมฺมกิเลสา}{mtk.110}
\bookmark[dest={mtk.110},level=1]{จตฺตาโร กมฺมกิเลสา}
\csromanheader{1}{\textbf{จตฺตาโรกมฺมกิเลสา}}
\proseitem{245}{กตมสฺส จตฺตาโร กมฺมกิเลสา ปหีนา โหนฺติ.\csromanspacer{} ปาณาติปาโต โข คหปติปุตฺต กมฺมกิเลโส, อทินฺนาทานํ กมฺมกิเลโส, กาเมสุมิจฺฉาจาโร กมฺมกิเลโส, มุสาวาโท กมฺมกิเลโส.\csromanspacer{} อิมสฺส จตฺตาโร กมฺมกิเลสา ปหีนา โหนฺตีติ.\csromanspacer{} อิทมโวจ ภควา, อิทํ วตฺวาน\footnote{อิทํ วตฺวา (สี, อิ) เอวมีทิเสสุ ฐาเนสุ.} สุคโต อถาปรํ เอตทโวจ สตฺถา\csromandash{}}

\csromangathameasure{ปาณาติปาโต อทินฺนาทานํ, มุสาวาโท จ วุจฺจติ. \\ ปรทารคมนญฺเจว, นปฺปสํสนฺติ ปณฺฑิตาติ.}
\csromangathagroup{\csromangathastanza{ปาณาติปาโต อทินฺนาทานํ, มุสาวาโท จ วุจฺจติ. \\ ปรทารคมนญฺเจว, นปฺปสํสนฺติ ปณฺฑิตาติ.}}

\csromanmatikaanchor{mtk.111}
\csromantocmarkref{section}{จตุฐาณํ}{mtk.111}
\bookmark[dest={mtk.111},level=1]{จตุฐาณํ}
\csromanheader{1}{\textbf{จตุฐานํ}}
\proseitem{246}{กตเมหิ จตูหิ ฐาเนหิ ปาปกมฺมํ น กโรติ.\csromanspacer{} ฉนฺทาคติํ คจฺฉนฺโต ปาปกมฺมํ กโรติ, โทสาคติํ คจฺฉนฺโต ปาปกมฺมํ กโรติ, โมหาคติํ คจฺฉนฺโต ปาปกมฺมํ กโรติ, ภยาคติํ คจฺฉนฺโต ปาปกมฺมํ กโรติ.\csromanspacer{} ยโต โข คหปติปุตฺต อริยสาวโก เนว ฉนฺทาคติํ คจฺฉติ, น โทสาคติํ คจฺฉติ, น โมหาคติํ คจฺฉติ, น ภยาคติํ คจฺฉติ.\csromanspacer{} อิเมหิ จตูหิ ฐาเนหิ ปาปกมฺมํ น กโรตีติ.\csromanspacer{} อิทมโวจ ภควา, อิทํ วตฺวาน สุคโต อถาปรํ เอตทโวจ สตฺถา\csromandash{}}

\csromangathameasure{ฉนฺทา โทสา ภยา โมหา, โย ธมฺมํ อติวตฺตติ. \\ นิหียติ ยโส ตสฺส, กาฬปกฺเขว จนฺทิมา. \\ ฉนฺทา โทสา ภยา โมหา, โย ธมฺมํ นาติวตฺตติ. \\ อาปูรติ ยโส ตสฺส3, สุกฺกปกฺเขว จนฺทิมาติ.}
\csromangathagroup{\csromangathastanza{ฉนฺทา โทสา ภยา โมหา, โย ธมฺมํ อติวตฺตติ. \\ นิหียติ ยโส ตสฺส\footnote{ตสฺส ยโส (พหูสุ, วินเยปิ)}, กาฬปกฺเขว จนฺทิมา.}\csromangathastanza{ฉนฺทา โทสา ภยา โมหา, โย ธมฺมํ นาติวตฺตติ. \\ อาปูรติ ยโส ตสฺส3, สุกฺกปกฺเขว\footnote{ชุณฺหปกฺเขว (ก)} จนฺทิมาติ.}}

\csromanmatikaanchor{mtk.112}
\csromantocmarkref{section}{ฉ-อปายมุขานิ}{mtk.112}
\bookmark[dest={mtk.112},level=1]{ฉ-อปายมุขานิ}
\csromanheader{1}{\textbf{ฉ อปายมุขานิ}}
\proseitem{247}{\csromanfolio{148}กตมานิ ฉ โภคานํ อปายมุขานิ น เสวติ.\csromanspacer{} สุราเมรยมชฺชปฺปมาทฏฺฐานานุโยโค โข คหปติปุตฺต โภคานํ อปายมุขํ, วิกาลวิสิขาจริยานุโยโค โภคานํ อปายมุขํ, สมชฺชภิจรณํ โภคานํ อปายมุขํ, ชูตปฺปมาทฏฺฐานานุโยโค โภคานํ อปายมุขํ, ปาปมิตฺตานุโยโค โภคานํ อปายมุขํ, อาลสฺยานุโยโค\footnote{อาลสฺสานุโยโค (สี, สฺยา, อิ)} โภคานํ อปายมุขํ.}

\csromanmatikaanchor{mtk.113}
\csromantocmarkref{section}{สุราเมรยสฺส ฉ อาทีนวา}{mtk.113}
\bookmark[dest={mtk.113},level=1]{สุราเมรยสฺส ฉ อาทีนวา}
\csromanheader{1}{\textbf{สุราเมรยสฺส ฉ อาทีนวา}}
\proseitem{248}{ฉ โขเม คหปติปุตฺต อาทีนวา สุราเมรยมชฺชปฺปมาทฏฺฐานานุโยเค.\csromanspacer{} สนฺทิฏฺฐิกา ธนชานิ\footnote{ธนญฺชานิ (สี, อิ)}, กลหปฺปวฑฺฒนี, โรคานํ อายตนํ, อกิตฺติสญฺชนนี, โกปีนนิทํสนี, ปญฺญาย ทุพฺพลิกรณีเตฺวว ฉฏฺฐํ ปทํ ภวติ.\csromanspacer{} อิเม โข คหปติปุตฺต ฉ อาทีนวา สุราเมรยมชฺชปฺปมาทฏฺฐานานุโยเค.}

\csromanmatikaanchor{mtk.114}
\csromantocmarkref{section}{วิกาลจริยาย ฉ อาทีนวา}{mtk.114}
\bookmark[dest={mtk.114},level=1]{วิกาลจริยาย ฉ อาทีนวา}
\csromanheader{1}{\textbf{วิกาลจริยาย ฉ อาทีนวา}}
\proseitem{249}{ฉ โขเม คหปติปุตฺต อาทีนวา วิกาลวิสิขาจริยานุโยเค.\csromanspacer{} อตฺตาปิสฺส อคุตฺโต อรกฺขิโต โหติ, ปุตฺตทาโรปิสฺส อคุตฺโต อรกฺขิโต โหติ, สาปเตยฺยํปิสฺส อคุตฺตํ อรกฺขิตํ โหติ, สงฺกิโย จ โหติ ปาปเกสุ ฐาเนสุ\footnote{เตสุ เตสุ ฐาเนสุ (สฺยา)}, อภูตวจนํ จ ตสฺมิํ รูหติ, พหูนญฺจ ทุกฺขธมฺมานํ ปุรกฺขโต โหติ.\csromanspacer{} อิเม โข คหปติปุตฺต ฉ อาทีนวา วิกาลวิสิขาจริยานุโยเค.}

\csromanmatikaanchor{mtk.115}
\csromantocmarkref{section}{สมชฺชาภิจรณสฺส ฉ อาทีนวา}{mtk.115}
\bookmark[dest={mtk.115},level=1]{สมชฺชาภิจรณสฺส ฉ อาทีนวา}
\csromanheader{1}{\textbf{สมชฺชาภิจรณสฺส ฉ อาทีนวา}}
\proseitem{250}{ฉ โขเม คหปติปุตฺต อาทีนวา สมชฺชาภิจรเณ.\csromanspacer{} กฺว\footnote{กุวํ (ก-สี, อิ)} นจฺจํ, กฺว คีตํ, กฺว วาทิตํ, กฺว อกฺขานํ, กฺว ปาณิสฺสรํ, กฺว กุมฺภถุนนฺติ.\csromanspacer{} อิเม โข คหปติปุตฺต ฉ อาทีนวา สมชฺชาภิจรเณ.}

\csromanmatikaanchor{mtk.116}
\csromantocmarkref{section}{ชูตปฺปมาทสฺส ฉ อาทีนวา}{mtk.116}
\bookmark[dest={mtk.116},level=1]{ชูตปฺปมาทสฺส ฉ อาทีนวา}
\csromanheader{1}{8. สิงฺคาลสุตฺต \textbf{ชูตปฺปมาทสฺส ฉ อาทีนวา}}
\proseitem{251}{\csromanfolio{149}ฉ โขเม คหปติปุตฺต อาทีนวา ชูตปฺปมาทฏฺฐานานุโยเค.\csromanspacer{} ชยํ เวรํ ปสวติ, ชิโน วิตฺตมนุโสจติ, สนฺทิฏฺฐิกา ธนชานิ, สภาคตสฺส\footnote{สภาเย ตสฺส (ก)} วจนํ น รูหติ, มิตฺตามจฺจานํ ปริภูโต โหติ, อาวาหวิวาหกานํ อปตฺถิโต โหติ “อกฺขธุตฺโต อยํ ปุริสปุคฺคโล นาลํ ทารภรณายา”ติ.\csromanspacer{} อิเม โข คหปติปุตฺต ฉ อาทีนวา ชูตปฺปมาทฏฺฐานานุโยเค.}

\csromanmatikaanchor{mtk.117}
\csromantocmarkref{section}{ปาปมิตฺตตาย ฉ อาทีนวา}{mtk.117}
\bookmark[dest={mtk.117},level=1]{ปาปมิตฺตตาย ฉ อาทีนวา}
\csromanheader{1}{\textbf{ปาปมิตฺตตาย ฉ อาทีนวา}}
\proseitem{252}{ฉ โขเม คหปติปุตฺต อาทีนวา ปาปมิตฺตานุโยเค.\csromanspacer{} เย ธุตฺตา, เย โสณฺฑา, เย ปิปาสา, เย เนกติกา, เย วญฺจนิกา, เย สาหสิกา.\csromanspacer{} ตฺยาสฺส มิตฺตา โหนฺติ เต สหายา.\csromanspacer{} อิเม โข คหปติปุตฺต ฉ อาทีนาวา ปาปมิตฺตานุโยเค.}

\csromanmatikaanchor{mtk.118}
\csromantocmarkref{section}{อาลสฺยสฺส ฉ อาทีนวา}{mtk.118}
\bookmark[dest={mtk.118},level=1]{อาลสฺยสฺส ฉ อาทีนวา}
\csromanheader{1}{\textbf{อาลสฺยสฺส ฉ อาทีนวา}}
\proseitem{253}{ฉ โขเม คหปติปุตฺต อาทีนวา อาลสฺยานุโยเค.\csromanspacer{} อติสีตนฺติ กมฺมํ น กโรติ, อติ-อุณฺหนฺติ กมฺมํ น กโรติ, อติสายนฺติ กมฺมํ น กโรติ, อติปาโตติ กมฺมํ น กโรติ, อติฉาโตสฺมีติ กมฺมํ น กโรติ, อติธาโตสฺมีติ กมฺมํ น กโรติ.\csromanspacer{} ตสฺส เอวํ กิจฺจาปเทสพหุลสฺส วิหรโต อนุปฺปนฺนา เจว โภคา นุปฺปชฺชนฺติ, อุปฺปนฺนา จ โภคา ปริกฺขยํ คจฺฉนฺติ.\csromanspacer{} อิเม โข คหปติปุตฺต ฉ อาทีนวา อาลสฺยานุโยเคติ.\csromanspacer{} อิทมโวจ ภควา, อิทํ วตฺวาน สุคโต อถาปรํ เอตทโวจ สตฺถา\csromandash{}}

\csromangathameasure{“โหติ ปานสขา นาม, \\ โหติ สมฺมิยสมฺมิโย. \\ โย จ อตฺเถสุ ชาเตสุ, \\ สหาโย โหติ โส สขา. \\ อุสฺสูรเสยฺยา ปรทารเสวนา, \\ เวรปฺปสโว จ อนตฺถตา จ. \\ ปาปา จ มิตฺตา สุกทริยตา จ, \\ เอเต ฉ ฐานา ปุริสํ ธํสยนฺติ. \\ ปาปมิตฺโต ปาปสโข, \\ ปาป-อาจารโคจโร. \\ อสฺมา โลกา ปรมฺหา จ, \\ อุภยา ธํสเต นโร. \\ อกฺขิตฺถิโย วารุณี นจฺจคีตํ, \\ ทิวา โสปฺปํ ปาริจริยา อกาเล. \\ ปาปา จ มิตฺตา สุกทริยตา จ, \\ เอเต ฉ ฐานา ปุริสํ ธํสยนฺติ. \\ อกฺเขหิ ทิพฺพนฺติ สุรํ ปิวนฺติ, \\ ยนฺติตฺถิโย ปาณสมา ปเรสํ. \\ นิหีนเสวี น จ วุทฺธเสวี, \\ นิหียเต กาฬปกฺเขว จนฺโท. \\ โย วารุณี อทฺธโน อกิญฺจโน, \\ ปิปาโส ปิวํ ปปาคโต. \\ อุทกมิว อิณํ วิคาหติ, \\ อกุลํ กาหิติ ขิปฺปมตฺตโน. \\ น ทิวา โสปฺปสีเลน, รตฺติมุฏฺฐานเทสฺสินา. \\ นิจฺจํ มตฺเตน โสณฺเฑน, สกฺกา อาวสิตุํ ฆรํ. \\ อติสีตํ อติ-อุณฺหํ, อติสายมิทํ อหุ. \\ อิติ วิสฺสฏฺฐากมฺมนฺเต, อตฺถา อจฺเจนฺติ มาณเว.}
\csromangathagroup{\csromangathastanza{“โหติ ปานสขา นาม, \\ โหติ สมฺมิยสมฺมิโย. \\ โย จ อตฺเถสุ ชาเตสุ, \\ สหาโย โหติ โส สขา.}\csromangathastanza{\csromanfolio{150}อุสฺสูรเสยฺยา ปรทารเสวนา, \\ เวรปฺปสโว\footnote{เวรปฺปสงฺโค (สี, สฺยา, อิ)} จ อนตฺถตา จ. \\ ปาปา จ มิตฺตา สุกทริยตา จ, \\ เอเต ฉ ฐานา ปุริสํ ธํสยนฺติ.}\csromangathastanza{ปาปมิตฺโต ปาปสโข, \\ ปาป-อาจารโคจโร. \\ อสฺมา โลกา ปรมฺหา จ, \\ อุภยา ธํสเต นโร.}\csromangathastanza{อกฺขิตฺถิโย วารุณี นจฺจคีตํ, \\ ทิวา โสปฺปํ ปาริจริยา อกาเล. \\ ปาปา จ มิตฺตา สุกทริยตา จ, \\ เอเต ฉ ฐานา ปุริสํ ธํสยนฺติ.}\csromangathastanza{อกฺเขหิ ทิพฺพนฺติ สุรํ ปิวนฺติ, \\ ยนฺติตฺถิโย ปาณสมา ปเรสํ. \\ นิหีนเสวี น จ วุทฺธเสวี\footnote{วุทฺธิเสวี (สฺยา), พุทฺธิเสวี (ก)}, \\ นิหียเต กาฬปกฺเขว จนฺโท.}\csromangathastanza{โย วารุณี อทฺธโน อกิญฺจโน, \\ ปิปาโส ปิวํ ปปาคโต\footnote{วิปาโสสิ อตฺถปาคโต (สฺยา), ปิปาโสปิ สมปฺปปาคโต (ก)}. \\ อุทกมิว อิณํ วิคาหติ, \\ อกุลํ\footnote{วิปาโสสิ อตฺถปาคโต (สฺยา), ปิปาโสปิ สมปฺปปาคโต (ก)} กาหิติ ขิปฺปมตฺตโน. \\ น ทิวา โสปฺปสีเลน, รตฺติมุฏฺฐานเทสฺสินา\footnote{รตฺตินุฏฺฐานทสฺสินา (สี, อิ), รตฺตินุฏฺฐานสีลินา (?)}. \\ นิจฺจํ มตฺเตน โสณฺเฑน, สกฺกา อาวสิตุํ ฆรํ.}\csromangathastanza{อติสีตํ อติ-อุณฺหํ, อติสายมิทํ อหุ. \\ อิติ วิสฺสฏฺฐากมฺมนฺเต, อตฺถา อจฺเจนฺติ มาณเว.}}

\csromanheader{2}{8. สิงฺคาลสุตฺต}
\csromangathameasure{โยธ สีตญฺจ อุณฺหญฺจ, ติณา ภิยฺโย น มญฺญติ. \\ กรํ ปุริสกิจฺจานิ, โส สุขํ น วิหายตี”ติ.}
\csromangathagroup{\csromangathastanza{\csromanfolio{151}โยธ สีตญฺจ อุณฺหญฺจ, ติณา ภิยฺโย น มญฺญติ. \\ กรํ ปุริสกิจฺจานิ, โส สุขํ\footnote{สุขา (สพฺพตฺถ) อฏฺฐกถา โอโลเกตพฺพา.} น วิหายตี”ติ.}}

\csromanmatikaanchor{mtk.119}
\csromantocmarkref{section}{มิตฺตปติรูปก}{mtk.119}
\bookmark[dest={mtk.119},level=1]{มิตฺตปติรูปก}
\csromanheader{1}{\textbf{มิตฺตปติรูปก}}
\proseitem{254}{จตฺตาโรเม คหปติปุตฺต อมิตฺตา มิตฺตปติรูปกา เวทิตพฺพา.\csromanspacer{} อญฺญทตฺถุหโร อมิตฺโต มิตฺตปติรูปโก เวทิตพฺโพ, วจีปรโม อมิตฺโต มิตฺตปติรูปโก เวทิตพฺโพ, อนุปฺปิยภาณี อมิตฺโต มิตฺตปติรูปโก เวทิตพฺโพ, อปายสหาโย อมิตฺโต มิตฺตปติรูปโก เวทิตพฺโพ.}

\proseitem{255}{จตูหิ โข คหปติปุตฺต ฐาเนหิ อญฺญทตฺถุหโร อมิตฺโต มิตฺตปติรูปโก เวทิตพฺโพ.}

\csromangathameasure{อญฺญทตฺถุหโร โหติ, อปฺเปน พหุมิจฺฉติ. \\ ภยสฺส กิจฺจํ กโรติ, เสวติ อตฺถการณา.}
\csromangathagroup{\csromangathastanza{อญฺญทตฺถุหโร โหติ, อปฺเปน พหุมิจฺฉติ. \\ ภยสฺส กิจฺจํ กโรติ, เสวติ อตฺถการณา.}}

\prose{อิเมหิ โข คหปติปุตฺต จตูหิ ฐาเนหิ อญฺญทตฺถุหโร อมิตฺโต มิตฺตปติรูปโก เวทิตพฺโพ.}

\proseitem{256}{จตูหิ โข คหปติปุตฺต ฐาเนหิ วจีปรโม อมิตฺโต มิตฺตปติรูปโก เวทิตพฺโพ.\csromanspacer{} อตีเตน ปฏิสนฺถรติ\footnote{ปฏิสนฺธรติ (ก)}, อนาคเตน ปฏิสนฺถรติ, นิรตฺถเกน สงฺคณฺหาติ, ปจฺจุปฺปนฺเนสุ กิจฺเจสุ พฺยสนํ ทสฺเสติ.\csromanspacer{} อิเมหิ โข คหปติปุตฺต จตูหิ ฐาเนหิ วจีปรโม อมิตฺโต มิตฺตปติรูปโก เวทิตพฺโพ.}

\proseitem{257}{จตูหิ โข คหปติปุตฺต ฐาเนหิ อนุปฺปิยภาณี อมิตฺโต มิตฺตปติรูปโก เวทิตพฺโพ.\csromanspacer{} ปาปกํปิสฺส\footnote{ปาปกมฺมํปิสฺส (สฺยา)} อนุชานาติ, กลฺยาณํปิสฺส อนุชานาติ, สมฺมุขาสฺส วณฺณํ ภาสติ, ปรมฺมุขาสฺส อวณฺณํ ภาสติ.\csromanspacer{} อิเมหิ โข คหปติปุตฺต จตูหิ ฐาเนหิ อนุปฺปิยภาณี อมิตฺโต มิตฺตปติรูปโก เวทิตพฺโพ.}

\proseitem{258}{\csromanfolio{152}จตูหิ โข คหปติปุตฺต ฐาเนหิ อปายสหาโย อมิตฺโต มิตฺตปติรูปโก เวทิตพฺโพ.\csromanspacer{} สุรา เมรย มชฺชปฺปมาทฏฺฐานา นุโยเค สหาโย โหติ, วิกาล วิสิขา จริยา นุโยเค สหาโย โหติ, สมชฺชาภิจรเณ สหาโย โหติ, ชูตปฺปมาทฏฺฐานานุโยเค สหาโย โหติ.\csromanspacer{} อิเมหิ โข คหปติปุตฺต จตูหิ ฐาเนหิ อปายสหาโย อมิตฺโต มิตฺตปติรูปโก เวทิตพฺโพติ.}

\proseitem{259}{อิทมโวจ ภควา, อิทํ วตฺวาน สุคโต อถาปรํ เอตทโวจ สตฺถา\csromandash{}}

\csromangathameasure{“อญฺญทตฺถุหโร มิตฺโต, โย จ มิตฺโต วจีปโร. \\ อนุปฺปิยญฺจ โย อาห, อปาเยสุ จ โย สขา. \\ เอเต อมิตฺเต จตฺตาโร, อิติ วิญฺญาย ปณฺฑิโต. \\ อารกา ปริวชฺเชยฺย, มคฺคํ ปฏิภยํ ยถา”ติ.}
\csromangathagroup{\csromangathastanza{“อญฺญทตฺถุหโร มิตฺโต, โย จ มิตฺโต วจีปโร\footnote{วจีปรโม (สฺยา)}. \\ อนุปฺปิยญฺจ โย อาห, อปาเยสุ จ โย สขา.}\csromangathastanza{เอเต อมิตฺเต จตฺตาโร, อิติ วิญฺญาย ปณฺฑิโต. \\ อารกา ปริวชฺเชยฺย, มคฺคํ ปฏิภยํ ยถา”ติ.}}

\csromanmatikaanchor{mtk.120}
\csromantocmarkref{section}{สุหทมิตฺต}{mtk.120}
\bookmark[dest={mtk.120},level=1]{สุหทมิตฺต}
\csromanheader{1}{\textbf{สุหทมิตฺต}}
\proseitem{260}{จตฺตาโรเม คหปติปุตฺต มิตฺตา สุหทา เวทิตพฺพา.\csromanspacer{} อุปกาโร\footnote{อุปการโก (สฺยา)} มิตฺโต สุหโท เวทิตพฺโพ, สมานสุขทุกฺโข มิตฺโต สุหโท เวทิตพฺโพ, อตฺถกฺขายี มิตฺโต สุหโท เวทิตพฺโพ, อนุกมฺปโก มิตฺโต สุหโท เวทิตพฺโพ.}

\proseitem{261}{จตูหิ โข คหปติปุตฺต ฐาเนหิ อุปกาโร มิตฺโต สุหโท เวทิตพฺโพ.\csromanspacer{} ปมตฺตํ รกฺขติ, ปมตฺตสฺส สาปเตยฺยํ รกฺขติ, ภีตสฺส สรณํ โหติ, อุปฺปนฺเนสุ กิจฺจกรณีเยสุ ตทฺทิคุณํ โภคํ อนุปฺปเทติ.\csromanspacer{} อิเมหิ โข คหปติปุตฺต จตูหิ ฐาเนหิ อุปกาโร มิตฺโต สุหโท เวทิตพฺโพ.}

\proseitem{262}{จตูหิ โข คหปติปุตฺต ฐาเนหิ สมานสุขทุกฺโข มิตฺโต สุหโท เวทิตพฺโพ.\csromanspacer{} คุยฺหมสฺส อาจิกฺขติ, คุยฺหมสฺส ปริคูหติ, อาปทาสุ น วิชหติ, ชีวิตํปิสฺส อตฺถาย ปริจฺจตฺตํ โหติ.\csromanspacer{} อิเมหิ โข คหปติปุตฺต จตูหิ ฐาเนหิ สมานสุขทุกฺโข มิตฺโต สุหโท เวทิตพฺโพ.}

\csromanheader{2}{8. สิงฺคาลสุตฺต}
\proseitem{263}{\csromanfolio{153}จตูหิ โข คหปติปุตฺต ฐาเนหิ อตฺถกฺขายี มิตฺโต สุหโท เวทิตพฺโพ.\csromanspacer{} ปาปา นิวาเรติ, กลฺยาเณ นิเวเสติ, อสฺสุตํ สาเวติ, สคฺคสฺส มคฺคํ อาจิกฺขติ.\csromanspacer{} อิเมหิ โข คหปติปุตฺต จตูหิ ฐาเนหิ อตฺถกฺขายี มิตฺโต สุหโท เวทิตพฺโพ.}

\proseitem{264}{จตูหิ โข คหปติปุตฺต ฐาเนหิ อนุกมฺปโก มิตฺโต สุหโท เวทิตพฺโพ.\csromanspacer{} อภเวนสฺส น นนฺทติ, ภเวนสฺส นนฺทติ, อวณฺณํ ภณมานํ นิวาเรติ, วณฺณํ ภณมานํ ปสํสติ.\csromanspacer{} อิเมหิ โข คหปติปุตฺต จตูหิ ฐาเนหิ อนุกมฺปโก มิตฺโต สุหโท เวทิตพฺโพติ.}

\proseitem{265}{อิทมโวจ ภควา, อิทํ วตฺวาน สุคโต อถาปรํ เอตทโวจ สตฺถา\csromandash{}}

\csromangathameasure{“อุปกาโร จ โย มิตฺโต, สุเข ทุกฺเข จ โย สขา. \\ อตฺถกฺขายี จ โย มิตฺโต, โย จ มิตฺตานุกมฺปโก. \\ เอเตปิ มิตฺเต จตฺตาโร, อิติ วิญฺญาย ปณฺฑิโต. \\ สกฺกจฺจํ ปยิรุปาเสยฺย, มาตา ปุตฺตํว โอรสํ. \\ ปณฺฑิโต สีลสมฺปนฺโน, ชลํ อคฺคีว ภาสติ. \\ โภเค สํหรมานสฺส, ภมรสฺเสว อิรียโต. \\ โภคา สนฺนิจยํ ยนฺติ, วมฺมิโกวุปจียติ. \\ เอวํ โภเค สมาหตฺวา, อลมตฺโต กุเล คิหี. \\ จตุธา วิภเช โภเค, ส เว มิตฺตานิ คนฺถติ. \\ เอเกน โภเค ภุญฺเชยฺย, ทฺวีหิ กมฺมํ ปโยชเย. \\ จตุตฺถญฺจ นิธาเปยฺย, อาปทาสุ ภวิสฺสตี”ติ.}
\csromangathagroup{\csromangathastanza{“อุปกาโร จ โย มิตฺโต, สุเข ทุกฺเข\footnote{สุขทุกฺโข (สฺยา, ก)} จ โย สขา\footnote{โย จ มิตฺโต สุเข ทุกฺเข (สี, อิ)}. \\ อตฺถกฺขายี จ โย มิตฺโต, โย จ มิตฺตานุกมฺปโก.}\csromangathastanza{เอเตปิ มิตฺเต จตฺตาโร, อิติ วิญฺญาย ปณฺฑิโต. \\ สกฺกจฺจํ ปยิรุปาเสยฺย, มาตา ปุตฺตํว โอรสํ.}\csromangathastanza{ปณฺฑิโต สีลสมฺปนฺโน, ชลํ อคฺคีว ภาสติ. \\ โภเค สํหรมานสฺส, ภมรสฺเสว อิรียโต.}\csromangathastanza{โภคา สนฺนิจยํ ยนฺติ, วมฺมิโกวุปจียติ. \\ เอวํ โภเค สมาหตฺวา\footnote{สมาหริตฺวา (สฺยา)}, อลมตฺโต กุเล คิหี.}\csromangathastanza{จตุธา วิภเช โภเค, ส เว มิตฺตานิ คนฺถติ. \\ เอเกน โภเค ภุญฺเชยฺย, ทฺวีหิ กมฺมํ ปโยชเย. \\ จตุตฺถญฺจ นิธาเปยฺย, อาปทาสุ ภวิสฺสตี”ติ.}}

\csromanmatikaanchor{mtk.121}
\csromantocmarkref{section}{ฉทฺทิสาปฏิจฺฉาทนกณฺฑ}{mtk.121}
\bookmark[dest={mtk.121},level=1]{ฉทฺทิสาปฏิจฺฉาทนกณฺฑ}
\csromanheader{1}{\textbf{ฉทฺทิสาปฏิจฺฉาทนกณฺฑ}}
\proseitem{266}{กถญฺจ คหปติปุตฺต อริยสาวโก ฉทฺทิสาปฏิจฺฉาที โหติ.\csromanspacer{} ฉ อิมา คหปติปุตฺต ทิสา เวทิตพฺพา.\csromanspacer{} ปุรตฺถิมา ทิสา มาตาปิตโร เวทิตพฺพา, ทกฺขิณา ทิสา อาจริยา เวทิตพฺพา, ปจฺฉิมา \csromanfolio{154}ทิสา ปุตฺตทารา เวทิตพฺพา, อุตฺตรา ทิสา มิตฺตามจฺจา เวทิตพฺพา, เหฏฺฐิมา ทิสา ทาสกมฺมกรา เวทิตพฺพา, อุปริมา ทิสา สมณพฺราหฺมณา เวทิตพฺพา.}

\proseitem{267}{ปญฺจหิ โข คหปติปุตฺต ฐาเนหิ ปุตฺเตน ปุรตฺถิมา ทิสา มาตาปิตโร ปจฺจุปฏฺฐาตพฺพา “ภโต เน\footnote{เนสํ (พหูสุ)} ภริสฺสามิ, กิจฺจํ เนสํ กริสฺสามิ, กุลวํสํ ฐเปสฺสามิ, ทายชฺชํ ปฏิปชฺชามิ, อถ วา ปน เปตานํ กาลงฺกตานํ ทกฺขิณํ อนุปฺปทสฺสามีติ.\csromanspacer{} อิเมหิ โข คหปติปุตฺต ปญฺจหิ ฐาเนหิ ปุตฺเตน ปุรตฺถิมา ทิสา มาตาปิตโร ปจฺจุปฏฺฐิตา ปญฺจหิ ฐาเนหิ ปุตฺตํ อนุกมฺปนฺติ, ปาปา นิวาเรนฺติ, กลฺยาเณ นิเวเสนฺติ, สิปฺปํ สิกฺขาเปนฺติ, ปติรูเปน ทาเรน สํโยเชนฺติ, สมเย ทายชฺชํ นิยฺยาเทนฺติ\footnote{นิยฺยาเตนฺติ (ก-สี)}.\csromanspacer{} อิเมหิ โข คหปติปุตฺต ปญฺจหิ ฐาเนหิ ปุตฺเตน ปุรตฺถิมา ทิสา มาตาปิตโร ปจฺจุปฏฺฐิตา อิเมหิ ปญฺจหิ ฐาเนหิ ปุตฺตํ อนุกมฺปนฺติ.\csromanspacer{} เอวมสฺส เอสา ปุรตฺถิมา ทิสา ปฏิจฺฉนฺนา โหติ เขมา อปฺปฏิภยา.}

\proseitem{268}{ปญฺจหิ โข คหปติปุตฺต ฐาเนหิ อนฺเตวาสินา ทกฺขิณา ทิสา อาจริยา ปจฺจุปฏฺฐาตพฺพา อุฏฺฐาเนน อุปฏฺฐาเนน สุสฺสุสาย ปาริจริยาย สกฺกจฺจํ สิปฺปปฏิคฺคหเณน\footnote{สิปฺปํ ปฏิคฺคหเณน (สฺยา), สิปฺป-อุคฺคหเณน (ก)}.\csromanspacer{} อิเมหิ โข คหปติปุตฺต ปญฺจหิ ฐาเนหิ อนฺเตวาสินา ทกฺขิณา ทิสา อาจริยา ปจฺจุปฏฺฐิตา ปญฺจหิ ฐาเนหิ อนฺเตวาสิํ อนุกมฺปนฺติ, สุวินีตํ วิเนนฺติ, สุคฺคหิตํ คาหาเปนฺติ, สพฺพสิปฺปสฺสุตํ สมกฺขายิโน ภวนฺติ, มิตฺตามจฺเจสุ ปฏิยาเทนฺติ\footnote{ปฏิเวเทนฺติ (สฺยา)}, ทิสาสุ ปริตฺตาณํ กโรนฺติ.\csromanspacer{} อิเมหิ โข คหปติปุตฺต ปญฺจหิ ฐาเนหิ อนฺเตวาสินา ทกฺขิณา ทิสา อาจริยา ปจฺจุปฏฺฐิตา อิเมหิ ปญฺจหิ ฐาเนหิ อนฺเตวาสิํ อนุกมฺปนฺติ.\csromanspacer{} เอวมสฺส เอสา ทกฺขิณา ทิสา ปฏิจฺฉนฺนา โหติ เขมา อปฺปฏิภยา.}

\proseitem{269}{ปญฺจหิ โข คหปติปุตฺต ฐาเนหิ สามิเกน ปจฺฉิมา ทิสา ภริยา ปจฺจุปฏฺฐาตพฺพา สมฺมานนาย อนวมานนาย\footnote{อวิมานนาย (สฺยา, อิ)} อนติจริยาย อิสฺสริยโวสฺสคฺเคน อลงฺการานุปฺปทาเนน.\csromanspacer{} อิเมหิ โข คหปติปุตฺต ปญฺจหิ ฐาเนหิ สามิเกน ปจฺฉิมา ทิสา ภริยา ปจฺจุปฏฺฐิตา ปญฺจหิ}

\vspace{\tipitakaparskip}
\csromancenter{สิงฺคาลสุตฺต ฐาเนหิ สามิกํ อนุกมฺปติ, สุสํวิหิตกมฺมนฺตา จ โหติ, สงฺคหิตปริชนา\footnote{สุสงฺคหิตปริชนา (สี, สฺยา, อิ)} จ, อนติจารินี จ, สมฺภตํ จ อนุรกฺขติ, ทกฺขา จ โหติ อนลสา สพฺพกิจฺเจสุ.\csromanspacer{} อิเมหิ โข คหปติปุตฺต ปญฺจหิ ฐาเนหิ สามิเกน ปจฺฉิมา ทิสา ภริยา ปจฺจุปฏฺฐิตา อิเมหิ ปญฺจหิ ฐาเนหิ สามิกํ อนุกมฺปติ.\csromanspacer{} เอวมสฺส เอสา ปจฺฉิมา ทิสา ปฏิจฺฉนฺนา โหติ เขมา อปฺปฏิภยา.}
\proseitem{270}{\csromanfolio{155}ปญฺจหิ โข คหปติปุตฺต ฐาเนหิ กุลปุตฺเตน อุตฺตรา ทิสา มิตฺตามจฺจา ปจฺจุปฏฺฐาตพฺพา ทาเนน เปยฺยวชฺเชน\footnote{ปิยวชฺเชน (สฺยา, ก)} อตฺถจริยาย สมานตฺตตาย อวิสํวาทนตาย.\csromanspacer{} อิเมหิ โข คหปติปุตฺต ปญฺจหิ ฐาเนหิ กุลปุตฺเตน อุตฺตรา ทิสา มิตฺตามจฺจา ปจฺจุปฏฺฐิตา ปญฺจหิ ฐาเนหิ กุลปุตฺตํ อนุกมฺปนฺติ, ปมตฺตํ รกฺขนฺติ, ปมตฺตสฺส สาปเตยฺยํ รกฺขนฺติ, ภีตสฺส สรณํ โหนฺติ, อาปทาสุ น วิชหนฺติ, อปรปชา จสฺส ปฏิปูเชนฺติ.\csromanspacer{} อิเมหิ โข คหปติปุตฺต ปญฺจหิ ฐาเนหิ กุลปุตฺเตน อุตฺตรา ทิสา มิตฺตามจฺจา ปจฺจุปฏฺฐิตา อิเมหิ ปญฺจหิ ฐาเนหิ กุลปุตฺตํ อนุกมฺปนฺติ.\csromanspacer{} เอวมสฺส เอสา อุตฺตรา ทิสา ปฏิจฺฉานฺนา โหติ เขมา อปฺปฏิภยา.}

\proseitem{271}{ปญฺจหิ โข คหปติปุตฺต ฐาเนหิ อยฺยิรเกน\footnote{อยิรเกน (สี, สฺยา, อิ)} เหฏฺฐิมา ทิสา ทาสกมฺมกรา ปจฺจุปฏฺฐาตพฺพา ยถาพลํ กมฺมนฺตสํวิธาเนน ภตฺตเวตนานุปฺปทาเนน คิลานุปฏฺฐาเนน อจฺฉริยานํ รสานํ สํวิภาเคน สมเย โวสฺสคฺเคน.\csromanspacer{} อิเมหิ โข คหปติปุตฺต ปญฺจหิ ฐาเนหิ อยฺยิรเกน เหฏฺฐิมา ทิสา ทาสกมฺมกรา ปจฺจุปฏฺฐิตา ปญฺจหิ ฐาเนหิ อยฺยิรกํ อนุกมฺปนฺติ, ปุพฺพุฏฺฐายิโน จ โหนฺติ, ปจฺฉา นิปาติโน จ, ทินฺนาทายิโน จ, สุกตกมฺมกรา จ, กิตฺติวณฺณหรา จ.\csromanspacer{} อิเมหิ โข คหปติปุตฺต ปญฺจหิ ฐาเนหิ อยฺยิรเกน เหฏฺฐิมา ทิสา ทาสกมฺมกรา ปจฺจุปฏฺฐิตา อิเมหิ ปญฺจหิ ฐาเนหิ อยฺยิรกํ อนุกมฺปนฺติ.\csromanspacer{} เอวมสฺส เอสา เหฏฺฐิมา ทิสา ปฏิจฺฉนฺนา โหติ เขมา อปฺปฏิภยา.}

\proseitem{272}{ปญฺจหิ โข คหปติปุตฺต ฐาเนหิ กุลปุตฺเตน อุปริมา ทิสา สมณพฺราหฺมณา ปจฺจุปฏฺฐาตพฺพา เมตฺเตน กายกมฺเมน เมตฺเตน วจีกมฺเมน \csromanfolio{156}เมตฺเตน มโนกมฺเมน อนาวฏทฺวารตาย อามิสานุปฺปทาเนน.\csromanspacer{} อิเมหิ โข คหปติปุตฺต ปญฺจหิ ฐาเนหิ กุลปุตฺเตน อุปริมา ทิสา สมณพฺราหฺมณา ปจฺจุปฏฺฐิตา ฉหิ ฐาเนหิ กุลปุตฺตํ อนุกมฺปนฺติ, ปาปา นิวาเรนฺติ, กลฺยาเณ นิเวเสนฺติ, กลฺยาเณน มนสา อนุกมฺปนฺติ, อสฺสุตํ สาเวนฺติ, สุตํ ปริโยทาเปนฺติ, สคฺคสฺส มคฺคํ อาจิกฺขนฺติ.\csromanspacer{} อิเมหิ โข คหปติปุตฺต ฉหิ ฐาเนหิ กุลปุตฺเตน อุปริมา ทิสา สมณพฺราหฺมณา ปจฺจุปฏฺฐิตา อิเมหิ ฉหิ ฐาเนหิ กุลปุตฺตํ อนุกมฺปนฺติ.\csromanspacer{} เอวมสฺส เอสา อุปริมา ทิสา ปฏิจฺฉนฺนา โหติ เขมา อปฺปฏิภยาติ.}

\proseitem{273}{อิทมโวจ ภควา, อิทํ วตฺวาน สุคโต อถาปรํ เอตทโวจ สตฺถา\csromandash{}}

\csromangathameasure{“มาตาปิตา ทิสา ปุพฺพา, อาจริยา ทกฺขิณา ทิสา. \\ ปุตฺตทารา ทิสา ปจฺฉา, มิตฺตามจฺจา จ อุตฺตรา. \\ ทาสกมฺมกรา เหฏฺฐา, อุทฺธํ สมณพฺราหฺมณา. \\ เอตา ทิสา นมสฺเสยฺย, อลมตฺโต กุเล คิหี. \\ ปณฺฑิโต สีลสมฺปนฺโน, สโณฺห จ ปฏิภานวา. \\ นิวาตวุตฺติ อตฺถทฺโธ, ตาทิโส ลภเต ยสํ. \\ อุฏฺฐานโก อนลโส, อาปทาสุ น เวธติ. \\ อจฺฉินฺนวุตฺติ เมธาวี, ตาทิโส ลภเต ยสํ. \\ สงฺคาหโก มิตฺตกโร, วทญฺญู วีตมจฺฉโร. \\ เนตา วิเนตา อนุเนตา, ตาทิโส ลภเต ยสํ. \\ ทานญฺจ เปยฺยวชฺชญฺจ, อตฺถจริยา จ ยา อิธ. \\ สมานตฺตตา จ ธมฺเมสุ, ตตฺถ ตตฺถ ยถารหํ. \\ เอเต โข สงฺคหา โลเก, รถสฺสาณีว ยายโต. \\ เอเต จ สงฺคหา นาสฺสุ, น มาตา ปุตฺตการณา. \\ ลเภถ มานํ ปูชํ วา, ปิตา วา ปุตฺตการณา. \\ ยสฺมา จ สงฺคหา เอเต, สมฺมเปกฺขนฺติ ปณฺฑิตา. \\ ตสฺมา มหตฺตํ ปปฺโปนฺติ, ปาสํสา จ ภวนฺติ เต”ติ.}
\csromangathagroup{\csromangathastanza{“มาตาปิตา ทิสา ปุพฺพา, อาจริยา ทกฺขิณา ทิสา. \\ ปุตฺตทารา ทิสา ปจฺฉา, มิตฺตามจฺจา จ อุตฺตรา.}\csromangathastanza{ทาสกมฺมกรา เหฏฺฐา, อุทฺธํ สมณพฺราหฺมณา. \\ เอตา ทิสา นมสฺเสยฺย, อลมตฺโต กุเล คิหี.}\csromangathastanza{ปณฺฑิโต สีลสมฺปนฺโน, สโณฺห จ ปฏิภานวา. \\ นิวาตวุตฺติ อตฺถทฺโธ, ตาทิโส ลภเต ยสํ.}\csromangathastanza{อุฏฺฐานโก อนลโส, อาปทาสุ น เวธติ. \\ อจฺฉินฺนวุตฺติ เมธาวี, ตาทิโส ลภเต ยสํ.}\csromangathastanza{สงฺคาหโก มิตฺตกโร, วทญฺญู วีตมจฺฉโร. \\ เนตา วิเนตา อนุเนตา, ตาทิโส ลภเต ยสํ.}\csromangathastanza{ทานญฺจ เปยฺยวชฺชญฺจ, อตฺถจริยา จ ยา อิธ. \\ สมานตฺตตา จ ธมฺเมสุ, ตตฺถ ตตฺถ ยถารหํ.}\csromangathastanza{เอเต โข สงฺคหา โลเก, รถสฺสาณีว ยายโต. \\ เอเต จ สงฺคหา นาสฺสุ, น มาตา ปุตฺตการณา.}\csromangathastanza{ลเภถ มานํ ปูชํ วา, ปิตา วา ปุตฺตการณา. \\ ยสฺมา จ สงฺคหา เอเต, สมฺมเปกฺขนฺติ\footnote{สมเวกฺขนฺติ (สี, อิ, ก)} ปณฺฑิตา. \\ ตสฺมา มหตฺตํ ปปฺโปนฺติ, ปาสํสา จ ภวนฺติ เต”ติ.}}

\csromanheader{2}{8. สิงฺคาลสุตฺต}
\proseitem{274}{\csromanfolio{157}เอวํ วุตฺเต สิงฺคาลโก คหปติปุตฺโต ภควนฺตํ เอตทโวจ “อภิกฺกนฺตํ ภนฺเต, อภิกฺกนฺตํ ภนฺเต, เสยฺยถาปิ ภนฺเต นิกฺกุชฺชิตํ วา อุกฺกุชฺเชยฺย, ปฏิจฺฉนฺนํ วา วิวเรยฺย, มูฬฺหสฺส วา มคฺคํ อาจิกฺเขยฺย, อนฺธกาเร วา เตลปชฺโชตํ ธาเรยฺย “จกฺขุมนฺโต รูปานิ ทกฺขนฺตี”ติ, เอวเมวํ ภควตา อเนกปริยาเยน ธมฺโม ปกาสิโต, เอสาหํ ภนฺเต ภควนฺตํ สรณํ คจฺฉามิ ธมฺมญฺจ ภิกฺขุสํฆญฺจ, อุปาสกํ มํ ภควา ธาเรตุ อชฺชตคฺเค ปาณุเปตํ สรณํ คตนฺ”ติ.}

\nitthitam{\textbf{สิงฺคาลสุตฺตํ}\footnote{สิงฺคาโลวาทสุตฺตนฺตํ (อิ)} \textbf{นิฏฺฐิตํ อฏฺฐมํ.}}
\csromanmatikaanchor{mtk.122}
\csromantocmarknopage{chapter}{9. อาฏานาฏิยสุตฺต}{mtk.122}
\bookmark[dest={mtk.122},level=0]{9. อาฏานาฏิยสุตฺต}
\chapterheadpageread{9. อาฏานาฏิยสุตฺต}
\csromanmatikaanchor{mtk.123}
\csromantocmarkref{section}{ปฐมภาณวาร}{mtk.123}
\bookmark[dest={mtk.123},level=1]{ปฐมภาณวาร}
\csromanheader{1}{\textbf{ปฐมภาณวาร}}
\proseitem{275}{\csromanfolio{158}เอวํ เม สุตํ\csromandash{}เอกํ สมยํ ภควา ราชคเห วิหรติ คิชฺฌกูเฏ ปพฺพเต.\csromanspacer{} อถ โข จตฺตาโร มหาราชา\footnote{มหาราชาโน (ก)} มหติยา จ ยกฺขเสนาย มหติยา จ คนฺธพฺพเสนาย มหติยา จ กุมฺภณฺฑเสนาย มหติยา จ นาคเสนาย จตุทฺทิสํ รกฺขํ ฐเปตฺวา จตุทฺทิสํ คุมฺพํ ฐเปตฺวา จตุทฺทิสํ โอวรณํ ฐเปตฺวา อภิกฺกนฺตาย รตฺติยา อภิกฺกนฺตวณฺณา เกวลกปฺปํ คิชฺฌกูฏํ ปพฺพตํ โอภาเสตฺวา\footnote{คิชฺฌกูฏํ โอภาเสตฺวา (สี, สฺยา, อิ)} เยน ภควา เตนุปสงฺกมิํสุ, อุปสงฺกมิตฺวา ภควนฺตํ อภิวาเทตฺวา เอกมนฺตํ นิสีทิํสุ.\csromanspacer{} เตปิ โข ยกฺขา อปฺเปกจฺเจ ภควนฺตํ อภิวาเทตฺวา เอกมนฺตํ นิสีทิํสุ, อปฺเปกจฺเจ ภควตา สทฺธิํ สมฺโมทิํสุ, สมฺโมทนียํ กถํ สารณียํ วีติสาเรตฺวา เอกมนฺตํ นิสีทิํสุ, อปฺเปกจฺเจ เยน ภควา เตนญฺชลิํ ปณาเมตฺวา เอกมนฺตํ นิสีทิํสุ, อปฺเปกจฺเจ นามโคตฺตํ สาเวตฺวา เอกมนฺตํ นิสีทิํสุ, อปฺเปกจฺเจ ตุณฺหีภูตา เอกมนฺตํ นิสีทิํสุ.}

\proseitem{276}{เอกมนฺตํ นิสินฺโน โข เวสฺสวโณ มหาราชา ภควนฺตํ เอตทโวจ “สนฺติ หิ ภนฺเต อุฬารา ยกฺขา ภควโต อปฺปสนฺนา, สนฺติ หิ ภนฺเต อุฬารา ยกฺขา ภควโต ปสนฺนา, สนฺติ หิ ภนฺเต มชฺฌิมา ยกฺขา ภควโต อปฺปสนฺนา, สนฺติ หิ ภนฺเต มชฺฌิมา ยกฺขา ภควโต ปสนฺนา, สนฺติ หิ ภนฺเต นีจา ยกฺขา ภควโต อปฺปสนฺนา, สนฺติ หิ ภนฺเต นีจา ยกฺขา ภควโต ปสนฺนา.\csromanspacer{} เยภุยฺเยน โข ปน ภนฺเต ยกฺขา อปฺปสนฺนาเยว ภควโต.\csromanspacer{} ตํ กิสฺส เหตุ, ภควา หิ ภนฺเต ปาณาติปาตา เวรมณิยา ธมฺมํ เทเสติ, อทินฺนาทานา เวรมณิยา ธมฺมํ เทเสติ, กาเมสุมิจฺฉาจารา เวรมณิยา ธมฺมํ เทเสติ, มุสาวาทา เวรมณิยา ธมฺมํ เทเสติ, สุราเมรยมชฺชปฺปมาทฏฺฐานา เวรมณิยา ธมฺมํ เทเสติ.\csromanspacer{} เยภุยฺเยน โข ปน ภนฺเต ยกฺขา อปฺปฏิวิรตาเยว ปาณาติปาตา, อปฺปฏิวิรตา อทินฺนาทานา, อปฺปฏิวิรตา กาเมสุมิจฺฉาจารา, อปฺปฏิวิรตา มุสาวาทา, อปฺปฏิวิรตา สุราเมรยมชฺชปฺปมาทฏฺฐานา.\csromanspacer{} เตสํ ตํ โหติ อปฺปิยํ อมนาปํ.\csromanspacer{} สนฺติ หิ ภนฺเต ภควโต สาวกา อรญฺญวนปตฺถานิ}

\vspace{\tipitakaparskip}
\csromancenter{อาฏานาฏิยสุตฺต ปนฺตานิ เสนาสนานิ ปฏิเสวนฺติ อปฺปสทฺทานิ อปฺปนิคฺโฆสานิ วิชนวาตานิ มนุสฺสราหสฺเสยฺยกานิ\footnote{มนุสฺสราหเสยฺยกานิ (สี, สฺยา, อิ)} ปฏิสลฺลานสารุปฺปานิ.\csromanspacer{} ตตฺถ สนฺติ อุฬารา ยกฺขา นิวาสิโน, เย อิมสฺมิํ ภควโต ปาวจเน อปฺปสนฺนา, เตสํ ปสาทาย อุคฺคณฺหาตุ ภนฺเต ภควา อาฏานาฏิยํ รกฺขํ ภิกฺขุนํ ภิกฺขุนีนํ อุปาสกานํ อุปาสิกานํ คุตฺติยา รกฺขาย อวิหิํสาย ผาสุวิหารายาติ.\csromanspacer{} อธิวาเสสิ ภควา ตุณฺหีภาเวน.}
\prose{\csromanfolio{159}อถ โข เวสฺสวโณ มหาราชา ภควโต อธิวาสนํ วิทิตฺวา ตายํ เวลายํ อิมํ อาฏานาฏิยํ รกฺขํ อภาสิ\csromandash{}}

\csromanhangingbegin
\hangingheaditem{277}{วิปสฺสิสฺส จ\footnote{อิเม จการา โปราณโปตฺถเกสุ นตฺถิ.} นมตฺถุ, จกฺขุมนฺตสฺส สิรีมโต.}
\hangingbody{สิขิสฺสปิ จ2 นมตฺถุ, สพฺพภูตานุกมฺปิโน.}
\hangingbody{เวสฺสภุสฺส จ2 นมตฺถุ, นฺหาตกสฺส ตปสฺสิโน.}
\hangingbody{นมตฺถุ กกุสนฺธสฺส, มารเสนาปมทฺทิโน.}
\hangingbody{โกณาคมนสฺส นมตฺถุ, พฺราหฺมณสฺส วุสีมโต.}
\hangingbody{กสฺสปสฺส จ2 นมตฺถุ, วิปฺปมุตฺตสฺส สพฺพธิ.}
\hangingbody{องฺคีรสสฺส นมตฺถุ, สกฺยปุตฺตสฺส สิรีมโต.}
\hangingbody{โย อิมํ ธมฺมํ เทเสสิ3, สพฺพทุกฺขาปนูทนํ.}
\hangingbody{เย จาปิ นิพฺพุตา โลเก, ยถาภูตํ วิปสฺสิสุํ.}
\hangingbody{เต ชนา อปิสุณาถ4, มหนฺตา วีตสารทา.}
\hangingbody{หิตํ เทวมนุสฺสานํ, ยํ นมสฺสนฺติ โคตมํ.}
\hangingbody{วิชฺชาจรณสมฺปนฺนํ, มหนฺตํ วีตสารทํ.}
\csromanhangingend

\csromanhangingbegin
\hangingheaditem{278}{ยโต อุคฺคจฺฉติ สูริโย\footnote{สุริโย (สี, สฺยา, อิ)}, อาทิจฺโจ มณฺฑลี มหา.}
\hangingbody{ยสฺส จุคฺคจฺฉมานสฺส, สํวรีปิ นิรุชฺฌติ.}
\hangingbody{ยสฺส จุคฺคเต สูริเย, ทิวโสติ ปวุจฺจติ.}
\csromanhangingend

\csromangathameasure{รหโทปิ ตตฺถ คมฺภีโร, สมุทฺโท สริโตทโก. \\ เอวํ ตํ ตตฺถ ชานนฺติ, สมุทฺโท สริโตทโก. \\ อิโต สา ปุริมา ทิสา, อิติ นํ อาจิกฺขตี ชโน. \\ ยํ ทิสํ อภิปาเลติ, มหาราชา ยสสฺสิ โส. \\ คนฺธพฺพานํ อธิปติ, ธตรฏฺโฐติ นามโส. \\ รมตี นจฺจคีเตหิ, คนฺธพฺเพหิ ปุรกฺขโต. \\ ปุตฺตาปิ ตสฺส พหโว, เอกนามาติ เม สุตํ. \\ อสีติ ทส เอโก จ, อินฺทนามา มหพฺพลา. \\ เต จาปิ พุทฺธํ ทิสฺวาน, พุทฺธํ อาทิจฺจพนฺธุนํ. \\ ทูรโตว นมสฺสนฺติ, มหนฺตํ วีตสารทํ. \\ นโม เต ปุริสาชญฺญ, นโม เต ปุริสุตฺตม. \\ กุสเลน สเมกฺขสิ, อมนุสฺสาปิ ตํ วนฺทนฺติ. \\ สุตํ เนตํ อภิณฺหโส, ตสฺมา เอวํ วเทมเส. \\ ชินํ วนฺทถ โคตมํ, ชินํ วนฺทาม โคตมํ. \\ วิชฺชาจรณสมฺปนฺนํ, พุทฺธํ วนฺทาม โคตมํ.}
\csromangathagroup{\csromangathastanza{\csromanfolio{160}รหโทปิ ตตฺถ คมฺภีโร, สมุทฺโท สริโตทโก. \\ เอวํ ตํ ตตฺถ ชานนฺติ, สมุทฺโท สริโตทโก.}\csromangathastanza{อิโต สา ปุริมา ทิสา, อิติ นํ อาจิกฺขตี ชโน. \\ ยํ ทิสํ อภิปาเลติ, มหาราชา ยสสฺสิ โส.}\csromangathastanza{คนฺธพฺพานํ อธิปติ\footnote{อาธิปติ (สี, สฺยา, อิ) เอวมุปริปิ,}, ธตรฏฺโฐติ นามโส. \\ รมตี นจฺจคีเตหิ, คนฺธพฺเพหิ ปุรกฺขโต.}\csromangathastanza{ปุตฺตาปิ ตสฺส พหโว, เอกนามาติ เม สุตํ. \\ อสีติ ทส เอโก จ, อินฺทนามา มหพฺพลา.}\csromangathastanza{เต จาปิ พุทฺธํ ทิสฺวาน, พุทฺธํ อาทิจฺจพนฺธุนํ. \\ ทูรโตว นมสฺสนฺติ, มหนฺตํ วีตสารทํ.}\csromangathastanza{นโม เต ปุริสาชญฺญ, นโม เต ปุริสุตฺตม. \\ กุสเลน สเมกฺขสิ, อมนุสฺสาปิ ตํ วนฺทนฺติ.}\csromangathastanza{สุตํ เนตํ อภิณฺหโส, ตสฺมา เอวํ วเทมเส. \\ ชินํ วนฺทถ โคตมํ, ชินํ วนฺทาม โคตมํ. \\ วิชฺชาจรณสมฺปนฺนํ, พุทฺธํ วนฺทาม โคตมํ.}}

\csromanhangingbegin
\hangingheaditem{279}{เยน เปตา ปวุจฺจนฺติ, ปิสุณา ปิฏฺฐิมํสิกา.}
\hangingbody{ปาณาติปาติโน ลุทฺทา2, โจรา เนกติกา ชนา.}
\hangingbody{อิโต สา ทกฺขิณา ทิสา, อิติ นํ อาจิกฺขตี ชโน.}
\hangingbody{ยํ ทิสํ อภิปาเลติ, มหาราชา ยสสฺสิ โส.}
\hangingbody{กุมฺภณฺฑานํ อธิปติ, วิรูโฬฺห-อิติ นามโส.}
\hangingbody{รมตี นจฺจคีเตหิ, กุมฺภณฺเฑหิ ปุรกฺขโต.}
\hangingbody{ปุตฺตาปิ ตสฺส พหโว, เอกนามาติ เม สุตํ.}
\hangingbody{อสีติ ทส เอโก จ, อินฺทนามา มหพฺพลา.}
\hangingbody{เต จาปิ พุทฺธํ ทิสฺวาน, พุทฺธํ อาทิจฺจพนฺธุนํ.}
\hangingbody{ทูรโตว นมสฺสนฺติ, มหนฺตํ วีตสารทํ.}
\csromanhangingend

\csromanheader{2}{9. อาฏานาฏิยสุตฺต}
\csromangathameasure{นโม เต ปุริสาชญฺญ, นโม เต ปุริสุตฺตม. \\ กุสเลน สเมกฺขสิ, อมนุสฺสาปิ ตํ วนฺทนฺติ. \\ สุตํ เนตํ อภิณฺหโส, ตสฺมา เอวํ วเทมเส. \\ ชินํ วนฺทถ โคตมํ, ชินํ วนฺทาม โคตมํ. \\ วิชฺชาจรณสมฺปนฺนํ, พุทฺธํ วนฺทาม โคตมํ.}
\csromangathagroup{\csromangathastanza{\csromanfolio{161}นโม เต ปุริสาชญฺญ, นโม เต ปุริสุตฺตม. \\ กุสเลน สเมกฺขสิ, อมนุสฺสาปิ ตํ วนฺทนฺติ.}\csromangathastanza{สุตํ เนตํ อภิณฺหโส, ตสฺมา เอวํ วเทมเส. \\ ชินํ วนฺทถ โคตมํ, ชินํ วนฺทาม โคตมํ. \\ วิชฺชาจรณสมฺปนฺนํ, พุทฺธํ วนฺทาม โคตมํ.}}

\csromanhangingbegin
\hangingheaditem{280}{ยตฺถ โจคฺคจฺฉติ สูริโย, อาทิจฺโจ มณฺฑลี มหา.}
\hangingbody{ยสฺส โจคฺคจฺฉมานสฺส, ทิวโสปิ นิรุชฺฌติ.}
\hangingbody{ยสฺส โจคฺคเต สูริเย, สํวรีติ ปวุจฺจติ.}
\hangingbody{รหโทปิ ตตฺถ คมฺภีโร, สมุทฺโท สริโตทโก.}
\hangingbody{เอวํ ตํ ตตฺถ ชานนฺติ, สมุทฺโท สริโตทโก.}
\hangingbody{อิโต สา ปจฺฉิมา ทิสา, อิติ นํ อาจิกฺขตี ชโน.}
\hangingbody{ยํ ทิสํ อภิปาเลติ, มหาราชา ยสสฺสิ โส.}
\hangingbody{นาคานญฺจ อธิปติ, วิรูปกฺโขติ นามโส.}
\hangingbody{รมตี นจฺจคีเตหิ, นาเคเหว ปุรกฺขโต.}
\hangingbody{ปุตฺตาปิ ตสฺส พหโว, เอกนามาติ เม สุตํ.}
\hangingbody{อสีติ ทส เอโก จ, อินฺทนามา มหพฺพลา.}
\hangingbody{เต จาปิ พุทฺธํ ทิสฺวาน, พุทฺธํ อาทิจฺจพนฺธุนํ.}
\hangingbody{ทูรโตว นมสฺสนฺติ, มหนฺตํ วีตสารทํ.}
\hangingbody{นโม เต ปุริสาชญฺญ, นโม เต ปุริสุตฺตม.}
\hangingbody{กุสเลน สเมกฺขสิ, อมนุสฺสาปิ ตํ วนฺทนฺติ.}
\hangingbody{สุตํ เนตํ อภิณฺหโส, ตสฺมา เอวํ เอเทมเส.}
\hangingbody{ชินํ วนฺทถ โคตมํ, ชินํ วนฺทาม โคตมํ.}
\hangingbody{วิชฺชาจรณสมฺปนฺนํ, พุทฺธํ วนฺทาม โคตมํ.}
\csromanhangingend

\csromanhangingbegin
\hangingheaditem{281}{เยน อุตฺตรกุรุโวฺห\footnote{อุตฺตรกุรู รมฺมา (สี, สฺยา, อิ)}, มหาเนรุ สุทสฺสโน.}
\hangingbody{มนุสฺสา ตตฺถ ชายนฺติ, อมมา อปริคฺคหา.}
\csromanhangingend

\csromangathameasure{น เต พีชํ ปวปนฺติ, นปิ นียนฺติ นงฺคลา. \\ อกฏฺฐปากิมํ สาลิํ, ปริภุญฺชนฺติ มานุสา. \\ อกณํ อถุสํ สุทฺธํ, สุคนฺธํ ตณฺฑุลปฺผลํ. \\ ตุณฺฑิกีเร ปจิตฺวาน, ตโต ภุญฺชนฺติ โภชนํ. \\ คาวิํ เอกขุรํ กตฺวา, อนุยนฺติ ทิโสทิสํ. \\ ปสุํ เอกขุรํ กตฺวา, อนุยนฺติ ทิโสทิสํ. \\ อิตฺถิํ วา วาหนํ กตฺวา, อนุยนฺติ ทิโสทิสํ. \\ ปุริสํ วาหนํ กตฺวา, อนุยนฺติ ทิโสทิสํ. \\ กุมาริํ วาหนํ กตฺวา, อนุยนฺติ ทิโสทิสํ. \\ กุมารํ วาหนํ กตฺวา, อนุยนฺติ ทิโสทิสํ. \\ เต ยาเน อภิรุหิตฺวา, \\ สพฺพา ทิสา อนุปริยายนฺติ. \\ ปจารา ตสฺส ราชิโน. \\ หตฺถิยานํ อสฺสยานํ, ทิพฺพํ ยานํ อุปฏฺฐิตํ. \\ ปาสาทา สิวิกา เจว, มหาราชสฺส ยสสฺสิโน. \\ ตสฺส จ นครา อหุ, \\ อนฺตลิกฺเข สุมาปิตา. \\ อาฏานาฏา กุสินาฏา ปรกุสินาฏา, \\ นาฏสุริยา ปรกุสิฏนาฏา. \\ อุตฺตเรน กสิวนฺโต, \\ ชโนฆมปเรน จ. \\ นวนวุติโย อมฺพร-อมฺพรวติโย, \\ อาฬกมนฺทา นาม ราชธานี. \\ กุเวรสฺส โข ปน มาริส มหาราชสฺส วิสาณา นาม ราชธานี, \\ ตสฺมา กุเวโร มหาราชา, เวสฺสวโณติ ปวุจฺจติ.}
\csromangathagroup{\csromangathastanza{\csromanfolio{162}น เต พีชํ ปวปนฺติ, นปิ นียนฺติ นงฺคลา. \\ อกฏฺฐปากิมํ สาลิํ, ปริภุญฺชนฺติ มานุสา.}\csromangathastanza{อกณํ อถุสํ สุทฺธํ, สุคนฺธํ ตณฺฑุลปฺผลํ. \\ ตุณฺฑิกีเร ปจิตฺวาน, ตโต ภุญฺชนฺติ โภชนํ.}\csromangathastanza{คาวิํ เอกขุรํ กตฺวา, อนุยนฺติ ทิโสทิสํ. \\ ปสุํ เอกขุรํ กตฺวา, อนุยนฺติ ทิโสทิสํ.}\csromangathastanza{อิตฺถิํ วา วาหนํ\footnote{อิตฺถี-วาหนํ (สี, อิ), อิตฺถิํ วาหนํ (สฺยา)} กตฺวา, อนุยนฺติ ทิโสทิสํ. \\ ปุริสํ วาหนํ กตฺวา, อนุยนฺติ ทิโสทิสํ.}\csromangathastanza{กุมาริํ วาหนํ กตฺวา, อนุยนฺติ ทิโสทิสํ. \\ กุมารํ วาหนํ กตฺวา, อนุยนฺติ ทิโสทิสํ.}\csromangathastanza{เต ยาเน อภิรุหิตฺวา, \\ สพฺพา ทิสา อนุปริยายนฺติ\footnote{อนุปริยนฺติ (สฺยา)}. \\ ปจารา ตสฺส ราชิโน.}\csromangathastanza{หตฺถิยานํ อสฺสยานํ, ทิพฺพํ ยานํ อุปฏฺฐิตํ. \\ ปาสาทา สิวิกา เจว, มหาราชสฺส ยสสฺสิโน.}\csromangathastanza{ตสฺส จ นครา อหุ, \\ อนฺตลิกฺเข สุมาปิตา. \\ อาฏานาฏา กุสินาฏา ปรกุสินาฏา, \\ นาฏสุริยา\footnote{นาฏปุริยา (สี, อิ), นาฏปริยา (สฺยา)} ปรกุสิฏนาฏา.}\csromangathastanza{อุตฺตเรน กสิวนฺโต\footnote{กปิวนฺโต (สี, สฺยา, อิ)}, \\ ชโนฆมปเรน จ. \\ นวนวุติโย อมฺพร-อมฺพรวติโย, \\ อาฬกมนฺทา นาม ราชธานี.}\csromangathastanza{กุเวรสฺส โข ปน มาริส มหาราชสฺส วิสาณา นาม ราชธานี,}\csromangathastanza{ตสฺมา กุเวโร มหาราชา, เวสฺสวโณติ ปวุจฺจติ.}}

\csromanheader{2}{9. อาฏานาฏิยสุตฺต}
\csromangathameasure{ปจฺเจสนฺโต ปกาเสนฺติ, ตโตลา ตตฺตลา ตโตตลา. \\ โอชสิ เตชสิ ตโตชสี, สูโร ราชา อริฏฺโฐ เนมิ. \\ รหโทปิ ตตฺถ ธรณี นาม, ยโต เมฆา ปวสฺสนฺติ. \\ วสฺสา ยโต ปตายนฺติ, สภาปิ ตตฺถ สาลวตี นาม. \\ ยตฺถ ยกฺขา ปยิรุปาสนฺติ, ตตฺถ นิจฺจผลา รุกฺขา. \\ นานา ทิชคณา ยุตา, มยูรโกญฺจาภิรุทา.}
\csromangathagroup{\csromangathastanza{\csromanfolio{163}ปจฺเจสนฺโต ปกาเสนฺติ, ตโตลา ตตฺตลา ตโตตลา. \\ โอชสิ เตชสิ ตโตชสี, สูโร ราชา อริฏฺโฐ เนมิ.}\csromangathastanza{รหโทปิ ตตฺถ ธรณี นาม, ยโต เมฆา ปวสฺสนฺติ. \\ วสฺสา ยโต ปตายนฺติ, สภาปิ ตตฺถ สาลวตี\footnote{ภคลวตี (สี, สฺยา, อิ)} นาม.}\csromangathastanza{ยตฺถ ยกฺขา ปยิรุปาสนฺติ, ตตฺถ นิจฺจผลา รุกฺขา. \\ นานา ทิชคณา ยุตา, มยูรโกญฺจาภิรุทา.}}

\prose{โกกิลาทีหิ วคฺคุหิ.}

\csromangathameasure{ชีวญฺชีวกสทฺเทตฺถ, อโถ โอฏฺฐวจิตฺตกา. \\ กุกฺกุฏกา กุฬีรกา, วเน โปกฺขรสาตกา. \\ สุกสาฬิกสทฺเทตฺถ, ทณฺฑมาณวกานิ จ. \\ โสภติ สพฺพกาลํ สา, กุเวรนฬินี สทา. \\ อิโต สา อุตฺตรา ทิสา, อิติ นํ อาจิกฺขตี ชโน. \\ ยํ ทิสํ อภิปาเลติ, มหาราชา ยสสฺสิ โส. \\ ยกฺขานญฺจ อธิปติ, กุเวโร อิติ นามโส. \\ รมตี นจฺจคีเตหิ, ยกฺเขเหว ปุรกฺขโต. \\ ปุตฺตาปิ ตสฺส พหโว, เอกนามาติ เม สุตํ. \\ อสีติ ทส เอโก จ, อินฺทนามา มหพฺพลา. \\ เต จาปิ พุทฺธํ ทิสฺวาน, พุทฺธํ อาทิจฺจพนฺธุนํ. \\ ทูรโตว นมสฺสนฺติ, มหนฺตํ วีตสารทํ. \\ นโม เต ปุริสาชญฺญ, นโม เต ปุริสุตฺตม. \\ กุสเลน สเมกฺขสิ, อมนุสฺสาปิ ตํ วนฺทนฺติ. \\ สุตํ เนตํ อภิณฺหโส, ตสฺมา เอวํ วเทมเส. \\ ชินํ วนฺทถ โคตมํ, ชินํ วนฺทาม โคตมํ. \\ วิชฺชาจรณสมฺปนฺนํ, พุทฺธํ วนฺทาม โคตมนฺติ.}
\csromangathagroup{\csromangathastanza{ชีวญฺชีวกสทฺเทตฺถ, อโถ โอฏฺฐวจิตฺตกา. \\ กุกฺกุฏกา\footnote{กุกุตฺถกา (สี, อิ)} กุฬีรกา, วเน โปกฺขรสาตกา.}\csromangathastanza{สุกสาฬิกสทฺเทตฺถ, ทณฺฑมาณวกานิ จ. \\ โสภติ สพฺพกาลํ สา, กุเวรนฬินี สทา.}\csromangathastanza{อิโต สา อุตฺตรา ทิสา, อิติ นํ อาจิกฺขตี ชโน. \\ ยํ ทิสํ อภิปาเลติ, มหาราชา ยสสฺสิ โส.}\csromangathastanza{ยกฺขานญฺจ อธิปติ, กุเวโร อิติ นามโส. \\ รมตี นจฺจคีเตหิ, ยกฺเขเหว ปุรกฺขโต.}\csromangathastanza{ปุตฺตาปิ ตสฺส พหโว, เอกนามาติ เม สุตํ. \\ อสีติ ทส เอโก จ, อินฺทนามา มหพฺพลา.}\csromangathastanza{เต จาปิ พุทฺธํ ทิสฺวาน, พุทฺธํ อาทิจฺจพนฺธุนํ. \\ ทูรโตว นมสฺสนฺติ, มหนฺตํ วีตสารทํ.}\csromangathastanza{นโม เต ปุริสาชญฺญ, นโม เต ปุริสุตฺตม. \\ กุสเลน สเมกฺขสิ, อมนุสฺสาปิ ตํ วนฺทนฺติ.}\csromangathastanza{สุตํ เนตํ อภิณฺหโส, ตสฺมา เอวํ วเทมเส. \\ ชินํ วนฺทถ โคตมํ, ชินํ วนฺทาม โคตมํ. \\ วิชฺชาจรณสมฺปนฺนํ, พุทฺธํ วนฺทาม โคตมนฺติ.}}

\prose{อยํ โข สา มาริส อาฏานาฏิยา รกฺขา ภิกฺขูนํ ภิกฺขุนีนํ อุปาสกานํ อุปาสิกานํ คุตฺติยา รกฺขาย อวิหิํสาย ผาสุวิหาราย.}

\proseitem{282}{\csromanfolio{164}ยสฺส กสฺสจิ มาริส ภิกฺขุสฺส วา ภิกฺขุนิยา วา อุปาสกสฺส วา อุปาสิกาย วา อยํ อาฏานาฏิยา รกฺขา สุคฺคหิตา ภวิสฺสติ สมตฺตา ปริยาปุตา\footnote{ปริยาปุฏา (ก)}.\csromanspacer{} ตํ เจ อมนุสฺโส ยกฺโข วา ยกฺขินี วา ยกฺขโปตโก วา ยกฺขโปติกา วา ยกฺขมหามตฺโต วา ยกฺขปาริสชฺโช วา ยกฺขปจาโร วา, คนฺธพฺโพ วา คนฺธพฺพี วา คนฺธพฺพโปตโก วา คนฺธพฺพโปติกา วา คนฺธพฺพมหามตฺโต วา คนฺธพฺพปาริสชฺโช วา คนฺธพฺพปจาโร วา, กุมฺภณฺโฑ วา กุมฺภณฺฑี วา กุมฺภณฺฑโปตโก วา กุมฺภณฺฑโปติกา วา กุมฺภณฺฑมหามตฺโต วา กุมฺภณฺฑปาริสชฺโช วา กุมฺภณฺฑปจาโร วา, นาโค วา นาคี วา นาคโปตโก วา นาคโปติกา วา นาคมหามตฺโต วา นาคปาริสชฺโช วา นาคปจาโร วา ปทุฏฺฐจิตฺโต ภิกฺขุํ วา ภิกฺขุนิํ วา อุปาสกํ วา อุปาสิกํ วา คจฺฉนฺตํ วา อนุคจฺเฉยฺย, ฐิตํ วา อุปติฏฺเฐยฺย, นิสินฺนํ วา อุปนิสีเทยฺย, นิปนฺนํ วา อุปนิปชฺเชยฺย.\csromanspacer{} น เม โส มาริส อมนุสฺโส ลเภยฺย คาเมสุ วา นิคเมสุ วา สกฺการํ วา ครุการํ วา.\csromanspacer{} น เม โส มาริส อมนุสฺโส ลเภยฺย อาฬกมนฺทาย นาม ราชธานิยา วตฺถุํ วา วาสํ วา.\csromanspacer{} น เม โส มาริส อมนุโส ลเภยฺย ยกฺขานํ สมิติํ คนฺตุํ.\csromanspacer{} อปิสฺสุ นํ มาริส อมนุสฺสา อนาวยฺหํปิ นํ กเรยฺยุํ อวิวยฺหํ.\csromanspacer{} อปิสฺสุ นํ มาริส อมนุสฺสา อตฺตาหิปิ ปริปุณฺณาหิ ปริภาสาหิ ปริภาเสยฺยุํ.\csromanspacer{} อปิสฺสุ นํ มาริส อมนุสฺสา ริตฺตํปิสฺส ปตฺตํ สีเส นิกฺกุชฺเชยฺยุํ.\csromanspacer{} อปิสฺสุ นํ มาริส อมนุสฺสา สตฺตธาปิสฺส มุทฺธํ ผาเลยฺยุํ.}

\prose{สนฺติ หิ มาริส อมนุสฺสา จณฺฑา รุทฺธา\footnote{รุทฺทา (สี, อิ)} รภสา, เต เนว มหาราชานํ อาทิยนฺติ, น มหาราชานํ ปุริสกานํ อาทิยนฺติ, น มหาราชานํ ปุริสกานํ ปุริสกานํ อาทิยนฺติ.\csromanspacer{} เตโข เต มาริส อมนุสฺสา มหาราชานํ อวรุทฺธานาม วุจฺจนฺติ.\csromanspacer{} เสยฺยถาปิ มาริส รญฺโญ มาคธสฺส วิชิเต มหาโจรา.\csromanspacer{} เต เนว รญฺโญ มาคธสฺส อาทิยนฺติ, น รญฺโญ มาคธสฺส ปุริสกานํ อาทิยนฺติ, น รญฺโญ มาคธสฺส ปุริสกานํ ปุริสกานํ อาทิยนฺติ.\csromanspacer{} เต โข เต มาริส มหาโจรา รญฺโญ มาคธสฺส อวรุทฺธา นาม วุจฺจนฺติ.\csromanspacer{} เอวเมว โข มาริส สนฺติ อมนุสฺสา จณฺฑา รุทฺธา รภสา, เต เนว มหาราชานํ อาทิยนฺติ, น มหาราชานํ ปุริสกานํ อาทิยนฺติ, น มหาราชานํ ปุริสกานํ ปุริสกานํ อาทิยนฺติ.}

\vspace{\tipitakaparskip}
\csromancenter{อาฏานาฏิยสุตฺต เต โข เต มาริส อมนุสฺสา มหาราชานํ อวรุทฺธา นาม วุจฺจนฺติ.\csromanspacer{} โย หิ โกจิ มาริส อมนุสฺโส ยกฺโข วา ยกฺขินี วา ฯเปฯ คนฺธพฺโพ วา คนฺธพฺพี วา ฯเปฯ กุมฺภณฺโฑ วา กุมฺภณฺฑี วา ฯเปฯ นาโค วา นาคี วา นาคโปตโก วา นาคโปติกา วา นาคมหามตฺโต วา นาคปาริสชฺโช วา นาคปจาโร วา ปทุฏฺฐจิตฺโต ภิกฺขุํ วา ภิกฺขุนิํ วา อุปาสกํ วา อุปาสิกํ วา คจฺฉนฺตํ วา อนุคจฺเฉยฺย, ฐิตํ วา อุปติฏฺเฐยฺย, นิสินฺนํ วา อุปนิสีเทยฺย, นิปนฺนํ วา อุปนิปชฺเชยฺย.\csromanspacer{} อิเมสํ ยกฺขานํ มหายกฺขานํ เสนาปตีนํ มหาเสนาปตีนํ อุชฺฌาเปตพฺพํ วิกฺกนฺทิตพฺพํ วิรวิตพฺพํ “อยํ ยกฺโข คณฺหาติ, อยํ ยกฺโข อาวิสติ, อยํ ยกฺโข เหเฐติ, อยํ ยกฺโข วิเหเฐติ, อยํ ยกฺโข หิํสติ, อยํ ยกฺโข วิหิํสติ, อยํ ยกฺโข น มุญฺจตี”ติ.}
\proseitem{283}{\csromanfolio{165}กตเมสํ ยกฺขานํ มหายกฺขานํ เสนาปตีนํ มหาเสนาปตีนํ.}

\csromangathameasure{อินฺโท โสโม วรุโณ จ, ภารทฺวาโช ปชาปติ. \\ จนฺทโน กามเสฏฺโฐ จ, กินฺนุฆณฺฑุ นิฆณฺฑุ จ. \\ ปนาโท โอปมญฺโญ จ, เทวสูโต จ มาตลิ. \\ จิตฺตเสโน จ คนฺธพฺโพ, นโฬ ราชา ชเนสโภ. \\ สาตาคิโร เหมวโต, ปุณฺณโก กรติโย คุโฬ. \\ สิวโก มุจลินฺโท จ, เวสฺสามิตฺโต ยุคนฺธโร. \\ โคปาโล สุปฺปโรโธ จ, หิริ เนตฺติ จ มนฺทิโย. \\ ปญฺจาลจณฺโฑ อาฬวโก, ปชฺชุนฺโน สุมโน สุมุโข. \\ ทธิมุโข มณิ มาณิวโร ทีโฆ, อโถ เสรีสโก สห.}
\csromangathagroup{\csromangathastanza{อินฺโท โสโม วรุโณ จ, ภารทฺวาโช ปชาปติ. \\ จนฺทโน กามเสฏฺโฐ จ, กินฺนุฆณฺฑุ นิฆณฺฑุ จ.}\csromangathastanza{ปนาโท โอปมญฺโญ จ, เทวสูโต จ มาตลิ. \\ จิตฺตเสโน จ คนฺธพฺโพ, นโฬ ราชา ชเนสโภ\footnote{ชโนสโภ (สฺยา)}.}\csromangathastanza{สาตาคิโร เหมวโต, ปุณฺณโก กรติโย คุโฬ. \\ สิวโก มุจลินฺโท จ, เวสฺสามิตฺโต ยุคนฺธโร.}\csromangathastanza{โคปาโล สุปฺปโรโธ จ\footnote{สุปฺปเคโธ จ (สี, สฺยา, อิ)}, หิริ เนตฺติ\footnote{หิรี เนตฺตี (สี, อิ)} จ มนฺทิโย. \\ ปญฺจาลจณฺโฑ อาฬวโก, ปชฺชุนฺโน สุมโน สุมุโข. \\ ทธิมุโข มณิ มาณิวโร\footnote{สุปฺปเคโธ จ (สี, สฺยา, อิ)} ทีโฆ, อโถ เสรีสโก สห.}}

\prose{อิเมสํ ยกฺขานํ มหายกฺขานํ เสนาปตีนํ มหาเสนาปตีนํ อุชฺฌาเปตพฺพํ วิกฺกนฺทิตพฺพํ วิรวิตพฺพํ “อยํ ยกฺโข คณฺหาติ, อยํ ยกฺโข อาวิสติ, อยํ ยกฺโข เหเฐติ, อยํ ยกฺโข วิเหเฐติ, อยํ ยกฺโข หิํสติ, อยํ ยกฺโข วิหิํสติ, อยํ ยกฺโข น มุญฺจตี”ติ.}

\prose{\csromanfolio{166}อยํ โข สา มาริส อาฏานาฏิยา รกฺขา ภิกฺขูนํ ภิกฺขุนีนํ อุปาสกานํ อุปาสิกานํ คุตฺติยา รกฺขาย อวิหิํสาย ผาสุวิหาราย.\csromanspacer{} หนฺท จ ทานิ มยํ มาริสา คจฺฉาม พหุกิจฺจา มยํ พหุกรณียาติ.\csromanspacer{} ยสฺสทานิ ตุเมฺห มหาราชาโน กาลํ มญฺญถาติ.}

\proseitem{284}{อถ โข จตฺตาโร มหาราชา อุฏฺฐายาสนา ภควนฺตํ อภิวาเทตฺวา ปทกฺขิณํ กตฺวา ตตฺเถวนฺตรธายิํสุ.\csromanspacer{} เตปิ โข ยกฺขา อุฏฺฐายาสนา อปฺเปกจฺเจ ภควนฺตํ อภิวาเทตฺวา ปทกฺขิณํ กตฺวา ตตฺเถวนฺตรธายิํสุ, อปฺเปกจฺเจ ภควตา สทฺธิํ สมฺโมทิํสุ, สมฺโมทนียํ กถํ สารณียํ วีติสาเรตฺวา ตตฺเถวนฺตรธายิํสุ, อปฺเปกจฺเจ เยน ภควา เตนญฺชลิํ ปณาเมตฺวา ตตฺเถวนฺตธายิํสุ, อปฺเปกจฺเจ นามโคตฺตํ สาเวตฺวา ตตฺเถวนฺตรธายิํสุ, อปฺเปกจฺเจ ตุณฺหีภูตา ตตฺเถวนฺตรธายิํสูติ.}

\nitthitamruled{ปฐมภาณวาโร นิฏฺฐิโต.}
\csromanmatikaanchor{mtk.124}
\csromantocmarkref{section}{ทุติยภาณวาร}{mtk.124}
\bookmark[dest={mtk.124},level=1]{ทุติยภาณวาร}
\csromanheader{1}{\textbf{ทุติยภาณวาร}}
\proseitem{285}{อถ โข ภควา ตสฺสา รตฺติยา อจฺจเยน ภิกฺขู อามนฺเตสิ\csromandash{}อิมํ ภิกฺขเว รตฺติํ จตฺตาโร มหาราชา มหติยา จ ยกฺขเสนาย มหติยา จ คนฺธพฺพเสนาย มหติยา จ กุมฺภณฺฑเสนาย มหติยา จ นาคเสนาย จตุทฺทิสํ รกฺขํ ฐเปตฺวา จตุทฺทิสํ คุมฺพํ ฐเปตฺวา จตุทฺทิสํ โอวรณํ ฐเปตฺวา อภิกฺกนฺตาย รตฺติยา อภิกฺกนฺตวณฺณา เกวลกปฺปํ คิชฺฌกูฏํ ปพฺพตํ โอภาเสตฺวา เยนาหํ เตนุปสงฺกมิํสุ, อุปสงฺกมิตฺวา มํ อภิวาเทตฺวา เอกมนฺตํ นิสีทิํสุ.\csromanspacer{} เตปิ โข ภิกฺขเว ยกฺขา อปฺเปกจฺเจ มํ อภิวาเทตฺวา เอกมนฺตํ นิสีทิํสุ, อปฺเปกจฺเจ มยา สทฺธิํ สมฺโมทิํสุ, สมฺโมทนียํ กถํ สารณียํ วีติสาเรตฺวา เอกมนฺตํ นิสีทิํสุ, อปฺเปกจฺเจ เยนาหํ เตนญฺชลิํ ปณาเมตฺวา เอกมนฺตํ นิสีทิํสุ, อปฺเปกจฺเจ นามโคตฺตํ สาเวตฺวา เอกมนฺตํ นิสีทิํสุ, อปฺเปกจฺเจ ตุณฺหีภูตา เอกมนฺตํ นิสีทิํสุ.}

\proseitem{286}{เอกมนฺตํ นิสินฺโน โข ภิกฺขเว เวสฺสวโณ มหาราชา มํ เอตทโวจ, สนฺติ หิ ภนฺเต อุฬารา ยกฺขา ภควโต อปฺปสนฺนา ฯเปฯ สนฺติ}

\vspace{\tipitakaparskip}
\csromancenter{อาฏานาฏิยสุตฺต หิ ภนฺเต.\csromanspacer{} นีจา ยกฺขา ภควโต ปสนฺนา.\csromanspacer{} เยภุยฺเยน โข ปน ภนฺเต ยกฺขา อปฺปสนฺนาเยว ภควโต.\csromanspacer{} ตํ กิสฺส เหตุ, ภควา หิ ภนฺเต ปาณาติปาตา เวรมณิยา ธมฺมํ เทเสติ ฯเปฯ สุราเมรยมชฺชปฺปมาทฏฺฐานา เวรมณิยา ธมฺมํ เทเสติ.\csromanspacer{} เยภุยฺเยน โข ปน ภนฺเต ยกฺขา อปฺปฏิวิรตาเยว ปาณาติปาตา ฯเปฯ อปฺปฏิวิรตา สุราเมรยมชฺชปฺปมาทฏฺฐานา.\csromanspacer{} เตสํ ตํ โหติ อปฺปิยํ อมนาปํ.\csromanspacer{} สนฺติ หิ ภนฺเต ภควโต สาวกา อรญฺญวนปตฺถานิ ปนฺตานิ เสนาสนานิ ปฏิเสวนฺติ อปฺปสทฺทานิ อปฺปนิคฺโฆสานิ วิชนวาตานิ มนุสฺสราหสฺเสยฺยกานิ ปฏิสลฺลานสารุปฺปานิ.\csromanspacer{} ตตฺถ สนฺติ อุฬารา ยกฺขา นิวาสิโน, เย อิมสฺมิํ ภควโต ปาวจเน อปฺปสนฺนา, เตสํ ปสาทาย อุคฺคณฺหาตุ ภนฺเต ภควา อาฏานาฏิยํ รกฺขํ ภิกฺขูนํ ภิกฺขุนีนํ อุปาสกานํ อุปาสิกานํ คุตฺติยา รกฺขาย อวิหิํสาย ผาสุวิหารายาติ.\csromanspacer{} อธิวาเสสิํ โข อหํ ภิกฺขเว ตุณฺหีภาเวน.\csromanspacer{} อถ โข ภิกฺขเว เวสฺสวโณ มหาราชา เม อธิวาสนํ วิทิตฺวา ตายํ เวลายํ อิมํ อาฏานาฏิยํ รกฺขํ อภาสิ.}
\csromanhangingbegin
\hangingheaditem{287}{\csromanfolio{167}วิปสฺสิสฺส จ นมตฺถุ, จกฺขุมนฺตสฺส สิรีมโต.}
\hangingbody{สิขิสฺสปิจ นมตฺถุ, สพฺพภูตานุกมฺปิโน.}
\hangingbody{เวสฺสภุสฺส จ นมตฺถุ, นฺหาตกสฺส ตปสฺสิโน.}
\hangingbody{นมตฺถุ กกุสนฺธสฺส, มารเสนาปมทฺทิโน.}
\hangingbody{โกณาคมนสฺส นมตฺถุ, พฺราหฺมณสฺส วุสีมโต.}
\hangingbody{กสฺสปสฺส จ นมตฺถุ, วิปฺปมุตฺตสฺส สพฺพธิ.}
\hangingbody{องฺคีรสสฺส นมตฺถุ, สกฺยปุตฺตสฺส สิรีมโต.}
\hangingbody{โย อิมํ ธมฺมํ เทเสสิ, สพฺพทุกฺขาปนูทนํ.}
\hangingbody{เย จาปิ นิพฺพุตา โลเก, ยถาภูตํ วิปสฺสิสุํ.}
\hangingbody{เต ชนา อปิสุณาถ, มหนฺตา วีตสารทา.}
\hangingbody{หิตํ เทวมนุสฺสานํ, ยํ นมสฺสนฺติ โคตมํ.}
\hangingbody{วิชฺชาจรณสมฺปนฺนํ, มหนฺตํ วีตสารทํ.}
\csromanhangingend

\csromanhangingbegin
\hangingheaditem{288}{ยโต อุคฺคจฺฉติ สูริโย, อาทิจฺโจ มณฺฑลี มหา.}
\hangingbody{ยสฺส จุคฺคจฺฉมานสฺส, สํวรีปิ นิรุชฺฌติ.}
\hangingbody{ยสฺส จุคฺคเต สูริเย, ทิวโสติ ปวุจฺจติ.}
\csromanhangingend

\csromangathameasure{รหโทปิ ตตฺถ คมฺภีโร, สมุทฺโท สริโตทโก. \\ เอวํ ตํ ตตฺถ ชานนฺติ, สมุทฺโท สริโตทโก. \\ อิโต สา ปุริมา ทิสา, อิติ นํ อาจิกฺขตี ชโน. \\ ยํ ทิสํ อภิปาเลติ, มหาราชา ยสสฺสิ โส. \\ คนฺธพฺพานํ อธิปติ, ธตรฏฺโฐติ นามโส. \\ รมตี นจฺจคีเตหิ, คนฺธพฺเพหิ ปุรกฺขโต. \\ ปุตฺตาปิ ตสฺส พหโว, เอกนามาติ เม สุตํ. \\ อสีติ ทส เอโก จ, อินฺทนามา มหพฺพลา. \\ เต จาปิ พุทฺธํ ทิสฺวาน, พุทฺธํ อาทิจฺจพนฺธุนํ. \\ ทูรโตว นมสฺสนฺติ, มหนฺตํ วีตสารทํ. \\ นโม เต ปุริสาชญฺญ, นโม เต ปุริสุตฺตม. \\ กุสเลน สเมกฺขสิ, อมนุสฺสาปิ ตํ วนฺทนฺติ. \\ สุตํ เนตํ อภิณฺหโส, ตสฺมา เอวํ วเทมเส. \\ ชินํ วนฺทถ โคตมํ, ชินํ วนฺทาม โคตมํ. \\ วิชฺชาจรณสมฺปนฺนํ, พุทฺธํ วนฺทาม โคตมํ.}
\csromangathagroup{\csromangathastanza{\csromanfolio{168}รหโทปิ ตตฺถ คมฺภีโร, สมุทฺโท สริโตทโก. \\ เอวํ ตํ ตตฺถ ชานนฺติ, สมุทฺโท สริโตทโก.}\csromangathastanza{อิโต สา ปุริมา ทิสา, อิติ นํ อาจิกฺขตี ชโน. \\ ยํ ทิสํ อภิปาเลติ, มหาราชา ยสสฺสิ โส.}\csromangathastanza{คนฺธพฺพานํ อธิปติ, ธตรฏฺโฐติ นามโส. \\ รมตี นจฺจคีเตหิ, คนฺธพฺเพหิ ปุรกฺขโต.}\csromangathastanza{ปุตฺตาปิ ตสฺส พหโว, เอกนามาติ เม สุตํ. \\ อสีติ ทส เอโก จ, อินฺทนามา มหพฺพลา.}\csromangathastanza{เต จาปิ พุทฺธํ ทิสฺวาน, พุทฺธํ อาทิจฺจพนฺธุนํ. \\ ทูรโตว นมสฺสนฺติ, มหนฺตํ วีตสารทํ.}\csromangathastanza{นโม เต ปุริสาชญฺญ, นโม เต ปุริสุตฺตม. \\ กุสเลน สเมกฺขสิ, อมนุสฺสาปิ ตํ วนฺทนฺติ.}\csromangathastanza{สุตํ เนตํ อภิณฺหโส, ตสฺมา เอวํ วเทมเส. \\ ชินํ วนฺทถ โคตมํ, ชินํ วนฺทาม โคตมํ. \\ วิชฺชาจรณสมฺปนฺนํ, พุทฺธํ วนฺทาม โคตมํ.}}

\csromanhangingbegin
\hangingheaditem{289}{เยน เปตา ปวุจฺจนฺติ, ปิสุณา ปิฏฺฐิมํสิกา.}
\hangingbody{ปาณาติปาติโน ลุทฺทา, โจรา เนกติกา ชนา.}
\hangingbody{อิโต สา ทกฺขินา ทิสา, อิติ นํ อาจิกฺขตี ชโน.}
\hangingbody{ยํ ทิสํ อภิปาเลติ, มหาราชา ยสสฺสิ โส.}
\hangingbody{กุมฺภณฺฑานํ อธิปติ, วิรูโฬฺห-อิติ นามโส.}
\hangingbody{รมตี นจฺจคีเตหิ, กุมฺภณฺเฑหิ ปุรกฺขโต.}
\hangingbody{ปุตฺตาปิ ตสฺส พหโว, เอกนามาติ เม สุตํ.}
\hangingbody{อสีติ ทส เอโก จ, อินฺทนามา มหพฺพลา.}
\hangingbody{เต จาปิ พุทฺธํ ทิสฺวาน, พุทฺธํ อาทิจฺจพนฺธุนํ.}
\hangingbody{ทูรโตว นมสฺสนฺติ, มหนฺตํ วีตสารทํ.}
\csromanhangingend

\csromanheader{2}{9. อาฏานาฏิยสุตฺต}
\csromangathameasure{นโม เต ปุริสาชญฺญ, นโม เต ปุริสุตฺตม. \\ กุสเลน สเมกฺขสิ, อมนุสฺสาปิ ตํ วนฺทนฺติ. \\ สุตํ เนตํ อภิณฺหโส, ตสฺมา เอวํ วเทมเส. \\ ชินํ วนฺทถ โคตมํ, ชินํ วนฺทฺนาม โคตมํ. \\ วิชฺชาจรณสมฺปนฺนํ, พุทฺธํ วนฺทาม โคตมํ.}
\csromangathagroup{\csromangathastanza{\csromanfolio{169}นโม เต ปุริสาชญฺญ, นโม เต ปุริสุตฺตม. \\ กุสเลน สเมกฺขสิ, อมนุสฺสาปิ ตํ วนฺทนฺติ.}\csromangathastanza{สุตํ เนตํ อภิณฺหโส, ตสฺมา เอวํ วเทมเส. \\ ชินํ วนฺทถ โคตมํ, ชินํ วนฺทฺนาม โคตมํ. \\ วิชฺชาจรณสมฺปนฺนํ, พุทฺธํ วนฺทาม โคตมํ.}}

\csromanhangingbegin
\hangingheaditem{290}{ยตฺถ โจคฺคจฺฉติ สูริโย, อาทิจฺโจ มณฺฑลี มหา.}
\hangingbody{ยสฺส โจคฺคจฺฉมานสฺส, ทิวโสปิ นิรุชฺฌติ.}
\hangingbody{ยสฺส โจคฺคเต สูริเย, สํวรีติ ปวุจฺจติ.}
\hangingbody{รหโทปิ ตตฺถ คมฺภีโร, สมุทฺโท สริโตทโก.}
\hangingbody{เอวํ ตํ ตตฺถ ชานนฺติ, สมุทฺโท สริโตทโก.}
\hangingbody{อิโต สา ปจฺฉิมา ทิสา, อิติ นํ อาจิกฺขตี ชโน.}
\hangingbody{ยํ ทิสํ อภิปาเลติ, มหาราชา ยสสฺสิ โส.}
\hangingbody{นาคานญฺจ อธิปติ, วิรูปกฺโขติ นามโส.}
\hangingbody{รมตี นจฺจคีเตหิ, นาเคเหว ปุรกฺขโต.}
\hangingbody{ปุตฺตาปิ ตสฺส พหโว, เอกนามาติ เม สุตํ.}
\hangingbody{อสีติ ทส เอโก จ, อินฺทนามา มหพฺพลา.}
\hangingbody{เตจาปิ พุทฺธํ ทิสฺวาน, พุทฺธํ อาทิจฺจพนฺธุนํ.}
\hangingbody{ทูรโตว นมสฺสนฺติ, มหนฺตํ วีตสารทํ.}
\hangingbody{นโม เต ปุริสาชญฺญ, นโม เต ปุริสุตฺตม.}
\hangingbody{กุสเลน สเมกฺขสิ, อมนุสฺสาปิ ตํ วนฺทนฺติ.}
\hangingbody{สุตํ เนตํ อภิณฺหโส, ตสฺมา เอวํ วเทมเส.}
\hangingbody{ชินํ วนฺทถ โคตมํ, ชินํ วนฺทาม โคตมํ.}
\hangingbody{วิชฺชาจรณสมฺปนฺนํ, พุทฺธํ วนฺทาม โคตมํ.}
\csromanhangingend

\csromanhangingbegin
\hangingheaditem{291}{เยน อุตฺตรกุรุโวฺห, มหาเนรุ สุทสฺสโน.}
\hangingbody{มนุสฺสา ตตฺถ ชายนฺติ, อมมา อปริคฺคหา.}
\csromanhangingend

\csromangathameasure{น เต พีชํ ปวปนฺติ, นาปิ นียนฺติ นงฺคลา. \\ อกฏฺฐปากิมํ สาลิํ, ปริภุญฺชนฺติ มานุสา. \\ อกณํ อถุสํ สุทฺธํ, สุคนฺธํ ตณฺฑุลปฺผลํ. \\ ตุณฺฑิกีเร ปจิตฺวาน, ตโต ภุญฺชนฺติ โภชนํ. \\ คาวิํ เอกขุรํ กตฺวา, อนุยนฺติ ทิโสทิสํ. \\ ปสุํ เอกขุรํ กตฺวา, อนุยนฺติ ทิโสทิสํ. \\ อิตฺถิํ วา วาหนํ กตฺวา, อนุยนฺติ ทิโสทิสํ. \\ ปุริสํ วาหนํ กตฺวา, อนุยนฺติ ทิโสทิสํ. \\ กุมาริํ วาหนํ กตฺวา, อนุยนฺติ ทิโสทิสํ. \\ กุมารํ วาหนํ กตฺวา, อนุยนฺติ ทิโสทิสํ. \\ เต ยาเน อภิรุหิตฺวา, \\ สพฺพา ทิสา อนุปริยายนฺติ. \\ ปจารา ตสฺส ราชิโน. \\ หตฺถิยานํ อสฺสยานํ, \\ ทิพฺพํ ยานํ อุปฏฺฐิตํ. \\ ปาสาทา สิวิกา เจว, \\ มหาราชสฺส ยสสฺสิโน. \\ ตสฺส จ นครา อหุ, \\ อนฺตลิกฺเข สุมาปิตา. \\ อาฏานาฏา กุสินาฏา ปรกุสินาฏา, \\ นาฏสุริยา ปรกุสิฏนาฏา. \\ อุตฺตเรน กสิวนฺโต, \\ ชโนฆมปเรน จ. \\ นวนวุติโย อมฺพร-อมฺพรวติโย, \\ อาฬกมนฺทา นาม ราชธานี. \\ กุเวรสฺส โข ปน มาริส มหาราชสฺส วิสาณา นาม ราชธานี, \\ ตสฺมา กุเวโร มหาราชา, เวสฺสวโณติ ปวุจฺจติ.}
\csromangathagroup{\csromangathastanza{\csromanfolio{170}น เต พีชํ ปวปนฺติ, นาปิ นียนฺติ นงฺคลา. \\ อกฏฺฐปากิมํ สาลิํ, ปริภุญฺชนฺติ มานุสา.}\csromangathastanza{อกณํ อถุสํ สุทฺธํ, สุคนฺธํ ตณฺฑุลปฺผลํ. \\ ตุณฺฑิกีเร ปจิตฺวาน, ตโต ภุญฺชนฺติ โภชนํ.}\csromangathastanza{คาวิํ เอกขุรํ กตฺวา, อนุยนฺติ ทิโสทิสํ. \\ ปสุํ เอกขุรํ กตฺวา, อนุยนฺติ ทิโสทิสํ.}\csromangathastanza{อิตฺถิํ วา วาหนํ กตฺวา, อนุยนฺติ ทิโสทิสํ. \\ ปุริสํ วาหนํ กตฺวา, อนุยนฺติ ทิโสทิสํ.}\csromangathastanza{กุมาริํ วาหนํ กตฺวา, อนุยนฺติ ทิโสทิสํ. \\ กุมารํ วาหนํ กตฺวา, อนุยนฺติ ทิโสทิสํ.}\csromangathastanza{เต ยาเน อภิรุหิตฺวา, \\ สพฺพา ทิสา อนุปริยายนฺติ. \\ ปจารา ตสฺส ราชิโน. \\ หตฺถิยานํ อสฺสยานํ,}\csromangathastanza{ทิพฺพํ ยานํ อุปฏฺฐิตํ. \\ ปาสาทา สิวิกา เจว, \\ มหาราชสฺส ยสสฺสิโน. \\ ตสฺส จ นครา อหุ,}\csromangathastanza{อนฺตลิกฺเข สุมาปิตา. \\ อาฏานาฏา กุสินาฏา ปรกุสินาฏา, \\ นาฏสุริยา ปรกุสิฏนาฏา. \\ อุตฺตเรน กสิวนฺโต,}\csromangathastanza{ชโนฆมปเรน จ. \\ นวนวุติโย อมฺพร-อมฺพรวติโย, \\ อาฬกมนฺทา นาม ราชธานี. \\ กุเวรสฺส โข ปน มาริส มหาราชสฺส วิสาณา นาม ราชธานี,}\csromangathastanza{ตสฺมา กุเวโร มหาราชา, เวสฺสวโณติ ปวุจฺจติ.}}

\csromanheader{2}{9. อาฏานาฏิยสุตฺต}
\csromangathameasure{ปจฺเจสนฺโต ปกาเสนฺติ, ตโตลา ตตฺตลา ตโตตลา. \\ โอชสิ เตชสิ ตโตชสี, สูโร ราชา อริฏฺโฐ เนมิ. \\ รหโทปิ ตตฺถ ธรณี นาม, ยโต เมฆา ปวสฺสนฺติ. \\ วสฺสา ยโต ปตายนฺติ, สภาปิ ตตฺถา สาลวตี นาม. \\ ยตฺถ ยกฺขา ปยิรุปาสนฺติ, ตตฺถ นิจฺจผลา รุกฺขา. \\ นานา ทิชคณา ยุตา, มยูรโกญฺจาภิรุทา.}
\csromangathagroup{\csromangathastanza{\csromanfolio{171}ปจฺเจสนฺโต ปกาเสนฺติ, ตโตลา ตตฺตลา ตโตตลา. \\ โอชสิ เตชสิ ตโตชสี, สูโร ราชา อริฏฺโฐ เนมิ.}\csromangathastanza{รหโทปิ ตตฺถ ธรณี นาม, ยโต เมฆา ปวสฺสนฺติ. \\ วสฺสา ยโต ปตายนฺติ, สภาปิ ตตฺถา สาลวตี นาม.}\csromangathastanza{ยตฺถ ยกฺขา ปยิรุปาสนฺติ, ตตฺถ นิจฺจผลา รุกฺขา. \\ นานา ทิชคณา ยุตา, มยูรโกญฺจาภิรุทา.}}

\prose{โกกิลาทีหิ วคฺคุหิ.}

\csromangathameasure{ชีวญฺชีวกสทฺเทตฺถ, อโถ โอฏฺฐวจิตฺตกา. \\ กุกฺกุฏกา กุฬีรกา, วเน โปกฺขรสาตกา. \\ สุกสาฬิก สทฺเทตฺถ, ทณฺฑมาณวกานิ จ. \\ โสภติ สพฺพกาลํ สา, กุเวรนฬินี สทา. \\ อิโต สา อุตฺตรา ทิสา, อิติ นํ อาจิกฺขตี ชโน. \\ ยํ ทิสํ อภิปาเลติ, มหาราชา ยสสฺสิ โส. \\ ยกฺขานญฺจ อธิปติ, กุเวโร อิติ นามโส. \\ รมตี นจฺจคีเตหิ, ยกฺเขเหว ปุรกฺขโต. \\ ปุตฺตาปิ ตสฺส พหโว, เอกนามาติ เม สุตํ. \\ อสีติ ทส เอโก จ, อินฺทนาม มหพฺพลา. \\ เต จาปิ พุทฺธํ ทิสฺวาน, พุทฺธํ อาทิจฺจพนฺธุนํ. \\ ทูรโตว นมสฺสนฺติ, มหนฺตํ วีตสารทํ. \\ นโม เต ปุริสาชญฺญ, นโม เต ปุริสุตฺตม. \\ กุสเลน สเมกฺขสิ, อมนุสฺสาปิ ตํ วนฺทนฺติ. \\ สุตํ เนตํ อภิณฺหโส, ตสฺมา เอวํ วเทมเส. \\ ชินํ วนฺทถ โคตมํ, ชินํ วนฺทาม โคตมํ. \\ วิชฺชาจรณสมฺปนฺนํ, พุทฺธํ วนฺทาม โคตมนฺติ.}
\csromangathagroup{\csromangathastanza{ชีวญฺชีวกสทฺเทตฺถ, อโถ โอฏฺฐวจิตฺตกา. \\ กุกฺกุฏกา กุฬีรกา, วเน โปกฺขรสาตกา.}\csromangathastanza{สุกสาฬิก สทฺเทตฺถ, ทณฺฑมาณวกานิ จ. \\ โสภติ สพฺพกาลํ สา, กุเวรนฬินี สทา.}\csromangathastanza{อิโต สา อุตฺตรา ทิสา, อิติ นํ อาจิกฺขตี ชโน. \\ ยํ ทิสํ อภิปาเลติ, มหาราชา ยสสฺสิ โส.}\csromangathastanza{ยกฺขานญฺจ อธิปติ, กุเวโร อิติ นามโส. \\ รมตี นจฺจคีเตหิ, ยกฺเขเหว ปุรกฺขโต.}\csromangathastanza{ปุตฺตาปิ ตสฺส พหโว, เอกนามาติ เม สุตํ. \\ อสีติ ทส เอโก จ, อินฺทนาม มหพฺพลา.}\csromangathastanza{เต จาปิ พุทฺธํ ทิสฺวาน, พุทฺธํ อาทิจฺจพนฺธุนํ. \\ ทูรโตว นมสฺสนฺติ, มหนฺตํ วีตสารทํ.}\csromangathastanza{นโม เต ปุริสาชญฺญ, นโม เต ปุริสุตฺตม. \\ กุสเลน สเมกฺขสิ, อมนุสฺสาปิ ตํ วนฺทนฺติ.}\csromangathastanza{สุตํ เนตํ อภิณฺหโส, ตสฺมา เอวํ วเทมเส. \\ ชินํ วนฺทถ โคตมํ, ชินํ วนฺทาม โคตมํ. \\ วิชฺชาจรณสมฺปนฺนํ, พุทฺธํ วนฺทาม โคตมนฺติ.}}

\proseitem{292}{อยํ โข สา มาริส อาฏานาฏิยา รกฺขา ภิกฺขูนํ ภิกฺขุนีนํ อุปาสกานํ อุปาสิกานํ คุตฺติยา รกฺขาย อวิหิํสาย ผาสุวิหาราย.\csromanspacer{} ยสฺส กสฺสจิ มาริส ภิกฺขุสฺส วา ภิกฺขุนิยา วา อุปาสกสฺส วา \csromanfolio{172}อุปาสิกาย วา อยํ อาฏานาฏิยา รกฺขา สุคฺคหิตา ภวิสฺสติ สมตฺตา ปริยาปุตา.\csromanspacer{} ตํ เจ อมนุสฺโส ยกฺโข วา ยกฺขินี วา ฯเปฯ คนฺธพฺโพ วา คนฺธพฺพี วา ฯเปฯ กุมฺภณฺโฑ วา กุมฺภณฺฑี วา ฯเปฯ นาโค วา นาคี วา นาคโปตโก วา นาคโปติกา วา นาคมหามตฺโต วา นาคปาริสชฺโช วา นาคปจาโร วา ปทุฏฺฐจิตฺโต ภิกฺขุํ วา ภิกฺขุนิํ วา อุปาสกํ วา อุปาสิกํ วา คจฺฉนฺตํ วา อนุคจฺเฉยฺย, ฐิตํ วา อุปติฏฺเฐยฺย, นิสินฺนํ วา อุปนิสีเทยฺย, นิปนฺนํ วา อุปนิปชฺเชยฺย.\csromanspacer{} น เม โส มาริส อมนุสฺโส ลเภยฺย คาเมสุ วา นิคเมสุ วา สกฺการํ วา ครุการํ วา.\csromanspacer{} น เม โส มาริส อมนุสฺโส ลเภยฺย อาฬกมนฺทาย นาม ราชธานิยา วตฺถุํ วา วาสํ วา.\csromanspacer{} น เม โส มาริส อมนุสฺโส ลเภยฺย ยกฺขานํ สมิติํ คนฺตุํ.\csromanspacer{} อปิสฺสุ นํ มาริส อมนุสฺสา อนาวยฺหํปิ นํ กเรยฺยุํ อวิวยฺหํ.\csromanspacer{} อปิสฺสุ นํ มาริส อมนุสฺสา อตฺตาหิ ปริปุณฺณาหิ ปริภาสาหิ ปริภาเสยฺยุํ.\csromanspacer{} อปิสฺสุ นํ มาริส อมนุสฺสา ริตฺตํปิสฺส ปตฺตํ สีเส นิกฺกุชฺเชยฺยุํ.\csromanspacer{} อปิสฺสุ นํ มาริส อมนุสฺสา สตฺตธาปิสฺส มุทฺธํ ผาเลยฺยุํ.\csromanspacer{} สนฺติ หิ มาริส อมนุสฺสา จณฺฑา รุทฺธา รภสา, เต เนว มหาราชานํ อาทิยนฺติ, น มหาราชานํ ปุริสกานํ อาทิยนฺติ, น มหาราชานํ ปุริสกานํ ปุริสกานํ อาทิยนฺติ.\csromanspacer{} เต โข เต มาริส อมนุสฺสา มหาราชานํ อวรุทฺธา นาม วุจฺจนฺติ.\csromanspacer{} เสยฺยถาปิ มาริส รญฺโญ มาคธสฺส วิชิเต มหาโจรา.\csromanspacer{} เต เนว รญฺโญ มาคธสฺส อาทิยนฺติ, น รญฺโญ มาคธสฺส ปุริสกานํ อาทิยนฺติ, น รญฺโญ มาคธสฺส ปุริสกานํ ปุริสกานํ อาทิยนฺติ.\csromanspacer{} เต โข เต มาริส มหาโจรา รญฺโญ มาคธสฺส อวรุทฺธา นาม วุจฺจนฺติ.\csromanspacer{} เอวเมว โข มาริส สนฺติ อมนุสฺสา จณฺฑา รุทฺธา รภสา, เต เนว มหาราชานํ อาทิยนฺติ, น มหาราชานํ ปุริสกานํ อาทิยนฺติ, น มหาราชานํ ปุริสกานํ ปุริสกานํ อาทิยนฺติ.\csromanspacer{} เต โข เต มาริส อมนุสฺสา มหาราชานํ อวรุทฺธา นาม วุจฺจนฺติ.\csromanspacer{} โย หิ โกจิ มาริส อมนุสฺโส ยกฺโข วา ยกฺขินี วา ฯเปฯ คนฺธพฺโพ วา คนฺธพฺพี วา ฯเปฯ กุมฺภณฺโฑ วา กุมฺภณฺฑี วา ฯเปฯ นาโค วา นาคี วา ฯเปฯ ปทุฏฺฐจิตฺโต ภิกฺขุํ วา ภิกฺขุนิํ วา อุปาสกํ วา อุปาสิกํ วา คจฺฉนฺตํ วา อุปคจฺเฉยฺย, ฐิตํ วา อุปติฏฺเฐยฺย, นิสินฺนํ วา อุปนิสีเทยฺย, นิปนฺนํ วา อุปนิปชฺเชยฺย.\csromanspacer{} อิเมสํ ยกฺขานํ มหายกฺขานํ เสนาปตีนํ มหาเสนาปตีนํ อุชฺฌาเปตพฺพํ วิกฺกนฺทิตพฺพํ วิรวิตพฺพํ “อยํ ยกฺโข คณฺหาติ, อยํ ยกฺโข อาวิสติ, อยํ}

\vspace{\tipitakaparskip}
\csromancenter{อาฏานาฏิยสุตฺต ยกฺโข เหเฐติ, อยํ ยกฺโข วิเหเฐติ, อยํ ยกฺโข หิํสติ, อยํ ยกฺโข วิหิํสติ, อยํ ยกฺโข น มุญฺจตี”ติ.}
\proseitem{293}{\csromanfolio{173}กตเมสํ ยกฺขานํ มหายกฺขานํ เสนาปตีนํ มหาเสนาปตีนํ.}

\csromangathameasure{อินฺโท โสโม วรุโณ จ, ภารทฺวาโช ปชาปติ. \\ จนฺทโน กามเสฏฺโฐ จ, กินฺนุฆณฺฑุ นิฆณฺฑุ จ. \\ ปนาโท โอปมญฺโญ จ, เทวสูโต จ มาตลิ. \\ จิตฺตเสโน จ คนฺธพฺโพ, นโฬ ราชา ชเนสโภ. \\ สาตาคิโร เหวมโต, ปุณฺณโก กรติโย คุโฬ. \\ สิวโก มุจลินฺโท จ, เวสฺสามิตฺโต ยุคนฺธโร. \\ โคปาโล สุปฺปโรโธ จ, หิริ เนตฺติ จ มนฺทิโย. \\ ปญฺจาลจณฺโฑ อาฬวโก, ปชฺชุนฺโน สุมโน สุมุโข. \\ ทธิมุโข มณิ มาณิวโร ทีโฆ, อโถ เสรีสโก สห.}
\csromangathagroup{\csromangathastanza{อินฺโท โสโม วรุโณ จ, ภารทฺวาโช ปชาปติ. \\ จนฺทโน กามเสฏฺโฐ จ, กินฺนุฆณฺฑุ นิฆณฺฑุ จ.}\csromangathastanza{ปนาโท โอปมญฺโญ จ, เทวสูโต จ มาตลิ. \\ จิตฺตเสโน จ คนฺธพฺโพ, นโฬ ราชา ชเนสโภ.}\csromangathastanza{สาตาคิโร เหวมโต, ปุณฺณโก กรติโย คุโฬ. \\ สิวโก มุจลินฺโท จ, เวสฺสามิตฺโต ยุคนฺธโร.}\csromangathastanza{โคปาโล สุปฺปโรโธ จ, หิริ เนตฺติ จ มนฺทิโย. \\ ปญฺจาลจณฺโฑ อาฬวโก, ปชฺชุนฺโน สุมโน สุมุโข. \\ ทธิมุโข มณิ มาณิวโร ทีโฆ, อโถ เสรีสโก สห.}}

\prose{อิเมสํ ยกฺขานํ มหายกฺขานํ เสนาปตีนํ มหาเสนาปตีนํ อุชฺฌาเปตพฺพํ วิกฺกนฺทิตพฺพํ วิรวิตพฺพํ “อยํ ยกฺโข คณฺหาติ, อยํ ยกฺโข อาวิสติ, อยํ ยกฺโข เหเฐติ, อยํ ยกฺโข วิเหเฐติ, อยํ ยกฺโข หิํสติ, อยํ ยกฺโข วิหิํสติ, อยํ ยกฺโข น มุญฺจตี”ติ.\csromanspacer{} อยํ โข มาริส อาฏานาฏิยา รกฺขา ภิกฺขูนํ ภิกฺขุนีนํ อุปาสกานํ อุปาสิกานํ คุตฺติยา รกฺขาย อวิหิํสาย ผาสุวิหาราย.\csromanspacer{} หนฺท จ ทานิ มยํ มาริส คจฺฉาม, พหุกิจฺจา มยํ พหุกรณียาติ.\csromanspacer{} ยสฺสทานิ ตุเมฺห มหาราชาโน กาลํ มญฺญถาติ.}

\proseitem{294}{อถ โข ภิกฺขเว จตฺตาโร มหาราชา อุฏฺฐายาสนา มํ อภิวาเทตฺวา ปทกฺขิณํ กตฺวา ตตฺเถวนฺตรธายิํสุ.\csromanspacer{} เตปิ โข ภิกฺขเว ยกฺขา อุฏฺฐายาสนา อปฺเปกจฺเจ มํ อภิวาเทตฺวา ปทกฺขิณํ กตฺวา ตตฺเถวนฺตรธายิํสุ, อปฺเปกจฺเจ มยา สทฺธิํ สมฺโมทิํสุ, สมฺโมทนียํ กถํ สารณียํ วีติสาเรตฺวา ตตฺเถวนฺตรธายิํสุ, อปฺเปกจฺเจ เยนาหํ เตนญฺชลิํ ปณาเมตฺวา ตตฺเถวนฺตรธายิํสุ, อปฺเปกจฺเจ \csromanfolio{174}นามโคตฺตํ สาเวตฺวา ตตฺเถวนฺตรธายิํสุ, อปฺเปกจฺเจ ตุณฺหีภูตา ตตฺเถวนฺตรธายิํสุ.}

\proseitem{295}{อุคฺคณฺหาถ ภิกฺขเว อาฏานาฏิยํ รกฺขํ, ปริยาปุณาถ ภิกฺขเว อาฏานาฏิยํ รกฺขํ, ธาเรถ ภิกฺขเว อาฏานาฏิยํ รกฺขํ, อตฺถสํหิตา\footnote{อตฺถสํหิตายํ (สฺยา)} ภิกฺขเว อาฏานาฏิยา รกฺขา ภิกฺขูนํ ภิกฺขุนีนํ อุปาสกานํ อุปาสิกานํ คุตฺติยา รกฺขาย อวิหิํสาย ผาสุวิหารายาติ.\csromanspacer{} อิทมโวจ ภควา.\csromanspacer{} อตฺตมนา เต ภิกฺขู ภควโต ภาสิตํ อภินนฺทุนฺติ.}

\nitthitam{\textbf{อาฏานาฏิยสุตฺตํ นิฏฺฐิตํ นวมํ.}}
\csromanmatikaanchor{mtk.125}
\csromantocmarknopage{chapter}{10. สงฺคีติสุตฺต}{mtk.125}
\bookmark[dest={mtk.125},level=0]{10. สงฺคีติสุตฺต}
\chapterheadpageread{10. สงฺคีติสุตฺต}
\proseitem{296}{\csromanfolio{175}เอวํ เม สุตํ\csromandash{}เอกํ สมยํ ภควา มลฺเลสุ จาริกํ จรมาโน มหตา ภิกฺขุสํเฆน สทฺธิํ ปญฺจมตฺเตหิ ภิกฺขุสเตหิ เยน ปาวา นาม มลฺลานํ นครํ ตทวสริ, ตตฺร สุทํ ภควา ปาวายํ วิหรติ จุนฺทสฺส กมฺมารปุตฺตสฺส อมฺพวเน.}

\csromanmatikaanchor{mtk.126}
\csromantocmarkref{section}{อุพฺภตกนวสนฺธาคาร}{mtk.126}
\bookmark[dest={mtk.126},level=1]{อุพฺภตกนวสนฺธาคาร}
\csromanheader{1}{\textbf{อุพฺภตกนวสนฺธาคาร}}
\proseitem{297}{เตน โข ปน สมเยน ปาเวยฺยกานํ มลฺลานํ อุพฺภตกํ นาม นวํ สนฺธาคารํ\footnote{สนฺถาคารํ (สี, อิ), สณฺฐาคารํ (สฺยา, กํ)} อจิรการิตํ โหติ อนชฺฌาวุฏฺฐํ\footnote{อนชฺฌาวุตฺถํ (สี, สฺยา, อิ, ก) 3. สพฺพสนฺถริํ สนฺถตํ (ก)} สมเณน วา พฺราหฺมเณน วา เกนจิ วา มนุสฺสภูเตน.\csromanspacer{} อสฺโสสุํ โข ปาเวยฺยกา มลฺลา “ภควา กิร มลฺเลสุ จาริกํ จรมาโน มหตา ภิกฺขุสํเฆน สทฺธิํ ปญฺจมตฺเตหิ ภิกฺขุสเตหิ ปาวํ อนุปฺปตฺโต ปาวายํ วิหรติ จุนฺทสฺส กมฺมารปุตฺตสฺส อมฺพวเน”ติ.\csromanspacer{} อถ โข ปาเวยฺยกา มลฺลา เยน ภควา เตนุปสงฺกมิํสุ, อุปสงฺกมิตฺวา ภควนฺตํ อภิวาเทตฺวา เอกมนฺตํ นิสีทิํสุ, เอกมนฺตํ นิสินฺนา โข ปาเวยฺยกา มลฺลา ภควนฺตํ เอตทโวจุํ “อิธ ภนฺเต ปาเวยฺยกานํ มลฺลานํ อุพฺภตกํ นาม นวํ สนฺธาคารํ อจิรการิตํ โหติ อนชฺฌาวุฏฺฐํ สมเณน วา พฺราหฺมเณน วา เกนจิ วา มนุสฺสภูเตน.\csromanspacer{} ตญฺจ โข ภนฺเต ภควา ปฐมํ ปริภุญฺชตุ, ภควตา ปฐมํ ปริภุตฺตํ ปจฺฉา ปาเวยฺยกา มลฺลา ปริภุญฺชิสฺสนฺติ.\csromanspacer{} ตทสฺส ปาเวยฺยกานํ มลฺลานํ ทีฆรตฺตํ หิตาย สุขายา”ติ.\csromanspacer{} อธิวาเสสิ โข ภควา ตุณฺหีภาเวน.}

\proseitem{298}{อถ โข ปาเวยฺยกา มลฺลา ภควโต อธิวาสนํ วิทิตฺวา อุฏฺฐายาสนา ภควนฺตํ อภิวาเทตฺวา ปทกฺขิณํ กตฺวา เยน สนฺธาคารํ เตนุปสงฺกมิํสุ, อุปสงฺกมิตฺวา สพฺพสนฺถริํ3 สนฺธาคารํ สนฺถริตฺวา ภควโต อาสนานิ ปญฺญาเปตฺวา อุทกมณิกํ ปติฏฺฐเปตฺวา เตลปทีปํ อาโรเปตฺวา เยน ภควา เตนุปสงฺกมิํสุ, อุปสงฺกมิตฺวา ภควนฺตํ อภิวาเทตฺวา เอกมนฺตํ อฏฺฐํสุ, เอกมนฺตํ ฐิตา โข เต ปาเวยฺยกา มลฺลา ภควนฺตํ เอตทโวจุํ “สพฺพสนฺถริสนฺถตํ\footnote{สพฺพสนฺถริํ สนฺถตํ (สี, อิ, ก)} ภนฺเต สนฺธาคารํ, ภควโต \csromanfolio{176}อาสนานิ ปญฺญตฺตานิ, อุทกมณิโก ปติฏฺฐาปิโต, เตลปทีโป อาโรปิโต, ยสฺสทานิ ภนฺเต ภควา กาลํ มญฺญตี”ติ.}

\proseitem{299}{อถ โข ภควา นิวาเสตฺวา ปตฺตจีวรมาทาย สทฺธิํ ภิกฺขุสํเฆน เยน สนฺธาคารํ เตนุปสงฺกมิ, อุปสงฺกมิตฺวา ปาเท ปกฺขาเลตฺวา สนฺธาคารํ ปวิสิตฺวา มชฺฌิมํ ถมฺภํ นิสฺสาย ปุรตฺถาภิมุโข นิสีทิ.\csromanspacer{} ภิกฺขุสํโฆปิ โข ปาเท ปกฺขาเลตฺวา สนฺธาคารํ ปวิสิตฺวา ปจฺฉิมํ ภิตฺติํ นิสฺสย ปุรตฺถาภิมุโข นิสีทิ ภควนฺตํเยว ปุรกฺขตฺวา.\csromanspacer{} ปาเวยฺยกาปิ โข มลฺลา ปาเท ปกฺขาเลตฺวา สนฺธาคารํ ปวิสิตฺวา ปุรตฺถิมํ ภิตฺติํ นิสฺสาย ปจฺฉิมาภิมุขา นิสีทิํสุ ภควนฺตํเยว ปุรกฺขตฺวา.\csromanspacer{} อถ โข ภควา ปาเวยฺยเก มลฺเล พหุเทว รตฺติํ ธมฺมิยา กถาย สนฺทสฺเสตฺวา สมาทเปตฺวา สมุตฺเตเชตฺวา สมฺปหํเสตฺวา อุยฺโยเชสิ “อภิกฺกนฺตา โข วาเสฏฺฐา รตฺติ, ยสฺสทานิ ตุเมฺห กาลํ มญฺญถา”ติ, “เอวํ ภนฺเต”ติ โข ปาเวยฺยกา มลฺลา ภควโต ปฏิสฺสุตฺวา อุฏฺฐายาสนา ภควนฺตํ อภิวาเทตฺวา ปทกฺขิณํ กตฺวา ปกฺกมิํสุ.}

\proseitem{300}{อถ โข ภควา อจิรปกฺกนฺเตสุ ปาเวยฺยเกสุ มลฺเลสุ ตุณฺหีภูตํ ตุณฺหีภูตํ ภิกฺขุสํฆํ อนุวิโลเกตฺวา อายสฺมนฺตํ สาริปุตฺตํ อามนฺเตสิ “วิคตถินมิทฺโธ\footnote{วิคตถีนมิทฺโธ (สี, สฺยา, กํ, อิ)} โข สาริปุตฺต ภิกฺขุสํโฆ, ปฏิภาตุ ตํ สาริปุตฺต ภิกฺขูนํ ธมฺมี กถา, ปิฏฺฐิ เม อาคิลายติ, ตมหํ อายมิสฺสามี”ติ\footnote{อายเมยฺยามีติ (สฺยา, กํ)}. “เอวํ ภนฺเต”ติ โข อายสฺมา สาริปุตฺโต ภควโต ปจฺจสฺโสสิ.\csromanspacer{} อถ โข ภควา จตุคฺคุณํ สํฆาฏิํ ปญฺญาเปตฺวา ทกฺขิเณน ปสฺเสน สีหเสยฺยํ กปฺเปสิ ปาเท ปาทํ อจฺจาธาย สโต สมฺปชาโน อุฏฺฐานสญฺญํ มนสิ กริตฺวา.}

\csromanmatikaanchor{mtk.127}
\csromantocmarkref{section}{ภินฺนนิคณฺฐวตฺถุ}{mtk.127}
\bookmark[dest={mtk.127},level=1]{ภินฺนนิคณฺฐวตฺถุ}
\csromanheader{1}{\textbf{ภินฺนนิคณฺฐวตฺถุ}}
\proseitem{301}{เตน โข ปน สมเยน นิคณฺโฐ นาฏปุตฺโต ปาวายํ อธุนากาลงฺกโต โหติ.\csromanspacer{} ตสฺส กาลงฺกิริยาย ภินฺนา นิคณฺฐา เทฺวธิกชาตา\footnote{เทฺวฬฺหกชาตา (สฺยา, กํ)} ภณฺฑนชาตา กลหชาตา วิวาทาปนฺนา}

\vspace{\tipitakaparskip}
\csromancenter{สงฺคีติสุตฺต อญฺญมญฺญํ มุขสตฺตีหิ วิตุทนฺตา วิหรนฺติ\footnote{วิจรนฺติ (สฺยา, กํ)} “น ตฺวํ อิมํ ธมฺมวินยํ อาชานาสิ, อหํ อิมํ ธมฺมวินยํ อาชานามิ, กิํ ตฺวํ อิมํ ธมฺมวินยํ อาชานิสฺสสิ.\csromanspacer{} มิจฺฉาปฏิปนฺโน ตฺวมสิ, อหมสฺมิ สมฺมาปฏิปนฺโน.\csromanspacer{} สหิตํ เม, อสหิตํ เต.\csromanspacer{} ปุเรวจนียํ ปจฺฉา อวจ, ปจฺฉาวจนียํ ปุเร อวจ.\csromanspacer{} อธิจิณฺณํ เต วิปราวตฺตํ, อาโรปิโต เต วาโท, นิคฺคหิโต ตฺวมสิ, จร วาทปฺปโมกฺขาย, นิพฺเพเฐหิ วา สเจ ปโหสี”ติ.\csromanspacer{} วโธเยว โข มญฺเญ นิคณฺเฐสุ นาฏปุตฺติเยสุ วตฺตติ.\csromanspacer{} เยปิ\footnote{เยปิ เต (สี, อิ)} นิคณฺฐสฺส นาฏปุตฺตสฺส สาวกา คิหี โอทาตวสนา, เตปิ นิคณฺเฐสุ นาฏปุตฺติเยสุ นิพฺพินฺนรูปา วิรตฺตรูปา ปฏิวานรูปา, ยถา ตํ ทุรกฺขาเต ธมฺมวินเย ทุปฺปเวทิเต อนิยฺยานิเก อนุปสมสํวตฺตนิเก อสมฺมาสมฺพุทฺธปฺปเวทิเต ภินฺนถูเป อปฺปฏิสรเณ.}
\proseitem{302}{\csromanfolio{177}อถ โข อายสฺมา สาริปุตฺโต ภิกฺขู อามนฺเตสิ “นิคณฺโฐ อาวุโส นาฏปุตฺโต ปาวายํ อธุนากาลงฺกโต, ตสฺส กาลงฺกิริยาย ภินฺนา นิคณฺฐา เทฺวธิกชาตา ฯเปฯ ภินฺนถูเป อปฺปฏิสรเณ.\csromanspacer{} เอวเญฺหตํ อาวุโส โหติ ทุรกฺขาเต ธมฺมวินเย ทุปฺปเวทิเต อนิยฺยานิเก อนุปสมสํวตฺตนิเก อสมฺมาสมฺพุทฺธปฺปเวทิเต.\csromanspacer{} อยํ โข ปนาวุโส อมฺหากํ\footnote{อสฺมากํ (อิ)} ภควตา\footnote{ภควโต (ก-สี)} ธมฺโม สฺวากฺขาโต สุปฺปเวทิโต นิยฺยานิโก อุปสมสํวตฺตนิโก สมฺมาสมฺพุทฺธปฺปเวทิโต.\csromanspacer{} ตตฺถ สพฺเพเหว สงฺคายิตพฺพํ, น วิวทิตพฺพํ, ยถยิทํ พฺรหฺมจริยํ อทฺธนิยํ อสฺส จิรฏฺฐิติกํ, ตทสฺส พหุชนหิตาย พหุชนสุขาย โลกานุกมฺปาย อตฺถาย หิตาย สุขาย เทวมนุสฺสานํ.}

\prose{กตโม จาวุโส อมฺหากํ ภควตา4 ธมฺโม สฺวากฺขาโต สุปฺปเวทิโต นิยฺยานิโก อุปสมสํวตฺตนิโก สมฺมาสมฺพุทฺธปฺปเวทิโต, ยตฺถ สพฺเพเหว สงฺคายิตพฺพํ, น วิวทิตพฺพํ, ยถยิทํ พฺรหฺมจริยํ อทฺธนิยํ อสฺส จิรฏฺฐิติกํ, ตทสฺส พหุชนหิตาย พหุชนสุขาย โลกานุกมฺปาย อตฺถาย หิตาย สุขาย เทวมนุสฺสานํ.}

\csromanmatikaanchor{mtk.128}
\csromantocmarkref{section}{เอกก}{mtk.128}
\bookmark[dest={mtk.128},level=1]{เอกก}
\csromanheader{1}{\textbf{เอกก}}
\proseitem{303}{\csromanfolio{178}อตฺถิ โข อาวุโส เตน ภควตา ชานตา ปสฺสตา อรหตา สมฺมาสมฺพุทฺเธน เอโก ธมฺโม สมฺมทกฺขาโต.\csromanspacer{} ตตฺถ สพฺเพเหว สงฺคายิตพฺพํ, น วิวทิตพฺพํ, ยถยิทํ พฺรหฺมจริยํ อทฺธนิยํ อสฺส จิรฏฺฐิติกํ, ตทสฺส พหุชนหิตาย พหุชนสุขาย โลกานุกมฺปาย อตฺถาย หิตาย สุขาย เทวมนุสฺสานํ.\csromanspacer{} กตโม เอโก ธมฺโม.\csromanspacer{} สพฺเพ สตฺตา อาหารฏฺฐิตากา.\csromanspacer{} สพฺเพ สตฺตา สงฺขารฏฺฐิติกา.\csromanspacer{} อยํ โข อาวุโส เตน ภควตา ชานตา ปสฺสตา อรหตา สมฺมาสมฺพุทฺเธน เอโก ธมฺโม สมฺมทกฺขาโต.\csromanspacer{} ตตฺถ สพฺเพเหว สงฺคายิตพฺพํ, น วิวทิตพฺพํ, ยถยิทํ พฺรหฺมจริยํ อทฺธนิยํ อสฺส จิรฏฺฐิติกํ, ตทสฺส พหุชนหิตาย พหุชนสุขาย โลกานุกมฺปาย อตฺถาย หิตาย สุขาย เทวมนุสฺสานํ.}

\csromanmatikaanchor{mtk.129}
\csromantocmarkref{section}{ทุก}{mtk.129}
\bookmark[dest={mtk.129},level=1]{ทุก}
\csromanheader{1}{\textbf{ทุก}}
\proseitem{304}{อตฺถิ โข อาวุโส เตน ภควตา ชานตา ปสฺสตา อรหตา สมฺมาสมฺพุทฺเธน เทฺว ธมฺมา สมฺมทกฺขาตา.\csromanspacer{} ตตฺถ สพฺเพเหว สงฺคายิตพฺพํ, น วิวทิตพฺพํ, ยถยิทํ พฺรหฺมจริยํ อทฺธนิยํ อสฺส จิรฏฺฐิติกํ, ตทสฺส พหุชนหิตาย พหุชนสุขาย โลกานุกมฺปาย อตฺถาย หิตาย สุขาย เทวมนุสฺสานํ.\csromanspacer{} กตเม เทฺว\footnote{เทฺว ธมฺมา (สฺยา, กํ) เอวมุปริปิ.}.}

\csromanheader{2}{นามญฺจ รูปญฺจ. (1)}
\csromanheader{2}{อวิชฺชา จ ภวตณฺหา จ. (2)}
\csromanheader{2}{ภวทิฏฺฐิ จ วิภวทิฏฺฐิ จ. (3)}
\csromanheader{2}{อหิริกญฺจ\footnote{อหิรีกญฺจ (กตฺถจิ)} อโนตฺตปฺปญฺจ. (4)}
\csromanheader{2}{หิรี จ โอตฺตปฺปญฺจ. (5)}
\csromanheader{2}{โทวจสฺสตา จ ปาปมิตฺตตา จ. (6)}
\csromanheader{2}{โสวจสฺสตา จ กลฺยาณมิตฺตตา จ. (7)}
\csromanheader{2}{อาปตฺติกุสลตา จ อาปตฺติวุฏฺฐานกุสลตา จ. (8)}
\csromanheader{2}{สมาปตฺติกุสลตา จ สมาปตฺติวุฏฺฐานกุสลตา จ. (9)}
\csromanheader{2}{10. สงฺคีติสุตฺต}
\csromanheader{2}{ธาตุกุสลตา จ มนสิการกุสลตา จ. (10)}
\csromanheader{2}{อายตนกุสลตา จ ปฏิจฺจสมุปฺปาทกุสลตา จ. (11)}
\csromanheader{2}{ฐานกุสลตา จ อฏฺฐานกุสลตา จ. (12)}
\csromanheader{2}{อชฺชวญฺจ ลชฺชวญฺจ. (13)}
\csromanheader{2}{ขนฺติ จ โสรจฺจญฺจ. (14)}
\csromanheader{2}{สาขลฺยญฺจ ปฏิสนฺถาโร จ. (15)}
\csromanheader{2}{อวิหิํสา จ โสเจยฺยญฺจ. (16)}
\csromanheader{2}{มุฏฺฐสฺสจฺจญฺจ อสมฺปชญฺญญฺจ. (17)}
\csromanheader{2}{สติ จ สมฺปชญฺญญฺจ. (18)}
\csromanheader{2}{อินฺทฺริเยสุ อคุตฺตทฺวารตา จ โภชเน อมตฺตญฺญุตา จ. (19)}
\csromanheader{2}{อินฺทฺริเยสุ คุตฺตทฺวารตา จ โภชเน มตฺตญฺญุตา จ. (20)}
\csromanheader{2}{ปฏิสงฺขานพลญฺจ\footnote{ปฏิสนฺธานพลญฺจ (สฺยา)} ภาวนาพลญฺจ. (21)}
\csromanheader{2}{สติพลญฺจ สมาธิพลญฺจ. (22)}
\csromanheader{2}{สมโถ จ วิปสฺสนา จ. (23)}
\csromanheader{2}{สมถนิมิตฺตญฺจ ปคฺคหนิมิตฺตญฺจ. (24)}
\csromanheader{2}{ปคฺคโห จ อวิกฺเขโป จ. (25)}
\csromanheader{2}{สีลวิปตฺติ จ ทิฏฺฐิวิปตฺติ จ. (26)}
\csromanheader{2}{สีลสมฺปทา จ ทิฏฺฐิสมฺปทา จ. (27)}
\csromanheader{2}{สีลวิสุทฺธิ จ ทิฏฺฐิวิสุทฺธิ จ. (28)}
\csromanheader{2}{ทิฏฺฐิวิสุทฺธิ โข ปน ยถา ทิฏฺฐิสฺส จ ปธานํ. (29)}
\prose{\csromanfolio{179}สํเวโค จ สํเวชนีเยสุ ฐาเนสุ สํวิคฺคสฺส จ โยนิโส ปธานํ. (30)}

\prose{อสนฺตุฏฺฐิตา จ กุสเลสุ ธมฺเมสุ อปฺปฏิวานิตา จ ปธานสฺมิํ. (31)}

\csromanheader{2}{วิชฺชา จ วิมุตฺติ จ. (32)}
\csromanheader{2}{ขเยญาณํ อนุปฺปาเทญาณํ. (33)}
\prose{\csromanfolio{180}อิเม โข อาวุโส เตน ภควตา ชานตา ปสฺสตา อรหตา สมฺมาสมฺพุทฺเธน เทฺว ธมฺมา สมฺมทกฺขาตา.\csromanspacer{} ตตฺถ สพฺเพเหว สงฺคายิตพฺพํ, น วิวทิตพฺพํ, ยถยิทํ พฺรหฺมจริยํ อทฺธนิยํ อสฺส จิรฏฺฐิติกํ, ตทสฺส พหุชนหิตาย พหุชนสุขาย โลกานุกมฺปาย อตฺถาย หิตาย สุขาย เทวมนุสฺสานํ.}

\csromanmatikaanchor{mtk.130}
\csromantocmarkref{section}{ติก}{mtk.130}
\bookmark[dest={mtk.130},level=1]{ติก}
\csromanheader{1}{\textbf{ติก}}
\proseitem{305}{อตฺถิ โข อาวุโส เตน ภควตา ชานตา ปสฺสตา อรหตา สมฺมาสมฺพุทฺเธน ตโย ธมฺมา สมฺมทกฺขาตา.\csromanspacer{} ตตฺถ สพฺเพเหว สงฺคายิตพฺพํ ฯเปฯ อตฺถาย หิตาย สุขาย เทวมนุสฺสานํ.\csromanspacer{} กตเม ตโย.}

\prose{ตีณิ อกุสลมูลานิ, โลโภ อกุสลมูลํ โทโส อกุสลมูลํ โมโห อกุสลมูลํ. (1)}

\prose{ตีณิ กุสลมูลานิ, อโลโภ กุสลมูลํ อโทโส กุสลมูลํ อโมโห กุสลมูลํ. (2)}

\prose{ตีณิ ทุจฺจริตานิ, กายทุจฺจริตํ วจีทุจฺจริตํ มโนทุจฺจริตํ. (3)}

\prose{ตีณิ สุจริตานิ, กายสุจริตํ วจีสุจริตํ มโนสุจริตํ. (4)}

\prose{ตโย อกุสลวิตกฺกา, กามวิตกฺโก พฺยาปาทวิตกฺโก วิหิํสาวิตกฺโก. (5)}

\prose{ตโย กุสลวิตกฺกา, เนกฺขมฺมวิตกฺโก อพฺยาปาทวิตกฺโก อวิหิํสาวิตกฺโก. (6)}

\prose{ตโย อกุสลสงฺกปฺปา, กามสงฺกปฺโป พฺยาปาทสงฺกปฺโป วิหิํสาสงฺกปฺโป. (7)}

\prose{ตโย กุสลสงฺกปฺปา, เนกฺขมฺมสงฺกปฺโป อพฺยาปาทสงฺกปฺโป อวิหิํสาสงฺกปฺโป. (8)}

\csromanheader{2}{10. สงฺคีติสุตฺต}
\prose{\csromanfolio{181}ติสฺโส อกุสลสญฺญา, กามสญฺญา พฺยาปทสญฺญา วิหิํสาสญฺญา. (9)}

\prose{ติสฺโส กุสลสญฺญา, เนกฺขมฺมสญฺญา อพฺยาปาทสญฺญา อวิหิํสาสญฺญา. (10)}

\prose{ติสฺโส อกุสลธาตุโย, กามธาตุ พฺยาปาทธาตุ วิหิํสาธาตุ. (11)}

\prose{ติสฺโส กุสลธาตุโย, เนกฺขมฺมธาตุ อพฺยาปาทธาตุ อวิหิํสาธาตุ. (12)}

\prose{อปราปิ ติสฺโส ธาตุโย, กามธาตุ รูปธาตุ อรูปธาตุ. (13)}

\prose{อปราปิ ติสฺโส ธาตุโย, รูปธาตุ อรูปธาตุ นิโรธธาตุ. (14)}

\prose{อปราปิ ติสฺโส ธาตุโย, หีนธาตุ มชฺฌิมธาตุ ปณีตธาตุ. (15)}

\prose{ติสฺโส ตณฺหา, กามตณฺหา ภวตณฺหา วิภวตณฺหา. (16)}

\prose{อปราปิ ติสฺโส ตณฺหา, กามตณฺหา รูปตณฺหา อรูปตณฺหา. (17)}

\prose{อปราปิ ติสฺโส ตณฺหา, รูปตณฺหา อรูปตณฺหา นิโรธตณฺหา. (18)}

\prose{ตีณิ สํโยชนานิ, สกฺกายทิฏฺฐิ วิจิกิจฺฉา สีลพฺพตปรามาโส. (19)}

\prose{ตโย อาสวา, กามาสโว ภวาสโว อวิชฺชาสโว. (20)}

\prose{ตโย ภวา, กามภโว รูปภโว อรูปภโว. (21)}

\prose{ติสฺโส เอสนา, กาเมสนา ภเวสนา พฺรหฺมจริเยสนา. (22)}

\prose{ติสฺโส วิธา, เสยฺโยหมสฺมีติ วิธา สทิโสหมสฺมีติ วิธา หีโนหมสฺมีติ วิธา. (23)}

\prose{ตโย อทฺธา, อตีโต อทฺธา อนาคโต อทฺธา ปจฺจุปฺปนฺโน อทฺธา. (24)}

\prose{ตโย อนฺตา, สกฺกาโย อนฺโต สกฺกายสมุทโย อนฺโต สกฺกายนิโรโธ อนฺโต. (25)}

\prose{ติสฺโส เวทนา, สุขา เวทนา ทุกฺขา เวทนา อทุกฺขมสุขา เวทนา. (26)}

\prose{ติสฺโส ทุกฺขตา, ทุกฺขทุกฺขตา สงฺขารทุกฺขตา วิปริณามทุกฺขตา. (27)}

\prose{\csromanfolio{182}ตโย ราสี, มิจฺฉตฺตนิยโต ราสิ สมฺมตฺตนิยโต ราสิ อนิยโต ราสิ. (28)}

\prose{ตโย ตมา\footnote{ติสฺโส กงฺขา (พหูสุ) อฏฺฐกถา โอโลเกตพฺพา.}, อตีตํ วา อทฺธานํ อารพฺภ กงฺขติ วิจิกิจฺฉติ นาธิมุจฺจติ น สมฺปสีทติ, อนาคตํ วา อทฺธานํ อารพฺภ กงฺขติ วิจิกิจฺฉติ นาธิมุจฺจติ น สมฺปสีทติ, เอตรหิ วา ปจฺจุปฺปนฺนํ อทฺธานํ อารพฺภ กงฺขติ วิจิกิจฺฉติ นาธิมุจฺจติ น สมฺปสีทติ. (29)}

\prose{ตีณิ ตถาคตสฺส อรกฺเขยฺยานิ, ปริสุทฺธกายสมาจาโร อาวุโส ตถาคโต, นตฺถิ ตถาคตสฺส กายทุจฺจริตํ, ยํ ตถาคโต รกฺเขยฺย “มา เม อิทํ ปโร อญฺญาสี”ติ.\csromanspacer{} ปริสุทฺธวจีสมาจาโร อาวุโส ตถาคโต, นตฺถิ ตถาคตสฺส วจีทุจฺจริตํ, ยํ ตถาคโต รกฺเขยฺย “มา เม อิทํ ปโร อญฺญาสี”ติ.\csromanspacer{} ปริสุทฺธมโนสมาจาโร อาวุโส ตถาคโต, นตฺถิ ตถาคตสฺส มโนทุจฺจริตํ, ยํ ตถาคโต รกฺเขยฺย “มา เม อิทํ ปโร อญฺญาสี”ติ. (30)}

\prose{ตโย กิญฺจนา, ราโค กิญฺจนํ โทโส กิญฺจนํ โมโห กิญฺจนํ. (31)}

\prose{ตโย อคฺคี, ราคคฺคิ โทสคฺคิ โมหคฺคิ. (32)}

\prose{อปเรปิ ตโย อคฺคี, อาหุเนยฺยคฺคิ คหปตคฺคิ ทกฺขิเณยฺยคฺคิ. (33)}

\prose{ติวิเธน รูปสงฺคโห, สนิทสฺสนสปฺปฏิฆํ รูปํ\footnote{สนิทสฺสนสปฺปฏิฆรูปํ (สฺยา, กํ) เอวมิตรทฺวเยปิ,} อนิทสฺสนสปฺปฏิฆํ รูปํ อนิทสฺสน-อปฺปฏิฆํ รูปํ. (34)}

\prose{ตโย สงฺขารา, ปุญฺญาภิสงฺขาโร อปุญฺญภิสงฺขาโร อาเนญฺชาภิสงฺขาโร. (35)}

\prose{ตโย ปุคฺคลา, เสกฺโข ปุคฺคโล อเสกฺโข ปุคฺคโล เนวเสกฺโขนาเสกฺโข ปุคฺคโล. (36)}

\prose{ตโย เถรา, ชาติเถโรธมฺมเถโร สมฺมุติเถโร\footnote{สมฺมติเถโร (สฺยา, กํ)}. (37)}

\prose{ตีณิ ปุญฺญกิริยวตฺถูนิ, ทานมยํ ปุญฺญกิริยวตฺถุ สีลมยํ ปุญฺญกิริยวตฺถุ ภาวนามยํ ปุญฺญกิริยวตฺถุ. (38)}

\prose{ตีณิ โจทนาวตฺถูนิ, ทิฏฺเฐน สุเตน ปริสงฺกาย. (39)}

\csromanheader{2}{10. สงฺคีติสุตฺต}
\prose{\csromanfolio{183}ติสฺโส กามูปปตฺติโย\footnote{กามุปฺปตฺติโย (สี), กามุปปตฺติโย (สฺยา, อิ, ก)}.\csromanspacer{} สนฺตาวุโส สตฺตา ปจฺจุปฏฺฐิตกามา, เต ปจฺจุปฏฺฐิเตสุ กาเมสุ วสํ วตฺเตนฺติ, เสยฺยถาปิ มนุสฺสา เอกจฺเจ จ เทวา เอกจฺเจ จ วินิปาติกา.\csromanspacer{} อยํ ปฐมา กามูปปปตฺติ.\csromanspacer{} สนฺตาวุโส สตฺตา นิมฺมิตกามา, เต นิมฺมินิตฺวา นิมฺมินิตฺวา กาเมสุ วสํ วตฺเตนฺติ, เสยฺยถาปิ เทวา นิมฺมานรตี.\csromanspacer{} อยํ ทุติยา กามูปปตฺติ.\csromanspacer{} สนฺตาวุโส สตฺตา ปรนิมฺมิตกามา, เต ปรนิมฺมิเตสุ กาเมสุ วสํ วตฺเตนฺติ, เสยฺยถาปิ เทวา ปรนิมฺมิตวสวตฺตี.\csromanspacer{} อยํ ตติยา กามูปปตฺติ. (40)}

\prose{ติสฺโส สุขูปปตฺติโย\footnote{สุขุปปตฺติโย (สฺยา, อิ, ก)}, สนฺตาวุโส สตฺตา\footnote{สตฺตา สุขํ (สฺยา, กํ)} อุปฺปาเทตฺวา อุปฺปาเทตฺวา สุขํ วิหรนฺติ, เสยฺยถาปิ เทวา พฺรหฺมกายิกา.\csromanspacer{} อยํ ปฐมา สุขูปปตฺติ.\csromanspacer{} สนฺตาวุโส สตฺตา สุเขน อภิสนฺนา ปริสนฺนา ปริปูรา ปริปฺผุฏา.\csromanspacer{} เต กทาจิ กรหจิ อุทานํ อุทาเนนฺติ “อโห สุขํ อโห สุขนฺ”ติ, เสยฺยถาปิ เทวา อาภสฺสรา.\csromanspacer{} อยํ ทุติยา สุขูปปตฺติ.\csromanspacer{} สนฺตาวุโส สตฺตา สุเขน อภิสนฺนา ปริสนฺนา ปริปูรา ปริปฺผุฏา.\csromanspacer{} เต สนฺตํเยว ตุสิตา\footnote{สนฺตุสิตา (สฺยา, กํ)} สุขํ\footnote{จิตฺตสุขํ (สฺยา, ก)} ปฏิสํเวเทนฺติ, เสยฺยถาปิ เทวา สุภกิณฺหา, อยํ ตติยา สุขูปปตฺติ. (41)}

\prose{ติสฺโส ปญฺญา, เสกฺขา ปญฺญา อเสกฺขา ปญฺญา เนวเสกฺขานาเสกฺขา ปญฺญา. (42)}

\prose{อปราปิ ติสฺโส ปญฺญา, จินฺตามยา ปญฺญา สุตมยา ปญฺญา ภาวนามยา ปญฺญา. (43)}

\prose{ตีณาวุธานิ, สุตาวุธํ ปวิเวกาวุธํ ปญฺญาวุธํ. (44)}

\prose{ตีณินฺทฺริยานิ, อนญฺญาตญฺญสฺสามีตินฺทฺริยํ อญฺญินฺทฺริยํ อญฺญาตาวินฺทฺริยํ. (45)}

\prose{ตีณิ จกฺขูนิ, มํสจกฺขุ ทิพฺพจกฺขุ ปญฺญาจกฺขุ. (46)}

\prose{ติสฺโส สิกฺขา, อธิสีลสิกฺขา อธิจิตฺตสิกฺขา อธิปญฺญาสิกฺขา. (47)}

\csromanheader{2}{ติสฺโส ภาวนา กายภาวนา จิตฺตภาวนา ปญฺญาภาวนา. (48)}
\prose{\csromanfolio{184}ตีณิ อนุตฺตริยานิ, ทสฺสนานุตฺตริยํ ปฏิปทานุตฺตริยํ วิมุตฺตานุตฺตริยํ. (49)}

\prose{ตโย สมาธี, สวิตกฺกสวิจาโร สมาธิ อวิตกฺกวิจารมตฺโต สมาธิ อวิตกฺก-อวิจาโร สมาธิ. (50)}

\prose{อปเรปิ ตโย สมาธี, สุญฺญโต สมาธิ อนิมิตฺโต สมาธิ อปฺปณิหิโต สมาธิ. (51)}

\prose{ตีณิ โสเจยฺยานิ, กายโสเจยฺยํ วจีโสเจยฺยํ มโนโสเจยฺยํ. (52)}

\prose{ตีณิ โมเนยฺยานิ, กายโมเนยฺยํ วจีโมเนยฺยํ มโนโมเนยฺยํ. (53)}

\prose{ตีณิ โกสลฺลานิ, อายโกสลฺลํ อปายโกสลฺลํ อุปายโกสลฺลํ. (54)}

\prose{ตโย มทา, อาโรคฺยมโท โยพฺพนมโท ชีวิตมโท. (55)}

\prose{ตีณิ อาธิปเตยฺยานิ, อตฺตาธิปเตยฺยํ โลกาธิปเตยฺยํ ธมฺมาธิปเตยฺยํ. (56)}

\prose{ตีณิ กถาวตฺถูนิ, อตีตํ วา อทฺธานํ อารพฺภ กถํ กเถยฺย “เอวํ อโหสิ อตีตมทฺธานนฺ”ติ, อนาคตํ วา อทฺธานํ อารพฺภ กถํ กเถยฺย “เอวํ ภวิสฺสติ อนาคตมทฺธานนฺ”ติ, เอตรหิ วา ปจฺจุปฺปนฺนํ อทฺธานํ อารพฺภ กถํ กเถยฺย “เอวํ โหติ เอตรหิ ปจฺจุปฺปนฺนํ อทฺธานนฺ”ติ. (57)}

\prose{ติสฺโส วิชฺชา, ปุพฺเพนิวาสานุสฺสติญาณํ วิชฺชา สตฺตานํ จุตูปปาเตญาณํ วิชฺชา อาสวานํ ขเยญาณํ วิชฺชา. (58)}

\prose{ตโย วิหารา, ทิพฺโพ วิหาโร พฺรหฺมา วิหาโร อริโย วิหาโร. (59)}

\prose{ตีณิ ปาฏิหาริยานิ, อิทฺธิปาฏิหาริยํ อาเทสนาปาฏิหาริยํ อนุสาสนีปาฏิหาริยํ. (60)}

\prose{อิเม โข อาวุโส เตน ภควตา ชานตา ปสฺสตา อรหตา สมฺมาสมฺพุทฺเธน ตโย ธมฺมา สมฺมทกฺขาตา.\csromanspacer{} ตตฺถ สพฺเพเหว สงฺคายิตพฺพํ ฯเปฯ อตฺถาย หิตาย สุขาย เทวมนุสฺสานํ.}

\csromanmatikaanchor{mtk.131}
\csromantocmarkref{section}{จตุกฺก}{mtk.131}
\bookmark[dest={mtk.131},level=1]{จตุกฺก}
\csromanheader{1}{10. สงฺคีติสุตฺต \textbf{จตุกฺก}}
\proseitem{306}{\csromanfolio{185}อตฺถิ โข อาวุโส เตน ภควตา ชานตา ปสฺสตา อรหตา สมฺมาสมฺพุทฺเธน จตฺตาโร ธมฺมา สมฺมทกฺขาตา.\csromanspacer{} ตตฺถ สพฺเพเหว สงฺคายิตพฺพํ, น วิวทิตพฺพํ ฯเปฯ อตฺถาย หิตาย สุขาย เทวมนุสฺสานํ.\csromanspacer{} กตเม จตฺตาโร.}

\prose{จตฺตาโร \textbf{สติปฏฺฐานา}.\csromanspacer{} อิธาวุโส ภิกฺขุ กาเย กายานุปสฺสี วิหรติ อาตาปี สมฺปชาโน สติมา วิเนยฺย โลเก อภิชฺฌาโทมนสฺสํ.\csromanspacer{} เวทนาสุ เวทนานุปสฺสี ฯเปฯ.\csromanspacer{} จิตฺเต จิตฺตานุปสฺสี ฯเปฯ.\csromanspacer{} ธมฺเมสุ ธมฺมานุปสฺสี วิหรติ อาตาปี สมฺปชาโน สติมา วิเนยฺย โลเก อภิชฺฌาโทมนสฺสํ. (1)}

\prose{จตฺตาโร \textbf{สมฺมปฺปธานา}.\csromanspacer{} อิธาวุโส ภิกฺขุ อนุปฺปนฺนานํ ปาปกานํ อกุสลานํ ธมฺมานํ อนุปฺปาทาย ฉนฺทํ ชเนติ วายมติ วีริยํ อารภติ จิตฺตํ ปคฺคณฺหาติ ปทหติ, อุปฺปนฺนานํ ปาปกานํ อกุสลานํ ธมฺมานํ ปหานาย ฉนฺทํ ชเนติ วายมติ วีริยํ อารภติ จิตฺตํ ปคฺคณฺหาติ ปทหติ, อนุปฺปนฺนานํ กุสลานํ ธมฺมานํ อุปฺปาทาย ฉนฺทํ ชเนติ วายมติ วีริยํ อารภติ จิตฺตํ ปคฺคณฺหาติ ปทหติ, อุปฺปนฺนานํ กุสลายํ ธมฺมานํ ฐิติยา อสมฺโมสาย ภิยฺโยภาวาย เวปุลฺลาย ภาวนาย ปาริปูริยา ฉนฺทํ ชเนติ วายมติ วีริยํ อารภติ จิตฺตํ ปคฺคณฺหาติ ปทหติ. (2)}

\prose{จตฺตาโร \textbf{อิทฺธิปาทา}.\csromanspacer{} อิธาวุโส ภิกฺขุ ฉนฺทสมาธิปธานสงฺขารสมนฺนาคตํ อิทฺธิปาทํ ภาเวติ, จิตฺตสมาธิปธานสงฺขารสมนฺนาคตํ อิทฺธิปาทํ ภาเวติ, วีริยสมาธิปธานสงฺขารสมนฺนาคตํ อิทฺธิปาทํ ภาเวติ, วีมํสาสมาธิปธานสงฺขารสมนฺนาคตํ อิทฺธิปาทํ ภาเวติ. (3)}

\prose{จตฺตาริ \textbf{ฌานานิ}.\csromanspacer{} อิธาวุโส ภิกฺขุ วิวิจฺเจว กาเมหิ วิวิจฺจ อกุสเลหิ ธมฺเมหิ สวิตกฺกํ สวิจารํ วิเวกชํ ปีติสุขํ ปฐมํ ฌานํ\footnote{ปฐมชฺฌานํ (สฺยา, กํ)} อุปสมฺปชฺช วิหรติ.\csromanspacer{} วิตกฺกวิจารานํ วูปสมา อชฺฌตฺตํ สมฺปสาทนํ เจตโส เอโกทิภาวํ อวิตกฺกํ อวิจารํ สมาธิชํ ปีติสุขํ ทุติยํ ฌานํ\footnote{ทุติยชฺฌานํ (สฺยา, กํ)} อุปสมฺปชฺช วิหรติ.\csromanspacer{} ปีติยา จ วิราคา อุเปกฺขโก จ วิหรติ สโต จ สมฺปชาโน, \csromanfolio{186}สุขญฺจ กาเยน ปฏิสํเวเทติ, ยํ ตํ อริยา อาจิกฺขนฺติ “อุเปกฺขโก สติมา สุขวิหารี”ติ, ตติยํ ฌานํ\footnote{ตติยชฺฌานํ (สฺยา, กํ)} อุปสมฺปชฺช วิหรติ.\csromanspacer{} สุขสฺส จ ปหานา ทุกฺขสฺส จ ปหานา ปุพฺเพว โสมนสฺสโทมนสฺสานํ อตฺถงฺคมา อทุกฺขมสุขํ อุเปกฺขาสติปาริสุทฺธิํ จตุตฺถํ ฌานํ\footnote{จตุตฺถชฺฌานํ (สฺยา, กํ)} อุปสมฺปชฺช วิหรติ. (4)}

\proseitem{307}{จตสฺโส \textbf{สมาธิภาวนา}.\csromanspacer{} อตฺถาวุโส \textbf{สมาธิภาวนา} ภาวิตา พหุลีกตา ทิฏฺฐธมฺมสุขวิหาราย สํวตฺตติ, อตฺถาวุโส \textbf{สมาธิภาวนา} ภาวิตา พหุลีกตา ญาณทสฺสนปฏิลาภาย สํวตฺตติ, อตฺถาวุโส \textbf{สมาธิภาวนา} ภาวิตา พหุลีกตา สติสมฺปชญฺญาย สํวตฺตติ, อตฺถาวุโส \textbf{สมาธิภาวนา} ภาวิตา พหุลีกตา อาสวานํ ขยาย สํวตฺตติ.}

\prose{กตมา จาวุโส \textbf{สมาธิภาวนา} ภาวิตา พหุลีกตา ทิฏฺฐธมฺมสุขวิหาราย สํวตฺตติ.\csromanspacer{} อิธาวุโส ภิกฺขุ วิวิจฺเจว กาเมหิ วิวิจฺจ อกุสเลหิ ธมฺเมหิ สวิตกฺกํ ฯเปฯ จตุตฺถํ ฌานํ อุปสมฺปชฺช วิหรติ.\csromanspacer{} อยํ อาวุโส \textbf{สมาธิภาวนา} ภาวิตา พหุลีกตา ทิฏฺฐธมฺมสุขวิหาราย สํวตฺตติ. (5-1)}

\prose{กตมา จาวุโส \textbf{สมาธิภาวนา} ภาวิตา พหุลีกตา ญาณทสฺสนปฏิลาภาย สํวตฺตติ.\csromanspacer{} อิธาวุโส ภิกฺขุ อาโลกสญฺญํ มนสิ กโรติ, ทิวาสญฺญํ อธิฏฺฐาติ ยถา ทิวา ตถา รตฺติํ, ยถา รตฺติํ ตถา ทิวา.\csromanspacer{} อิติ วิวเฏน เจตสา อปริโยนทฺเธน สปฺปภาสํ จิตฺตํ ภาเวติ.\csromanspacer{} อยํ อาวุโส \textbf{สมาธิภาวนา} ภาวิตา พหุลีกตา ญาณทสฺสนปฏิลาภาย สํวตฺตติ. (5-2)}

\prose{กตมา จาวุโส \textbf{สมาธิภาวนา} ภาวิตา พหุลีกตา สติสมฺปชญฺญาย สํวตฺตติ.\csromanspacer{} อิธาวุโส ภิกฺขุโน วิทิตา เวทนา อุปฺปชฺชนฺติ, วิทิตา อุปฏฺฐหนฺติ, วิทิตา อพฺภตฺถํ คจฺฉนฺติ.\csromanspacer{} วิทิตา สญฺญา อุปฺปชฺชนฺติ, วิทิตา อุปฏฺฐหนฺติ, วิทิตา อพฺภตฺถํ คจฺฉนฺติ.\csromanspacer{} วิทิตา วิตกฺกา อุปฺปชฺชนฺติ, วิทิตา อุปฏฺฐหนฺติ, วิทิตา อพฺภตฺถํ คจฺฉนฺติ.\csromanspacer{} อยํ อาวุโส \textbf{สมาธิภาวนา} ภาวิตา พหุลีกตา สติสมฺปชญฺญาย สํวตฺตติ. (5-3)}

\csromanheader{2}{10. สงฺคีติสุตฺต}
\prose{\csromanfolio{187}กตมา จาวุโส สมาธิภาวนา ภาวิตา พหุลีกตา อาสวานํ ขยาย สํวตฺตติ.\csromanspacer{} อิธาวุโส ภิกฺขุ ปญฺจสุ อุปาทานกฺขนฺเธสุ อุทยพฺพยานุปสฺสี วิหรติ.\csromanspacer{} อิติ รูปํ, อิติ รูปสฺส สมุทโย, อิติ รูปสฺส อตฺถงฺคโม.\csromanspacer{} อิติ เวทนา.\csromanspacer{} อิติ สญฺญา.\csromanspacer{} อิติ สงฺขารา.\csromanspacer{} อิติ วิญฺญาณํ, อิติ วิญฺญาณสฺส สมุทโย, อิติ วิญฺญาณสฺส อตฺถงฺคโม.\csromanspacer{} อยํ อาวุโส สมาธิภาวนา ภาวิตา พหุลีกตา อาสวานํ ขยาย สํวตฺตติ. (5-4)}

\proseitem{308}{จตสฺโส \textbf{อปฺปมญฺญา}.\csromanspacer{} อิธาวุโส ภิกฺขุ เมตฺตาสหคเตน เจตสา เอกํ ทิสํ ผริตฺวา วิหรติ.\csromanspacer{} ตถา ทุติยํ.\csromanspacer{} ตถา ตติยํ.\csromanspacer{} ตถา จตุตฺถํ.\csromanspacer{} อิติ อุทฺธมโธ ติริยํ สพฺพธิ สพฺพตฺตตาย สพฺพาวนฺตํ โลกํ เมตฺตาสหคเตน เจตสา วิปุเลน มหคฺคเตน อปฺปมาเณน อเวเรน อพฺยาปชฺเชน\footnote{อพฺยาปชฺเฌน (สี, สฺยา, กํ, อิ)} ผริตฺวา วิหรติ.\csromanspacer{} กรุณาสหคเตน เจตสา ฯเปฯ.\csromanspacer{} มุทิตาสหคเตน เจตสา ฯเปฯ.\csromanspacer{} อุเปกฺขาสหคเตน เจตสา เอกํ ทิสํ ผริตฺวา วิหรติ.\csromanspacer{} ตถา ทุติยํ.\csromanspacer{} ตถา ตติยํ.\csromanspacer{} ตถา จตุตฺถํ.\csromanspacer{} อิติ อุทฺธมโธ ติริยํ สพฺพธิ สพฺพตฺตตาย สพฺพาวนฺตํ โลกํ อุเปกฺขาสหคเตน เจตสา วิปุเลน มหคฺคเตน อปฺปมาเณน อเวเรน อพฺยาปชฺเชน ผริตฺวา วิหรติ. (6)}

\prose{จตฺตาโร \textbf{อารุปฺปา}\footnote{อรูปา (สฺยา, กํ, อิ)}.\csromanspacer{} อิธาวุโส ภิกฺขุ สพฺพโส รูปสญฺญานํ สมติกฺกมา ปฏิฆสญฺญานํ อตฺถงฺคมา นานตฺตสญฺญานํ อมนสิการา “อนนฺโต อากาโส”ติ อากาสานญฺจายตนํ อุปสมฺปชฺช วิหรติ, สพฺพโส อากาสานญฺจายตนํ สมติกฺกมฺม “อนนฺตํ วิญฺญาณนฺ”ติ วิญฺญาณญฺจายตนํ อุปสมฺปชฺช วิหรติ, สพฺพโส วิญฺญาณญฺจายตนํ สมติกฺกมฺม “นตฺถิ กิญฺจี”ติ อากิญฺจญฺญายตนํ อุปสมฺปชฺช วิหรติ, สพฺพโส อากิญฺจญฺญายตนํ สมติกฺกมฺม เนวสญฺญานาสญฺญายตนํ อุปสมฺปชฺช วิหรติ. (7)}

\prose{จตฺตาริ \textbf{อปสฺเสนานิ}.\csromanspacer{} อิธาวุโส ภิกฺขุ สงฺขาเยกํ ปฏิเสวติ, สงฺขาเยกํ อธิวาเสติ, สงฺขาเยกํ ปริวชฺเชติ, สงฺขาเยกํ วิโนเทติ. (8)}

\proseitem{309}{\csromanfolio{188}จตฺตาโร \textbf{อริยวํสา}.\csromanspacer{} อิธาวุโส ภิกฺขุ สนฺตุฏฺโฐ โหติ อิตรีตเรน จีวเรน, อิตรีตรจีวรสนฺตุฏฺฐิยา จ วณฺณวาที, น จ จีวรเหตุ อเนสนํ อปฺปติรูปํ อาปชฺชติ, อลทฺธา จ จีวรํ น ปริตสฺสติ, ลทฺธา จ จีวรํ อคธิโต\footnote{อคถิโต (สี, อิ)} อมุจฺฉิโต อนชฺฌาปนฺโน อาทีนวทสฺสาวี นิสฺสรณปญฺโญ ปริภุญฺชติ, ตาย จ ปน อิตรีตรจีวรสนฺตุฏฺฐิยา เนวตฺตานุกฺกํเสติ น ปรํ วมฺเภติ.\csromanspacer{} โย หิ ตตฺถ ทกฺโข โหติ อนลโส สมฺปชาโน ปฏิสฺสโต, อยํ วุจฺจตาวุโส ภิกฺขุ โปราเณ อคฺคญฺเญ อริยวํเส ฐิโต. (9-1)}

\prose{ปุน จปรํ อาวุโส ภิกฺขุ สนฺตุฏฺโฐ โหติ อิตรีตเรน ปิณฺฑปาเตน, อิตรีตรปิณฺฑปาตสนฺตุฏฺฐิยา จ วณฺณวาที, น จ ปิณฺฑปาตเหตุ อเนสนํ อปฺปติรูปํ อาปชฺชติ, อลทฺธา จ ปิณฺฑปาตํ น ปริตสฺสติ, ลทฺธา จ ปิณฺฑปาตํ อคธิโต อมุจฺฉิโต อนชฺฌาปนฺโน อาทีนวทสฺสาวี นิสฺสรณปญฺโญ ปริภุญฺชติ, ตาย จ ปน อิตรีตรปิณฺฑปาตสนฺตุฏฺฐิยา เนวตฺตานุกฺกํเสติ น ปรํ วมฺเภติ.\csromanspacer{} โย หิ ตตฺถ ทกฺโข อนลโส สมฺปชาโน ปฏิสฺสโต, อยํ วุจฺจตาวุโส ภิกฺขุ โปราเณ อคฺคญฺเญ อริยวํเส ฐิโต. (9-2)}

\prose{ปุน จปรํ อาวุโส ภิกฺขุ สนฺตุฏฺโฐ โหติ อิตรีตเรน เสนาสเนน, อิตรีตรเสนาสนสนฺตุฏฺฐิยา จ วณฺณวาที, น จ เสนาสนเหตุ อเนสนํ อปฺปติรูปํ อาปชฺชติ, อลทฺธา จ เสนาสนํ น ปริตสฺสติ, ลทฺธา จ เสนาสนํ อคธิโต อมุจฺฉิโต อนชฺฌาปนฺโน อาทีนวทสฺสาวี นิสฺสรณปญฺโญ ปริภุญฺชติ, ตาย จ ปน อิตรีตรเสนาสนสนฺตุฏฺฐิยา เนวตฺตานุกฺกํเสติ น ปรํ วมฺเภติ.\csromanspacer{} โย หิ ตตฺถ ทกฺโข อนลโส สมฺปชาโน ปฏิสฺสโต, อยํ วุจฺจตาวุโส ภิกฺขุ โปราเณ อคฺคญฺเญ อริยวํเส ฐิโต. (9-3)}

\prose{ปุน จปรํ อาวุโส ภิกฺขุ ปหานาราโม โหติ ปหานรโต, ภาวนาราโม โหติ ภาวนารโต, ตาย จ ปน ปหานารามตาย ปหานรติยา ภาวนารามตาย ภาวนารติยา เนวตฺตานุกฺกํเสติ น ปรํ วมฺเภติ.\csromanspacer{} โย หิ ตตฺถ ทกฺโข อนลโส สมฺปชาโน ปฏิสฺสโต, อยํ วุจฺจตาวุโส ภิกฺขุ โปราเณ อคฺคญฺเญ อริยวํเส ฐิโตติ. (9-4)}

\csromanheader{2}{10. สงฺคีติสุตฺต}
\proseitem{310}{\csromanfolio{189}จตฺตาริ \textbf{ปธานานิ}.\csromanspacer{} สํวรปธานํ ปหานปธานํ ภาวนาปธานํ\footnote{ภาวนาปฺปธานํ (สฺยา)} อนุรกฺขณาปธานํ\footnote{อนุรกฺขนาปฺปธานํ (สฺยา)}.\csromanspacer{} กตมญฺจาวุโส สํวรปธานํ.\csromanspacer{} อิธาวุโส ภิกฺขุ จกฺขุนา รูปํ ทิสฺวา น นิมิตฺตคฺคาหี โหติ นานุพฺยญฺชนคฺคาหี, ยตฺวาธิกรณเมนํ จกฺขุนฺทฺริยํ อสํวุตํ วิหรนฺตํ อภิชฺฌาโทมนสฺสา ปาปกา อกุสลา ธมฺมา อนฺวาสเวยฺยุํ, ตสฺส สํวราย ปฏิปชฺชติ รกฺขติ จกฺขุนฺทฺริยํ, จกฺขุนฺทฺริเย สํวรํ อาปชฺชติ.\csromanspacer{} โสเตน สทฺทํ สุตฺวา.\csromanspacer{} ฆาเนน คนฺธํ ฆายิตฺวา.\csromanspacer{} ชิวฺหาย รสํ สายิตฺวา.\csromanspacer{} กาเยน โผฏฺฐพฺพํ ผุสิตฺวา.\csromanspacer{} มนสา ธมฺมํ วิญฺญาย น นิมิตฺตคฺคาหี โหติ นานุพฺยญฺชนคฺคาหี, ยตฺวาธิกรณเมนํ มนินฺทฺริยํ อสํวุตํ วิหรนฺตํ อภิชฺฌาโทมนสฺสา ปาปกา อกุสลา ธมฺมา อนฺวาสฺสเวยฺยุํ, ตสฺส สํวราย ปฏิปชฺชติ มนินฺทฺริยํ, มนินฺทฺริเย สํวรํ อาปชฺชติ.\csromanspacer{} อิทํ วุจฺจตาวุโส สํวรปธานํ. (10-1)}

\prose{กตมญฺจาวุโส ปหานปธานํ.\csromanspacer{} อิธาวุโส ภิกฺขุ อุปฺปนฺนํ กามวิตกฺกํ นาธิวาเสติ ปชหติ วิโนเทติ พฺยนฺติํ กโรติ\footnote{พฺยนฺตี กโรติ (สฺยา, กํ)} อนภาวํ คเมติ.\csromanspacer{} อุปฺปนฺนํ พฺยาปาทวิตกฺกํ ฯเปฯ.\csromanspacer{} อุปฺปนฺนํ วิหิํสาวิตกฺกํ ฯเปฯ.\csromanspacer{} อุปฺปนฺนุปฺปนฺเน ปาปเก อกุสเล ธมฺเม นาธิวาเสติ ปชหติ วิโนเทติ พฺยนฺติํ กโรติ อนภาวํ คเมติ.\csromanspacer{} อิทํ วุจฺจตาวุโส ปหานปธานํ. (10-2)}

\prose{กตมญฺจาวุโส ภาวนาปธานํ.\csromanspacer{} อิธาวุโส ภิกฺขุ สติสมฺโพชฺฌงฺคํ ภาเวติ วิเวกนิสฺสิตํ วิราคนิสฺสิตํ นิโรธนิสฺสิตํ โวสฺสคฺคปริณามิํ.\csromanspacer{} ธมฺมวิจยสมฺโพชฺฌงฺคํ ภาเวติ.\csromanspacer{} วีริยสมฺโพชฺฌงฺคํ ภาเวติ.\csromanspacer{} ปีติสมฺโพชฺฌงฺคํ ภาเวติ.\csromanspacer{} ปสฺสทฺธิสมฺโพชฺฌงฺคํ ภาเวติ.\csromanspacer{} สมาธิสมฺโพชฺฌงฺคํ ภาเวติ.\csromanspacer{} อุเปกฺขาสมฺโพชฺฌงฺคํ ภาเวติ วิเวกนิสฺสิตํ วิราคนิสฺสิตํ นิโรธ- นิสฺสิตํ โวสฺสคฺคปริณามิํ.\csromanspacer{} อิทํ วุจฺจตาวุโส ภาวนาปธานํ. (10-3)}

\prose{กตมญฺจาวุโส อนุรกฺขณาปธานํ.\csromanspacer{} อิธาวุโส ภิกฺขุ อุปฺปนฺนํ ภทฺรกํ\footnote{ภทฺทกํ (สฺยา, กํ, อิ)} สมาธินิมิตฺตํ อนุรกฺขติ อฏฺฐิกสญฺญํ ปุฬุวกสญฺญํ\footnote{ปุฬวกสญฺญํ (สี, อิ)} วินีลกสญฺญํ วิจฺฉิทฺทกสญฺญํ อุทฺธุมาตกสญฺญํ.\csromanspacer{} อิทํ วุจฺจตาวุโส อนุรกฺขณาปธานํ. (10-4)}

\prose{จตฺตาริ \textbf{ญาณานิ}, ธมฺเม ญาณํ อนฺวเย ญาณํ ปริเย\footnote{ปริจฺเจ (สี, ก), ปริจฺเฉเท (สฺยา, อิ, ก) ฏีกา โอโลเกตพฺพา.} ญาณํ สมฺมุติยา ญาณํ\footnote{สมฺมติญาณํ (สฺยา, กํ)}. (11)}

\prose{\csromanfolio{190}อปรานิปิ จตฺตาริ \textbf{ญาณานิ}, ทุกฺเข ญาณํ ทุกฺขสมุทเย ญาณํ ทุกฺขนิโรเธ ญาณํ ทุกฺขนิโรธคามินิยา ปฏิปทาย ญาณํ. (12)}

\proseitem{311}{จตฺตาริ \textbf{โสตาปตฺติยงฺคานิ}, สปฺปุริสสํเสโว สทฺธมฺมสฺสวนํ โยนิโสมนสิกาโร ธมฺมานุธมฺมปฺปฏิปตฺติ. (13)}

\prose{จตฺตาริ \textbf{โสตาปนฺนสฺส องฺคานิ}.\csromanspacer{} อิธาวุโส อริยสาวโก พุทฺเธ อเวจฺจปฺปสาเทน สมนฺนาคโต โหติ “อิติปิ โส ภควา อรหํ สมฺมาสมฺพุทฺโธ วิชฺชาจรณสมฺปนฺโน สุคโต โลกวิทู อนุตฺตโร ปุริสทมฺมสารถิ สตฺถา เทวมนุสฺสานํ พุทฺโธ ภควา”ติ, ธมฺเม อเวจฺจปฺปสาเทน สมนฺนาคโต โหติ “สฺวากฺขาโต ภควตา ธมฺโม สนฺทิฏฺฐิโก อกาลิโก เอหิปสฺสิโก โอปเนยฺยิโก\footnote{โอปนยิโก (สฺยา, กํ)} ปจฺจตฺตํ เวทิตพฺโพ วิญฺญูหี”ติ, สํเฆ อเวจฺจปฺปสาเทน สมนฺนาคโต โหติ “สุปฺปฏิปนฺโน ภควโต สาวกสํโฆ อุชุปฺปฏิปนฺโน ภควโต สาวกสํโฆ ญายปฺปฏิปนฺโน ภควโต สาวกสํโฆ สามีจิปฺปฏิปนฺโน ภควโต สาวกสํโฆ ยทิทํ จตฺตาริ ปุริสยุคานิ อฏฺฐ ปุริสปุคฺคลา เอส ภควโต สาวกสํโฆ อาหุเนยฺโย ปาหุเนยฺโย ทกฺขิเณยฺโย อญฺชลิกรณีโย อนุตฺตรํ ปุญฺญกฺเขตฺตํ โลกสฺสา”ติ, อริยกนฺเตหิ สีเลหิ สมนฺนาคโต โหติ อขณฺเฑหิ อจฺฉิทฺเทหิ อสพเลหิ อกมฺมาเสหิ ภุชิสฺเสหิ วิญฺญุปฺปสตฺเถหิ อปรามฏฺเฐหิ สมาธิสํวตฺตนิเกหิ. (14)}

\prose{จตฺตาริ \textbf{สามญฺญผลานิ}, โสตาปตฺติผลํ สกทาคามิผลํ อนาคามิผลํ อรหตฺตผลํ. (15)}

\prose{จตสฺโส \textbf{ธาตุโย}, ปถวีธาตุ อาโปธาตุ เตโชธาตุ วาโยธาตุ. (16)}

\prose{จตฺตาโร \textbf{อาหารา}, กพฬีกาโร อาหาโร โอฬาริโก วา สุขุโม วา, ผสฺโส ทุติโย, มโนสญฺเจตนา ตติยา, วิญฺญาณํ จตุตฺถํ. (17)}

\prose{จตสฺโส \textbf{วิญฺญาณฏฺฐิติโย}.\csromanspacer{} รูปูปายํ วา อาวุโส วิญฺญาณํ ติฏฺฐมานํ ติฏฺฐติ รูปารมฺมณํ\footnote{รูปารมณํ (?)} รูปปฺปติฏฺฐํ นนฺทูปเสจนํ วุทฺธิํ วิรูฬฺหิํ เวปุลฺลํ อาปชฺชติ,}

\vspace{\tipitakaparskip}
\csromancenter{สงฺคีติสุตฺต เวทนูปายํ วา อาวุโส.\csromanspacer{} สญฺญูปายํ วา อาวุโส.\csromanspacer{} สงฺขารูปายํ วา อาวุโส วิญฺญาณํ ติฏฺฐมานํ ติฏฺฐติ สงฺขารารมฺมณํ สงฺขารปฺปติฏฺฐํ นนฺทูปเสจนํ วุทฺธิํ วิรูฬฺหิํ เวปุลฺลํ อาปชฺชติ. (18)}
\prose{\csromanfolio{191}จตฺตาริ \textbf{อคติคมานานิ}.\csromanspacer{} ฉนฺทาคติํ คจฺฉติ, โทสาคติํ คจฺฉติ, โมหาคติํ คจฺฉติ, ภยาคติํ คจฺฉติ. (19)}

\prose{จตฺตาโร \textbf{ตณฺหุปฺปาทา}.\csromanspacer{} จีวรเหตุ วา อาวุโส ภิกฺขุโน ตณฺหา อุปฺปชฺชมานา อุปฺปชฺชติ, ปิณฺฑปาตเหตุ วา อาวุโส ภิกฺขุโน ตณฺหา อุปฺปชฺชมานา อุปฺปชฺชติ, เสนาสนเหตุ วา อาวุโส ภิกฺขุโน ตณฺหา อุปฺปชฺชมานา อุปฺปชฺชติ, อิติภวาภวเหตุ วา อาวุโส ภิกฺขุโน ตณฺหา อุปฺปชฺชมานา อุปฺปชฺชติ. (20)}

\prose{จตสฺโส \textbf{ปฏิปทา}, ทุกฺขา \textbf{ปฏิปทา} ทนฺธาภิญฺญา, ทุกฺขา \textbf{ปฏิปทา} ขิปฺปาภิญฺญา, สุขา \textbf{ปฏิปทา} ทนฺธาภิญฺญา, สุขา \textbf{ปฏิปทา} ขิปฺปาภิญฺญา. (21)}

\prose{อปราปิ จตสฺโส \textbf{ปฏิปทา}, อกฺขมา \textbf{ปฏิปทา}, ขมา \textbf{ปฏิปทา}, ทมา \textbf{ปฏิปทา}, สมา \textbf{ปฏิปทา}. (22)}

\prose{จตฺตาริ \textbf{ธมฺมปทานิ}, อนภิชฺฌา ธมฺมปทํ, อพฺยาปาโท ธมฺมปทํ, สมฺมาสติ ธมฺมปทํ, สมฺมาสมาธิ ธมฺมปทํ. (23)}

\prose{จตฺตาริ \textbf{ธมฺมสมาทานานิ}.\csromanspacer{} อตฺถาวุโส ธมฺมสมาทานํ ปจฺจุปฺปนฺนทุกฺขญฺเจว อายติํ จ ทุกฺขวิปากํ, อตฺถาวุโส ธมฺมสมาทานํ ปจฺจุปฺปนฺนทุกฺขํ อายติํ สุขวิปากํ, อตฺถาวุโส ธมฺมสมาทานํ ปจฺจุปฺปนฺนสุขํ อายติํ ทุกฺขวิปากํ, อตฺถาวุโส ธมฺมสมาทานํ ปจฺจุปฺปนฺนสุขญฺเจว อายติํ จ สุขวิปากํ. (24)}

\prose{จตฺตาโร \textbf{ธมฺมกฺขนฺธา}, สีลกฺขนฺโธ สมาธิกฺขนฺโธ ปญฺญากฺขนฺโธ วิมุตฺติกฺขนฺโธ. (25)}

\prose{จตฺตาริ \textbf{พลานิ}, วีริยพลํ สติพลํ สมาธิพลํ ปญฺญาพลํ. (26)}

\prose{จตฺตาริ \textbf{อธิฏฺฐานานิ}, ปญฺญาธิฏฺฐานํ, สจฺจาธิฏฺฐานํ, จาคาธิฏฺฐานํ, อุปสมาธิฏฺฐานํ. (27)}

\proseitem{312}{\csromanfolio{192}จตฺตาริ \textbf{ปญฺหพฺยากรณานิ}\footnote{จตฺตาโร ปญฺหาพฺยากรณา (สี, สฺยา, กํ, อิ)}, เอกํสพฺยากรณีโย ปโญฺห, ปฏิปุจฺฉาพฺยากรณีโย ปโญฺห, วิภชฺชพฺยากรณีโย ปโญฺห, ฐปนีโย ปโญฺห. (28)}

\prose{จตฺตาริ \textbf{กมฺมานิ}.\csromanspacer{} อตฺถาวุโส กมฺมํ กณฺหํ กณฺหวิปากํ, อตฺถาวุโส กมฺมํ สุกฺกํ สุกฺกวิปากํ, อตฺถาวุโส กมฺมํ กณฺหสุกฺกํ กณฺหสุกฺกวิปากํ, อตฺถาวุโส กมฺมํ อกณฺห-อสุกฺกํ อกณฺห-อสุกฺกวิปากํ กมฺมกฺขยาย สํวตฺตติ. (29)}

\prose{จตฺตาโร \textbf{สจฺฉิกรณียา} \textbf{ธมฺมา}.\csromanspacer{} ปุพฺเพนิวาโส สติยา สจฺฉิกรณีโย, สตฺตานํ จุตูปปาโต จกฺขุนา สจฺฉิกรณีโย, อฏฺฐวิโมกฺขา กาเยน \textbf{สจฺฉิกรณียา}, อาสวานํ ขโย ปญฺญาย สจฺฉิกรณีโย. (30)}

\prose{จตฺตาโร \textbf{โอฆา}, กาโมโฆ ภโวโฆ ทิฏฺโฐโฆ อวิชฺโชโฆ. (31)}

\prose{จตฺตาโร \textbf{โยคา}, กามโยโค ภวโยโค ทิฏฺฐิโยโค อวิชฺชาโยโค. (32)}

\prose{จตฺตาโร \textbf{วิสญฺโญคา}, กามโยควิสญฺโญโค ภวโยควิสญฺโญโค ทิฏฺฐิโยควิสญฺโญโค อวิชฺชาโยควิสญฺโญโค. (33)}

\prose{จตฺตาโร \textbf{คนฺถา}, อภิชฺฌา กายคนฺโถ พฺยาปาโท กายคนฺโถ สีลพฺพตปรามาโส กายคนฺโถ อิทํสจฺจาภินิเวโส กายคนฺโถ. (34)}

\prose{จตฺตาริ \textbf{อุปาทานานิ}, กามุปาทานํ\footnote{กามูปาทานํ (สี, อิ) เอวมิตเรสุปิ.} ทิฏฺฐุปาทานํ สีลพฺพตุปาทานํ อตฺตวาทุปาทานํ. (35)}

\prose{จตสฺโส \textbf{โยนิโย}, อณฺฑชโยนิ ชลาพุชโยนิ สํเสทชโยนิ โอปปาติกโยนิ. (36)}

\prose{จตสฺโส \textbf{คพฺภาวกฺกนฺติโย}.\csromanspacer{} อิธาวุโส เอกจฺโจ อสมฺปชาโน มาตุกุจฺฉิํ โอกฺกมติ, อสมฺปชาโน มาตุกุจฺฉิสฺมิํ ฐาติ, อสมฺปชาโน มาตุกุจฺฉิมฺหา นิกฺขมติ, อยํ ปฐมา คพฺภาวกฺกนฺติ.\csromanspacer{} ปุน จปรํ อาวุโส}

\vspace{\tipitakaparskip}
\csromancenter{สงฺคีติสุตฺต อิเธกจฺโจ สมฺปชาโน มาตุกุจฺฉิํ โอกฺกมติ, อสมฺปชาโน มาตุกุจฺฉิสฺมิํ ฐาติ, อสมฺปชาโน มาตุกุจฺฉิมฺหา นิกฺขมติ, อยํ ทุติยา คพฺภาวกฺกนฺติ.\csromanspacer{} ปุน จปรํ อาวุโส อิเธกจฺโจ สมฺปชาโน มาตุกุจฺฉิํ โอกฺกมติ, สมฺปชาโน มาตุกุจฺฉิสฺมิํ ฐาติ, อสมฺปชาโน มาตุกุจฺฉิมฺหา นิกฺขมติ, อยํ ตติยา คพฺภาวกฺกนฺติ.\csromanspacer{} ปุน จปรํ อาวุโส อิเธกจฺโจ สมฺปชาโน มาตุกุจฺฉิํ โอกฺกมติ, สมฺปชาโน มาตุกุจฺฉิสฺมิํ ฐาติ, สมฺปชาโน มาตุกุจฺจิมฺหา นิกฺขมติ, อยํ จตุตฺถา คพฺภาวกฺกนฺติ. (37)}
\prose{\csromanfolio{193}จตฺตาโร \textbf{อตฺตภาวปฏิลาภา}.\csromanspacer{} อตฺถาวุโส อตฺตภาวปฏิลาโภ, ยสฺมิํ อตฺตภาวปฏิลาเภ อตฺตสญฺเจตนาเยว กมติ, โน ปรสญฺเจตนา.\csromanspacer{} อตฺถาวุโส อตฺตภาวปฏิลาโภ, ยสฺมิํ อตฺตภาวปฏิลาเภ ปรสญฺเจตนาเยว กมติ, โน อตฺตสญฺเจตนา.\csromanspacer{} อตฺถาวุโส อตฺตภาวปฏิลาโภ, ยสฺมิํ อตฺตภาวปฏิลาเภ อตฺตสญฺเจตนา เจว กมติ ปรสญฺเจตนา จ.\csromanspacer{} อตฺถาวุโส อตฺตภาวปฏิลาโภ, ยสฺมิํ อตฺตภาวปฏิลาเภ เนว อตฺตสญฺเจตนา กมติ, โน ปรสญฺเจตนา. (38)}

\proseitem{313}{จตสฺโส \textbf{ทกฺขิณาวิสุทฺธิโย}.\csromanspacer{} อตฺถาวุโส ทกฺขิณา ทายกโต วิสุชฺฌติ โน ปฏิคฺคาหกโต, อตฺถาวุโส ทกฺขิณา ปฏิคฺคาหกโต วิสุชฺฌติ โน ทายกโต, อตฺถาวุโส ทกฺขิณา เนว ทายกโต วิสุชฺฌติ โน ปฏิคฺคาหกโต, อตฺถาวุโส ทกฺขิณา ทายกโต เจว วิสุชฺฌติ ปฏิคฺคาหกโต จ. (39)}

\prose{จตฺตาริ \textbf{สงฺคหวตฺถูนิ}, ทานํ เปยฺยวชฺชํ\footnote{ปิยวชฺชํ (สฺยา, กํ, ก)} อตฺถจริยา สมานตฺตตา. (40)}

\prose{จตฺตาโร \textbf{อนริยโวหารา}, มุสาวาโท ปิสุณาวาจา ผรุสาวาจา สมฺผปฺปลาโป. (41)}

\prose{จตฺตาโร \textbf{อริยโวหารา}, มุสาวาทา เวรมณี\footnote{เวรมณิ (ก)} ปิสุณาย วาจาย เวรมณี ผรุสาย วาจาย เวรมณี สมฺผปฺปลาปา เวรมณี. (42)}

\prose{อปเรปิ จตฺตาโร \textbf{อนริยโวหารา}, อทิฏฺเฐ ทิฏฺฐวาทิตา อสฺสุเต สุตวาทิตา อมุเต มุตวาทิตา อวิญฺญาเต วิญฺญาตวาทิตา. (43)}

\prose{\csromanfolio{194}อปเรปิ จตฺตาโร \textbf{อริยโวหารา}, อทิฏฺเฐ อทิฏฺฐวาทิตา อสฺสุเต อสฺสุตวาทิตา อมุเต อมุตวาทิตา อวิญฺญาเต อวิญฺญาตวาทิตา. (44)}

\prose{อปเรปิ จตฺตาโร \textbf{อนริยโวหารา}, ทิฏฺเฐ อทิฏฺฐวาทิตา สุเต อสฺสุตวาทิตา, มุเต อมุตวาทิตา, วิญฺญาเต อวิญฺญาตวาทิตา. (45)}

\prose{อปเรปิ จตฺตาโร \textbf{อริยโวหารา}, ทิฏฺเฐ ทิฏฺฐวาทิตา สุเต สุตวาทิตา มุเต มุตวาทิตา วิญฺญาเต วิญฺญาตวาทิตา. (46)}

\proseitem{314}{จตฺตาโร \textbf{ปุคฺคลา}.\csromanspacer{} อิธาวุโส เอกจฺโจ ปุคฺคโล อตฺตนฺตโป โหติ อตฺตปริตาปนานุโยคมนุยุตฺโต.\csromanspacer{} อิธาวุโส เอกจฺโจ ปุคฺคโล ปรนฺตโป โหติ ปรปริตาปนานุโยคมนุยุตฺโต.\csromanspacer{} อิธาวุโส เอกจฺโจ ปุคฺคโล อตฺตนฺตโป จ โหติ อตฺตปริตาปนานุโยคมนุยุตฺโต, ปรนฺตโป จ ปรปริตาปนานุโยคมนุยุตฺโต.\csromanspacer{} อิธปนาวุโส เอกจฺโจ ปุคฺคโล เนว อตฺตนฺตโป โหติ น อตฺตปริตาปนานุโยคมนุยุตฺโต, น ปรนฺตโป น ปรปริตาปนานุโยคมนุยุตฺโต, โส อนตฺตนฺตโป อปรนฺตโป ทิฏฺเฐว ธมฺเม นิจฺฉาโต นิพฺพุโต สีตีภูโต\footnote{สีติภูโต (ก)} สุขปฺปฏิสํเวที พฺรหฺมภูเตน อตฺตนา วิหรติ. (47)}

\prose{อปเรปิ จตฺตาโร \textbf{ปุคฺคลา}.\csromanspacer{} อิธาวุโส เอกจฺโจ ปุคฺคโล อตฺตหิตาย ปฏิปนฺโน โหติ โน ปรหิตาย, อิธาวุโส เอกจฺโจ ปุคฺคโล ปรหิตาย ปฏิปนฺโน โหติ โน อตฺตหิตาย, อิธาวุโส เอกจฺโจ ปุคฺคโล เนว อตฺตหิตาย ปฏิปนฺโน โหติ โน ปรหิตาย, อิธาวุโส เอกจฺโจ ปุคฺคโล อตฺตหิตาย เจว ปฏิปนฺโน โหติ ปรหิตาย จ. (48)}

\prose{อปเรปิ จตฺตาโร \textbf{ปุคฺคลา}, ตโม ตมปรายโน, ตโม โชติปรายโน, โชติ ตมปรายโน, โชติ โชติปรายโน. (49)}

\prose{อปเรปิ จตฺตาโร \textbf{ปุคฺคลา}, สมณมจโล สมณปทุโม สมณปุณฺฑรีโก สมเณสุ สมณสุขุมาโล. (50)}

\csromanheader{2}{10. สงฺคีติสุตฺต}
\prose{\csromanfolio{195}อิเม โข อาวุโส เตน ภควตา ชานตา ปสฺสตา อรหตา สมฺมาสมฺพุทฺเธน จตฺตาโร ธมฺมา สมฺมทกฺขาตา, ตตฺถ สพฺเพเหว สงฺคายิตพฺพํ ฯเปฯ อตฺถาย หิตาย สุขาย เทวมนุสฺสานํ.}

\nitthitamruled{ปฐมภาณวาโร นิฏฺฐิโต.}
\csromanmatikaanchor{mtk.132}
\csromantocmarkref{section}{ปญฺจก}{mtk.132}
\bookmark[dest={mtk.132},level=1]{ปญฺจก}
\csromanheader{1}{\textbf{ปญฺจก}}
\proseitem{315}{อตฺถิ โข อาวุโส เตน ภควตา ชานตา ปสฺสตา อรหตา สมฺมาสมฺพุทฺเธน ปญฺจ ธมฺมา สมฺมทกฺขาตา, ตตฺถ สพฺเพเหว สงฺคายิตพฺพํ ฯเปฯ อตฺถาย หิตาย สุขาย เทวมนุสฺสานํ.\csromanspacer{} กตเม ปญฺจ.}

\prose{\textbf{ปญฺจกฺขนฺธา}, รูปกฺขนฺโธ เวทนากฺขนฺโธ สญฺญากฺขนฺโธ สงฺขารกฺขนฺโธ วิญฺญาณกฺขนฺโธ. (1)}

\prose{\textbf{ปญฺจุปาทานกฺขนฺธา}, รูปุปาทานกฺขนฺโธ\footnote{รูปูปาทานกฺขนฺโธ (สี, สฺยา, กํ, อิ) เอวมิตเรสุปิ,} เวทนุปาทานกฺขนฺโธ สญฺญุปาทานกฺขนฺโธ สงฺขารุปาทานกฺขนฺโธ วิญฺญาณุปาทานกฺขนฺโธ. (2)}

\prose{ปญฺจ \textbf{กามคุณา}, จกฺขุวิญฺเญยฺยา รูปา อิฏฺฐา กนฺตา มนาปา ปิยรูปา กามูปสํหิตา รชนียา, โสตวิญฺเญยฺยา สทฺทา.\csromanspacer{} ฆานวิญฺเญยฺยา คนฺธา.\csromanspacer{} ชิวฺหาวิญฺเญยฺยา รสา.\csromanspacer{} กายวิญฺเญยฺยา โผฏฺฐพฺพา อิฏฺฐา กนฺตา มนาปา ปิยรูปา กามูปสํหิตา รชนียา. (3)}

\prose{ปญฺจ \textbf{คติโย}, นิรโย ติรจฺฉานโยนิ เปตฺติวิสโย มนุสฺสา เทวา. (4)}

\prose{ปญฺจ \textbf{มจฺฉริยานิ}, อาวาสมจฺฉริยํ กุลมจฺฉริยํ ลาภมจฺฉริยํ วณฺณมจฺฉริยํ ธมฺมมจฺฉริยํ. (5)}

\prose{ปญฺจ \textbf{นีวรณานิ}, กามจฺฉนฺทนีวรณํ พฺยาปาทนีวรณํ ถินมิทฺธนีวรณํ อุทฺธจฺจกุกฺกุจฺจนีวรณํ วิจิกิจฺฉานีวรณํ. (6)}

\prose{ปญฺจ \textbf{โอรมฺภาคิยานิ} สญฺโญชนานิ, สกฺกายทิฏฺฐิ วิจิกิจฺฉา สีลพฺพตปรามาโส กามจฺฉนฺโท พฺยาปาโท. (7)}

\prose{\csromanfolio{196}ปญฺจ \textbf{อุทฺธมฺภาคิยานิ} สญฺโญชนานิ, รูปราโค อรูปราโค มาโน อุทฺธจฺจํ อวิชฺชา. (8)}

\prose{ปญฺจ \textbf{สิกฺขาปทานิ}, ปาณาติปาตา เวรมณี อทินฺนาทานา เวรมณี กาเมสุมิจฺฉาจารา เวรมณี มุสาวาทา เวรมณี สุราเมรยมชฺชปฺปมาทฏฺฐานา เวรมณี. (9)}

\proseitem{316}{ปญฺจ \textbf{อภพฺพฏฺฐานานิ}.\csromanspacer{} อภพฺโพ อาวุโส ขีณาสโว ภิกฺขุ สญฺจิจฺจ ปาณํ ชีวิตา โวโรเปตุํ, อภพฺโพ ขีณาสโว ภิกฺขุ อทินฺนํ เถยฺยสงฺขาตํ อาทิยิตุํ\footnote{อาทาตุํ (สฺยา, กํ, อิ)}, อภพฺโพ ขีณาสโว ภิกฺขุ เมถุนํ ธมฺมํ ปฏิเสวิตุํ, อภพฺโพ ขีณาสโว ภิกฺขุ สมฺปชานมุสา ภาสิตุํ, อภพฺโพ ขีณาสโว ภิกฺขุ สนฺนิธิการกํ กาเม ปริภุญฺชิตุํ, เสยฺยถาปิ ปุพฺเพ อาคาริกภูโต. (10)}

\prose{ปญฺจ \textbf{พฺยสนานิ}, ญาติพฺยสนํ โภคพฺยสนํ โรคพฺยสนํ สีลพฺยสนํ ทิฏฺฐิพฺยสนํ.\csromanspacer{} นาวุโส สตฺตา ญาติพฺยสนเหตุ วา โภคพฺยสนเหตุ วา โรคพฺยสนเหตุ วา กายสฺส เภทา ปรํ มรณา อปายํ ทุคฺคติํ วินิปาตํ นิรยํ อุปปชฺชนฺติ, สีลพฺยสนเหตุ วา อาวุโส สตฺตา ทิฏฺฐิพฺยสนเหตุ วา กายสฺส เภทา ปรํ มรณา อปายํ ทุคฺคติํ วินิปาตํ นิรยํ อุปปชฺชนฺติ. (11)}

\prose{ปญฺจ \textbf{สมฺปทา}, ญาติ\textbf{สมฺปทา} โภค\textbf{สมฺปทา} อาโรคฺย\textbf{สมฺปทา} สีล\textbf{สมฺปทา} ทิฏฺฐิ\textbf{สมฺปทา}.\csromanspacer{} นาวุโส สตฺตา ญาติ\textbf{สมฺปทา}เหตุ วา โภค\textbf{สมฺปทา}เหตุ วา อาโรคฺย\textbf{สมฺปทา}เหตุ วา กายสฺส เภทา ปรํ มรณา สุคติํ สคฺคํ โลกํ อุปปชฺชนฺติ, สีล\textbf{สมฺปทา}เหตุ วา อาวุโส สตฺตา ทิฏฺฐิ\textbf{สมฺปทา}เหตุ วา กายสฺส เภทา ปรํ มรณา สุคติํ สคฺคํ โลกํ อุปปชฺชนฺติ. (12)}

\prose{ปญฺจ \textbf{อาทีนวา} \textbf{ทุสฺสีลสฺส} \textbf{สีลวิปตฺติยา}.\csromanspacer{} อิธาวุโส ทุสฺสีโล สีลวิปนฺโน ปมาทาธิกรณํ มหติํ โภคชานิํ นิคจฺฉติ, อยํ ปฐโม อาทีนโว \textbf{ทุสฺสีลสฺส} \textbf{สีลวิปตฺติยา}.\csromanspacer{} ปุน จปรํ อาวุโส \textbf{ทุสฺสีลสฺส} สีลวิปนฺนสฺส ปาปโก กิตฺติสทฺโท อพฺภุคฺคจฺฉติ, อยํ ทุติโย อาทีนโว \textbf{ทุสฺสีลสฺส} \textbf{สีลวิปตฺติยา}.\csromanspacer{} ปุน จปรํ อาวุโส ทุสฺสีโล}

\vspace{\tipitakaparskip}
\csromancenter{สงฺคีติสุตฺต สีลวิปนฺโน ยญฺญเทว ปริสํ อุปสงฺกมติ ยทิ ขตฺติยปริสํ ยทิ พฺราหฺมณปริสํ ยทิ คหปติปริสํ ยทิ สมณปริสํ, อวิสารโท อุปสงฺกมติ มงฺกุภูโต, อยํ ตติโย อาทีนโว ทุสฺสีลสฺส สีลวิปตฺติยา.\csromanspacer{} ปุน จปรํ อาวุโส ทุสฺสีโล สีลวิปนฺโน สมฺมูโฬฺห กาลํ กโรติ, อยํ จตุตฺโถ อาทีนโว ทุสฺสีลสฺส สีลวิปตฺติยา.\csromanspacer{} ปุน จปรํ อาวุโส ทุสฺสีโล สีลวิปนฺโน กายสฺส เภทา ปรํ มรณา อปายํ ทุคฺคติํ วินิปาตํ นิรยํ อุปปชฺชติ, อยํ \textbf{ปญฺจ}โม อาทีนโว ทุสฺสีลสฺส สีลวิปตฺติยา. (13)}
\prose{\csromanfolio{197}ปญฺจ \textbf{อานิสํสา} \textbf{สีลวโต} \textbf{สีลสมฺปทาย}.\csromanspacer{} อิธาวุโส สีลวา สีลสมฺปนฺโน อปฺปมาทาธิกรณํ มหนฺตํ โภคกฺขนฺนํ อธิคจฺฉติ, อยํ ปฐโม อานิสํโส \textbf{สีลวโต} \textbf{สีลสมฺปทาย}.\csromanspacer{} ปุน จปรํ อาวุโส \textbf{สีลวโต} สีลสมฺปนฺนสฺส กลฺยาโณ กิตฺติสทฺโท อพฺภุคฺคจฺฉติ, อยํ ทุติโย อานิสํโส \textbf{สีลวโต} \textbf{สีลสมฺปทาย}.\csromanspacer{} ปุน จปรํ อาวุโส สีลวา สีลสมฺปนฺโน ยญฺญเทว ปริสํ อุปสงฺกมติ ยทิ ขตฺติยปริสํ ยทิ พฺราหฺมณปริสํ ยทิ คหปติปริสํ ยทิ สมณปริสํ, วิสารโท อุปสงฺกมติ อมงฺกุภูโต, อยํ ตติโย อานิสํโส \textbf{สีลวโต} \textbf{สีลสมฺปทาย}.\csromanspacer{} ปุน จปรํ อาวุโส สีลวา สีลสมฺปนฺโน อสมฺมูโฬฺห กาลํ กโรติ, อยํ จตุตฺโถ อานิสํโส \textbf{สีลวโต} \textbf{สีลสมฺปทาย}.\csromanspacer{} ปุน จปรํ อาวุโส สีลวา สีลสมฺปนฺโน กายสฺส เภทา ปรํ มรณา สุคติํ สคฺคํ โลกํ อุปปชฺชติ, อยํ \textbf{ปญฺจ}โม อานิสํโส \textbf{สีลวโต} \textbf{สีลสมฺปทาย}. (14)}

\prose{โจทเกน อาวุโส ภิกฺขุนา ปรํ โจเทตุกาเมน \textbf{ปญฺจ} \textbf{ธมฺเม} อชฺฌตฺตํ อุปฏฺฐเปตฺวา ปโร โจเทตพฺโพ.\csromanspacer{} กาเลน วกฺขามิ โน อกาเลน, ภูเตน วกฺขามิ โน อภูเตน, สเณฺหน วกฺขามิ โน ผรุเสน, อตฺถสํหิเตน วกฺขามิ โน อนตฺถสํหิเตน, เมตฺตจิตฺเตน\footnote{เมตฺตาจิตฺเตน (กตฺถจิ)} วกฺขามิ โน โทสนฺตเรนาติ.\csromanspacer{} โจทเกน อาวุโส ภิกฺขุนา ปรํ โจเทตุกาเมน อิเม \textbf{ปญฺจ} \textbf{ธมฺเม} อชฺฌตฺตํ อุปฏฺฐเปตฺวา ปโร โจเทตพฺโพ. (15)}

\proseitem{317}{\csromanfolio{198}ปญฺจ \textbf{ปธานิยงฺคานิ}.\csromanspacer{} อิธาวุโส ภิกฺขุ สทฺโธ โหติ, สทฺทหติ ตถาคตสฺส โพธิํ “อิติปิ โส ภควา อรหํ สมฺมาสมฺพุทฺโธ วิชฺชาจรณสมฺปนฺโน สุคโต โลกวิทู อนุตฺตโร ปุริสทมฺมสารถิ สตฺถา เทวมนุสฺสานํ พุทฺโธ ภควา”ติ.\csromanspacer{} อปฺปาพาโธ โหติ อปฺปาตงฺโก, สมเวปากินิยา คหณิยา สมนฺนาคโต นาติสีตาย นาจฺจุณฺหาย มชฺฌิมาย ปธานกฺขมาย.\csromanspacer{} อสโฐ โหติ อมายาวี, ยถาภูตํ อตฺตานํ อาวิกตฺตา สตฺถริ วา วิญฺญูสุ วา สพฺรหฺมจารีสุ.\csromanspacer{} อารทฺธวีริโย วิหรติ อกุสลานํ ธมฺมานํ ปหานาย กุสลานํ ธมฺมานํ อุปสมฺปทาย ถามวา ทฬฺหปรกฺกโม อนิกฺขิตฺตธุโร กุสเลสุ ธมฺเมสุ.\csromanspacer{} ปญฺญวา โหติ อุทยตฺถคามินิยา ปญฺญาย สมนฺนาคโต อริยาย นิพฺเพธิกาย สมฺมาทุกฺขกฺขยคามินิยา. (16)}

\proseitem{318}{ปญฺจ \textbf{สุทฺธาวาสา}, อวิหา อตปฺปา สุทสฺสา สุทสฺสี อกนิฏฺฐา. (17)}

\prose{ปญฺจ \textbf{อนาคามิโน}, อนฺตราปรินิพฺพายี อุปหจฺจปรินิพฺพายี อสงฺขารปรินิพฺพายี สสงฺขารปรินิพฺพายี อุทฺธํโสโต-อกนิฏฺฐคามี. (18)}

\proseitem{319}{ปญฺจ \textbf{เจโตขิลา}.\csromanspacer{} อิธาวุโส ภิกฺขุ สตฺถริ กงฺขติ วิจิกิจฺฉติ นาธิมุจฺจติ น สมฺปสีทติ, โย โส อาวุโส ภิกฺขุ สตฺถริ กงฺขติ วิจิกิจฺฉติ นาธิมุจฺจติ น สมฺปสีทติ, ตสฺส จิตฺตํ น นมติ อาตปฺปาย อนุโยคาย สาตจฺจาย ปธานาย, ยสฺส จิตฺตํ น นมติ อาตปฺปาย อนุโยคาย สาตจฺจาย ปธานาย, อยํ ปฐโม เจโตขิโล.\csromanspacer{} ปุน จปรํ อาวุโส ภิกฺขุ ธมฺเม กงฺขติ วิจิกิจฺฉติ ฯเปฯ สํเฆ กงฺขติ วิจิกิจฺฉติ.\csromanspacer{} สิกฺขาย กงฺขติ วิจิกิจฺฉติ.\csromanspacer{} สพฺรหฺมจารีสุ กุปิโต โหติ อนตฺตมโน อาหตจิตฺโต ขิลชาโต, โย โส อาวุโส ภิกฺขุ สพฺรหฺมจารีสุ กุปิโต โหติ อนตฺตมโน อาหตจิตฺโต ขิลชาโต, ตสฺส จิตฺตํ น นมติ อาตปฺปาย อนุโยคาย สาตจฺจาย ปธานาย, ยสฺส จิตฺตํ น นมติ อาตปฺปาย อนุโยคาย สาตจฺจาย ปธานาย, อยํ ปญฺจโม เจโตขิโล. (19)}

\proseitem{320}{ปญฺจ \textbf{เจตโสวินิพนฺธา}.\csromanspacer{} อิธาวุโส ภิกฺขุ กาเมสุ อวีตราโค โหติ อวิคตจฺฉนฺโท อวิคตเปโม อวิคตปิปาโส}

\vspace{\tipitakaparskip}
\csromancenter{สงฺคีติสุตฺต อวิคตปริฬาโห อวิคตตโณฺห.\csromanspacer{} โย โส อาวุโส ภิกฺขุ กาเมสุ อวีตราโค โหติ อวิคตจฺฉนฺโท อวิคตเปโม อวิคตปิปาโส อวิคตปริฬาโห อวิคตตโณฺห, ตสฺส จิตฺตํ น นมติ อาตปฺปาย อนุโยคาย สาตจฺจาย ปธานาย, ยสฺส จิตฺตํ น นมติ อาตปฺปาย อนุโยคาย สาตจฺจาย ปธานาย, อยํ ปฐโม เจตโสวินิพนฺโธ.\csromanspacer{} ปุน จปรํ อาวุโส ภิกฺขุ กาเย อวีตราโค โหติ ฯเปฯ รูเป อวีตราโค โหติ ฯเปฯ.\csromanspacer{} ปุน จปรํ อาวุโส ภิกฺขุ ยาวทตฺถํ อุทราวเทหกํ ภุญฺชิตฺวา เสยฺยสุขํ ปสฺสสุขํ มิทฺธสุขํ อนุยุตฺโต วิหรติ ฯเปฯ.\csromanspacer{} ปุน จปรํ อาวุโส ภิกฺขุ อญฺญตรํ เทวนิกายํ ปณิธาย พฺรหฺมจริยํ จรติ “อิมินาหํ สีเลน วา วเตน วา ตเปน วา พฺรหฺมจริเยน วา เทโว วา ภวิสฺสามิ เทวญฺญตโร วา”ติ, โย โส อาวุโส ภิกฺขุ อญฺญตรํ เทวนิกายํ ปณิธาย พฺรหฺมจริยํ จรติ “อิมินาหํ สีเลน วา วเตน วา ตเปน วา พฺรหฺมจริเยน วา เทโว วา ภวิสฺสามิ เทวญฺญตโร วา”ติ, ตสฺส จิตฺตํ น นมติ อาตปฺปาย อนุโยคาย สาตจฺจาย ปธานาย, ยสฺส จิตฺตํ น นมติ อาตปฺปาย อนุโยคาย สาตจฺจาย ปธานาย, อยํ ปญฺจโม เจตโสวินิพนฺโธ. (20)}
\prose{\csromanfolio{199}\textbf{ปญฺจินฺทฺริยานิ}, จกฺขุนฺทฺริยํ โสตินฺทฺริยํ ฆานินฺทฺริยํ ชิวฺหินฺทฺริยํ กายินฺทฺริยํ. (2\textbf{1})}

\prose{อปรานิปิ \textbf{ปญฺจินฺทฺริยานิ}, สุขินฺทฺริยํ ทุกฺขินฺทฺริยํ โสมนสฺสินฺทฺริยํ โทมนสฺสินฺทฺริยํ อุเปกฺขินฺทฺริยํ. (22)}

\prose{อปรานิปิ \textbf{ปญฺจินฺทฺริยานิ}, สทฺธินฺทฺริยํ วีริยินฺทฺริยํ สตินฺทฺริยํ สมาธินฺทฺริยํ ปญฺญินฺทฺริยํ. (23)}

\proseitem{321}{ปญฺจ \textbf{นิสฺสรณิยา}\footnote{นิสฺสารณียา (สี, สฺยา, กํ, อิ) ฏีกา โอโลเกตพฺพา.} \textbf{ธาตุโย}.\csromanspacer{} อิธาวุโส ภิกฺขุโน กาเม มนสิกโรโต กาเมสุ จิตฺตํ น ปกฺขนฺทติ น ปสีทติ น สนฺติฏฺฐติ น วิมุจฺจติ, เนกฺขมฺมํ โข ปนสฺส มนสิกโรโต เนกฺขมฺเม จิตฺตํ ปกฺขนฺทติ ปสีทติ สนฺติฏฺฐติ วิมุจฺจติ, ตสฺส ตํ จิตฺตํ สุคตํ สุภาวิตํ สุวุฏฺฐิตํ สุวิมุตฺตํ วิสํยุตฺตํ กาเมหิ, เย จ กามปจฺจยา อุปฺปชฺชนฺติ อาสวา วิฆาตา ปริฬาหา\footnote{วิฆาตปริฬาหา (สฺยา, กํ)}, มุตฺโต โส เตหิ, น โส ตํ เวทนํ เวเทติ.\csromanspacer{} อิทมกฺขาตํ กามานํ นิสฺสรณํ. (24-\textbf{1})}

\prose{\csromanfolio{200}ปุน จปรํ อาวุโส ภิกฺขุโน พฺยาปาทํ มนสิกโรโต พฺยาปาเท จิตฺตํ น ปกฺขนฺทติ น ปสีทติ น สนฺติฏฺฐติ น วิมุจฺจติ, อพฺยาปาทํ โข ปนสฺส มนสิกโรโต อพฺยาปาเท จิตฺตํ ปกฺขนฺทติ ปสีทติ สนฺติฏฺฐติ วิมุจฺจติ, ตสฺส ตํ จิตฺตํ สุคตํ สุภาวิตํ สุวุฏฺฐิตํ สุวิมุตฺตํ วิสํยุตฺตํ พฺยาปาเทน, เย จ พฺยาปาทปจฺจยา อุปฺปชฺชนฺติ อาสวา วิฆาตา ปริฬาหา, มุตฺโต โส เตหิ, น โส ตํ เวทนํ เวเทติ.\csromanspacer{} อิทมกฺขาตํ พฺยาปาทสฺส นิสฺสรณํ. (24-2)}

\prose{ปุน จปรํ อาวุโส ภิกฺขุโน วิเหสํ มนสิกโรโต วิเหสาย จิตฺตํ น ปกฺขนฺทติ น ปสีทติ น สนฺติฏฺฐติ น วิมุจฺจติ, อวิเหสํ โข ปนสฺส มนสิกโรโต อวิเหสาย จิตฺตํ ปกฺขนฺทติ ปสีทติ สนฺติฏฺฐติ วิมุจฺจติ, ตสฺส ตํ จิตฺตํ สุคตํ สุภาวิตํ สุวุฏฺฐิตํ สุวิมุตฺตํ วิสํยุตฺตํ วิเหสาย, เย จ วิเหสาปจฺจยา อุปฺปชฺชนฺติ อาสวา วิฆาตา ปริฬาหา, มุตฺโต โส เตหิ, น โส ตํ เวทนํ เวเทติ.\csromanspacer{} อิทมกฺขาตํ วิเหสาย นิสฺสรณํ. (24-3)}

\prose{ปุน จปรํ อาวุโส ภิกฺขุโน รูเป มนสิกโรโต รูเปสุ จิตฺตํ น ปกฺขนฺทติ น ปสีทติ น สนฺติฏฺฐติ น วิมุจฺจติ, อรูปํ โข ปนสฺส มนสิกโรโต อรูเป จิตฺตํ ปกฺขนฺทติ ปสีทติ สนฺติฏฺฐติ วิมุจฺจติ, ตสฺส ตํ จิตฺตํ สุคตํ สุภาวิตํ สุวุฏฺฐิตํ สุวิมุตฺตํ วิสํยุตฺตํ รูเปหิ, เย จ รูปปจฺจยา อุปฺปชฺชนฺติ อาสวา วิฆาตา ปริฬาหา, มุตฺโต โส เตหิ, น โส ตํ เวทนํ เวเทติ.\csromanspacer{} อิทมกฺขาตํ รูปานํ นิสฺสรณํ. (24-4)}

\prose{ปุน จปรํ อาวุโส ภิกฺขุโน สกฺกายํ มนสิกโรโต สกฺกาเย จิตฺตํ น ปกฺขนฺทติ น ปสีทติ น สนฺติฏฺฐติ น วิมุจฺจติ, สกฺกายนิโรธํ โข ปนสฺส มนสิกโรโต สกฺกายนิโรเธ จิตฺตํ ปกฺขนฺทติ ปสีทติ สนฺติฏฺฐติ วิมุจฺจติ, ตสฺส ตํ จิตฺตํ สุคตํ สุภาวิตํ สุวุฏฺฐิตํ สุวิมุตฺตํ วิสํยุตฺตํ สกฺกาเยน, เย จ สกฺกายปจฺจยา อุปฺปชฺชนฺติ อาสวา วิฆาตา ปริฬาหา, มุตฺโต โส เตหิ, น โส ตํ เวทนํ เวเทติ.\csromanspacer{} อิทมกฺขาตํ สกฺกายสฺส นิสฺสรณํ. (24-5)}

\proseitem{322}{ปญฺจ \textbf{วิมุตฺตายตนานิ.}\csromanspacer{} อิธาวุโส ภิกฺขุโน สตฺถา ธมฺมํ เทเสติ อญฺญตโร วา ครุฏฺฐานิโก สพฺรหฺมจารี, ยถา ยถา อาวุโส ภิกฺขุโน สตฺถา ธมฺมํ เทเสติ อญฺญตโร วา ครุฏฺฐานิโย}

\vspace{\tipitakaparskip}
\csromancenter{สงฺคีติสุตฺต สพฺรหฺมจารี, ตถา ตถา โส ตสฺมิํ ธมฺเม อตฺถปฏิสํเวที จ โหติ ธมฺมปฏิสํเวที จ, ตสฺส อตฺถปฏิสํเวทิโน ธมฺมปฏิสํเวทิโน ปาโมชฺชํ ชายติ, ปมุทิตสฺส ปีติ ชายติ, ปีติมนสฺส กาโย ปสฺสมฺภติ, ปสฺสทฺธกาโย สุขํ เวเทติ, สุขิโน จิตฺตํ สมาธิยติ.\csromanspacer{} อิทํ ปฐมํ วิมุตฺตายตนํ.}
\prose{\csromanfolio{201}ปุน จปรํ อาวุโส ภิกฺขุโน น เหว โข สตฺถา ธมฺมํ เทเสติ อญฺญตโร วา ครุฏฺฐานิโก สพฺรหฺมจารี, อปิ จ โข ยถาสุตํ ยถาปริยตฺตํ ธมฺมํ วิตฺถาเรน ปเรสํ เทเสติ ฯเปฯ อปิ จ โข ยถาสุตํ ยถาปริยตฺตํ ธมฺมํ วิตฺถาเรน สชฺฌายํ กโรติ ฯเปฯ อปิ จ โข ยถาสุตํ ยถาปริยตฺตํ ธมฺมํ เจตสา อนุวิตกฺเกติ อนุวิจาเรติ มนสานุเปกฺขติ ฯเปฯ อปิ จ ขฺวสฺส อญฺญตรํ สมาธินิมิตฺตํ สุคหิตํ โหติ สุมนสิกตํ สูปธาริตํ สุปฺปฏิวิทฺธํ ปญฺญาย, ยถา ยถา อาวุโส ภิกฺขุโน อญฺญตรํ สมาธินิมิตฺตํ สุคหิตํ โหติ สุมนสิกตํ สูปธาริตํ สุปฺปฏิวิทฺธํ ปญฺญาย, ตถา ตถา โส ตสฺมิํ ธมฺเม อตฺถปฏิสํเวที จ โหติ ธมฺมปฏิสํเวที จ, ตสฺส อตฺถปฏิสํเวทิโน ธมฺมปฏิสํเวทิโน ปาโมชฺชํ ชายติ, ปมุทิตสฺส ปีติ ชายติ, ปีติมนสฺส กาโย ปสฺสมฺภติ, ปสฺสทฺธกาโย สุขํ เวเทติ, สุขิโน จิตฺตํ สมาธิยติ.\csromanspacer{} อิทํ ปญฺจมํ วิมุตฺตายตนํ. (25)}

\prose{ปญฺจ \textbf{วิมุตฺติปริปาจนียา} \textbf{สญฺญา}, อนิจฺจ\textbf{สญฺญา} อนิจฺเจ ทุกฺข\textbf{สญฺญา} ทุกฺเข อนตฺต\textbf{สญฺญา} ปหาน\textbf{สญฺญา} วิราค\textbf{สญฺญา}. (26)}

\prose{อิเม โข อาวุโส เตน ภควตา ชานตา ปสฺสตา อรหตา สมฺมาสมฺพุทฺเธน ปญฺจ ธมฺมา สมฺมทกฺขาตา, ตตฺถ สพฺเพเหว สงฺคายิตพฺพํ ฯเปฯ อตฺถาย หิตาย สุขาย เทวมนุสฺสานํ\footnote{สงฺคิติยปญฺจกํ นิฏฺฐิตํ (สฺยา, กํ)}.}

\csromanmatikaanchor{mtk.133}
\csromantocmarkref{section}{ฉกฺก}{mtk.133}
\bookmark[dest={mtk.133},level=1]{ฉกฺก}
\csromanheader{1}{\textbf{ฉกฺก}}
\proseitem{323}{อตฺถิ โข อาวุโส เตน ภควตา ชานตา ปสฺสตา อรหตา สมฺมาสมฺพุทฺเธน ฉ ธมฺมา สมฺมทกฺขาตา, ตตฺถ สพฺเพเหว สงฺคายิตพฺพํ ฯเปฯ อตฺถาย หิตาย สุขาย เทวมนุสฺสานํ.\csromanspacer{} กตเม ฉ.}

\prose{\csromanfolio{202}ฉ \textbf{อชฺฌตฺติกานิ} \textbf{อายตนานิ}, จกฺขายตนํ โสตายตนํ ฆานายตนํ ชิวฺหายตนํ กายายตนํ มนายตนํ. (1)}

\prose{ฉ \textbf{พาหิรานิ} \textbf{อายตนานิ}, รูปายตนํ สทฺทายตนํ คนฺธายตนํ รสายตนํ โผฏฺฐพฺพายตนํ ธมฺมายตนํ. (2)}

\prose{ฉ \textbf{วิญฺญาณกายา}, จกฺขุวิญฺญาณํ โสตวิญฺญาณํ ฆานวิญฺญาณํ ชิวฺหาวิญฺญาณํ กายวิญฺญาณํ มโนวิญฺญาณํ. (3)}

\prose{ฉ \textbf{ผสฺสกายา}, จกฺขุสมฺผสฺโส โสตสมฺผสฺโส ฆานสมฺผสฺโส ชิวฺหาสมฺผสฺโส กายสมฺผสฺโส มโนสมฺผสฺโส. (4)}

\prose{ฉ \textbf{เวทนากายา}, จกฺขุสมฺผสฺสชา เวทนา โสตสมฺผสฺสชา เวทนา ฆานสมฺผสฺสชา เวทนา ชิวฺหาสมฺผสฺสชา เวทนา กายสมฺผสฺสชา เวทนา มโนสมฺผสฺสชา เวทนา. (5)}

\prose{ฉ \textbf{สญฺญากายา}, รูปสญฺญา สทฺทสญฺญา คนฺธสญฺญา รสสญฺญา โผฏฺฐพฺพสญฺญา ธมฺมสญฺญา. (6)}

\prose{ฉ \textbf{สญฺเจตนากายา}, รูปสญฺเจตนา สทฺทสญฺเจตนา คนฺธสญฺเจตนา รสสญฺเจตนา โผฏฺฐพฺพสญฺเจตนา ธมฺมสญฺเจตนา. (7)}

\prose{ฉ \textbf{ตณฺหากายา}, รูปตณฺหา สทฺทตณฺหา คนฺธตณฺหา รสตณฺหา โผฏฺฐพฺพตณฺหา ธมฺมตณฺหา. (8)}

\proseitem{324}{ฉ \textbf{อคารวา}.\csromanspacer{} อิธาวุโส ภิกฺขุ สตฺถริ อคารโว วิหรติ อปฺปติสฺโส, ธมฺเม อคารโว วิหรติ อปฺปติสฺโส, สํเฆ อคารโว วิหรติ อปฺปติสฺโส, สิกฺขาย อคารโว วิหรติ อปฺปติสฺโส, อปฺปมาเท อคารโว วิหรติ อปฺปติสฺโส, ปฏิสนฺถาเร\footnote{ปฏิสนฺธาเร (ก)} อคารโว วิหรติ อปฺปติสฺโส. (9)}

\prose{ฉฺ\textbf{อ คารวา}.\csromanspacer{} อิธาวุโส ภิกฺขุ สตฺถริ สคารโว วิหรติ สปฺปติสฺโส, ธมฺเม สคารโว วิหรติ สปฺปติสฺโส, สํเฆ สคารโว วิหรหิ สปฺปติสฺโส, สิกฺขาย สคารโว วิหรติ สปฺปติสฺโส, อปฺปมาเท สคารโว วิหรติ สปฺปติสฺโส, ปฏิสนฺถาเร สคารโว วิหรติ สปฺปติสฺโส. (10)}

\csromanheader{2}{10. สงฺคีติสุตฺต}
\prose{\csromanfolio{203}ฉ \textbf{โสมนสฺสูปวิจารา}.\csromanspacer{} จกฺขุนา รูปํ ทิสฺวา โสมนสฺสฏฺฐานิยํ รูปํ อุปวิจรติ, โสเตน สทฺทํ สุตฺวา.\csromanspacer{} ฆาเนน คนฺธํ ฆายิตฺวา.\csromanspacer{} ชิวฺหาย รสํ สายิตฺวา.\csromanspacer{} กาเยน โผฏฺฐพฺพํ ผุสิตฺวา.\csromanspacer{} มนสา ธมฺมํ วิญฺญาย โสมนสฺสฏฺฐานิยํ ธมฺมํ อุปวิจรติ. (11)}

\prose{ฉ \textbf{โทมนสฺสูปวิจารา}.\csromanspacer{} จกฺขุนา รูปํ ทิสฺวา โทมนสฺสฏฺฐานิยํ รูปํ อุปวิจรติ ฯเปฯ.\csromanspacer{} มนสา ธมฺมํ วิญฺญาย โทมนสฺสฏฺฐานิยํ ธมฺมํ อุปวิจรติ. (12)}

\prose{ฉ \textbf{อุเปกฺขูปวิจารา}.\csromanspacer{} จกฺขุนา รูปํ ทิสฺวา อุเปกฺขาฏฺฐานิยํ\footnote{อุเปกฺขาฐานิยํ (ก)} รูปํ อุปวิจรติ ฯเปฯ.\csromanspacer{} มนสา ธมฺมํ วิญฺญาย อุเปกฺขาฏฺฐานิยํ ธมฺมํ อุปวิจรติ. (13)}

\prose{ฉ \textbf{สารณียา} \textbf{ธมฺมา}.\csromanspacer{} อิธาวุโส ภิกฺขุโน เมตฺตํ กายกมฺมํ ปจฺจุปฏฺฐิตํ โหติ สพฺรหฺมจารีสุ อาวิ\footnote{อาวี (ก-สี, อิ, ก)} เจว รโห จ, อยมฺปิ ธมฺโม สารณีโย ปิยกรโณ ครุกรโณ สงฺคหาย อวิวาทาย สามคฺคิยา เอกีภาวาย สํวตฺตติ. (14-1)}

\prose{ปุน จปรํ อาวุโส ภิกฺขุโน เมตฺตํ วจีกมฺมํ ปจฺจุปฏฺฐิตํ โหติ สพฺรหฺมจารีสุ อาวิ เจว รโห จ, อยมฺปิ ธมฺโม สารณีโย ฯเปฯ เอกีภาวาย สํวตฺตติ. (14-2)}

\prose{ปุน จปรํ อาวุโส ภิกฺขุโน เมตฺตํ มโนกมฺมํ ปจฺจุปฏฺฐิตํ โหติ สพฺรหฺมจารีสุ อาวิ เจว รโห จ, อยมฺปิ ธมฺโม สารณีโย ฯเปฯ เอกีภาวาย สํวตฺตติ. (14-3)}

\prose{ปุน จปรํ อาวุโส ภิกฺขุ เย เต ลาภา ธมฺมิกา ธมฺมลทฺธา อนฺตมโส ปตฺตปริยาปนฺนมตฺตมฺปิ, ตถารูเปหิ ลาเภหิ อปฺปฏิวิภตฺตโภคี โหติ สีลวนฺเตหิ สพฺรหฺมจารีหิ สาธารณโภคี, อยมฺปิ ธมฺโม สารณีโย ฯเปฯ เอกีภาวาย สํวตฺตติ. (14-4)}

\prose{ปุน จปรํ อาวุโส ภิกฺขุ ยานิ ตานิ สีลานิ อขณฺฑานิ อจฺฉิทฺทานิ อสพลานิ อกมฺมาสานิ ภุชิสฺสานิ วิญฺญุปฺปสตฺถานิ อปรามฏฺฐานิ สมาธิสํวตฺตนิกานิ, ตถารูเปสุ สีเลสุ สีลสามญฺญคโต วิหรติ สพฺรหฺมจารีหิ อาวิ เจว รโห จ, อยมฺปิ ธมฺโม สารณีโย ฯเปฯ เอกีภาวาย สํวตฺตติ. (14-5)}

\prose{\csromanfolio{204}ปุน จปรํ อาวุโส ภิกฺขุ ยายํ ทิฏฺฐิ อริยา นิยฺยานิกา นิยฺยาติ ตกฺกรสฺส สมฺมา ทุกฺขกฺขยาย, ตถารูปาย ทิฏฺฐิยา ทิฏฺฐิสามญฺญคโต วิหรติ สพฺรหฺมจารีหิ อาวิ เจว รโห จ, อยมฺปิ ธมฺโม สารณีโย ปิยกรโณ ครุกรโณ สงฺคหาย อวิวาทาย สามคฺคิยา เอกีภาวาย สํวตฺตติ. (14-6)}

\proseitem{325}{ฉ \textbf{วิวาทมูลานิ}.\csromanspacer{} อิธาวุโส ภิกฺขุ โกธโน โหติ อุปนาหี.\csromanspacer{} โย โส อาวุโส ภิกฺขุ โกธโน โหติ อุปนาหี, โส สตฺถริปิ อคารโว วิหรติ อปฺปติสฺโส, ธมฺเมปิ อคารโว วิหรติ อปฺปติสฺโส, สํเฆปิ อคารโว วิหรติ อปฺปติสฺโส, สิกฺขายปิ น ปริปูรการี\footnote{ปริปูรีการี (สฺยา, กํ)} โหติ.\csromanspacer{} โย โส อาวุโส ภิกฺขุ สตฺถริ อคารโว วิหรติ อปฺปติสฺโส, ธมฺเม อคารโว วิหรติ อปฺปติสฺโส, สํเฆ อคารโว วิหรติ อปฺปติสฺโส, สิกฺขาย น ปริปูรการี, โส สํเฆ วิวาทํ ชเนติ.\csromanspacer{} โย โส โหติ วิวาโท พหุชนหิตาย พหุชนสุขาย อนตฺถาย อหิตาย ทุกฺขาย เทวมนุสฺสานํ.\csromanspacer{} เอวรูปํ เจ ตุเมฺห อาวุโส วิวาทมูลํ อชฺฌตฺตํ วา พหิทฺธา วา สมนุปสฺเสยฺยาถ, ตตฺร ตุเมฺห อาวุโส ตสฺเสว ปาปกสฺส วิวาทมูลสฺส ปหานาย วายเมยฺยาถ.\csromanspacer{} เอวรูปํ เจ ตุเมฺห อาวุโส วิวาทมูลํ อชฺฌตฺตํ วา พหิทฺธา วา น สมนุปสฺเสยฺยาถ, ตตฺร ตุเมฺห อาวุโส ตสฺเสว ปาปกสฺส วิวาทมูลสฺส อายติํ อนวสฺสวาย ปฏิปชฺเชยฺยาถ.\csromanspacer{} เอวเมตสฺส ปาปกสฺส วิวาทมูลสฺส ปหานํ โหติ, เอวเมตสฺส ปาปกสฺส วิวาทมูลสฺส อายติํ อนวสฺสโว โหติ.}

\prose{ปุน จปรํ อาวุโส ภิกฺขุ มกฺขี โหติ ปฬาสี ฯเปฯ อิสฺสุกี โหติ มจฺฉรี.\csromanspacer{} สโฐ โหติ มายาวี.\csromanspacer{} ปาปิจฺโฉ โหติ มิจฺฉาทิฏฺฐี.\csromanspacer{} สนฺทิฏฺฐิปรามาสี โหติ อาธานคฺคาหี ทุปฺปฏินิสฺสคฺคี.\csromanspacer{} โย โส อาวุโส ภิกฺขุ สนฺทิฏฺฐิปรามาสี โหติ อาธานคฺคาหี ทุปฺปฏินิสฺสคฺคี, โส สตฺถริปิ อคารโว วิหรติ อปฺปติสฺโส, ธมฺเมปิ อคารโว วิหรติ อปฺปติสฺโส, สํเฆปิ อคารโว วิหรติ อปฺปติสฺโส, สิกฺขายปิ น ปริปูรการี โหติ.\csromanspacer{} โย โส อาวุโส ภิกฺขุ สตฺถริ อคารโว วิหรติ อปฺปติสฺโส, ธมฺเม อคารโว วิหรติ อปฺปติสฺโส, สํเฆ อคารโว วิหรติ}

\vspace{\tipitakaparskip}
\csromancenter{สงฺคีติสุตฺต อปฺปติสฺโส, สิกฺขาย น ปริปูรการี, โส สํเฆ วิวาทํ ชเนติ, โย โส โหติ วิวาโท พหุชนหิตาย พหุชนสุขาย อนตฺถาย อหิตาย ทุกฺขาย เทวมนุสฺสานํ.\csromanspacer{} เอวรูปํ เจ ตุเมฺห อาวุโส วิวาทมูลํ อชฺฌตฺตํ วา พหิทฺธา วา สมนุปสฺเสยฺยาถ, ตตฺร ตุเมฺห อาวุโส ตสฺเสว ปาปกสฺส วิวาทมูลสฺส ปหานาย วายเมยฺยาถ.\csromanspacer{} เอวรูปํ เจ ตุเมฺห อาวุโส วิวาทมูลํ อชฺฌตฺตํ วา พหิทฺธา วา น สมนุปสฺเสยฺยาถ, ตตฺร ตุเมฺห อาวุโส ตสฺเสว ปาปกสฺส วิวาทมูลสฺส อายติํ อนวสฺสวาย ปฏิปชฺเชยฺยาถ.\csromanspacer{} เอวเมตสฺส ปาปกสฺส วิวาทมูลสฺส ปหานํ โหติ, เอวเมตสฺส ปาปกสฺส วิวาทมูลสฺส อายติํ อนวสฺสโว โหติ. (15)}
\prose{\csromanfolio{205}ฉ \textbf{ธาตุโย}.\csromanspacer{} ปถวีธาตุ อาโปธาตุ เตโชธาตุ วาโยธาตุ อากาสธาตุ วิญฺญาณธาตุ. (16)}

\proseitem{326}{ฉ \textbf{นิสฺสรณิยา} \textbf{ธาตุโย}.\csromanspacer{} อิธาวุโส ภิกฺขุ เอวํ วเทยฺย “เมตฺตา หิ โข เม เจโตวิมุตฺติ ภาวิตา พหุลีกตา ยานีกตา วตฺถุกตา อนุฏฺฐิตา ปริจิตา สุสมารทฺธา, อถ จ ปน เม พฺยาปาโท จิตฺตํ ปริยาทาย ติฏฺฐตี”ติ.\csromanspacer{} โส “มา เหวนฺ”ติสฺส วจนีโย “มายสฺมา เอวํ อวจ, มา ภควนฺตํ อพฺภาจิกฺขิ, น หิ สาธุ ภควโต อพฺภกฺขานํ, น หิ ภควา เอวํ วเทยฺย.\csromanspacer{} อฏฺฐานเมตํ อาวุโส อนวกาโส, ยํ เมตฺตาย เจโตวิมุตฺติยา ภาวิตาย พหุลีกตาย ยานีกตาย วตฺถุกตาย อนุฏฺฐิตาย ปริจิตาย สุสมารทฺธาย, อถ จ ปนสฺส พฺยาปาโท จิตฺตํ ปริยาทาย ฐสฺสติ, เนตํ ฐานํ วิชฺชติ.\csromanspacer{} นิสฺสรณํ เหตํ อาวุโส พฺยาปาทสฺส, ยทิทํ เมตฺตา เจโตวิมุตฺตี”ติ. (17-1)}

\prose{อิธ ปนาวุโส ภิกฺขุ เอวํ วเทยฺย “กรุณา หิ โข เม เจโตวิมุตฺติ ภาวิตา พหุลีกตา ยานีกตา วตฺถุกตา อนุฏฺฐิตา ปริจิตา สุสมารทฺธา, อถ จ ปน เม วิเหสา จิตฺตํ ปริยาทาย ติฏฺฐตี”ติ.\csromanspacer{} โส “มา เหวนฺ”ติสฺส วจนีโย “มายสฺมา เอวํ อวจ, มา ภควนฺตํ อพฺภาจิกฺขิ, น หิ สาธุ ภควโต อพฺภกฺขานํ, น หิ ภควา เอวํ วเทยฺย.\csromanspacer{} อฏฺฐานเมตํ อาวุโส อนวกาโส, ยํ กรุณาย เจโตวิมุตฺติยา ภาวิตาย พหุลีกตาย ยานีกตาย \csromanfolio{206}วตฺถุกตาย อนุฏฺฐิตาย ปริจิตาย สุสมารทฺธาย, อถ จ ปนสฺส วิเหสา จิตฺตํ ปริยาทาย ฐสฺสติ, เนตํ ฐานํ วิชฺชติ.\csromanspacer{} นิสฺสรณํ เหตํ อาวุโส วิเหสาย, ยทิทํ กรุณา เจโตวิมุตฺตี”ติ. (17-2)}

\prose{อิธ ปนาวุโส ภิกฺขุ เอวํ วเทยฺย “มุทิตา หิ โข เม เจโตวิมุตฺติ ภาวิตา พหุลีกตา ยานีกตา วตฺถุกตา อนุฏฺฐิตา ปริจิตา สุสมารทฺธา, อถ จ ปน เม อรติ จิตฺตํ ปริยาทาย ติฏฺฐตี”ติ.\csromanspacer{} โส “มา เหวนฺ”ติสฺส วจนีโย “มายสฺมา เอวํ อวจ, มา ภควนฺตํ อพฺภาจิกฺขิ, น หิ สาธุ ภควโต อพฺภกฺขานํ, น หิ ภควา เอวํ วเทยฺย.\csromanspacer{} อฏฺฐานเมตํ อาวุโส อนวกาโส, ยํ มุทิตาย เจโตวิมุตฺติยา ภาวิตาย พหุลีกตาย ยานีกตาย วตฺถุกตาย อนุฏฺฐิตาย ปริจิตาย สุสมารทฺธาย, อถ จ ปนสฺส อรติ จิตฺตํ ปริยาทาย ฐสฺสติ, เนตํ ฐานํ วิชฺชติ.\csromanspacer{} นิสฺสรณํ เหตํ อาวุโส อรติยา, ยทิทํ มุทิตา เจโตวิมุตฺตี”ติ. (17-3)}

\prose{อิธ ปนาวุโส ภิกฺขุ เอวํ วเทยฺย “อุเปกฺขา หิ โข เม เจโตวิมุตฺติ ภาวิตา พหุลีกตา ยานีกตา วตฺถุกตา อนุฏฺฐิตา ปริจิตา สุสมารทฺธา, อถ จ ปน เม ราโค จิตฺตํ ปริยาทาย ติฏฺฐตี”ติ.\csromanspacer{} โส “มา เหวนฺ”ติสฺส วจนีโย “มายสฺมา เอวํ อวจ, มา ภควนฺตํ อพฺภาจิกฺขิ, น หิ สาธุ ภควโต อพฺภกฺขานํ, น หิ ภควา เอวํ วเทยฺย.\csromanspacer{} อฏฺฐานเมตํ อาวุโส อนวกาโส, ยํ อุเปกฺขาย เจโตวิมุตฺติยา ภาวิตาย พหุลีกตาย ยานีกตาย วตฺถุกตาย อนุฏฺฐิตาย ปริจิตาย สุสมารทฺธาย, อถ จ ปนสฺส ราโค จิตฺตํ ปริยาทาย ฐสฺสติ, เนตํ ฐานํ วิชฺชติ.\csromanspacer{} นิสฺสรณํ เหตํ อาวุโส ราคสฺส, ยทิทํ อุเปกฺขา เจโตวิมุตฺตี”ติ. (17-4)}

\prose{อิธ ปนาวุโส ภิกฺขุ เอวํ วเทยฺย “อนิมิตฺตา หิ โข เม เจโตวิมุตฺติ ภาวิตา พหุลีกตา ยานีกตา วตฺถุกตา อนุฏฺฐิตา ปริจิตา สุสมารทฺธา, อถ จ ปน เม นิมิตฺตานุสาริ วิญฺญาณํ โหตี”ติ.\csromanspacer{} โส “มา เหวนฺ”ติสฺส วจนีโย “มายสฺมา เอวํ อวจ, มา ภควนฺตํ อพฺภาจิกฺขิ, น หิ สาธุ ภควโต อพฺภกฺขานํ, น หิ ภควา เอวํ วเทยฺย, อฏฺฐานเมตํ อาวุโส อนวกาโส, ยํ อนิมิตฺตาย เจโตวิมุตฺติยา ภาวิตาย พหุลีกตาย ยานีกตาย วตฺถุกตาย อนุฏฺฐิตาย ปริจิตาย สุสมารทฺธาย, อถ จ ปนสฺส}

\vspace{\tipitakaparskip}
\csromancenter{สงฺคีติสุตฺต นิมิตฺตานุสาริ วิญฺญาณํ ภวิสฺสติ, เนตํ ฐานํ วิชฺชติ.\csromanspacer{} นิสฺสรณํ เหตํ อาวุโส สพฺพนิมิตฺตานํ, ยทิทํ อนิมิตฺตา เจโตวิมุตฺตี”ติ. (17-5)}
\prose{\csromanfolio{207}อิธ ปนาวุโส ภิกฺขุ เอวํ วเทยฺย “อสฺมีติ โข เม วิคตํ\footnote{วิฆาตํ (สี, อิ), วิคเต (สฺยา, ก)}, อยมหมสฺมีติ น สมนุปสฺสามิ, อถ จ ปน เม วิจิกิจฺฉากถํกถาสลฺลํ จิตฺตํ ปริยาทาย ติฏฺฐตี”ติ.\csromanspacer{} โส “มา เหวนฺ”ติสฺส วจนีโย, มายสฺมา เอวํ อวจ, มา ภควนฺตํ อพฺภาจิกฺขิ, น หิ สาธุ ภควโต อพฺภกฺขานํ, น หิ ภควา เอวํ วเทยฺย.\csromanspacer{} อฏฺฐานเมตํ อาวุโส อนวกาโส, ยํ อสฺมีติ วิคเต\footnote{วิฆาเต (สี, อิ)} อยมหมสฺมีติ อสมนุปสฺสโต, อถ จ ปนสฺส วิจิกิจฺฉากถํกถาสลฺลํ จิตฺตํ ปริยาทาย ฐสฺสติ, เนตํ ฐานํ วิชฺชติ.\csromanspacer{} นิสฺสรณํ เหตํ อาวุโส วิจิกิจฺฉากถํกถาสลฺลสฺส, ยทิทํ อสฺมิมานสมุคฺฆาโตติ. (17-6)}

\proseitem{327}{ฉ \textbf{อนุตฺตริยานิ}, ทสฺสนานุตฺตริยํ สวนานุตฺตริยํ ลาภานุตฺตริยํ สิกฺขานุตฺตริยํ ปาริจริยานุตฺตริยํ อนุสฺสตานุตฺตริยํ. (18)}

\prose{ฉ \textbf{อนุสฺสติฏฺฐานานิ}, พุทฺธานุสฺสติ ธมฺมานุสฺสติ สํฆานุสฺสติ สีลานุสฺสติ จาคานุสฺสติ เทวตานุสฺสติ. (19)}

\proseitem{328}{ฉ \textbf{สตตวิหารา}.\csromanspacer{} อิธาวุโส ภิกฺขุ จกฺขุนา รูปํ ทิสฺวา เนว สุมโน โหติ น ทุมฺมโน, อุเปกฺขโก\footnote{อุเปกฺขโก จ (สฺยา, ก)} วิหรติ สโต สมฺปชาโน.\csromanspacer{} โสเตน สทฺทํ สุตฺวา ฯเปฯ.\csromanspacer{} มนสา ธมฺมํ วิญฺญาย เนว สุมโน โหติ น ทุมฺมโน, อุเปกฺขโก วิหรติ สโต สมฺปชาโน. (20)}

\proseitem{329}{\textbf{ฉฬาภิชาติโย}.\csromanspacer{} อิธาวุโส เอกจฺโจ กณฺหาภิชาติโก สมาโน กณฺหํ ธมฺมํ อภิชายติ, อิธ ปนาวุโส เอกจฺโจ กณฺหาภิชาติโก สมาโน สุกฺกํ ธมฺมํ อภิชายติ, อิธ ปนาวุโส เอกจฺโจ กณฺหาภิชาติโก สมาโน อกณฺหํ อสุกฺกํ นิพฺพานํ อภิชายติ, อิธ ปนาวุโส เอกจฺโจ สุกฺกาภิชาติโก สมาโน สุกฺกํ ธมฺมํ อภิชายติ, อิธ ปนาวุโส เอกจฺโจ สุกฺกาภิชาติโก สมาโน กณฺหํ ธมฺมํ อภิชายติ, อิธ ปนาวุโส เอกจฺโจ สุกฺกาภิชาติโก \csromanfolio{208}สมาโน อกณฺหํ อสุกฺกํ นิพฺพานํ อภิชายติ. (21)}

\prose{ฉ \textbf{นิพฺเพธภาคิยา} สญฺญา\footnote{นิพฺเพธภาคิยสญฺญา (สฺยา, กํ)}, อนิจฺจสญฺญา อนิจฺเจ ทุกฺขสญฺญา ทุกฺเข อนตฺตสญฺญา ปหานสญฺญา วิราคสญฺญา นิโรธสญฺญา. (22)}

\prose{อิเม โข อาวุโส เตน ภควตา ชานตา ปสฺสตา อรหตา สมฺมาสมฺพุทฺเธน ฉ ธมฺมา สมฺมทกฺขาตา, ตตฺถ สพฺเพเหว สงฺคายิตพฺพํ ฯเปฯ อตฺถาย หิตาย สุขาย เทวมนุสฺสานํ.}

\csromanmatikaanchor{mtk.134}
\csromantocmarkref{section}{สตฺตก}{mtk.134}
\bookmark[dest={mtk.134},level=1]{สตฺตก}
\csromanheader{1}{\textbf{สตฺตก}}
\proseitem{330}{อตฺถิ โข อาวุโส เตน ภควตา ชานตา ปสฺสตา อรหตา สมฺมาสมฺพุทฺเธน สตฺต ธมฺมา สมฺมทกฺขาตา, ตตฺถ สพฺเพเหว สงฺคายิตพฺพํ ฯเปฯ อตฺถาย หิตาย สุขาย เทวมนุสฺสานํ.\csromanspacer{} กเตเม สตฺต.}

\prose{สตฺต \textbf{อริยธนานิ}, สทฺธาธนํ สีลธนํ หิริธนํ โอตฺตปฺปธนํ สุตธนํ จาคธนํ ปญฺญาธนํ. (1)}

\prose{สตฺต \textbf{โพชฺฌงฺคา}, สติสมฺโพชฺฌงฺโค ธมฺมวิจยสมฺโพชฺฌงฺโค วีริยสมฺโพชฺฌงฺโค ปีติสมฺโพชฺฌงฺโค ปสฺสทฺธิสมฺโพชฺฌงฺโค สมาธิสมฺโพชฺฌงฺโค อุเปกฺขาสมฺโพชฺฌงฺโค. (2)}

\prose{สตฺต \textbf{สมาธิปริกฺขารา}, สมฺมาทิฏฺฐิ สมฺมาสงฺกปฺโป สมฺมาวาจา สมากมฺมนฺโต สมฺมา-อาชีโว สมฺมาวายาโม สมฺมาสติ. (3)}

\prose{สตฺต \textbf{อสทฺธมฺมา}.\csromanspacer{} อิธาวุโส ภิกฺขุ อสฺสทฺโธ โหติ, อหิริโก โหติ, อโนตฺตปฺปี โหติ, อปฺปสฺสุโต โหติ, กุสีโต โหติ, มุฏฺฐสฺสติ โหติ, ทุปฺปญฺโญ โหติ. (4)}

\prose{สตฺตฺ\textbf{อ สทฺธมฺมา}.\csromanspacer{} อิธาวุโส ภิกฺขุ สทฺโธ โหติ, หิริมา โหติ, โอตฺตปฺปี โหติ, พหุสฺสุโต โหติ, อารทฺธวีริโย โหติ, อุปฏฺฐิตสฺสติ โหติ, ปญฺญวา โหติ. (5)}

\prose{สตฺต \textbf{สปฺปุริสธมฺมา}.\csromanspacer{} อิธาวุโส ภิกฺขุ ธมฺมญฺญู จ โหติ อตฺถญฺญู จ อตฺตญฺญู จ มตฺตญฺญู จ กาลญฺญู จ ปริสญฺญู จ ปุคฺคลญฺญู จ. (6)}

\csromanheader{2}{10. สงฺคีติสุตฺต}
\proseitem{331}{\csromanfolio{209}สตฺต นิทฺทสวตฺถูนิ.\csromanspacer{} อิธาวุโส ภิกฺขุ สิกฺขาสมาทาเน ติพฺพจฺฉนฺโท โหติ, อายติํ จ สิกฺขาสมาทาเน อวิคตเปโม.\csromanspacer{} ธมฺมนิสนฺติยา ติพฺพจฺฉนฺโท โหติ, อายติํ จ ธมฺมนิสนฺติยา อวิคตเปโม.\csromanspacer{} อิจฺฉาวินเย ติพฺพจฺฉนฺโท โหติ, อายติํ จ อิจฺฉาวินเย อวิคตเปโม.\csromanspacer{} ปฏิสลฺลาเน ติพฺพจฺฉนฺโท โหติ, อายติํ จ ปฏิสลฺลาเน อวิคตเปโม.\csromanspacer{} วีริยารมฺเภ ติพฺพจฺฉนฺโท โหติ, อายติํ จ วีริยารมฺเภ อวิคตเปโม.\csromanspacer{} สติเนปกฺเก ติพฺพจฺฉนฺโท โหติ, อายติํ จ สติเนปกฺเก อวิคตเปโม.\csromanspacer{} ทิฏฺฐิปฏิเวเธ ติพฺพจฺฉนฺโท โหติ, อายติํ จ ทิฏฺฐิปฏิเวเธ อวิคตเปโม. (7)}

\prose{สตฺต \textbf{สญฺญา}, อนิจฺจ\textbf{สญฺญา} อนตฺต\textbf{สญฺญา} อสุภ\textbf{สญฺญา} อาทีนว\textbf{สญฺญา} ปหาน\textbf{สญฺญา} วิราค\textbf{สญฺญา} นิโรธ\textbf{สญฺญา}. (8)}

\prose{สตฺต \textbf{พลานิ}, สทฺธาพลํ วีริยพลํ หิริพลํ โอตฺตปฺปพลํ สติพลํ สมาธิพลํ ปญฺญาพลํ. (9)}

\proseitem{332}{สตฺต \textbf{วิญฺญาณฏฺฐิติโย}.\csromanspacer{} สนฺตาวุโส สตฺตา นานตฺตกายา นานตฺตสญฺญิโน เสยฺยถาปิ มนุสฺสา เอกจฺเจ จ เทวา เอกจฺเจ จ วินิปาติกา, อยํ ปฐมา วิญฺญาณฏฺฐิติ.}

\prose{สนฺตาวุโส สตฺตา นานตฺตกายา เอกตฺตสญฺญิโน เสยฺยถาปิ เทวา พฺรหฺมกายิกา ปฐมาภินิพฺพตฺตา, อยํ ทุติยา วิญฺญาณฏฺฐิติ.}

\prose{สนฺตาวุโส สตฺตา เอกตฺตกายา นานตฺตสญฺญิโน เสยฺยถาปิ เทวา อาภสฺสรา, อยํ ตติยา วิญฺญาณฏฺฐิติ.}

\prose{สนฺตาวุโส สตฺตา เอกตฺตกายา เอกตฺตสญฺญิโน เสยฺยถาปิ เทวา สุภกิณฺหา, อยํ จตุตฺถี วิญฺญาณฏฺฐิติ.}

\prose{สนฺตาวุโส สตฺตา สพฺพโส รูป\textbf{สญฺญา}นํ สมติกฺกมา ปฏิฆ\textbf{สญฺญา}นํ อตฺถงฺคมา นานตฺต\textbf{สญฺญา}นํ อมนสิการา “อนนฺโต อากาโส”ติ อากาสานญฺจายตนูปคา, อยํ ปญฺจมี วิญฺญาณฏฺฐิติ.}

\prose{สนฺตาวุโส สตฺตา สพฺพโส อากาสานญฺจายตนํ สมติกฺกมฺม “อนนฺตํ วิญฺญาณนฺ”ติ วิญฺญาณญฺจายตนูปคา, อยํ ฉฏฺฐี วิญฺญาณฏฺฐิติ.}

\prose{\csromanfolio{210}สนฺตาวุโส สตฺตา สพฺพโส วิญฺญาณญฺจายตนํ สมติกฺกมฺม “นตฺถิ กิญฺจี”ติ อากิญฺจญฺญายตนูปคา, อยํ สตฺตมี วิญฺญาณฏฺฐิติ. (10)}

\prose{สตฺต ปุคฺคลา \textbf{ทกฺขิเณยฺยา}, อุภโตภาควิมุตฺโต ปญฺญาวิมุตฺโต กายสกฺขิ ทิฏฺฐิปฺปตฺโต สทฺธาวิมุตฺโต ธมฺมานุสารี สทฺธานุสารี. (11)}

\prose{สตฺต \textbf{อนุสยา}, กามราคานุสโย ปฏิฆานุสโย ทิฏฺฐานุสโย วิจิกิจฺฉานุสโย มานานุสโย ภวราคานุสโย อวิชฺชานุสโย. (12)}

\prose{สตฺต \textbf{สญฺโญชนานิ}, อนุนยสญฺโญชนํ\footnote{กามสญฺโญชนํ (สฺยา, กํ)} ปฏิฆสญฺโญชนํ ทิฏฺฐิสญฺโญชนํ วิจิกิจฺฉาสญฺโญชนํ มานสญฺโญชนํ ภวราคสญฺโญชนํ อวิชฺชาสญฺโญชนํ. (13)}

\prose{สตฺต \textbf{อธิกรณสมถา}, อุปฺปนฺนุปฺปนฺนานํ อธิกรณานํ สมถาย วูปสมาย สมฺมุขาวินโย ทาตพฺโพ, สติวินโย ทาตพฺโพ, อมูฬฺหวินโย ทาตพฺโพ, ปฏิญฺญาย กาเรตพฺพํ, เยภุยฺยสิกา, ตสฺสปาปิยสิกา, ติณวตฺถารโก. (14)}

\prose{อิเม โข อาวุโส เตน ภควตา ชานตา ปสฺสตา อรหตา สมฺมาสมฺพุทฺเธน สตฺต ธมฺมา สมฺมทกฺขาตา, ตตฺถ สพฺเพเหว สงฺคายิตพฺพํ ฯเปฯ อตฺถาย หิตาย สุขาย เทวมนุสฺสานํ.}

\nitthitamruled{ทุติยภาณวาโร นิฏฺฐิโต.}
\csromanmatikaanchor{mtk.135}
\csromantocmarkref{section}{อฏฺฐก}{mtk.135}
\bookmark[dest={mtk.135},level=1]{อฏฺฐก}
\csromanheader{1}{\textbf{อฏฺฐก}}
\proseitem{333}{อตฺถิ โข อาวุโส เตน ภควตา ชานตา ปสฺสตา อรหตา สมฺมาสมฺพุทฺเธน อฏฺฐ ธมฺมา สมฺมทกฺขาตา, ตตฺถ สพฺเพเหว สงฺคายิตพฺพํ ฯเปฯ อตฺถาย หิตาย สุขาย เทวมนุสฺสานํ.\csromanspacer{} กตเม อฏฺฐ.}

\prose{อฏฺฐ \textbf{มิจฺฉตฺตา}, มิจฺฉาทิฏฺฐิ มิจฺฉาสงฺกปฺโป มิจฺฉาวาจา มิจฺฉากมฺมนฺโต มิจฺฉา-อาชีโว มิจฺฉาวายาโม มิจฺฉาสติ มิจฺฉาสมาธิ. (1)}

\csromanheader{2}{10. สงฺคีติสุตฺต}
\prose{\csromanfolio{211}อฏฺฐ \textbf{สมฺมตฺตา}, สมฺมาทิฏฺฐิ สมฺมาสงฺกปฺโป สมฺมาวาจา สมฺมากมฺมนฺโต สมฺมา-อาชีโว สมฺมาวายาโม สมฺมาสติ สมฺมาสมาธิ. (2)}

\prose{อฏฺฐ ปุคฺคลา \textbf{ทกฺขิเณยฺยา}, โสตาปนฺโน, โสตาปตฺติผลสจฺฉิกิริยาย ปฏิปนฺโน, สกทาคามี, สกทาคามิผลสจฺฉิกิริยาย ปฏิปนฺโน, อนาคามี, อนาคามิผลสจฺฉิกิริยาย ปฏิปนฺโน, อรหา, อรหตฺตผลสจฺฉิกิริยาย ปฏิปนฺโน. (3)}

\proseitem{334}{อฏฺฐ \textbf{กุสีตวตฺถูนิ}.\csromanspacer{} อิธาวุโส ภิกฺขุนา กมฺมํ กาตพฺพํ โหติ, ตสฺส เอวํ โหติ “กมฺมํ โข เม กาตพฺพํ ภวิสฺสติ, กมฺมํ โข ปน เม กโรนฺตสฺส กาโย กิลมิสฺสติ, หนฺทาหํ นิปชฺชามี”ติ.\csromanspacer{} โส นิปชฺชติ น วีริยํ อารภติ อปฺปตฺตสฺส ปตฺติยา อนธิคตสฺส อธิคมาย อสจฺฉิกตสฺส สจฺฉิกิริยาย, อิทํ ปฐมํ กุสีตวตฺถุ.}

\prose{ปุน จปรํ อาวุโส ภิกฺขุนา กมฺมํ กตํ โหติ, ตสฺส เอวํ โหติ “อหํ โข กมฺมํ อกาสิํ, กมฺมํ โข ปน เม กโรนฺตสฺส กาโย กิลนฺโต, หนฺทาหํ นิปชฺชามี”ติ.\csromanspacer{} โส นิปชฺชติ น วีริยํ อารภติ ฯเปฯ อิทํ ทุติยํ กุสีตวตฺถุ.}

\prose{ปุน จปรํ อาวุโส ภิกฺขุนา มคฺโค คนฺตพฺโพ โหติ, ตสฺส เอวํ โหติ “มคฺโค โข เม คนฺตพฺโพ ภวิสฺสติ, มคฺคํ โข ปน เม คจฺฉนฺตสฺส กาโย กิลมิสฺสติ, หนฺทาหํ นิปชฺชามี”ติ.\csromanspacer{} โส นิปชฺชติ น วีริยํ อารภติ, อิทํ ตติยํ กุสีตวตฺถุ.}

\prose{ปุน จปรํ อาวุโส ภิกฺขุนา มคฺโค คโต โหติ, ตสฺส เอวํ โหติ “อหํ โข มคฺคํ อคมาสิํ, มคฺคํ โข ปน เม คจฺฉนฺตสฺส กาโย กิลนฺโต, หนฺทาหํ นิปชฺชามี”ติ.\csromanspacer{} โส นิปชฺชติ น วีริยํ อารภติ, อิทํ จตุตฺถํ กุสีตวตฺถุ.}

\prose{ปุน จปรํ อาวุโส ภิกฺขุ คามํ วา นิคมํ วา ปิณฺฑาย จรนฺโต น ลภติ ลูขสฺส วา ปณีตสฺส วา โภชนสฺส ยาวทตฺถํ ปาริปูริํ, ตสฺส เอวํ โหติ “อหํ โข คามํ วา นิคมํ วา ปิณฺฑาย จรนฺโต นาลตฺถํ ลูขสฺส วา ปณีตสฺส วา โภชนสฺส ยาวทตฺถํ ปาริปูริํ, ตสฺส เม กาโย กิลนฺโต อกมฺมญฺโญ, หนฺทาหํ นิปชฺชามี”ติ.\csromanspacer{} โส นิปชฺชติ น วีริยํ อารภติ, อิทํ ปญฺจมํ กุสีตวตฺถุ.}

\prose{\csromanfolio{212}ปุน จปรํ อาวุโส ภิกฺขุ คามํ วา นิคมํ วา ปิณฺฑาย จรนฺโต ลภติ ลูขสฺส วา ปณีตสฺส วา โภชนสฺส ยาวทตฺถํ ปาริปูริํ, ตสฺส เอวํ โหติ “อหํ โข คามํ วานิคมํ วา ปิณฺฑาย จรนฺโต อลตฺถํ ลูขสฺส วา ปณีตสฺส วา โภชนสฺส ยาวทตฺถํ ปาริปูริํ, ตสฺส เม กาโย ครุโก อกมฺมญฺโญ, มาสาจิตํ มญฺเญ, หนฺทาหํ นิปชฺชามี”ติ.\csromanspacer{} โส นิปชฺชติ น วีริยํ อารภติ, อิทํ ฉฏฺฐํ กุสีตวตฺถุ.}

\prose{ปุน จปรํ อาวุโส ภิกฺขุโน อุปฺปนฺโน โหติ อปฺปมตฺตโก อาพาโธ, ตสฺส เอวํ โหติ “อุปฺปนฺโน โข เม อยํ อปฺปมตฺตโก อาพาโธ, อตฺถิ กปฺโป นิปชฺชิตุํ, หนฺทาหํ นิปชฺชามี”ติ.\csromanspacer{} โส นิปชฺชติ น วีริยํ อารภติ, อิทํ สตฺตมํ กุสีตวตฺถุ.}

\prose{ปุน จปรํ อาวุโส ภิกฺขุ คิลานวุฏฺฐิโต\footnote{คิลานา วุฏฺฐิโต (สพฺพตฺถ) อฏฺฐกถา โอโลเกตพฺพา.} โหติ อจิรวุฏฺฐิโต เคลญฺญา, ตสฺส เอวํ โหติ “อหํ โข คิลานวุฏฺฐิโต อจิรวุฏฺฐิโต เคลญฺญา, ตสฺส เม กาโย ทุพฺพโล อกมฺมญฺโญ, หนฺทาหํ นิปชฺชามี”ติ.\csromanspacer{} โส นิปชฺชติ น วีริยํ อารภติ อปฺปตฺตสฺส ปตฺติยา อนธิคตสฺส อธิคมาย อสจฺฉิกตสฺส สจฺฉิกิริยาย, อิทํ อฏฺฐมํ กุสีตวตฺถุ. (4)}

\proseitem{335}{อฏฺฐ \textbf{อารมฺภวตฺถูนิ}.\csromanspacer{} อิธาวุโส ภิกฺขุนา กมฺมํ กาตพฺพํ โหติ, ตสฺส เอวํ โหติ “กมฺมํ โข เม กาตพฺพํ ภวิสฺสติ, กมฺมํ โข ปน เม กโรนฺเตน น สุกรํ พุทฺธานํ สาสนํ มนสิ กาตุํ, หนฺทาหํ วีริยํ อารภามิ อปฺปตฺตสฺส ปตฺติยา อนธิคตสฺส อธิคมาย อสจฺฉิกตสฺส สจฺฉิกิริยายา”ติ.\csromanspacer{} โส วีริยํ อารภติ อปฺปตฺตสฺส ปตฺติยา อนธิคตสฺส อธิคมาย อสจฺฉิกตสฺส สจฺฉิกิริยาย, อิทํ ปฐมํ อารมฺภวตฺถุ.}

\prose{ปุน จปรํ อาวุโส ภิกฺขุนา กมฺมํ กตํ โหติ, ตสฺส เอวํ โหติ “อหํ โข กมฺมํ อกาสิํ, กมฺมํ โข ปนาหํ กโรนฺโต นาสกฺขิํ พุทฺธานํ สาสนํ มนสิ กาตุํ, หนฺทาหํ วีริยํ อารภามี”ติ.\csromanspacer{} โส วีริยํ อารภติ ฯเปฯ อิทํ ทุติยํ อารมฺภวตฺถุ.}

\csromanheader{2}{10. สงฺคีติสุตฺต}
\prose{\csromanfolio{213}ปุน จปรํ อาวุโส ภิกฺขุนา มคฺโค คนฺตพฺโพ โหติ.\csromanspacer{} ตสฺส เอวํ โหติ “มคฺโค โข เม คนฺตพฺโพ ภวิสฺสติ, มคฺคํ โข ปน เม คจฺฉนฺเตน น สุกรํ พุทฺธานํ สาสนํ มนสิ กาตุํ.\csromanspacer{} หนฺทาหํ วีริยํ อารภามี”ติ.\csromanspacer{} โส วีริยํ อารภติ, อิทํ ตติยํ อารมฺภวตฺถุ.}

\prose{ปุน จปรํ อาวุโส ภิกฺขุนา มคฺโค คโต โหติ, ตสฺส เอวํ โหติ “อหํ โข มคฺคํ อคมาสิํ, มคฺคํ โข ปนาหํ คจฺฉนฺโต นาสกฺขิํ พุทฺธานํ สาสนํ มนสิ กาตุํ, หนฺทาหํ วีริยํ อารภามี”ติ.\csromanspacer{} โส วีริยํ อารภติ, อิทํ จตุตฺถํ อารมฺภวตฺถุ.}

\prose{ปุน จปรํ อาวุโส ภิกฺขุ คามํ วา นิคมํ วา ปิณฺฑาย จรนฺโต น ลภติ ลูขสฺส วา ปณีตสฺส วา โภชนสฺส ยาวทตฺถํ ปาริปูริํ, ตสฺส เอวํ โหติ “อหํ โข คามํ วา นิคมํ วา ปิณฺฑาย จรนฺโต นาลตฺถํ ลูขสฺส วา ปณีตสฺส วา โภชนสฺส ยาวทตฺถํ ปาริปูริํ, ตสฺส เม กาโย ลหุโก กมฺมญฺโญ, หนฺทาหํ วีริยํ อารภามี”ติ.\csromanspacer{} โส วีริยํ อารภติ, อิทํ ปญฺจมํ อารมฺภวตฺถุ.}

\prose{ปุน จปรํ อาวุโส ภิกฺขุ คามํ วา นิคมํ วา ปิณฺฑาย จรนฺโต ลภติ ลูขสฺส วา ปณีตสฺส วา โภชนสฺส ยาวทตฺถํ ปาริปูริํ, ตสฺส เอวํ โหติ “อหํ โข คามํ วา นิคมํ วา ปิณฺฑาย จรนฺโต อลตฺถํ ลูขสฺส วา ปณีตสฺส วา โภชนสฺส ยาวทตฺถํ ปาริปูริํ, ตสฺส เม กาโย พลวา กมฺมญฺโญ, หนฺทาหํ วีริยํ อารภามี”ติ.\csromanspacer{} โส วีริยํ อารภติ, อิทํ ฉฏฺฐํ อารมฺภวตฺถุ.}

\prose{ปุน จปรํ อาวุโส ภิกฺขุโน อุปฺปนฺโน โหติ อปฺปมตฺตโก อาพาโธ, ตสฺส เอวํ โหติ “อุปฺปนฺโน โข เม อยํ อปฺปมตฺตโก อาพาโธ, ฐานํ โข ปเนตํ วิชฺชติ ยํ เม อาพาโธ ปวฑฺเฒยฺย, หนฺทาหํ วีริยํ อารภามี”ติ.\csromanspacer{} โส วีริยํ อารภติ, อิทํ สตฺตมํ อารมฺภวตฺถุ.}

\prose{ปุน จปรํ อาวุโส ภิกฺขุ คิลานวุฏฺฐิโต โหติ อจิรวุฏฺฐิโต เคลญฺญา, ตสฺส เอวํ โหติ “อหํ โข คิลานวุฏฺฐิโต อจิรวุฏฺฐิโต เคลญฺญา, ฐานํ โข ปเนตํ วิชฺชติ ยํ เม อาพาโธ ปจฺจุทาวตฺเตยฺย, หนฺทาหํ วีริยํ อารภามิ อปฺปตฺตสฺส ปตฺติยา อนธิคตสฺส อธิคมาย อสจฺฉิกตสฺส สจฺฉิกิริยายา”ติ.\csromanspacer{} โส วีริยํ อารภติ อปฺปตฺตสฺส}

\prose{\csromanfolio{214}ปตฺติยา อนธิคตสฺส อธิคมาย อสจฺฉิกตสฺส สจฺฉิกิริยาย, อิทํ อฏฺฐมํ อารมฺภวตฺถุ. (5)}

\proseitem{336}{อฏฺฐ \textbf{ทานวตฺถูนิ}.\csromanspacer{} อาสชฺช ทานํ เทติ, ภยา ทานํ เทติ, “อทาสิ เม”ติ ทานํ เทติ, “ทสฺสติ เม”ติ ทานํ เทติ, “สาหุ ทานนฺ”ติ ทานํ เทติ, “อหํ ปจามิ, อิเม น ปจนฺติ, นารหามิ ปจนฺโต อปจนฺตานํ ทานํ น ทาตุนฺ”ติ ทานํ เทติ, “อิทํ เม ทานํ ททโต กลฺยาโณ กิตฺติสทฺโท อพฺภุคฺคจฺฉตี”ติ ทานํ เทติ, จิตฺตาลงฺการจิตฺตปริกฺขารตฺถํ ทานํ เทติ. (6)}

\proseitem{337}{อฏฺฐ \textbf{ทานูปปตฺติโย.}\csromanspacer{} อิธาวุโส เอกจฺโจ ทานํ เทติ สมณสฺส วา พฺราหฺมณสฺส วา อนฺนํ ปานํ วตฺถํ ยานํ มาลาคนฺธวิเลปนํ เสยฺยาวสถปทีเปยฺยํ.\csromanspacer{} โส ยํ เทติ ตํ ปจฺจาสีสติ\footnote{ปจฺจาสิํสติ (สี, สฺยา, กํ, อิ)}.\csromanspacer{} โส ปสฺสติ ขตฺติยมหาสาลํ วา พฺราหฺมณมหาสาลํ วา คหปติมหาสาลํ วา ปญฺจหิ กามคุเณหิ สมปฺปิตํ สมงฺคีภูตํ ปริจารยมานํ, ตสฺส เอวํ โหติ “อโห วตาหํ กายสฺส เภทา ปรํ มรณา ขตฺติยมหาสาลานํ วา พฺราหฺมณมหาสาลานํ วา คหปติมหาสาลานํ วา สหพฺยตํ อุปปชฺเชยฺยนฺ”ติ.\csromanspacer{} โส ตํ จิตฺตํ ทหติ, ตํ จิตฺตํ อธิฏฺฐาติ, ตํ จิตฺตํ ภาเวติ, ตสฺส ตํ จิตฺตํ หีเน วิมุตฺตํ อุตฺตริ อภาวิตํ ตตฺรูปปตฺติยา สํวตฺตติ.\csromanspacer{} ตญฺจ โข สีลวโต วทามิ โน ทุสฺสีลสฺส.\csromanspacer{} อิชฺฌตาวุโส สีลวโต เจโตปณิธิ วิสุทฺธตฺตา. (7-1)}

\prose{ปุน จปรํ อาวุโส อิเธกจฺโจ ทานํ เทติ สมณสฺส วา พฺราหฺมณสฺส วา อนฺนํ ปานํ ฯเปฯ เสยฺยาวสถปทีเปยฺยํ.\csromanspacer{} โส ยํ เทติ ตํ ปจฺจาสีสติ.\csromanspacer{} ตสฺส สุตํ โหติ “จาตุมหาราชิกา\footnote{จาตุมฺมหาราชิกา (สี, สฺยา, อิ)} เทวา ทีฆายุกา วณฺณวนฺโต สุขพหุลา”ติ.\csromanspacer{} ตสฺส เอวํ โหติ “อโห วตาหํ กายสฺส เภทา ปรํ มรณา จาตุมหาราชิกานํ เทวานํ สหพฺยตํ อุปปชฺเชยฺยนฺ”ติ.\csromanspacer{} โส ตํ จิตฺตํ ทหติ, ตํ จิตฺตํ อธิฏฺฐาติ, ตํ จิตฺตํ ภาเวติ, ตสฺส ตํ จิตฺตํ หีเน วิมุตฺตํ อุตฺตริ อภาวิตํ ตตฺรูปปตฺติยา สํวตฺตติ.\csromanspacer{} ตญฺจ โข สีลวโต วทามิ โน ทุสฺสีลสฺส.\csromanspacer{} อิชฺฌตาวุโส สีลวโต เจโตปณิธิ วิสุทฺธตฺตา. (7-2)}

\csromanheader{2}{10. สงฺคีติสุตฺต}
\prose{\csromanfolio{215}ปุน จปรํ อาวุโส อิเธกจฺโจ ทานํ เทติ สมณสฺส วา พฺราหฺมณสฺส วา อนฺนํ ปานํ ฯเปฯ เสยฺยาวสถปทีเปยฺยํ.\csromanspacer{} โส ยํ เทติ ตํ ปจฺจาสีสติ.\csromanspacer{} ตสฺส สุตํ โหติ “ตาวติํสา เทวา ฯเปฯ ยามา เทวา ฯเปฯ ตุสิตา เทวา ฯเปฯ นิมฺมานรตี เทวา ฯเปฯ ปรนิมฺมิตวสวตฺตี เทวา ทีฆายุกา วณฺณวนฺโต สุขพหุลา”ติ.\csromanspacer{} ตสฺส เอวํ โหติ “อโห วตาหํ กายสฺส เภทา ปรํ มรณา ปรนิมฺมิตวสวตฺตีนํ เทวานํ สหพฺยตํ อุปปชฺเชยฺยนฺ”ติ.\csromanspacer{} โส ตํ จิตฺตํ ทหติ, ตํ จิตฺตํ อธิฏฺฐาติ, ตํ จิตฺตํ ภาเวติ, ตสฺส ตํ จิตฺตํ หีเน วิมุตฺตํ อุตฺตริ อภาวิตํ ตตฺรูปปตฺติยา สํวตฺตติ.\csromanspacer{} ตญฺจ โข สีลวโต วทามิ โน ทุสฺสีลสฺส.\csromanspacer{} อิชฺฌตาวุโส สีลวโต เจโตปณิธิ วิสุทฺธตฺตา. (7-7)}

\prose{ปุน จปรํ อาวุโส อิเธกจฺโจ ทานํ เทติ สมณสฺส วา พฺราหฺมณสฺส วา อนฺนํ ปานํ วตฺถํ ยานํ มาลาคนฺธวิเลปนํ เสยฺยาวสถปทีเปยฺยํ.\csromanspacer{} โส ยํ เทติ ตํ ปจฺจาสีสติ.\csromanspacer{} ตสฺส สุตํ โหติ “พฺรหฺมกายิกา เทวา ทีฆายุกา วณฺณวนฺโต สุขพหุลา”ติ.\csromanspacer{} ตสฺส เอวํ โหติ “อโห วตาหํ กายสฺส เภทา ปรํ มรณา พฺรหฺมกายิกานํ เทวานํ สหพฺยตํ อุปปชฺเชยฺยนฺ”ติ.\csromanspacer{} โส ตํ จิตฺตํ ทหติ, ตํ จิตฺตํ อธิฏฺฐาติ, ตํ จิตฺตํ ภาเวติ, ตสฺส ตํ จิตฺตํ หีเน วิมุตฺตํ อุตฺตริ อภาวิตํ ตตฺรูปปตฺติยา สํวตฺตติ.\csromanspacer{} ตญฺจ โข สีลวโต วทามิ โน ทุสฺสีลสฺส, วีตราคสฺส โน สราคสฺส.\csromanspacer{} อิชฺฌตาวุโส สีลวโต เจโตปณิธิ วีตราคตฺตา. (7-8)}

\prose{อฏฺฐ \textbf{ปริสา}, ขตฺติย\textbf{ปริสา} พฺราหฺมณ\textbf{ปริสา} คหปติ\textbf{ปริสา} สมณ\textbf{ปริสา} จาตุมหาราชิก\textbf{ปริสา} ตาวติํส\textbf{ปริสา} มาร\textbf{ปริสา} พฺรหฺม\textbf{ปริสา}. (8)}

\prose{อฏฺฐ \textbf{โลกธมฺมา}, ลาโภ จ อลาโภ จ ยโส จ อยโส จ นินฺทา จ ปสํสา จ สุขญฺจ ทุกฺขญฺจ. (9)}

\proseitem{338}{อฏฺฐ \textbf{อภิภายตนานิ}.\csromanspacer{} อชฺฌตฺตํ รูปสญฺญี เอโก พหิทฺธา รูปานิ ปสฺสติ ปริตฺตานิ สุวณฺณทุพฺพณฺณานิ, “ตานิ อภิภุยฺย ชานามิ ปสฺสามี”ติ เอวํสญฺญี โหติ, อิทํ ปฐมํ อภิภายตนํ.}

\prose{\csromanfolio{216}อชฺฌตฺตํ รูปสญฺญี เอโก พหิทฺธา รูปานิ ปสฺสติ อปฺปมาณานิ สุวณฺณทุพฺพณฺณานิ, “ตานิ อภิภุยฺย ชานามิ ปสฺสามี”ติ เอวํสญฺญี โหติ, อิทํ ทุติยํ อภิภายตนํ.}

\prose{อชฺฌตฺตํ อรูปสญฺญี เอโก พหิทฺธา รูปานิ ปสฺสติ ปริตฺตานิ สุวณฺณทุพฺพณฺณานิ, “ตานิ อภิภุยฺย ชานามิ ปสฺสามี”ติ เอวํสญฺญี โหติ, อิทํ ตติยํ อภิภายตนํ.}

\prose{อชฺฌตฺตํ อรูปสญฺญี เอโก พหิทฺธา รูปานิ ปสฺสติ อปฺปมาณานิ สุวณฺณทุพฺพณฺณานิ, “ตานิ อภิภุยฺย ชานามิ ปสฺสามี”ติ เอวํ สญฺญี โหติ, อิทํ จตุตฺถํ อภิภายตนํ.}

\prose{อชฺฌตฺตํ อรูปสญฺญี เอโก พหิทฺธา รูปานิ ปสฺสติ นีลานิ นีลวณฺณานิ นีลนิทสฺสนานิ นีลนิภาสานิ.\csromanspacer{} เสยฺยถาปิ นาม อุมาปุปฺผํ นีลํ นีลวณฺณํ นีลนิทสฺสนํ นีลนิภาสํ, เสยฺยถา วา ปน ตํ วตฺถํ พาราณเสยฺยกํ อุภโตภาควิมฏฺฐํ นีลํ นีลวณฺณํ นีลนิทสฺสนํ นีลนิภาสํ, เอวเมว\footnote{เอวเมวํ (ก)} อชฺฌตฺตํ อรูปสญฺญี เอโก พหิทฺธา รูปานิ ปสฺสติ นีลานิ นีลวณฺณานิ นีลนิทสฺสนานิ นีลนิภาสานิ, “ตานิ อภิภุยฺย ชานามิ ปสฺสามี”ติ เอวํสญฺญี โหติ, อิทํ ปญฺจมํ อภิภายตนํ.}

\prose{อชฺฌตฺตํ อรูปสญฺญี เอโก พหิทฺธา รูปานิ ปสฺสติ ปีตานิ ปีตวณฺณานิ ปีตนิทสฺสนานิ ปีตนิภาสานิ.\csromanspacer{} เสยฺยถาปิ นาม กณิการปุปฺผํ\footnote{กณฺณิการปุปฺผํ (สฺยา, กํ)} ปีตํ ปีตวณฺณํ ปีตนิทสฺสนํ ปีตนิภาสํ, เสยฺยถา วา ปน ตํ วตฺถํ พาราณเสยฺยกํ อุภโตภาควิมฏฺฐํ ปีตํ ปีตวณฺณํ ปีตนิทสฺสนํ ปีตนิภาสํ, เอวเมว อชฺฌตฺตํ อรูปสญฺญี เอโก พหิทฺธา รูปานิ ปสฺสติ ปีตานิ ปีตวณฺณานิ ปีตนิทสฺสนานิ ปีตนิภาสานิ, “ตานิ อภิภุยฺย ชานามิ ปสฺสามี”ติ เอวํสญฺญี โหติ, อิทํ ฉฏฺฐํ อภิภายตนํ.}

\prose{อชฺฌตฺตํ อรูปสญฺญี เอโก พหิทฺธา รูปานิ ปสฺสติ โลหิตกานิ โลหิตกวณฺณานิ โลหิตกนิทสฺสนานิ โลหิตกนิภาสานิ.\csromanspacer{} เสยฺยถาปิ นาม พนฺธุชีวกปุปฺผํ โลหิตกํ โลหิตกวณฺณํ โลหิตกนิทสฺสนํ โลหิตกนิภาสํ, เสยฺยถา วา ปน ตํ วตฺถํ พาราณเสยฺยกํ อุภโตภาควิมฏฺฐํ โลหิตกํ โลหิตกวณฺณํ โลหิตกนิทสฺสนํ โลหิตกนิภาสํ, เอวเมว อชฺฌตฺตํ อรูปสญฺญี}

\vspace{\tipitakaparskip}
\csromancenter{สงฺคีติสุตฺต เอโก พหิทฺธา รูปานิ ปสฺสติ โลหิตกานิ โลหิตกวณฺณานิ โลหิตกนิทสฺสนานิ โลหิตกนิภาสานิ, “ตานิ อภิภุยฺย ชานามิ ปสฺสามี”ติ เอวํสญฺญี โหติ, อิทํ สตฺตมํ อภิภายตนํ.}
\prose{\csromanfolio{217}อชฺฌตฺตํ อรูปสญฺญี เอโก พหิทฺธา รูปานิ ปสฺสติ โอทาตานิ โอทาตวณฺณานิ โอทาตนิทสฺสนานิ โอทาตนิภาสานิ.\csromanspacer{} เสยฺยถาปิ นาม โอสธิตารกา โอทาตา โอทาตวณฺณา โอทาตนิทสฺสนา โอทาตนิภาสา, เสยฺยถา วา ปน ตํ วตฺถํ พาราณเสยฺยกํ อุภโตภาควิมฏฺฐํ โอทาตํ โอทาตวณฺณํ โอทาตนิทสฺสนํ โอทาตนิภาสํ, เอวเมว อชฺฌตฺตํ อรูปสญฺญี เอโก พหิทฺธา รูปานิ ปสฺสติ โอทาตานิ โอทาตวณฺณานิ โอทาตนิทสฺสนานิ โอทาตนิภาสานิ, “ตานิ อภิภุยฺย ชานามิ ปสฺสามี”ติ เอวํสญฺญี โหติ, อิทํ อฏฺฐมํ อภิภายตนํ. (10)}

\proseitem{339}{อฏฺฐ \textbf{วิโมกฺขา}.\csromanspacer{} รูปี รูปานิ ปสฺสติ, อยํ ปฐโม วิโมกฺโข.}

\prose{อชฺฌตฺตํ อรูปสญฺญี เอโก พหิทฺธา รูปานิ ปสฺสติ, อยํ ทุติโย วิโมกฺโข.}

\prose{สุภนฺเตว อธิมุตฺโต โหติ, อยํ ตติโย วิโมกฺโข.}

\prose{สพฺพโส รูปสญฺญานํ สมติกฺกมา ปฏิฆสญฺญานํ อตฺถงฺคมา นานตฺตสญฺญานํ อมนสิการา “อนนฺโต อากาโส”ติ อากาสานญฺจายตนํ อุปสมฺปชฺช วิหรติ, อยํ จตุตฺโถ วิโมกฺโข.}

\prose{สพฺพโส อากาสานญฺจายตนํ สมติกฺกมฺม “อนนฺตํ วิญฺญาณนฺ”ติ วิญฺญาณญฺจายตนํ อุปสมฺปชฺช วิหรติ, อยํ ปญฺจโม วิโมกฺโข.}

\prose{สพฺพโส วิญฺญาณญฺจายตนํ สมติกฺกมฺม “นตฺถิ กิญฺจี”ติ อากิญฺจญฺญายตนํ อุปสมฺปชฺช วิหรติ, อยํ ฉฏฺโฐ วิโมกฺโข.}

\prose{สพฺพโส อากิญฺจญฺญายตนํ สมติกฺกมฺม เนวสญฺญานาสญฺญายตนํ อุปสมฺปชฺช วิหรติ, อยํ สตฺตโม วิโมกฺโข.}

\prose{สพฺพโส เนวสญฺญานาสญฺญายตนํ สมติกฺกมฺม สญฺญาเวทยิตนิโรธํ อุปสมฺปชฺช วิหรติ, อยํ อฏฺฐโม วิโมกฺโข. (11)}

\prose{\csromanfolio{218}อิเม โข อาวุโส เตน ภควตา ชานตา ปสฺสตา อรหตา สมฺมาสมฺพุทฺเธน อฏฺฐ ธมฺมา สมฺมทกฺขาตา, ตตฺถ สพฺเพเหว สงฺคายิตพฺพํ ฯเปฯ อตฺถาย หิตาย สุขาย เทวมนุสฺสานํ.}

\csromanmatikaanchor{mtk.136}
\csromantocmarkref{section}{นวก}{mtk.136}
\bookmark[dest={mtk.136},level=1]{นวก}
\csromanheader{1}{\textbf{นวก}}
\proseitem{340}{อตฺถิ โข อาวุโส เตน ภควตา ชานตา ปสฺสตา อรหตา สมฺมาสมฺพุทฺเธน นว ธมฺมา สมฺมทกฺขาตา, ตตฺถ สพฺเพเหว สงฺคายิตพฺพํ ฯเปฯ อตฺถาย หิตาย สุขาย เทวมนุสฺสานํ.\csromanspacer{} กตเม นว.}

\prose{นว \textbf{อาฆาตวตฺถูนิ}. “อนตฺถํ เม อจรี”ติ อาฆาตํ พนฺธติ, “อนตฺถํ เม จรตี”ติ อาฆาตํ พนฺธติ, “อนตฺถํ เม จริสฺสตี”ติ อาฆาตํ พนฺธติ, “ปิยสฺส เม มนาปสฺส อนตฺถํ อจรี”ติ อาฆาตํ พนฺธติ ฯเปฯ อนตฺถํ จรตี”ติ อาฆาตํ พนฺธติ ฯเปฯ อนตฺถํ จริสฺสตี”ติ อาฆาตํ พนฺธติ, “อปฺปิยสฺส เม อมนาปสฺส อตฺถํ อจรี”ติ อาฆาตํ พนฺธติ ฯเปฯ อตฺถํ จรตี”ติ อาฆาตํ พนฺธติ ฯเปฯ อตฺถํ จริสฺสตี”ติ อาฆาตํ พนฺธติ. (1)}

\prose{นว \textbf{อาฆาตปฏิวินยา}. “อนตฺถํ เม อจริ\footnote{อจรีติ (สฺยา, ก) เอวํ “จรติ-จริสฺสติ” ปเทสุปิ.}, ตํ กุเตตฺถ ลพฺภา”ติ อาฆาตํ ปฏิวิเนติ, “อนตฺถํ เม จรติ, ตํ กุเตตฺถ ลพฺภา”ติ อาฆาตํ ปฏิวิเนติ, “อนตฺถํ เม จริสฺสติ, ตํ กุเตตฺถ ลพฺภา”ติ อาฆาตํ ปฏิวิเนติ, “ปิยสฺส เม มนาปสฺส อนตฺถํ อจริ ฯเปฯ อนตฺถํ จรติ ฯเปฯ อนตฺถํ จริสฺสติ, ตํ กุเตตฺถ ลพฺภา”ติ อาฆาตํ ปฏิวิเนติ, “อปฺปิยสฺส เม อมนาปสฺส อตฺถํ อจริ ฯเปฯ อตฺถํ จรติ ฯเปฯ อตฺถํ จริสฺสติ, ตํ กุเตตฺถ ลพฺภา”ติ อาฆาตํ ปฏิวิเนติ. (2)}

\proseitem{341}{นว \textbf{สตฺตาวาสา}.\csromanspacer{} สนฺตาวุโส สตฺตา นานตฺตกายา นานตฺตสญฺญิโน, เสยฺยถาปิ มนุสฺสา เอกจฺเจ จ เทวา เอกจฺเจ จ วินิปาติกา, อยํ ปฐโม สตฺตาวาโส.}

\prose{สนฺตาวุโส สตฺตา นานตฺตกายา เอกตฺตสญฺญิโน, เสยฺยถาปิ เทวา พฺรหฺมกายิกา ปฐมาภินิพฺพตฺตา, อยํ ทุติโย สตฺตาวาโส.}

\prose{สนฺตาวุโส สตฺตา เอกตฺตกายา นานตฺตสญฺญิโน, เสยฺยถาปิ เทวา อาภสฺสรา, อยํ ตติโย สตฺตาวาโส.}

\csromanheader{2}{10. สงฺคีติสุตฺต}
\prose{\csromanfolio{219}สนฺตาวุโส สตฺตา เอกตฺตกายา เอกตฺตสญฺญิโน, เสยฺยถาปิ เทวา สุภกิณฺหา, อยํ จตุตฺโถ สตฺตาวาโส.}

\prose{สนฺตาวุโส สตฺตา อสญฺญิโน อปฺปฏิสํเวทิโน, เสยฺยถาปิ เทวา อสญฺญสตฺตา\footnote{อสญฺญิสตฺตา (สฺยา, กํ)}, อยํ ปญฺจโม สตฺตาวาโส.}

\prose{สนฺตาวุโส สตฺตา สพฺพโส รูปสญฺญานํ สมติกฺกมา ปฏิฆสญฺญานํ อตฺถงฺคมา นานตฺตสญฺญานํ อมนสิการา “อนนฺโต อากาโส”ติ อากาสานญฺจายตนูปคา, อยํ ฉฏฺโฐ สตฺตาวาโส.}

\prose{สนฺตาวุโส สตฺตา สพฺพโส อากาสานญฺจายตนํ สมติกฺกมฺม “อนนฺตํ วิญฺญาณนฺ”ติ วิญฺญาณญฺจายตนูปคา, อยํ สตฺตโม สตฺตาวาโส.}

\prose{สนฺตาวุโส สตฺตา สพฺพโส วิญฺญาณญฺจายตนํ สมติกฺกมฺม “นตฺถิ กิญฺจี”ติ อากิญฺจญฺญายตนูปคา, อยํ อฏฺฐโม สตฺตาวาโส.}

\prose{สนฺตาวุโส สตฺตา สพฺพโส อากิญฺจญฺญายตนํ สมติกฺกมฺม\footnote{สมติกฺกมฺม สนฺตเมตํ ปณีตเมตนฺติ (สฺยา, กํ)} เนวสญฺญานาสญฺญายตนูปคา, อยํ นวโม สตฺตาวาโส. (3)}

\proseitem{342}{นว \textbf{อกฺขณา อสมยา} พฺรหฺมจริยวาสาย.\csromanspacer{} อิธาวุโส ตถาคโต จ โลเก อุปฺปนฺโน โหติ อรหํ สมฺมาสมฺพุทฺโธ, ธมฺโม จ เทสิยติ โอปสมิโก ปรินิพฺพานิโก สมฺโพธคามี สุคตปฺปเวทิโต, อยญฺจ ปุคฺคโล นิรยํ อุปปนฺโน โหติ, อยํ ปฐโม อกฺขโณ อสมโย พฺรหฺมจริยวาสาย.}

\prose{ปุน จปรํ อาวุโส ตถาคโต จ โลเก อุปฺปนฺโน โหติ อรหํ สมฺมาสมฺพุทฺโธ, ธมฺโม จ เทสิยติ โอปสมิโก ปรินิพฺพานิโก สมฺโพธคามี สุคตปฺปเวทิโต, อยญฺจ ปุคฺคโล ติรจฺฉานโยนิํ อุปปนฺโน โหติ, อยํ ทุติโย อกฺขโณ อสมโย พฺรหฺมจริยวาสาย.}

\prose{ปุน จปรํ ฯเปฯ เปตฺติวิสยํ อุปปนฺโน โหติ, อยํ ตติโย อกฺขโณ อสมโย พฺรหฺมจริยวาสาย.}

\prose{\csromanfolio{220}ปุน จปรํ ฯเปฯ อสุรกายํ อุปปนฺโน โหติ, อยํ จตุตฺโถ อกฺขโณ อสมโย พฺรหฺมจริยวาสาย.}

\prose{ปุน จปรํ ฯเปฯ อญฺญตรํ ทีฆายุกํ เทวนิกายํ อุปปนฺโน โหติ, อยํ ปญฺจโม อกฺขโณ อสมโย พฺรหฺมจริยวาสาย.}

\prose{ปุน จปรํ ฯเปฯ ปจฺจนฺติเมสุ ชนปเทสุ ปจฺจาชาโต โหติ มิลกฺเขสุ\footnote{มิลกฺขเกสุ (สฺยา, กํ), มิลกฺขูสุ (ก)} อวิญฺญาตาเรสุ, ยตฺถ นตฺถิ คติ ภิกฺขูนํ ภิกฺขุนีนํ อุปาสกานํ อุปาสิกานํ, อยํ ฉฏฺโฐ อกฺขโณ อสมโย พฺรหฺมจริยวาสาย.}

\prose{ปุน จปรํ ฯเปฯ มชฺฌิเมสุ ชนปเทสุ ปจฺจาชาโต โหติ, โส จ โหติ มิจฺฉาทิฏฺฐิโก วิปรีตทสฺสโน “นตฺถิ ทินฺนํ, นตฺถิ ยิฏฺฐํ, นตฺถิ หุตํ, นตฺถิ สุกตทุกฺกฏานํ\footnote{สุกฏ ทุกฺกฏานํ (สี, อิ)} กมฺมานํ ผลํ วิปาโก, นตฺถิ อยํ โลโก, นตฺถิ ปโร โลโก, นตฺถิ มาตา, นตฺถิ ปิตา, นตฺถิ สตฺตา โอปปาติกา, นตฺถิ โลเก สมณพฺราหฺมณา สมฺมคฺคตา สมฺมาปฏิปนฺนา เย อิมญฺจ โลกํ ปรญฺจ โลกํ สยํ อภิญฺญา สจฺฉิกตฺวา ปเวเทนฺตี”ติ, อยํ สตฺตโม อกฺขโณ อสมโย พฺรหฺมจริยวาสาย.}

\prose{ปุน จปรํ ฯเปฯ มชฺฌิเมสุ ชนปเทสุ ปจฺจาชาโต โหติ, โส จ โหติ ทุปฺปญฺโญ ชโฬ เอฬมูโค, นปฺปฏิพโล สุภาสิตทุพฺภาสิตานมตฺถมญฺญาตุํ, อยํ อฏฺฐโม อกฺขโณ อสมโย พฺรหฺมจริยวาสาย.}

\prose{ปุน จปรํ อาวุโส ตถาคโต จ โลเก น\footnote{กตฺถจิ นกาโร น ทิสฺสติ,} อุปฺปนฺโน โหติ อรหํ สมฺมาสมฺพุทฺโธ, ธมฺโม จ น เทสิยติ โอปสมิโก ปรินิพฺพานิโก สมฺโพธคามี สุคตปฺปเวทิโต, อยญฺจ ปุคฺคโล มชฺฌิเมสุ ชนปเทสุ ปจฺจาชาโต โหติ, โส จ โหติ ปญฺญวา อชโฬ อเนฬมูโค, ปฏิพโล สุภาสิตทุพฺภาสิตานมตฺถมญฺญาตุํ, อยํ นวโม อกฺขโณ อสมโย พฺรหฺมจริยวาสาย. (4)}

\proseitem{343}{นว \textbf{อนุปุพฺพวิหารา}.\csromanspacer{} อิธาวุโส ภิกฺขุ วิวิจฺเจว กาเมหิ วิวิจฺจ อกุสเลหิ ธมฺเมหิ สวิตกฺกํ สวิจารํ วิเวกชํ ปีติสุขํ ปฐมํ ฌานํ อุปสมฺปชฺช วิหรติ.\csromanspacer{} วิตกฺกวิจารานํ วูปสมา ฯเปฯ ทุติยํ ฌานํ อุปสมฺปชฺช วิหรติ.\csromanspacer{} ปีติยา จ วิราคา ฯเปฯ ตติยํ ฌานํ อุปสมฺปชฺช วิหรติ.\csromanspacer{} สุขสฺส จ}

\vspace{\tipitakaparskip}
\csromancenter{สงฺคีติสุตฺต ปหานา ฯเปฯ จตุตฺถํ ฌานํ อุปสมฺปชฺช วิหรติ.\csromanspacer{} สพฺพโส รูปสญฺญานํ สมติกฺกมา ฯเปฯ อากาสานญฺจายตนํ อุปสมฺปชฺช วิหรติ.\csromanspacer{} สพฺพโส อากาสานญฺจายตนํ สมติกฺกมฺม “อนนฺตํ วิญฺญาณนฺ”ติ วิญฺญาณญฺจายตนํ อุปสมฺปชฺช วิหรติ.\csromanspacer{} สพฺพโส วิญฺญาณญฺจายตนํ สมติกฺกมฺม “นตฺถิ กิญฺจี”ติ อากิญฺจญฺญายตนํ อุปสมฺปชฺช วิหรติ.\csromanspacer{} สพฺพโส อากิญฺจญฺญายตนํ สมติกฺกมฺม เนวสญฺญานาสญฺญายตนํ อุปสมฺปชฺช วิหรติ.\csromanspacer{} สพฺพโส เนวสญฺญานาสญฺญายตนํ สมติกฺกมฺม สญฺญาเวทยิตนิโรธํ อุปสมฺปชฺช วิหรติ. (5)}
\proseitem{344}{\csromanfolio{221}นว \textbf{อนุปุพฺพนิโรธา}.\csromanspacer{} ปฐมํ ฌานํ สมาปนฺนสฺส กามสญฺญา นิรุทฺธา โหติ.\csromanspacer{} ทุติยํ ฌานํ สมาปนฺนสฺส วิตกฺกวิจารา นิรุทฺธา โหนฺติ.\csromanspacer{} ตติยํ ฌานํ สมาปนฺนสฺส ปีติ นิรุทฺธา โหติ.\csromanspacer{} จตุตฺถํ ฌานํ สมาปนฺนสฺส อสฺสาสปสฺสาสา นิรุทฺธา โหนฺติ.\csromanspacer{} อากาสานญฺจายตนํ สมาปนฺนสฺส รูปสญฺญา นิรุทฺธา โหติ.\csromanspacer{} วิญฺญาณญฺจายตนํ สมาปนฺนสฺส อากาสานญฺจายตนสญฺญา นิรุทฺธา โหติ.\csromanspacer{} อากิญฺจญฺญายตนํ สมาปนฺนสฺส วิญฺญาณญฺจายตนสญฺญา นิรุทฺธา โหติ.\csromanspacer{} เนวสญฺญานาสญฺญายตนํ สมาปนฺนสฺส อากิญฺจญฺญายตนสญฺญา นิรุทฺธา โหติ.\csromanspacer{} สญฺญาเวทยิตนิโรธํ สมาปนฺนสฺส สญฺญา จ เวทนา จ นิรุทฺธา โหนฺติ. (6)}

\prose{อิเม โข อาวุโส เตน ภควตา ชานตา ปสฺสตา อรหตา สมฺมาสมฺพุทฺเธน นว ธมฺมา สมฺมทกฺขาตา.\csromanspacer{} ตตฺถ สพฺเพเหว สงฺคายิตพฺพํ ฯเปฯ อตฺถาย หิตาย สุขาย เทวมนุสฺสานํ.}

\csromanmatikaanchor{mtk.137}
\csromantocmarkref{section}{ทสก}{mtk.137}
\bookmark[dest={mtk.137},level=1]{ทสก}
\csromanheader{1}{\textbf{ทสก}}
\proseitem{345}{อตฺถิ โข อาวุโส เตน ภควตา ชานตา ปสฺสตา อรหตา สมฺมาสมฺพุทฺเธน ทส ธมฺมา สมฺมทกฺขาตา.\csromanspacer{} ตตฺถ สพฺเพเหว สงฺคายิตพฺพํ ฯเปฯ อตฺถาย หิตาย สุขาย เทวมนุสฺสานํ.\csromanspacer{} กตเม ทส.}

\prose{ทส \textbf{นาถกรณา} ธมฺมา.\csromanspacer{} อิธาวุโส ภิกฺขุ สีลวา โหติ, ปาติโมกฺขสํวรสํวุโต วิหรติ อาจารโคจรสมฺปนฺโน, อณุมตฺเตสุ วชฺเชสุ ภยทสฺสาวี สมาทาย สิกฺขติ สิกฺขาปเทสุ.\csromanspacer{} ยํปาวุโส \csromanfolio{222}ภิกฺขุ สีลวา โหติ, ปาติโมกฺขสํวรสํวุโส วิหรติ อาจารโคจรสมฺปนฺโน, อณุมตฺเตสุ วชฺเชสุ ภยทสฺสาวี สมาทาย สิกฺขติ สิกฺขาปเทสุ.\csromanspacer{} อยํปิ ธมฺโม นาถกรโณ. (\footnote{สาตฺถํ สพฺยญฺชนํ (สี, สฺยา, อิ)}-1)}

\prose{ปุน จปรํ อาวุโส ภิกฺขุ พหุสฺสุโต โหติ สุตธโร สุตสนฺนิจโย, เย เต ธมฺมา อาทิกลฺยาณา มชฺเฌกลฺยาณา ปริโยสานกลฺยาณา สาตฺถา สพฺยญฺชนา1 เกวลปริปุณฺณํ ปริสุทฺธํ พฺรหฺมจริยํ อภิวทนฺติ, ตถารูปาสฺส ธมฺมา พหุสฺสุตา โหนฺติ ธาตา\footnote{ธตา (ก-สี, สฺยา, กํ)} วจสา ปริจิตา มนสานุเปกฺขิตา ทิฏฺฐิยา สุปฺปฏิวิทฺธา.\csromanspacer{} ยํปาวุโส ภิกฺขุ พหุสฺสุโต โหติ ฯเปฯ ทิฏฺฐิยา สุปฺปฏิวิทฺธา.\csromanspacer{} อยํปิ ธมฺโม นาถกรโณ. (1-2)}

\prose{ปุน จปรํ อาวุโส ภิกฺขุ กลฺยาณมิตฺโต โหติ กลฺยาณสหาโย กลฺยาณสมฺปวงฺโก.\csromanspacer{} ยํปาวุโส ภิกฺขุ กลฺยาณมิตฺโต โหติ กลฺยาณสหาโย กลฺยาณสมฺปวงฺโก.\csromanspacer{} อยํปิ ธมฺโม นาถกรโณ. (1-3)}

\prose{ปุน จปรํ อาวุโส ภิกฺขุ สุวโจ โหติ โสวจสฺสกรเณหิ ธมฺเมหิ สมนฺนาคโต ขโม ปทกฺขิณคฺคาหี อนุสาสนิํ.\csromanspacer{} ยํปาวุโส ภิกฺขุ สุวโจ โหติ ฯเปฯ ปทกฺขิณคฺคาหี อนุสาสนิํ.\csromanspacer{} อยํปิ ธมฺโม นาถกรโณ. (1-4)}

\prose{ปุน จปรํ อาวุโส ภิกฺขุ ยานิ ตานิ สพฺรหฺมจารีนํ อุจฺจาวจานิ กิํกรณียานิ, ตตฺถ ทกฺโข โหติ อนลโส ตตฺรุปายาย วีมํสาย สมนฺนาคโต, อลํ กาตุํ อลํ สํวิธาตุํ.\csromanspacer{} ยํปาวุโส ภิกฺขุ ยานิ ตานิ สพฺรหฺมจารีนํ ฯเปฯ อลํ สํวิธาตุํ.\csromanspacer{} อยํปิ ธมฺโม นาถกรโณ. (1-5)}

\prose{ปุน จปรํ อาวุโส ภิกฺขุ ธมฺมกาโม โหติ ปิยสมุทาหาโร, อภิธมฺเม อภิวินเย อุฬารปาโมชฺโช\footnote{อุฬารปามุชฺโช (สี, อิ), โอฬารปาโมชฺโช (สฺยา, กํ)}.\csromanspacer{} ยํปาวุโส ภิกฺขุ ธมฺมกาโม โหติ ฯเปฯ อุฬารปาโมชฺโช3.\csromanspacer{} อยํปิ ธมฺโม นาถกรโณ. (1-6)}

\csromanheader{2}{10. สงฺคีติสุตฺต}
\prose{\csromanfolio{223}ปุน จปรํ อาวุโส ภิกฺขุ สนฺตุฏฺโฐ โหติ อิตรีตเรหิ จีวรปิณฺฑปาตเสนาสนคิลานปฺปจฺจยเภสชฺชปริกฺขาเรหิ.\csromanspacer{} ยํปาวุโส ภิกฺขุ สนฺตุฏฺโฐ โหติ ฯเปฯ ปริกฺขาเรหิ.\csromanspacer{} อยํปิ ธมฺโม นาถกรโณ. (1-7)}

\prose{ปุน จปรํ อาวุโส ภิกฺขุ อารทฺธวีริโย วิหรติ อกุสลานํ ธมฺมานํ ปหานาย กุสลานํ ธมฺมานํ อุปสมฺปทาย, ถามวา ทฬฺหปรกฺกโม อนิกฺขิตฺตธุโร กุสเลสุ ธมฺเมสุ.\csromanspacer{} ยํปาวุโส ภิกฺขุ อารทฺธวีริโย วิหรติ ฯเปฯ อนิกฺขิตฺตธุโร กุสเลสุ ธมฺเมสุ.\csromanspacer{} อยํปิ ธมฺโม นาถกรโณ. (1-8)}

\prose{ปุน จปรํ อาวุโส ภิกฺขุ สติมา โหติ ปรเมน สติเนปกฺเกน สมนฺนาคโต จิรกตํปิ จิรภาสิตํปิ สริตา อนุสฺสริตา.\csromanspacer{} ยํปาวุโส ภิกฺขุ สติมา โหติ ฯเปฯ สริตา อนุสฺสริตา.\csromanspacer{} อยํปิ ธมฺโม นาถกรโณ. (1-9)}

\prose{ปุน จปรํ อาวุโส ภิกฺขุ ปญฺญวา โหติ อุทยตฺถคามินิยา ปญฺญาย สมนฺนาคโต อริยาย นิพฺเพธิกาย สมฺมาทุกฺขกฺขยคามินิยา.\csromanspacer{} ยํปาวุโส ภิกฺขุ ปญฺญวา โหติ ฯเปฯ สมฺมาทุกฺขกฺขยคามินิยา.\csromanspacer{} อยํปิ ธมฺโม นาถกรโณ. (1-10)}

\proseitem{346}{ทส \textbf{กสิณายตนานิ}.\csromanspacer{} ปถวีกสิณเมโก สญฺชานาติ, อุทฺธํ อโธ ติริยํ อทฺวยํ อปฺปมาณํ.\csromanspacer{} อาโปกสิณเมโก สญฺชานาติ ฯเปฯ.\csromanspacer{} เตโชกสิณเมโก สญฺชานาติ.\csromanspacer{} วาโยกสิณเมโก สญฺชานาติ.\csromanspacer{} นีลกสิณเมโก สญฺชานาติ.\csromanspacer{} ปีตกสิณเมโก สญฺชานาติ.\csromanspacer{} โลหิตกสิณเมโก สญฺชานาติ.\csromanspacer{} โอทาตกสิณเมโก สญฺชานาติ.\csromanspacer{} อากาสกสิณเมโก สญฺชานาติ.\csromanspacer{} วิญฺญาณกสิณเมโก สญฺชานาติ, อุทฺธํ อโธ ติริยํ อทฺวยํ อปฺปมาณํ. (2)}

\proseitem{347}{ทส \textbf{อกุสลกมฺมปถา}, ปาณาติปาโต อทินฺนาทานํ กาเมสุมิจฺฉาจาโร มุสาวาโท ปิสุณา วาจา ผรุสา วาจา สมฺผปฺปลาโป อภิชฺฌา พฺยาปาโท มิจฺฉาทิฏฺฐิ. (3)}

\prose{\csromanfolio{224}ทส \textbf{กุสลกมฺมปถา}, ปาณาติปาตา เวรมณี อทินฺนาทานา เวรมณี กาเมสุมิจฺฉาจารา เวรมณี มุสาวาทา เวรมณี ปิสุณาย วาจาย เวรมณี ผรุสาย วาจาย เวรมณี สมฺผปฺปลาปา เวรมณี อนภิชฺฌา อพฺยาปาโท สมฺมาทิฏฺฐิ. (4)}

\proseitem{348}{ทส \textbf{อริยวาสา.}\csromanspacer{}\textbf{ อิ}ธาวุโส ภิกฺขุ ปญฺจงฺควิปฺปหีโน โหติ, ฉฬงฺคสมนฺนาคโต, เอการกฺโข จตุราปสฺเสโน, ปณุนฺนปจฺเจกสจฺโจ, สมวยสฏฺเฐสโน, อนาวิลสงฺกปฺโป, ปสฺสทฺธกายสงฺขาโร, สุวิมุตฺตจิตฺโต, สุวิมุตฺตปญฺโญ.}

\prose{กถญฺจาวุโส ภิกฺขุ ปญฺจงฺควิปฺปหีโน โหติ.\csromanspacer{} อิธาวุโส ภิกฺขุโน กามจฺฉนฺโท ปหีโน โหติ, พฺยาปาโท ปหีโน โหติ, ถินมิทฺธํ ปหีนํ โหติ, อุทฺธจฺจกุกฺกุจฺจํ ปหีนํ โหติ, วิจิกิจฺฉา ปหีนา โหติ.\csromanspacer{} เอวํ โข อาวุโส ภิกฺขุ ปญฺจงฺควิปฺปหีโน โหติ. (5-1)}

\prose{กถญฺจาวุโส ภิกฺขุ ฉฬงฺคสมนฺนาคโต โหติ.\csromanspacer{} อิธาวุโส ภิกฺขุ จกฺขุนา รูปํ ทิสฺวา เนว สุมโน โหติ น ทุมฺมโน, อุเปกฺขโก วิหรติ สโต สมฺปชาโน.\csromanspacer{} โสเตน สทฺทํ สุตฺวา ฯเปฯ.\csromanspacer{} มนสา ธมฺมํ วิญฺญาย เนว สุมโน โหติ น ทุมฺมโน, อุเปกฺขโก วิหรติ สโต สมฺปชาโน.\csromanspacer{} เอวํ โข อาวุโส ภิกฺขุ ฉฬงฺคสมนฺนาคโต โหติ. (5-2)}

\prose{กถญฺจาวุโส ภิกฺขุ เอการกฺโข โหติ.\csromanspacer{} อิธาวุโส ภิกฺขุ สตารกฺเขน เจตสา สมนฺนาคโต โหติ.\csromanspacer{} เอวํ โข อาวุโส ภิกฺขุ เอการกฺโข โหติ. (5-3)}

\prose{กถญฺจาวุโส ภิกฺขุ จตุราปสฺเสโน โหติ.\csromanspacer{} อิธาวุโส ภิกฺขุ สงฺขาเยกํ ปฏิเสวติ, สงฺขาเยกํ อธิวาเสติ, สงฺขาเยกํ วิโนเทติ, สงฺขาเยกํ ปริวชฺเชติ.\csromanspacer{} เอวํ โข อาวุโส ภิกฺขุ จตุราปสฺเสโน โหติ. (5-4)}

\prose{กถญฺจาวุโส ภิกฺขุ ปณุนฺนปจฺเจกสจฺโจ โหติ.\csromanspacer{} อิธาวุโส ภิกฺขุโน ยานิ ตานิ ปุถุสมณพฺราหฺมณานํ ปุถุปจฺเจกสจฺจานิ, สพฺพานิ ตานิ นุนฺนานิ โหนฺติ ปณุนฺนานิ จตฺตานิ วนฺตานิ มุตฺตานิ ปหีนานิ ปฏินิสฺสฏฺฐานิ.\csromanspacer{} เอวํ โข อาวุโส ภิกฺขุ ปณุนฺนปจฺเจกสจฺโจ โหติ. (5-5)}

\csromanheader{2}{10. สงฺคีติสุตฺต}
\prose{\csromanfolio{225}กถญฺจาวุโส ภิกฺขุ สมวยสฏฺเฐสโน โหติ.\csromanspacer{} อิธาวุโส ภิกฺขุโน กาเมสนา ปหีนา โหติ, ภเวสนา ปหีนา โหติ, พฺรหฺมจริเยสนา ปฏิปฺปสฺสทฺธา.\csromanspacer{} เอวํ โข อาวุโส ภิกฺขุ สมวยสฏฺเฐสโน โหติ. (5-6)}

\prose{กถญฺจาวุโส ภิกฺขุ อนาวิลสงฺกปฺโป โหติ.\csromanspacer{} อิธาวุโส ภิกฺขุโน กามสงฺกปฺโป ปหีโน โหติ, พฺยาปาทสงฺกปฺโป ปหีโน โหติ, วิหิํสาสงฺกปฺโป ปหีโน โหติ.\csromanspacer{} เอวํ โข อาวุโส ภิกฺขุ อนาวิลสงฺกปฺโป โหติ. (5-7)}

\prose{กถญฺจาวุโส ภิกฺขุ ปสฺสทฺธกายสงฺขาโร โหติ.\csromanspacer{} อิธาวุโส ภิกฺขุ สุขสฺส จ ปหานา ทุกฺขสฺส จ ปหานา ปุพฺเพว โสมนสฺสโทมนสฺสานํ อตฺถงฺคมา อทุกฺขมสุขํ อุเปกฺขาสติปาริสุทฺธิํ จตุตฺถํ ฌานํ อุปสมฺปชฺช วิหรติ.\csromanspacer{} เอวํ โข อาวุโส ภิกฺขุ ปสฺสทฺธกายสงฺขาโร โหติ. (5-8)}

\prose{กถญฺจาวุโส ภิกฺขุ สุวิมุตฺตจิตฺโต โหติ.\csromanspacer{} อิธาวุโส ภิกฺขุโน ราคา จิตฺตํ วิมุตฺตํ โหติ, โทสา จิตฺตํ วิมุตฺตํ โหติ, โมหา จิตฺตํ วิมุตฺตํ โหติ.\csromanspacer{} เอวํ โข อาวุโส ภิกฺขุ สุวิมุตฺตจิตฺโต โหติ. (5-9)}

\prose{กถญฺจาวุโส ภิกฺขุ สุวิมุตฺตปญฺโญ โหติ.\csromanspacer{} อิธาวุโส ภิกฺขุ “ราโค เม ปหีโน อุจฺฉินฺนมูโล ตาลาวตฺถุกโต อนภาวํกโต อายติํ อนุปฺปาทธมฺโม”ติ ปชานาติ. “โทโส เม ปหีโน อุจฺฉินฺนมูโล ตาลาวตฺถุกโต อนภาวํกโต อายติํ อนุปฺปาทธมฺโม”ติ ปชานาติ. “โมโห เม ปหีโน อุจฺฉินฺนมูโล ตาลาวตฺถุกโต อนภาวํกโต อายติํ อนุปฺปาทธมฺโม”ติ ปชานาติ.\csromanspacer{} เอวํ โข อาวุโส ภิกฺขุ สุวิมุตฺตปญฺโญ โหติ. (5-10)}

\prose{ทส \textbf{อเสกฺขา ธมฺมา.}\csromanspacer{} อเสกฺขา สมฺมาทิฏฺฐิ, อเสกฺโข สมฺมาสงฺกปฺโป, อเสกฺขา สมฺมาวาจา, อเสกฺโข สมฺมากมฺมนฺโต, อเสกฺโข สมฺมา-อาชีโว, อเสกฺโข สมฺมาวายาโม, อเสกฺขา สมฺมาสติ, อเสกฺโข สมฺมาสมาธิ, อเสกฺขํ สมฺมาญาณํ, อเสกฺขา สมฺมาวิมุตฺติ. (6)}

\prose{อิเม โข อาวุโส เตน ภควตา ชานตา ปสฺสตา อรหตา สมฺมาสมฺพุทฺเธน ทส ธมฺมา สมฺมทกฺขาตา, ตตฺถ สพฺเพเหว สงฺคายิตพฺพํ}

\prose{\csromanfolio{226}น วิวทิตพฺพํ, ยถยิทํ พฺรหฺมจริยํ อทฺธนิยํ อสฺส จิรฏฺฐิติกํ, ตทสฺส พหุชนหิตาย พหุชนสุขาย โลกานุกมฺปาย อตฺถาย หิตาย สุขาย เทวมนุสฺสานนฺติ.}

\proseitem{349}{อถ โข ภควา อุฏฺฐหิตฺวา อายสฺมนฺตํ สาริปุตฺตํ อามนฺเตสิ “สาธุ สาธุ สาริปุตฺต, สาธุ โข ตฺวํ สาริปุตฺต ภิกฺขูนํ สงฺคีติปริยายํ อภาสี”ติ.\csromanspacer{} อิทมโวจายสฺมา สาริปุตฺโต, สมนุญฺโญ สตฺถา อโหสิ.\csromanspacer{} อตฺตมนา จ เต ภิกฺขู อายสฺมโต สาริปุตฺตสฺส ภาสิตํ อภินนฺทุนฺติ.}

\nitthitam{\textbf{สงฺคีติสุตฺตํ นิฏฺฐิตํ ทสมํ.}}
\csromanmatikaanchor{mtk.138}
\csromantocmarknopage{chapter}{11. ทสุตฺตรสุตฺต}{mtk.138}
\bookmark[dest={mtk.138},level=0]{11. ทสุตฺตรสุตฺต}
\chapterheadpageread{11. ทสุตฺตรสุตฺต}
\csromanmatikaanchor{mtk.139}
\csromantocmarkref{section}{เอโกธมฺโม}{mtk.139}
\bookmark[dest={mtk.139},level=1]{เอโกธมฺโม}
\proseitem{350}{\csromanfolio{227}เอวํ เม สุตํ\csromandash{}เอกํ สมยํ ภควา จมฺปายํ วิหรติ คคฺคราย โปกฺขรณิยา ตีเร มหตา ภิกฺขุสํเฆน สทฺธิํ ปญฺจมตฺเตหิ ภิกฺขุสเตหิ.\csromanspacer{} ตตฺร โข อายสฺมา สาริปุตฺโต ภิกฺขู อามนฺเตสิ “อาวุโส ภิกฺขเว”ติ. “อาวุโส”ติ โข เต ภิกฺขู อายสฺมโต สาริปุตฺตสฺส ปจฺจสฺโสสุํ.\csromanspacer{} อายสฺมา สาริปุตฺโต เอตทโวจ\csromandash{}}

\csromangathameasure{ทสุตฺตรํ ปวกฺขามิ, ธมฺมํ นิพฺพานปตฺติยา. \\ ทุกฺขสฺสนฺต กิริยาย, สพฺพคนฺถปฺปโมจนํ.}
\csromangathagroup{\csromangathastanza{ทสุตฺตรํ ปวกฺขามิ, ธมฺมํ นิพฺพานปตฺติยา. \\ ทุกฺขสฺสนฺต กิริยาย, สพฺพคนฺถปฺปโมจนํ.}}

\csromanheader{2}{\textbf{เอโก ธมฺโม}}
\proseitem{351}{เอโก อาวุโส ธมฺโม พหุกาโร, เอโก ธมฺโม \textbf{ภาเวตพฺโพ}, เอโก ธมฺโม \textbf{ปริญฺเญยฺโย}, เอโก ธมฺโม \textbf{ปหาตพฺโพ}, เอโก ธมฺโม \textbf{หานภาคิโย}, เอโก ธมฺโม \textbf{วิเสสภาคิโย}, เอโก ธมฺโม ทุปฺปฏิวิชฺโฌ, เอโก ธมฺโม อุปฺปาเทตพฺโพ, เอโก ธมฺโม อภิญฺเญยฺโย, เอโก ธมฺโม สจฺฉิกาตพฺโพ.}

\prose{(ก) กตโม เอโก ธมฺโม \textbf{พหุกาโร.}\csromanspacer{} อปฺปมาโท กุสเลสุ ธมฺเมสุ.\csromanspacer{} อยํ เอโก ธมฺโม \textbf{พหุกาโร.}}

\prose{(ข) กตโม เอโก ธมฺโม \textbf{ภาเวตพฺโพ}.\csromanspacer{} กายคตาสติ สาตสหคตา.\csromanspacer{} อยํ เอโก ธมฺโม \textbf{ภาเวตพฺโพ}.}

\prose{(ค) กตโม เอโก ธมฺโม \textbf{ปริญฺเญยฺโย}.\csromanspacer{} ผสฺโส สาสโว อุปาทานิโย.\csromanspacer{} อยํ เอโก ธมฺโม \textbf{ปริญฺเญยฺโย}.}

\prose{(ฆ) กตโม เอโก ธมฺโม \textbf{ปหาตพฺโพ}.\csromanspacer{} อสฺมิมาโน.\csromanspacer{} อยํ เอโก ธมฺโม \textbf{ปหาตพฺโพ}.}

\prose{(\^{}อ) กตโม เอโก ธมฺโม \textbf{หานภาคิโย}.\csromanspacer{} อโยนิโส มนสิกาโร.\csromanspacer{} อยํ เอโก ธมฺโม \textbf{หานภาคิโย}.}

\prose{(จ) กตโม เอโก ธมฺโม \textbf{วิเสสภาคิโย}.\csromanspacer{} โยนิโส มนสิกาโร.\csromanspacer{} อยํ เอโก ธมฺโม \textbf{วิเสสภาคิโย}.}

\csromanmatikaanchor{mtk.140}
\csromantocmarkref{section}{เทฺวธมฺมา}{mtk.140}
\bookmark[dest={mtk.140},level=1]{เทฺวธมฺมา}
\prose{\csromanfolio{228}(ฉ) กตโม เอโก ธมฺโม \textbf{ทุปฺปฏิวิชฺโฌ}.\csromanspacer{} อานนฺตริโก เจโตสมาธิ.\csromanspacer{} อยํ เอโก ธมฺโม \textbf{ทุปฺปฏิวิชฺโฌ}.}

\prose{(ช) กตโม เอโก ธมฺโม \textbf{อุปฺปาเทตพฺโพ}.\csromanspacer{} อกุปฺปํ ญาณํ.\csromanspacer{} อยํ เอโก ธมฺโม \textbf{อุปฺปาเทตพฺโพ}.}

\prose{(ฌ) กตโม เอโก ธมฺโม \textbf{อภิญฺเญยฺโย}.\csromanspacer{} สพฺเพ สตฺตา อาหารฏฺฐิติกา.\csromanspacer{} อยํ เอโก ธมฺโม \textbf{อภิญฺเญยฺโย}.}

\prose{(ญ) กตโม เอโก ธมฺโม \textbf{สจฺฉิกาตพฺโพ}.\csromanspacer{} อกุปฺปา เจโตวิมุตฺติ.\csromanspacer{} อยํ เอโก ธมฺโม \textbf{สจฺฉิกาตพฺโพ}.}

\prose{อิติ อิเม ทส ธมฺมา ภูตา ตจฺฉา ตถา อวิตถา อนญฺญถา สมฺมา ตถาคเตน อภิสมฺพุทฺธา.}

\csromanheader{2}{\textbf{เทฺว ธมฺมา}}
\vspace{\tipitakaparskip}
\csromancenter{\textbf{เทฺว ธมฺมา} \textbf{พหุการา}, เทฺว ธมฺมา \textbf{ภาเวตพฺพา}, เทฺว ธมฺมา \textbf{ปริญฺเญยฺยา}, เทฺว ธมฺมา \textbf{ปหาตพฺพา}, เทฺว ธมฺมา \textbf{หานภาคิยา}, เทฺว ธมฺมา \textbf{วิเสสภาคิยา}, เทฺว ธมฺมา ทุปฺปฏิวิชฺฌา, เทฺว ธมฺมา อุปฺปาเทตพฺพา, เทฺว ธมฺมา อภิญฺเญยฺยา, เทฺว ธมฺมา สจฺฉิกาตพฺพา.}
\prose{(ก) กตเม เทฺว ธมฺมา \textbf{พหุการา}.\csromanspacer{} สติ จ สมฺปชญฺญญฺจ.\csromanspacer{} อิเม เทฺว ธมฺมา \textbf{พหุการา}.}

\prose{(ข) กตเม เทฺว ธมฺมา \textbf{ภาเวตพฺพา}.\csromanspacer{} สมโถ จ วิปสฺสนา จ.\csromanspacer{} อิเม เทฺว ธมฺมา \textbf{ภาเวตพฺพา}.}

\prose{(ค) กตเม เทฺว ธมฺมา \textbf{ปริญฺเญยฺยา}.\csromanspacer{} นามญฺจ รูปญฺจ.\csromanspacer{} อิเม เทฺว ธมฺมา \textbf{ปริญฺเญยฺยา}.}

\prose{(ฆ) กตเม เทฺว ธมฺมา \textbf{ปหาตพฺพา}.\csromanspacer{} อวิชฺชา จ ภวตณฺหา จ.\csromanspacer{} อิเม เทฺว ธมฺมา \textbf{ปหาตพฺพา}.}

\prose{(\^{}อ) กตเม เทฺว ธมฺมา \textbf{หานภาคิยา}.\csromanspacer{} โทวจสฺสตา จ ปาปมิตฺตตา จ.\csromanspacer{} อิเม เทฺว ธมฺมา \textbf{หานภาคิยา}.}

\prose{(จ) กตเม เทฺว ธมฺมา \textbf{วิเสสภาคิยา}.\csromanspacer{} โสวจสฺสตา จ กลฺยาณมิตฺตตา จ.\csromanspacer{} อิเม เทฺว ธมฺมา \textbf{วิเสสภาคิยา}.}

\csromanheader{2}{11. ทสุตฺตรสุตฺต}
\csromanmatikaanchor{mtk.141}
\csromantocmarkref{section}{ตโยธมฺมา}{mtk.141}
\bookmark[dest={mtk.141},level=1]{ตโยธมฺมา}
\prose{\csromanfolio{229}(ฉ) กตเม เทฺว ธมฺมา \textbf{ทุปฺปฏิวิชฺฌา}.\csromanspacer{} โย จ เหตุ โย จ ปจฺจโย สตฺตานํ สํกิเลสาย, โย จ เหตุ โย จ ปจฺจโย สตฺตานํ วิสุทฺธิยา.\csromanspacer{} อิเม เทฺว ธมฺมา \textbf{ทุปฺปฏิวิชฺฌา}.}

\prose{(ช) กตเม เทฺว ธมฺมา \textbf{อุปฺปาเทตพฺพา}.\csromanspacer{} เทฺว ญาณานิ ขเย ญาณํ อนุปฺปาเท ญาณํ.\csromanspacer{} อิเม เทฺว ธมฺมา \textbf{อุปฺปาเทตพฺพา}.}

\prose{(ฌ) กตเม เทฺว ธมฺมา \textbf{อภิญฺเญยฺยา}.\csromanspacer{} เทฺว ธาตุโย สงฺขตา จ ธาตุ อสงฺขตา จ ธาตุ.\csromanspacer{} อิเม เทฺว ธมฺมา \textbf{อภิญฺเญยฺยา}.}

\prose{(ญ) กตเม เทฺว ธมฺมา \textbf{สจฺฉิกาตพฺพา}.\csromanspacer{} วิชฺชา จ วิมุตฺติ จ.\csromanspacer{} อิเม เทฺว ธมฺมา \textbf{สจฺฉิกาตพฺพา}.}

\prose{อิติ อิเม วีสติ ธมฺมา ภูตา ตจฺฉา ตถา อวิตถา อนญฺญถา สมฺมา ตถาคเตน อภิสมฺพุทฺธา.}

\csromanheader{2}{\textbf{ตโย ธมฺมา}}
\vspace{\tipitakaparskip}
\csromancenter{\textbf{ตโย ธมฺมา} \textbf{พหุการา}, ตโย ธมฺมา \textbf{ภาเวตพฺพา} ฯเปฯ ตโย ธมฺมา \textbf{สจฺฉิกาตพฺพา}.}
\prose{(ก) กตเม ตโย ธมฺมา \textbf{พหุการา}.\csromanspacer{} สปฺปุริสสํเสโว, สทฺธมฺมสฺสวนํ, ธมฺมานุธมฺมปฺปฏิปตฺติ.\csromanspacer{} อิเม ตโย ธมฺมา \textbf{พหุการา}.}

\prose{(ข) กตเม ตโย ธมฺมา \textbf{ภาเวตพฺพา}.\csromanspacer{} ตโย สมาธี, สวิตกฺโก สวิจาโร สมาธิ อวิตกฺโก วิจารมตฺโต สมาธิ อวิตกฺโก อวิจาโร สมาธิ.\csromanspacer{} อิเม ตโย ธมฺมา \textbf{ภาเวตพฺพา}.}

\prose{(ค) กตเม ตโย ธมฺมา \textbf{ปริญฺเญยฺยา}.\csromanspacer{} ติสฺโส เวทนา, สุขา เวทนา ทุกฺขา เวทนา อทุกฺขมสุขา เวทนา.\csromanspacer{} อิเม ตโย ธมฺมา \textbf{ปริญฺเญยฺยา}.}

\prose{(ฆ) กตเม ตโย ธมฺมา \textbf{ปหาตพฺพา}.\csromanspacer{} ติสฺโส ตณฺหา, กามตณฺหา ภวตณฺหา วิภวตณฺหา.\csromanspacer{} อิเม ตโย ธมฺมา \textbf{ปหาตพฺพา}.}

\prose{(\^{}อ) กตเม ตโย ธมฺมา \textbf{หานภาคิยา}.\csromanspacer{} ตีณิ อกุสลมูลานิ, โลโภ อกุสลมูลํ โทโส อกุสลมูลํ โมโห อกุสลมูลํ.\csromanspacer{} อิเม ตโย ธมฺมา \textbf{หานภาคิยา}.}

\csromanmatikaanchor{mtk.142}
\csromantocmarkref{section}{จตฺตาโรธมฺมา}{mtk.142}
\bookmark[dest={mtk.142},level=1]{จตฺตาโรธมฺมา}
\prose{\csromanfolio{230}(จ) กตเม ตโย ธมฺมา \textbf{วิเสสภาคิยา}.\csromanspacer{} ตีณิ กุสลมูลานิ, อโลโภ กุสลมูลํ อโทโส กุสลมูลํ อโมโห กุสลมูลํ.\csromanspacer{} อิเม ตโย ธมฺมา \textbf{วิเสสภาคิยา}.}

\prose{(ฉ) กตเม ตโย ธมฺมา \textbf{ทุปฺปฏิวิชฺฌา}.\csromanspacer{} ติสฺโส นิสฺสรณิยา ธาตุโย, กามานเมตํ นิสฺสรณํ ยทิทํ เนกฺขมฺมํ, รูปานเมตํ นิสฺสรณํ ยทิทํ อรูปํ, ยํ โข ปน กิญฺจิ ภูตํ สงฺขตํ ปฏิจฺจสมุปฺปนฺนํ, นิโรโธ ตสฺส นิสฺสรณํ.\csromanspacer{} อิเม ตโย ธมฺมา \textbf{ทุปฺปฏิวิชฺฌา}.}

\prose{(ช) กตเม ตโย ธมฺมา \textbf{อุปฺปาเทตพฺพา}.\csromanspacer{} ตีณิ ญาณานิ, อตีตํเส ญาณํ อนาคตํเส ญาณํ ปจฺจุปฺปนฺนํเส ญาณํ.\csromanspacer{} อิเม ตโย ธมฺมา \textbf{อุปฺปาเทตพฺพา}.}

\prose{(ฌ) กตเม ตโย ธมฺมา \textbf{อภิญฺเญยฺยา}.\csromanspacer{} ติสฺโส ธาตุโย, กามธาตุ รูปาธาตุ อรูปธาตุ.\csromanspacer{} อิเม ตโย ธมฺมา \textbf{อภิญฺเญยฺยา}.}

\prose{(ญ) กตเม ตโย ธมฺมา \textbf{สจฺฉิกาตพฺพา}.\csromanspacer{} ติสฺโส วิชฺชา, ปุพฺเพนิวาสานุสฺสติญาณํ วิชฺชา สตฺตานํ จุตูปปาเต ญาณํ วิชฺชา อาสวานํ ขเย ญาณํ วิชฺชา.\csromanspacer{} อิเม ตโย ธมฺมา \textbf{สจฺฉิกาตพฺพา}.}

\prose{อิติ อิเม ติํสธมฺมา ภูตา ตจฺฉา ตถา อวิตถา อนญฺญถา สมฺมา ตถาคเตน อภิสมฺพุทฺทา.}

\csromanheader{2}{\textbf{จตฺตาโร ธมฺมา}}
\vspace{\tipitakaparskip}
\csromancenter{\textbf{จตฺตาโร ธมฺมา} \textbf{พหุการา}, จตฺตาโร ธมฺมา \textbf{ภาเวตพฺพา} ฯเปฯ จตฺตาโร ธมฺมา \textbf{สจฺฉิกาตพฺพา}.}
\prose{(ก) กตเม จตฺตาโร ธมฺมา \textbf{พหุการา}.\csromanspacer{} จตฺตาริ จกฺกานิ, ปติรูปเทสวาโส สปฺปุริสูปนิสฺสโย\footnote{สปฺปุริสุปสฺสโย (สฺยา, กํ)} อตฺตสมฺมาปณิธิ ปุพฺเพ จ กตปุญฺญตา.\csromanspacer{} อิเม จตฺตาโร ธมฺมา \textbf{พหุการา}.}

\prose{(ข) กตเม จตฺตาโร ธมฺมา \textbf{ภาเวตพฺพา}.\csromanspacer{} จตฺตาโร สติปฏฺฐานา, อิธาวุโส ภิกฺขุ กาเย กายานุปสฺสี วิหรติ อาตาปี สมฺปชาโน สติมา วิเนยฺย โลเก อภิชฺฌาโทมนสฺสํ.\csromanspacer{} เวทนาสุ ฯเปฯ.\csromanspacer{} จิตฺเต ฯเปฯ.}

\vspace{\tipitakaparskip}
\csromancenter{ทสุตฺตรสุตฺต ธมฺเมสุ ธมฺมานุปสฺสี วิหรติ อาตาปี สมฺปชาโน สติมา วิเนยฺย โลเก อภิชฺฌาโทมนสฺสํ.\csromanspacer{} อิเม จตฺตาโร ธมฺมา ภาเวตพฺพา.}
\prose{\csromanfolio{231}(ค) กตเม จตฺตาโร ธมฺมา \textbf{ปริญฺเญยฺยา}.\csromanspacer{} จตฺตาโร อาหารา, กพฬีกาโร\footnote{กวฬีกาโร (สฺยา, กํ)} อาหาโร โอฬาริโก วา สุขุโม วา, ผสฺโส ทุติโย, มโนสญฺเจตนา ตติยา, วิญฺญาณํ จตุตฺถํ.\csromanspacer{} อิเม จตฺตาโร ธมฺมา \textbf{ปริญฺเญยฺยา}.}

\prose{(ฆ) กตเม จตฺตาโร ธมฺมา \textbf{ปหาตพฺพา}.\csromanspacer{} จตฺตาโร โอฆา, กาโมโฆ ภโวโฆ ทิฏฺโฐโฆ อวิชฺโชโฆ.\csromanspacer{} อิเม จตฺตาโร ธมฺมา \textbf{ปหาตพฺพา}.}

\prose{(\^{}อ) กตเม จตฺตาโร ธมฺมา \textbf{หานภาคิยา}.\csromanspacer{} จตฺตาโร โยคา, กามโยโค ภวโยโค ทิฏฺฐิโยโค อวิชฺชาโยโค.\csromanspacer{} อิเม จตฺตาโร ธมฺมา \textbf{หานภาคิยา}.}

\prose{(จ) กตเม จตฺตาโร ธมฺมา \textbf{วิเสสภาคิยา}.\csromanspacer{} จตฺตาโร วิสํโยคา, กามโยควิสํโยโค ภวโยควิสํโยโค ทิฏฺฐิโยควิสํโยโค อวิชฺชาโยควิสํโยโค.\csromanspacer{} อิเม จตฺตาโร ธมฺมา \textbf{วิเสสภาคิยา}.}

\prose{(ฉ) กตเม จตฺตาโร ธมฺมา \textbf{ทุปฺปฏิวิชฺฌา}.\csromanspacer{} จตฺตาโร สมาธี, หานภาคิโย สมาธิ ฐิติภาคิโย สมาธิ วิเสสภาคิโย สมาธิ นิพฺเพธภาคิโย สมาธิ.\csromanspacer{} อิเม จตฺตาโร ธมฺมา \textbf{ทุปฺปฏิวิชฺฌา}.}

\prose{(ช) กตเม จตฺตาโร ธมฺมา \textbf{อุปฺปาเทตพฺพา}.\csromanspacer{} จตฺตาริ ญาณานิ, ธมฺเม ญาณํ อนฺวเย ญาณํ ปริเย ญาณํ สมฺมุติยา ญาณํ.\csromanspacer{} อิเม จตฺตาโร ธมฺมา \textbf{อุปฺปาเทตพฺพา}.}

\prose{(ฌ) กตเม จตฺตาโร ธมฺมา \textbf{อภิญฺเญยฺยา}.\csromanspacer{} จตฺตาริ อริยสจฺจานิ, ทุกฺขํ อริยสจฺจํ ทุกฺขสมุทยํ\footnote{ทุกฺขสมุทโย (สฺยา, กํ)} อริยสจฺจํ ทุกฺขนิโรธํ\footnote{ทุกฺขนิโรโธ (สฺยา, กํ)} อริยสจฺจํ ทุกฺขนิโรธาคามินี ปฏิปทา อริยสจฺจํ.\csromanspacer{} อิเม จตฺตาโร ธมฺมา \textbf{อภิญฺเญยฺยา}.}

\prose{(ญ) กตเม จตฺตาโร ธมฺมา \textbf{สจฺฉิกาตพฺพา}.\csromanspacer{} จตฺตาริ สามญฺญผลานิ, โสตาปตฺติผลํ สกทาคามิผลํ อนาคามิผลํ อรหตฺตผลํ.\csromanspacer{} อิเม จตฺตาโร ธมฺมา \textbf{สจฺฉิกาตพฺพา}.}

\csromanmatikaanchor{mtk.143}
\csromantocmarkref{section}{ปญฺจธมฺมา}{mtk.143}
\bookmark[dest={mtk.143},level=1]{ปญฺจธมฺมา}
\prose{\csromanfolio{232}อิติ อิเม จตฺตารีสธมฺมา ภูตา ตจฺฉา ตถา อวิตถา อนญฺญถา สมฺมา ตถาคเตน อภิสมฺพุทฺธา.}

\csromanheader{2}{\textbf{ปญฺจ ธมฺมา}}
\vspace{\tipitakaparskip}
\csromancenter{\textbf{ปญฺจ ธมฺมา} \textbf{พหุการา} ฯเปฯ ปญฺจ ธมฺมา สจฺฉิกาตพฺพา.}
\prose{(ก) กตเม ปญฺจ ธมฺมา \textbf{พหุการา}.\csromanspacer{} ปญฺจ ปธานิยงฺคานิ, อิธาวุโส ภิกฺขุ สทฺโธ โหติ, สทฺทหติ ตถาคตสฺส โพธิํ “อิติปิ โส ภควา อรหํ สมฺมาสมฺพุทฺโธ วิชฺชาจรณสมฺปนฺโน สุคโต โลกวิทู อนุตฺตโร ปุริสทมฺมสารถิ สตฺถา เทวมนุสฺสานํ พุทฺโธ ภควา”ติ.\csromanspacer{} อปฺปาพาโธ โหติ อปฺปาตงฺโก สมเวปากินิยา คหณิยา สมนฺนาคโต นาติสีตาย นาจฺจุณฺหาย มชฺฌิมาย ปธานกฺขมาย.\csromanspacer{} อสโฐ โหติ อมายาวี ยถาภูตมตฺตานํ อาวีกตฺตา สตฺถริ วา วิญฺญูสุ วา สพฺรหฺมจารีสุ.\csromanspacer{} อารทฺธวีริโย วิหรติ อกุสลานํ ธมฺมานํ ปหานาย กุสลานํ ธมฺมานํ อุปสมฺปทาย, ถามวา ทฬฺหปรกฺกโม อนิกฺขิตฺตธุโร กุสเลสุ ธมฺเมสุ.\csromanspacer{} ปญฺญวา โหติ อุทยตฺถคามินิยา ปญฺญาย สมนฺนาคโต อริยาย นิพฺเพธิกาย สมฺมา ทุกฺขกฺขยคามินิยา.\csromanspacer{} อิเม ปญฺจ ธมฺมา \textbf{พหุการา}.}

\prose{(ข) กตเม ปญฺจ ธมฺมา \textbf{ภาเวตพฺพา}.\csromanspacer{} ปญฺจงฺคิโก สมฺมาสมาธิ, ปีติผรณตา สุขผรณตา เจโตผรณตา อาโลกผรณตา ปจฺจเวกฺขณนิมิตฺตํ\footnote{ปจฺจเวกฺขณานิมิตฺตํ (สฺยา, กํ)}.\csromanspacer{} อิเม ปญฺจ ธมฺมา \textbf{ภาเวตพฺพา}.}

\prose{(ค) กตเม ปญฺจ ธมฺมา \textbf{ปริญฺเญยฺยา}.\csromanspacer{} ปญฺจุปาทานกฺขนฺธา * รูปุปาทานกฺขนฺโธ เวทนุปาทานกฺขนฺโธ สญฺญุปาทานกฺขนฺโธ สงฺขารุปาทานกฺขนฺโธ วิญฺญาณุปาทานกฺขนฺโธ.\csromanspacer{} อิเม ปญฺจ ธมฺมา \textbf{ปริญฺเญยฺยา}.}

\prose{(ฆ) กตเม ปญฺจ ธมฺมา \textbf{ปหาตพฺพา}.\csromanspacer{} ปญฺจ นีวรณานิ, กามจฺฉนฺทนีวรณํ พฺยาปาทนีวรณํ ถินมิทฺธนีวรณํ อุทฺธจฺจกุกฺกุจฺจนีวรณํ วิจิกิจฺฉานีวรณํ.\csromanspacer{} อิเม ปญฺจ ธมฺมา \textbf{ปหาตพฺพา}.}

\prose{(\^{}อ) กตเม ปญฺจ ธมฺมา \textbf{หานภาคิยา}.\csromanspacer{} ปญฺจ เจโตขิลา, อิธาวุโส ภิกฺขุ สตฺถริ กงฺขติ วิจิกิจฺฉติ นาธิมุจฺจติ น สมฺปสีทติ, โย โส อาวุโส ภิกฺขุ สตฺถริ กงฺขติ วิจิกิจฺฉติ นาธิมุจฺจติ น สมฺปสีทติ,}

\vspace{\tipitakaparskip}
\csromancenter{ทสุตฺตรสุตฺต ตสฺส จิตฺตํ น นมติ อาตปฺปาย อนุโยคาย สาตจฺจาย ปธานาย, ยสฺส จิตฺตํ น นมติ อาตปฺปาย อนุโยคาย สาตจฺจาย ปธานาย.\csromanspacer{} อยํ ปฐโม เจโตขิโล.\csromanspacer{} ปุน จปรํ อาวุโส ภิกฺขุ ธมฺเม กงฺขติ วิจิกิจฺฉติ ฯเปฯ สํเฆ กงฺขติ วิจิกิจฺฉติ ฯเปฯ สิกฺขาย กงฺขติ วิจิกิจฺฉติ ฯเปฯ สพฺรหฺมจารีสุ กุปิโต โหติ อนตฺตมโน อาหตจิตฺโต ขิลชาโต, โย โส อาวุโส ภิกฺขุ สพฺรหฺมจารีสุ กุปิโต โหติ อนตฺตมโน อาหตจิตฺโต ขิลชาโต, ตสฺส จิตฺตํ น นมติ อาตปฺปาย อนุโยคาย สาตจฺจาย ปธานาย, ยสฺส จิตฺตํ น นมติ อาตปฺปาย อนุโยคาย สาตจฺจาย ปธานาย.\csromanspacer{} อยํ ปญฺจโม เจโตขิโล.\csromanspacer{} อิเม ปญฺจ ธมฺมา หานภาคิยา.}
\prose{\csromanfolio{233}(จ) กตเม ปญฺจ ธมฺมา \textbf{วิเสสภาคิยา}.\csromanspacer{} ปญฺจินฺทฺริยานิ, สทฺธินฺทฺริยํ วีริยินฺทฺริยํ สตินฺทฺริยํ สมาธินฺทฺริยํ ปญฺญินฺทฺริยํ.\csromanspacer{} อิเม ปญฺจ ธมฺมา \textbf{วิเสสภาคิยา}.}

\prose{(ฉ) กตเม ปญฺจ ธมฺมา \textbf{ทุปฺปฏิวิชฺฌา}.\csromanspacer{} ปญฺจ นิสฺสรณิยา ธาตุโย.\csromanspacer{} อิธาวุโส ภิกฺขุโน กาเม มนสิกโรโต กาเมสุ จิตฺตํ น ปกฺขนฺทติ น ปสีทติ น สนฺติฏฺฐติ น วิมุจฺจติ.\csromanspacer{} เนกฺขมฺมํ โข ปนสฺส มนสิกโรโต เนกฺขมฺเม จิตฺตํ ปกฺขนฺทติ ปสีทติ สนฺติฏฺฐติ วิมุจฺจติ.\csromanspacer{} ตสฺส ตํ จิตฺตํ สุคตํ สุภาวิตํ สุวุฏฺฐิตํ สุวิมุตฺตํ วิสํยุตฺตํ กาเมหิ.\csromanspacer{} เย จ กามปจฺจยา อุปฺปชฺชนฺติ อาสวา วิฆาตา ปริฬาหา, มุตฺโต โส เตหิ.\csromanspacer{} น โส ตํ เวทนํ เวเทติ.\csromanspacer{} อิทมกฺขาตํ กามานํ นิสฺสรณํ. (1)}

\prose{ปุน จปรํ อาวุโส ภิกฺขุโน พฺยาปาทํ มนสิกโรโต พฺยาปาเท จิตฺตํ น ปกฺขนฺทติ น ปสีทติ น สนฺติฏฺฐติ น วิมุจฺจติ.\csromanspacer{} อพฺยาปาทํ โข ปนสฺส มนสิกโรโต อพฺยาปาเท จิตฺตํ ปกฺขนฺทติ ปสีทติ สนฺติฏฺฐติ วิมุจฺจติ.\csromanspacer{} ตสฺส ตํ จิตฺตํ สุคตํ สุภาวิตํ สุวุฏฺฐิตํ สุวิมุตฺตํ วิสํยุตฺตํ พฺยาปาเทน.\csromanspacer{} เย จ พฺยาปาทปจฺจยา อุปฺปชฺชนฺติ อาสวา วิฆาตา ปริฬาหา, มุตฺโต โส เตหิ.\csromanspacer{} น โส ตํ เวทนํ เวเทติ.\csromanspacer{} อิทมกฺขาตํ พฺยาปาทสฺส นิสฺสรณํ. (2)}

\prose{ปุน จปรํ อาวุโส ภิกฺขุโน วิเหสํ มนสิกโรโต วิเหสาย จิตฺตํ น ปกฺขนฺทติ น ปสีทติ น สนฺติฏฺฐติ น วิมุจฺจติ.\csromanspacer{} อวิเหสํ โข ปนสฺส มนสิกโรโต อวิเหสาย จิตฺตํ ปกฺขนฺทติ ปสีทติ สนฺติฏฺฐติ \csromanfolio{234}วิมุจฺจติ.\csromanspacer{} ตสฺส ตํ จิตฺตํ สุคตํ สุภาวิตํ สุวุฏฺฐิตํ สุวิมุตฺตํ วิสํยุตฺตํ วิเหสาย.\csromanspacer{} เย จ วิเหสาปจฺจยา อุปฺปชฺชนฺติ อาสวา วิฆาตา ปริฬาหา, มุตฺโต โส เตหิ.\csromanspacer{} น โส ตํ เวทนํ เวเทติ.\csromanspacer{} อิทมกฺขาตํ วิเหสาย นิสฺสรณํ. (3)}

\prose{ปุน จปรํ อาวุโส ภิกฺขุโน รูเป มนสิกโรโต รูเปสุ จิตฺตํ น ปกฺขนฺทติ น ปสีทติ น สนฺติฏฺฐติ น วิมุจฺจติ.\csromanspacer{} อรูปํ โข ปนสฺส มนสิกโรโต อรูเป จิตฺตํ ปกฺขนฺทติ ปสีทติ สนฺติฏฺฐติ วิมุจฺจติ.\csromanspacer{} ตสฺส ตํ จิตฺตํ สุคตํ สุภาวิตํ สุวุฏฺฐิตํ สุวิมุตฺตํ วิสํยุตฺตํ รูเปหิ.\csromanspacer{} เย จ รูปปจฺจยา อุปฺปชฺชนฺติ อาสวา วิฆาตา ปริฬาหา, มุตฺโต โส เตหิ.\csromanspacer{} น โส ตํ เวทนํ เวเทติ.\csromanspacer{} อิทมกฺขาตํ รูปานํ นิสฺสรณํ. (4)}

\prose{ปุน จปรํ อาวุโส ภิกฺขุโน สกฺกายํ มนสิกโรโต สกฺกาเย จิตฺตํ น ปกฺขนฺทติ น ปสีทติ น สนฺติฏฺฐติ น วิมุจฺจติ.\csromanspacer{} สกฺกายนิโรธํ โข ปนสฺส มนสิกโรโต สกฺกายนิโรเธ จิตฺตํ ปกฺขนฺทติ ปสีทติ สนฺติฏฺฐติ วิมุจฺจติ.\csromanspacer{} ตสฺส ตํ จิตฺตํ สุคตํ สุภาวิตํ สุวุฏฺฐิตํ สุวิมุตฺตํ วิสํยุตฺตํ สกฺกาเยน.\csromanspacer{} เย จ สกฺกายปจฺจยา อุปฺปชฺชนฺติ อาสวา วิฆาตา ปริฬาหา, มุตฺโต โส เตหิ.\csromanspacer{} น โส ตํ เวทนํ เวเทติ.\csromanspacer{} อิทมกฺขาตํ สกฺกายสฺส นิสฺสรณํ.\csromanspacer{} อิเม ปญฺจ ธมฺม ทุปฺปฏิวิชฺฌา. (5)}

\prose{(ช) กตเม ปญฺจ ธมฺมา \textbf{อุปฺปาเทตพฺพา}.\csromanspacer{} ปญฺจ ญาณิโก สมฺมาสมาธิ. “อยํ สมาธิ ปจฺจุปฺปนฺนสุโข เจว อายติํ จ สุขวิปาโก”ติ ปจฺจตฺตํเยว ญาณํ อุปฺปชฺชติ. “อยํ สมาธิ อริโย นิรามิโส”ติ ปจฺจตฺตํเยว ญาณํ อุปฺปชฺชติ. “อยํ สมาธิ อกาปุริสเสวิโต”ติ ปจฺจตฺตํเยว ญาณํ อุปฺปชฺชติ. “อยํ สมาธิ สนฺโต ปณีโต ปฏิปฺปสฺสทฺธลทฺโธ เอโกทิภาวาธิคโต, น สสงฺขารนิคฺคยฺหวาริตคโต”ติ\footnote{น จ สสงฺขารนิคฺคยฺห วาริตวโตติ (สี, สฺยา, กํ, อิ), น สสงฺขารนิคฺคยฺหวาริวาวโต (ก), น สสงฺขารนิคฺคยฺหวาริยาธิคโต (?)} ปจฺจตฺตํเยว ญาณํ อุปฺปชฺชติ. “โส โข ปนาหํ อิมํ สมาธิํ สโตว สมาปชฺชามิ, สโต วุฏฺฐหามี”ติ ปจฺจตฺตํเยว ญาณํ อุปฺปชฺชติ.\csromanspacer{} อิเม ปญฺจ ธมฺมา \textbf{อุปฺปาเทตพฺพา}.}

\prose{(ฌ) กตเม ปญฺจ ธมฺมา \textbf{อภิญฺเญยฺยา}.\csromanspacer{} ปญฺจ วิมุตฺตายตนานิ.\csromanspacer{} อิธาวุโส ภิกฺขุโน สตฺถา ธมฺมํ เทเสติ, อญฺญตโร วา ครุฏฺฐานิโย สพฺรหฺมจารี.}

\vspace{\tipitakaparskip}
\csromancenter{ทสุตฺตรสุตฺต ยถา ยถาวุโส ภิกฺขุโน สตฺถา ธมฺมํ เทเสติ, อญฺญตโร วา ครุฏฺฐานิโย สพฺรหฺมจารี, ตถา ตถา โส\footnote{ภิกฺขุ (สฺยา, กํ)} ตสฺมิํ ธมฺเม อตฺถปฺปฏิสํเวที จ โหติ ธมฺมปฺปฏิสํเวที จ, ตสฺส อตฺถปฺปฏิสํเวทิโน ธมฺมปฺปฏิสํเวทิโน ปาโมชฺชํ ชายติ, ปมุทิตสฺส ปีติ ชายติ, ปีติมนสฺส กาโย ปสฺสมฺภติ, ปสฺสทฺธกาโย สุขํ เวเทติ, สุขิโน จิตฺตํ สมาธิยติ, อิทํ ปฐมํ วิมุตฺตายตนํ. (1)}
\prose{\csromanfolio{235}ปุน จปรํ อาวุโส ภิกฺขุโน น เหว โข สตฺถา ธมฺมํ เทเสติ, อญฺญตโร วา ครุฏฺฐานิโย สพฺรหฺมจารี, อปิ จ โข ยถาสุตํ ยถาปริยตฺตํ ธมฺมํ วิตฺถาเรน ปเรสํ เทเสติ.\csromanspacer{} ยถา ยถาวุโส ภิกฺขุ ยถาสุตํ ยถาปริยตฺตํ ธมฺมํ วิตฺถาเรน ปเรสํ เทเสติ.\csromanspacer{} ตถา ตถา โส ตสฺมิํ ธมฺเม อตฺถปฺปฏิสํเวที จ โหติ ธมฺมปฺปฏิสํเวที จ, ตสฺส อตฺถปฺปฏิสํเวทิโน ธมฺมปฺปฏิสํเวทิโน ปาโมชฺชํ ชายติ, ปมุทิตสฺส ปีติ ชายติ, ปีติมนสฺส กาโย ปสฺสมฺภติ, ปสฺสทฺธกาโย สุขํ เวเทติ, สุขิโน จิตฺตํ สมาธิยติ, อิทํ ทุติยํ วิมุตฺตายตนํ. (2)}

\prose{ปุน จปรํ อาวุโส ภิกฺขุโน น เหว โข สตฺถา ธมฺมํ เทเสติ, อญฺญตโร วา ครุฏฺฐานิโย สพฺรหฺมจารี, นาปิ ยถาสุตํ ยถาปริยตฺตํ ธมฺมํ วิตฺถาเรน ปเรสํ เทเสติ, อปิ จ โข ยถาสุตํ ยถาปริยตฺตํ ธมฺมํ วิตฺถาเรน สชฺฌายํ กโรติ.\csromanspacer{} ยถา ยถาวุโส ภิกฺขุ ยถาสุตํ ยถาปริยตฺตํ ธมฺมํ วิตฺถาเรน สชฺฌายํ กโรติ.\csromanspacer{} ตถา ตถา โส ตสฺมิํ ธมฺเม อตฺถปฺปฏิสํเวที จ โหติ ธมฺมปฺปฏิสํเวที จ, ตสฺส อตฺถปฺปฏิสํเวทิโน ธมฺมปฺปฏิสํเวทิโน ปาโมชฺชํ ชายติ, ปมุทิตสฺส ปีติ ชายติ, ปีติมนสฺส กาโย ปสฺสมฺภติ, ปสฺสทฺธกาโย สุขํ เวเทติ, สุขิโน จิตฺตํ สมาธิยติ, อิทํ ตติยํ วิมุตฺตายตนํ. (3)}

\csromanmatikaanchor{mtk.144}
\csromantocmarkref{section}{ฉธมฺมา}{mtk.144}
\bookmark[dest={mtk.144},level=1]{ฉธมฺมา}
\prose{ปุน จปรํ อาวุโส ภิกฺขุโน น เหว โข สตฺถา ธมฺมํ เทเสติ, อญฺญตโร วา ครุฏฺฐานิโย สพฺรหฺมจารี, นาปิ ยถาสุตํ ยถาปริยตฺตํ ธมฺมํ วิตฺถาเรน ปเรสํ เทเสติ, นาปิ ยถาสุตํ ยถาปริยตฺตํ ธมฺมํ วิตฺถาเรน สชฺฌายํ กโรติ, อปิ จ โข ยถาสุตํ ยถาปริยตฺตํ ธมฺมํ เจตสา อนุวิตกฺเกติ อนุวิจาเรติ มนสานุเปกฺขติ.\csromanspacer{} ยถา \csromanfolio{236}ยถาวุโส ภิกฺขุ ยถาสุตํ ยถาปริยตฺตํ ธมฺมํ เจตสา อนุวิตกฺเกติ อนุวิจาเรติ มนสานุเปกฺขติ.\csromanspacer{} ตถา ตถา โส ตสฺมิํ ธมฺเม อตฺถปฺปฏิสํเวที จ โหติ ธมฺมปฺปฏิสํเวที จ, ตสฺส อตฺถปฺปฏิสํเวทิโน ธมฺมปฺปฏิสํเวทิโน ปาโมชฺชํ ชายติ, ปมุทิตสฺส ปีติ ชายติ, ปีติมนสฺส กาโย ปสฺสมฺภติ, ปสฺสทฺธกาโย สุขํ เวเทติ, สุขิโน จิตฺตํ สมาธิยติ, อิทํ จตุตฺถํ วิมุตฺตายตนํ. (4)}

\prose{ปุน จปรํ อาวุโส ภิกฺขุโน น เหว โข สตฺถา ธมฺมํ เทเสติ, อญฺญตโร วา ครุฏฺฐานิโย สพฺรหฺมจารี, นาปิ ยถาสุตํ ยถาปริยตฺตํ ธมฺมํ วิตฺถาเรน ปเรสํ เทเสติ, นาปิ ยถาสุตํ ยถาปริยตฺตํ ธมฺมํ วิตฺถาเรน สชฺฌายํ กโรติ, นาปิ ยถาสุตํ ยถาปริยตฺตํ ธมฺมํ เจตสา อนุวิตกฺเกติ อนุวิจาเรติ มนสานุเปกฺขติ, อปิ จ ขฺวสฺส อญฺญตรํ สมาธินิมิตฺตํ สุคฺคหิตํ โหติ สุมนสิกตํ สูปธาริตํ สุปฺปฏิวิทฺธํ ปญฺญาย.\csromanspacer{} ยถา ยถาวุโส ภิกฺขุโน อญฺญตรํ สมาธินิมิตฺตํ สุคฺคหิตํ โหติ สุมนสิกตํ สูปธาริตํ สุปฺปฏิวิทฺธํ ปญฺญาย.\csromanspacer{} ตถา ตถา โส ตสฺมิํ ธมฺเม อตฺถปฺปฏิสํเวที จ โหติ ธมฺมปฺปฏิสํเวที จ, ตสฺส อตฺถปฺปฏิสํเวทิโน ธมฺมปฺปฏิสํเวทิโน ปาโมชฺชํ ชายติ, ปมุทิตสฺส ปีติ ชายติ, ปีติมนสฺส กาโย ปสฺสมฺภติ, ปสฺสทฺธกาโย สุขํ เวเทติ, สุขิโน จิตฺตํ สมาธิยติ, อิทํ ปญฺจมํ วิมุตฺตายตนํ.\csromanspacer{} อิเม ปญฺจ ธมฺมา อภิญฺเญยฺยา. (5)}

\prose{(ญ) กตเม ปญฺจ ธมฺมา \textbf{สจฺฉิกาตพฺพา}.\csromanspacer{} ปญฺจ ธมฺมกฺขนฺธา, สีลกฺขนฺโธ สมาธิกฺขนฺโธ ปญฺญากฺขนฺโธ วิมุตฺติกฺขนฺโธ วิมุตฺติญาณทสฺสนกฺขนฺโธ.\csromanspacer{} อิเม ปญฺจ ธมฺมา \textbf{สจฺฉิกาตพฺพา}.}

\prose{อิติ อิเม ปญฺญาสธมฺมา ภูตา ตจฺฉา ตถา อวิตถา อนญฺญถา สมฺมา ตถาคเตน อภิสมฺพุทฺธา.}

\csromanheader{2}{\textbf{ฉ ธมฺมา}}
\vspace{\tipitakaparskip}
\csromancenter{\textbf{ฉ ธมฺมา} พหุการา ฯเปฯ ฉ ธมฺมา \textbf{สจฺฉิกาตพฺพา}.}
\prose{(ก) กตเม ฉ ธมฺมา \textbf{พหุการา,} ฉ สารณียา ธมฺมา.\csromanspacer{} อิธาวุโส ภิกฺขุโน เมตฺตํ กายกมฺมํ ปจฺจุปฏฺฐิตํ โหติ สพฺรหฺมจารีสุ อาวิ เจว}

\vspace{\tipitakaparskip}
\csromancenter{ทสุตฺตรสุตฺต รโห จ, อยํปิ ธมฺโม สารณีโย ปิยกรโณ ครุกรโณ สงฺคหาย อวิวาทาย สามคฺคิยา เอกีภาวาย สํวตฺตติ. (1)}
\prose{\csromanfolio{237}ปุน จปรํ อาวุโส ภิกฺขุโน เมตฺตํ วจีกมฺมํ ฯเปฯ เอกีภาวาย สํวตฺตติ. (2)}

\prose{ปุน จปรํ อาวุโส ภิกฺขุโน เมตฺตํ มโนกมฺมํ ฯเปฯ เอกีภาวาย สํวตฺตติ. (3)}

\prose{ปุน จปรํ อาวุโส ภิกฺขุ เย เต ลาภา ธมฺมิกา ธมฺมลทฺธา อนฺตมโส ปตฺตปริยาปนฺนมตฺตมฺปิ, ตถารูเปหิ ลาเภหิ อปฺปฏิวิภตฺตโภคี โหติ สีลวนฺเตหิ สพฺรหฺมจารีหิ สาธารณโภคี, อยํปิ ธมฺโม สารณีโย ฯเปฯ เอกีภาวาย สํวตฺตติ. (4)}

\prose{ปุน จปรํ อาวุโส ภิกฺขุ ยานิ ตานิ สีลานิ อขณฺฑานิ อจฺฉิทฺทานิ อสพลานิ อกมฺมาสานิ ภุชิสฺสานิ วิญฺญุปฺปสตฺถานิ อปรามฏฺฐานิ สมาธิสํวตฺตนิกานิ, ตถารูเปสุ สีเลสุ สีลสามญฺญคโต วิหรติ สพฺรหฺมจารีหิ อาวิ เจว รโห จ, อยํปิ ธมฺโม สารณีโย ฯเปฯ เอกีภาวาย สํวตฺตติ. (5)}

\prose{ปุน จปรํ อาวุโส ภิกฺขุ ยายํ ทิฏฺฐิ อริยา นิยฺยานิกา นิยฺยาติ ตกฺกรสฺส สมฺมา ทุกฺขกฺขยาย, ตถารูปาย ทิฏฺฐิยา ทิฏฺฐิสามญฺญคโต วิหรติ สพฺรหฺมจารีหิ อาวิ เจว รโห จ, อยํปิ ธมฺโม สารณีโย ปิยกรโณ ครุกรโณ สงฺคหาย อวิวาทาย สามคฺคิยา เอกีภาวาย สํวตฺตติ.\csromanspacer{} อิเม ฉ ธมฺมา พหุการา. (6)}

\prose{(ข) กตเม ฉ ธมฺมา \textbf{ภาเวตพฺพา}.\csromanspacer{} ฉ อนุสฺสติฏฺฐานานิ, พุทฺธานุสฺสติ ธมฺมานุสฺสติ สํฆานุสฺสติ สีลานุสฺสติ จาคานุสฺสติ เทวตานุสฺสติ.\csromanspacer{} อิเม ฉ ธมฺมา \textbf{ภาเวตพฺพา}.}

\prose{(ค) กตเม ฉ ธมฺมา \textbf{ปริญฺเญยฺยา}.\csromanspacer{} ฉ อชฺฌตฺติกานิ อายตนานิ, จกฺขายตนํ โสตายตนํ ฆานายตนํ ชิวฺหายตนํ กายายตนํ มนายตนํ.\csromanspacer{} อิเม ฉ ธมฺมา \textbf{ปริญฺเญยฺยา}.}

\prose{(ฆ) กตเม ฉ ธมฺมา \textbf{ปหาตพฺพา}.\csromanspacer{} ฉ ตณฺหากายา, รูปตณฺหา สทฺทตณฺหา คนฺธตณฺหา รสตณฺหา โผฏฺฐพฺพตณฺหา ธมฺมตณฺหา.\csromanspacer{} อิเม ฉ ธมฺมา \textbf{ปหาตพฺพา}.}

\prose{\csromanfolio{238}(\^{}อ) กตเม ฉ ธมฺมา \textbf{หานภาคิยา}.\csromanspacer{} ฉ อคารวา.\csromanspacer{} อิธาวุโส ภิกฺขุ สตฺถริ อคารโว วิหรติ อปฺปติสฺโส.\csromanspacer{} ธมฺเม ฯเปฯ.\csromanspacer{} สํเฆ.\csromanspacer{} สิกฺขาย.\csromanspacer{} อปฺปมาเท.\csromanspacer{} ปติสนฺถาเร อคารโว วิหรติ อปฺปติสฺโส.\csromanspacer{} อิเม ฉ ธมฺมา \textbf{หานภาคิยา}.}

\prose{(จ) กตเม ฉ ธมฺมา \textbf{วิเสสภาคิยา}.\csromanspacer{} ฉ คารวา.\csromanspacer{} อิธาวุโส ภิกฺขุ สตฺถริ สคารโว วิหรติ สปฺปติสฺโส.\csromanspacer{} ธมฺเม ฯเปฯ.\csromanspacer{} สํเฆ.\csromanspacer{} สิกฺขาย.\csromanspacer{} อปฺปมาเท.\csromanspacer{} ปฏิสนฺถาเร สคารโว วิหรติ สปฺปติสฺโส.\csromanspacer{} อิเม ฉ ธมฺมา \textbf{วิเสสภาคิยา}.}

\prose{(ฉ) กตเม ฉ ธมฺมา \textbf{ทุปฺปฏิวิชฺฌา}.\csromanspacer{} ฉ นิสฺสรณิยา ธาตุโย.\csromanspacer{} อิธาวุโส ภิกฺขุ เอวํ วเทยฺย “เมตฺตา หิ โข เม อาวุโส เจโตวิมุตฺติ ภาวิตา พหุลีกตา ยานีกตา วตฺถุกตา อนุฏฺฐิตา ปริจิตา สุสมารทฺธา, อถ จ ปน เม พฺยาปาโท จิตฺตํ ปริยาทาย ติฏฺฐตี”ติ.\csromanspacer{} โส “มา เหวนฺ”ติสฺส วจนีโย, “มายสฺมา เอวํ อวจ, มา ภควนฺตํ อพฺภาจิกฺขิ, น หิ สาธุ ภควโต อพฺภกฺขานํ, น หิ ภควา เอวํ วเทยฺย.\csromanspacer{} อฏฺฐานเมตํ อาวุโส อนวกาโส ยํ เมตฺตาย เจโตวิมุตฺติยา ภาวิตาย พหุลีกตาย ยานีกตาย วตฺถุกตาย อนุฏฺฐิตาย ปริจิตาย สุสมารทฺธาย, อถ จ ปนสฺส พฺยาปาโท จิตฺตํ ปริยาทาย ฐสฺสติ, เนตํ ฐานํ วิชฺชติ.\csromanspacer{} นิสฺสรณเญฺหตํ อาวุโส พฺยาปาทสฺส, ยทิทํ เมตฺตาเจโตวิมุตฺตี”ติ. (1)}

\prose{อิธ ปนาวุโส ภิกฺขุ เอวํ วเทยฺย “กรุณา หิ โข เม เจโตวิมุตฺติ ภาวิตา พหุลีกตา ยานีกตา วตฺถุกตา อนุฏฺฐิตา ปริจิตา สุสมารทฺธา, อถ จ ปน เม วิเหสา จิตฺตํ ปริยาทาย ติฏฺฐตี”ติ.\csromanspacer{} โส “มา เหวนฺ”ติสฺส วจนีโย, “มายสฺมา เอวํ อวจ, มา ภควนฺตํ อพฺภาจิกฺขิ ฯเปฯ.\csromanspacer{} นิสฺสรณเญฺหตํ อาวุโส วิเหสาย, ยทิทํ กรุณาเจโตวิมุตฺตี”ติ. (2)}

\prose{อิธ ปนาวุโส ภิกฺขุ เอวํ วเทยฺย “มุทิตา หิ โข เม เจโตวิมุตฺติ ภาวิตา ฯเปฯ อถ จ ปน เม อรติ จิตฺตํ ปริยาทาย ติฏฺฐตี”ติ.\csromanspacer{} โส “มา เหวนฺ”ติสฺส วจนีโย, “มายสฺมา เอวํ อวจ ฯเปฯ นิสฺสรณเญฺหตํ อาวุโส อรติยา, ยทิทํ มุทิตาเจโตวิมุตฺตี”ติ. (3)}

\csromanheader{2}{11. ทสุตฺตรสุตฺต}
\prose{\csromanfolio{239}อิธ ปนาวุโส ภิกฺขุ เอวํ วเทยฺย “อุเปกฺขา หิ โข เม เจโตวิมุตฺติ ภาวิตา ฯเปฯ อถ จ ปน เม ราโค จิตฺตํ ปริยาทาย ติฏฺฐตี”ติ.\csromanspacer{} โส “มา เหวนฺ”ติสฺส วจนีโย “มายสฺมา เอวํ อวจ ฯเปฯ.\csromanspacer{} นิสฺสรณเญฺหตํ อาวุโส ราคสฺส, ยทิทํ อุเปกฺขาเจโตวิมุตฺตี”ติ. (4)}

\prose{อิธ ปนาวุโส ภิกฺขุ เอวํ วเทยฺย “อนิมิตฺตา หิ โข เม เจโตวิมุตฺติ ภาวิตา ฯเปฯ อถ จ ปน เม นิมิตฺตานุสาริ วิญฺญาณํ โหตี”ติ.\csromanspacer{} โส “มา เหวนฺ”ติสฺส วจนีโย “มายสฺมา เอวํ อวจ ฯเปฯ.\csromanspacer{} นิสฺสรณเญฺหตํ อาวุโส สพฺพนิมิตฺตานํ, ยทิทํ อนิมิตฺตา เจโตวิมุตฺตี”ติ. (5)}

\prose{อิธ ปนาวุโส ภิกฺขุ เอวํ วเทยฺย “อสฺมีติ โข เม วิคตํ, อยมหมสฺมีติ น สมนุปสฺสามิ, อถ จ ปน เม วิจิกิจฺฉากถํกถาสลฺลํ จิตฺตํ ปริยาทาย ติฏฺฐตี”ติ.\csromanspacer{} โส “มา เหวนฺ”ติสฺส วจนีโย, “มายสฺมา เอวํ อวจ, มา ภควนฺตํ อพฺภาจิกฺขิ, น หิ สาธุ ภควโต อพฺภกฺขานํ, น หิ ภควา เอวํ วเทยฺย.\csromanspacer{} อฏฺฐานเมตํ อาวุโส อนวกาโส ยํ อสฺมีติ วิคเต อยมหมสฺมีติ อสมนุปสฺสโต, อถ จ ปนสฺส วิจิกิจฺฉากถํกถาสลฺลํ จิตฺตํ ปริยาทาย ฐสฺสติ, เนตํ ฐานํ วิชฺชติ.\csromanspacer{} นิสฺสรณเญฺหตํ อาวุโส วิจิกิจฺฉากถํกถาสลฺลสฺส, ยทิทํ อสฺมิมานสมุคฺฆาโฏติ.\csromanspacer{} อิเม ฉ ธมฺมา ทุปฺปฏิวิชฺฌา. (6)}

\prose{(ช) กตเม ฉ ธมฺมา \textbf{อุปฺปาเทตพฺพา}.\csromanspacer{} ฉ สตตวิหารา.\csromanspacer{} อิธาวุโส ภิกฺขุ จกฺขุนา รูปํ ทิสฺวา เนว สุมโน โหติ น ทุมฺมโน, อุเปกฺขโก วิหรติ สโต สมฺปชาโน.\csromanspacer{} โสเตน สทฺทํ สุตฺวา ฯเปฯ.\csromanspacer{} ฆาเนน คนฺธํ ฆายิตฺวา ฯเปฯ.\csromanspacer{} ชิวฺหาย รสํ สายิตฺวา ฯเปฯ.\csromanspacer{} กาเยน โผฏฺฐพฺพํ ผุสิตฺวา ฯเปฯ.\csromanspacer{} มนสา ธมฺมํ วิญฺญาย เนว สุมโน โหติ น ทุมฺมโน, อุเปกฺขโก วิหรติ สโต สมฺปชาโน.\csromanspacer{} อิเม ฉ ธมฺมา \textbf{อุปฺปาเทตพฺพา}.}

\prose{(ฌ) กตเม ฉ ธมฺมา \textbf{อภิญฺเญยฺยา}.\csromanspacer{} ฉ อนุตฺตริยานิ, ทสฺสานานุตฺตริยํ สวนานุตฺตริยํ ลาภานุตฺตริยํ สิกฺขานุตฺตริยํ ปาริจริยานุตฺตริยํ อนุสฺสตานุตฺตริยํ.\csromanspacer{} อิเม ฉ ธมฺมา \textbf{อภิญฺเญยฺยา}.}

\prose{(ญ) กตเม ฉ ธมฺมา \textbf{สจฺฉิกาตพฺพา}.\csromanspacer{} ฉ อภิญฺญา.\csromanspacer{} อิธาวุโส ภิกฺขุ อเนกวิหิตํ อิทฺธิวิธํ ปจฺจนุโภติ เอโกปิ หุตฺวา พหุธา โหติ, พหุธาปิ หุตฺวา เอโก โหติ, อาวิภาวํ ติโรภาวํ, ติโรกุฏฺฏํ ติโรปาการํ ติโรปพฺพตํ อสชฺชมาโน คจฺฉติ เสยฺยาถาปิ}

\csromanmatikaanchor{mtk.145}
\csromantocmarkref{section}{สตฺตธมฺมา}{mtk.145}
\bookmark[dest={mtk.145},level=1]{สตฺตธมฺมา}
\prose{\csromanfolio{240}อากาเส, ปถวิยาปิ อุมฺมุชฺชนิมุชฺชํ กโรติ เสยฺยถาปิ อุทเก, อุทเกปิ อภิชฺชมาเน คจฺฉติ เสยฺยถาปิ ปถวิยํ, อากาเสปิ ปลฺลงฺเกน จงฺกมติ เสยฺยาถาปิ ปกฺขี สกุโณ, อิเมปิ จนฺทิมสูริเย เอวํมหิทฺธิเก เอวํมหานุภาเว ปาณินา ปรามสติ ปริมชฺชติ, ยาวพฺรหฺมโลกาปิ กาเยน วสํ วตฺเตติ. (1)}

\prose{ทิพฺพาย โสตธาตุยา วิสุทฺธาย อติกฺกนฺตมานุสิกาย อุโภ สทฺเท สุณาติ ทิพฺเพ จ มานุเส จ เย ทูเร สนฺติเก จ. (2)}

\prose{ปรสตฺตานํ ปรปุคฺคลานํ เจตสา เจโต ปริจฺจ ปชานาติ\footnote{ชานาติ (สฺยา, กํ)}, สราคํ วา จิตฺตํ สราคํ จิตฺตนฺติ ปชานาติ ฯเปฯ อวิมุตฺตํ วา จิตฺตํ อวิมุตฺตํ จิตฺตนฺติ ปชานาติ. (3)}

\prose{โส อเนกวิหิตํ ปุพฺเพนิวาสํ อนุสฺสรติ, เสยฺยถิทํ, เอกมฺปิ ชาติํ ฯเปฯ อิติ สาการํ ส-อุทฺเทสํ อเนกวิหิตํ ปุพฺเพนิวาสํ อนุสฺสรติ. (4)}

\prose{ทิพฺเพน จกฺขุนา วิสุทฺเธน อติกฺกนฺตมานุสเกน สตฺเต ปสฺสติ จวมาเน อุปปชฺชมาเน หีเน ปณีเต สุวณฺเณ ทุพฺพณฺเณ สุคเต ทุคฺคเต ยถากมฺมูปเค สตฺเต ปชานาติ ฯเปฯ. (5)}

\prose{อาสวานํ ขยา อนาสวํ เจโตวิมุตฺติํ ปญฺญาวิมุตฺติํ ทิฏฺเฐวธมฺเม สยํ อภิญฺญา สจฺฉิกตฺวา อุปสมฺปชฺช วิหรติ.\csromanspacer{} อิเม ฉ ธมฺมา สจฺฉิกาตพฺพา. (6)}

\prose{อิติ อิเม สฏฺฐิ ธมฺมา ภูตา ตจฺฉา ตถา อวิตถา อนญฺญถา สมฺมา ตถาคเตน อภิสมฺพุทฺธา.}

\csromanheader{2}{\textbf{สตฺต ธมฺมา}}
\vspace{\tipitakaparskip}
\csromancenter{\textbf{สตฺต ธมฺมา} \textbf{พหุการา} ฯเปฯ สตฺต ธมฺมา สจฺฉิกาตพฺพา.}
\prose{(ก) กตเม สตฺต ธมฺมา \textbf{พหุการา}.\csromanspacer{} สตฺต อริยธนานิ, สทฺธาธนํ สีลธนํ หิริธนํ โอตฺตปฺปธนํ สุตธนํ จาคธนํ ปญฺญาธนํ.\csromanspacer{} อิเม สตฺต ธมฺมา \textbf{พหุการา}.}

\csromanheader{2}{11. ทสุตฺตรสุตฺต}
\prose{\csromanfolio{241}(ข) กตเม สตฺต ธมฺมา \textbf{ภาเวตพฺพา.}\csromanspacer{} สตฺต สมฺโพชฺฌงฺคา, สติสมฺโพชฺฌงฺโค ธมฺมวิจยสมฺโพชฺฌงฺโค วีริยสมฺโพชฺฌงฺโค ปีติสมฺโพชฺฌงฺโค ปสฺสทฺธิสมฺโพชฺฌงฺโค สมาธิสมฺโพชฺฌงฺโค อุเปกฺขาสมฺโพชฺฌงฺโค.\csromanspacer{} อิเม สตฺต ธมฺมา \textbf{ภาเวตพฺพา.}}

\prose{(ค) กตเม สตฺต ธมฺมา \textbf{ปริญฺเญยฺยา}.\csromanspacer{} สตฺต วิญฺญาณฏฺฐิติโย.\csromanspacer{} สนฺตาวุโส สตฺตา นานตฺตกายา นานตฺตสญฺญิโน, เสยฺยถาปิ มนุสฺสา เอกจฺเจ จ เทวา เอกจฺเจ จ วินิปาติกา, อยํ ปฐมา วิญฺญาณฏฺฐิติ. (1)}

\prose{สนฺตาวุโส สตฺตา นานตฺตกายา เอกตฺตสญฺญิโน, เสยฺยถาปิ เทวา พฺรหฺมกายิกา ปฐมาภินิพฺพตฺตา, อยํ ทุติยา วิญฺญาณฏฺฐิติ. (2)}

\prose{สนฺตาวุโส สตฺตา เอกตฺตกายา นานตฺตสญฺญิโน, เสยฺยถาปิ เทวา อาภสฺสรา, อยํ ตติยา วิญฺญาณฏฺฐิติ. (3)}

\prose{สนฺตาวุโส สตฺตา เอกตฺตกายา เอกตฺตสญฺญิโน, เสยฺยถาปิ เทวา สุภกิณฺหา, อยํ จตุตฺถี วิญฺญาณฏฺฐิติ. (4)}

\prose{สนฺตาวุโส สตฺตา สพฺพโส รูปสญฺญานํ สมติกฺกมา ฯเปฯ “อนนฺโต อากาโส”ติ อากาสานญฺจายตนูปคา, อยํ ปญฺจมี วิญฺญาณฏฺฐิติ. (5)}

\prose{สนฺตาวุโส สตฺตา สพฺพโส อากาสานญฺจายตนํ สมติกฺกมฺม “อนนฺตํ วิญฺญาณนฺ”ติ วิญฺญาณญฺจายตนูปคา, อยํ ฉฏฺฐี วิญฺญาณฏฺฐิติ. (6)}

\prose{สนฺตาวุโส สตฺตา สพฺพโส วิญฺญาณญฺจายตนํ สมติกฺกมฺม “นตฺถิ กิญฺจี”ติ อากิญฺจญฺญายตนูปคา, อยํ สตฺตมี วิญฺญาณฏฺฐิติ.\csromanspacer{} อิเม สตฺต ธมฺมา \textbf{ปริญฺเญยฺยา}. (7)}

\prose{(ฆ) กตเม สตฺต ธมฺมา \textbf{ปหาตพฺพา}.\csromanspacer{} สตฺตานุสยา, กามราคานุสโย ปฏิฆานุสโย ทิฏฺฐานุสโย วิจิกิจฺฉานุสโย มานานุสโย ภวราคานุสโย อวิชฺชานุสโย.\csromanspacer{} อิเม สตฺต ธมฺมา \textbf{ปหาตพฺพา}.}

\prose{(\^{}อ) กตเม สตฺต ธมฺมา \textbf{หานภาคิยา}.\csromanspacer{} สตฺต อสทฺธมฺมา.\csromanspacer{} อิธาวุโส ภิกฺขุ อสฺสทฺโธ โหติ, อหิริโก โหติ, อโนตฺตปฺปี โหติ,}

\prose{\csromanfolio{242}อปฺปสฺสุโต โหติ, กุสีโต โหติ, มุฏฺฐสฺสติ โหติ, ทุปฺปญฺโญ โหติ.\csromanspacer{} อิเม สตฺต ธมฺมา หานภาคิยา.}

\prose{(จ) กตเม สตฺต ธมฺมา \textbf{วิเสสภาคิยา}.\csromanspacer{} สตฺต สทฺธมฺมา.\csromanspacer{} อิธาวุโส ภิกฺขุ สทฺโธ โหติ, หิริมา\footnote{หิริโก (สฺยา, กํ)} โหติ, โอตฺตปฺปี โหติ, พหุสฺสุโต โหติ, อารทฺธวีริโย โหติ, อุปฏฺฐิตสฺสติ โหติ, ปญฺญวา โหติ.\csromanspacer{} อิเม สตฺต ธมฺมา \textbf{วิเสสภาคิยา}.}

\prose{(ฉ) กตเม สตฺต ธมฺมา \textbf{ทุปฺปฏิวิชฺฌา}.\csromanspacer{} สตฺต สปฺปุริสธมฺมา.\csromanspacer{} อิธาวุโส ภิกฺขุ ธมฺมญฺญู จ โหติ อตฺถญฺญู จ อตฺตญฺญู จ มตฺตญฺญู จ กาลญฺญู จ ปริสญฺญู จ ปุคฺคลญฺญู จ.\csromanspacer{} อิเม สตฺต ธมฺมา \textbf{ทุปฺปฏิวิชฺฌา}.}

\prose{(ช) กตเม สตฺต ธมฺมา \textbf{อุปฺปาเทตพฺพา}.\csromanspacer{} สตฺต สญฺญา, อนิจฺจสญฺญา อนตฺตสญฺญา อสุภสญฺญา อาทีนวสญฺญา ปหานสญฺญา วิราคสญฺญา นิโรธสญฺญา.\csromanspacer{} อิเม สตฺต ธมฺมา \textbf{อุปฺปาเทตพฺพา}.}

\prose{(ฌ) กตเม สตฺต ธมฺมา \textbf{อภิญฺเญยฺยา}.\csromanspacer{} สตฺต นิทฺทสวตฺถูนิ.\csromanspacer{} อิธาวุโส ภิกฺขุ สิกฺขาสมาทาเน ติพฺพจฺฉนฺโท โหติ, อายติํ จ สิกฺขาสมาทาเน อวิคตเปโม.\csromanspacer{} ธมฺมนิสนฺติยา ติพฺพจฺฉนฺโท โหติ, อายติํ จ ธมฺมนิสนฺติยา อวิคตเปโม.\csromanspacer{} อิจฺฉาวินเย ติพฺพจฺฉนฺโท โหติ, อายติํ จ อิจฺฉาวินเย อวิคตเปโม.\csromanspacer{} ปฏิสลฺลาเน ติพฺพจฺฉนฺโท โหติ, อายติํ จ ปฏิสลฺลาเน อวิคตเปโม.\csromanspacer{} วีริยารมฺเภ ติพฺพจฺฉนฺโท โหติ, อายติํ จ วีริยารมฺเภ อวิคตเปโม.\csromanspacer{} สติเนปกฺเก ติพฺพจฺฉนฺโท โหติ, อายติํ จ สติเนปกฺเก อวิคตเปโม.\csromanspacer{} ทิฏฺฐิปฏิเวเธ ติพฺพจฺฉนฺโท โหติ, อายติํ จ ทิฏฺฐิปฏิเวเธ อวิคตเปโม.\csromanspacer{} อิเม สตฺต ธมฺมา \textbf{อภิญฺเญยฺยา}.}

\prose{(ญ) กตเม สตฺต ธมฺมา \textbf{สจฺฉิกาตพฺพา.}\csromanspacer{} สตฺต ขีณาสวพลานิ.\csromanspacer{} อิธาวุโส ขีณาสวสฺส ภิกฺขุโน อนิจฺจโต สพฺเพ สงฺขารา ยถาภูตํ สมฺมปฺปญฺญาย สุทิฏฺฐา โหนฺติ.\csromanspacer{} ยํปาวุโส ขีณาสวสฺส ภิกฺขุโน อนิจฺจโต สพฺเพ สงฺขารา ยถาภูตํ สมฺมปฺปญฺญาย สุทิฏฺฐา โหนฺติ, อิทมฺปิ ขีณาสวสฺส ภิกฺขุโน พลํ โหติ, ยํ พลํ อาคมฺม ขีณาสโว ภิกฺขุ อาสวานํ ขยํ ปฏิชานาติ “ขีณา เม อาสวา”ติ. (1)}

\csromanheader{2}{11. ทสุตฺตรสุตฺต}
\csromanmatikaanchor{mtk.146}
\csromantocmarkref{section}{อฏฺฐธมฺมา}{mtk.146}
\bookmark[dest={mtk.146},level=1]{อฏฺฐธมฺมา}
\prose{\csromanfolio{243}ปุน จปรํ อาวุโส ขีณาสวสฺส ภิกฺขุโน องฺคารกาสูปมา กามา ยถาภูตํ สมฺมปฺปญฺญาย สุทิฏฺฐา โหนฺติ.\csromanspacer{} ยํปาวุโส ฯเปฯ “ขีณา เม อาสวา”ติ. (2)}

\prose{ปุน จปรํ อาวุโส ขีณาสวสฺส ภิกฺขุโน วิเวกนินฺนํ จิตฺตํ โหติ วิเวกโปณํ วิเวกปพฺภารํ วิเวกฏฺฐํ เนกฺขมฺมาภิรตํ พฺยนฺตีภูตํ สพฺพโส อาสวฏฺฐานิเยหิ ธมฺเมหิ.\csromanspacer{} ยํปาวุโส ฯเปฯ “ขีณา เม อาสวา”ติ. (3)}

\prose{ปุน จปรํ อาวุโส ขีณาสวสฺส ภิกฺขุโน จตฺตาโร สติปฏฺฐานา ภาวิตา โหนฺติ สุภาวิตา.\csromanspacer{} ยํปาวุโส ฯเปฯ “ขีณา เม อาสวา”ติ. (4)}

\prose{ปุน จปรํ อาวุโส ขีณาสวสฺส ภิกฺขุโน ปญฺจินฺทฺริยานิ ภาวิตานิ โหนฺติ สุภาวิตานิ.\csromanspacer{} ยํปาวุโส ฯเปฯ “ขีณา เม อาสวา”ติ. (5)}

\prose{ปุน จปรํ อาวุโส ขีณาสวสฺส ภิกฺขุโน สตฺต โพชฺฌงฺคา ภาวิตา โหนฺติ สุภาวิตา.\csromanspacer{} ยํปาวุโส ฯเปฯ “ขีณา เม อาสวา”ติ. (6)}

\prose{ปุน จปรํ อาวุโส ขีณาสวสฺส ภิกฺขุโน อริโย อฏฺฐงฺคิโก มคฺโค ภาวิโต โหติ สุภาวิโต.\csromanspacer{} ยํปาวุโส ขีณาสวสฺส ภิกฺขุโน อริโย อฏฺฐงฺคิโก มคฺโค ภาวิโต โหติ สุภาวิโต, อิทมฺปิ ขีณาสวสฺส ภิกฺขุโน พลํ โหติ, ยํ พลํ อาคมฺม ขีณาสโว ภิกฺขุ อาสวานํ ขยํ ปฏิชานาติ “ขีณา เม อาสวา”ติ.\csromanspacer{} อิเม สตฺต ธมฺมา สจฺฉิกาตพฺพา. (7)}

\prose{อิติเม สตฺตติ ธมฺมา ภูตา ตจฺฉา ตถา อวิตถา อนญฺญถา สมฺมา ตถาคเตน อภิสมฺพุทฺธา.}

\nitthitamruled{ปฐมภาณวาโร นิฏฺฐิโต.}
\csromanheader{2}{\textbf{อฏฺฐ ธมฺมา}}
\vspace{\tipitakaparskip}
\csromancenter{\textbf{อฏฺฐ ธมฺมา} \textbf{พหุการา} ฯเปฯ อฏฺฐ ธมฺมา สจฺฉิกาตพฺพา.}
\prose{(ก) กตเม อฏฺฐ ธมฺมา \textbf{พหุการา}.\csromanspacer{} อฏฺฐ เหตู อฏฺฐ ปจฺจยา อาทิพฺรหฺมจริยิกาย ปญฺญาย อปฺปฏิลทฺธาย ปฏิลาภาย ปฏิลทฺธาย \csromanfolio{244}ภิยฺโยภาวาย เวปุลฺลาย ภาวนาย ปาริปูริยา สํวตฺตนฺติ.\csromanspacer{} กตเม อฏฺฐ.\csromanspacer{} อิธาวุโส ภิกฺขุ สตฺถารํ\footnote{สตฺถารํ วา (สฺยา, ก)} อุปนิสฺสาย วิหรติ อญฺญตรํ วา ครุฏฺฐานิยํ สพฺรหฺมจาริํ, ยตฺถสฺส ติพฺพํ หิโรตฺตปฺปํ ปจฺจุปฏฺฐิตํ โหติ เปมญฺจ คารโว จ.\csromanspacer{} อยํ ปฐโม เหตุ ปฐโม ปจฺจโย อาทิพฺรหฺมจริยิกาย ปญฺญาย อปฺปฏิลทฺธาย ปฏิลาภาย, ปฏิลทฺธาย ภิยฺโยภาวาย เวปุลฺลาย ภาวนาย ปาริปูริยา สํวตฺตติ. (1)}

\prose{ตํ โข ปน สตฺถารํ อุปนิสฺสาย วิหรติ อญฺญตรํ วา ครุฏฺฐานิยํ สพฺรหฺมจาริํ, ยตฺถสฺส ติพฺพํ หิโรตฺตปฺปํ ปจฺจุปฏฺฐิตํ โหติ เปมญฺจ คารโว จ.\csromanspacer{} เต กาเลน กาลํ อุปสงฺกมิตฺวา ปริปุจฺฉติ ปริปญฺหติ “อิทํ ภนฺเต กถํ, อิมสฺส โก อตฺโถ”ติ.\csromanspacer{} ตสฺส เต อายสฺมนฺโต อวิวฏญฺเจว วิวรนฺติ, อนุตฺตานีกตญฺจ อุตฺตานี\footnote{อนุตฺตานิกตญฺจ อุตฺตานิํ (ก)} กโรนฺติ, อเนกวิหิเตสุ จ กงฺขาฏฺฐานิเยสุ ธมฺเมสุ กงฺขํ ปฏิวิโนเทนฺติ.\csromanspacer{} อยํ ทุติโย เหตุ ทุติโย ปจฺจโย อาทิพฺรหฺมจริยิกาย ปญฺญาย อปฺปฏิลทฺธาย ปฏิลาภาย, ปฏิลทฺธาย ภิยฺโยภาวาย เวปุลฺลาย ภาวนาย ปาริปูริยา สํวตฺตติ. (2)}

\prose{ตํ โข ปน ธมฺมํ สุตฺวา ทฺวเยน วูปกาเสน สมฺปาเทติ กายวูปกาเสน จ จิตฺตวูปกาเสน จ.\csromanspacer{} อยํ ตติโย เหตุ ตติโย ปจฺจโย อาทิพฺรหฺมจริยิกาย ปญฺญาย อปฺปฏิลทฺธาย ปฏิลาภาย, ปฏิลทฺธาย ภิยฺโยภาวาย เวปุลฺลาย ภาวนาย ปาริปูริยา สํวตฺตติ. (3)}

\prose{ปุน จปรํ อาวุโส ภิกฺขุ สีลวา โหติ, ปาติโมกฺขสํวรสํวุโต วิหรติ อาจารโคจรสมฺปนฺโน, อณุมตฺเตสุ วชฺเชสุ ภยทสฺสาวี สมาทาย สิกฺขติ สิกฺขาปเทสุ.\csromanspacer{} อยํ จตุตฺโถ เหตุ จตุตฺโถ ปจฺจโย อาทิพฺรหฺมจริยิกาย ปญฺญาย อปฺปฏิลทฺธาย ปฏิลาภาย, ปฏิลทฺธาย ภิยฺโยภาวาย เวปุลฺลาย ภาวนาย ปาริปูริยา สํวตฺตติ. (4)}

\prose{ปุน จปรํ อาวุโส ภิกฺขุ พหุสฺสุโต โหติ สุตธโร สุตสนฺนิจโย, เย เต ธมฺมา อาทิกลฺยาณา มชฺเฌกลฺยาณา ปริโยสานกลฺยาณา สาตฺถา สพฺยญฺชนา เกวลปริปุณฺณํ ปริสุทฺธํ พฺรหฺมจริยํ อภิวทนฺติ, ตถารูปาสฺส ธมฺมา พหุสฺสุตา โหนฺติ ธาตา}

\vspace{\tipitakaparskip}
\csromancenter{ทสุตฺตรสุตฺต วจสา ปริจิตา มนสานุเปกฺขิตา ทิฏฺฐิยา สุปฺปฏิวิทฺธา.\csromanspacer{} อยํ ปญฺจโม เหตุ ปญฺจโม ปจฺจโย อาทิพฺรหฺมจริยิกาย ปญฺญาย อปฺปฏิลทฺธาย ปฏิลาภาย, ปฏิลทฺธาย ภิยฺโยภาวาย เวปุลฺลาย ภาวนาย ปาริปูริยา สํวตฺตติ. (5)}
\prose{\csromanfolio{245}ปุน จปรํ อาวุโส ภิกฺขุ อารทฺธวีริโย วิหรติ อกุสลานํ ธมฺมานํ ปหานาย กุสลานํ ธมฺมานํ อุปสมฺปทาย ถามวา ทฬฺหปรกฺกโม อนิกฺขิตฺตธุโร กุสเลสุ ธมฺเมสุ.\csromanspacer{} อยํ ฉฏฺโฐ เหตุ ฉฏฺโฐ ปจฺจโย อาทิพฺรหฺมจริยิกาย ปญฺญาย อปฺปฏิลทฺธาย ปฏิลาภาย, ปฏิลทฺธาย ภิยฺโยภาวาย เวปุลฺลาย ภาวนาย ปาริปูริยา สํวตฺตติ. (6)}

\prose{ปุน จปรํ อาวุโส ภิกฺขุ สติมา โหติ ปรเมน สติเนปกฺเกน สมนฺนาคโต จิรกตมฺปิ จิรภาสิตมฺปิ สริตา อนุสฺสริตา.\csromanspacer{} อยํ สตฺตโม เหตุ สตฺตโม ปจฺจโย อาทิพฺรหฺมจริยิกาย ปญฺญาย อปฺปฏิลทฺธาย ปฏิลาภาย, ปฏิลทฺธาย ภิยฺโยภาวาย เวปุลฺลาย ภาวนาย ปาริปูริยา สํวตฺตติ. (7)}

\prose{ปุน จปรํ อาวุโส ภิกฺขุ ปญฺจสุ อุปาทานกฺขนฺเธสุ อุทยพฺพยานุปสฺสี วิหรติ “อิติ รูปํ อิติ รูปสฺส สมุทโย อิติ รูปสฺส อตฺถงฺคโม, อิติ เวทนา อิติ เวทนาย สมุทโย อิติ เวทนาย อตฺถงฺคโม, อิติ สญฺญา อิติ สญฺญาย สมุทโย อิติ สญฺญาย อตฺถงฺคโม, อิติ สงฺขารา อิติ สงฺขารานํ สมุทโย อิติ สงฺขารานํ อตฺถงฺคโม, อิติ วิญฺญาณํ อิติ วิญฺญาณสฺส สมุทโย อิติ วิญฺญาณสฺส อตฺถงฺคโม”ติ.\csromanspacer{} อยํ อฏฺฐโม เหตุ อฏฺฐโม ปจฺจโย อาทิพฺรหฺมจริยิกาย ปญฺญาย อปฺปฏิลทฺธาย ปฏิลาภาย, ปฏิลทฺธาย ภิยฺโยภาวาย เวปุลฺลาย ภาวนาย ปาริปูริยา สํวตฺตติ.\csromanspacer{} อิเม อฏฺฐ ธมฺมา พหุการา. (8)}

\prose{(ข) กตเม อฏฺฐ ธมฺมา \textbf{ภาเวตพฺพา}.\csromanspacer{} อริโย อฏฺฐงฺคิโก มคฺโค, เสยฺยถิทํ, สมฺมาทิฏฺฐิ สมฺมาสงฺกปฺโป สมฺมาวาจา สมฺมากมฺมนฺโต สมฺมา-อาชีโว สมฺมาวายาโม สมฺมาสติ สมฺมาสมาธิ.\csromanspacer{} อิเม อฏฺฐ ธมฺมา \textbf{ภาเวตพฺพา}.}

\prose{(ค) กตเม อฏฺฐ ธมฺมา \textbf{ปริญฺเญยฺยา}.\csromanspacer{} อฏฺฐ โลกธมฺมา, ลาโภ จ อลาโภ จ ยโส จ อยโส จ นินฺทา จ ปสํสา จ สุขญฺจ ทุกฺขญฺจ.\csromanspacer{} อิเม อฏฺฐ ธมฺมา \textbf{ปริญฺเญยฺยา}.}

\prose{\csromanfolio{246}(ฆ) กตเม อฏฺฐ ธมฺมา \textbf{ปหาตพฺพา}.\csromanspacer{} อฏฺฐ มิจฺฉตฺตา, มิจฺฉาทิฏฺฐิ มิจฺฉาสงฺกปฺโป มิจฺฉาวาจา มิจฺฉากมฺมนฺโต มิจฺฉา-อาชีโว มิจฺฉาวายาโม มิจฺฉาสติ มิจฺฉาสมาธิ.\csromanspacer{} อิเม อฏฺฐ ธมฺมา \textbf{ปหาตพฺพา}.}

\prose{(\^{}อ) กตเม อฏฺฐ ธมฺมา \textbf{หานภาคิยา}.\csromanspacer{} อฏฺฐ กุสีตวตฺถูนิ.\csromanspacer{} อิธาวุโส ภิกฺขุนา กมฺมํ กาตพฺพํ โหติ, ตสฺส เอวํ โหติ “กมฺมํ โข เม กาตพฺพํ ภวิสฺสติ, กมฺมํ โข ปน เม กโรนฺตสฺส กาโย กิลมิสฺสติ, หนฺทาหํ นิปชฺชามี”ติ.\csromanspacer{} โส นิปชฺชติ, น วีริยํ อารภติ อปฺปตฺตสฺส ปตฺติยา อนธิคตสฺส อธิคมาย อสจฺฉิกตสฺส สจฺฉิกิริยาย, อิทํ ปฐมํ กุสีตวตฺถุ.}

\prose{ปุน จปรํ อาวุโส ภิกฺขุนา กมฺมํ กตํ โหติ, ตสฺส เอวํ โหติ “อหํ โข กมฺมํ อกาสิํ, กมฺมํ โข ปน เม กโรนฺตสฺส กาโย กิลนฺโต, หนฺทาหํ นิปชฺชามี”ติ.\csromanspacer{} โส นิปชฺชติ, น วีริยํ อารภติ ฯเปฯ อิทํ ทุติยํ กุสีตวตฺถุ.}

\prose{ปุน จปรํ อาวุโส ภิกฺขุนา มคฺโค คนฺตพฺโพ โหติ, ตสฺส เอวํ โหติ “มคฺโค โข เม คนฺตพฺโพ ภวิสฺสติ, มคฺคํ โข ปน เม คจฺฉนฺตสฺส กาโย กิลมิสฺสติ, หนฺทาหํ นิปชฺชามี”ติ.\csromanspacer{} โส นิปชฺชติ, น วีริยํ อารภติ ฯเปฯ อิทํ ตติยํ กุสีตวตฺถุ.}

\prose{ปุน จปรํ อาวุโส ภิกฺขุนา มคฺโค คโต โหติ, ตสฺส เอวํ โหติ “อหํ โข มคฺคํ อคมาสิํ, มคฺคํ โข ปน เม คจฺฉนฺตสฺส กาโย กิลนฺโต, หนฺทาหํ นิปชฺชามี”ติ.\csromanspacer{} โส นิปชฺชติ, น วีริยํ อารภติ ฯเปฯ อิทํ จตุตฺถํ กุสีตวตฺถุ.}

\prose{ปุน จปรํ อาวุโส ภิกฺขุ คามํ วา นิคมํ วา ปิณฺฑาย จรนฺโต น ลภติ ลูขสฺส วา ปณีตสฺส วา โภชนสฺส ยาวทตฺถํ ปาริปูริํ, ตสฺส เอวํ โหติ “อหํ โข คามํ วา นิคมํ วา ปิณฺฑาย จรนฺโต นาลตฺถํ ลูขสฺส วา ปณีตสฺส วา โภชนสฺส ยาวทตฺถํ ปาริปูริํ, ตสฺส เม กาโย กิลนฺโต อกมฺมญฺโญ, หนฺทาหํ นิปชฺชามี”ติ ฯเปฯ อิทํ ปญฺจมํ กุสีตวตฺถุ.}

\prose{ปุน จปรํ อาวุโส ภิกฺขุ คามํ วา นิคมํ วา ปิณฺฑาย จรนฺโต ลภติ ลูขสฺส วา ปณีตสฺส วา โภชนสฺส ยาวทตฺถํ ปาริปูริํ, ตสฺส เอวํ โหติ “อหํ โข คามํ วา นิคมํ วา ปิณฺฑาย จรนฺโต อลตฺถํ}

\vspace{\tipitakaparskip}
\csromancenter{ทสุตฺตรสุตฺต ลูขสฺส วา ปณีตสฺส วา โภชนสฺส ยาวทตฺถํ ปาริปูริํ, ตสฺส เม กาโย ครุโก อกมฺมญฺโญ, มาสาจิตํ มญฺเญ, หนฺทาหํ นิปชฺชามี”ติ.\csromanspacer{} โส นิปชฺชติ ฯเปฯ อิทํ ฉฏฺฐํ กุสีตวตฺถุ.}
\prose{\csromanfolio{247}ปุน จปรํ อาวุโส ภิกฺขุโน อุปฺปนฺโน โหติ อปฺปมตฺตโก อาพาโธ, ตสฺส เอวํ โหติ “อุปฺปนฺโน โข เม อปฺปมตฺตโก อาพาโธ อตฺถิ กปฺโป นิปชฺชิตุํ, หนฺทาหํ นิปชฺชามี”ติ.\csromanspacer{} โส นิปชฺชติ ฯเปฯ อิทํ สตฺตมํ กุสีตวตฺถุ.}

\prose{ปุน จปรํ อาวุโส ภิกฺขุ คิลานวุฏฺฐิโต โหติ อจิรวุฏฺฐิโต เคลญฺญา.\csromanspacer{} ตสฺส เอวํ โหติ “อหํ โข คิลานวุฏฺฐิโต อจิรวุฏฺฐิโต เคลญฺญา.\csromanspacer{} ตสฺส เม กาโย ทุพฺพโล อกมฺมญฺโญ อตฺถิ กปฺโป นิปชฺชิตุํ, หนฺทาหํ นิปชฺชามี”ติ.\csromanspacer{} โส นิปชฺชติ ฯเปฯ อิทํ อฏฺฐมํ กุสีตวตฺถุ.\csromanspacer{} อิเม อฏฺฐ ธมฺมา หานภาคิยา.}

\prose{(จ) กตเม อฏฺฐ ธมฺมา \textbf{วิเสสภาคิยา}.\csromanspacer{} อฏฺฐ อารมฺภวตฺถูนิ.\csromanspacer{} อิธาวุโส ภิกฺขุนา กมฺมํ กาตพฺพํ โหติ, ตสฺส เอวํ โหติ “กมฺมํ โข เม กาตพฺพํ ภวิสฺสติ, กมฺมํ โข ปน เม กโรนฺเตน น สุกรํ พุทฺธานํ สาสนํ มนสิกาตุํ, หนฺทาหํ วีริยํ อารภามิ อปตฺตสฺส ปตฺติยา อนธิคตสฺส อธิคมาย อสจฺฉิกตสฺส สจฺฉิกิริยายา”ติ.\csromanspacer{} โส วีริยํ อารภติ อปฺปตฺตสฺส ปตฺติยา อนธิคตสฺส อธิคมาย อสจฺฉิกตสฺส สจฺฉิกิริยาย, อิทํ ปฐมํ อารมฺภวตฺถุ.}

\prose{ปุน จปรํ อาวุโส ภิกฺขุนา กมฺมํ กตํ โหติ, ตสฺส เอวํ โหติ “อหํ โข กมฺมํ อกาสิํ, กมฺมํ โข ปนาหํ กโรนฺโต นาสกฺขิํ พุทฺธานํ สาสนํ มนสิกาตุํ, หนฺทาหํ วีริยํ อารภามิ ฯเปฯ อิทํ ทุติยํ อารมฺภวตฺถุ.}

\prose{ปุน จปรํ อาวุโส ภิกฺขุนา มคฺโค คนฺตพฺโพ โหติ, ตสฺส เอวํ โหติ “มคฺโค โข เม คนฺตพฺโพ ภวิสฺสติ, มคฺคํ โข ปน เม คจฺฉนฺเตน น สุกรํ พุทฺธานํ สาสนํ มนสิกาตุํ, หนฺทาหํ วีริยํ อารภามิ ฯเปฯ อิทํ ตติยํ อารมฺภวตฺถุ.}

\prose{ปุน จปรํ อาวุโส ภิกฺขุนา มคฺโค คโต โหติ, ตสฺส เอวํ โหติ “อหํ โข มคฺคํ อคมาสิํ, มคฺคํ โข ปนาหํ คจฺฉนฺโต นาสกฺขิํ พุทฺธานํ สาสนํ มนสิกาตุํ, หนฺทาหํ วีริยํ อารภามิ ฯเปฯ อิทํ จตุตฺถํ อารมฺภวตฺถุ.}

\prose{\csromanfolio{248}ปุน จปรํ อาวุโส ภิกฺขุ คามํ วา นิคมํ วา ปิณฺฑาย จรนฺโต น ลภติ ลูขสฺส วา ปณีตสฺส วา โภชนสฺส ยาวทตฺถํ ปาริปูริํ, ตสฺส เอวํ โหติ “อหํ โข คามํ วา นิคมํ วา ปิณฺฑาย จรนฺโต นาลตฺถํ ลูขสฺส วา ปณีตสฺส วา โภชนสฺส ยาวทตฺถํ ปาริปูริํ, ตสฺส เม กาโย ลหุโก กมฺมญฺโญ, หนฺทาหํ วีริยํ อารภามิ ฯเปฯ อิทํ ปญฺจมํ อารมฺภวตฺถุ.}

\prose{ปุน จปรํ อาวุโส ภิกฺขุ คามํ วา นิคมํ วา ปิณฺฑาย จรนฺโต ลภติ ลูขสฺส วา ปณีตสฺส วา โภชนสฺส ยาวทตฺถํ ปาริปูริํ, ตสฺส เอวํ โหติ “อหํ โข คามํ วา นิคมํ วา ปิณฺฑาย จรนฺโต อลตฺถํ ลูขสฺส วา ปณีตสฺส วา โภชนสฺส ยาวทตฺถํ ปาริปูริํ, ตสฺส เม กาโย พลวา กมฺมญฺโญ, หนฺทาหํ วีริยํ อารภามิ ฯเปฯ อิทํ ฉฏฺฐํ อารมฺภวตฺถุ.}

\prose{ปุน จปรํ อาวุโส ภิกฺขุโน อุปฺปนฺโน โหติ อปฺปมตฺตโก อาพาโธ, ตสฺส เอวํ โหติ “อุปฺปนฺโน โข เม อยํ อปฺปมตฺตโก อาพาโธ ฐานํ โข ปเนตํ วิชฺชติ, ยํ เม อาพาโธ ปวฑฺเฒยฺย, หนฺทาหํ วีริยํ อารภามิ ฯเปฯ อิทํ สตฺตมํ อารมฺภวตฺถุ.}

\prose{ปุน จปรํ อาวุโส ภิกฺขุ คิลานวุฏฺฐิโต โหติ อจิรวุฏฺฐิโต เคลญฺญา, ตสฺส เอวํ โหติ “อหํ โข คิลานวุฏฺฐิโต อจิรวุฏฺฐิโต เคลญฺญา, ฐานํ โข ปเนตํ วิชฺชติ, ยํ เม อาพาโธ ปจฺจุทาวตฺเตยฺย, หนฺทาหํ วีริยํ อารภามิ อปฺปตฺตสฺส ปตฺติยา อนธิคตสฺส อธิคมาย อสจฺฉิกตสฺส สจฺฉิกิริยายา”ติ.\csromanspacer{} โส วีริยํ อารภติ อปฺปตฺตสฺส ปตฺติยา อนธิคตสฺส อธิคมาย อสจฺฉิกตสฺส สจฺฉิกิริยาย, อิทํ อฏฺฐมํ อารมฺภวตฺถุ.\csromanspacer{} อิเม อฏฺฐ ธมฺมา วิเสสภาคิยา.}

\prose{(ฉ) กตเม อฏฺฐ ธมฺมา \textbf{ทุปฺปฏิวิชฺฌา}.\csromanspacer{} อฏฺฐ อกฺขณา อสมยา พฺรหฺมจริยวาสาย.\csromanspacer{} อิธาวุโส ตถาคโต จ โลเก อุปฺปนฺโน โหติ อรหํ สมฺมาสมฺพุทฺโธ, ธมฺโม จ เทสิยติ โอปสมิโก ปรินิพฺพานิโก สมฺโพธคามี สุคตปฺปเวทิโต.\csromanspacer{} อยญฺจ ปุคฺคโล นิรยํ อุปปนฺโน โหติ.\csromanspacer{} อยํ ปฐโม อกฺขโณ อสมโย พฺรหฺมจริยวาสาย.}

\prose{ปุน จปรํ อาวุโส ตถาคโต จ โลเก อุปฺปนฺโน โหติ อรหํ สมฺมาสมฺพุทฺโธ, ธมฺโม จ เทสิยติ โอปสมิโก ปรินิพฺพานิโก}

\vspace{\tipitakaparskip}
\csromancenter{ทสุตฺตรสุตฺต สมฺโพธคามี สุคตปฺปเวทิโต.\csromanspacer{} อยญฺจ ปุคฺคโล ติรจฺฉานโยนิํ อุปปนฺโน โหติ.\csromanspacer{} อยํ ทุติโย อกฺขโณ อสมโย พฺรหฺมจริยวาสาย.}
\prose{\csromanfolio{249}ปุน จปรํ ฯเปฯ เปตฺติวิสยํ อุปปนฺโน โหติ.\csromanspacer{} อยํ ตติโย อกฺขโณ อสมโย พฺรหฺมจริยวาสาย.}

\prose{ปุน จปรํ ฯเปฯ อญฺญตรํ ทีฆายุกํ เทวนิกายํ อุปปนฺโน โหติ.\csromanspacer{} อยํ จตุตฺโถ อกฺขโณ อสมโย พฺรหฺมจริยวาสาย.}

\prose{ปุน จปรํ ฯเปฯ ปจฺจนฺติเมสุ ชนปเทสุ ปจฺจาชาโต โหติ มิลกฺเขสุ อวิญฺญาตาเรสุ, ยตฺถ นตฺถิ คติ ภิกฺขูนํ ภิกฺขุนีนํ อุปาสกานํ อุปาสิกานํ.\csromanspacer{} อยํ ปญฺจโม อกฺขโณ อสมโย พฺรหฺมจริยวาสาย.}

\prose{ปุน จปรํ ฯเปฯ อยญฺจ ปุคฺคโล มชฺฌิเมสุ ชนปเทสุ ปจฺจาชาโต โหติ, โส จ โหติ มิจฺฉาทิฏฺฐิโก วิปรีตทสฺสโน “นตฺถิ ทินฺนํ นตฺถิ ยิฏฺฐํ, นตฺถิ หุตํ, นตฺถิ สุกตทุกฺกฏานํ กมฺมานํ ผลํ วิปาโก, นตฺถิ อยํ โลโก, นตฺถิ ปโร โลโก, นตฺถิ มาตา, นตฺถิ ปิตา, นตฺถิ สตฺตา โอปปาติกา, นตฺถิ โลเก สมณพฺราหฺมณา สมฺมคฺคตา สมฺมาปฏิปนฺนา เย อิมญฺจ โลกํ ปรญฺจ โลกํ สยํ อภิญฺญา สจฺฉิกตฺวา ปเวเทนฺตี”ติ.\csromanspacer{} อยํ ฉฏฺโฐ อกฺขโณ อสมโย พฺรหฺมจริยวาสาย.}

\prose{ปุน จปรํ ฯเปฯ อยญฺจ ปุคฺคโล มชฺฌิเมสุ ชนปเทสุ ปจฺจาชาโต โหติ, โส จ โหติ ทุปฺปญฺโญ ชโฬ เอฬมูโค, นปฺปฏิพโล สุภาสิตทุพฺภาสิตานมตฺถมญฺญาตุํ.\csromanspacer{} อยํ สตฺตโม อกฺขโณ อสมโย พฺรหฺมจริยวาสาย.}

\prose{ปุน จปรํ ฯเปฯ อยญฺจ ปุคฺคโล มชฺฌิเมสุ ชนปเทสุ ปจฺจาชาโต โหติ, โส จ โหติ ปญฺญวา อชโฬ อเนฬมูโค, ปฏิพโล สุภาสิตทุพฺภาสิตานมตฺถมญฺญาตุํ.\csromanspacer{} อยํ อฏฺฐโม อกฺขโณ อสมโย พฺรหฺมจริยวาสาย.\csromanspacer{} อิเม อฏฺฐ ธมฺมา ทุปฺปฏิวิชฺฌา.}

\prose{(ช) กตเม อฏฺฐ ธมฺมา \textbf{อุปฺปาเทตพฺพา}.\csromanspacer{} อฏฺฐ มหาปุริสวิตกฺกา.\csromanspacer{} อปฺปิจฺฉสฺสายํ ธมฺโม, นายํ ธมฺโม มหิจฺฉสฺส.\csromanspacer{} สนฺตุฏฺฐสฺสายํ ธมฺโม, นายํ ธมฺโม อสนฺตุฏฺฐสฺส.\csromanspacer{} ปวิวิตฺตสฺสยํ ธมฺโม, นายํ ธมฺโม สงฺคณิการามสฺส.\csromanspacer{} อารทฺธวีริยสฺสายํ ธมฺโม, นายํ ธมฺโม กุสีตสฺส.\csromanspacer{} อุปฏฺฐิตสติสฺสายํ ธมฺโม, นายํ ธมฺโม มุฏฺฐสฺสติสฺส.\csromanspacer{} สมาหิตสฺสายํ ธมฺโม, นายํ ธมฺโม \csromanfolio{250}อสมาหิตสฺส.\csromanspacer{} ปญฺญวโต\footnote{ปญฺญาวโต (สี, อิ)} อยํ ธมฺโม, นายํ ธมฺโม ทุปฺปญฺญสฺส.\csromanspacer{} นิปฺปปญฺจสฺสายํ ธมฺโม, นายํ ธมฺโม ปปญฺจารามสฺสาติ\footnote{นิปฺปปญฺจารามสฺส อยํ ธมฺโม นิปฺปปญฺจรติโน, นายํ ธมฺโม ปปญฺจารามสฺส ปปญฺจรติโนติ (สี, สฺยา, อิ) องฺคุตฺตเรปิ ตเถว ทิสฺสติ. อฏฺฐกถาฏีกา ปน โอโลเกตพฺพา.}.\csromanspacer{} อิเม อฏฺฐ ธมฺมา อุปฺปาเทตพฺพา.}

\prose{(ฌ) กตเม อฏฺฐ ธมฺมา \textbf{อภิญฺเญยฺยา.}\csromanspacer{} อฏฺฐ อภิภายตนานิ.\csromanspacer{} อชฺฌตฺตํ รูปสญฺญี เอโก พหิทฺธา รูปานิ ปสฺสติ ปริตฺตานิ สุวณฺณทุพฺพณฺณานิ, “ตานิ อภิภุยฺย ชานามิ ปสฺสามี”ติ เอวํสญฺญี โหติ, อิทํ ปฐมํ อภิภายตนํ. อชฺฌตฺตํ รูปสญฺญี เอโก พหิทฺธา รูปานิ ปสฺสติ อปฺปมาณานิ สุวณฺณทุพฺพณฺณานิ, “ตานิ อภิภุยฺย ชานามิ ปสฺสามี”ติ เอวํสญฺญี โหติ, อิทํ ทุติยํ อภิภายตนํ.}

\prose{อชฺฌตฺตํ อรูปสญฺญี เอโก พหิทฺธา รูปานิ ปสฺสติ ปริตฺตานิ สุวณฺณทุพฺพณฺณานิ, “ตานิ อภิภุยฺย ชานามิ ปสฺสามี”ติ เอวํสญฺญี โหติ, อิทํ ตติยํ อภิภายตนํ.}

\prose{อชฺฌตฺตํ อรูปสญฺญี เอโก พหิทฺธา รูปานิ ปสฺสติ อปฺปมาณานิ สุวณฺณทุพฺพณฺณานิ, “ตานิ อภิภุยฺย ชานามิ ปสฺสามี”ติ เอวํสญฺญี โหติ, อิทํ จตุตฺถํ อภิภายตนํ.}

\prose{อชฺฌตฺตํ อรูปสญฺญี เอโก พหิทฺธา รูปานิ ปสฺสติ นีลานิ นีลวณฺณานิ นีลนิทสฺสนานิ นีลนิภาสานิ.\csromanspacer{} เสยฺยถาปิ นาม อุมาปุปฺผํ นีลํ นีลวณฺณํ นีลนิทสฺสนํ นีลนิภาสํ, เสยฺยถาปิ วา ปน ตํ วตฺถํ พาราณเสยฺยกํ อุภโตภาควิมฏฺฐํ นีลํ นีลวณฺณํ นีลนิทสฺสนํ นีลนิภาสํ, เอวเมว อชฺฌตฺตํ อรูปสญฺญี เอโก พหิทฺธา รูปานิ ปสฺสติ นีลานิ นีลวณฺณานิ นีลนิทสฺสนานิ นีลนิภาสานิ, “ตานิ อภิภุยฺย ชานามิ ปสฺสามี”ติ เอวํสญฺญี โหติ, อิทํ ปญฺจมํ อภิภายตนํ.}

\prose{อชฺฌตฺตํ อรูปสญฺญี เอโก พหิทฺธา รูปานิ ปสฺสติ ปีตานิ ปีตวณฺณานิ ปีตนิทสฺสนานิ ปีตนิภาสานิ.\csromanspacer{} เสยฺยถาปิ นาม กณิการปุปฺผํ ปีตํ ปีตวณฺณํ ปีตนิทสฺสนํ ปีตนิภาสํ, เสยฺยถาปิ วา ปน ตํ วตฺถํ พาราณ-}

\vspace{\tipitakaparskip}
\csromancenter{ทสุตฺตรสุตฺต เสยฺยกํ อุภโตภาควิมฏฺฐํ ปีตํ ปีตวณฺณํ ปีตนิทสฺสนํ ปีตนิภาสํ, เอวเมว อชฺฌตฺตํ อรูปสญฺญี เอโก พหิทฺธา รูปานิ ปสฺสติ ปีตานิ ปีตวณฺณานิ ปีตนิทสฺสนานิ ปีตนิภาสานิ, “ตานิ อภิภุยฺย ชานามิ ปสฺสามี”ติ เอวํสญฺญี โหติ, อิทํ ฉฏฺฐํ อภิภายตนํ.}
\prose{\csromanfolio{251}อชฺฌตฺตํ อรูปสญฺญี เอโก พหิทฺธา รูปานิ ปสฺสติ โลหิตกานิ โลหิตกวณฺณานิ โลหิตกนิทสฺสนานิ โลหิตกนิภาสานิ.\csromanspacer{} เสยฺยถาปิ นาม พนฺธุชีวกปุปฺผํ โลหิตกํ โลหิตกวณฺณํ โลหิตกนิทสฺสนํ โลหิตกนิภาสํ, เสยฺยถาปิ วา ปน ตํ วตฺถํ พาราณเสยฺยกํ อุภโตภาควิมฏฺฐํ โลหิตกํ โลหิตกวณฺณํ โลหิตกนิทสฺสนํ โลหิตกนิภาสํ, เอวเมว อชฺฌตฺตํ อรูปสญฺญี เอโก พหิทฺธา รูปานิ ปสฺสติ โลหิตกานิ โลหิตกวณฺณานิ โลหิตกนิทสฺสนานิ โลหิตกนิภาสานิ, “ตานิ อภิภุยฺย ชานามิ ปสฺสามี”ติ เอวํ สญฺญี โหติ, อิทํ สตฺตมํ อภิภายตนํ.}

\prose{อชฺฌตฺตํ อรูปสญฺญี เอโก พหิทฺธา รูปานิ ปสฺสติ โอทาตานิ โอทาตวณฺณานิ โอทาตนิทสฺสนานิ โอทาตนิภาสานิ.\csromanspacer{} เสยฺยถาปิ นาม โอสธิตารกา โอทาตา โอทาตวณฺณา โอทาตนิทสฺสนา โอทาตนิภาสา, เสยฺยถาปิ วา ปน ตํ วตฺถํ พาราณเสยฺยกํ อุภโตภาควิมฏฺฐํ โอทาตํ โอทาตวณฺณํ โอทาตนิทสฺสนํ โอทาตนิภาสํ, เอวเมว อชฺฌตฺตํ อรูปสญฺญี เอโก พหิทฺธา รูปานิ ปสฺสติ โอทาตานิ โอทาตวณฺณานิ โอทาตนิทสฺสนานิ โอทาตนิภาสานิ, “ตานิ อภิภุยฺย ชานามิ ปสฺสามี”ติ เอวํสญฺญี โหติ, อิทํ อฏฺฐมํ อภิภายตนํ.\csromanspacer{} อิเม อฏฺฐธมฺมา อภิญฺเญยฺยา.}

\prose{(ญ) กตเม อฏฺฐ ธมฺมา \textbf{สจฺฉิกาตพฺพา}.\csromanspacer{} อฏฺฐ วิโมกฺขา.\csromanspacer{} รูปี รูปานิ ปสฺสติ, อยํ ปฐโม วิโมกฺโข.}

\prose{อชฺฌตฺตํ อรูปสญฺญี เอโก พหิทฺธา รูปานิ ปสฺสติ, อยํ ทุติโย วิโมกฺโข.}

\prose{สุภนฺเตว อธิมุตฺโต โหติ, อยํ ตติโย วิโมกฺโข.}

\prose{สพฺพโส รูปสญฺญานํ สมติกฺกมา ปฏิฆสญฺญานํ อตฺถงฺคมา นานตฺตสญฺญานํ อมนสิการา “อนนฺโต อากาโส”ติ อากาสานญฺจายตนํ อุปสมฺปชฺช วิหรติ, อยํ จตุตฺโถ วิโมกฺโข.}

\csromanmatikaanchor{mtk.147}
\csromantocmarkref{section}{โนวธมฺมา}{mtk.147}
\bookmark[dest={mtk.147},level=1]{โนวธมฺมา}
\prose{\csromanfolio{252}สพฺพโส อากาสานญฺจยตนํ สมติกฺกมฺม “อนนฺตํ วิญฺญาณนฺ”ติ วิญฺญาณญฺจายตนํ อุปสมฺปชฺช วิหรติ, อยํ ปญฺจโม วิโมกฺโข.}

\prose{สพฺพโส วิญฺญาณญฺจายตนํ สมติกฺกมฺม “นตฺถิ กิญฺจี”ติ อากิญฺจญฺญายตนํ อุปสมฺปชฺช วิหรติ, อยํ ฉฏฺโฐ วิโมกฺโข.}

\prose{สพฺพโส อากิญฺจญฺญายตนํ สมติกฺกมฺม เนวสญฺญานาสญฺญายตนํ อุปสมฺปชฺช วิหรติ, อยํ สตฺตโม วิโมกฺโข.}

\prose{สพฺพโส เนวสญฺญานาสญฺญายตนํ สมติกฺกมฺม สญฺญาเวทยิตนิโรธํ อุปสมฺปชฺช วิหรติ, อยํ อฏฺฐโม วิโมกฺโข.\csromanspacer{} อิเม อฏฺฐ ธมฺมา สจฺฉิกาตพฺพา.}

\prose{อิติ อิเม อสีติ ธมฺมา ภูตา ตจฺฉา ตถา อวิตถา อนญฺญถา สมฺมา ตถาคเตน อภิสมฺพุทฺธา.}

\csromanheader{2}{\textbf{นว ธมฺมา}}
\vspace{\tipitakaparskip}
\csromancenter{\textbf{นว ธมฺมา} \textbf{พหุการา} ฯเปฯ นว ธมฺมา สจฺฉิกาตพฺพา.}
\prose{(ก) กตเม นว ธมฺมา \textbf{พหุการา}.\csromanspacer{} นว โยนิโสมนสิการมูลกา ธมฺมา.\csromanspacer{} โยนิโสมนสิกโรโต ปาโมชฺชํ ชายติ, ปมุทิตสฺส ปีติ ชายติ, ปีติมนสฺส กาโย ปสฺสมฺภติ, ปสฺสทฺธกาโย สุขํ เวเทติ, สุขิโน จิตฺตํ สมาธิยติ, สมาหิเต จิตฺเต ยถาภูตํ ชานาติ ปสฺสติ, ยถาภูตํ ชานํ ปสฺสํ นิพฺพินฺทติ, นิพฺพินฺทํ วิรชฺชติ, วิราคา วิมุจฺจติ.\csromanspacer{} อิเม นว ธมฺมา \textbf{พหุการา}.}

\prose{(ข) กตเม นว ธมฺมา \textbf{ภาเวตพฺพา}.\csromanspacer{} นว ปาริสุทฺธิ ปธานิยงฺคานิ, สีลวิสุทฺธิ ปาริสุทฺธิปธานิยงฺคํ, จิตฺตวิสุทฺธิ ปาริสุทฺธิปธานิยงฺคํ, ทิฏฺฐิวิสุทฺธิ ปาริสุทฺธิปธานิยงฺคํ, กงฺขาวิตรณวิสุทฺธิ ปาริสุทฺธิปธานิยงฺคํ, มคฺคามคฺคญาณทสฺสนวิสุทฺธิ ปาริสุทฺธิปธานิยงฺคํ, ปฏิปทาญาณทสฺสนวิสุทฺธิ ปาริสุทฺธิปธานิยงฺคํ, ญาณทสฺสนวิสุทฺธิ ปาริสุทฺธิปธานิยงฺคํ, ปญฺญาวิสุทฺธิ ปาริสุทฺธิปธานิยางฺคํ, วิมุตฺติวิสุทฺธิ ปาริสุทฺธิปธานิยงฺคํ.\csromanspacer{} อิเม นว ธมฺมา \textbf{ภาเวตพฺพา}.}

\prose{(ค) กตเม นว ธมฺมา \textbf{ปริญฺเญยฺยา}.\csromanspacer{} นว สตฺตาวาสา.\csromanspacer{} สนฺตาวุโส สตฺตา นานตฺตกายา นานตฺตสญฺญิโน, เสยฺยถาปิ มนุสฺสา เอกจฺเจ จ เทวา เอกจฺเจ จ วินิปาติกา, อยํ ปฐโม สตฺตาวาโส.}

\csromanheader{2}{11. ทสุตฺตรสุตฺต}
\prose{\csromanfolio{253}สนฺตาวุโส สตฺตา นานตฺตกายา เอกตฺตสญฺญิโน, เสยฺยถาปิ เทวา พฺรหฺมกายิกา ปฐมาภินิพฺพตฺตา, อยํ ทุติโย สตฺตาวาโส.}

\prose{สนฺตาวุโส สตฺตา เอกตฺตกายา นานตฺตสญฺญิโน, เสยฺยถาปิ เทวา อาภสฺสรา, อยํ ตติโย สตฺตาวาโส.}

\prose{สนฺตาวุโส สตฺตา เอกตฺตกายา เอกตฺตสญฺญิโน, เสยฺยถาปิ เทวา สุภกิณฺหา, อยํ จตุตฺโถ สตฺตาวาโส.}

\prose{สนฺตาวุโส สตฺตา อสญฺญิโน อปฺปฏิสํเวทิโน, เสยฺยถาปิ เทวา อสญฺญสตฺตา, อยํ ปญฺจโม สตฺตาวาโส.}

\prose{สนฺตาวุโส สตฺตา สพฺพโส รูปสญฺญานํ สมติกฺกมา ปฏิฆสญฺญานํ อตฺถงฺคมา นานตฺตสญฺญานํ อมนสิการา “อนนฺโต อากาโส”ติ อากาสานญฺจายตนูปคา, อยํ ฉฏฺโฐ สตฺตาวาโส.}

\prose{สนฺตาวุโส สตฺตา สพฺพโส อากาสานญฺจายตนํ สมติกฺกมฺม “อนนฺตํ วิญฺญาณนฺ”ติ วิญฺญาณญฺจายตนูปคา, อยํ สตฺตโม สตฺตาวาโส.}

\prose{สนฺตาวุโส สตฺตา สพฺพโส วิญฺญาณญฺจายตนํ สมติกฺกมฺม “นตฺถิ กิญฺจี”ติ อากิญฺจญฺญายตนูปคา, อยํ อฏฺฐโม สตฺตาวาโส.}

\prose{สนฺตาวุโส สตฺตา สพฺพโส อากิญฺจญฺญายตนํ สมติกฺกมฺม เนวสญฺญานาสญฺญายตนูปคา, อยํ นวโม สตฺตาวาโส.\csromanspacer{} อิเม นว ธมฺมา ปริญฺเญยฺยา.}

\prose{(ฆ) กตเม นว ธมฺมา \textbf{ปหาตพฺพา}.\csromanspacer{} นว ตณฺหามูลกา ธมฺมา.\csromanspacer{} ตณฺหํ ปฏิจฺจ ปริเยสนา, ปริเยสนํ ปฏิจฺจ ลาโภ, ลาภํ ปฏิจฺจ วินิจฺฉโย, วินิจฺฉยํ ปฏิจฺจ ฉนฺทราโค, ฉนฺทราคํ ปฏิจฺจ อชฺโฌสานํ, อชฺโฌสานํ ปฏิจฺจ ปริคฺคโห, ปริคฺคหํ ปฏิจฺจ มจฺฉริยํ, มจฺฉริยํ ปฏิจฺจ อารกฺโข, อารกฺขาธิกรณํ\footnote{อารกฺขาธิกรณํ ปฏิจฺจ (สฺยา, อิ, ก)} ทณฺฑาทาน สตฺถาทาน กลห วิคฺคห วิวาท ตุวํตุวํ เปสุญฺญ มุสาวาทา อเนเก ปาปกา อกุสลา ธมฺมา สํวตฺตนฺติ.\csromanspacer{} อิเม นว ธมฺมา \textbf{ปหาตพฺพา}.}

\prose{\csromanfolio{254}(\^{}อ) กตเม นว ธมฺมา \textbf{หานภาคิยา}.\csromanspacer{} นว อาฆาตวตฺถูนิ. “อนตฺถํ เม อจรี”ติ อาฆาตํ พนฺธติ, “อนตฺถํ เม จรตี”ติ อาฆาตํ พนฺธติ, “อนตฺถํ เม จริสฺสตี”ติ อาฆาตํ พนฺธติ, “ปิยสฺส เม มนาปสฺส อนตฺถํ อจรี”ติ อาฆาตํ พนฺธติ ฯเปฯ “อนตฺถํ จรตี”ติ อาฆาตํ พนฺธติ ฯเปฯ อนตฺถํ จริสฺสตี”ติ อาฆาตํ พนฺธติ, “อปฺปิยสฺส เม อมนาปสฺส อตฺถํ อจรี”ติ อาฆาตํ พนฺธติ ฯเปฯ “อตฺถํ จรตี”ติ อาฆาตํ พนฺธติ ฯเปฯ อตฺถํ จริสฺสตี”ติ อาฆาตํ พนฺธติ.\csromanspacer{} อิเม นว ธมฺมา \textbf{หานภาคิยา}.}

\prose{(จ) กตเม นว ธมฺมา \textbf{วิเสสภาคิยา}.\csromanspacer{} นว อาฆาตปฏิวินยา. “อนตฺถํ เม อจริ, ตํ กุเตตฺถ ลพฺภา”ติ อาฆาตํ ปฏิวิเนติ, “อนตฺถํ เม จรติ, ตํ กุเตตฺถ ลพฺภา”ติ อาฆาตํ ปฏิวิเนติ, “อนตฺถํ เม จริสฺสติ, ตํ กุเตตฺถ ลพฺภา”ติ อาฆาตํ ปฏิวิเนติ, “ปิยสฺส เม มนาปสฺส อนตฺถํ อจริ ฯเปฯ อนตฺถํ จรติ ฯเปฯ อนตฺถํ จริสฺสติ, ตํ กุเตตฺถ ลพฺภา”ติ อาฆาตํ ปฏิวิเนติ, “อปฺปิยสฺส เม อมนาปสฺส อตฺถํ อจริ ฯเปฯ อตฺถํ จรติ ฯเปฯ อตฺถํ จริสฺสติ, ตํ กุเตตฺถ ลพฺภา”ติ อาฆาตํ ปฏิวิเนติ.\csromanspacer{} อิเม นว ธมฺมา \textbf{วิเสสภาคิยา}.}

\prose{(ฉ) กตเม นว ธมฺมา \textbf{ทุปฺปฏิวิชฺฌา}.\csromanspacer{} นว นานตฺตา.\csromanspacer{} ธาตุนานตฺตํ ปฏิจฺจ อุปฺปชฺชติ ผสฺสนานตฺตํ, ผสฺสนานตฺตํ ปฏิจฺจ อุปฺปชฺชติ เวทนานานตฺตํ, เวทนานานตฺตํ ปฏิจฺจ อุปฺปชฺชติ สญฺญานานตฺตํ, สญฺญานานตฺตํ ปฏิจฺจ อุปฺปชฺชติ สงฺกปฺปนานตฺตํ, สงฺกปฺปนานตฺตํ ปฏิจฺจ อุปฺปชฺชติ ฉนฺทนานตฺตํ, ฉนฺทนานตฺตํ ปฏิจฺจ อุปฺปชฺชติ ปริฬาหนานตฺตํ, ปริฬาหนานตฺตํ ปฏิจฺจ อุปฺปชฺชติ ปริเยสนานานตฺตํ, ปริเยสนานานตฺตํ ปฏิจฺจ อุปฺปชฺชติ ลาภนานตฺตํ.\csromanspacer{} อิเม นว ธมฺมา \textbf{ทุปฺปฏิวิชฺฌา}.}

\prose{(ช) กตเม นว ธมฺมา \textbf{อุปฺปาเทตพฺพา}.\csromanspacer{} นว สญฺญา, อสุภสญฺญา มรณสญฺญา อาหาเรปฏิกูลสญฺญา สพฺพโลเก-อนภิรติสญฺญา\footnote{อนภิรตสญฺญา (สฺยา, ก)} อนิจฺจสญฺญา อนิจฺเจ ทุกฺขสญฺญา ทุกฺเข อนตฺตสญฺญา ปหานสญฺญา วิราคสญฺญา.\csromanspacer{} อิเม นว ธมฺมา \textbf{อุปฺปาเทตพฺพา}.}

\prose{(ฌ) กตเม นว ธมฺมา \textbf{อภิญฺเญยฺยา}.\csromanspacer{} นว อนุปุพฺพวิหารา.\csromanspacer{} อิธาวุโส ภิกฺขุ วิวิจฺเจว กาเมหิ วิวิจฺจ อกุสเลหิ ธมฺเมหิ สวิตกฺกํ สวิจารํ}

\vspace{\tipitakaparskip}
\csromancenter{ทสุตฺตรสุตฺต วิเวกชํ ปีติสุขํ ปฐมํ ฌานํ อุปสมฺปชฺช วิหรติ.\csromanspacer{} วิตกฺกวิจารานํ วูปสมา ฯเปฯ ทุติยํ ฌานํ อุปสมฺปชฺช วิหรติ.\csromanspacer{} ปีติยาจ วิราคา ฯเปฯ ตติยํ ฌานํ อุปสมฺปชฺช วิหรติ.\csromanspacer{} สุขสฺสจ ปหานา ฯเปฯ จตุตฺถํ ฌานํ อุปสมฺปชฺช วิหรติ.\csromanspacer{} สพฺพโส รูปสญฺญานํ สมติกฺกมา ฯเปฯ อากาสานญฺจายตนํ อุปสมฺปชฺช วิหรติ.\csromanspacer{} สพฺพโส อากาสานญฺจายตนํ สมติกฺกมฺม “อนนฺตํ วิญฺญาณนฺ”ติ วิญฺญาณญฺจายตนํ อุปสมฺปชฺช วิหรติ.\csromanspacer{} สพฺพโส วิญฺญาณญฺจายตนํ สมติกฺกมฺม “นตฺถิ กิญฺจี”ติ อากิญฺจญฺญายตนํ อุปสมฺปชฺช วิหรติ.\csromanspacer{} สพฺพโส อากิญฺจญฺญายตนํ สมติกฺกมฺม เนวสญฺญานาสญฺญายตนํ อุปสมฺปชฺช วิหรติ.\csromanspacer{} สพฺพโส เนวสญฺญานาสญฺญายตนํ สมติกฺกมฺม สญฺญาเวทยิตนิโรธํ อุปสมฺปชฺช วิหรติ.\csromanspacer{} อิเม นว ธมฺมา อภิญฺเญยฺยา.}
\csromanmatikaanchor{mtk.148}
\csromantocmarkref{section}{ทสธมฺมา}{mtk.148}
\bookmark[dest={mtk.148},level=1]{ทสธมฺมา}
\prose{\csromanfolio{255}(ญ) กตเม นว ธมฺมา \textbf{สจฺฉิกาตพฺพา}.\csromanspacer{} นว อนุปุพฺพนิโรธา.\csromanspacer{} ปฐมํ ฌานํ สมาปนฺนสฺส กามสญฺญา นิรุทฺธา โหติ, ทุติยํ ฌานํ สมาปนฺนสฺส วิตกฺกวิจารา นิรุทฺธา โหนฺติ, ตติยํ ฌานํ สมาปนฺนสฺส ปีติ นิรุทฺธา โหติ, จตุตฺถํ ฌานํ สมาปนฺนสฺส อสฺสาสปสฺสาสา นิรุทฺธา โหนฺติ, อากาสานญฺจายตนํ สมาปนฺนสฺส รูปสญฺญา นิรุทฺธา โหติ, วิญฺญาณญฺจายตนํ สมาปนฺนสฺส อากาสานญฺจายตนสญฺญา นิรุทฺธา โหติ, อากิญฺจญฺญายตนํ สมาปนฺนสฺส วิญฺญาณญฺจายตนสญฺญา นิรุทฺธา โหติ, เนวสญฺญานาสญฺญายตนํ สมาปนฺนสฺส อากิญฺจญฺญายตนสญฺญา นิรุทฺธา โหติ, สญฺญาเวทยิตนิโรธํ สมาปนฺนสฺส สญฺญา จ เวทนา จ นิรุทฺธา โหนฺติ.\csromanspacer{} อิเม นว ธมฺมา \textbf{สจฺฉิกาตพฺพา}.}

\prose{อิติ อิเม นวุติ ธมฺมา ภูตา ตจฺฉา ตถา อวิตถา อนญฺญถา สมฺมา ตถาคเตน อภิสมฺพุทฺธา.}

\csromanheader{2}{\textbf{ทส ธมฺมา}}
\vspace{\tipitakaparskip}
\csromancenter{\textbf{ทส ธมฺมา} \textbf{พหุการา} ฯเปฯ ทส ธมฺมา \textbf{สจฺฉิกาตพฺพา}.}
\prose{(ก) กตเม ทส ธมฺมา \textbf{พหุการา}.\csromanspacer{} ทส นาถกรณธมฺมา.\csromanspacer{} อิธาวุโส ภิกฺขุ สีลวา โหติ, ปาติโมกฺขสํวรสํวุโต วิหรติ อาจารโคจรสมฺปนฺโน, อณุมตฺเตสุ วชฺเชสุ ภยทสฺสาวี สมาทาย สิกฺขติ สิกฺขาปเทสุ, ยํปาวุโส ภิกฺขุ สีลวา โหติ ฯเปฯ สิกฺขติ สิกฺขาปเทสุ.\csromanspacer{} อยํปิ ธมฺโม นาถกรโณ. (1)}

\prose{\csromanfolio{256}ปุน จปรํ อาวุโส ภิกฺขุ พหุสฺสุโต ฯเปฯ ทิฏฺฐิยา สุปฺปฏิวิทฺธา, ยํปาวุโส ภิกฺขุ พหุสฺสุโต ฯเปฯ.\csromanspacer{} อยํปิ ธมฺโม นาถกรโณ. (2)}

\prose{ปุน จปรํ อาวุโส ภิกฺขุ กลฺยาณมิตฺโต โหติ กลฺยาณสหาโย กลฺยาณสมฺปวงฺโก.\csromanspacer{} ยํปาวุโส ภิกฺขุ ฯเปฯ กลฺยาณสมฺปวงฺโก.\csromanspacer{} อยํปิ ธมฺโม นาถกรโณ. (3)}

\prose{ปุน จปรํ อาวุโส ภิกฺขุ สุวโจ โหติ โสวจสฺสกรเณหิ ธมฺเมหิ สมนฺนาคโต, ขโม ปทกฺขิณคฺคาหี อนุสาสนิํ.\csromanspacer{} ยํปาวุโส ภิกฺขุ ฯเปฯ อนุสาสนิํ.\csromanspacer{} อยํปิ ธมฺโม นาถกรโณ. (4)}

\prose{ปุน จปรํ อาวุโส ภิกฺขุ ยานิ ตานิ สพฺรหฺมจารีนํ อุจฺจาวจานิ กิํกรณียานิ ตตฺถ ทกฺโข โหติ อนลโส ตตฺรุปายาย วีมํสาย สมนฺนาคโต, อลํ กาตุํ, อลํ สํวิธาตุํ.\csromanspacer{} ยํปาวุโส ภิกฺขุ ฯเปฯ อลํ สํวิธาตุํ.\csromanspacer{} อยํปิ ธมฺโม นาถกรโณ. (5)}

\prose{ปุน จปรํ อาวุโส ภิกฺขุ ธมฺมกาโม โหติ ปิยสมุทาหาโร อภิธมฺเม อภิวินเย อุฬารปาโมชฺโช.\csromanspacer{} ยํปาวุโส ภิกฺขุ ฯเปฯ อุฬารปาโมชฺโช.\csromanspacer{} อยํปิ ธมฺโม นาถกรโณ. (6)}

\prose{ปุน จปรํ อาวุโส ภิกฺขุ สนฺตุฏฺโฐ โหติ อิตรีตรจีวรปิณฺฑปาต เสนาสน คิลานปฺปจฺจย เภสชฺช ปริกฺขาเรหิ.\csromanspacer{} ยํปาวุโส ภิกฺขุ ฯเปฯ.\csromanspacer{} อยํปิ ธมฺโม นาถกรโณ. (7)}

\prose{ปุน จปรํ อาวุโส ภิกฺขุ อารทฺธวีริโย วิหรติ ฯเปฯ กุสเลสุ ธมฺเมสุ.\csromanspacer{} ยํปาวุโส ภิกฺขุ ฯเปฯ.\csromanspacer{} อยํปิ ธมฺโม นาถกรโณ. (8)}

\prose{ปุน จปรํ อาวุโส ภิกฺขุ สติมา โหติ ปรเมน สติเนปกฺเกน สมนฺนาคโต, จิรกตมฺปิ จิรภาสิตมฺปิ สริตา อนุสฺสริตา.\csromanspacer{} ยํปาวุโส ภิกฺขุ ฯเปฯ.\csromanspacer{} อยํปิ ธมฺโม นาถกรโณ. (9)}

\prose{ปุน จปรํ อาวุโส ภิกฺขุ ปญฺญวา โหติ อุทยตฺถคามินิยา ปญฺญาย สมนฺนาคโต, อริยาย นิพฺเพธิกาย สมฺมา ทุกฺขกฺขยคามินิยา.\csromanspacer{} ยํปาวุโส ภิกฺขุ ฯเปฯ.\csromanspacer{} อยํปิ ธมฺโม นาถกรโณ.\csromanspacer{} อิเม ทส ธมฺมา พหุการา. (10)}

\csromanheader{2}{11. ทสุตฺตรสุตฺต}
\prose{\csromanfolio{257}(ข) กตเม ทส ธมฺมา \textbf{ภาเวตพฺพา}.\csromanspacer{} ทส กสิณายตนานิ.\csromanspacer{} ปถวีกสิณเมโก สญฺชานาติ อุทฺธํ อโธ ติริยํ อทฺวยํ อปฺปมาณํ, อาโปกสิณเมโก สญฺชานาติ ฯเปฯ เตโชกสิณเมโก สญฺชานาติ.\csromanspacer{} วาโยกสิณเมโก สญฺชานาติ.\csromanspacer{} นีลกสิณเมโก สญฺชานาติ.\csromanspacer{} ปีตกสิณเมโก สญฺชานาติ.\csromanspacer{} โลหิตกสิณเมโก สญฺชานาติ.\csromanspacer{} โอทาตกสิณเมโก สญฺชานาติ.\csromanspacer{} อากาสกสิณเมโก สญฺชานาติ.\csromanspacer{} วิญฺญาณกสิณเมโก สญฺชานาติ อุทฺธํ อโธ ติริยํ อทฺวยํ อปฺปมาณํ.\csromanspacer{} อิเม ทส ธมฺมา \textbf{ภาเวตพฺพา}.}

\prose{(ค) กตเม ทส ธมฺมา \textbf{ปริญฺเญยฺยา}.\csromanspacer{} ทสายตนานิ, จกฺขายตนํ รูปายตนํ โสตายตนํ สทฺทายตนํ ฆานายตนํ คนฺธายตนํ ชิวฺหายตนํ รสายตนํ กายายตนํ โผฏฺฐพฺพายตนํ.\csromanspacer{} อิเม ทส ธมฺมา \textbf{ปริญฺเญยฺยา}.}

\prose{(ฆ) กตเม ทส ธมฺมา \textbf{ปหาตพฺพา}.\csromanspacer{} ทส มิจฺฉตฺตา, มิจฺฉาทิฏฺฐิ มิจฺฉาสงฺกปฺโป มิจฺฉาวาจา มิจฺฉากมฺมนฺโต มิจฺฉา-อาชีโว มิจฺฉาวายโม มิจฺฉาสติ มิจฺฉาสมาธิ มิจฺฉาญาณํ มิจฺฉาวิมุตฺติ.\csromanspacer{} อิเม ทส ธมฺมา \textbf{ปหาตพฺพา}.}

\prose{(\^{}อ) กตเม ทส ธมฺมา \textbf{หานภาคิยา}.\csromanspacer{} ทส อกุสลกมฺมปถา, ปาณติปาโต อทินฺนาทานํ กาเมสุมิจฺฉาจาโร มุสาวาโท ปิสุณา วาจา ผรุสา วาจา สมฺผปฺปลาโป อภิชฺฌา พฺยาปาโท มิจฺฉาทิฏฺฐิ.\csromanspacer{} อิเม ทส ธมฺมา \textbf{หานภาคิยา}.}

\prose{(จ) กตเม ทส ธมฺมา \textbf{วิเสสภาคิยา}.\csromanspacer{} ทส กุสลกมฺมปถา, ปาณติปาตา เวรมณี, อทินฺนาทานา เวรมณี, กาเมสุมิจฺฉาจารา เวรมณี, มุสาวาทา เวรมณี, ปิสุณาย วาจาย เวรมณี, ผรุสาย วาจาย เวรมณี, สมฺผปฺปลาปา เวรมณี, อนภิชฺฌา, อพฺยาปาโท, สมฺมาทิฏฺฐิ.\csromanspacer{} อิเม ทส ธมฺมา \textbf{วิเสสภาคิยา}.}

\prose{(ฉ) กตเม ทส ธมฺมา \textbf{ทุปฺปฏิวิชฺฌา}.\csromanspacer{} ทส อริยวาสา.\csromanspacer{} อิธาวุโส ภิกฺขุ ปญฺจงฺควิปฺปหีโน โหติ, ฉฬงฺคสมนฺนาคโต, เอการกฺโข, จตุราปสฺเสโน, ปณุนฺนปจฺเจกสจฺโจ, สมวยสฏฺเฐสโน, อนาวิลสงฺกปฺโป, ปสฺสทฺธกายสงฺขาโร, สุวิมุตฺตจิตฺโต, สุวิมุตฺตปญฺโญ.}

\prose{\csromanfolio{258}กถญฺจาวุโส ภิกฺขุ ปญฺจงฺควิปฺปหีโน โหติ.\csromanspacer{} อิธาวุโส ภิกฺขุโน กามจฺฉนฺโท ปหีโน โหติ, พฺยาปาโท ปหีโน โหติ, ถินมิทฺธํ ปหีนํ โหติ, อุทฺธจฺจกุกฺกุจฺจํ ปหีนํ โหติ, วิจิกิจฺฉา ปหีนา โหติ.\csromanspacer{} เอวํ โข อาวุโส ภิกฺขุ ปญฺจงฺควิปฺปหีโน โหติ. (1)}

\prose{กถญฺจาวุโส ภิกฺขุ ฉฬงฺคสมนฺนาคโต โหติ.\csromanspacer{} อิธาวุโส ภิกฺขุ จกฺขุนา รูปํ ทิสฺวา เนว สุมโน โหติ น ทุมฺมโน, อุเปกฺขโก วิหรติ สโต สมฺปชาโน.\csromanspacer{} โสเตน สทฺทํ สุตฺวา, ฆาเนน คนฺธํ ฆายิตฺวา, ชิวฺหาย รสํ สายิตฺวา, กาเยน โผฏฺฐพฺพํ ผุสิตฺวา, มนสา ธมฺมํ วิญฺญาย เนว สุมโน โหติ น ทุมฺมโน, อุเปกฺขโก วิหรติ สโต สมฺปชาโน.\csromanspacer{} เอวํ โข อาวุโส ภิกฺขุ ฉฬงฺคสมนฺนาคโต โหติ. (2)}

\prose{กถญฺจาวุโส ภิกฺขุ เอการกฺโข โหติ.\csromanspacer{} อิธาวุโส ภิกฺขุ สตารกฺเขน เจตสา สมนฺนาคโต โหติ.\csromanspacer{} เอวํ โข อาวุโส ภิกฺขุ เอการกฺโข โหติ. (3)}

\prose{กถญฺจาวุโส ภิกฺขุ จตุราปสฺเสโน โหติ.\csromanspacer{} อิธาวุโส ภิกฺขุ สงฺขาเยกํ ปฏิเสวติ, สงฺขาเยกํ อธิวาเสติ, สงฺขาเยกํ ปริวชฺเชติ, สงฺขาเยกํ วิโนเทติ.\csromanspacer{} เอวํ โข อาวุโส ภิกฺขุ จตุราปสฺเสโน โหติ. (4)}

\prose{กถญฺจาวุโส ภิกฺขุ ปณุนฺนปจฺเจกสจฺโจ โหติ.\csromanspacer{} อิธาวุโส ภิกฺขุโน ยานิ ตานิ ปุถุสมณพฺราหฺมณานํ ปุถุปจฺเจกสจฺจานิ, สพฺพานิ ตานิ นุนฺนานิ โหนฺติ ปณุนฺนานิ จตฺตานิ วนฺตานิ มุตฺตานิ ปหีนานิ ปฏินิสฺสฏฺฐานิ.\csromanspacer{} เอวํ โข อาวุโส ภิกฺขุ ปณุนฺนปจฺเจกสจฺโจ โหติ. (5)}

\prose{กถญฺจาวุโส ภิกฺขุ สมวยสฏฺเฐสโน โหติ.\csromanspacer{} อิธาวุโส ภิกฺขุโน กาเมสนา ปหีนา โหติ, ภเวสนา ปหีนา โหติ, พฺรหฺมจริเยสนา ปฏิปฺปสฺสทฺธา.\csromanspacer{} เอวํ โข อาวุโส ภิกฺขุ สมวยสฏฺเฐสโน โหติ. (6)}

\csromanheader{2}{11. ทสุตฺตรสุตฺต}
\prose{\csromanfolio{259}กถญฺจาวุโส ภิกฺขุ อนาวิลสงฺกปฺโป โหติ.\csromanspacer{} อิธาวุโส ภิกฺขุโน กามสงฺกปฺโป ปหีโน โหติ, พฺยาปาทสงฺกปฺโป ปหีโน โหติ, วิหิํสาสงฺกปฺโป ปหีโน โหติ.\csromanspacer{} เอวํ โข อาวุโส ภิกฺขุ อนาวิลสงฺกปฺโป โหติ. (7)}

\prose{กถญฺจาวุโส ภิกฺขุ ปสฺสทฺธกายสงฺขาโร โหติ.\csromanspacer{} อิธาวุโส ภิกฺขุ สุขสฺส จ ปหานา ทุกฺขสฺส จ ปหานา ปุพฺเพว โสมนสฺสโทมนสฺสานํ อตฺถงฺคมา อทุกฺขมสุขํ อุเปกฺขาสติปาริสุทฺธิํ จตุตฺถํ ฌานํ อุปสมฺปชฺช วิหรติ.\csromanspacer{} เอวํ โข อาวุโส ภิกฺขุ ปสฺสทฺธกายสงฺขาโร โหติ. (8)}

\prose{กถญฺจาวุโส ภิกฺขุ สุวิมุตฺตจิตฺโต โหติ.\csromanspacer{} อิธาวุโส ภิกฺขุโน ราคา จิตฺตํ วิมุตฺตํ โหติ, โทสา จิตฺตํ วิมุตฺตํ โหติ, โมหา จิตฺตํ วิมุตฺตํ โหติ.\csromanspacer{} เอวํ โข อาวุโส ภิกฺขุ สุวิมุตฺตจิตฺโต โหติ. (9)}

\prose{กถญฺจาวุโส ภิกฺขุ สุวิมุตฺตปญฺโญ โหติ.\csromanspacer{} อิธาวุโส ภิกฺขุ “ราโค เม ปหีโน อุจฺฉินฺนมูโล ตาลาวตฺถุกโต อนภาวํกโต อายติํ อนุปฺปาทธมฺโม”ติ ปชานาติ, “โทโส เม ปหีโน ฯเปฯ อายติํ อนุปฺปาทธมฺโม”ติ ปชานาติ. “โมโห เม ปหีโน ฯเปฯ อายติํ อนุปฺปาทธมฺโม”ติ ปชานาติ.\csromanspacer{} เอวํ โข อาวุโส ภิกฺขุ สุวิมุตฺตปญฺโญ โหติ.\csromanspacer{} อิเม ทส ธมฺมา ทุปฺปฏิวิชฺฌา. (10)}

\prose{(ช) กตเม ทส ธมฺมา \textbf{อุปฺปาเทตพฺพา}.\csromanspacer{} ทส สญฺญา, อสุภสญฺญา มรณสญฺญา อาหาเรปฏิกูลสญฺญา สพฺพโลเก-อนภิรติสญฺญา อนิจฺจสญฺญา อนิจฺเจ ทุกฺขสญฺญา ทุกฺเข อนตฺตสญฺญา ปหานสญฺญา วิราคสญฺญา นิโรธสญฺญา.\csromanspacer{} อิเม ทส ธมฺมา \textbf{อุปฺปาเทตพฺพา}.}

\prose{(ฌ) กตเม ทส ธมฺมา \textbf{อภิญฺเญยฺยา}.\csromanspacer{} ทส นิชฺชรวตฺถูนิ.\csromanspacer{} สมฺมาทิฏฺฐิสฺส มิจฺฉาทิฏฺฐิ นิชฺชิณฺณา โหติ.\csromanspacer{} เย จ มิจฺฉาทิฏฺฐิปจฺจยา อเนเก ปาปกา อกุสลา ธมฺมา สมฺภวนฺติ, เต จสฺส นิชฺชิณฺณา โหนฺติ.\csromanspacer{} สมฺมาสงฺกปฺปสฺส มิจฺฉาสงฺกปฺโป ฯเปฯ.\csromanspacer{} สมฺมาวาจสฺส มิจฺฉาวาจา ฯเปฯ.\csromanspacer{} สมฺมากมฺมนฺตสฺส มิจฺฉากมฺมนฺโต ฯเปฯ.\csromanspacer{} สมฺมา-อาชีวสฺส มิจฺฉา-อาชีโว ฯเปฯ.\csromanspacer{} สมฺมาวายามสฺส มิจฺฉาวายาโม ฯเปฯ.\csromanspacer{} สมฺมาสติสฺส มิจฺฉาสติ ฯเปฯ.\csromanspacer{} สมฺมาสมาธิสฺส มิจฺฉาสมาธิ ฯเปฯ.\csromanspacer{} สมฺมาญาณสฺส มิจฺฉาญาณํ นิชฺชิณฺณํ โหติ.\csromanspacer{} สมฺมาวิมุตฺติสฺส}

\csromanmatikaanchor{mtk.149}
\csromantocmarkref{section}{อุทฺทานคาถา}{mtk.149}
\bookmark[dest={mtk.149},level=1]{อุทฺทานคาถา}
\prose{\csromanfolio{260}มิจฺฉาวิมุตฺติ นิชฺชิณฺณา โหติ.\csromanspacer{} เย จ มิจฺฉาวิมุตฺติปจฺจยา อเนเก ปาปกา อกุสลา ธมฺมา สมฺภวนฺติ, เต จสฺส นิชฺชิณฺณา โหนฺติ.\csromanspacer{} อิเม ทส ธมฺมา อภิญฺเญยฺยา.}

\prose{(ญ) กตเม ทส ธมฺมา \textbf{สจฺฉิกาตพฺพา.}\csromanspacer{} ทส อเสกฺขา ธมฺมา, อเสกฺขา สมฺมาทิฏฺฐิ อเสกฺโข สมฺมาสงฺกปฺโป อเสกฺขา สมฺมาวาจา อเสกฺโข สมฺมากมฺมนฺโต อเสกฺโข สมฺมา-อาชีโว อเสกฺโข สมฺมาวายาโม อเสกฺขา สมฺมาสติ อเสกฺโข สมฺมาสมาธิ อเสกฺขํ สมฺมา ญาณํ อเสกฺขา สมฺมาวิมุตฺติ.\csromanspacer{} อิเม ทส ธมฺมา \textbf{สจฺฉิกาตพฺพา.}}

\prose{อิติ อิเม สตธมฺมา ภูตา ตจฺฉา ตถา อวิตถา อนญฺญถา สมฺมา ตถาคเตน อภิสมฺพุทฺธาติ.\csromanspacer{} อิทมโวจายสฺมา สาริปุตฺโต.\csromanspacer{} อตฺตมนา เต ภิกฺขู อายสฺมโต สาริปุตฺตสฺส ภาสิตํ อภินนฺทุนฺติ.}

\nitthitam{ทสุตฺตรสุตฺตํ นิฏฺฐิตํ เอกาทสมํ}
\nitthitam{ปาถิกวคฺโค\footnote{ปาฏิกวคฺโค (สี, สฺยา, อิ)} นิฏฺฐิโต.}
\titlehead{\textbf{ตสฺสุทฺทานํ}}
\csromangathameasure{ปาถิโก จ อุทุมฺพรํ, จกฺกวตฺติ อคฺคญฺญกํ. \\ สมฺปสาทนปาสาทํ, มหาปุริสลกฺขณํ. \\ สิงฺคาลาฏานาฏิยกํ, สงฺคีติ จ ทสุตฺตรํ. \\ เอกาทสหิ สุตฺเตหิ, ปาถิกวคฺโคติ วุจฺจติ.}
\csromangathagroup{\csromangathastanza{ปาถิโก จ\footnote{ปาฏิกญฺจ (สฺยา, กํ)} อุทุมฺพรํ\footnote{ปาฏิโกทุมฺพรีเจว (สี, อิ)}, จกฺกวตฺติ อคฺคญฺญกํ. \\ สมฺปสาทนปาสาทํ\footnote{ปาฏิกญฺจ (สฺยา, กํ)}, มหาปุริสลกฺขณํ.}\csromangathastanza{สิงฺคาลาฏานาฏิยกํ, สงฺคีติ จ ทสุตฺตรํ. \\ เอกาทสหิ สุตฺเตหิ, ปาถิกวคฺโคติ วุจฺจติ.}}

\nitthitam{\textbf{ปาถิกวคฺคปาฬิ นิฏฺฐิตา.}}
\csromanheader{2}{ตีหิ วคฺเคหิ ปฏิมณฺฑิโต สกโล}
\vspace{\tipitakaparskip}
\csromancenter{ทีฆนิกาโย สมตฺโต.}
\orphannote{เสยฺยถีทํ (สี, สฺยา, กํ, อิ)}

